\documentclass{situationist-preprint}

\usepackage{longtable,array}
\newenvironment{errlist}{\small\setlength{\LTpre}{2pt}\setlength{\LTpost}{4pt}%
\begin{longtable}{@{}>{\raggedleft}p{1.1cm}@{\hspace{0.8em}}>{\raggedright}p{3.6cm}@{\hspace{0.8em}}>{\raggedright\arraybackslash}p{\dimexpr\linewidth-4.7cm-1.6em\relax}@{}}}%
{\end{longtable}}
\newenvironment{skeltab}{\small\setlength{\LTpre}{2pt}\setlength{\LTpost}{8pt}%
\begin{longtable}{@{}r@{\hspace{0.8em}}p{\dimexpr\linewidth-3.6cm\relax}@{\hspace{0.8em}}r@{}}}%
{\end{longtable}}

\title{Corrigenda to \enquote{On the Existence of God}, Volumes III--XI}
\msid{REL-01-h}
\msfamily{REL-01}
\msdiscipline{Philosophy of Religion}
\mstype{Erratum}
\mskeywords{corrigenda; ontological argument; possibility premise; Lemma 2417; divine attributes; modal logic}

\begin{document}
\maketitle

\begin{abstract}
We list the corrections to Volumes III to XI of \emph{On the Existence of God} (REL-01), pages 791 to 4180. Eighteen corrections are substantive: they change the printed statement of a definition, lemma or proposition, a displayed premise, or a citation on which an argument relies. For each we give the printed text, the corrected text, and the results that use it, and we show that every such result stands. The most important concern Lemma~2417 of Volume~X, whose statement is unchanged but whose proof cites an idle principle and whose accompanying remark understates what the lemma shows; the definitions of omnipotence and omniscience in Volume~IV; the relation between the possibility premise and the conclusion in Volume~V; and a lemma on frames in Volume~X. The remaining corrections, about six thousand, are misprints and slips of reference, listed by volume, chapter and page. No correction affects the conclusion of the work, the argument of Volume~III for the possibility premise, or any result that Volume~I formalizes.
\end{abstract}

\section{Preface}

\subsection{Scope}

\emph{On the Existence of God} appeared in eleven volumes with continuous pagination. Volume~I (REL-01-a), pages 1 to 388, sets out the modal argument and its formalization; Volume~II, pages 389 to 790, treats dependence and the cosmological argument. The present list covers the other nine volumes: the defence of the possibility premise (Volume~III), the characterization of the maximally great being and the replies to objections (Volumes~IV to IX), the formal appendices (Volume~X) and the bibliography and indices (Volume~XI). Corrections to Volumes~I and~II are not included. Volume~I is in any case the volume least exposed to error of the kinds listed here, since its main argument is checked in the Lean module that accompanies it.

Section~\ref{s:volumes} gives the chapter structure of the nine volumes with page ranges, so that every entry can be located without the volumes at hand. Three conventions of the work bear on how entries are stated. First, chapters are numbered afresh in each volume, so that \enquote{Chapter~5} without further indication means the fifth chapter of the volume under which the entry is listed, and a result of another volume is cited as, for example, \enquote{Lemma~3.3 of Volume~I}. Second, in Volumes~III to IX definitions, lemmas, propositions, theorems, corollaries and remarks share one counter within each chapter (Definition~5.1, Lemma~5.2, and so on), and displayed formulas are numbered separately within each chapter, as (5.1), (5.2), and so on. Third, Volume~X numbers its lemmas consecutively through the whole volume, from Lemma~1 in Chapter~1 to Lemma~3361 in Chapter~8; the range of lemmas in each chapter is shown in the table of Section~\ref{s:volumes}. Chapter~5 of Volume~X, which collects the formal material for Volume~III, contains Lemmas 1802 to 2650 and so Lemma~2417.

\subsection{How the list was compiled}

The list has four sources. The first is a complete check of cross-references. Every reference in Volumes~III to X to a numbered item, of the same volume or of another, was compared with the item it names and with the Index of Numbered Results in Volume~XI; every page reference was compared with the page on which the item it names is printed. Where the reference named an existing item that was not the one the argument needs, the entry gives the item the argument uses. Where the index itself was wrong, the correction is listed under Volume~XI. The second source is a comparison of formulas. Each displayed formula in Volumes~III to IX that restates a definition or a result of Volume~I or of Volume~X was compared with its source, and each formula in Volume~X with the corresponding statement in the modules described in Chapter~8 of that volume. The third is a check of every entry of the bibliography against the published work, and of every citation in the text against the bibliography. The fourth is the reports of readers and our own rereading, which account for most of the purely typographical corrections and for several of the substantive ones.

The first three sources are systematic; the fourth is not. We therefore believe that the list is complete for cross-references, for formulas that restate results proved elsewhere in the work, and for the bibliography, and we do not claim completeness for typographical errors in running text.

\subsection{Policy}

Corrections are of two kinds, and the list treats them differently.

A correction is \emph{substantive} if it changes what a numbered statement says, what a displayed premise says, or which work an argument relies on, so that a reader who took the printed text at its word would be misled about what is claimed or about what supports it. Substantive corrections are given in full in Section~\ref{s:substantive}, numbered S1 to S18. Each gives the printed text, the corrected text, the reason for the change, and every place where the corrected item is used, with an account of why the use is unaffected. In five cases (S4, S5, S6, S10 and S17) a statement as printed is false or unproved and the correction adds a missing hypothesis or replaces a wrong one; in each case every use of the statement already satisfies the corrected hypothesis. In two cases (S8 and S13) a statement as printed asserts the opposite of what the surrounding argument uses, and the correction restores the statement the argument uses. The remaining eleven correct citations of results or of works in places where an argument depends on them. At the place of each substantive correction, the list by volume carries a pointer to it.

All other corrections are listed by volume in Sections~\ref{s:vol3} to \ref{s:vol11}. They are misprints of words, symbols and numbers, omitted or repeated words, wrong cross-references, wrong page references, and errors in the citation of works and in the indices. None of them changes what a numbered result states. Many occur inside proofs or inside the statement of a result; these are marked with a dagger (\dag). For every marked entry we have checked that the proof is valid when the correction is made, and that the corrected reading is the one the proof uses: a misprinted variable, a reversed accessibility relation in a step that the next line uses correctly, or a reference to a neighboring lemma in place of the one whose content the step applies. Where the corrected reading appears elsewhere in the same chapter, the entry says where, so that the reader can verify that the misprint is a misprint and not a second version of the claim.

No correction, substantive or not, affects the main conclusions of the work. The conclusion that a maximally great being exists rests on Theorem~3.4 of Volume~I and on the possibility premise, Assumption~3.1 of Volume~I. Theorem~3.4 is formalized, and nothing in Volumes~III to XI is a premise of it. The case for the possibility premise against its converse is made in Volume~III, Chapters~5 to~7. Its formal core is Lemma~2417 of Volume~X, whose statement is unchanged (S2); the correction to the remark that follows it shows that the lemma holds under a weaker hypothesis than it states, which can only strengthen the transfer from coherence to possibility that Volume~III makes, as REL-01-f argues. The corrections to Volume~III itself (S1, S3) concern citations in the proofs of results whose statements are unchanged. The characterization of the maximally great being in Volume~IV (S4 to S7) is unchanged in substance: the corrected definitions of omnipotence and omniscience are the ones on which its theorems were proved. The replies to objections in Volumes~V to VII, the treatment of uniqueness, simplicity and indiscernible godlike beings in Volume~VIII, and the conclusion of Volume~IX that the main result leaves the causal role of the maximally great being undetermined, all stand as printed. In particular, the claims of the work on which companion manuscripts rely are unaffected: the transfer of Volume~III's support to the ancestral premise of REL-01-b, the relation between godlikeness and maximal greatness treated in Volume~IV and discussed in REL-01-c, and the analysis of Lemma~2417 in REL-01-f.

\subsection{Form of the entries}

Each entry gives the page, the location on the page, and the correction, in the form \enquote{\textit{for} X \textit{read} Y}. Lines are counted from the top of the text block, not counting the running head; \enquote{l.~4 from foot} counts upward from the last line of text, not counting footnotes; \enquote{n.~7} is a footnote; a displayed formula is located by its number; an entry inside a proof names the result whose proof it is. Entries are ordered by page and, within a page, by position. Where one correction applies at several places (a running head, for example), the entry gives the range of pages. Entries in the bibliography of Volume~XI are located by the heading of the entry corrected, and entries in its indices by the heading under which the page number is listed (\enquote{s.v.}). Short notes after a semicolon or in parentheses say what the corrected reference names, where the correct form is printed, or which result a step uses; they are part of the correction, not commentary.

The list contains the following numbers of entries.

\begin{center}\small
% generated by generate.py -- do not edit
\begin{tabular}{@{}lr@{}}
\toprule
Typographical errors (spelling, running heads) & 1675 \\
Wrong word or phrase & 80 \\
Omitted or repeated words & 259 \\
Misprinted formulas and symbols & 875 \\
Cross-references and page references & 2562 \\
Citations in the text & 229 \\
Bibliography (Volume XI) & 72 \\
Indices (Volume XI) & 275 \\
\midrule
Total & 6027 \\
\bottomrule
\end{tabular}
\hspace{2em}
\begin{tabular}{@{}lr@{}}
\toprule
Volume III & 1065 \\
Volume IV & 624 \\
Volume V & 781 \\
Volume VI & 667 \\
Volume VII & 503 \\
Volume VIII & 502 \\
Volume IX & 566 \\
Volume X & 980 \\
Volume XI & 348 \\
\bottomrule
\end{tabular}

\end{center}

The proportions differ by volume as one would expect. Volume~X, which consists of formal material, has the largest share of misprinted formulas and of wrong references to neighboring lemmas, since its lemmas are short and numerous (eight to the page in places) and each proof cites several. Volumes~VI, VII and IX have proportionally more errors of citation and fewer of formulas. Volume~XI has few entries for its length: most of its bibliography is correct, and the indices were checked mechanically against the page numbers of the volumes.

A few entries illustrate the kinds of error the list contains. In Volume~III, p.~955, display (5.3) prints the rigidity of maximal greatness with the accessibility relation reversed; under the S5 conditions the two readings are equivalent, which is why the slip went unnoticed in proof, but the chapter does not use the S5 conditions, and the formula is used in its correct form on p.~958. In Volume~V, p.~1817, display (3.4) prints necessary maximal excellence with an existential quantifier where the universal is meant; the correct form is printed in (3.5) on the next page. In Volume~X, p.~3539, the Contents give the first page of Chapter~6 as 3843 for 3833. In Volume~XI, p.~4135, the Index of Subjects refers under \enquote{maximal excellence} to p.~1824 for p.~1724, which is in Volume~IV, Chapter~8, where godlikeness and maximal greatness are compared.

\section{The volumes and their chapters}\label{s:volumes}

The following table gives, for each of the nine volumes, its title, its page range and the title and pages of each of its chapters. For Volume~X it also gives the range of lemma numbers in each chapter. Every entry in Sections~\ref{s:vol3} to~\ref{s:vol11} falls within the page range of the chapter under which it is listed.

\medskip
% generated by generate.py -- do not edit
\noindent\textbf{Volume III: The Possibility Premise} \hfill pp.~791--1401\par\nopagebreak
\begin{skeltab}
1 & Introduction: the problem of the two premises & 795--810 \\
2 & Conceivability and possibility & 811--857 \\
3 & The epistemology of modal premises & 858--901 \\
4 & Positive and negative conceivability & 902--950 \\
5 & From coherence to possibility & 951--1013 \\
6 & The converse premise and the evidence of absence & 1014--1077 \\
7 & Symmetry-breaking: the asymmetry of support & 1078--1138 \\
8 & Parodies: islands, devils and the maximally evil being & 1139--1195 \\
9 & The ancestral premise and weaker modal logics & 1196--1250 \\
10 & Leibniz's problem and the compossibility of perfections & 1251--1301 \\
11 & Probabilistic assessments of the possibility premise & 1302--1370 \\
12 & Conclusion of Volume III & 1371--1401 \\
\end{skeltab}

\noindent\textbf{Volume IV: The Divine Attributes} \hfill pp.~1402--1757\par\nopagebreak
\begin{skeltab}
1 & Introduction & 1405--1411 \\
2 & Maximal excellence and its components & 1412--1446 \\
3 & Omnipotence & 1447--1495 \\
4 & Omniscience & 1496--1540 \\
5 & Moral perfection & 1541--1582 \\
6 & The compossibility of the attributes & 1583--1630 \\
7 & Necessary existence and aseity & 1631--1671 \\
8 & Godlikeness and maximal greatness & 1672--1728 \\
9 & The coherence of maximal excellence: summary & 1729--1757 \\
\end{skeltab}

\noindent\textbf{Volume V: Objections from Logic and Semantics} \hfill pp.~1758--2204\par\nopagebreak
\begin{skeltab}
1 & Introduction & 1761--1769 \\
2 & Existence as a predicate & 1770--1811 \\
3 & The question-begging objection & 1812--1874 \\
4 & Definitions and existential import & 1875--1912 \\
5 & Actualism, necessitism and the Barcan formula & 1913--1967 \\
6 & The S5 objection revisited & 1968--2028 \\
7 & Modal collapse & 2029--2080 \\
8 & Semantic objections: rigidity and reference to God & 2081--2126 \\
9 & Paraconsistent and intuitionistic variants & 2127--2188 \\
10 & Summary & 2189--2204 \\
\end{skeltab}

\noindent\textbf{Volume VI: Objections from Evil and Hiddenness} \hfill pp.~2205--2602\par\nopagebreak
\begin{skeltab}
1 & Introduction & 2208--2214 \\
2 & The logical problem of evil & 2215--2258 \\
3 & The structure of logical space and the free-will defence & 2259--2311 \\
4 & The evidential problem of evil & 2312--2360 \\
5 & The noseeum inference & 2361--2413 \\
6 & Divine hiddenness & 2414--2458 \\
7 & Hiddenness as a signalling equilibrium & 2459--2510 \\
8 & Evil, the possibility premise and the converse & 2511--2580 \\
9 & Summary & 2581--2602 \\
\end{skeltab}

\noindent\textbf{Volume VII: Objections from the Sciences} \hfill pp.~2603--2914\par\nopagebreak
\begin{skeltab}
1 & Introduction & 2606--2612 \\
2 & Naturalism and the modal argument & 2613--2650 \\
3 & Cosmology and the beginning of the universe & 2651--2698 \\
4 & Fine-tuning and the multiverse & 2699--2747 \\
5 & Evolutionary debunking of theistic belief & 2748--2796 \\
6 & The cognitive science of religion & 2797--2839 \\
7 & Probability, evidence and a priori premises & 2840--2896 \\
8 & Summary & 2897--2914 \\
\end{skeltab}

\noindent\textbf{Volume VIII: Uniqueness and Simplicity} \hfill pp.~2915--3203\par\nopagebreak
\begin{skeltab}
1 & Introduction & 2918--2923 \\
2 & The uniqueness of the maximally great being & 2924--2966 \\
3 & Indiscernible godlike beings & 2967--3010 \\
4 & Identity of indiscernibles in modal contexts & 3011--3057 \\
5 & Divine simplicity: the doctrine & 3058--3101 \\
6 & Simplicity and the attributes & 3102--3140 \\
7 & The objection from modal collapse to simplicity & 3141--3188 \\
8 & Summary & 3189--3203 \\
\end{skeltab}

\noindent\textbf{Volume IX: Divine Agency} \hfill pp.~3204--3537\par\nopagebreak
\begin{skeltab}
1 & Introduction & 3207--3213 \\
2 & What the main result says about action & 3214--3250 \\
3 & Models agreeing on maximal greatness & 3251--3297 \\
4 & Creation and conservation & 3298--3345 \\
5 & The freedom of the maximally great being & 3346--3392 \\
6 & Must a maximally great being create the best? & 3393--3439 \\
7 & Omniscience and foreknowledge & 3440--3487 \\
8 & Particular providence and the limits of the result & 3488--3520 \\
9 & Summary & 3521--3537 \\
\end{skeltab}

\noindent\textbf{Volume X: Formal Appendices} \hfill pp.~3538--3939\par\nopagebreak
\begin{skeltab}
1 & The modal language and its semantics \textit{(Lemmas 1--386)} & 3541--3581 \\
2 & Formal material for Volume I \textit{(Lemmas 387--790)} & 3582--3627 \\
3 & Higher-order modal logic and property abstraction \textit{(Lemmas 791--1302)} & 3628--3682 \\
4 & Formal material for Volume II \textit{(Lemmas 1303--1801)} & 3683--3728 \\
5 & Formal material for Volume III \textit{(Lemmas 1802--2650)} & 3729--3832 \\
6 & Formal material for Volumes IV and V \textit{(Lemmas 2651--3044)} & 3833--3874 \\
7 & Formal material for Volumes VI to IX \textit{(Lemmas 3045--3297)} & 3875--3912 \\
8 & The Lean modules \textit{(Lemmas 3298--3361)} & 3913--3939 \\
\end{skeltab}

\noindent\textbf{Volume XI: Bibliography and Indices} \hfill pp.~3940--4180\par\nopagebreak
\begin{skeltab}
 & Bibliography & 3943--4062 \\
 & Index of names & 4063--4101 \\
 & Index of subjects & 4102--4157 \\
 & Index of numbered results & 4158--4180 \\
\end{skeltab}


\section{Substantive corrigenda}\label{s:substantive}

\paragraph{S1. Volume~III, Chapter~5, p.~967, l.~4, and proof of Theorem~5.6, p.~983.} Chapter~5 argues that the support Volume~IV gives to the coherence of maximal excellence carries over to the possibility premise. The step from coherence to possibility is Lemma~2417 of Volume~X. It is cited twice in the chapter, on p.~967 and in the proof of Theorem~5.6, and both times printed as \enquote{Lemma~2471 of Volume~X}. \textit{For} Lemma~2471 \textit{read} Lemma~2417, in both places. Lemma~2471 is a lemma of the same chapter of Volume~X on another subject, and it does not deliver the step. Theorem~5.6 states that whatever supports the hypothesis of Lemma~2417 supports the possibility premise at least as strongly. Its proof has two steps: by Lemma~2417 the hypothesis entails the premise, and support is not lost under entailment. With the correction the first step cites the result it uses. Theorem~5.6 stands as printed. Since the hypothesis of Lemma~2417 can be weakened (S2), the theorem also holds with the weaker hypothesis in place of the stronger one, but we have not altered its statement.

\paragraph{S2. Volume~X, Chapter~5, p.~3803: Lemma~2417, its proof, and the following paragraph.} The lemma is printed as follows, and its statement is correct.

\begin{quote}\small
\textbf{Lemma 2417.} Let $\mathcal{P}$ be the family of great-making properties. If $\mathcal{P}$ is a proper filter, that is, if it is closed under conjunction and under entailment and does not contain the empty property, and if necessary existence and necessary maximal excellence belong to $\mathcal{P}$, then the possibility premise holds.
\end{quote}

Two corrections concern the text around it. First, the proof (pp.~3803--3804) opens with \enquote{By Lemma~3.3 of Volume~I}. Lemma~3.3 of Volume~I is the characteristic principle of S5, that whatever is possibly necessary is necessary. The proof does not use it: it uses closure under conjunction to place the conjunction of necessary existence and necessary maximal excellence in $\mathcal{P}$, properness to conclude that this conjunction is instantiated by some being at some world accessible from the actual world, and Definition~2.4 of Volume~I to identify that conjunction with maximal greatness. \textit{Delete} \enquote{By Lemma~3.3 of Volume~I}. The proof is otherwise correct, and as corrected it uses neither the possibility premise nor any of the S5 conditions, as REL-01-f observes.

Second, the paragraph that follows the proof (p.~3804, ll.~6--14) says that \enquote{none of the three filter conditions can be dropped}. This is not so. The proof uses properness only for the conjunction of two properties, and closure under entailment not at all; what it uses is that this one conjunction is instantiated. REL-01-f, Proposition~2.2, shows that the conclusion holds whenever $\mathcal{P}$ contains necessary existence and necessary maximal excellence and has the finite intersection property, and its Proposition~2.3 that every proper filter has that property but not conversely. The paragraph should read: \enquote{The hypothesis is stronger than the proof requires. It suffices that $\mathcal{P}$ have the finite intersection property, that is, that every finite subfamily of $\mathcal{P}$ be jointly instantiated; closure under conjunction and under entailment are not used.} The paragraph is not cited anywhere in the work. Its use in Volume~III, through Theorem~5.6 (S1), needs only that the hypothesis of the lemma entail the possibility premise, and a weaker hypothesis that entails it makes the transfer stronger, not weaker.

\paragraph{S3. Volume~III, Chapter~6, proof of Proposition~6.3, p.~1029.} Proposition~6.3 states that the possibility premise and the converse premise are not both true. The proof cites \enquote{Theorem~3.4 and Corollary~3.5 of Volume~I}. \textit{For} Corollary~3.5 \textit{read} Proposition~4.1. The proof runs as follows: if the possibility premise holds, then by the argument of Theorem~3.4 of Volume~I some being is maximally great at every world accessible from the actual world, and so, by reflexivity, at the actual world; if the converse premise holds, then by Proposition~4.1 of Volume~I no being is maximally great at the actual world. Corollary~3.5 of Volume~I, that some being exists and is maximally excellent at the actual world, does not contradict the converse premise by itself, and the proof does not use it. Proposition~6.3 stands as printed, and so does its use in Chapter~7, where the two premises are compared as exclusive alternatives.

\paragraph{S4. Volume~IV, Chapter~3, Definition~3.2, p.~1452.} The definition is printed:

\begin{quote}\small
\textbf{Definition 3.2.} A being $x$ is \emph{omnipotent} at a world $w$ if and only if $x$ can, at $w$, perform every task.
\end{quote}

It should read:

\begin{quote}\small
\textbf{Definition 3.2 (corrected).} A being $x$ is \emph{omnipotent} at a world $w$ if and only if $x$ can, at $w$, perform every task whose description does not quantify over $x$'s own powers.
\end{quote}

The restriction is stated and motivated on p.~1449, in the discussion that introduces the definition, and it is the restriction defended in REL-04. Without it the definition is unsatisfiable, since the task of making a stone its maker cannot lift is one that no being can perform if it can perform every task \citep{Mavrodes1963,Savage1967}. Definition~3.2 is used in three places. Theorem~3.9 (p.~1478), that the paradox of the stone imposes no limitation on an omnipotent being, would be false under the printed definition; its proof begins by observing that the description of the stone task quantifies over the maker's power, and so already uses the corrected definition. Lemma~6.4 (p.~1601), that omnipotence and moral perfection are compossible, applies the definition only to tasks described without reference to the agent's powers. Chapter~9 uses it only through Theorem~3.9 and Lemma~6.4. All three stand.

\paragraph{S5. Volume~IV, Chapter~4, Definition~4.1, p.~1499.} Omniscience is defined as knowledge of \enquote{every member of the set of propositions true at $w$}. \textit{Read}: \enquote{$x$ is omniscient at $w$ if and only if, for every proposition $p$ true at $w$, $x$ knows $p$ at $w$}. The printed definition presupposes that there is a set of all truths. Grim's Cantorian argument shows that there is none: for any set of truths, the truths about which of its subsets contain a given truth outnumber its members \citep{Grim1983,Grim1991}. The corrected definition quantifies over propositions without collecting them into a set. Theorem~4.6 (p.~1517) states that Grim's argument does not show omniscience to be incoherent, and its proof assumes the corrected form, since its point is that omniscience requires no set of all truths; under the printed definition the theorem would have no application. The other uses of the definition, in Chapters~6 and~9 and in Volume~IX, Chapter~7, quantify over propositions one at a time and are unaffected.

\paragraph{S6. Volume~IV, Chapter~8, Proposition~8.5, p.~1694.} The proposition is printed as \enquote{Every godlike being, in the sense of \citet{Anderson1990}, is maximally great.} \textit{Read}: \enquote{If necessary existence and each of omnipotence, omniscience and moral perfection are positive, then every godlike being, in the sense of \citet{Anderson1990}, is maximally great.} A godlike being in Anderson's sense has necessarily exactly the positive properties. Nothing in his axioms makes the components of maximal excellence positive, and without that hypothesis the proof (pp.~1694--1696), which moves from \enquote{$F$ is positive} to \enquote{$x$ has $F$ at every accessible world}, has nothing to apply to. The proof is correct under the hypothesis, which it states in its first line. Proposition~8.5 is used in Theorem~8.7 (p.~1703), which compares godlikeness and maximal greatness and which includes the positivity hypothesis in its statement. Theorem~8.7 stands. The corrected proposition is the connection that REL-01-c describes as belonging to Volume~IV: the two notions are related through the positivity of necessary existence and of the components of maximal excellence.

\paragraph{S7. Volume~IV, Chapter~6, p.~1589, l.~12.} The argument that an omnipotent being who cannot sin is not omnipotent is cited as \enquote{Pike (1965)}. \textit{For} Pike (1965) \textit{read} Pike (1969). The argument is in \citet{Pike1969}. \citet{Pike1965} is the paper on divine foreknowledge and voluntary action, discussed in Volume~IX, Chapter~7. Both are correctly listed in the bibliography of Volume~XI. The discussion on pp.~1589--1597 follows \citet{Pike1969} correctly and is unaffected.

\paragraph{S8. Volume~V, Chapter~3, Proposition~3.4, p.~1838.} Chapter~3 answers the objection that the modal argument begs the question. Proposition~3.4 is printed:

\begin{quote}\small
\textbf{Proposition 3.4.} The possibility premise is not S5-equivalent to the conclusion of Theorem~3.4 of Volume~I.
\end{quote}

This is false. Under the S5 conditions the possibility premise entails the conclusion (Theorem~3.4 of Volume~I), and the conclusion entails the premise by reflexivity; the two are equivalent, as the chapter itself states on p.~1841. The proposition should read:

\begin{quote}\small
\textbf{Proposition 3.4 (corrected).} The possibility premise is S5-equivalent to the conclusion of Theorem~3.4 of Volume~I; but the grounds offered for it in Volume~III do not presuppose that conclusion.
\end{quote}

The reply to the objection in Chapter~3 uses only the second clause: an argument whose premise is equivalent to its conclusion in the background logic begs the question only if the premise can be supported in no way that presupposes the conclusion, and Volume~III supports it through the coherence of maximal excellence and Lemma~2417. The proof of Proposition~3.4 (pp.~1838--1840) establishes both clauses of the corrected statement; its first paragraph proves the equivalence that the printed statement denies. Nothing else in the volume depends on the printed form.

\paragraph{S9. Volume~V, Chapter~2, p.~1783, l.~20.} Kant's statement that being is not a real predicate is cited at \enquote{A599/B627}. \textit{Read} A598/B626 \citep{Kant}. The surrounding discussion is unaffected.

\paragraph{S10. Volume~V, Chapter~5, Lemma~5.3, p.~1926.} The lemma is printed: \enquote{On every frame with increasing domains (if $wRv$ then $D_w \subseteq D_v$), the Barcan formula $\forall x\,\nec Fx \to \nec\forall x\,Fx$ is valid.} \textit{For} increasing domains (if $wRv$ then $D_w \subseteq D_v$) \textit{read} decreasing domains (if $wRv$ then $D_v \subseteq D_w$). Increasing domains validate the converse Barcan formula, not the Barcan formula; the Barcan formula is valid on a frame exactly when domains do not grow along accessibility \citep{HughesCresswell1996}. The proof on p.~1927 in fact assumes decreasing domains: it takes an object of $D_v$ and uses that it belongs to $D_w$. The lemma is used in Theorem~5.5 (p.~1934), that the constant-domain framework of Volume~I validates both the Barcan formula and its converse. A constant domain is both increasing and decreasing, so the theorem follows from the corrected lemma together with its counterpart for the converse formula, which is stated correctly as Lemma~5.4. The formal version in Volume~X, Chapter~6, is stated correctly.

\paragraph{S11. Volume~VI, Chapter~4, display (4.2), p.~2318.} Display (4.2) reproduces the second premise of the evidential argument of \citet{Rowe1979} without its qualifying clause, as \enquote{an omniscient, wholly good being would prevent the occurrence of any intense suffering it could}. \textit{Read}: an omniscient, wholly good being would prevent the occurrence of any intense suffering it could, unless it could not do so without thereby losing some greater good or permitting some evil equally bad or worse. Without the clause the premise is one that no theist accepts and the argument would be of no interest. Chapter~4 discusses the argument with the clause throughout, and Chapter~5's reply, that the first premise rests on an unwarranted noseeum inference \citep{Wykstra1984}, presupposes the qualified second premise. Both chapters are unaffected.

\paragraph{S12. Volume~VI, Chapter~7, Lemma~7.5, p.~2483.} Chapter~7 models divine hiddenness as a signalling game between God and free creatures. Lemma~7.5 states that the strategy of non-revelation is an equilibrium of the game. It omits a hypothesis: that the creatures' payoff from a freely chosen response exceeds their payoff from a response that unambiguous revelation would make assured. \textit{After} \enquote{is an equilibrium} \textit{insert} \enquote{provided that the value to the creatures of a free response exceeds the value of an assured one}. Without the hypothesis, revelation weakly dominates and the lemma fails. The proof (pp.~2483--2485) uses the hypothesis in its second step, where it compares the two payoffs, and the condition is stated on p.~2477, where the payoffs are introduced. Theorem~7.7 (p.~2490), that hiddenness is the equilibrium a perfectly loving God would choose, includes the condition in its statement. It stands, as does its agreement with REL-03.

\paragraph{S13. Volume~VII, Chapter~5, display (5.6), p.~2771.} Chapter~5 asks whether the existence of an evolutionary explanation of theistic belief disconfirms theism. Display (5.6) states the condition for evidence $E$ to disconfirm theism $T$ as $P(E \mid T)/P(E \mid \neg T) > 1$. \textit{For} $>$ \textit{read} $<$. Evidence disconfirms a hypothesis when it is less probable on the hypothesis than on its negation. The text on p.~2772 applies the correct condition, and argues that the evolutionary explanation is about equally probable on theism and on its negation, so that the ratio is close to one; it argues further that such evidence bears on how belief is caused and not on the a priori possibility premise. The conclusion of the chapter is unaffected.

\paragraph{S14. Volume~VIII, Chapter~3, Lemma~3.6, p.~2989.} Chapter~3 treats the question whether indiscernible godlike beings are identical. Lemma~3.6 is printed: \enquote{If $x$ and $y$ have the same great-making properties, then $x = y$.} \textit{After} \enquote{the same great-making properties} \textit{insert} \enquote{at every world accessible from the actual world}. The proof (pp.~2989--2991) compares $x$ and $y$ at each accessible world in turn and uses the hypothesis at every one of them; sharing the great-making properties at a single world is not enough for it. The lemma is used in Theorem~3.8 (p.~2996), that any two maximally great beings are identical, where it is applied to beings that exist and have every great-making property at every accessible world, by Definition~2.4 of Volume~I. The corrected hypothesis is therefore satisfied, and Theorem~3.8 stands.

\paragraph{S15. Volume~VIII, Chapter~5, p.~3066, l.~25.} The article in which Aquinas asks whether God is altogether simple is cited as \enquote{\emph{Summa Theologiae} I, q.~3, a.~8}. \textit{Read} I, q.~3, a.~7 \citep{Aquinas}. Article~8 of the same question asks whether God enters into composition with other things, which is discussed separately on p.~3071, where it is cited correctly.

\paragraph{S16. Volume~IX, Chapter~3, proof of Theorem~3.5, p.~3280.} Theorem~3.5 states that there are two models of the framework of Volume~I which agree on which beings are maximally great at every world but differ on whether the maximally great being causes a given contingent event. The proof constructs the second model, $\mathcal{M}_2$, from the first, $\mathcal{M}_1$, by removing the causal relation between the maximally great being and the event, and prints the accessibility relation of $\mathcal{M}_2$ as $R_2 = \{(w,v) : w \neq v\}$. This relation is not reflexive, so $\mathcal{M}_2$ does not satisfy Assumption~2.1 of Volume~I. \textit{For} $R_2 = \{(w,v) : w \neq v\}$ \textit{read} $R_2 = W \times W$, which is the accessibility relation of $\mathcal{M}_1$. With the correction, $\mathcal{M}_2$ satisfies the S5 conditions; since the two models have the same worlds, the same accessibility relation, the same existence predicate and the same maximal excellence predicate, they agree on maximal greatness at every world, and they differ on the causal relation by construction. Theorem~3.5 stands, and with it the conclusion of Volume~IX, stated in REL-01-a and in REL-01-e, that the main result leaves the causal role of the maximally great being undetermined.

\paragraph{S17. Volume~X, Chapter~1, Lemma~214, p.~3563.} The lemma is printed: \enquote{A relation is an equivalence relation if and only if it is serial and Euclidean.} \textit{For} serial \textit{read} reflexive. A serial and Euclidean relation need not be reflexive: on the worlds $a$ and $b$, the relation $\{(a,b), (b,b)\}$ is serial and Euclidean, and $a$ does not see itself. Seriality and the Euclidean property characterize the frames of KD45, not of S5 \citep{Chellas1980,HughesCresswell1996}. The corrected lemma, that a relation is an equivalence relation if and only if it is reflexive and Euclidean, is the standard one. It is used in Chapter~2, where the frames of Volume~I are shown to be S5 frames from Assumption~2.1 of Volume~I, which includes reflexivity explicitly. That use is unaffected, and no other lemma of the volume uses Lemma~214.

\paragraph{S18. Volume~XI, Bibliography, p.~3976, entry Gödel (1995).} The entry places Gödel's ontological proof in \enquote{\emph{Collected Works, Volume II}}. \textit{Read} \emph{Collected Works, Volume III: Unpublished Essays and Lectures}, pp.~403--404 \citep{Godel1995}. Volume~II of the \emph{Collected Works} contains Gödel's publications from 1938 to 1974, and the ontological proof was not published in his lifetime. The year of the entry, 1995, is correct, and every citation of it in Volumes~III to X is correct.

\section{Volume III: The Possibility Premise}\label{s:vol3}
% generated by generate.py -- do not edit
\subsection*{Front matter \hfill \normalfont\small pp.~791--794}
\begin{errlist}
792 & Contents, Chapter~5 & \textit{for} 950 \textit{read} 951 \\
\end{errlist}

\subsection*{Chapter~1. Introduction: the problem of the two premises \hfill \normalfont\small pp.~795--810}
\begin{errlist}
796 & l.~21 & \textit{for} sequecne \textit{read} sequence \\
796 & l.~30 & \textit{delete} the second occurrence of \emph{of} \\
797 & l.~10 & \textit{for} van Inwagen (1976) \textit{read} van Inwagen (1977) \\
797 & l.~6 from foot & \textit{for} atlernative \textit{read} alternative \\
797 & l.~5 from foot & \textit{for} Definition~2.2 of Volume~I (maximal greatness) \textit{read} Definition~2.4 of Volume~I (maximal greatness) \\
798 & l.~33 & \textit{for} extensoin \textit{read} extension \\
799 & l.~18 & \textit{for} organiaztion \textit{read} organization \\
800 & l.~4 & \textit{for} naturaly \textit{read} naturally \\
800 & l.~23 & \textit{for} Lemma~2202 of Volume~X \textit{read} Lemma~2201 of Volume~X \\
800 & l.~24 & \textit{for} occrured \textit{read} occurred \\
800 & l.~35 & \textit{for} precisly \textit{read} precisely \\
800 & n.~2 & \textit{for} Hartshorne (1961) \textit{read} Hartshorne (1962) \\
801 & l.~7 & \textit{for} Volume~III, Chapter~4 \textit{read} Volume~IV, Chapter~4 (omniscience) \\
801 & l.~6 from foot & \textit{for} somehting \textit{read} something \\
801 & l.~2 from foot & \textit{for} thoroguh \textit{read} thorough \\
802 & l.~9 & \textit{for} oriignal \textit{read} original \\
803 & l.~13 & \textit{for} Volume~V, Chapter~4 \textit{read} Volume~IV, Chapter~4 (omniscience) \\
803 & l.~18 & \textit{delete} the second occurrence of \emph{that} \\
803 & l.~19 & \textit{for} certianly \textit{read} certainly \\
804 & l.~7 & \textit{for} immediattely \textit{read} immediately \\
805 & l.~33 & \textit{for} recieve \textit{read} receive \\
805 & l.~37 & \textit{for} implcit \textit{read} implicit \\
807 & l.~6 & \textit{for} demmonstrate \textit{read} demonstrate \\
807 & l.~26 & \textit{for} Volume~X, Chapter~7 \textit{read} Volume~X, Chapter~5 (the formal material for Volume III, including Lemma~2417) \\
807 & l.~37 & \textit{for} rfeerence \textit{read} reference \\
808 & l.~27 & \textit{for} reefrence \textit{read} reference \\
808 & l.~8 from foot & \textit{for} Plantinga (1973) \textit{read} Plantinga (1974b) \\
810 & l.~2 & \textit{for} exstence \textit{read} existence \\
\end{errlist}

\subsection*{Chapter~2. Conceivability and possibility \hfill \normalfont\small pp.~811--857}
\begin{errlist}
812 & l.~28 & \textit{for} beileve \textit{read} believe \\
813 & l.~12 & \textit{for} asssesment \textit{read} assessment \\
813 & l.~26 & \textit{for} $P(\nec\neg G \mid A) = P(\pos\neg G)$ \textit{read} $P(\pos\neg G \mid A) = P(\pos\neg G)$ \\
814 & l.~3 from foot & \textit{for} empiriacl \textit{read} empirical \\
815 & l.~14 & \textit{for} Lemma~2524 of Volume~X \textit{read} Lemma~2518 of Volume~X; the result cited is on p.~3816 \\
815 & display (2.2) & \textit{for} $\neg\mathrm{Conc}( p) \wedge \mathrm{Conc}(p)$ \textit{read} $\neg\mathrm{Conc}(\neg p) \wedge \mathrm{Conc}(p)$ \\
815 & l.~27 & \textit{for} explanasion \textit{read} explanation \\
815 & l.~30 & \textit{for} Volume~III, Chapter~6 \textit{read} Volume~IV, Chapter~6 (the compossibility of the attributes) \\
815 & l.~31 & \textit{for} postion \textit{read} position \\
815 & l.~6 from foot & \textit{for} Willaimson (2013) \textit{read} Williamson (2013) \\
816 & l.~8 & \textit{for} crittical \textit{read} critical \\
816 & l.~10 & \textit{delete} the second occurrence of \emph{that} \\
817 & l.~15 & \textit{for} instaed \textit{read} instead \\
818 & l.~12 & \textit{for} methd \textit{read} method \\
818 & l.~2 from foot & \textit{for} seciton \textit{read} section \\
819 & l.~2 & \textit{for} pirncipal \textit{read} principal \\
\dag\,819 & proof of Corollary~2.4, l.~26 & \textit{for} Definition~2.5 \textit{read} Definition~2.1 \\
821 & l.~5 & \textit{for} challnege \textit{read} challenge \\
821 & l.~8 & \textit{for} intuiitve \textit{read} intuitive \\
822 & l.~26 & \textit{for} bears on the possibility premise \textit{read} bears on the converse premise \\
824 & l.~9 & \textit{for} quatnifier \textit{read} quantifier \\
\dag\,824 & proof of Proposition~2.8, l.~11 & \textit{for} Theorem~2.11 \textit{read} Theorem~2.7 \\
\dag\,824 & proof of Proposition~2.8, (2.4) & \textit{for} $\mathrm{Conc}(r) \to \nec r$ \textit{read} $\mathrm{Conc}(r) \to \pos r$; the formula is used in its correct form on p.~826 \\
825 & l.~3 & \textit{for} Volume~X, Chapter~6 \textit{read} Volume~X, Chapter~8 (the Lean modules) \\
825 & l.~10 & \textit{for} tradtiion \textit{read} tradition \\
\dag\,825 & proof of Proposition~2.8, l.~16 & \textit{for} Lemma~2604 of Volume~X \textit{read} Lemma~2605 of Volume~X \\
826 & l.~2 from foot & \textit{delete} the second occurrence of \emph{to} \\
827 & l.~18 & \textit{for} simpliicty \textit{read} simplicity \\
828 & l.~36 & \textit{for} sgugestion \textit{read} suggestion \\
829 & l.~9 & \textit{for} Volume~VII, Chapter~3 \textit{read} Volume~VII, Chapter~5 (evolutionary debunking) \\
830 & display (2.5) & \textit{for} $P(\nec G) > P(\pos\neg G)$ \textit{read} $P(\pos G) > P(\pos\neg G)$; the formula is used in its correct form on p.~834 \\
830 & l.~16 & \textit{for} Lemma~1891 of Volume~X \textit{read} Lemma~1887 of Volume~X \\
831 & l.~18 & \textit{for} Williamson (2012) \textit{read} Williamson (2013) \\
834 & l.~32 & \textit{for} difefrence \textit{read} difference \\
834 & l.~3 from foot & \textit{for} $P(q \mid r) = P(r \mid q)\,P(r)/P(r)$ \textit{read} $P(q \mid r) = P(r \mid q)\,P(q)/P(r)$ \\
835 & l.~7 & \textit{for} princpial \textit{read} principal \\
835 & l.~8 & \textit{for} beneficail \textit{read} beneficial \\
836 & l.~30 & \textit{for} p.~831 \textit{read} p.~832 (Definition~2.10) \\
837 & l.~4 & \textit{after} conceivable \textit{insert} on ideal reflection \\
837 & l.~8 & \textit{for} simlpicity \textit{read} simplicity \\
838 & l.~4 & \textit{delete} the second occurrence of \emph{a} \\
838 & l.~35 & \textit{for} Lemma~2382 of Volume~X \textit{read} Lemma~2388 of Volume~X \\
840 & l.~12 & \textit{for} neverthelses \textit{read} nevertheless \\
840 & l.~3 from foot & \textit{for} ultimtely \textit{read} ultimately \\
841 & l.~5 & \textit{for} wheraes \textit{read} whereas \\
841 & l.~11 & \textit{for} arguemnt \textit{read} argument \\
841 & l.~26 & \textit{for} sqeuence \textit{read} sequence \\
841 & l.~31 & \textit{for} arugably \textit{read} arguably \\
\dag\,842 & proof of Theorem~2.11, l.~4 & \textit{for} Definition~2.5 \textit{read} Definition~2.10 \\
843 & l.~3 & \textit{for} therrefore \textit{read} therefore \\
\dag\,843 & proof of Theorem~2.11, l.~11 & \textit{for} Definition~2.9 \textit{read} Definition~2.5 \\
844 & l.~12 & \textit{for} understadning \textit{read} understanding \\
844 & l.~14 & \textit{for} ordinry \textit{read} ordinary \\
844 & display (2.9) & \textit{for} $P(\pos G) + P(\pos\neg G) \geq 1$ \textit{read} $P(\pos G) + P(\pos\neg G) \leq 1$ \\
844 & l.~16 & \textit{for} theorrem \textit{read} theorem \\
844 & display (2.10) & \textit{for} $P(G) = P(\nec G)$ \textit{read} $P(G) = P(\pos G)$ \\
844 & l.~4 from foot & \textit{for} Lemma~1978 of Volume~X \textit{read} Lemma~1976 of Volume~X \\
845 & l.~6 & \textit{for} intreest \textit{read} interest \\
845 & l.~23 & \textit{for} anaolgous \textit{read} analogous \\
845 & l.~27 & \textit{for} trdaition \textit{read} tradition \\
846 & l.~7 & \textit{for} smiplicity \textit{read} simplicity \\
\dag\,847 & proof of Proposition~2.14, l.~15 & \textit{for} $P(\pos\neg G \mid A) = P(\pos G)$ \textit{read} $P(\pos\neg G \mid A) = P(\pos\neg G)$ \\
\dag\,847 & proof of Proposition~2.14, l.~28 & \textit{for} Proposition~2.8 \textit{read} Proposition~2.3 \\
848 & l.~2 from foot & \textit{for} Lemma~1952 of Volume~X \textit{read} Lemma~1954 of Volume~X \\
849 & l.~7 & \textit{for} negatively conceivable \textit{read} positively conceivable \\
849 & l.~10 & \textit{for} p.~845 \textit{read} p.~844 (Definition~2.12) \\
849 & l.~17 & \textit{for} availlable \textit{read} available \\
850 & l.~6 from foot & \textit{for} the converse premise is better supported \textit{read} the possibility premise is better supported \\
851 & l.~2 & \textit{for} Kripke (1979) \textit{read} Kripke (1980) \\
851 & n.~15 & \textit{for} Krikpe (1980) \textit{read} Kripke (1980) \\
852 & l.~4 & \textit{for} diferent \textit{read} different \\
853 & l.~31 & \textit{for} appraoch \textit{read} approach \\
854 & l.~5 & \textit{for} cirucmstance \textit{read} circumstance \\
854 & l.~7 & \textit{for} p.~846 \textit{read} p.~847 (Proposition~2.14) \\
854 & l.~11 & \textit{for} $\neg\mathrm{Conc}( r) \wedge \mathrm{Conc}(r)$ \textit{read} $\neg\mathrm{Conc}(\neg r) \wedge \mathrm{Conc}(r)$ \\
854 & l.~14 & \textit{for} Descartes, Meditation~III (the supremely perfect being whose essence includes existence) \textit{read} Descartes, Meditation~V (the supremely perfect being whose essence includes existence) \\
855 & l.~28 & \textit{delete} the second occurrence of \emph{a} \\
855 & l.~2 from foot & \textit{for} miantain \textit{read} maintain \\
856 & l.~19 & \textit{for} positive conceivability \textit{read} negative conceivability \\
\dag\,856 & proof of Lemma~2.16, l.~4 from foot & \textit{for} $P(\pos G) > P(\pos G)$ \textit{read} $P(\pos G) > P(\pos\neg G)$ \\
857 & l.~2 & \textit{for} desriable \textit{read} desirable \\
\dag\,857 & proof of Lemma~2.16, l.~5 & \textit{for} Lemma~2197 of Volume~X \textit{read} Lemma~2202 of Volume~X \\
\end{errlist}

\subsection*{Chapter~3. The epistemology of modal premises \hfill \normalfont\small pp.~858--901}
\begin{errlist}
859 & display (3.1) & \textit{for} $w_1 \models \nec q \iff \exists w'\,(w_1Rw' \wedge w' \models q)$ \textit{read} $w_1 \models \pos q \iff \exists w'\,(w_1Rw' \wedge w' \models q)$; the formula is used in its correct form on p.~862 \\
859 & l.~16 & \textit{for} apparant \textit{read} apparent \\
860 & l.~21 & \textit{for} dilmema \textit{read} dilemma \\
860 & l.~40 & \textit{for} descirption \textit{read} description \\
861 & l.~7 & \textit{for} preicsely \textit{read} precisely \\
861 & l.~19 & \textit{for} reasonalbe \textit{read} reasonable \\
861 & l.~23 & \textit{for} van Inwagen (1978) \textit{read} van Inwagen (1977) \\
862 & l.~31 & \textit{for} preciesly \textit{read} precisely \\
862 & display (3.2) & \textit{for} $v' \models \pos r \iff \exists w_2\,(w_2Rv' \wedge w_2 \models r)$ \textit{read} $v' \models \pos r \iff \exists w_2\,(v'Rw_2 \wedge w_2 \models r)$ \\
863 & l.~22 & \textit{for} Volume~V, Chapter~8 \textit{read} Volume~IV, Chapter~8 (godlikeness and maximal greatness) \\
864 & l.~15 & \textit{for} nottaion \textit{read} notation \\
865 & l.~5 & \textit{for} croresponding \textit{read} corresponding \\
865 & l.~6 & \textit{for} agrued \textit{read} argued \\
866 & l.~1 & \textit{for} Volume~X, Chapter~3 \textit{read} Volume~X, Chapter~5 (the formal material for Volume III, including Lemma~2417) \\
866 & l.~17 & \textit{for} the axiom T, $p \to \nec\pos p$ \textit{read} the axiom B, $p \to \nec\pos p$ \\
866 & l.~34 & \textit{for} judgemment \textit{read} judgement \\
867 & l.~5 & \textit{for} resisatnce \textit{read} resistance \\
869 & l.~31 & \textit{for} $\nec(p \to r) \to (\nec p \to \pos r)$ \textit{read} $\nec(p \to r) \to (\nec p \to \nec r)$ \\
869 & l.~8 from foot & \textit{for} $\nec(p \leftrightarrow q) \to (\nec p \to \nec q)$ \textit{read} $\nec(p \to q) \to (\nec p \to \nec q)$ \\
870 & l.~4 & \textit{for} Lemma~2216 of Volume~X \textit{read} Lemma~2210 of Volume~X; the result cited is on p.~3779 \\
870 & l.~10 & \textit{for} Lemma~2405 of Volume~X \textit{read} Lemma~2399 of Volume~X \\
870 & display (3.11) & \textit{for} $\nec p \leftrightarrow \nec\nec p$ \textit{read} $\nec p \to \nec\nec p$; correctly printed in (3.14), p.~876 \\
871 & display (3.12) & \textit{for} $r \to \nec\nec r$ \textit{read} $r \to \nec\pos r$; the formula is used in its correct form on p.~872 \\
871 & l.~21 & \textit{for} Williamson (2013) \textit{read} van Inwagen (1977) (the symmetry of the modal ontological argument) \\
871 & l.~7 from foot & \textit{for} idnependent \textit{read} independent \\
874 & l.~4 & \textit{for} circumtsance \textit{read} circumstance \\
874 & l.~6 from foot & \textit{delete} the second occurrence of \emph{that} \\
875 & l.~11 & \textit{for} slighlty \textit{read} slightly \\
875 & l.~22 & \textit{for} contnt \textit{read} content \\
876 & l.~23 & \textit{for} orgiinal \textit{read} original \\
877 & l.~11 & \textit{delete} the second occurrence of \emph{is} \\
878 & l.~26 & \textit{for} p.~823 \textit{read} p.~822 (Lemma~2.6) \\
879 & l.~2 & \textit{for} satisifed \textit{read} satisfied \\
879 & l.~27 & \textit{for} disupte \textit{read} dispute \\
881 & l.~7 & \textit{for} Swinburne (2005) \textit{read} Swinburne (2004) \\
881 & l.~22 & \textit{for} presmuably \textit{read} presumably \\
881 & l.~27 & \textit{for} ambigouus \textit{read} ambiguous \\
882 & l.~15 & \textit{for} judgmeent \textit{read} judgement \\
\dag\,882 & proof of Lemma~3.6, l.~20 & \textit{for} Theorem~3.7 \textit{read} Theorem~3.5 \\
882 & l.~22 & \textit{for} posssesion \textit{read} possession \\
882 & l.~28 & \textit{for} capactiy \textit{read} capacity \\
883 & l.~2 & \textit{for} Swinburne (2003) \textit{read} Swinburne (2004) \\
\dag\,883 & proof of Theorem~3.7, l.~9 & \textit{for} Lemma~3.6 \textit{read} Lemma~3.2 \\
884 & l.~2 & \textit{for} mroeover \textit{read} moreover \\
884 & l.~5 & \textit{after} accessible \textit{insert} from $w'$ \\
884 & l.~21 & \textit{for} p.~882 \textit{read} p.~883 (Theorem~3.7) \\
885 & l.~5 & \textit{for} moroever \textit{read} moreover \\
886 & l.~6 & \textit{for} remarkble \textit{read} remarkable \\
886 & l.~18 & \textit{for} nautrally \textit{read} naturally \\
888 & l.~13 & \textit{for} immediaetly \textit{read} immediately \\
889 & l.~3 & \textit{for} evidnt \textit{read} evident \\
889 & l.~4 & \textit{for} hypothessis \textit{read} hypothesis \\
889 & l.~25 & \textit{for} commtiment \textit{read} commitment \\
889 & n.~21 & \textit{for} Swinbrune (2004) \textit{read} Swinburne (2004) \\
890 & l.~13 & \textit{for} ambguity \textit{read} ambiguity \\
890 & l.~28 & \textit{for} Lemma~2402 of Volume~X \textit{read} Lemma~2401 of Volume~X; the result cited is on p.~3801 \\
891 & l.~15 & \textit{for} Lemma~2324 of Volume~X \textit{read} Lemma~2325 of Volume~X \\
891 & l.~6 from foot & \textit{for} mehtod \textit{read} method \\
892 & l.~4 & \textit{for} the axiom 4, $\pos q \to \nec\pos q$ \textit{read} the axiom 5, $\pos q \to \nec\pos q$ \\
892 & l.~21 & \textit{for} definittion \textit{read} definition \\
893 & l.~1 & \textit{for} eviednce \textit{read} evidence \\
893 & l.~14 & \textit{for} asesrtion \textit{read} assertion \\
893 & l.~24 & \textit{delete} the second occurrence of \emph{to} \\
893 & l.~31 & \textit{for} wehther \textit{read} whether \\
\dag\,894 & proof of Proposition~3.13, l.~10 & \textit{for} Lemma~3.6 \textit{read} Lemma~3.9; the step uses Lemma~3.9 (p.~888) \\
894 & l.~11 & \textit{for} naturrally \textit{read} naturally \\
\dag\,894 & proof of Corollary~3.11, l.~18 & \textit{for} Lemma~3.12 \textit{read} Lemma~3.9; the step uses Lemma~3.9 (p.~888) \\
894 & l.~36 & \textit{delete} the second occurrence of \emph{is} \\
895 & l.~2 & \textit{after} true at every world \textit{insert} accessible from $w$ \\
895 & l.~5 & \textit{for} fraemwork \textit{read} framework \\
895 & l.~20 & \textit{for} mdoality \textit{read} modality \\
\dag\,895 & proof of Proposition~3.13, l.~23 & \textit{for} Lemma~3.2 \textit{read} Lemma~3.3; the step uses Lemma~3.3 (p.~867) \\
897 & l.~3 & \textit{for} Volume~I, Chapter~8 \textit{read} Volume~I, Chapter~6 (maximal greatness and the Barcan formula) \\
900 & l.~9 & \textit{for} ratehr \textit{read} rather \\
900 & l.~12 & \textit{for} etnirely \textit{read} entirely \\
900 & l.~6 from foot & \textit{for} p.~891 \textit{read} p.~819 (Definition~2.5) \\
900 & l.~3 from foot & \textit{for} $vRw_1 \to w_2Rv$ \textit{read} $vRw_1 \to w_1Rv$ \\
\end{errlist}

\subsection*{Chapter~4. Positive and negative conceivability \hfill \normalfont\small pp.~902--950}
\begin{errlist}
903 & l.~37 & \textit{for} Volume~VIII, Chapter~1 \textit{read} Volume~VIII, Chapter~3 (indiscernible godlike beings) \\
903 & l.~3 from foot & \textit{for} pargaraph \textit{read} paragraph \\
904 & l.~3 & \textit{for} vairation \textit{read} variation \\
906 & l.~17 & \textit{for} consideratoin \textit{read} consideration \\
908 & l.~3 & \textit{for} Lemma~2052 of Volume~X \textit{read} Lemma~2054 of Volume~X \\
908 & l.~40 & \textit{for} careuflly \textit{read} carefully \\
909 & l.~30 & \textit{for} knowlegde \textit{read} knowledge \\
910 & l.~22 & \textit{for} hypotehses \textit{read} hypotheses \\
910 & l.~40 & \textit{for} reocgnize \textit{read} recognize \\
911 & n.~2 & \textit{for} Leibniz, \emph{Monadology}~\S 54 (the necessary being exists if it is possible) \textit{read} Leibniz, \emph{Monadology}~\S 45 (the necessary being exists if it is possible) \\
912 & l.~2 & \textit{for} $P(\nec G \mid C) > P(\pos G)$ \textit{read} $P(\pos G \mid C) > P(\pos G)$ \\
912 & display (4.2) & \textit{for} $\mathrm{Conc}(p) \leftrightarrow \pos p$ \textit{read} $\mathrm{Conc}(p) \to \pos p$ \\
912 & l.~33 & \textit{for} p.~891 \textit{read} p.~819 (Corollary~2.4) \\
\dag\,913 & proof of Proposition~4.6, l.~21 & \textit{for} Lemma~4.3 \textit{read} Lemma~4.4 \\
914 & l.~20 & \textit{for} transtion \textit{read} transition \\
914 & l.~31 & \textit{for} qusetionable \textit{read} questionable \\
915 & l.~5 & \textit{for} restricsion \textit{read} restriction \\
915 & l.~13 & \textit{for} probabliity \textit{read} probability \\
915 & l.~19 & \textit{for} exetrnal \textit{read} external \\
916 & l.~15 & \textit{for} differecne \textit{read} difference \\
916 & l.~17 & \textit{for} altohugh \textit{read} although \\
916 & l.~27 & \textit{for} a proiri \textit{read} a priori \\
916 & l.~39 & \textit{for} prposal \textit{read} proposal \\
917 & l.~28 & \textit{for} Lemma~2544 of Volume~X \textit{read} Lemma~2542 of Volume~X; the result cited is on p.~3818 \\
\dag\,918 & proof of Corollary~4.7, (4.3) & \textit{for} $P(\pos G) + P(\nec\neg G) \leq 1$ \textit{read} $P(\pos G) + P(\pos\neg G) \leq 1$ \\
919 & l.~5 & \textit{for} instnce \textit{read} instance \\
\dag\,920 & proof of Proposition~4.8, l.~16 & \textit{for} Proposition~4.9 \textit{read} Proposition~4.6 \\
920 & l.~20 & \textit{for} p.~822 \textit{read} p.~812 (Definition~2.1) \\
920 & l.~4 from foot & \textit{for} latetr \textit{read} latter \\
921 & l.~21 & \textit{for} treatmnt \textit{read} treatment \\
\dag\,921 & proof of Proposition~4.9, (4.5) & \textit{for} $P(r \mid p) = P(p \mid r)\,P(p)/P(p)$ \textit{read} $P(r \mid p) = P(p \mid r)\,P(r)/P(p)$; the formula is used in its correct form on p.~925 \\
922 & l.~13 & \textit{for} porceed \textit{read} proceed \\
923 & l.~20 & \textit{for} conecrned \textit{read} concerned \\
924 & l.~4 & \textit{for} underimne \textit{read} undermine \\
925 & l.~8 & \textit{for} Volume~I, Chapter~8 \textit{read} Volume~I, Chapter~7 (the choice of S5) \\
925 & l.~35 & \textit{for} Lemma~2366 of Volume~X \textit{read} Lemma~2365 of Volume~X; the result cited is on p.~3797 \\
926 & l.~19 & \textit{for} mechnism \textit{read} mechanism \\
926 & l.~7 from foot & \textit{for} depsite \textit{read} despite \\
926 & l.~1 from foot & \textit{for} Lemma~2072 of Volume~X \textit{read} Lemma~2070 of Volume~X; the result cited is on p.~3762 \\
927 & l.~14 & \textit{for} coheernt \textit{read} coherent \\
927 & l.~18 & \textit{for} Volume~IX, Chapter~2 \textit{read} Volume~IX, Chapter~4 (creation and conservation) \\
\dag\,928 & proof of Lemma~4.11, l.~12 & \textit{for} Lemma~2578 of Volume~X \textit{read} Lemma~2581 of Volume~X \\
\dag\,928 & proof of Lemma~4.11, l.~20 & \textit{for} Proposition~4.8 \textit{read} Proposition~4.9 \\
928 & l.~22 & \textit{for} assupmtion \textit{read} assumption \\
\dag\,929 & proof of Lemma~4.12, (4.6) & \textit{for} $P(\pos G) + P(\pos G) \leq 1$ \textit{read} $P(\pos G) + P(\pos\neg G) \leq 1$ \\
\dag\,930 & proof of Lemma~4.12, l.~12 & \textit{for} Lemma~4.13 \textit{read} Lemma~4.11; the step uses Lemma~4.11 (p.~928) \\
931 & l.~6 & \textit{for} p.~814 \textit{read} p.~824 (Definition~2.9) \\
932 & l.~5 & \textit{for} througohut \textit{read} throughout \\
932 & l.~14 & \textit{for} doctirne \textit{read} doctrine \\
932 & l.~30 & \textit{for} cotnrary \textit{read} contrary \\
932 & l.~3 from foot & \textit{delete} the second occurrence of \emph{is} \\
932 & l.~40 & \textit{delete} the second occurrence of \emph{that} \\
934 & l.~1 & \textit{for} respectvely \textit{read} respectively \\
934 & l.~8 & \textit{for} evidence of absence \textit{read} absence of evidence \\
934 & l.~9 & \textit{for} idantify \textit{read} identify \\
\dag\,935 & proof of Lemma~4.13, l.~10 & \textit{for} Definition~4.1 \textit{read} Definition~4.2 \\
935 & l.~15 & \textit{for} whteher \textit{read} whether \\
935 & l.~32 & \textit{for} morevoer \textit{read} moreover \\
936 & l.~10 & \textit{for} effecitvely \textit{read} effectively \\
936 & l.~33 & \textit{delete} the second occurrence of \emph{the} \\
937 & l.~4 & \textit{for} occsaionally \textit{read} occasionally \\
937 & l.~22 & \textit{for} Lemma~1981 of Volume~X \textit{read} Lemma~1983 of Volume~X \\
938 & l.~2 & \textit{for} nothng \textit{read} nothing \\
938 & l.~16 & \textit{for} inded \textit{read} indeed \\
\dag\,938 & proof of Proposition~4.15, l.~27 & \textit{for} Definition~4.10 \textit{read} Definition~4.14 \\
939 & l.~30 & \textit{for} abstarct \textit{read} abstract \\
939 & l.~6 from foot & \textit{for} oppoennt \textit{read} opponent \\
941 & l.~12 & \textit{for} effctively \textit{read} effectively \\
\dag\,941 & proof of Proposition~4.16, (4.7) & \textit{for} $\mathrm{Conc}(\neg p) \wedge \mathrm{Conc}(p)$ \textit{read} $\neg\mathrm{Conc}(\neg p) \wedge \mathrm{Conc}(p)$; correctly printed in (4.9), p.~947 \\
942 & l.~28 & \textit{for} Oppy (1996) \textit{read} Oppy (1995) \\
943 & l.~6 & \textit{for} p.~938 \textit{read} p.~928 (Lemma~4.11) \\
943 & l.~17 & \textit{for} immeidate \textit{read} immediate \\
943 & l.~34 & \textit{for} throguhout \textit{read} throughout \\
944 & l.~6 & \textit{for} definision \textit{read} definition \\
944 & l.~29 & \textit{for} p.~808 \textit{read} p.~818 (Proposition~2.3) \\
946 & l.~4 & \textit{for} cnoverse \textit{read} converse \\
946 & display (4.8) & \textit{for} $P(\pos G \mid C) < P(\pos G)$ \textit{read} $P(\pos G \mid C) > P(\pos G)$; the formula is used in its correct form on p.~948 \\
946 & l.~5 & \textit{for} governmnet \textit{read} government \\
946 & l.~15 & \textit{for} hitehrto \textit{read} hitherto \\
946 & l.~7 from foot & \textit{for} obatin \textit{read} obtain \\
947 & display (4.10) & \textit{for} $\neg\mathrm{Conc}(\neg p) \vee \mathrm{Conc}(p)$ \textit{read} $\neg\mathrm{Conc}(\neg p) \wedge \mathrm{Conc}(p)$; correctly printed in (4.11), p.~949 \\
948 & l.~8 from foot & \textit{for} premse \textit{read} premise \\
949 & l.~3 & \textit{for} metpahysical \textit{read} metaphysical \\
949 & l.~38 & \textit{for} lgoically \textit{read} logically \\
949 & n.~17 & \textit{for} Oppy (1994) \textit{read} Oppy (1995) \\
950 & l.~15 & \textit{for} Volume~IV, Chapter~5 \textit{read} Volume~III, Chapter~5 (the transfer from coherence to possibility) \\
\end{errlist}

\subsection*{Chapter~5. From coherence to possibility \hfill \normalfont\small pp.~951--1013}
\begin{errlist}
952 & running heads, pp.~952--1012 (even) & \textit{for} Positive and negative conceivability \textit{read} From coherence to possibility \\
952 & display (5.1) & \textit{for} $\mathrm{Pos}({\chi}) \wedge \mathrm{Pos}({\psi}) \leftrightarrow \mathrm{Pos}({\chi} \sqcap {\psi})$ \textit{read} $\mathrm{Pos}({\chi}) \wedge \mathrm{Pos}({\psi}) \to \mathrm{Pos}({\chi} \sqcap {\psi})$; the formula is used in its correct form on p.~953 \\
953 & display (5.2) & \textit{for} $\emptyset \in \mathcal{P}$ \textit{read} $\emptyset \notin \mathcal{P}$ \\
954 & l.~9 & \textit{for} Anderson (1991) \textit{read} Anderson (1990) \\
955 & l.~5 & \textit{for} Lemma~1870 of Volume~X \textit{read} Lemma~1867 of Volume~X; the result cited is on p.~3737 \\
955 & display (5.3) & \textit{for} $\mathrm{MG}(z,v) \wedge w'Rv \to \mathrm{MG}(z,w')$ \textit{read} $\mathrm{MG}(z,v) \wedge vRw' \to \mathrm{MG}(z,w')$; the formula is used in its correct form on p.~958 \\
958 & l.~2 & \textit{for} necessary existence and maximal excellence \textit{read} necessary existence and necessary maximal excellence \\
958 & l.~12 & \textit{for} Lemma~2425 of Volume~X \textit{read} Lemma~2424 of Volume~X \\
958 & l.~38 & \textit{for} p.~931 \textit{read} p.~913 (Proposition~4.6) \\
959 & l.~4 & \textit{for} occurence \textit{read} occurrence \\
959 & l.~2 from foot & \textit{for} Volume~IV, Chapter~8 \textit{read} Volume~IV, Chapter~9 (the coherence of maximal excellence) \\
960 & l.~6 & \textit{for} p.~832 \textit{read} p.~842 (Theorem~2.11) \\
960 & l.~10 & \textit{for} Lemma~2061 of Volume~X \textit{read} Lemma~2057 of Volume~X; the result cited is on p.~3760 \\
961 & l.~4 & \textit{for} ${\psi} \in \mathcal{P} \wedge \nec\exists z\,({\psi}z \to {\chi}z) \Rightarrow {\chi} \in \mathcal{P}$ \textit{read} ${\psi} \in \mathcal{P} \wedge \nec\forall z\,({\psi}z \to {\chi}z) \Rightarrow {\chi} \in \mathcal{P}$ \\
961 & l.~20 & \textit{for} presuppostion \textit{read} presupposition \\
961 & l.~22 & \textit{for} Lemma~2299 of Volume~X \textit{read} Lemma~2295 of Volume~X; the result cited is on p.~3789 \\
961 & l.~27 & \textit{for} geenrally \textit{read} generally \\
962 & display (5.5) & \textit{for} $\mathrm{Pos}({\phi}) \leftrightarrow \nec\mathrm{Pos}({\phi})$ \textit{read} $\mathrm{Pos}({\phi}) \to \nec\mathrm{Pos}({\phi})$; correctly printed in (5.11), p.~967 \\
963 & l.~5 & \textit{after} great-making properties \textit{insert} jointly \\
963 & display (5.7) & \textit{for} $\mathrm{G}x \leftrightarrow \forall {\psi}\,(\mathrm{Pos}({\psi}) \leftrightarrow \pos {\psi}x)$ \textit{read} $\mathrm{G}x \leftrightarrow \forall {\psi}\,(\mathrm{Pos}({\psi}) \leftrightarrow \nec {\psi}x)$ \\
963 & l.~23 & \textit{delete} the second occurrence of \emph{a} \\
963 & display (5.6) & \textit{for} $\nec\forall y\,\mathrm{G}y$ \textit{read} $\nec\exists y\,\mathrm{G}y$; correctly printed in (5.8), p.~964 \\
\dag\,965 & proof of Theorem~5.4, (5.10) & \textit{for} $\exists x\,\mathrm{MG}(x,v) \to \exists x\,(E(x,v) \vee \mathrm{ME}(x,v))$ \textit{read} $\exists x\,\mathrm{MG}(x,v) \to \exists x\,(E(x,v) \wedge \mathrm{ME}(x,v))$ \\
\dag\,965 & proof of Theorem~5.4, (5.9) & \textit{for} $\neg(\pos G \wedge \pos G)$ \textit{read} $\neg(\pos G \wedge \pos\neg G)$; correctly printed in (5.13), p.~971 \\
\dag\,966 & proof of Proposition~5.3, l.~3 & \textit{for} Proposition~5.7 \textit{read} Proposition~5.2 \\
\dag\,966 & proof of Theorem~5.4, l.~8 & \textit{for} Proposition~5.3 \textit{read} Proposition~5.2; the step uses Proposition~5.2 (p.~956) \\
966 & l.~12 & \textit{for} stiuation \textit{read} situation \\
966 & l.~16 & \textit{for} Volume~IX, Chapter~2 \textit{read} Volume~IX, Chapter~3 (the causal role of the maximally great being) \\
967 & l.~4 & \textit{see} Corrigendum~S1 \\
968 & l.~25 & \textit{for} mxaimal \textit{read} maximal \\
968 & l.~28 & \textit{for} $\exists x\,\mathrm{MG}(x,v') \to \exists x\,(E(x,v') \vee \mathrm{ME}(x,v'))$ \textit{read} $\exists x\,\mathrm{MG}(x,v') \to \exists x\,(E(x,v') \wedge \mathrm{ME}(x,v'))$ \\
968 & l.~3 from foot & \textit{for} inetrsection \textit{read} intersection \\
969 & l.~3 & \textit{for} $\pos G \leftrightarrow \neg\pos\neg G$ \textit{read} $\pos G \to \neg\pos\neg G$ \\
969 & display (5.12) & \textit{for} $\pos(G \to \nec G)$ \textit{read} $\nec(G \to \nec G)$; correctly printed in (5.16), p.~973 \\
969 & l.~31 & \textit{for} hypothesses \textit{read} hypotheses \\
971 & l.~4 & \textit{for} undesrtanding \textit{read} understanding \\
971 & l.~17 & \textit{for} recmmend \textit{read} recommend \\
971 & l.~18 & \textit{for} $@ \in \pos\,\exists y\, \mathrm{MG}(y,\cdot)$ \textit{read} $@ \in \nec\,\exists y\, \mathrm{MG}(y,\cdot)$ \\
971 & display (5.14) & \textit{for} $\mathrm{MG}(x,w) \iff \exists v'\,(wRv' \to E(x,v') \wedge \mathrm{ME}(x,v'))$ \textit{read} $\mathrm{MG}(x,w) \iff \forall v'\,(wRv' \to E(x,v') \wedge \mathrm{ME}(x,v'))$; the formula is used in its correct form on p.~974 \\
971 & l.~29 & \textit{for} the converse premise, Assumption~3.1 of Volume~I \textit{read} the possibility premise, Assumption~3.1 of Volume~I \\
972 & l.~22 & \textit{for} interpretattion \textit{read} interpretation \\
972 & l.~24 & \textit{for} closed under entailment and under conjunction (the finite intersection property) \textit{read} jointly instantiable in every finite subfamily (the finite intersection property) \\
972 & display (5.15) & \textit{for} $\mathcal{P}_0 = \{\mathrm{NE}, \mathrm{NME}\} \supseteq \mathcal{P}$ \textit{read} $\mathcal{P}_0 = \{\mathrm{NE}, \mathrm{NME}\} \subseteq \mathcal{P}$ \\
972 & l.~37 & \textit{for} $\mathrm{ME}(y,w_1) \iff \mathrm{Op}(y,w_2) \wedge \mathrm{Os}(y,w_2) \wedge \mathrm{MP}(y,w_2)$ \textit{read} $\mathrm{ME}(y,w_2) \iff \mathrm{Op}(y,w_2) \wedge \mathrm{Os}(y,w_2) \wedge \mathrm{MP}(y,w_2)$ \\
973 & l.~16 & \textit{delete} the second occurrence of \emph{of} \\
973 & l.~19 & \textit{for} Lemma~2394 of Volume~X \textit{read} Lemma~2390 of Volume~X; the result cited is on p.~3800 \\
974 & l.~9 & \textit{for} every infinite subfamily \textit{read} every finite subfamily \\
975 & l.~20 & \textit{for} jointly instantiated at some accessible world \textit{read} jointly instantiated at some one accessible world \\
975 & l.~21 & \textit{for} detrminate \textit{read} determinate \\
975 & l.~32 & \textit{for} p.~908 \textit{read} p.~907 (Lemma~4.4) \\
976 & display (5.17) & \textit{for} $@ \notin \nec\,\exists y\, \mathrm{MG}(y,\cdot)$ \textit{read} $@ \in \nec\,\exists y\, \mathrm{MG}(y,\cdot)$; the formula is used in its correct form on p.~980 \\
976 & l.~25 & \textit{for} necessraily \textit{read} necessarily \\
978 & l.~10 & \textit{for} $\mathrm{ME}(x,w') \iff \mathrm{Op}(x,w') \vee \mathrm{Os}(x,w') \wedge \mathrm{MP}(x,w')$ \textit{read} $\mathrm{ME}(x,w') \iff \mathrm{Op}(x,w') \wedge \mathrm{Os}(x,w') \wedge \mathrm{MP}(x,w')$ \\
979 & display (5.18) & \textit{for} $\mathrm{MG}(y,v') \to \mathrm{NE}(y,v') \vee \mathrm{NME}(y,v')$ \textit{read} $\mathrm{MG}(y,v') \to \mathrm{NE}(y,v') \wedge \mathrm{NME}(y,v')$ \\
980 & l.~14 & \textit{for} Lemma~2234 of Volume~X \textit{read} Lemma~2230 of Volume~X \\
980 & l.~16 & \textit{after} exists \textit{insert} at every world accessible from $v$ \\
981 & display (5.19) & \textit{for} $\mathrm{Pos}(\neg {\psi}) \to \neg\mathrm{Pos}({\psi})$ \textit{read} $\mathrm{Pos}(\neg {\psi}) \leftrightarrow \neg\mathrm{Pos}({\psi})$; correctly printed in (5.20), p.~983 \\
981 & l.~5 from foot & \textit{for} Lemma~1902 of Volume~X \textit{read} Lemma~1901 of Volume~X \\
983 & proof of Theorem~5.6 & \textit{see} Corrigendum~S1 \\
985 & l.~7 & \textit{for} Lemma~2451 of Volume~X \textit{read} Lemma~2449 of Volume~X \\
985 & l.~22 & \textit{for} Volume~VII, Chapter~5 \textit{read} Volume~VII, Chapter~4 (fine-tuning) \\
986 & l.~27 & \textit{for} condtiional \textit{read} conditional \\
\dag\,987 & proof of Proposition~5.7, (5.22) & \textit{for} ${\psi} \in \mathcal{P} \wedge \nec\forall x\,({\psi}x \leftrightarrow {\chi}x) \Rightarrow {\chi} \in \mathcal{P}$ \textit{read} ${\psi} \in \mathcal{P} \wedge \nec\forall x\,({\psi}x \to {\chi}x) \Rightarrow {\chi} \in \mathcal{P}$; the formula is used in its correct form on p.~989 \\
\dag\,988 & proof of Proposition~5.7, l.~15 & \textit{for} Definition~5.10 \textit{read} Definition~5.5 \\
988 & l.~23 & \textit{for} $\mathrm{MG}(y,w_0) \leftrightarrow \mathrm{NE}(y,w_0) \wedge \mathrm{NME}(y,w_0)$ \textit{read} $\mathrm{MG}(y,w_0) \to \mathrm{NE}(y,w_0) \wedge \mathrm{NME}(y,w_0)$ \\
988 & l.~25 & \textit{for} $\pos\exists z\,\mathrm{G}z$ \textit{read} $\nec\exists z\,\mathrm{G}z$ \\
989 & l.~20 & \textit{for} Lemma~1880 of Volume~X \textit{read} Lemma~1876 of Volume~X \\
990 & l.~16 & \textit{for} cnosequently \textit{read} consequently \\
990 & l.~31 & \textit{for} verison \textit{read} version \\
991 & l.~5 & \textit{for} Lemma~1890 of Volume~X \textit{read} Lemma~1888 of Volume~X \\
\dag\,991 & proof of Corollary~5.9, l.~15 & \textit{for} Proposition~5.8 \textit{read} Proposition~5.3 \\
\dag\,991 & proof of Corollary~5.9, l.~28 & \textit{for} Theorem~5.4 \textit{read} Theorem~5.6 \\
993 & l.~7 & \textit{for} analyss \textit{read} analysis \\
993 & display (5.25) & \textit{for} $\nec(G \leftrightarrow \nec G)$ \textit{read} $\nec(G \to \nec G)$ \\
\dag\,994 & proof of Corollary~5.12, l.~4 & \textit{for} Proposition~5.3 \textit{read} Proposition~5.7; the step uses Proposition~5.7 (p.~987) \\
\dag\,994 & proof of Corollary~5.12, (5.26) & \textit{for} $\mathrm{Pos}({\psi}) \leftrightarrow \nec\mathrm{Pos}({\psi})$ \textit{read} $\mathrm{Pos}({\psi}) \to \nec\mathrm{Pos}({\psi})$ \\
994 & l.~14 & \textit{for} $\forall z\,\mathrm{MG}(z,w_0) \to \exists z\,(E(z,w_0) \wedge \mathrm{ME}(z,w_0))$ \textit{read} $\exists z\,\mathrm{MG}(z,w_0) \to \exists z\,(E(z,w_0) \wedge \mathrm{ME}(z,w_0))$ \\
994 & l.~27 & \textit{for} ${\phi} \in \mathcal{P} \wedge \nec\forall y\,({\phi}y \leftrightarrow {\chi}y) \Rightarrow {\chi} \in \mathcal{P}$ \textit{read} ${\phi} \in \mathcal{P} \wedge \nec\forall y\,({\phi}y \to {\chi}y) \Rightarrow {\chi} \in \mathcal{P}$ \\
995 & l.~4 & \textit{for} commited \textit{read} committed \\
\dag\,995 & proof of Corollary~5.12, (5.27) & \textit{for} $\nec\forall z\,\mathrm{G}z$ \textit{read} $\nec\exists z\,\mathrm{G}z$ \\
\dag\,995 & proof of Corollary~5.12, (5.28) & \textit{for} $\nec G \to \neg\pos\neg G$ \textit{read} $\pos G \to \neg\pos\neg G$; the formula is used in its correct form on p.~997 \\
996 & display (5.29) & \textit{for} $\mathrm{MG}(y,w_1) \vee w_1Rw' \to \mathrm{MG}(y,w')$ \textit{read} $\mathrm{MG}(y,w_1) \wedge w_1Rw' \to \mathrm{MG}(y,w')$ \\
997 & l.~26 & \textit{for} Gödel (1994) \textit{read} Gödel (1995) \\
\dag\,998 & proof of Proposition~5.16, l.~13 & \textit{for} $\mathrm{ME}(z,w_0) \iff \mathrm{Op}(z,w_0) \wedge \mathrm{Os}(z,w_0) \vee \mathrm{MP}(z,w_0)$ \textit{read} $\mathrm{ME}(z,w_0) \iff \mathrm{Op}(z,w_0) \wedge \mathrm{Os}(z,w_0) \wedge \mathrm{MP}(z,w_0)$ \\
\dag\,999 & proof of Proposition~5.16, l.~10 & \textit{for} Definition~5.14 \textit{read} Definition~5.13 \\
999 & l.~15 & \textit{for} evidecne \textit{read} evidence \\
999 & l.~17 & \textit{for} excellnece \textit{read} excellence \\
999 & l.~25 & \textit{for} knoweldge \textit{read} knowledge \\
1000 & display (5.30) & \textit{for} $\mathrm{NE}, \mathrm{NME} \notin \mathcal{P}$ \textit{read} $\mathrm{NE}, \mathrm{NME} \in \mathcal{P}$ \\
1000 & display (5.31) & \textit{for} $@ \notin \pos\,\exists x\, \mathrm{MG}(x,\cdot)$ \textit{read} $@ \in \pos\,\exists x\, \mathrm{MG}(x,\cdot)$ \\
1000 & l.~8 from foot & \textit{for} $@ \in \nec\,\forall y\, \mathrm{MG}(y,\cdot)$ \textit{read} $@ \in \nec\,\exists y\, \mathrm{MG}(y,\cdot)$ \\
1001 & display (5.33) & \textit{for} $\neg(\nec G \wedge \pos\neg G)$ \textit{read} $\neg(\pos G \wedge \pos\neg G)$; correctly printed in (5.38), p.~1005 \\
1002 & l.~36 & \textit{for} paradgim \textit{read} paradigm \\
1002 & l.~2 from foot & \textit{delete} the second occurrence of \emph{of} \\
\dag\,1003 & proof of Lemma~5.17, (5.35) & \textit{for} $\mathrm{NE}(x,w) \iff \forall w_0\,(w_0Rw \to E(x,w_0))$ \textit{read} $\mathrm{NE}(x,w) \iff \forall w_0\,(wRw_0 \to E(x,w_0))$ \\
\dag\,1004 & proof of Lemma~5.17, (5.36) & \textit{for} $\mathrm{Pos}({\chi}) \vee \mathrm{Pos}({\psi}) \to \mathrm{Pos}({\chi} \sqcap {\psi})$ \textit{read} $\mathrm{Pos}({\chi}) \wedge \mathrm{Pos}({\psi}) \to \mathrm{Pos}({\chi} \sqcap {\psi})$ \\
\dag\,1005 & proof of Lemma~5.19, (5.37) & \textit{for} $\pos\neg\exists z\,\mathrm{MG}(z) \to \exists z\,\mathrm{MG}(z)$ \textit{read} $\pos\neg\exists z\,\mathrm{MG}(z) \to \neg\exists z\,\mathrm{MG}(z)$ \\
1005 & l.~25 & \textit{for} Volume~IX, Chapter~9 \textit{read} Volume~IX, Chapter~7 (foreknowledge) \\
\dag\,1005 & proof of Lemma~5.19, l.~27 & \textit{for} Lemma~2336 of Volume~X \textit{read} Lemma~2340 of Volume~X; the result cited is on p.~3794 \\
1005 & l.~33 & \textit{for} intsead \textit{read} instead \\
1006 & l.~7 & \textit{for} p.~942 \textit{read} p.~941 (Proposition~4.16) \\
\dag\,1006 & proof of Lemma~5.19, l.~30 & \textit{for} Corollary~5.12 \textit{read} Corollary~5.9 \\
1008 & l.~14 & \textit{for} materail \textit{read} material \\
1009 & l.~4 & \textit{for} Proposition~4.1 of Volume~I, the main result \textit{read} Theorem~3.4 of Volume~I, the main result \\
\dag\,1010 & proof of Lemma~5.21, l.~21 & \textit{for} $\mathrm{Pos}({\phi}) \to \pos\mathrm{Pos}({\phi})$ \textit{read} $\mathrm{Pos}({\phi}) \to \nec\mathrm{Pos}({\phi})$ \\
1010 & l.~32 & \textit{for} Soebl (1987) \textit{read} Sobel (1987) \\
1010 & n.~20 & \textit{for} Anderson (1989) \textit{read} Anderson (1990) \\
1011 & l.~28 & \textit{for} contorversy \textit{read} controversy \\
1012 & l.~17 & \textit{for} paargraph \textit{read} paragraph \\
1013 & l.~12 & \textit{for} $\mathrm{MG}(z,w') \wedge w'Rw \to \mathrm{MG}(z,w')$ \textit{read} $\mathrm{MG}(z,w') \wedge w'Rw \to \mathrm{MG}(z,w)$ \\
1013 & l.~21 & \textit{for} conrtary \textit{read} contrary \\
1013 & l.~24 & \textit{delete} the second occurrence of \emph{is} \\
1013 & l.~28 & \textit{for} Volume~III, Chapter~8 \textit{read} Volume~IV, Chapter~8 (godlikeness and maximal greatness) \\
\end{errlist}

\subsection*{Chapter~6. The converse premise and the evidence of absence \hfill \normalfont\small pp.~1014--1077}
\begin{errlist}
1015 & l.~17 & \textit{for} advantge \textit{read} advantage \\
1017 & l.~13 & \textit{for} phenmenon \textit{read} phenomenon \\
1017 & l.~19 & \textit{for} p.~819 \textit{read} p.~809 (Lemma~1.2) \\
1017 & l.~21 & \textit{for} differnece \textit{read} difference \\
1017 & l.~39 & \textit{for} continegnt \textit{read} contingent \\
1019 & l.~20 & \textit{for} rahter \textit{read} rather \\
1020 & l.~12 & \textit{for} Volume~VII, Chapter~2 \textit{read} Volume~VI, Chapter~2 (the logical problem of evil) \\
1020 & n.~23 & \textit{for} Oppy (2005) \textit{read} Oppy (2006) \\
1021 & l.~8 & \textit{for} perference \textit{read} preference \\
1021 & l.~16 & \textit{for} Lemma~1902 of Volume~X \textit{read} Lemma~1905 of Volume~X \\
1021 & l.~28 & \textit{for} Plantigna (1974b) \textit{read} Plantinga (1974b) \\
1021 & l.~6 from foot & \textit{for} itnerpretation \textit{read} interpretation \\
1022 & l.~25 & \textit{for} aplication \textit{read} application \\
1023 & l.~24 & \textit{for} Lemma~1871 of Volume~X \textit{read} Lemma~1866 of Volume~X \\
1023 & l.~26 & \textit{for} dsepite \textit{read} despite \\
1024 & l.~24 & \textit{for} philosophicl \textit{read} philosophical \\
1024 & l.~31 & \textit{for} impotrance \textit{read} importance \\
1025 & l.~37 & \textit{for} Lemma~2085 of Volume~X \textit{read} Lemma~2091 of Volume~X; the result cited is on p.~3764 \\
1027 & l.~15 & \textit{for} appropriatte \textit{read} appropriate \\
1027 & l.~29 & \textit{for} Volume~VIII, Chapter~5 \textit{read} Volume~VIII, Chapter~3 (indiscernible godlike beings) \\
1028 & l.~14 & \textit{for} convrese \textit{read} converse \\
1028 & l.~40 & \textit{for} p.~989 \textit{read} p.~899 (Theorem~3.14) \\
1029 & proof of Proposition~6.3 & \textit{see} Corrigendum~S3 \\
\dag\,1030 & proof of Corollary~6.4, (6.5) & \textit{for} $\pos G \to \pos\neg G$ \textit{read} $\pos G \to \neg\pos\neg G$; correctly printed in (6.7), p.~1033 \\
1030 & l.~26 & \textit{for} $\mathrm{ME}(y,w_1) \iff \mathrm{Op}(y,w_1) \wedge \mathrm{Os}(y,w_2) \wedge \mathrm{MP}(y,w_1)$ \textit{read} $\mathrm{ME}(y,w_1) \iff \mathrm{Op}(y,w_1) \wedge \mathrm{Os}(y,w_1) \wedge \mathrm{MP}(y,w_1)$ \\
\dag\,1031 & proof of Corollary~6.4, (6.6) & \textit{for} $\mathrm{MG}(x,w) \vee wRv' \to \mathrm{MG}(x,v')$ \textit{read} $\mathrm{MG}(x,w) \wedge wRv' \to \mathrm{MG}(x,v')$ \\
1031 & l.~5 & \textit{for} Rowe (1976) \textit{read} Rowe (1975) \\
1031 & l.~16 & \textit{for} geniune \textit{read} genuine \\
\dag\,1031 & proof of Corollary~6.4, l.~27 & \textit{for} Lemma~2523 of Volume~X \textit{read} Lemma~2526 of Volume~X \\
1032 & l.~4 & \textit{for} Volume~I, Chapter~9 \textit{read} Volume~I, Chapter~8 (the comparison with Gödel's argument) \\
1034 & l.~27 & \textit{for} srtictly \textit{read} strictly \\
1035 & l.~3 & \textit{for} entirley \textit{read} entirely \\
1035 & l.~32 & \textit{for} coherecne \textit{read} coherence \\
1036 & l.~8 & \textit{for} occsion \textit{read} occasion \\
1037 & l.~15 & \textit{for} reprsent \textit{read} represent \\
1039 & l.~14 & \textit{for} p.~928 \textit{read} p.~929 (Lemma~4.12) \\
1039 & l.~26 & \textit{for} Definition~2.2 of Volume~I \textit{read} Definition~2.4 of Volume~I \\
1040 & l.~24 & \textit{for} Plantniga (1974b) \textit{read} Plantinga (1974b) \\
1041 & l.~14 & \textit{for} repersent \textit{read} represent \\
1041 & l.~39 & \textit{for} mainatin \textit{read} maintain \\
1042 & l.~7 & \textit{for} simialrly \textit{read} similarly \\
1042 & l.~11 & \textit{for} mtehod \textit{read} method \\
\dag\,1042 & proof of Proposition~6.8, (6.9) & \textit{for} $\mathrm{NE}(x,w_0) \iff \exists w_2\,(w_0Rw_2 \to E(x,w_2))$ \textit{read} $\mathrm{NE}(x,w_0) \iff \forall w_2\,(w_0Rw_2 \to E(x,w_2))$ \\
1042 & l.~17 & \textit{for} menaing \textit{read} meaning \\
\dag\,1043 & proof of Proposition~6.8, (6.11) & \textit{for} $\exists y\,\mathrm{MG}(y,w) \to \exists y\,(E(y,w) \vee \mathrm{ME}(y,w))$ \textit{read} $\exists y\,\mathrm{MG}(y,w) \to \exists y\,(E(y,w) \wedge \mathrm{ME}(y,w))$ \\
1043 & l.~29 & \textit{for} cmomitted \textit{read} committed \\
1043 & l.~34 & \textit{for} apparetnly \textit{read} apparently \\
1044 & display (6.13) & \textit{for} $\neg(\pos G \vee \pos\neg G)$ \textit{read} $\neg(\pos G \wedge \pos\neg G)$ \\
1044 & display (6.12) & \textit{for} $P(q \mid q) = P(q \mid r)\,P(r)/P(q)$ \textit{read} $P(r \mid q) = P(q \mid r)\,P(r)/P(q)$ \\
1044 & l.~20 & \textit{for} regrdless \textit{read} regardless \\
1044 & l.~25 & \textit{for} Lemma~2315 of Volume~X \textit{read} Lemma~2312 of Volume~X; the result cited is on p.~3791 \\
1044 & l.~32 & \textit{for} restrcition \textit{read} restriction \\
1045 & l.~20 & \textit{for} namley \textit{read} namely \\
1045 & l.~25 & \textit{for} perahps \textit{read} perhaps \\
1045 & display (6.15) & \textit{for} $\mathrm{MG}(y,w') \iff \forall v'\,(w'Rv' \leftrightarrow E(y,v') \wedge \mathrm{ME}(y,v'))$ \textit{read} $\mathrm{MG}(y,w') \iff \forall v'\,(w'Rv' \to E(y,v') \wedge \mathrm{ME}(y,v'))$; correctly printed in (6.17), p.~1046 \\
1046 & l.~23 & \textit{for} pressumably \textit{read} presumably \\
\dag\,1046 & proof of Lemma~6.9, l.~27 & \textit{for} Proposition~6.13 \textit{read} Proposition~6.8 \\
\dag\,1047 & proof of Lemma~6.9, (6.19) & \textit{for} $\pos\neg\exists z\,\mathrm{MG}(z) \leftrightarrow \neg\exists z\,\mathrm{MG}(z)$ \textit{read} $\pos\neg\exists z\,\mathrm{MG}(z) \to \neg\exists z\,\mathrm{MG}(z)$ \\
1047 & l.~15 & \textit{for} rperesent \textit{read} represent \\
1050 & l.~36 & \textit{for} probbaility \textit{read} probability \\
1051 & l.~5 & \textit{for} prefeernce \textit{read} preference \\
1051 & l.~27 & \textit{after} maximally excellent \textit{insert} at $w'$ \\
1052 & l.~33 & \textit{for} retsriction \textit{read} restriction \\
1053 & l.~3 & \textit{for} presmably \textit{read} presumably \\
1053 & l.~15 & \textit{for} threfore \textit{read} therefore \\
1053 & l.~30 & \textit{for} significnt \textit{read} significant \\
1054 & l.~4 & \textit{for} exisetnce \textit{read} existence \\
1055 & l.~13 & \textit{for} interst \textit{read} interest \\
\dag\,1055 & proof of Proposition~6.13, l.~19 & \textit{for} Definition~6.6 \textit{read} Definition~6.7 \\
1055 & l.~22 & \textit{for} presance \textit{read} presence \\
1056 & l.~7 & \textit{for} prcatical \textit{read} practical \\
\dag\,1056 & proof of Proposition~6.13, l.~11 & \textit{for} Definition~6.1 \textit{read} Definition~6.6 \\
\dag\,1056 & proof of Proposition~6.13, l.~19 & \textit{for} Proposition~6.8 \textit{read} Proposition~6.3 \\
\dag\,1056 & proof of Proposition~6.13, l.~25 & \textit{for} Lemma~2282 of Volume~X \textit{read} Lemma~2287 of Volume~X; the result cited is on p.~3788 \\
1057 & l.~6 & \textit{for} sequnece \textit{read} sequence \\
1057 & l.~24 & \textit{for} desiralbe \textit{read} desirable \\
1057 & l.~30 & \textit{for} anaylsis \textit{read} analysis \\
1058 & l.~16 & \textit{for} depndence \textit{read} dependence \\
1059 & display (6.28) & \textit{for} $\mathrm{NE}(y,w) \iff \forall w_0\,(w_0Rw \to E(y,w_0))$ \textit{read} $\mathrm{NE}(y,w) \iff \forall w_0\,(wRw_0 \to E(y,w_0))$; correctly printed in (6.32), p.~1063 \\
1060 & l.~18 & \textit{for} stattement \textit{read} statement \\
1060 & l.~40 & \textit{for} Lemma~2527 of Volume~X \textit{read} Lemma~2521 of Volume~X; the result cited is on p.~3816 \\
1062 & l.~2 & \textit{for} intuittion \textit{read} intuition \\
1062 & l.~40 & \textit{for} Lemma~2400 of Volume~X \textit{read} Lemma~2404 of Volume~X; the result cited is on p.~3802 \\
1063 & l.~12 & \textit{for} simplicty \textit{read} simplicity \\
1063 & l.~27 & \textit{for} $\mathrm{Conc}(p) \to \nec p$ \textit{read} $\mathrm{Conc}(p) \to \pos p$ \\
1063 & display (6.33) & \textit{for} $\mathrm{NME}(y,w') \iff \exists w_2\,(w'Rw_2 \to \mathrm{ME}(y,w_2))$ \textit{read} $\mathrm{NME}(y,w') \iff \forall w_2\,(w'Rw_2 \to \mathrm{ME}(y,w_2))$; correctly printed in (6.34), p.~1069 \\
1063 & l.~36 & \textit{for} equlaly \textit{read} equally \\
1063 & l.~40 & \textit{for} $\mathrm{MG}(x,v) \to \mathrm{NE}(x,v) \vee \mathrm{NME}(x,v)$ \textit{read} $\mathrm{MG}(x,v) \to \mathrm{NE}(x,v) \wedge \mathrm{NME}(x,v)$ \\
1064 & l.~8 from foot & \textit{for} sbusequent \textit{read} subsequent \\
1065 & l.~16 & \textit{for} suppsoe \textit{read} suppose \\
1065 & l.~24 & \textit{for} represant \textit{read} represent \\
1065 & l.~7 from foot & \textit{for} especailly \textit{read} especially \\
1065 & l.~40 & \textit{for} Lemma~2253 of Volume~X \textit{read} Lemma~2247 of Volume~X \\
1067 & l.~8 & \textit{for} Oppy (2007) \textit{read} Oppy (2006) \\
1067 & l.~28 & \textit{for} unelss \textit{read} unless \\
1067 & l.~34 & \textit{for} possibly maximally excellent \textit{read} possibly maximally great \\
1069 & l.~5 & \textit{for} effectievly \textit{read} effectively \\
\dag\,1069 & proof of Proposition~6.16, l.~2 from foot & \textit{for} Lemma~2477 of Volume~X \textit{read} Lemma~2476 of Volume~X \\
\dag\,1070 & proof of Proposition~6.16, l.~17 & \textit{for} Lemma~6.9 \textit{read} Lemma~6.11 \\
\dag\,1070 & proof of Corollary~6.17, (6.35) & \textit{for} $\exists x\,\mathrm{MG}(x,w_2) \to \forall x\,(E(x,w_2) \wedge \mathrm{ME}(x,w_2))$ \textit{read} $\exists x\,\mathrm{MG}(x,w_2) \to \exists x\,(E(x,w_2) \wedge \mathrm{ME}(x,w_2))$; the formula is used in its correct form on p.~1072 \\
1070 & l.~25 & \textit{for} claiemd \textit{read} claimed \\
1070 & l.~8 from foot & \textit{for} dotcrine \textit{read} doctrine \\
1071 & l.~3 & \textit{for} statemment \textit{read} statement \\
1071 & l.~29 & \textit{for} p.~851 \textit{read} p.~861 (Definition~3.1) \\
1072 & display (6.36) & \textit{for} $@ \notin \pos\,\exists z\, \mathrm{MG}(z,\cdot)$ \textit{read} $@ \in \pos\,\exists z\, \mathrm{MG}(z,\cdot)$ \\
\dag\,1074 & proof of Proposition~6.18, l.~4 & \textit{for} Lemma~6.12 \textit{read} Lemma~6.11 \\
\dag\,1074 & proof of Lemma~6.19, l.~10 & \textit{for} Definition~6.1 \textit{read} Definition~6.5; the step uses Definition~6.5 (p.~1035) \\
\dag\,1075 & proof of Lemma~6.19, l.~1 & \textit{for} Proposition~6.13 \textit{read} Proposition~6.15 \\
1075 & l.~25 & \textit{for} p.~883 \textit{read} p.~893 (Corollary~3.11) \\
1075 & l.~1 from foot & \textit{for} applicaiton \textit{read} application \\
1077 & l.~22 & \textit{for} striclty \textit{read} strictly \\
1077 & l.~24 & \textit{for} argment \textit{read} argument \\
1077 & l.~1 from foot & \textit{for} p.~1054 \textit{read} p.~1064 (Proposition~6.15) \\
\end{errlist}

\subsection*{Chapter~7. Symmetry-breaking: the asymmetry of support \hfill \normalfont\small pp.~1078--1138}
\begin{errlist}
1079 & l.~22 & \textit{for} encounetr \textit{read} encounter \\
1079 & l.~23 & \textit{for} individaul \textit{read} individual \\
1079 & display (7.2) & \textit{for} $\mathrm{NE}(x,w_2) \iff \forall w\,(w_1Rw \to E(x,w))$ \textit{read} $\mathrm{NE}(x,w_2) \iff \forall w\,(w_2Rw \to E(x,w))$ \\
1079 & l.~4 from foot & \textit{for} Lemma~2504 of Volume~X \textit{read} Lemma~2500 of Volume~X \\
1080 & l.~15 & \textit{for} versoin \textit{read} version \\
1080 & l.~37 & \textit{for} diffrence \textit{read} difference \\
1081 & l.~13 & \textit{for} $\mathrm{ME}(z,u) \iff \mathrm{Op}(z,u) \vee \mathrm{Os}(z,u) \wedge \mathrm{MP}(z,u)$ \textit{read} $\mathrm{ME}(z,u) \iff \mathrm{Op}(z,u) \wedge \mathrm{Os}(z,u) \wedge \mathrm{MP}(z,u)$ \\
1081 & l.~27 & \textit{for} intrpretation \textit{read} interpretation \\
1081 & l.~29 & \textit{for} respnose \textit{read} response \\
1083 & l.~1 & \textit{for} combinnation \textit{read} combination \\
1083 & l.~7 & \textit{for} spuport \textit{read} support \\
\dag\,1083 & proof of Lemma~7.2, (7.4) & \textit{for} $@ \in \pos\,\exists x\, \mathrm{MG}(x,\cdot)$ \textit{read} $@ \in \nec\,\exists x\, \mathrm{MG}(x,\cdot)$; the formula is used in its correct form on p.~1086 \\
1083 & l.~25 & \textit{for} respnse \textit{read} response \\
1083 & l.~26 & \textit{for} considertion \textit{read} consideration \\
1083 & l.~37 & \textit{for} prmeise \textit{read} premise \\
1085 & l.~5 & \textit{for} perhps \textit{read} perhaps \\
1085 & l.~23 & \textit{for} somethng \textit{read} something \\
\dag\,1086 & proof of Corollary~7.3, l.~2 & \textit{for} Lemma~2565 of Volume~X \textit{read} Lemma~2564 of Volume~X \\
1086 & l.~23 & \textit{for} concluison \textit{read} conclusion \\
\dag\,1086 & proof of Corollary~7.3, (7.6) & \textit{for} $P(p \mid r) = P(r \mid p)\,P(p)/P(p)$ \textit{read} $P(p \mid r) = P(r \mid p)\,P(p)/P(r)$ \\
\dag\,1086 & proof of Corollary~7.3, l.~1 from foot & \textit{for} Lemma~2498 of Volume~X \textit{read} Lemma~2496 of Volume~X \\
\dag\,1087 & proof of Corollary~7.3, l.~5 & \textit{for} Lemma~7.7 \textit{read} Lemma~7.2 \\
1088 & l.~5 & \textit{for} pervious \textit{read} previous \\
1088 & l.~16 & \textit{for} contemproary \textit{read} contemporary \\
1088 & l.~28 & \textit{for} ineded \textit{read} indeed \\
1089 & l.~26 & \textit{for} itnroduce \textit{read} introduce \\
1089 & l.~37 & \textit{for} scenarrio \textit{read} scenario \\
1091 & l.~1 & \textit{for} p.~985 \textit{read} p.~995 (Definition~5.14) \\
1091 & l.~3 & \textit{for} fundamntal \textit{read} fundamental \\
1091 & l.~9 & \textit{for} maxmially \textit{read} maximally \\
1091 & l.~10 & \textit{for} apparrent \textit{read} apparent \\
1091 & l.~16 & \textit{for} sentnce \textit{read} sentence \\
1092 & l.~29 & \textit{for} $\mathrm{ME}(x,w') \iff \mathrm{Op}(x,w') \wedge \mathrm{Os}(x,w') \vee \mathrm{MP}(x,w')$ \textit{read} $\mathrm{ME}(x,w') \iff \mathrm{Op}(x,w') \wedge \mathrm{Os}(x,w') \wedge \mathrm{MP}(x,w')$ \\
1092 & l.~34 & \textit{for} reasonnable \textit{read} reasonable \\
1093 & l.~7 from foot & \textit{for} rleevant \textit{read} relevant \\
1093 & l.~5 from foot & \textit{for} definitly \textit{read} definitely \\
1093 & n.~6 & \textit{for} Soebl (2004) \textit{read} Sobel (2004) \\
1094 & l.~20 & \textit{for} idntify \textit{read} identify \\
1094 & l.~25 & \textit{for} demonstratte \textit{read} demonstrate \\
1094 & l.~28 & \textit{for} qauntifier \textit{read} quantifier \\
1094 & l.~39 & \textit{for} Sobel (2003) \textit{read} Sobel (2004) \\
1095 & l.~4 & \textit{for} niether \textit{read} neither \\
1095 & l.~7 from foot & \textit{for} mdoest \textit{read} modest \\
1095 & l.~40 & \textit{for} presecne \textit{read} presence \\
1096 & display (7.13) & \textit{for} $\exists x\,\mathrm{MG}(x,w_0) \leftrightarrow \exists x\,(E(x,w_0) \wedge \mathrm{ME}(x,w_0))$ \textit{read} $\exists x\,\mathrm{MG}(x,w_0) \to \exists x\,(E(x,w_0) \wedge \mathrm{ME}(x,w_0))$ \\
1097 & l.~14 & \textit{for} episetmic \textit{read} epistemic \\
1097 & l.~25 & \textit{for} conseuqently \textit{read} consequently \\
\dag\,1098 & proof of Proposition~7.6, l.~29 & \textit{for} Corollary~7.3 \textit{read} Corollary~7.5 \\
1099 & l.~9 & \textit{for} presuasive \textit{read} persuasive \\
1099 & display (7.18) & \textit{for} $\mathrm{NE}(x,w_2) \iff \forall v\,(w_1Rv \to E(x,v))$ \textit{read} $\mathrm{NE}(x,w_1) \iff \forall v\,(w_1Rv \to E(x,v))$; correctly printed in (7.19), p.~1099 \\
1099 & l.~27 & \textit{for} Volume~VII, Chapter~2 \textit{read} Volume~VIII, Chapter~2 (uniqueness) \\
1100 & l.~11 & \textit{for} prboable \textit{read} probable \\
1100 & l.~16 & \textit{for} inevtable \textit{read} inevitable \\
1101 & l.~1 & \textit{for} addrss \textit{read} address \\
1101 & l.~33 & \textit{for} meainng \textit{read} meaning \\
1104 & l.~10 & \textit{for} conclusoin \textit{read} conclusion \\
\dag\,1104 & proof of Lemma~7.7, (7.24) & \textit{for} $\nec(G \to \pos G)$ \textit{read} $\nec(G \to \nec G)$ \\
\dag\,1105 & proof of Lemma~7.10, l.~8 & \textit{for} Lemma~7.4 \textit{read} Lemma~7.2; the step uses Lemma~7.2 (p.~1083) \\
\dag\,1105 & proof of Lemma~7.10, (7.27) & \textit{for} $@ \in \nec\,\forall x\, \mathrm{MG}(x,\cdot)$ \textit{read} $@ \in \nec\,\exists x\, \mathrm{MG}(x,\cdot)$; the formula is used in its correct form on p.~1109 \\
1106 & l.~7 & \textit{delete} the second occurrence of \emph{the} \\
\dag\,1106 & proof of Lemma~7.10, l.~14 & \textit{for} Definition~7.11 \textit{read} Definition~7.8 \\
1107 & l.~1 & \textit{for} Lemma~1919 of Volume~X \textit{read} Lemma~1913 of Volume~X \\
1108 & l.~31 & \textit{for} Volume~IV, Chapter~7 \textit{read} Volume~IV, Chapter~6 (the compossibility of the attributes) \\
1109 & l.~12 & \textit{for} adderss \textit{read} address \\
1110 & l.~17 & \textit{for} combintaion \textit{read} combination \\
1110 & l.~35 & \textit{for} p.~813 \textit{read} p.~823 (Theorem~2.7) \\
1111 & l.~37 & \textit{for} emhpasis \textit{read} emphasis \\
1111 & l.~38 & \textit{for} otbain \textit{read} obtain \\
1112 & l.~5 & \textit{for} p.~1004 \textit{read} p.~1104 (Definition~7.8) \\
1112 & l.~22 & \textit{for} situtaion \textit{read} situation \\
1112 & l.~7 from foot & \textit{for} Volume~III, Chapter~3 \textit{read} Volume~III, Chapter~2 (conceivability and possibility) \\
1113 & l.~8 & \textit{for} conjecturre \textit{read} conjecture \\
1114 & l.~6 from foot & \textit{for} Lemma~1817 of Volume~X \textit{read} Lemma~1811 of Volume~X \\
1115 & display (7.29) & \textit{for} $\mathrm{MG}(y,w') \vee w'Rw_0 \to \mathrm{MG}(y,w_0)$ \textit{read} $\mathrm{MG}(y,w') \wedge w'Rw_0 \to \mathrm{MG}(y,w_0)$; the formula is used in its correct form on p.~1119 \\
1115 & l.~5 from foot & \textit{for} omittd \textit{read} omitted \\
1117 & l.~2 & \textit{for} Volume~IV, Chapter~8 \textit{read} Volume~III, Chapter~8 (parodies of the argument) \\
1117 & l.~3 from foot & \textit{for} symmtry \textit{read} symmetry \\
1119 & l.~13 & \textit{for} prctical \textit{read} practical \\
1121 & l.~28 & \textit{for} refrered \textit{read} referred \\
1122 & l.~30 & \textit{for} plauisble \textit{read} plausible \\
1124 & l.~24 & \textit{for} Volume~IV, Chapter~6 \textit{read} Volume~III, Chapter~6 (the converse premise) \\
1125 & l.~2 & \textit{for} exists and is maximally excellent at some world accessible from $v$ \textit{read} exists and is maximally excellent at every world accessible from $v$ \\
\dag\,1125 & proof of Corollary~7.15, (7.32) & \textit{for} $\mathrm{Conc}(q) \to \nec q$ \textit{read} $\mathrm{Conc}(q) \to \pos q$ \\
1125 & l.~24 & \textit{for} importnce \textit{read} importance \\
1126 & l.~3 & \textit{for} cnosiderable \textit{read} considerable \\
1126 & l.~12 & \textit{for} conceivbaility \textit{read} conceivability \\
1126 & l.~15 & \textit{for} cappacity \textit{read} capacity \\
1126 & l.~16 & \textit{delete} the second occurrence of \emph{the} \\
1126 & l.~25 & \textit{for} objcetive \textit{read} objective \\
1127 & l.~20 & \textit{for} Lemma~2108 of Volume~X \textit{read} Lemma~2104 of Volume~X; the result cited is on p.~3766 \\
1127 & l.~2 from foot & \textit{for} seprate \textit{read} separate \\
1128 & l.~10 & \textit{for} metaphysiacl \textit{read} metaphysical \\
1128 & display (7.35) & \textit{for} $\pos\neg\exists x\,\mathrm{MG}(x) \to \exists x\,\mathrm{MG}(x)$ \textit{read} $\pos\neg\exists x\,\mathrm{MG}(x) \to \neg\exists x\,\mathrm{MG}(x)$ \\
1128 & display (7.34) & \textit{for} $\mathrm{ME}(y,w_0) \iff \mathrm{Op}(y,w_0) \wedge \mathrm{Os}(y,w_0) \vee \mathrm{MP}(y,w_0)$ \textit{read} $\mathrm{ME}(y,w_0) \iff \mathrm{Op}(y,w_0) \wedge \mathrm{Os}(y,w_0) \wedge \mathrm{MP}(y,w_0)$ \\
1129 & l.~3 & \textit{for} Sobel (2005) \textit{read} Sobel (2004) \\
\dag\,1129 & proof of Proposition~7.17, l.~20 & \textit{for} $P(\pos G \mid C) > P(\nec G)$ \textit{read} $P(\pos G \mid C) > P(\pos G)$ \\
1129 & l.~24 & \textit{for} inedpendently \textit{read} independently \\
\dag\,1130 & proof of Lemma~7.19, l.~4 & \textit{for} Lemma~1867 of Volume~X \textit{read} Lemma~1866 of Volume~X; the result cited is on p.~3737 \\
1130 & l.~36 & \textit{for} aparently \textit{read} apparently \\
1132 & l.~17 & \textit{for} prefernce \textit{read} preference \\
1132 & l.~22 & \textit{for} probabiilty \textit{read} probability \\
1133 & l.~11 & \textit{for} recognze \textit{read} recognize \\
1133 & l.~33 & \textit{for} oppossite \textit{read} opposite \\
1134 & l.~37 & \textit{for} p.~919 \textit{read} p.~918 (Corollary~4.7) \\
1136 & l.~16 & \textit{for} Lemma~2630 of Volume~X \textit{read} Lemma~2624 of Volume~X \\
1138 & l.~9 & \textit{for} p.~1082 \textit{read} p.~1081 (Definition~7.1) \\
1138 & l.~31 & \textit{for} essentail \textit{read} essential \\
\end{errlist}

\subsection*{Chapter~8. Parodies: islands, devils and the maximally evil being \hfill \normalfont\small pp.~1139--1195}
\begin{errlist}
1141 & l.~1 & \textit{for} impliict \textit{read} implicit \\
1141 & l.~18 & \textit{for} whther \textit{read} whether \\
\dag\,1141 & proof of Lemma~8.4, l.~29 & \textit{for} $\mathrm{ME}(z,w') \iff \mathrm{Op}(z,w') \wedge \mathrm{Os}(z,w') \vee \mathrm{MP}(z,w')$ \textit{read} $\mathrm{ME}(z,w') \iff \mathrm{Op}(z,w') \wedge \mathrm{Os}(z,w') \wedge \mathrm{MP}(z,w')$ \\
1142 & l.~9 & \textit{for} p.~1183 \textit{read} p.~1083 (Lemma~7.2) \\
1142 & l.~1 from foot & \textit{for} Plantinga (1975) \textit{read} Plantinga (1974a) \\
1143 & l.~22 & \textit{for} Plantinga (1973) \textit{read} Plantinga (1974a) \\
1143 & l.~26 & \textit{for} frmaework \textit{read} framework \\
1144 & l.~12 & \textit{for} p.~868 \textit{read} p.~867 (Lemma~3.3) \\
1144 & l.~16 & \textit{for} adeqaute \textit{read} adequate \\
1144 & l.~26 & \textit{for} $\exists z\,\mathrm{MG}(z,w_1) \to \forall z\,(E(z,w_1) \wedge \mathrm{ME}(z,w_1))$ \textit{read} $\exists z\,\mathrm{MG}(z,w_1) \to \exists z\,(E(z,w_1) \wedge \mathrm{ME}(z,w_1))$ \\
1144 & l.~1 from foot & \textit{for} excellecne \textit{read} excellence \\
1145 & l.~1 & \textit{for} amibguous \textit{read} ambiguous \\
1145 & l.~5 & \textit{for} Lemma~2523 of Volume~X \textit{read} Lemma~2517 of Volume~X; the result cited is on p.~3815 \\
1145 & l.~12 & \textit{for} proces \textit{read} process \\
1146 & l.~11 & \textit{for} desiarble \textit{read} desirable \\
1146 & l.~15 & \textit{for} occurrecne \textit{read} occurrence \\
1146 & l.~34 & \textit{for} Theorem~3.4 of Volume~I, the converse argument \textit{read} Proposition~4.1 of Volume~I, the converse argument \\
1147 & l.~18 & \textit{for} techncial \textit{read} technical \\
1148 & l.~11 & \textit{for} bleieve \textit{read} believe \\
1149 & l.~7 & \textit{for} adquate \textit{read} adequate \\
\dag\,1151 & proof of Corollary~8.5, l.~3 & \textit{for} Definition~8.7 \textit{read} Definition~8.3 \\
\dag\,1151 & proof of Corollary~8.5, l.~20 & \textit{for} Lemma~8.2 \textit{read} Lemma~8.4 \\
1151 & l.~27 & \textit{for} Volume~IX, Chapter~8 \textit{read} Volume~X, Chapter~8 (the Lean modules) \\
\dag\,1151 & proof of Theorem~8.6, (8.9) & \textit{for} $\mathrm{NME}(y,w') \iff \forall v\,(w'Rv \to \mathrm{ME}(y,w'))$ \textit{read} $\mathrm{NME}(y,w') \iff \forall v\,(w'Rv \to \mathrm{ME}(y,v))$; correctly printed in (8.13), p.~1157 \\
1151 & l.~38 & \textit{for} prefernece \textit{read} preference \\
1152 & l.~25 & \textit{after} supported \textit{insert} independently of the conclusion \\
1154 & l.~10 & \textit{for} ebmarrassment \textit{read} embarrassment \\
1154 & l.~11 & \textit{for} p.~1091 \textit{read} p.~1090 (Corollary~7.5) \\
1156 & l.~1 & \textit{for} wherreas \textit{read} whereas \\
1156 & l.~11 & \textit{for} Lemma~2018 of Volume~X \textit{read} Lemma~2016 of Volume~X \\
1156 & l.~17 & \textit{for} intutiive \textit{read} intuitive \\
1156 & l.~35 & \textit{for} ceratinly \textit{read} certainly \\
1156 & l.~4 from foot & \textit{for} cocneivable \textit{read} conceivable \\
1157 & l.~8 & \textit{delete} the second occurrence of \emph{the} \\
1157 & l.~17 & \textit{for} seriuos \textit{read} serious \\
1157 & l.~24 & \textit{for} p.~993 \textit{read} p.~983 (Theorem~5.6) \\
1159 & l.~14 & \textit{for} dependant \textit{read} dependent \\
1161 & l.~7 & \textit{for} possessoin \textit{read} possession \\
1161 & l.~36 & \textit{for} philosophiacl \textit{read} philosophical \\
1162 & l.~7 & \textit{for} omniptent \textit{read} omnipotent \\
1162 & display (8.15) & \textit{for} $@ \in \nec\,\exists x\, \mathrm{MG}(x,\cdot)$ \textit{read} $@ \in \pos\,\exists x\, \mathrm{MG}(x,\cdot)$; correctly printed in (8.21), p.~1166 \\
1163 & l.~10 & \textit{for} $\mathrm{MG}(x,w) \iff \forall u\,(wRu \to E(x,u) \vee \mathrm{ME}(x,u))$ \textit{read} $\mathrm{MG}(x,w) \iff \forall u\,(wRu \to E(x,u) \wedge \mathrm{ME}(x,u))$ \\
\dag\,1163 & proof of Corollary~8.8, l.~11 & \textit{for} Definition~8.3 \textit{read} Definition~8.1; the step uses Definition~8.1 (p.~1140) \\
1164 & l.~13 & \textit{for} chllenge \textit{read} challenge \\
1164 & l.~38 & \textit{for} singificant \textit{read} significant \\
1165 & l.~11 & \textit{for} emiprical \textit{read} empirical \\
1165 & l.~16 & \textit{for} derivaiton \textit{read} derivation \\
1166 & l.~2 from foot & \textit{for} antciipate \textit{read} anticipate \\
1167 & l.~8 & \textit{for} improtance \textit{read} importance \\
1169 & l.~29 & \textit{for} ecxlusion \textit{read} exclusion \\
1169 & l.~38 & \textit{for} occrred \textit{read} occurred \\
1170 & l.~5 & \textit{for} conevrse \textit{read} converse \\
1170 & l.~9 & \textit{for} remmarkable \textit{read} remarkable \\
1171 & l.~10 & \textit{for} Volume~IV, Chapter~2 \textit{read} Volume~V, Chapter~2 (existence as a predicate) \\
1172 & l.~18 & \textit{for} exculsion \textit{read} exclusion \\
\dag\,1173 & proof of Lemma~8.12, l.~25 & \textit{for} Definition~8.7 \textit{read} Definition~8.11 \\
1173 & l.~26 & \textit{for} tradiiton \textit{read} tradition \\
1173 & l.~5 from foot & \textit{for} desrciption \textit{read} description \\
1173 & l.~1 from foot & \textit{for} governmant \textit{read} government \\
1173 & n.~3 & \textit{for} \emph{Proslogion}~4 (that God cannot be conceived not to exist) \textit{read} \emph{Proslogion}~3 (that God cannot be conceived not to exist) \\
1174 & l.~17 & \textit{for} uniofrm \textit{read} uniform \\
1174 & l.~22 & \textit{for} constrcution \textit{read} construction \\
1175 & l.~31 & \textit{for} Definition~2.3 of Volume~I \textit{read} Definition~2.4 of Volume~I \\
1176 & l.~23 & \textit{for} Lemma~2418 of Volume~X \textit{read} Lemma~2415 of Volume~X; the result cited is on p.~3803 \\
1176 & l.~6 from foot & \textit{for} separatte \textit{read} separate \\
1176 & l.~36 & \textit{for} surrely \textit{read} surely \\
\dag\,1177 & proof of Proposition~8.13, l.~20 & \textit{for} Lemma~8.4 \textit{read} Lemma~8.9 \\
1177 & l.~26 & \textit{for} conrtast \textit{read} contrast \\
\dag\,1178 & proof of Proposition~8.13, l.~5 & \textit{for} Lemma~1940 of Volume~X \textit{read} Lemma~1939 of Volume~X \\
\dag\,1178 & proof of Proposition~8.13, l.~29 & \textit{for} Lemma~2078 of Volume~X \textit{read} Lemma~2079 of Volume~X \\
\dag\,1179 & proof of Proposition~8.15, (8.27) & \textit{for} $\mathrm{MG}(z,w_2) \wedge uRw_2 \to \mathrm{MG}(z,u)$ \textit{read} $\mathrm{MG}(z,w_2) \wedge w_2Ru \to \mathrm{MG}(z,u)$; correctly printed in (8.29), p.~1184 \\
\dag\,1179 & proof of Proposition~8.15, l.~27 & \textit{for} Corollary~8.8 \textit{read} Corollary~8.5; the step uses Corollary~8.5 (p.~1151) \\
1180 & l.~30 & \textit{for} Plntinga (1974a) \textit{read} Plantinga (1974a) \\
1180 & l.~31 & \textit{for} occasoin \textit{read} occasion \\
\dag\,1181 & proof of Theorem~8.16, l.~13 & \textit{for} Lemma~8.9 \textit{read} Lemma~8.12; the step uses Lemma~8.12 (p.~1173) \\
1182 & l.~5 & \textit{for} conesquently \textit{read} consequently \\
1182 & l.~23 & \textit{for} undremine \textit{read} undermine \\
1183 & l.~40 & \textit{for} later \textit{read} latter \\
1184 & l.~3 & \textit{for} p.~1310 \textit{read} p.~1130 (Lemma~7.19) \\
1184 & l.~22 & \textit{for} Volume~III, Chapter~7 \textit{read} Volume~III, Chapter~9 (the ancestral premise) \\
\dag\,1184 & proof of Theorem~8.17, (8.30) & \textit{for} $\nec\neg\exists y\,\mathrm{MG}(y) \to \neg\exists y\,\mathrm{MG}(y)$ \textit{read} $\pos\neg\exists y\,\mathrm{MG}(y) \to \neg\exists y\,\mathrm{MG}(y)$; correctly printed in (8.31), p.~1186 \\
1186 & l.~28 & \textit{for} quantifeir \textit{read} quantifier \\
1186 & l.~4 from foot & \textit{for} p.~892 \textit{read} p.~891 (Definition~3.10) \\
1187 & l.~8 from foot & \textit{for} advatnage \textit{read} advantage \\
1187 & l.~38 & \textit{for} Lemma~2430 of Volume~X \textit{read} Lemma~2431 of Volume~X \\
1188 & l.~24 & \textit{for} knolwedge \textit{read} knowledge \\
1189 & l.~3 from foot & \textit{for} withrdawn \textit{read} withdrawn \\
1190 & l.~5 & \textit{for} ifnerence \textit{read} inference \\
1190 & l.~10 & \textit{for} annalogous \textit{read} analogous \\
1190 & l.~11 & \textit{for} Lemma~2265 of Volume~X \textit{read} Lemma~2268 of Volume~X; the result cited is on p.~3785 \\
1190 & l.~17 & \textit{for} p.~1075 \textit{read} p.~1074 (Proposition~6.18) \\
1190 & l.~20 & \textit{for} icnluding \textit{read} including \\
1190 & l.~23 & \textit{for} oposite \textit{read} opposite \\
1191 & l.~17 & \textit{for} hithreto \textit{read} hitherto \\
1191 & l.~24 & \textit{for} $\mathrm{MG}(x,u) \to \mathrm{NE}(x,u) \vee \mathrm{NME}(x,u)$ \textit{read} $\mathrm{MG}(x,u) \to \mathrm{NE}(x,u) \wedge \mathrm{NME}(x,u)$ \\
1191 & l.~32 & \textit{delete} the second occurrence of \emph{a} \\
1192 & l.~10 & \textit{for} excatly \textit{read} exactly \\
1193 & l.~17 & \textit{for} secsion \textit{read} section \\
1194 & l.~15 & \textit{for} dneial \textit{read} denial \\
1195 & l.~6 & \textit{for} Gaunilo, \emph{Pro insipiente}~5 (the lost island) \textit{read} Gaunilo, \emph{Pro insipiente}~6 (the lost island) \\
\end{errlist}

\subsection*{Chapter~9. The ancestral premise and weaker modal logics \hfill \normalfont\small pp.~1196--1250}
\begin{errlist}
1198 & running heads, pp.~1198--1250 (even) & \textit{for} Parodies \textit{read} The ancestral premise \\
1198 & l.~2 & \textit{for} Volume~V, Chapter~3 \textit{read} Volume~V, Chapter~5 (the Barcan formula) \\
1198 & l.~11 & \textit{for} soundnss \textit{read} soundness \\
1198 & display (9.1) & \textit{for} $p \leftrightarrow \nec\pos p$ \textit{read} $p \to \nec\pos p$ \\
1200 & l.~13 & \textit{for} the axiom 4, $\nec q \to q$ \textit{read} the axiom T, $\nec q \to q$ \\
1200 & l.~1 from foot & \textit{for} $\nec(q \to p) \to (\pos q \to \nec p)$ \textit{read} $\nec(q \to p) \to (\nec q \to \nec p)$ \\
1201 & display (9.4) & \textit{for} $vRw_0 \wedge w_0Rw \leftrightarrow vRw$ \textit{read} $vRw_0 \wedge w_0Rw \to vRw$; correctly printed in (9.6), p.~1202 \\
1201 & l.~6 from foot & \textit{for} paralel \textit{read} parallel \\
1202 & display (9.5) & \textit{for} $(w'Rv' \vee w'Rw_1) \to v'Rw_1$ \textit{read} $(w'Rv' \wedge w'Rw_1) \to v'Rw_1$; correctly printed in (9.9), p.~1208 \\
1202 & l.~12 & \textit{for} argualby \textit{read} arguably \\
1203 & l.~8 & \textit{for} purpsoe \textit{read} purpose \\
1203 & l.~17 & \textit{for} setcion \textit{read} section \\
1203 & l.~19 & \textit{for} condtion \textit{read} condition \\
1203 & l.~20 & \textit{for} demonstarte \textit{read} demonstrate \\
1203 & l.~21 & \textit{for} immediatte \textit{read} immediate \\
1204 & l.~16 & \textit{for} maxiaml \textit{read} maximal \\
1204 & l.~25 & \textit{for} $\mathrm{MG}(y,v) \to \mathrm{NE}(y,v) \vee \mathrm{NME}(y,v)$ \textit{read} $\mathrm{MG}(y,v) \to \mathrm{NE}(y,v) \wedge \mathrm{NME}(y,v)$ \\
1205 & display (9.8) & \textit{for} $\nec\exists z\,\mathrm{MG}(z) \to \nec\exists z\,\mathrm{MG}(z)$ \textit{read} $\pos\exists z\,\mathrm{MG}(z) \to \nec\exists z\,\mathrm{MG}(z)$; correctly printed in (9.10), p.~1210 \\
1206 & l.~8 & \textit{for} prpoosal \textit{read} proposal \\
1208 & l.~4 & \textit{for} $\mathrm{MG}(z,u) \iff \forall u\,(uRv' \to E(z,v') \wedge \mathrm{ME}(z,v'))$ \textit{read} $\mathrm{MG}(z,u) \iff \forall v'\,(uRv' \to E(z,v') \wedge \mathrm{ME}(z,v'))$ \\
1208 & l.~17 & \textit{for} $\nec\nec q \to \nec q$ \textit{read} $\pos\nec q \to \nec q$ \\
1208 & l.~22 & \textit{for} prsumably \textit{read} presumably \\
1208 & l.~26 & \textit{for} $\mathrm{MG}(y,w_2) \to \mathrm{NE}(y,w_2) \vee \mathrm{NME}(y,w_2)$ \textit{read} $\mathrm{MG}(y,w_2) \to \mathrm{NE}(y,w_2) \wedge \mathrm{NME}(y,w_2)$ \\
1210 & display (9.11) & \textit{for} $\nec(r \wedge p) \to (\nec r \wedge \nec p)$ \textit{read} $\nec(r \wedge p) \leftrightarrow (\nec r \wedge \nec p)$ \\
1211 & l.~9 & \textit{for} $\mathrm{NME}(y,v') \iff \exists w'\,(v'Rw' \to \mathrm{ME}(y,w'))$ \textit{read} $\mathrm{NME}(y,v') \iff \forall w'\,(v'Rw' \to \mathrm{ME}(y,w'))$ \\
1211 & l.~13 & \textit{for} Hughes and Cresswell (1996) \textit{read} Chandler (1976) (the failure of S4 for metaphysical modality) \\
1211 & l.~17 & \textit{for} Definition~2.4 of Volume~I (maximal excellence) \textit{read} Definition~2.2 of Volume~I (maximal excellence) \\
1211 & l.~23 & \textit{for} $\nec\nec r \to r$ \textit{read} $\pos\nec r \to r$ \\
1212 & l.~1 & \textit{for} $uRw_1$ by symmetry \textit{read} $uRw_1$ by transitivity \\
1212 & l.~11 & \textit{for} Lemma~2065 of Volume~X \textit{read} Lemma~2059 of Volume~X; the result cited is on p.~3760 \\
1212 & l.~15 & \textit{for} vesrion \textit{read} version \\
1213 & l.~17 & \textit{for} prceed \textit{read} proceed \\
1213 & n.~14 & \textit{for} Hughes and Cresswell (1997) \textit{read} Hughes and Cresswell (1996) \\
1214 & l.~11 & \textit{for} devleop \textit{read} develop \\
1214 & l.~1 from foot & \textit{for} $\mathrm{NME}(y,w') \iff \exists w\,(w'Rw \to \mathrm{ME}(y,w))$ \textit{read} $\mathrm{NME}(y,w') \iff \forall w\,(w'Rw \to \mathrm{ME}(y,w))$ \\
1217 & display (9.13) & \textit{for} $\forall y\,\mathrm{MG}(y,w') \to \exists y\,(E(y,w') \wedge \mathrm{ME}(y,w'))$ \textit{read} $\exists y\,\mathrm{MG}(y,w') \to \exists y\,(E(y,w') \wedge \mathrm{ME}(y,w'))$ \\
1217 & l.~25 & \textit{for} lagnuage \textit{read} language \\
1218 & l.~4 & \textit{for} Lemma~2431 of Volume~X \textit{read} Lemma~2432 of Volume~X; the result cited is on p.~3805 \\
1218 & l.~6 & \textit{for} permisisble \textit{read} permissible \\
1218 & display (9.14) & \textit{for} $\pos(p \vee q) \to (\pos p \vee \pos q)$ \textit{read} $\pos(p \vee q) \leftrightarrow (\pos p \vee \pos q)$ \\
1218 & l.~21 & \textit{for} tendecny \textit{read} tendency \\
\dag\,1219 & proof of Theorem~9.7, (9.15) & \textit{for} $w'Rv \vee vRw_0 \to w'Rw_0$ \textit{read} $w'Rv \wedge vRw_0 \to w'Rw_0$; correctly printed in (9.22), p.~1224 \\
1219 & l.~11 & \textit{for} p.~1019 \textit{read} p.~1109 (Definition~7.11) \\
1219 & l.~5 from foot & \textit{delete} the second occurrence of \emph{to} \\
\dag\,1220 & proof of Theorem~9.6, l.~23 & \textit{for} Definition~9.1 \textit{read} Definition~9.2 \\
1222 & display (9.18) & \textit{for} $\mathrm{ME}(y,v') \iff \mathrm{Op}(y,v') \vee \mathrm{Os}(y,v') \wedge \mathrm{MP}(y,v')$ \textit{read} $\mathrm{ME}(y,v') \iff \mathrm{Op}(y,v') \wedge \mathrm{Os}(y,v') \wedge \mathrm{MP}(y,v')$ \\
1222 & display (9.19) & \textit{for} $@ \in \nec\,\exists z\, \mathrm{MG}(z,\cdot)$ \textit{read} $@ \in \pos\,\exists z\, \mathrm{MG}(z,\cdot)$ \\
1222 & display (9.17) & \textit{for} $q \to \pos\pos q$ \textit{read} $q \to \nec\pos q$; the formula is used in its correct form on p.~1223 \\
1223 & display (9.20) & \textit{for} $v' \models \nec p \iff \exists u\,(v'Ru \wedge u \models p)$ \textit{read} $v' \models \pos p \iff \exists u\,(v'Ru \wedge u \models p)$ \\
1224 & l.~2 & \textit{for} cohrence \textit{read} coherence \\
1224 & l.~18 & \textit{for} $\exists y\,\mathrm{MG}(y,w_1) \to \exists y\,(E(y,w_1) \vee \mathrm{ME}(y,w_1))$ \textit{read} $\exists y\,\mathrm{MG}(y,w_1) \to \exists y\,(E(y,w_1) \wedge \mathrm{ME}(y,w_1))$ \\
1224 & display (9.21) & \textit{for} $\pos\exists x\,\mathrm{MG}(x) \to \pos\exists x\,\mathrm{MG}(x)$ \textit{read} $\pos\exists x\,\mathrm{MG}(x) \to \nec\exists x\,\mathrm{MG}(x)$ \\
1224 & l.~25 & \textit{for} applcation \textit{read} application \\
1225 & l.~7 & \textit{for} inculding \textit{read} including \\
\dag\,1225 & proof of Lemma~9.9, l.~23 & \textit{for} $@ \in \pos\,\forall y\, \mathrm{MG}(y,\cdot)$ \textit{read} $@ \in \pos\,\exists y\, \mathrm{MG}(y,\cdot)$ \\
\dag\,1225 & proof of Lemma~9.9, l.~25 & \textit{for} Lemma~2034 of Volume~X \textit{read} Lemma~2036 of Volume~X \\
1226 & l.~6 & \textit{for} nottion \textit{read} notion \\
1226 & l.~28 & \textit{for} altrenative \textit{read} alternative \\
1227 & l.~1 & \textit{for} $\pos\pos q \leftrightarrow \pos q$ \textit{read} $\pos\pos q \to \pos q$ \\
1227 & l.~35 & \textit{for} $\neg\pos q \to \nec\neg q$ \textit{read} $\neg\pos q \leftrightarrow \nec\neg q$ \\
1229 & l.~8 & \textit{for} the axiom B, $\nec r \to \nec\nec r$ \textit{read} the axiom 4, $\nec r \to \nec\nec r$ \\
\dag\,1229 & proof of Proposition~9.10, l.~20 & \textit{for} Definition~9.1 \textit{read} Definition~9.5; the step uses Definition~9.5 (p.~1212) \\
1230 & l.~23 & \textit{for} suffiicent \textit{read} sufficient \\
1230 & display (9.23) & \textit{for} $(vRw_2 \wedge vRv') \leftrightarrow w_2Rv'$ \textit{read} $(vRw_2 \wedge vRv') \to w_2Rv'$ \\
1231 & display (9.24) & \textit{for} $\mathrm{NE}(x,w_2) \iff \forall w'\,(w_1Rw' \to E(x,w'))$ \textit{read} $\mathrm{NE}(x,w_2) \iff \forall w'\,(w_2Rw' \to E(x,w'))$ \\
1231 & l.~24 & \textit{for} $\mathrm{ME}(y,w_1) \iff \mathrm{Op}(y,w_1) \wedge \mathrm{Os}(y,w_1) \vee \mathrm{MP}(y,w_1)$ \textit{read} $\mathrm{ME}(y,w_1) \iff \mathrm{Op}(y,w_1) \wedge \mathrm{Os}(y,w_1) \wedge \mathrm{MP}(y,w_1)$ \\
\dag\,1233 & proof of Lemma~9.11, l.~9 & \textit{for} Lemma~9.13 \textit{read} Lemma~9.9 \\
\dag\,1233 & proof of Lemma~9.11, l.~1 from foot & \textit{for} Lemma~2361 of Volume~X \textit{read} Lemma~2360 of Volume~X \\
\dag\,1234 & proof of Lemma~9.11, l.~19 & \textit{for} $\mathrm{MG}(z,w_1) \wedge w_2Rw \to \mathrm{MG}(z,w)$ \textit{read} $\mathrm{MG}(z,w_2) \wedge w_2Rw \to \mathrm{MG}(z,w)$ \\
1234 & l.~25 & \textit{for} conrtoversy \textit{read} controversy \\
1235 & l.~19 & \textit{for} explaantion \textit{read} explanation \\
1235 & display (9.27) & \textit{for} $\nec(p \vee q) \leftrightarrow (\pos p \vee \pos q)$ \textit{read} $\pos(p \vee q) \leftrightarrow (\pos p \vee \pos q)$; the formula is used in its correct form on p.~1236 \\
1235 & display (9.26) & \textit{for} $\nec r \wedge \pos p \to \pos(r \vee p)$ \textit{read} $\nec r \wedge \pos p \to \pos(r \wedge p)$; correctly printed in (9.28), p.~1238 \\
1236 & l.~2 & \textit{for} $w'Rw_2 \to w'Rw_2$ \textit{read} $w'Rw_2 \to w_2Rw'$ \\
1236 & l.~3 & \textit{for} opposte \textit{read} opposite \\
1236 & l.~6 from foot & \textit{for} $uRv'$ by reflexivity \textit{read} $uRu$ by reflexivity \\
1237 & l.~3 & \textit{for} entierly \textit{read} entirely \\
1238 & l.~8 & \textit{for} unlss \textit{read} unless \\
1238 & l.~11 & \textit{for} Lemma~2275 of Volume~X \textit{read} Lemma~2270 of Volume~X; the result cited is on p.~3786 \\
1238 & l.~4 from foot & \textit{for} Euclidean and symmetric \textit{read} Euclidean and reflexive \\
1240 & display (9.29) & \textit{for} $\nec(r \to p) \to (\nec r \leftrightarrow \nec p)$ \textit{read} $\nec(r \to p) \to (\nec r \to \nec p)$ \\
1240 & l.~18 & \textit{for} $\mathrm{NE}(x,w) \iff \forall w\,(wRw' \to E(x,w'))$ \textit{read} $\mathrm{NE}(x,w) \iff \forall w'\,(wRw' \to E(x,w'))$ \\
1242 & l.~8 & \textit{after} S5 conditions \textit{insert} on $R$ \\
1242 & l.~19 & \textit{for} traetment \textit{read} treatment \\
1242 & l.~7 from foot & \textit{for} Lemma~2269 of Volume~X \textit{read} Lemma~2275 of Volume~X \\
\dag\,1243 & proof of Lemma~9.13, l.~4 & \textit{for} Definition~9.5 \textit{read} Definition~9.2; the step uses Definition~9.2 (p.~1202) \\
\dag\,1244 & proof of Proposition~9.16, l.~19 & \textit{for} $\mathrm{MG}(y,w_1) \iff \forall w'\,(w_1Rw' \to E(y,w') \vee \mathrm{ME}(y,w'))$ \textit{read} $\mathrm{MG}(y,w_1) \iff \forall w'\,(w_1Rw' \to E(y,w') \wedge \mathrm{ME}(y,w'))$ \\
\dag\,1246 & proof of Lemma~9.17, l.~29 & \textit{for} Definition~9.8 \textit{read} Definition~9.5 \\
1247 & display (9.31) & \textit{for} $\neg\pos p \leftrightarrow \nec p$ \textit{read} $\neg\pos p \leftrightarrow \nec\neg p$ \\
1247 & l.~39 & \textit{for} inconsisetnt \textit{read} inconsistent \\
\dag\,1248 & proof of Proposition~9.18, l.~12 & \textit{for} Theorem~9.7 \textit{read} Theorem~9.6; the step uses Theorem~9.6 (p.~1219) \\
\dag\,1248 & proof of Proposition~9.18, (9.32) & \textit{for} $\mathrm{MG}(y,w') \wedge w'Ru \leftrightarrow \mathrm{MG}(y,u)$ \textit{read} $\mathrm{MG}(y,w') \wedge w'Ru \to \mathrm{MG}(y,u)$ \\
1248 & l.~22 & \textit{for} p.~1234 \textit{read} p.~1243 (Definition~9.14) \\
1248 & l.~36 & \textit{for} Lemma~2034 of Volume~X \textit{read} Lemma~2033 of Volume~X \\
\dag\,1249 & proof of Proposition~9.18, (9.33) & \textit{for} $\pos\pos q \to \nec q$ \textit{read} $\pos\pos q \to \pos q$; the formula is used in its correct form on p.~1250 \\
\dag\,1249 & proof of Proposition~9.18, l.~27 & \textit{for} Lemma~9.13 \textit{read} Lemma~9.17 \\
1250 & l.~12 & \textit{for} admitetdly \textit{read} admittedly \\
\end{errlist}

\subsection*{Chapter~10. Leibniz's problem and the compossibility of perfections \hfill \normalfont\small pp.~1251--1301}
\begin{errlist}
1252 & l.~12 & \textit{for} p.~1211 \textit{read} p.~1111 (Theorem~7.12) \\
1252 & l.~20 & \textit{after} every task \textit{insert} whose description does not quantify over the agent's own powers \\
1252 & l.~26 & \textit{for} deefnder \textit{read} defender \\
1252 & l.~36 & \textit{for} proponant \textit{read} proponent \\
1254 & l.~21 & \textit{for} Lemma~2368 of Volume~X \textit{read} Lemma~2373 of Volume~X; the result cited is on p.~3798 \\
1257 & l.~6 & \textit{for} Lemma~2525 of Volume~X \textit{read} Lemma~2524 of Volume~X \\
1257 & l.~14 & \textit{for} reasnoable \textit{read} reasonable \\
1258 & l.~5 & \textit{for} Gödel (1996) \textit{read} Gödel (1995) \\
\dag\,1258 & proof of Corollary~10.3, l.~9 & \textit{for} $\mathcal{P}_0 \supseteq \mathcal{P} \Rightarrow \exists x\,\exists w\,\forall {\chi} \in \mathcal{P}_0\; {\chi}(x,w)$ \textit{read} $\mathcal{P}_0 \subseteq \mathcal{P} \Rightarrow \exists x\,\exists w\,\forall {\chi} \in \mathcal{P}_0\; {\chi}(x,w)$ \\
1258 & l.~25 & \textit{for} the paradox of the stone shows a limitation \textit{read} the paradox of the stone shows no limitation \\
1259 & l.~1 & \textit{for} witdhrawn \textit{read} withdrawn \\
1259 & display (10.5) & \textit{for} ${\chi}, {\phi} \notin \mathcal{P} \Rightarrow {\chi} \sqcap {\phi} \in \mathcal{P}$ \textit{read} ${\chi}, {\phi} \in \mathcal{P} \Rightarrow {\chi} \sqcap {\phi} \in \mathcal{P}$ \\
1260 & l.~24 & \textit{for} eipstemic \textit{read} epistemic \\
1261 & l.~17 & \textit{for} essenital \textit{read} essential \\
1262 & l.~9 & \textit{for} convntion \textit{read} convention \\
1262 & l.~13 & \textit{for} compossibilty \textit{read} compossibility \\
1262 & l.~5 from foot & \textit{for} Morris (1988) \textit{read} Morris (1987) \\
1262 & l.~4 from foot & \textit{for} omnipotnce \textit{read} omnipotence \\
1262 & l.~2 from foot & \textit{for} wehnever \textit{read} whenever \\
1263 & l.~24 & \textit{for} asseity \textit{read} aseity \\
1265 & l.~2 & \textit{for} maening \textit{read} meaning \\
1265 & l.~24 & \textit{for} supose \textit{read} suppose \\
1265 & l.~3 from foot & \textit{for} conventoin \textit{read} convention \\
1267 & l.~5 & \textit{delete} the second occurrence of \emph{to} \\
1267 & l.~14 & \textit{for} psoition \textit{read} position \\
1267 & l.~20 & \textit{after} proper filter \textit{insert} of properties \\
\dag\,1267 & proof of Proposition~10.5, l.~24 & \textit{for} $K_{z}q \leftrightarrow K_{z}K_{z}q$ \textit{read} $K_{z}q \to K_{z}K_{z}q$ \\
1268 & l.~2 & \textit{for} accoridngly \textit{read} accordingly \\
1268 & l.~15 & \textit{for} necessray \textit{read} necessary \\
1269 & l.~6 from foot & \textit{for} Lemma~2418 of Volume~X \textit{read} Lemma~2421 of Volume~X \\
\dag\,1270 & proof of Proposition~10.9, (10.11) & \textit{for} $\mathrm{NE}(x) \leftrightarrow \neg\exists z\,D(x,z)$ \textit{read} $\mathrm{NE}(x) \to \neg\exists z\,D(x,z)$; correctly printed in (10.14), p.~1275 \\
\dag\,1271 & proof of Proposition~10.9, l.~7 & \textit{for} Definition~10.1 \textit{read} Definition~10.4; the step uses Definition~10.4 (p.~1265) \\
1271 & l.~12 & \textit{for} omniscience entails omnipotence \textit{read} omnipotence entails omniscience \\
1271 & l.~29 & \textit{for} remakrable \textit{read} remarkable \\
1272 & l.~26 & \textit{for} can do wrong but never does \textit{read} cannot do wrong \\
1272 & l.~28 & \textit{for} prefreence \textit{read} preference \\
1274 & l.~26 & \textit{for} p.~945 \textit{read} p.~935 (Lemma~4.13) \\
1275 & l.~15 & \textit{for} curently \textit{read} currently \\
1275 & l.~22 & \textit{for} restritcion \textit{read} restriction \\
1276 & l.~21 & \textit{for} modst \textit{read} modest \\
1276 & l.~29 & \textit{for} Harthsorne (1962) \textit{read} Hartshorne (1962) \\
1276 & l.~31 & \textit{for} omniscinece \textit{read} omniscience \\
1278 & display (10.17) & \textit{for} $K_{y}p \leftrightarrow p$ \textit{read} $K_{y}p \to p$ \\
1278 & l.~11 & \textit{for} reecive \textit{read} receive \\
1278 & l.~20 & \textit{for} considerasion \textit{read} consideration \\
1278 & l.~23 & \textit{for} conditoin \textit{read} condition \\
1278 & l.~27 & \textit{for} $\mathrm{Os}(y) \wedge \mathrm{Os}(z) \to \exists p\,(K_{y}p \leftrightarrow K_{z}p)$ \textit{read} $\mathrm{Os}(y) \wedge \mathrm{Os}(z) \to \forall p\,(K_{y}p \leftrightarrow K_{z}p)$ \\
1279 & l.~27 & \textit{for} Lemma~1875 of Volume~X \textit{read} Lemma~1880 of Volume~X; the result cited is on p.~3739 \\
1279 & l.~7 from foot & \textit{for} disagrement \textit{read} disagreement \\
1280 & l.~14 & \textit{for} inteerst \textit{read} interest \\
1280 & display (10.18) & \textit{for} $\mathrm{Pos}({\phi}) \wedge \mathrm{Pos}({\chi}) \leftrightarrow \mathrm{Pos}({\phi} \sqcap {\chi})$ \textit{read} $\mathrm{Pos}({\phi}) \wedge \mathrm{Pos}({\chi}) \to \mathrm{Pos}({\phi} \sqcap {\chi})$ \\
1281 & l.~2 & \textit{for} carrefully \textit{read} carefully \\
1281 & l.~18 & \textit{for} variuos \textit{read} various \\
1281 & l.~39 & \textit{for} attributte \textit{read} attribute \\
1283 & l.~4 & \textit{for} dliemma \textit{read} dilemma \\
1283 & display (10.20) & \textit{for} $\nec(\mathrm{Op}(y) \vee \mathrm{Os}(y) \wedge \mathrm{MP}(y))$ \textit{read} $\nec(\mathrm{Op}(y) \wedge \mathrm{Os}(y) \wedge \mathrm{MP}(y))$ \\
1283 & l.~18 & \textit{for} Lemma~2359 of Volume~X \textit{read} Lemma~2353 of Volume~X; the result cited is on p.~3796 \\
1285 & display (10.23) & \textit{for} $K_{y}r \leftrightarrow K_{y}K_{y}r$ \textit{read} $K_{y}r \to K_{y}K_{y}r$ \\
1286 & l.~19 & \textit{for} comparisson \textit{read} comparison \\
1286 & l.~31 & \textit{for} existance \textit{read} existence \\
1287 & l.~8 & \textit{for} recgonize \textit{read} recognize \\
1288 & l.~34 & \textit{for} p.~1026 \textit{read} p.~1036 (Definition~6.6) \\
1288 & l.~37 & \textit{for} inclding \textit{read} including \\
1289 & l.~18 & \textit{for} subseuqent \textit{read} subsequent \\
1290 & l.~12 & \textit{for} Hartshrne (1962) \textit{read} Hartshorne (1962) \\
\dag\,1290 & proof of Lemma~10.10, l.~21 & \textit{for} Proposition~10.2 \textit{read} Proposition~10.5 \\
1290 & l.~40 & \textit{for} parragraph \textit{read} paragraph \\
1291 & l.~26 & \textit{for} godlike in Gödel's sense \textit{read} godlike in Anderson's sense \\
\dag\,1292 & proof of Proposition~10.12, l.~16 & \textit{for} Proposition~10.8 \textit{read} Proposition~10.5 \\
\dag\,1292 & proof of Proposition~10.12, (10.27) & \textit{for} $\mathrm{Pos}(\neg {\psi}) \leftrightarrow \mathrm{Pos}({\psi})$ \textit{read} $\mathrm{Pos}(\neg {\psi}) \leftrightarrow \neg\mathrm{Pos}({\psi})$ \\
1292 & l.~1 from foot & \textit{for} closed under disjunction \textit{read} closed under conjunction \\
1293 & l.~28 & \textit{for} neverthelss \textit{read} nevertheless \\
1294 & l.~5 & \textit{for} p.~946 \textit{read} p.~936 (Definition~4.14) \\
1294 & l.~14 & \textit{for} occaison \textit{read} occasion \\
1294 & display (10.29) & \textit{for} $\mathrm{MP}(x) \to \pos\,\mathrm{Wrong}(x,a)$ \textit{read} $\mathrm{MP}(x) \to \neg\pos\,\mathrm{Wrong}(x,a)$ \\
1295 & l.~28 & \textit{for} logiclaly \textit{read} logically \\
1295 & display (10.30) & \textit{for} $ Q(t,z) \wedge \mathrm{Task}(t) \to \mathrm{Can}(z,t)$ \textit{read} $\neg Q(t,z) \wedge \mathrm{Task}(t) \to \mathrm{Can}(z,t)$ \\
\dag\,1296 & proof of Theorem~10.14, l.~4 & \textit{for} Proposition~10.8 \textit{read} Proposition~10.9 \\
1296 & l.~15 & \textit{for} propossal \textit{read} proposal \\
\dag\,1296 & proof of Theorem~10.14, l.~21 & \textit{for} Lemma~1858 of Volume~X \textit{read} Lemma~1860 of Volume~X \\
1296 & l.~27 & \textit{for} tasks whose description quantifies over the agent's own powers \textit{read} tasks whose description does not quantify over the agent's own powers \\
1296 & l.~31 & \textit{for} hypotheiss \textit{read} hypothesis \\
1296 & l.~35 & \textit{for} exsitence \textit{read} existence \\
1297 & l.~12 & \textit{for} p.~1167 \textit{read} p.~1177 (Proposition~8.13) \\
1297 & l.~16 & \textit{for} innevitable \textit{read} inevitable \\
1297 & l.~30 & \textit{for} itnrinsic \textit{read} intrinsic \\
1298 & l.~8 & \textit{for} p.~1293 \textit{read} p.~1292 (Definition~10.13) \\
1298 & l.~35 & \textit{for} concrened \textit{read} concerned \\
1299 & l.~30 & \textit{for} critrion \textit{read} criterion \\
\end{errlist}

\subsection*{Chapter~11. Probabilistic assessments of the possibility premise \hfill \normalfont\small pp.~1302--1370}
\begin{errlist}
1303 & l.~4 & \textit{for} conclussion \textit{read} conclusion \\
1304 & l.~13 & \textit{for} p.~913 \textit{read} p.~903 (Definition~4.1) \\
1304 & l.~18 & \textit{for} p.~903 \textit{read} p.~913 (Lemma~4.5) \\
1304 & l.~32 & \textit{for} smymetry \textit{read} symmetry \\
1306 & l.~20 & \textit{for} p.~1115 \textit{read} p.~1151 (Corollary~8.5) \\
1306 & l.~21 & \textit{for} importacne \textit{read} importance \\
1307 & l.~28 & \textit{for} Lemma~2569 of Volume~X \textit{read} Lemma~2567 of Volume~X \\
1308 & display (11.2) & \textit{for} $P(\nec G) + P(\pos\neg G) \leq 1$ \textit{read} $P(\pos G) + P(\pos\neg G) \leq 1$ \\
1308 & l.~18 & \textit{for} contribtion \textit{read} contribution \\
1308 & l.~19 & \textit{for} simmilarly \textit{read} similarly \\
1308 & l.~5 from foot & \textit{for} dependnce \textit{read} dependence \\
1309 & l.~28 & \textit{for} absrtact \textit{read} abstract \\
1310 & l.~14 & \textit{for} crurently \textit{read} currently \\
1311 & l.~10 & \textit{for} Volume~X, Chapter~6 \textit{read} Volume~IX, Chapter~6 (whether a maximally great being must create the best) \\
1311 & l.~32 & \textit{for} sucessful \textit{read} successful \\
1311 & l.~6 from foot & \textit{for} hpyotheses \textit{read} hypotheses \\
1312 & l.~8 & \textit{for} Lemma~2583 of Volume~X \textit{read} Lemma~2589 of Volume~X; the result cited is on p.~3824 \\
1312 & l.~28 & \textit{for} p.~1280 \textit{read} p.~1270 (Proposition~10.8) \\
1313 & l.~23 & \textit{for} vraious \textit{read} various \\
1313 & l.~40 & \textit{for} Lemma~2287 of Volume~X \textit{read} Lemma~2285 of Volume~X; the result cited is on p.~3788 \\
1314 & l.~29 & \textit{for} respctively \textit{read} respectively \\
1314 & l.~31 & \textit{for} variasion \textit{read} variation \\
1314 & l.~5 from foot & \textit{for} conrtibution \textit{read} contribution \\
1315 & l.~11 & \textit{for} transfered \textit{read} transferred \\
1315 & l.~22 & \textit{for} incnsistent \textit{read} inconsistent \\
1315 & l.~24 & \textit{for} paralell \textit{read} parallel \\
1316 & l.~12 & \textit{for} acocunt \textit{read} account \\
1316 & l.~32 & \textit{for} Volume~VI, Chapter~2 \textit{read} Volume~VI, Chapter~4 (the evidential problem of evil) \\
1318 & l.~7 & \textit{for} dependnt \textit{read} dependent \\
1318 & l.~16 & \textit{for} cnotroversial \textit{read} controversial \\
1320 & l.~6 & \textit{for} mantion \textit{read} mention \\
1320 & l.~40 & \textit{for} suceed \textit{read} succeed \\
1321 & l.~15 & \textit{for} cnotext \textit{read} context \\
1321 & l.~30 & \textit{for} Lemma~1874 of Volume~X \textit{read} Lemma~1872 of Volume~X \\
\dag\,1322 & proof of Proposition~11.4, l.~4 & \textit{for} Proposition~11.3 \textit{read} Proposition~11.2; the step uses Proposition~11.2 (p.~1314) \\
1322 & l.~16 & \textit{for} developmant \textit{read} development \\
1322 & l.~17 & \textit{for} attirbute \textit{read} attribute \\
\dag\,1322 & proof of Proposition~11.4, l.~28 & \textit{for} Lemma~1951 of Volume~X \textit{read} Lemma~1946 of Volume~X; the result cited is on p.~3747 \\
1322 & l.~7 from foot & \textit{for} emphaiss \textit{read} emphasis \\
1323 & l.~16 & \textit{for} mension \textit{read} mention \\
1324 & l.~1 & \textit{for} aanlysis \textit{read} analysis \\
1325 & l.~12 & \textit{for} probabillity \textit{read} probability \\
1325 & l.~17 & \textit{for} unedrmine \textit{read} undermine \\
1325 & l.~26 & \textit{for} Lemma~2433 of Volume~X \textit{read} Lemma~2428 of Volume~X \\
1327 & l.~5 & \textit{for} relevnt \textit{read} relevant \\
1328 & l.~21 & \textit{for} Volume~V, Chapter~4 \textit{read} Volume~V, Chapter~5 (the Barcan formula) \\
1328 & l.~24 & \textit{for} lnaguage \textit{read} language \\
1328 & l.~6 from foot & \textit{for} admttedly \textit{read} admittedly \\
1329 & l.~6 & \textit{for} literrature \textit{read} literature \\
1329 & l.~8 & \textit{for} p.~1047 \textit{read} p.~1048 (Lemma~6.11) \\
\dag\,1329 & proof of Corollary~11.5, l.~13 & \textit{for} Proposition~11.2 \textit{read} Proposition~11.4 \\
1329 & l.~17 & \textit{for} spearate \textit{read} separate \\
1329 & l.~3 from foot & \textit{for} maitnain \textit{read} maintain \\
1330 & l.~39 & \textit{for} tenedncy \textit{read} tendency \\
1331 & l.~6 & \textit{for} Lemma~2333 of Volume~X \textit{read} Lemma~2329 of Volume~X \\
1331 & l.~7 & \textit{for} interset \textit{read} interest \\
1333 & l.~22 & \textit{for} intuitoin \textit{read} intuition \\
1335 & l.~17 & \textit{for} prmise \textit{read} premise \\
1335 & l.~3 from foot & \textit{for} attribtue \textit{read} attribute \\
\dag\,1336 & proof of Proposition~11.8, (11.7) & \textit{for} $\mathrm{Conc}(q) \leftrightarrow \pos q$ \textit{read} $\mathrm{Conc}(q) \to \pos q$ \\
1337 & l.~23 & \textit{delete} the second occurrence of \emph{of} \\
1338 & l.~17 & \textit{for} decsribe \textit{read} describe \\
1338 & l.~24 & \textit{for} bteween \textit{read} between \\
1341 & l.~31 & \textit{for} introudce \textit{read} introduce \\
1341 & l.~5 from foot & \textit{for} descritpion \textit{read} description \\
1342 & display (11.9) & \textit{for} $P(q \mid q) = P(r \mid q)\,P(q)/P(r)$ \textit{read} $P(q \mid r) = P(r \mid q)\,P(q)/P(r)$; the formula is used in its correct form on p.~1343 \\
\dag\,1343 & proof of Corollary~11.11, l.~26 & \textit{for} Lemma~2111 of Volume~X \textit{read} Lemma~2117 of Volume~X \\
1344 & l.~12 & \textit{for} scennario \textit{read} scenario \\
1344 & l.~22 & \textit{for} inlcuding \textit{read} including \\
1345 & l.~5 & \textit{for} Volume~IV, Chapter~9 \textit{read} Volume~IV, Chapter~7 (necessary existence and aseity) \\
1346 & l.~9 & \textit{for} accordngly \textit{read} accordingly \\
1347 & l.~15 & \textit{for} craefully \textit{read} carefully \\
1347 & l.~26 & \textit{for} specifc \textit{read} specific \\
1348 & l.~3 & \textit{for} transfrered \textit{read} transferred \\
1348 & l.~29 & \textit{for} strattegy \textit{read} strategy \\
\dag\,1349 & proof of Theorem~11.13, l.~25 & \textit{for} Proposition~11.3 \textit{read} Proposition~11.8 \\
\dag\,1349 & proof of Theorem~11.13, (11.13) & \textit{for} $\mathrm{Conc}(r) \leftrightarrow \pos r$ \textit{read} $\mathrm{Conc}(r) \to \pos r$; the formula is used in its correct form on p.~1353 \\
1350 & l.~23 & \textit{for} abmiguous \textit{read} ambiguous \\
1350 & l.~30 & \textit{for} expalnation \textit{read} explanation \\
\dag\,1351 & proof of Lemma~11.16, l.~10 & \textit{for} Proposition~11.4 \textit{read} Proposition~11.8; the step uses Proposition~11.8 (p.~1335) \\
\dag\,1351 & proof of Lemma~11.16, l.~14 & \textit{for} Lemma~2485 of Volume~X \textit{read} Lemma~2480 of Volume~X \\
1351 & l.~20 & \textit{for} propsect \textit{read} prospect \\
1352 & l.~11 & \textit{for} accordnig \textit{read} according \\
1352 & l.~26 & \textit{for} rseponse \textit{read} response \\
1352 & l.~28 & \textit{for} p.~1128 \textit{read} p.~1129 (Proposition~7.17) \\
1353 & l.~11 & \textit{for} regardles \textit{read} regardless \\
1353 & l.~12 & \textit{for} Lemma~2357 of Volume~X \textit{read} Lemma~2351 of Volume~X \\
\dag\,1353 & proof of Lemma~11.19, l.~18 & \textit{for} Corollary~11.9 \textit{read} Corollary~11.11 \\
1354 & l.~16 & \textit{for} remiander \textit{read} remainder \\
1354 & l.~26 & \textit{for} considration \textit{read} consideration \\
1355 & l.~3 from foot & \textit{for} controversail \textit{read} controversial \\
1356 & l.~20 & \textit{for} p.~1224 \textit{read} p.~1124 (Corollary~7.15) \\
1356 & l.~23 & \textit{for} propnent \textit{read} proponent \\
1356 & l.~25 & \textit{for} evenutally \textit{read} eventually \\
1356 & display (11.14) & \textit{for} $P(\pos\neg G \mid A) = P(\nec\neg G)$ \textit{read} $P(\pos\neg G \mid A) = P(\pos\neg G)$ \\
1357 & l.~9 & \textit{for} unform \textit{read} uniform \\
1357 & l.~19 & \textit{for} corectly \textit{read} correctly \\
1359 & l.~25 & \textit{for} especilaly \textit{read} especially \\
1360 & l.~28 & \textit{for} rseistance \textit{read} resistance \\
\dag\,1361 & proof of Lemma~11.20, l.~15 & \textit{for} Lemma~2193 of Volume~X \textit{read} Lemma~2197 of Volume~X \\
\dag\,1361 & proof of Lemma~11.20, l.~23 & \textit{for} Lemma~1857 of Volume~X \textit{read} Lemma~1853 of Volume~X \\
1361 & l.~24 & \textit{for} acount \textit{read} account \\
\dag\,1362 & proof of Lemma~11.20, l.~2 & \textit{for} Definition~11.15 \textit{read} Definition~11.17 \\
\dag\,1362 & proof of Corollary~11.21, l.~11 & \textit{for} Lemma~11.20 \textit{read} Lemma~11.19; the step uses Lemma~11.19 (p.~1353) \\
\dag\,1362 & proof of Corollary~11.21, l.~25 & \textit{for} Corollary~11.5 \textit{read} Corollary~11.9 \\
1364 & l.~14 & \textit{for} remainedr \textit{read} remainder \\
1364 & l.~25 & \textit{for} funadmental \textit{read} fundamental \\
1365 & l.~31 & \textit{for} Lemma~1890 of Volume~X \textit{read} Lemma~1886 of Volume~X \\
1366 & l.~9 & \textit{for} hithetro \textit{read} hitherto \\
1366 & l.~32 & \textit{for} proviedd \textit{read} provided \\
1367 & l.~8 & \textit{for} sittuation \textit{read} situation \\
1367 & l.~19 & \textit{for} Lemma~2580 of Volume~X \textit{read} Lemma~2585 of Volume~X \\
1367 & l.~23 & \textit{for} specfiic \textit{read} specific \\
1367 & l.~26 & \textit{for} prcoedure \textit{read} procedure \\
1368 & l.~21 & \textit{for} p.~1330 \textit{read} p.~1340 (Definition~11.10) \\
1368 & l.~8 from foot & \textit{for} decsription \textit{read} description \\
1369 & l.~4 & \textit{for} Lemma~2542 of Volume~X \textit{read} Lemma~2543 of Volume~X \\
1369 & l.~8 from foot & \textit{for} Volume~I, Chapter~2 \textit{read} Volume~I, Chapter~4 (what the formalization establishes) \\
1370 & l.~9 & \textit{for} possssion \textit{read} possession \\
1370 & l.~15 & \textit{for} Lemma~1981 of Volume~X \textit{read} Lemma~1978 of Volume~X \\
1370 & l.~16 & \textit{for} throuhgout \textit{read} throughout \\
\end{errlist}

\subsection*{Chapter~12. Conclusion of Volume III \hfill \normalfont\small pp.~1371--1401}
\begin{errlist}
1372 & l.~6 & \textit{for} carefluly \textit{read} carefully \\
1372 & l.~19 & \textit{for} meanning \textit{read} meaning \\
1372 & l.~23 & \textit{for} philosophcial \textit{read} philosophical \\
1372 & l.~4 from foot & \textit{for} Lemma~1983 of Volume~X \textit{read} Lemma~1980 of Volume~X; the result cited is on p.~3751 \\
1373 & l.~9 & \textit{for} neevrtheless \textit{read} nevertheless \\
1373 & l.~8 from foot & \textit{for} p.~1098 \textit{read} p.~1088 (Lemma~7.4) \\
1375 & l.~3 from foot & \textit{for} Volume~IV, Chapter~2 \textit{read} Volume~IV, Chapter~3 (omnipotence) \\
1377 & l.~22 & \textit{for} clasical \textit{read} classical \\
1377 & l.~31 & \textit{for} immeidately \textit{read} immediately \\
1377 & l.~38 & \textit{for} Lemma~2253 of Volume~X \textit{read} Lemma~2256 of Volume~X \\
1378 & l.~22 & \textit{for} Lemma~2333 of Volume~X \textit{read} Lemma~2336 of Volume~X \\
1379 & l.~31 & \textit{for} Lemma~2424 of Volume~X \textit{read} Lemma~2429 of Volume~X; the result cited is on p.~3805 \\
1380 & l.~38 & \textit{for} Volume~III, Chapter~3 \textit{read} Volume~III, Chapter~5 (the transfer from coherence to possibility) \\
1382 & l.~12 & \textit{for} ratinoal \textit{read} rational \\
1382 & l.~7 from foot & \textit{for} Volume~V, Chapter~5 \textit{read} Volume~V, Chapter~7 (modal collapse) \\
1382 & l.~36 & \textit{for} corrcetly \textit{read} correctly \\
1383 & l.~32 & \textit{for} vresion \textit{read} version \\
1384 & l.~4 & \textit{for} Lemma~2519 of Volume~X \textit{read} Lemma~2521 of Volume~X; the result cited is on p.~3816 \\
1384 & l.~5 from foot & \textit{for} beleive \textit{read} believe \\
1385 & l.~26 & \textit{for} althoguh \textit{read} although \\
1385 & l.~29 & \textit{for} betwen \textit{read} between \\
1385 & l.~30 & \textit{for} Lemma~2271 of Volume~X \textit{read} Lemma~2265 of Volume~X \\
1389 & l.~15 & \textit{for} paraidgm \textit{read} paradigm \\
1389 & l.~31 & \textit{for} appropraite \textit{read} appropriate \\
1391 & l.~2 & \textit{for} Lemma~2303 of Volume~X \textit{read} Lemma~2308 of Volume~X \\
\dag\,1391 & proof of Proposition~12.4, l.~10 & \textit{for} Lemma~2539 of Volume~X \textit{read} Lemma~2541 of Volume~X \\
1392 & l.~1 & \textit{for} Lemma~2419 of Volume~X \textit{read} Lemma~2420 of Volume~X; the result cited is on p.~3804 \\
1392 & l.~9 & \textit{for} tradision \textit{read} tradition \\
1393 & l.~13 & \textit{for} deppendent \textit{read} dependent \\
1393 & l.~18 & \textit{for} whetehr \textit{read} whether \\
1394 & l.~13 & \textit{for} p.~975 \textit{read} p.~965 (Proposition~5.3) \\
1394 & l.~18 & \textit{for} Volume~III, Chapter~10 \textit{read} Volume~III, Chapter~9 (the ancestral premise) \\
1394 & l.~20 & \textit{for} sattement \textit{read} statement \\
1394 & l.~31 & \textit{for} Volume~V, Chapter~9 \textit{read} Volume~V, Chapter~7 (modal collapse) \\
1394 & l.~6 from foot & \textit{for} committd \textit{read} committed \\
1395 & l.~1 & \textit{for} Lemma~2575 of Volume~X \textit{read} Lemma~2571 of Volume~X \\
1395 & l.~24 & \textit{for} cnocerned \textit{read} concerned \\
1395 & l.~2 from foot & \textit{for} knwledge \textit{read} knowledge \\
1396 & l.~6 & \textit{for} differnt \textit{read} different \\
1396 & l.~28 & \textit{for} Lemma~2462 of Volume~X \textit{read} Lemma~2461 of Volume~X \\
1397 & l.~6 & \textit{for} vloume \textit{read} volume \\
1397 & l.~23 & \textit{for} judgment \textit{read} judgement \\
1397 & l.~32 & \textit{for} proposision \textit{read} proposition \\
1397 & l.~7 from foot & \textit{for} idneed \textit{read} indeed \\
1398 & l.~6 & \textit{for} cretainly \textit{read} certainly \\
1398 & l.~25 & \textit{for} accessble \textit{read} accessible \\
1398 & l.~32 & \textit{for} silghtly \textit{read} slightly \\
1400 & l.~9 & \textit{for} expresion \textit{read} expression \\
1400 & l.~13 & \textit{for} Volume~VIII, Chapter~6 \textit{read} Volume~VIII, Chapter~5 (divine simplicity) \\
1400 & l.~31 & \textit{for} conisstent \textit{read} consistent \\
1400 & l.~37 & \textit{for} thoruogh \textit{read} thorough \\
1401 & l.~26 & \textit{for} Volume~IV, Chapter~4 \textit{read} Volume~IV, Chapter~5 (moral perfection) \\
1401 & l.~31 & \textit{for} dfeended \textit{read} defended \\
1401 & l.~2 from foot & \textit{for} Lemma~2301 of Volume~X \textit{read} Lemma~2297 of Volume~X; the result cited is on p.~3789 \\
\end{errlist}


\section{Volume IV: The Divine Attributes}\label{s:vol4}
% generated by generate.py -- do not edit
\subsection*{Front matter \hfill \normalfont\small pp.~1402--1404}
\begin{errlist}
1403 & Contents, Chapter~9 & \textit{for} 1719 \textit{read} 1729 \\
\end{errlist}

\subsection*{Chapter~1. Introduction \hfill \normalfont\small pp.~1405--1411}
\begin{errlist}
1406 & l.~16 & \textit{for} entriely \textit{read} entirely \\
1406 & l.~31 & \textit{for} intrnisic \textit{read} intrinsic \\
1407 & l.~6 from foot & \textit{for} Wierrenga (1989) \textit{read} Wierenga (1989) \\
1408 & l.~8 & \textit{for} theologicl \textit{read} theological \\
1408 & l.~11 & \textit{for} Volume~IX, Chapter~9 \textit{read} Volume~IX, Chapter~7 (foreknowledge) \\
1408 & n.~18 & \textit{for} Wiernga (1989) \textit{read} Wierenga (1989) \\
1409 & l.~14 & \textit{for} Volume~VII, Chapter~2 \textit{read} Volume~VII, Chapter~4 (fine-tuning) \\
1410 & l.~9 & \textit{for} purposse \textit{read} purpose \\
1411 & l.~19 & \textit{for} idenitcal \textit{read} identical \\
1411 & l.~23 & \textit{for} perspectvie \textit{read} perspective \\
1411 & l.~30 & \textit{for} availbale \textit{read} available \\
1411 & n.~14 & \textit{for} Wieernga (1989) \textit{read} Wierenga (1989) \\
\end{errlist}

\subsection*{Chapter~2. Maximal excellence and its components \hfill \normalfont\small pp.~1412--1446}
\begin{errlist}
1414 & running heads, pp.~1414--1446 (even) & \textit{for} Omnipotence \textit{read} Maximal excellence \\
1414 & l.~5 & \textit{for} $K_{y}r \leftrightarrow r$ \textit{read} $K_{y}r \to r$ \\
1415 & l.~1 & \textit{for} merley \textit{read} merely \\
1415 & l.~16 & \textit{for} tsak \textit{read} task \\
1415 & l.~1 from foot & \textit{for} knows every false proposition \textit{read} knows no false proposition \\
1416 & display (2.2) & \textit{for} $\pos\exists z\,\mathrm{MG}(z) \leftrightarrow \nec\exists z\,\mathrm{MG}(z)$ \textit{read} $\pos\exists z\,\mathrm{MG}(z) \to \nec\exists z\,\mathrm{MG}(z)$ \\
1416 & display (2.3) & \textit{for} $\mathrm{MG}(x,w_0) \iff \forall v'\,(w_0Rv' \leftrightarrow E(x,v') \wedge \mathrm{ME}(x,v'))$ \textit{read} $\mathrm{MG}(x,w_0) \iff \forall v'\,(w_0Rv' \to E(x,v') \wedge \mathrm{ME}(x,v'))$ \\
1418 & l.~1 & \textit{for} surley \textit{read} surely \\
1418 & l.~2 & \textit{for} simplcity \textit{read} simplicity \\
1418 & l.~40 & \textit{for} theoloigcal \textit{read} theological \\
\dag\,1420 & proof of Lemma~2.4, (2.6) & \textit{for} $\mathrm{MG}(y,w') \iff \forall w_1\,(w_1Rw' \to E(y,w_1) \wedge \mathrm{ME}(y,w_1))$ \textit{read} $\mathrm{MG}(y,w') \iff \forall w_1\,(w'Rw_1 \to E(y,w_1) \wedge \mathrm{ME}(y,w_1))$ \\
1420 & l.~34 & \textit{for} thereforre \textit{read} therefore \\
1421 & l.~3 & \textit{for} conitngent \textit{read} contingent \\
1421 & l.~10 & \textit{delete} the second occurrence of \emph{that} \\
1422 & l.~3 & \textit{after} every task \textit{insert} whose description does not quantify over the agent's own powers \\
1423 & l.~6 & \textit{for} Lemma~2821 of Volume~X \textit{read} Lemma~2824 of Volume~X \\
1423 & l.~25 & \textit{for} $\mathrm{Can}(z,t) \to \nec\,\mathrm{Does}(z,t)$ \textit{read} $\mathrm{Can}(z,t) \to \pos\,\mathrm{Does}(z,t)$ \\
1424 & l.~4 & \textit{for} $K_{z}r \leftrightarrow K_{z}K_{z}r$ \textit{read} $K_{z}r \to K_{z}K_{z}r$ \\
1424 & l.~7 & \textit{for} proposiiton \textit{read} proposition \\
1424 & l.~24 & \textit{for} prspect \textit{read} prospect \\
1424 & l.~5 from foot & \textit{for} prefect \textit{read} perfect \\
1424 & l.~39 & \textit{for} similalry \textit{read} similarly \\
\dag\,1425 & proof of Proposition~2.6, l.~25 & \textit{for} Proposition~2.5 \textit{read} Proposition~2.2 \\
1426 & l.~22 & \textit{for} formluation \textit{read} formulation \\
\dag\,1426 & proof of Proposition~2.5, l.~28 & \textit{for} Lemma~2.8 \textit{read} Lemma~2.4 \\
1427 & l.~4 & \textit{for} abslutely \textit{read} absolutely \\
1428 & l.~1 from foot & \textit{delete} the second occurrence of \emph{that} \\
1429 & l.~10 & \textit{for} maximl \textit{read} maximal \\
1429 & l.~24 & \textit{for} p.~1417 \textit{read} p.~1407 (Proposition~1.2) \\
1429 & l.~27 & \textit{for} p.~1407 \textit{read} p.~1406 (Definition~1.1) \\
1429 & l.~32 & \textit{for} invriant \textit{read} invariant \\
1429 & l.~33 & \textit{for} indepednence \textit{read} independence \\
1430 & l.~13 & \textit{for} conscoius \textit{read} conscious \\
\dag\,1430 & proof of Corollary~2.9, l.~24 & \textit{for} Proposition~2.5 \textit{read} Proposition~2.6 \\
1430 & l.~8 from foot & \textit{for} Lemma~2974 of Volume~X \textit{read} Lemma~2976 of Volume~X \\
1430 & n.~22 & \textit{for} Hoffman and Rosenkrantz (2001) \textit{read} Hoffman and Rosenkrantz (2002) \\
1431 & l.~4 & \textit{for} requiremment \textit{read} requirement \\
1432 & l.~23 & \textit{for} $\mathrm{MP}(z) \to \neg\nec\,\mathrm{Wrong}(z,a)$ \textit{read} $\mathrm{MP}(z) \to \neg\pos\,\mathrm{Wrong}(z,a)$ \\
1432 & l.~30 & \textit{for} howveer \textit{read} however \\
1433 & l.~36 & \textit{for} prefection \textit{read} perfection \\
\dag\,1435 & proof of Lemma~2.10, l.~28 & \textit{for} Definition~2.3 \textit{read} Definition~2.1; the step uses Definition~2.1 (p.~1417) \\
\dag\,1435 & proof of Lemma~2.10, (2.11) & \textit{for} $\mathrm{Op}(z) \leftrightarrow \forall t\,(\mathrm{Task}(t) \wedge  Q(t,z) \to \mathrm{Can}(z,t))$ \textit{read} $\mathrm{Op}(z) \leftrightarrow \forall t\,(\mathrm{Task}(t) \wedge \neg Q(t,z) \to \mathrm{Can}(z,t))$ \\
1436 & l.~3 & \textit{for} p.~1428 \textit{read} p.~1427 (Lemma~2.8) \\
\dag\,1436 & proof of Lemma~2.10, (2.12) & \textit{for} $\mathrm{MG}(z,w_1) \to \mathrm{NE}(z,w_1) \vee \mathrm{NME}(z,w_1)$ \textit{read} $\mathrm{MG}(z,w_1) \to \mathrm{NE}(z,w_1) \wedge \mathrm{NME}(z,w_1)$; correctly printed in (2.18), p.~1441 \\
1436 & l.~38 & \textit{after} every true proposition \textit{insert} (plurally quantified) \\
1437 & l.~6 & \textit{for} crorectly \textit{read} correctly \\
1437 & l.~33 & \textit{for} Lemma~2957 of Volume~X \textit{read} Lemma~2958 of Volume~X; the result cited is on p.~3865 \\
1438 & l.~38 & \textit{for} philosopihcal \textit{read} philosophical \\
1439 & l.~18 & \textit{for} fammiliar \textit{read} familiar \\
\dag\,1439 & proof of Lemma~2.11, (2.13) & \textit{for} $\mathrm{Can}(y,t) \leftrightarrow \pos\,\mathrm{Does}(y,t)$ \textit{read} $\mathrm{Can}(y,t) \to \pos\,\mathrm{Does}(y,t)$ \\
1440 & l.~4 & \textit{after} maximally excellent \textit{insert} at $v'$ \\
1440 & display (2.15) & \textit{for} $\nec(\mathrm{Op}(x) \vee \mathrm{Os}(x) \wedge \mathrm{MP}(x))$ \textit{read} $\nec(\mathrm{Op}(x) \wedge \mathrm{Os}(x) \wedge \mathrm{MP}(x))$; correctly printed in (2.19), p.~1441 \\
1440 & l.~26 & \textit{for} deifnition \textit{read} definition \\
1441 & display (2.16) & \textit{for} $\pos(\mathrm{Op}(z) \wedge \mathrm{Os}(z) \wedge \mathrm{MP}(z))$ \textit{read} $\nec(\mathrm{Op}(z) \wedge \mathrm{Os}(z) \wedge \mathrm{MP}(z))$ \\
1442 & l.~26 & \textit{for} Hoffman and Rosenkrantz (2003) \textit{read} Hoffman and Rosenkrantz (2002) \\
1442 & l.~7 from foot & \textit{delete} the second occurrence of \emph{that} \\
1442 & l.~3 from foot & \textit{for} exists and is maximally excellent at some world accessible from $u$ \textit{read} exists and is maximally excellent at every world accessible from $u$ \\
1443 & l.~5 from foot & \textit{for} santence \textit{read} sentence \\
1444 & l.~1 & \textit{for} assumtpion \textit{read} assumption \\
1444 & l.~6 from foot & \textit{for} Lemma~2751 of Volume~X \textit{read} Lemma~2749 of Volume~X; the result cited is on p.~3843 \\
1446 & l.~10 & \textit{for} availble \textit{read} available \\
1446 & l.~14 & \textit{for} p.~1448 \textit{read} p.~1438 (Lemma~2.11) \\
\end{errlist}

\subsection*{Chapter~3. Omnipotence \hfill \normalfont\small pp.~1447--1495}
\begin{errlist}
1448 & l.~31 & \textit{for} accommodtae \textit{read} accommodate \\
1448 & l.~32 & \textit{delete} the second occurrence of \emph{is} \\
1448 & n.~18 & \textit{for} Frankfurt (1965) \textit{read} Frankfurt (1964) \\
1449 & l.~15 & \textit{for} strcitly \textit{read} strictly \\
1450 & l.~11 & \textit{for} equaly \textit{read} equally \\
1450 & l.~6 from foot & \textit{for} immedaite \textit{read} immediate \\
1451 & l.~11 & \textit{for} $\mathrm{Os}(x) \wedge \mathrm{Os}(x) \to \forall p\,(K_{x}p \leftrightarrow K_{z}p)$ \textit{read} $\mathrm{Os}(x) \wedge \mathrm{Os}(z) \to \forall p\,(K_{x}p \leftrightarrow K_{z}p)$ \\
1451 & l.~20 & \textit{delete} the second occurrence of \emph{to} \\
1452 & Definition~3.2 & \textit{see} Corrigendum~S4 \\
1453 & l.~20 & \textit{for} alternattive \textit{read} alternative \\
1453 & l.~31 & \textit{for} compariosn \textit{read} comparison \\
1455 & n.~5 & \textit{for} \emph{Summa Theologiae} I, q.~25, a.~4 (whether God is omnipotent) \textit{read} \emph{Summa Theologiae} I, q.~25, a.~3 (whether God is omnipotent) \\
1456 & l.~15 & \textit{for} \emph{Summa Theologiae} I, q.~3, a.~2 (the five ways) \textit{read} \emph{Summa Theologiae} I, q.~2, a.~3 (the five ways) \\
1456 & l.~33 & \textit{for} proopsition \textit{read} proposition \\
1457 & l.~9 & \textit{for} $\mathrm{Op}(z) \leftrightarrow \forall t\,(\mathrm{Task}(t) \vee \neg Q(t,z) \to \mathrm{Can}(z,t))$ \textit{read} $\mathrm{Op}(z) \leftrightarrow \forall t\,(\mathrm{Task}(t) \wedge \neg Q(t,z) \to \mathrm{Can}(z,t))$ \\
1457 & l.~27 & \textit{for} hoewver \textit{read} however \\
1457 & l.~30 & \textit{for} corerctly \textit{read} correctly \\
1458 & l.~15 & \textit{for} wheneevr \textit{read} whenever \\
1460 & l.~39 & \textit{delete} the second occurrence of \emph{of} \\
1462 & display (3.4) & \textit{for} $\mathrm{NE}(y) \leftrightarrow \neg\exists x\,D(y,x)$ \textit{read} $\mathrm{NE}(y) \to \neg\exists x\,D(y,x)$; correctly printed in (3.6), p.~1466 \\
1463 & l.~6 & \textit{for} Lemma~2709 of Volume~X \textit{read} Lemma~2715 of Volume~X; the result cited is on p.~3840 \\
1463 & l.~38 & \textit{for} $K_{y}p \leftrightarrow K_{y}K_{y}p$ \textit{read} $K_{y}p \to K_{y}K_{y}p$ \\
1464 & l.~10 & \textit{for} conculsion \textit{read} conclusion \\
1464 & l.~12 & \textit{for} inevitalbe \textit{read} inevitable \\
1464 & l.~15 & \textit{for} indepenednce \textit{read} independence \\
1465 & display (3.5) & \textit{for} $\mathrm{Can}(x,t) \to \nec\,\mathrm{Does}(x,t)$ \textit{read} $\mathrm{Can}(x,t) \to \pos\,\mathrm{Does}(x,t)$ \\
1465 & n.~19 & \textit{for} Mavrodes (1964) \textit{read} Mavrodes (1963) \\
\dag\,1466 & proof of Proposition~3.6, l.~11 & \textit{for} Proposition~3.5 \textit{read} Proposition~3.3 \\
\dag\,1466 & proof of Proposition~3.5, l.~15 & \textit{for} Proposition~3.6 \textit{read} Proposition~3.3; the step uses Proposition~3.3 (p.~1457) \\
1467 & l.~1 & \textit{for} evaulate \textit{read} evaluate \\
\dag\,1467 & proof of Proposition~3.6, l.~14 & \textit{for} Lemma~2656 of Volume~X \textit{read} Lemma~2659 of Volume~X \\
1467 & l.~16 & \textit{for} challegne \textit{read} challenge \\
\dag\,1467 & proof of Proposition~3.5, l.~17 & \textit{for} Definition~3.1 \textit{read} Definition~3.2 \\
1467 & l.~32 & \textit{for} imposible \textit{read} impossible \\
1468 & l.~13 & \textit{for} Volume~I, Chapter~6 \textit{read} Volume~I, Chapter~7 (the choice of S5) \\
1472 & l.~3 & \textit{for} depednence \textit{read} dependence \\
\dag\,1472 & proof of Theorem~3.7, l.~29 & \textit{for} Definition~3.4 \textit{read} Definition~3.2 \\
1473 & l.~8 & \textit{for} preliminnary \textit{read} preliminary \\
1473 & l.~20 & \textit{for} whethr \textit{read} whether \\
1474 & l.~5 & \textit{for} aprpoach \textit{read} approach \\
1475 & l.~28 & \textit{for} Lemma~2727 of Volume~X \textit{read} Lemma~2726 of Volume~X \\
1475 & l.~29 & \textit{delete} the second occurrence of \emph{is} \\
1475 & l.~32 & \textit{for} p.~1425 \textit{read} p.~1452 (Definition~3.2) \\
1476 & l.~3 & \textit{for} circumstacne \textit{read} circumstance \\
1476 & display (3.8) & \textit{for} $\nec(\mathrm{Op}(y) \wedge \mathrm{Os}(y) \vee \mathrm{MP}(y))$ \textit{read} $\nec(\mathrm{Op}(y) \wedge \mathrm{Os}(y) \wedge \mathrm{MP}(y))$ \\
1476 & l.~16 & \textit{for} maetrial \textit{read} material \\
1479 & l.~18 & \textit{for} posesssion \textit{read} possession \\
1480 & l.~18 & \textit{for} indepndent \textit{read} independent \\
1480 & l.~24 & \textit{for} condtiion \textit{read} condition \\
1480 & l.~30 & \textit{for} expilcit \textit{read} explicit \\
1480 & l.~40 & \textit{for} matterial \textit{read} material \\
1481 & l.~8 from foot & \textit{for} manifset \textit{read} manifest \\
1482 & l.~13 & \textit{for} Mvarodes (1963) \textit{read} Mavrodes (1963) \\
1482 & l.~24 & \textit{for} contnet \textit{read} content \\
1482 & l.~30 & \textit{for} Hoffman and Rosenkrantz (2002) \textit{read} Savage (1967) (the paradox of the stone) \\
\dag\,1483 & proof of Lemma~3.11, l.~16 & \textit{for} Proposition~3.15 \textit{read} Proposition~3.10; the step uses Proposition~3.10 (p.~1479) \\
1483 & l.~19 & \textit{for} betewen \textit{read} between \\
1484 & l.~9 & \textit{for} unnceessary \textit{read} unnecessary \\
1484 & l.~15 & \textit{for} atribute \textit{read} attribute \\
1484 & l.~28 & \textit{for} clraify \textit{read} clarify \\
1486 & l.~7 & \textit{for} Lemma~2841 of Volume~X \textit{read} Lemma~2843 of Volume~X; the result cited is on p.~3853 \\
1486 & l.~11 & \textit{for} compossbility \textit{read} compossibility \\
\dag\,1486 & proof of Theorem~3.13, l.~20 & \textit{for} Proposition~3.5 \textit{read} Proposition~3.6 \\
\dag\,1486 & proof of Theorem~3.13, l.~30 & \textit{for} Lemma~2936 of Volume~X \textit{read} Lemma~2935 of Volume~X \\
1487 & l.~7 & \textit{for} Lemma~2787 of Volume~X \textit{read} Lemma~2789 of Volume~X \\
1487 & l.~12 & \textit{for} cosntraint \textit{read} constraint \\
1487 & l.~13 & \textit{for} varoius \textit{read} various \\
1487 & l.~21 & \textit{for} challnge \textit{read} challenge \\
1487 & display (3.10) & \textit{for} $\mathrm{Os}(y) \wedge \mathrm{Os}(z) \to \exists q\,(K_{y}q \leftrightarrow K_{z}q)$ \textit{read} $\mathrm{Os}(y) \wedge \mathrm{Os}(z) \to \forall q\,(K_{y}q \leftrightarrow K_{z}q)$; the formula is used in its correct form on p.~1488 \\
1488 & l.~16 & \textit{for} sattisfied \textit{read} satisfied \\
1488 & l.~25 & \textit{for} prncipal \textit{read} principal \\
1489 & l.~7 & \textit{for} contngent \textit{read} contingent \\
1489 & l.~21 & \textit{for} identcal \textit{read} identical \\
\dag\,1489 & proof of Proposition~3.14, l.~7 from foot & \textit{for} Lemma~2871 of Volume~X \textit{read} Lemma~2874 of Volume~X; the result cited is on p.~3856 \\
1490 & l.~8 from foot & \textit{for} evntually \textit{read} eventually \\
1491 & l.~2 from foot & \textit{delete} the second occurrence of \emph{of} \\
1492 & l.~10 & \textit{for} connetcion \textit{read} connection \\
1492 & l.~16 & \textit{for} extnsion \textit{read} extension \\
1492 & l.~1 from foot & \textit{for} strcture \textit{read} structure \\
1493 & l.~4 & \textit{for} spupose \textit{read} suppose \\
\dag\,1493 & proof of Proposition~3.15, l.~19 & \textit{for} Lemma~3.12 \textit{read} Lemma~3.11; the step uses Lemma~3.11 (p.~1483) \\
1493 & l.~30 & \textit{for} cosntruction \textit{read} construction \\
1493 & n.~15 & \textit{for} Savgae (1967) \textit{read} Savage (1967) \\
1494 & l.~1 & \textit{for} prelimnary \textit{read} preliminary \\
1494 & l.~3 & \textit{for} p.~1583 \textit{read} p.~1483 (Lemma~3.11) \\
\dag\,1494 & proof of Proposition~3.15, (3.12) & \textit{for} $\mathrm{Op}(x) \leftrightarrow \exists t\,(\mathrm{Task}(t) \wedge \neg Q(t,x) \to \mathrm{Can}(x,t))$ \textit{read} $\mathrm{Op}(x) \leftrightarrow \forall t\,(\mathrm{Task}(t) \wedge \neg Q(t,x) \to \mathrm{Can}(x,t))$ \\
1495 & l.~5 & \textit{for} Mavrodes (1963) \textit{read} Frankfurt (1964) (the paradox of the stone) \\
\end{errlist}

\subsection*{Chapter~4. Omniscience \hfill \normalfont\small pp.~1496--1540}
\begin{errlist}
1497 & display (4.1) & \textit{for} $\mathrm{Op}(z) \leftrightarrow \forall t\,(\mathrm{Task}(t) \wedge \neg Q(t,z) \leftrightarrow \mathrm{Can}(z,t))$ \textit{read} $\mathrm{Op}(z) \leftrightarrow \forall t\,(\mathrm{Task}(t) \wedge \neg Q(t,z) \to \mathrm{Can}(z,t))$; the formula is used in its correct form on p.~1499 \\
1497 & l.~36 & \textit{for} $\mathrm{Os}(y) \leftrightarrow \forall r\,(r \leftrightarrow K_{y}r)$ \textit{read} $\mathrm{Os}(y) \leftrightarrow \forall r\,(r \to K_{y}r)$ \\
1498 & l.~6 & \textit{for} dependecne \textit{read} dependence \\
1498 & l.~28 & \textit{for} Definition~2.4 of Volume~I \textit{read} Definition~2.2 of Volume~I \\
1499 & Definition~4.1 & \textit{see} Corrigendum~S5 \\
1500 & l.~21 & \textit{for} Swinbrne (1993) \textit{read} Swinburne (1993) \\
1500 & l.~25 & \textit{for} excpetion \textit{read} exception \\
1500 & l.~6 from foot & \textit{for} euqally \textit{read} equally \\
1501 & l.~20 & \textit{for} hyoptheses \textit{read} hypotheses \\
1501 & l.~22 & \textit{for} seirous \textit{read} serious \\
1503 & l.~15 & \textit{for} advanatge \textit{read} advantage \\
1503 & l.~20 & \textit{for} unncessary \textit{read} unnecessary \\
1503 & l.~31 & \textit{for} carefuly \textit{read} carefully \\
1504 & l.~10 & \textit{for} accommdoate \textit{read} accommodate \\
1504 & l.~16 & \textit{for} Kretzmann (1965) \textit{read} Kretzmann (1966) \\
1506 & l.~8 & \textit{for} p.~1407 \textit{read} p.~1417 (Definition~2.1) \\
1506 & l.~32 & \textit{for} guaranete \textit{read} guarantee \\
1507 & l.~23 & \textit{for} impilcit \textit{read} implicit \\
1508 & display (4.5) & \textit{for} $\neg Q(t,y) \vee \mathrm{Task}(t) \to \mathrm{Can}(y,t)$ \textit{read} $\neg Q(t,y) \wedge \mathrm{Task}(t) \to \mathrm{Can}(y,t)$ \\
1510 & l.~29 & \textit{for} frotiori \textit{read} fortiori \\
1511 & l.~8 & \textit{delete} the second occurrence of \emph{to} \\
1511 & l.~11 & \textit{for} idnicate \textit{read} indicate \\
1511 & l.~25 & \textit{for} p.~1415 \textit{read} p.~1425 (Proposition~2.5) \\
1511 & l.~29 & \textit{for} explanaiton \textit{read} explanation \\
1512 & display (4.6) & \textit{for} $ K_{z}r \to K_{z}\neg K_{z}r$ \textit{read} $\neg K_{z}r \to K_{z}\neg K_{z}r$ \\
1513 & display (4.7) & \textit{for} $\mathrm{Os}(z) \to \forall r\,(r \to K_{z}r)$ \textit{read} $\mathrm{Os}(z) \leftrightarrow \forall r\,(r \to K_{z}r)$; the formula is used in its correct form on p.~1516 \\
1513 & l.~21 & \textit{for} Lemma~2703 of Volume~X \textit{read} Lemma~2707 of Volume~X; the result cited is on p.~3839 \\
1513 & l.~32 & \textit{for} cotnent \textit{read} content \\
1514 & l.~4 & \textit{for} develpoment \textit{read} development \\
1514 & l.~11 & \textit{for} Lemma~2740 of Volume~X \textit{read} Lemma~2736 of Volume~X \\
1514 & l.~12 & \textit{for} corretcly \textit{read} correctly \\
1514 & l.~30 & \textit{delete} the second occurrence of \emph{a} \\
1515 & l.~2 & \textit{for} Swniburne (1993) \textit{read} Swinburne (1993) \\
1515 & l.~29 & \textit{for} illustarte \textit{read} illustrate \\
1516 & l.~12 & \textit{for} throough \textit{read} thorough \\
1516 & l.~4 from foot & \textit{for} Theorem~5.6 of Volume~III \textit{read} Theorem~5.4 of Volume~III; the result cited is on p.~965 \\
1518 & l.~31 & \textit{for} claimmed \textit{read} claimed \\
1519 & l.~23 & \textit{for} acnkowledge \textit{read} acknowledge \\
1520 & l.~1 & \textit{for} denail \textit{read} denial \\
\dag\,1520 & proof of Proposition~4.8, l.~29 & \textit{for} Definition~4.1 \textit{read} Definition~4.4 \\
\dag\,1521 & proof of Lemma~4.9, (4.8) & \textit{for} $\mathrm{NE}(y) \to \exists x\,D(y,x)$ \textit{read} $\mathrm{NE}(y) \to \neg\exists x\,D(y,x)$; the formula is used in its correct form on p.~1522 \\
1521 & l.~29 & \textit{for} deevlop \textit{read} develop \\
1521 & l.~5 from foot & \textit{for} parrallel \textit{read} parallel \\
1522 & l.~7 & \textit{for} dfeinition \textit{read} definition \\
1522 & l.~15 & \textit{for} p.~1467 \textit{read} p.~1466 (Proposition~3.6) \\
1522 & l.~6 from foot & \textit{for} Swinburne (1994) \textit{read} Swinburne (1993) \\
1522 & l.~1 from foot & \textit{for} Lemma~2890 of Volume~X \textit{read} Lemma~2891 of Volume~X \\
1523 & l.~3 & \textit{for} $\mathrm{NE}(z) \to \neg\forall x\,D(z,x)$ \textit{read} $\mathrm{NE}(z) \to \neg\exists x\,D(z,x)$ \\
1523 & l.~8 & \textit{for} aesity \textit{read} aseity \\
1523 & l.~24 & \textit{for} Grim (1983) \textit{read} Kretzmann (1966) (omniscience and change) \\
1524 & l.~11 & \textit{for} cnotrast \textit{read} contrast \\
1524 & l.~14 & \textit{for} criticl \textit{read} critical \\
1524 & l.~6 from foot & \textit{for} resmblance \textit{read} resemblance \\
1526 & l.~18 & \textit{for} contxet \textit{read} context \\
1528 & l.~2 & \textit{for} currntly \textit{read} currently \\
1529 & l.~12 & \textit{for} p.~1593 \textit{read} p.~1493 (Proposition~3.15) \\
1529 & l.~21 & \textit{for} generlly \textit{read} generally \\
1529 & l.~8 from foot & \textit{delete} the second occurrence of \emph{is} \\
1529 & l.~1 from foot & \textit{for} Kretzmann (1966) \textit{read} Grim (1983) (the Cantorian argument against omniscience) \\
1530 & l.~2 & \textit{delete} the second occurrence of \emph{of} \\
1530 & l.~6 from foot & \textit{delete} the second occurrence of \emph{the} \\
\dag\,1531 & proof of Proposition~4.12, l.~7 & \textit{for} Definition~4.1 \textit{read} Definition~4.5 \\
\dag\,1531 & proof of Proposition~4.12, (4.9) & \textit{for} $K_{x}q \leftrightarrow q$ \textit{read} $K_{x}q \to q$ \\
1531 & l.~29 & \textit{for} epistemmic \textit{read} epistemic \\
\dag\,1532 & proof of Proposition~4.12, (4.10) & \textit{for} $\mathrm{MP}(x) \leftrightarrow \neg\pos\,\mathrm{Wrong}(x,a)$ \textit{read} $\mathrm{MP}(x) \to \neg\pos\,\mathrm{Wrong}(x,a)$; correctly printed in (4.11), p.~1533 \\
\dag\,1532 & proof of Proposition~4.12, l.~4 & \textit{for} Theorem~4.7 \textit{read} Theorem~4.6; the step uses Theorem~4.6 (p.~1517) \\
1532 & l.~13 & \textit{for} pesrpective \textit{read} perspective \\
1533 & l.~14 & \textit{for} consiedrable \textit{read} considerable \\
\dag\,1534 & proof of Proposition~4.14, l.~5 & \textit{for} Definition~4.11 \textit{read} Definition~4.13 \\
1534 & l.~30 & \textit{for} Assumption~2.1 of Volume~I \textit{read} Assumption~3.1 of Volume~I \\
1535 & l.~23 & \textit{for} defintiely \textit{read} definitely \\
1535 & l.~38 & \textit{for} permise \textit{read} premise \\
1535 & n.~13 & \textit{for} Kretzmann (1967) \textit{read} Kretzmann (1966) \\
1536 & l.~12 & \textit{for} p.~1548 \textit{read} p.~1458 (Definition~3.4) \\
1537 & l.~23 & \textit{for} p.~1520 \textit{read} p.~1510 (Lemma~4.3) \\
1538 & l.~13 & \textit{for} p.~1634 \textit{read} p.~1534 (Definition~4.13) \\
1538 & l.~29 & \textit{for} clarrify \textit{read} clarify \\
1539 & l.~15 & \textit{for} consderable \textit{read} considerable \\
1539 & l.~20 & \textit{for} rmearkable \textit{read} remarkable \\
\dag\,1540 & proof of Corollary~4.15, l.~21 & \textit{for} Definition~4.4 \textit{read} Definition~4.5 \\
\end{errlist}

\subsection*{Chapter~5. Moral perfection \hfill \normalfont\small pp.~1541--1582}
\begin{errlist}
1542 & l.~3 from foot & \textit{for} purpse \textit{read} purpose \\
1543 & l.~4 & \textit{for} combinatoin \textit{read} combination \\
1543 & l.~10 & \textit{for} ipmortance \textit{read} importance \\
1543 & l.~24 & \textit{for} dircetly \textit{read} directly \\
1543 & l.~31 & \textit{for} mreely \textit{read} merely \\
1544 & l.~7 & \textit{for} $\neg K_{z}q \to K_{z} K_{z}q$ \textit{read} $\neg K_{z}q \to K_{z}\neg K_{z}q$ \\
1544 & l.~10 & \textit{for} Lemma~2747 of Volume~X \textit{read} Lemma~2752 of Volume~X; the result cited is on p.~3844 \\
1544 & l.~20 & \textit{for} Lemma~2851 of Volume~X \textit{read} Lemma~2852 of Volume~X; the result cited is on p.~3854 \\
1547 & l.~28 & \textit{for} implicattion \textit{read} implication \\
1547 & l.~31 & \textit{for} sufficinet \textit{read} sufficient \\
1548 & l.~8 & \textit{delete} the second occurrence of \emph{the} \\
1548 & l.~20 & \textit{for} observatoin \textit{read} observation \\
1549 & l.~9 & \textit{for} elsewhree \textit{read} elsewhere \\
1549 & l.~17 & \textit{for} inclduing \textit{read} including \\
1549 & l.~21 & \textit{for} sufficiant \textit{read} sufficient \\
1549 & l.~31 & \textit{for} primairly \textit{read} primarily \\
\dag\,1550 & proof of Proposition~5.4, (5.2) & \textit{for} $\mathrm{Os}(y) \wedge \mathrm{Os}(y) \to \forall p\,(K_{x}p \leftrightarrow K_{y}p)$ \textit{read} $\mathrm{Os}(x) \wedge \mathrm{Os}(y) \to \forall p\,(K_{x}p \leftrightarrow K_{y}p)$ \\
1551 & l.~18 & \textit{for} indepndence \textit{read} independence \\
1552 & l.~35 & \textit{for} mechansim \textit{read} mechanism \\
1553 & l.~29 & \textit{for} referance \textit{read} reference \\
\dag\,1553 & proof of Proposition~5.6, (5.3) & \textit{for} $ Q(t,y) \wedge \mathrm{Task}(t) \to \mathrm{Can}(y,t)$ \textit{read} $\neg Q(t,y) \wedge \mathrm{Task}(t) \to \mathrm{Can}(y,t)$; the formula is used in its correct form on p.~1556 \\
1555 & l.~7 & \textit{for} Lemma~2697 of Volume~X \textit{read} Lemma~2698 of Volume~X; the result cited is on p.~3838 \\
1555 & l.~14 & \textit{for} persupposition \textit{read} presupposition \\
1555 & l.~18 & \textit{for} consisetnt \textit{read} consistent \\
\dag\,1555 & proof of Corollary~5.7, l.~26 & \textit{for} Lemma~2812 of Volume~X \textit{read} Lemma~2816 of Volume~X; the result cited is on p.~3850 \\
\dag\,1555 & proof of Corollary~5.7, l.~27 & \textit{for} Definition~5.5 \textit{read} Definition~5.1 \\
1556 & l.~16 & \textit{for} $\pos(\mathrm{Op}(x) \wedge \mathrm{Os}(x) \wedge \mathrm{MP}(x))$ \textit{read} $\nec(\mathrm{Op}(x) \wedge \mathrm{Os}(x) \wedge \mathrm{MP}(x))$ \\
1557 & l.~8 & \textit{for} deifnitely \textit{read} definitely \\
1557 & l.~19 & \textit{for} compraison \textit{read} comparison \\
1559 & l.~3 from foot & \textit{for} Pike (1968) \textit{read} Pike (1969) \\
1561 & l.~9 & \textit{for} condtional \textit{read} conditional \\
1562 & l.~27 & \textit{for} partiuclar \textit{read} particular \\
\dag\,1562 & proof of Proposition~5.9, l.~29 & \textit{for} Proposition~5.6 \textit{read} Proposition~5.4 \\
\dag\,1562 & proof of Proposition~5.9, l.~34 & \textit{for} Lemma~2671 of Volume~X \textit{read} Lemma~2666 of Volume~X \\
\dag\,1562 & proof of Proposition~5.9, (5.5) & \textit{for} $\mathrm{MP}(z) \to \pos\,\mathrm{Wrong}(z,a)$ \textit{read} $\mathrm{MP}(z) \to \neg\pos\,\mathrm{Wrong}(z,a)$ \\
\dag\,1563 & proof of Proposition~5.9, l.~16 & \textit{for} Proposition~5.4 \textit{read} Proposition~5.3; the step uses Proposition~5.3 (p.~1548) \\
1563 & l.~30 & \textit{for} particullarly \textit{read} particularly \\
1563 & l.~32 & \textit{for} formulaiton \textit{read} formulation \\
1564 & display (5.6) & \textit{for} $\nec(\mathrm{Op}(z) \vee \mathrm{Os}(z) \wedge \mathrm{MP}(z))$ \textit{read} $\nec(\mathrm{Op}(z) \wedge \mathrm{Os}(z) \wedge \mathrm{MP}(z))$ \\
1565 & l.~3 & \textit{delete} the second occurrence of \emph{that} \\
1566 & l.~10 & \textit{for} p.~1505 \textit{read} p.~1504 (Proposition~4.2) \\
1567 & display (5.7) & \textit{for} $\neg Q(t,z) \vee \mathrm{Task}(t) \to \mathrm{Can}(z,t)$ \textit{read} $\neg Q(t,z) \wedge \mathrm{Task}(t) \to \mathrm{Can}(z,t)$ \\
1567 & l.~7 & \textit{for} modset \textit{read} modest \\
1567 & l.~7 from foot & \textit{for} unecessary \textit{read} unnecessary \\
1568 & l.~7 & \textit{for} $\mathrm{Os}(z) \wedge \mathrm{Os}(z) \to \forall p\,(K_{z}p \leftrightarrow K_{y}p)$ \textit{read} $\mathrm{Os}(z) \wedge \mathrm{Os}(y) \to \forall p\,(K_{z}p \leftrightarrow K_{y}p)$ \\
1569 & l.~7 & \textit{for} prnciple \textit{read} principle \\
1569 & l.~10 & \textit{for} p.~1658 \textit{read} p.~1568 (Definition~5.10) \\
\dag\,1569 & proof of Proposition~5.11, l.~17 & \textit{for} Corollary~5.2 \textit{read} Corollary~5.7 \\
1569 & l.~18 & \textit{for} premissible \textit{read} permissible \\
\dag\,1569 & proof of Proposition~5.11, (5.8) & \textit{for} $K_{z}q \leftrightarrow q$ \textit{read} $K_{z}q \to q$ \\
1569 & l.~32 & \textit{for} p.~1527 \textit{read} p.~1517 (Definition~4.5) \\
1571 & l.~5 & \textit{for} prbable \textit{read} probable \\
\dag\,1574 & proof of Lemma~5.12, l.~2 & \textit{for} Lemma~2845 of Volume~X \textit{read} Lemma~2844 of Volume~X; the result cited is on p.~3853 \\
1574 & l.~11 & \textit{for} benefiical \textit{read} beneficial \\
1574 & l.~21 & \textit{for} p.~1521 \textit{read} p.~1520 (Lemma~4.9) \\
\dag\,1574 & proof of Lemma~5.12, l.~23 & \textit{for} Proposition~5.9 \textit{read} Proposition~5.6 \\
1574 & l.~32 & \textit{for} unnecssary \textit{read} unnecessary \\
\dag\,1575 & proof of Lemma~5.12, l.~25 & \textit{for} Lemma~2911 of Volume~X \textit{read} Lemma~2910 of Volume~X \\
1575 & l.~4 from foot & \textit{for} srtucture \textit{read} structure \\
1577 & l.~22 & \textit{for} cocnerned \textit{read} concerned \\
1577 & l.~35 & \textit{for} Lemma~2996 of Volume~X \textit{read} Lemma~2999 of Volume~X; the result cited is on p.~3869 \\
1578 & l.~6 & \textit{for} trnsition \textit{read} transition \\
1578 & l.~29 & \textit{for} assumptoin \textit{read} assumption \\
1579 & l.~3 & \textit{for} obtian \textit{read} obtain \\
1579 & l.~8 & \textit{for} Morris (1987) \textit{read} Adams (1972) (whether God must create the best) \\
1579 & l.~23 & \textit{for} Volume~IV, Chapter~5 \textit{read} Volume~IV, Chapter~7 (necessary existence and aseity) \\
1579 & display (5.10) & \textit{for} $\neg K_{x}r \leftrightarrow K_{x}\neg K_{x}r$ \textit{read} $\neg K_{x}r \to K_{x}\neg K_{x}r$ \\
1580 & l.~19 & \textit{for} p.~1472 \textit{read} p.~1471 (Theorem~3.7) \\
1580 & l.~36 & \textit{for} Lemma~2873 of Volume~X \textit{read} Lemma~2877 of Volume~X \\
1581 & l.~4 & \textit{for} Adams (1971) \textit{read} Adams (1972) \\
1581 & l.~32 & \textit{for} p.~1556 \textit{read} p.~1546 (Corollary~5.2) \\
1582 & l.~1 & \textit{for} tranistion \textit{read} transition \\
1582 & l.~10 & \textit{for} Pike (1970) \textit{read} Pike (1969) \\
1582 & l.~25 & \textit{for} Lemma~2711 of Volume~X \textit{read} Lemma~2715 of Volume~X; the result cited is on p.~3840 \\
\end{errlist}

\subsection*{Chapter~6. The compossibility of the attributes \hfill \normalfont\small pp.~1583--1630}
\begin{errlist}
1584 & running heads, pp.~1584--1630 (even) & \textit{for} Moral perfection \textit{read} Compossibility of the attributes \\
1584 & l.~17 & \textit{for} throughuot \textit{read} throughout \\
1585 & l.~3 & \textit{for} intrdouce \textit{read} introduce \\
1585 & l.~19 & \textit{for} exmaple \textit{read} example \\
1587 & l.~24 & \textit{for} Lemma~2734 of Volume~X \textit{read} Lemma~2731 of Volume~X; the result cited is on p.~3842 \\
1589 & l.~12 & \textit{see} Corrigendum~S7 \\
1590 & l.~37 & \textit{for} fortiroi \textit{read} fortiori \\
1591 & display (6.5) & \textit{for} $\neg K_{x}p \leftrightarrow K_{x}\neg K_{x}p$ \textit{read} $\neg K_{x}p \to K_{x}\neg K_{x}p$ \\
1591 & l.~13 & \textit{for} conceivbale \textit{read} conceivable \\
1592 & l.~20 & \textit{for} $\mathrm{Pos}({\chi}) \to \pos\mathrm{Pos}({\chi})$ \textit{read} $\mathrm{Pos}({\chi}) \to \nec\mathrm{Pos}({\chi})$ \\
1592 & l.~22 & \textit{for} acecssible \textit{read} accessible \\
1594 & l.~4 & \textit{for} siimlarly \textit{read} similarly \\
1595 & l.~27 & \textit{for} nceessary \textit{read} necessary \\
1596 & l.~7 & \textit{for} Lemma~2907 of Volume~X \textit{read} Lemma~2913 of Volume~X; the result cited is on p.~3860 \\
1596 & l.~28 & \textit{for} intrnsic \textit{read} intrinsic \\
1597 & l.~11 & \textit{for} framwork \textit{read} framework \\
1598 & l.~1 from foot & \textit{for} disagerement \textit{read} disagreement \\
1599 & l.~3 & \textit{for} infrence \textit{read} inference \\
1599 & l.~5 & \textit{for} porspect \textit{read} prospect \\
1599 & l.~12 & \textit{for} accesible \textit{read} accessible \\
1599 & l.~21 & \textit{for} interprettaion \textit{read} interpretation \\
1599 & l.~7 from foot & \textit{for} posittive \textit{read} positive \\
1600 & l.~32 & \textit{for} undrestanding \textit{read} understanding \\
1600 & l.~2 from foot & \textit{for} speciifc \textit{read} specific \\
1603 & l.~8 & \textit{for} fitler \textit{read} filter \\
1603 & l.~28 & \textit{delete} the second occurrence of \emph{to} \\
1603 & display (6.10) & \textit{for} $\mathrm{Pos}(\neg {\chi}) \leftrightarrow \mathrm{Pos}({\chi})$ \textit{read} $\mathrm{Pos}(\neg {\chi}) \leftrightarrow \neg\mathrm{Pos}({\chi})$ \\
1604 & l.~11 & \textit{for} suggsetion \textit{read} suggestion \\
1604 & l.~24 & \textit{for} scuceed \textit{read} succeed \\
1604 & display (6.11) & \textit{for} $\neg Q(t,x) \wedge \mathrm{Task}(t) \leftrightarrow \mathrm{Can}(x,t)$ \textit{read} $\neg Q(t,x) \wedge \mathrm{Task}(t) \to \mathrm{Can}(x,t)$; correctly printed in (6.12), p.~1604 \\
1604 & l.~6 from foot & \textit{for} discssion \textit{read} discussion \\
1607 & l.~6 & \textit{for} inutitive \textit{read} intuitive \\
1607 & l.~18 & \textit{for} accrdingly \textit{read} accordingly \\
1607 & l.~37 & \textit{for} p.~1532 \textit{read} p.~1531 (Proposition~4.12) \\
1608 & l.~2 & \textit{for} hyopthesis \textit{read} hypothesis \\
1608 & l.~35 & \textit{for} Morris (1987) \textit{read} Pike (1969) (omnipotence and the ability to sin) \\
1609 & l.~2 & \textit{for} p.~1646 \textit{read} p.~1546 (Definition~5.1) \\
1609 & l.~12 & \textit{for} prevoius \textit{read} previous \\
\dag\,1609 & proof of Lemma~6.6, (6.17) & \textit{for} $\mathrm{MP}(z) \leftrightarrow \neg\pos\,\mathrm{Wrong}(z,a)$ \textit{read} $\mathrm{MP}(z) \to \neg\pos\,\mathrm{Wrong}(z,a)$; correctly printed in (6.18), p.~1611 \\
1611 & l.~13 & \textit{for} wrongdiong \textit{read} wrongdoing \\
1611 & l.~36 & \textit{for} addtiional \textit{read} additional \\
1612 & l.~2 & \textit{for} evaluatte \textit{read} evaluate \\
1613 & display (6.19) & \textit{for} $\mathrm{Os}(x) \wedge \mathrm{Os}(z) \to \forall q\,(K_{x}q \to K_{z}q)$ \textit{read} $\mathrm{Os}(x) \wedge \mathrm{Os}(z) \to \forall q\,(K_{x}q \leftrightarrow K_{z}q)$; correctly printed in (6.24), p.~1618 \\
1615 & l.~5 & \textit{for} graet-making \textit{read} great-making \\
1615 & l.~8 & \textit{for} rsepectively \textit{read} respectively \\
1615 & display (6.21) & \textit{for} ${\phi}, {\psi} \notin \mathcal{P} \Rightarrow {\phi} \sqcap {\psi} \in \mathcal{P}$ \textit{read} ${\phi}, {\psi} \in \mathcal{P} \Rightarrow {\phi} \sqcap {\psi} \in \mathcal{P}$; correctly printed in (6.22), p.~1617 \\
1616 & l.~4 & \textit{for} primarliy \textit{read} primarily \\
1616 & l.~20 & \textit{delete} the second occurrence of \emph{of} \\
1616 & l.~40 & \textit{for} eaxmple \textit{read} example \\
1617 & l.~7 & \textit{for} p.~1502 \textit{read} p.~1512 (Definition~4.4) \\
\dag\,1617 & proof of Proposition~6.10, (6.23) & \textit{for} $\nec(\mathrm{Op}(x) \wedge \mathrm{Os}(x) \vee \mathrm{MP}(x))$ \textit{read} $\nec(\mathrm{Op}(x) \wedge \mathrm{Os}(x) \wedge \mathrm{MP}(x))$ \\
\dag\,1617 & proof of Proposition~6.10, l.~12 & \textit{for} Lemma~2945 of Volume~X \textit{read} Lemma~2944 of Volume~X \\
1617 & l.~15 & \textit{for} invarinat \textit{read} invariant \\
\dag\,1617 & proof of Proposition~6.9, l.~16 & \textit{for} Definition~6.7 \textit{read} Definition~6.8 \\
1617 & l.~19 & \textit{delete} the second occurrence of \emph{the} \\
1618 & l.~3 & \textit{for} rigoruos \textit{read} rigorous \\
1618 & l.~19 & \textit{for} terminoolgy \textit{read} terminology \\
1618 & l.~30 & \textit{for} propoistion \textit{read} proposition \\
1618 & l.~32 & \textit{for} exlpicit \textit{read} explicit \\
\dag\,1619 & proof of Corollary~6.11, (6.25) & \textit{for} $\mathrm{Os}(z) \wedge \mathrm{Os}(y) \to \exists q\,(K_{z}q \leftrightarrow K_{y}q)$ \textit{read} $\mathrm{Os}(z) \wedge \mathrm{Os}(y) \to \forall q\,(K_{z}q \leftrightarrow K_{y}q)$; the formula is used in its correct form on p.~1623 \\
1619 & l.~29 & \textit{for} konwledge \textit{read} knowledge \\
1620 & l.~2 & \textit{for} Lemma~2788 of Volume~X \textit{read} Lemma~2784 of Volume~X; the result cited is on p.~3847 \\
1620 & l.~16 & \textit{for} observasion \textit{read} observation \\
1620 & l.~28 & \textit{for} separte \textit{read} separate \\
\dag\,1622 & proof of Lemma~6.12, l.~11 & \textit{for} Proposition~6.13 \textit{read} Proposition~6.9 \\
1622 & l.~33 & \textit{for} identfy \textit{read} identify \\
1623 & l.~11 & \textit{for} the empty property belongs to $\mathcal{P}$ \textit{read} the empty property does not belong to $\mathcal{P}$ \\
1623 & l.~31 & \textit{for} whennever \textit{read} whenever \\
1624 & l.~1 & \textit{for} Lemma~2718 of Volume~X \textit{read} Lemma~2724 of Volume~X; the result cited is on p.~3841 \\
1624 & l.~26 & \textit{for} p.~1501 \textit{read} p.~1601 (Lemma~6.4) \\
1624 & l.~27 & \textit{for} perspectve \textit{read} perspective \\
1624 & l.~3 from foot & \textit{for} great-amking \textit{read} great-making \\
1625 & l.~3 & \textit{for} poistive \textit{read} positive \\
1625 & l.~20 & \textit{for} embararssment \textit{read} embarrassment \\
1625 & l.~8 from foot & \textit{for} deciisve \textit{read} decisive \\
1626 & l.~36 & \textit{for} p.~1537 \textit{read} p.~1527 (Lemma~4.10) \\
1626 & l.~3 from foot & \textit{for} Lemma~3029 of Volume~X \textit{read} Lemma~3027 of Volume~X; the result cited is on p.~3872 \\
1627 & l.~27 & \textit{for} ocucrrence \textit{read} occurrence \\
1627 & l.~32 & \textit{for} nammely \textit{read} namely \\
1628 & l.~14 & \textit{for} p.~1518 \textit{read} p.~1618 (Corollary~6.11) \\
\dag\,1628 & proof of Lemma~6.16, l.~20 & \textit{for} Lemma~2780 of Volume~X \textit{read} Lemma~2777 of Volume~X; the result cited is on p.~3846 \\
\dag\,1628 & proof of Lemma~6.16, l.~21 & \textit{for} Lemma~6.6 \textit{read} Lemma~6.4 \\
1628 & l.~31 & \textit{for} mettaphysical \textit{read} metaphysical \\
1629 & l.~3 & \textit{for} seaprate \textit{read} separate \\
1629 & l.~6 & \textit{for} joitnly \textit{read} jointly \\
1629 & l.~6 from foot & \textit{for} ${\psi} \notin \mathcal{P} \wedge \nec\forall z\,({\psi}z \to {\phi}z) \Rightarrow {\phi} \in \mathcal{P}$ \textit{read} ${\psi} \in \mathcal{P} \wedge \nec\forall z\,({\psi}z \to {\phi}z) \Rightarrow {\phi} \in \mathcal{P}$ \\
\end{errlist}

\subsection*{Chapter~7. Necessary existence and aseity \hfill \normalfont\small pp.~1631--1671}
\begin{errlist}
1632 & running heads, pp.~1632--1670 (even) & \textit{for} Godlikeness and maximal greatness \textit{read} Necessary existence and aseity \\
1632 & l.~9 & \textit{for} p.~1439 \textit{read} p.~1539 (Corollary~4.15) \\
1632 & l.~7 from foot & \textit{for} unneecssary \textit{read} unnecessary \\
1633 & l.~1 from foot & \textit{for} Plantnga (1980) \textit{read} Plantinga (1980) \\
1634 & display (7.1) & \textit{for} $wRw_0 \leftrightarrow D_{w} \subseteq D_{w_0}$ \textit{read} $wRw_0 \to D_{w} \subseteq D_{w_0}$ \\
1635 & l.~8 & \textit{for} stregnth \textit{read} strength \\
1635 & l.~17 & \textit{for} formultion \textit{read} formulation \\
1635 & l.~26 & \textit{for} contrversy \textit{read} controversy \\
1637 & l.~8 & \textit{for} $\mathrm{Op}(y) \leftrightarrow \forall t\,(\mathrm{Task}(t) \wedge \neg Q(t,y) \leftrightarrow \mathrm{Can}(y,t))$ \textit{read} $\mathrm{Op}(y) \leftrightarrow \forall t\,(\mathrm{Task}(t) \wedge \neg Q(t,y) \to \mathrm{Can}(y,t))$ \\
1638 & display (7.3) & \textit{for} $\forall y\,\nec {\phi}y \leftrightarrow \nec\forall y\,{\phi}y$ \textit{read} $\forall y\,\nec {\phi}y \to \nec\forall y\,{\phi}y$ \\
1638 & l.~26 & \textit{for} Plantinga (1979) \textit{read} Plantinga (1980) \\
1639 & l.~19 & \textit{for} implictaion \textit{read} implication \\
1640 & n.~15 & \textit{for} Prus (2006) \textit{read} Pruss (2006) \\
1641 & l.~27 & \textit{for} ordnary \textit{read} ordinary \\
1642 & l.~5 & \textit{for} p.~1652 \textit{read} p.~1552 (Definition~5.5) \\
1643 & l.~2 & \textit{for} intuitvie \textit{read} intuitive \\
1643 & l.~15 & \textit{for} $E(y,v') \leftrightarrow y \notin D_{v'}$ \textit{read} $E(y,v') \leftrightarrow y \in D_{v'}$ \\
1644 & l.~1 & \textit{for} parargaph \textit{read} paragraph \\
1644 & l.~27 & \textit{for} Fine (1993) \textit{read} Fine (1994) \\
1644 & display (7.6) & \textit{for} $\neg K_{z}r \to K_{z} K_{z}r$ \textit{read} $\neg K_{z}r \to K_{z}\neg K_{z}r$ \\
1645 & display (7.8) & \textit{for} $\exists z\,\pos {\psi}z \to \pos\forall z\,{\psi}z$ \textit{read} $\exists z\,\pos {\psi}z \to \pos\exists z\,{\psi}z$; correctly printed in (7.9), p.~1645 \\
1646 & l.~10 & \textit{for} cetrainly \textit{read} certainly \\
1646 & l.~38 & \textit{for} ocurred \textit{read} occurred \\
1647 & l.~7 & \textit{for} withdarwn \textit{read} withdrawn \\
\dag\,1648 & proof of Lemma~7.3, l.~26 & \textit{for} Definition~7.5 \textit{read} Definition~7.1 \\
\dag\,1648 & proof of Lemma~7.3, l.~30 & \textit{for} Definition~7.1 \textit{read} Definition~7.2 \\
1649 & l.~5 & \textit{for} providd \textit{read} provided \\
1649 & l.~18 & \textit{for} Lemma~2686 of Volume~X \textit{read} Lemma~2682 of Volume~X \\
1649 & l.~36 & \textit{for} Pruss (2006) \textit{read} Fine (1994) (essence and modality) \\
1650 & l.~22 & \textit{for} independnt \textit{read} independent \\
1651 & l.~7 & \textit{for} sytsematic \textit{read} systematic \\
\dag\,1651 & proof of Proposition~7.6, (7.10) & \textit{for} $\mathrm{NE}(z) \to \neg\forall y\,D(z,y)$ \textit{read} $\mathrm{NE}(z) \to \neg\exists y\,D(z,y)$; correctly printed in (7.11), p.~1652 \\
1651 & l.~8 from foot & \textit{for} transittion \textit{read} transition \\
1651 & n.~6 & \textit{for} Leftw (2012) \textit{read} Leftow (2012) \\
1652 & l.~12 & \textit{for} increasing domains validate the Barcan formula \textit{read} increasing domains validate the converse Barcan formula \\
\dag\,1652 & proof of Proposition~7.7, l.~30 & \textit{for} Lemma~3018 of Volume~X \textit{read} Lemma~3023 of Volume~X; the result cited is on p.~3871 \\
\dag\,1653 & proof of Proposition~7.7, l.~19 & \textit{for} Definition~7.5 \textit{read} Definition~7.2; the step uses Definition~7.2 (p.~1643) \\
\dag\,1653 & proof of Proposition~7.7, l.~23 & \textit{for} Proposition~7.11 \textit{read} Proposition~7.6 \\
1653 & l.~28 & \textit{for} Fine (1995) \textit{read} Fine (1994) \\
1653 & l.~5 from foot & \textit{for} parllel \textit{read} parallel \\
1654 & l.~18 & \textit{for} Plantinga (1981) \textit{read} Plantinga (1980) \\
1654 & l.~24 & \textit{for} utlimately \textit{read} ultimately \\
1658 & l.~5 & \textit{after} exists \textit{insert} at $w_1$ \\
1658 & l.~18 & \textit{for} contsant \textit{read} constant \\
1658 & l.~20 & \textit{for} isntead \textit{read} instead \\
1658 & l.~30 & \textit{for} develpment \textit{read} development \\
1659 & l.~18 & \textit{for} a variable domain validates necessitism \textit{read} a constant domain validates necessitism \\
1661 & l.~14 & \textit{for} unniversal \textit{read} universal \\
1661 & display (7.17) & \textit{for} $\mathrm{Can}(z,t) \leftrightarrow \pos\,\mathrm{Does}(z,t)$ \textit{read} $\mathrm{Can}(z,t) \to \pos\,\mathrm{Does}(z,t)$; the formula is used in its correct form on p.~1665 \\
1661 & l.~5 from foot & \textit{for} procss \textit{read} process \\
1662 & l.~9 & \textit{for} problm \textit{read} problem \\
1662 & l.~24 & \textit{for} p.~1487 \textit{read} p.~1587 (Definition~6.1) \\
\dag\,1662 & proof of Theorem~7.12, l.~25 & \textit{for} Proposition~7.11 \textit{read} Proposition~7.7 \\
1663 & l.~31 & \textit{for} Volume~IX, Chapter~5 \textit{read} Volume~IX, Chapter~6 (whether a maximally great being must create the best) \\
\dag\,1663 & proof of Theorem~7.12, (7.21) & \textit{for} $\nec(\exists x\,\mathrm{MG}(x) \leftrightarrow E(g))$ \textit{read} $\nec(\exists x\,\mathrm{MG}(x) \to E(g))$ \\
1664 & l.~33 & \textit{for} strctly \textit{read} strictly \\
1665 & l.~13 & \textit{for} apparnetly \textit{read} apparently \\
1665 & display (7.22) & \textit{for} $w_0Rw_2 \to D_{w_0} \subseteq D_{w_1}$ \textit{read} $w_0Rw_2 \to D_{w_0} \subseteq D_{w_2}$ \\
1665 & l.~32 & \textit{for} srtaightforward \textit{read} straightforward \\
1665 & l.~35 & \textit{for} Lemma~2910 of Volume~X \textit{read} Lemma~2909 of Volume~X; the result cited is on p.~3860 \\
1666 & l.~9 & \textit{for} guaarntee \textit{read} guarantee \\
1666 & l.~20 & \textit{for} Letfow (2012) \textit{read} Leftow (2012) \\
1666 & l.~31 & \textit{after} domain \textit{insert} of the world $w_1$ \\
1666 & l.~34 & \textit{for} de dicto \textit{read} de re \\
1666 & l.~2 from foot & \textit{for} p.~1722 \textit{read} p.~1622 (Lemma~6.12) \\
1667 & l.~1 & \textit{for} Pruss (2007) \textit{read} Pruss (2006) \\
1667 & l.~13 & \textit{for} familliar \textit{read} familiar \\
1668 & l.~3 & \textit{for} $\forall z\,\nec {\phi}z \leftrightarrow \nec\forall z\,{\phi}z$ \textit{read} $\forall z\,\nec {\phi}z \to \nec\forall z\,{\phi}z$ \\
1668 & l.~7 & \textit{delete} the second occurrence of \emph{the} \\
1668 & display (7.24) & \textit{for} $\nec\exists x\,{\phi}x \to \exists x\,\pos {\phi}x$ \textit{read} $\pos\exists x\,{\phi}x \to \exists x\,\pos {\phi}x$ \\
1668 & l.~29 & \textit{for} independnece \textit{read} independence \\
1670 & l.~17 & \textit{for} virtaully \textit{read} virtually \\
1671 & l.~17 & \textit{for} serrious \textit{read} serious \\
\end{errlist}

\subsection*{Chapter~8. Godlikeness and maximal greatness \hfill \normalfont\small pp.~1672--1728}
\begin{errlist}
1674 & running heads, pp.~1674--1728 (even) & \textit{for} Necessary existence and aseity \textit{read} Godlikeness and maximal greatness \\
1674 & display (8.1) & \textit{for} $\mathrm{Pos}({\psi}) \to \nec\exists x\,{\psi}x$ \textit{read} $\mathrm{Pos}({\psi}) \to \pos\exists x\,{\psi}x$ \\
1675 & l.~3 & \textit{for} perefction \textit{read} perfection \\
1675 & l.~17 & \textit{for} combiantion \textit{read} combination \\
1675 & l.~30 & \textit{for} Lemma~2849 of Volume~X \textit{read} Lemma~2848 of Volume~X \\
1676 & l.~10 & \textit{for} tehcnical \textit{read} technical \\
1676 & l.~16 & \textit{for} admitteldy \textit{read} admittedly \\
1676 & l.~32 & \textit{for} tehrefore \textit{read} therefore \\
1677 & display (8.3) & \textit{for} $\mathrm{ME}(x,w_2) \iff \mathrm{Op}(x,w_1) \wedge \mathrm{Os}(x,w_1) \wedge \mathrm{MP}(x,w_1)$ \textit{read} $\mathrm{ME}(x,w_1) \iff \mathrm{Op}(x,w_1) \wedge \mathrm{Os}(x,w_1) \wedge \mathrm{MP}(x,w_1)$ \\
1677 & l.~8 & \textit{delete} the second occurrence of \emph{a} \\
1677 & l.~15 & \textit{for} Lemma~2842 of Volume~X \textit{read} Lemma~2837 of Volume~X; the result cited is on p.~3852 \\
1678 & l.~2 & \textit{for} defintely \textit{read} definitely \\
1678 & l.~17 & \textit{for} proeprty \textit{read} property \\
1678 & l.~31 & \textit{for} Lemma~2653 of Volume~X \textit{read} Lemma~2658 of Volume~X; the result cited is on p.~3834 \\
1679 & l.~16 & \textit{for} procdeure \textit{read} procedure \\
\dag\,1679 & proof of Lemma~8.2, (8.4) & \textit{for} $\pos\forall x\,\mathrm{MG}(x) \to \nec\exists x\,\mathrm{MG}(x)$ \textit{read} $\pos\exists x\,\mathrm{MG}(x) \to \nec\exists x\,\mathrm{MG}(x)$ \\
1679 & l.~2 from foot & \textit{for} obejction \textit{read} objection \\
1680 & l.~11 & \textit{for} $\mathcal{P}_0 \supseteq \mathcal{P} \Rightarrow \exists y\,\exists w_0\,\forall {\psi} \in \mathcal{P}_0\; {\psi}(y,w_0)$ \textit{read} $\mathcal{P}_0 \subseteq \mathcal{P} \Rightarrow \exists y\,\exists w_0\,\forall {\psi} \in \mathcal{P}_0\; {\psi}(y,w_0)$ \\
1684 & l.~4 & \textit{for} rigrous \textit{read} rigorous \\
1684 & l.~25 & \textit{for} convitcion \textit{read} conviction \\
1684 & l.~34 & \textit{for} Lemma~2892 of Volume~X \textit{read} Lemma~2890 of Volume~X \\
1685 & l.~8 & \textit{for} demonsrtate \textit{read} demonstrate \\
1685 & display (8.7) & \textit{for} $\mathrm{NME}(x,v') \iff \forall w_0\,(v'Rw_0 \leftrightarrow \mathrm{ME}(x,w_0))$ \textit{read} $\mathrm{NME}(x,v') \iff \forall w_0\,(v'Rw_0 \to \mathrm{ME}(x,w_0))$ \\
1687 & l.~13 & \textit{for} nmaely \textit{read} namely \\
1687 & l.~17 & \textit{for} maxmal \textit{read} maximal \\
1687 & l.~30 & \textit{for} ${\chi} \in \mathcal{P} \wedge \nec\exists x\,({\chi}x \to {\psi}x) \Rightarrow {\psi} \in \mathcal{P}$ \textit{read} ${\chi} \in \mathcal{P} \wedge \nec\forall x\,({\chi}x \to {\psi}x) \Rightarrow {\psi} \in \mathcal{P}$ \\
1687 & l.~3 from foot & \textit{for} categroy \textit{read} category \\
1687 & l.~1 from foot & \textit{for} Lemma~2679 of Volume~X \textit{read} Lemma~2678 of Volume~X \\
1688 & l.~15 & \textit{for} recevie \textit{read} receive \\
1690 & l.~12 & \textit{delete} the second occurrence of \emph{to} \\
1690 & l.~20 & \textit{for} $\mathrm{G}y \to \forall {\psi}\,(\mathrm{Pos}({\psi}) \leftrightarrow \nec {\psi}y)$ \textit{read} $\mathrm{G}y \leftrightarrow \forall {\psi}\,(\mathrm{Pos}({\psi}) \leftrightarrow \nec {\psi}y)$ \\
1690 & l.~8 from foot & \textit{for} aplpication \textit{read} application \\
\dag\,1691 & proof of Lemma~8.3, l.~21 & \textit{for} Lemma~2803 of Volume~X \textit{read} Lemma~2805 of Volume~X \\
1691 & l.~24 & \textit{for} hypotehsis \textit{read} hypothesis \\
\dag\,1691 & proof of Lemma~8.3, (8.10) & \textit{for} ${\chi}, {\phi} \in \mathcal{P} \Rightarrow {\chi} \sqcap {\phi} \notin \mathcal{P}$ \textit{read} ${\chi}, {\phi} \in \mathcal{P} \Rightarrow {\chi} \sqcap {\phi} \in \mathcal{P}$ \\
\dag\,1692 & proof of Lemma~8.3, l.~10 & \textit{for} Definition~8.6 \textit{read} Definition~8.1 \\
1694 & Proposition~8.5 & \textit{see} Corrigendum~S6 \\
1695 & l.~2 & \textit{for} esence \textit{read} essence \\
\dag\,1695 & proof of Proposition~8.5, l.~3 & \textit{for} Lemma~8.2 \textit{read} Lemma~8.3; the step uses Lemma~8.3 (p.~1691) \\
1695 & l.~9 & \textit{for} pracitcal \textit{read} practical \\
1695 & l.~6 from foot & \textit{for} adequatte \textit{read} adequate \\
\dag\,1695 & proof of Proposition~8.5, l.~5 from foot & \textit{for} Lemma~2972 of Volume~X \textit{read} Lemma~2976 of Volume~X \\
1697 & l.~8 & \textit{for} combinaiton \textit{read} combination \\
1697 & l.~37 & \textit{for} Lemma~2940 of Volume~X \textit{read} Lemma~2941 of Volume~X \\
1700 & display (8.14) & \textit{for} $\mathrm{Pos}(\neg {\phi}) \to \neg\mathrm{Pos}({\phi})$ \textit{read} $\mathrm{Pos}(\neg {\phi}) \leftrightarrow \neg\mathrm{Pos}({\phi})$; correctly printed in (8.16), p.~1703 \\
1700 & l.~19 & \textit{for} naturlaly \textit{read} naturally \\
1700 & l.~28 & \textit{for} definietly \textit{read} definitely \\
1700 & l.~4 from foot & \textit{for} immdeiate \textit{read} immediate \\
1701 & l.~11 & \textit{for} Lemma~2773 of Volume~X \textit{read} Lemma~2777 of Volume~X \\
1702 & l.~27 & \textit{for} $\mathrm{Pos}({\psi}) \vee \mathrm{Pos}({\chi}) \to \mathrm{Pos}({\psi} \sqcap {\chi})$ \textit{read} $\mathrm{Pos}({\psi}) \wedge \mathrm{Pos}({\chi}) \to \mathrm{Pos}({\psi} \sqcap {\chi})$ \\
1702 & l.~29 & \textit{for} $\pos\exists z\,\mathrm{MG}(z) \to \nec\forall z\,\mathrm{MG}(z)$ \textit{read} $\pos\exists z\,\mathrm{MG}(z) \to \nec\exists z\,\mathrm{MG}(z)$ \\
1704 & l.~24 & \textit{for} trnasferred \textit{read} transferred \\
1705 & display (8.17) & \textit{for} ${\phi} \in \mathcal{P} \wedge \nec\forall z\,({\phi}z \to {\chi}z) \Rightarrow {\chi} \notin \mathcal{P}$ \textit{read} ${\phi} \in \mathcal{P} \wedge \nec\forall z\,({\phi}z \to {\chi}z) \Rightarrow {\chi} \in \mathcal{P}$ \\
1705 & display (8.18) & \textit{for} $\mathrm{Pos}({\phi}) \leftrightarrow \pos\exists x\,{\phi}x$ \textit{read} $\mathrm{Pos}({\phi}) \to \pos\exists x\,{\phi}x$; the formula is used in its correct form on p.~1706 \\
1705 & l.~31 & \textit{for} $\exists x\,\mathrm{MG}(x,w) \to \exists x\,(E(x,w) \vee \mathrm{ME}(x,w))$ \textit{read} $\exists x\,\mathrm{MG}(x,w) \to \exists x\,(E(x,w) \wedge \mathrm{ME}(x,w))$ \\
1706 & l.~18 & \textit{delete} the second occurrence of \emph{a} \\
\dag\,1706 & proof of Theorem~8.10, (8.19) & \textit{for} $\mathrm{NME}(z,w_2) \iff \forall w_1\,(w_1Rw_2 \to \mathrm{ME}(z,w_1))$ \textit{read} $\mathrm{NME}(z,w_2) \iff \forall w_1\,(w_2Rw_1 \to \mathrm{ME}(z,w_1))$ \\
1707 & l.~29 & \textit{for} definiiton \textit{read} definition \\
1708 & l.~8 & \textit{for} essnece \textit{read} essence \\
\dag\,1709 & proof of Proposition~8.11, (8.20) & \textit{for} ${\psi}, {\phi} \notin \mathcal{P} \Rightarrow {\psi} \sqcap {\phi} \in \mathcal{P}$ \textit{read} ${\psi}, {\phi} \in \mathcal{P} \Rightarrow {\psi} \sqcap {\phi} \in \mathcal{P}$; correctly printed in (8.22), p.~1714 \\
1710 & l.~5 & \textit{for} parctical \textit{read} practical \\
1711 & l.~14 & \textit{for} rellevance \textit{read} relevance \\
\dag\,1711 & proof of Lemma~8.13, l.~25 & \textit{for} Theorem~8.7 \textit{read} Theorem~8.10; the step uses Theorem~8.10 (p.~1706) \\
1711 & l.~27 & \textit{for} $\forall z\,\mathrm{MG}(z,v) \to \exists z\,(E(z,v) \wedge \mathrm{ME}(z,v))$ \textit{read} $\exists z\,\mathrm{MG}(z,v) \to \exists z\,(E(z,v) \wedge \mathrm{ME}(z,v))$ \\
\dag\,1711 & proof of Lemma~8.13, l.~28 & \textit{for} Definition~8.12 \textit{read} Definition~8.8 \\
1713 & display (8.21) & \textit{for} $\mathrm{G}y \to \forall {\psi}\,(\mathrm{Pos}({\psi}) \to {\psi}y)$ \textit{read} $\mathrm{G}y \leftrightarrow \forall {\psi}\,(\mathrm{Pos}({\psi}) \to {\psi}y)$ \\
1713 & l.~20 & \textit{for} parlalel \textit{read} parallel \\
1714 & l.~10 & \textit{for} Lemma~3033 of Volume~X \textit{read} Lemma~3035 of Volume~X \\
1714 & l.~19 & \textit{for} occasionlaly \textit{read} occasionally \\
1715 & l.~9 & \textit{for} partcular \textit{read} particular \\
\dag\,1715 & proof of Lemma~8.17, l.~28 & \textit{for} Definition~8.15 \textit{read} Definition~8.14; the step uses Definition~8.14 (p.~1712) \\
1717 & l.~11 & \textit{for} Gödel (1995) \textit{read} Anderson (1990) (the emended positivity axioms) \\
1717 & display (8.25) & \textit{for} $\mathcal{P}_0 \subseteq \mathcal{P} \Rightarrow \exists x\,\exists w_2\,\forall {\phi} \in \mathcal{P}_0\; {\phi}(x,w_1)$ \textit{read} $\mathcal{P}_0 \subseteq \mathcal{P} \Rightarrow \exists x\,\exists w_1\,\forall {\phi} \in \mathcal{P}_0\; {\phi}(x,w_1)$; correctly printed in (8.31), p.~1723 \\
1717 & l.~27 & \textit{for} $\mathrm{MG}(z,u) \iff \forall v'\,(v'Ru \to E(z,v') \wedge \mathrm{ME}(z,v'))$ \textit{read} $\mathrm{MG}(z,u) \iff \forall v'\,(uRv' \to E(z,v') \wedge \mathrm{ME}(z,v'))$ \\
1717 & l.~34 & \textit{for} p.~1525 \textit{read} p.~1625 (Theorem~6.14) \\
1717 & display (8.26) & \textit{for} $\mathrm{MG}(y,v) \leftrightarrow \mathrm{NE}(y,v) \wedge \mathrm{NME}(y,v)$ \textit{read} $\mathrm{MG}(y,v) \to \mathrm{NE}(y,v) \wedge \mathrm{NME}(y,v)$; correctly printed in (8.27), p.~1719 \\
1718 & l.~24 & \textit{for} defnded \textit{read} defended \\
1718 & l.~26 & \textit{for} $\mathrm{MG}(x,w') \to \mathrm{NE}(x,w') \vee \mathrm{NME}(x,w')$ \textit{read} $\mathrm{MG}(x,w') \to \mathrm{NE}(x,w') \wedge \mathrm{NME}(x,w')$ \\
1718 & l.~7 from foot & \textit{for} relevnace \textit{read} relevance \\
1719 & l.~24 & \textit{for} literatture \textit{read} literature \\
1719 & l.~4 from foot & \textit{for} requirment \textit{read} requirement \\
1720 & display (8.28) & \textit{for} $\mathrm{G}z \leftrightarrow \exists {\phi}\,(\mathrm{Pos}({\phi}) \to {\phi}z)$ \textit{read} $\mathrm{G}z \leftrightarrow \forall {\phi}\,(\mathrm{Pos}({\phi}) \to {\phi}z)$ \\
1720 & l.~21 & \textit{for} opopsite \textit{read} opposite \\
1721 & l.~23 & \textit{for} Lemma~3015 of Volume~X \textit{read} Lemma~3013 of Volume~X \\
1721 & l.~7 from foot & \textit{for} Hajek (2002) \textit{read} Hájek (2002) \\
1722 & l.~23 & \textit{for} p.~1618 \textit{read} p.~1617 (Proposition~6.10) \\
1722 & l.~39 & \textit{for} p.~1456 \textit{read} p.~1457 (Proposition~3.3) \\
1723 & l.~30 & \textit{for} disptue \textit{read} dispute \\
1723 & l.~39 & \textit{for} $\mathrm{MG}(x,u) \wedge v'Ru \to \mathrm{MG}(x,u)$ \textit{read} $\mathrm{MG}(x,v') \wedge v'Ru \to \mathrm{MG}(x,u)$ \\
1723 & n.~1 & \textit{for} Fittng (2002) \textit{read} Fitting (2002) \\
1724 & l.~6 & \textit{for} p.~1543 \textit{read} p.~1553 (Proposition~5.6) \\
1724 & l.~16 & \textit{for} invitable \textit{read} inevitable \\
1724 & l.~28 & \textit{for} $\mathrm{MG}(x,v) \wedge wRv \to \mathrm{MG}(x,w)$ \textit{read} $\mathrm{MG}(x,v) \wedge vRw \to \mathrm{MG}(x,w)$ \\
1724 & l.~29 & \textit{for} p.~1598 \textit{read} p.~1589 (Corollary~6.3) \\
1725 & l.~2 & \textit{for} inconsitsent \textit{read} inconsistent \\
\dag\,1725 & proof of Lemma~8.19, (8.32) & \textit{for} $\mathrm{Pos}( {\phi}) \leftrightarrow \neg\mathrm{Pos}({\phi})$ \textit{read} $\mathrm{Pos}(\neg {\phi}) \leftrightarrow \neg\mathrm{Pos}({\phi})$ \\
\dag\,1725 & proof of Lemma~8.19, l.~15 & \textit{for} Definition~8.12 \textit{read} Definition~8.15 \\
\dag\,1725 & proof of Lemma~8.19, l.~24 & \textit{for} Definition~8.8 \textit{read} Definition~8.6; the step uses Definition~8.6 (p.~1694) \\
\dag\,1725 & proof of Lemma~8.19, l.~27 & \textit{for} Definition~8.14 \textit{read} Definition~8.12 \\
\dag\,1726 & proof of Lemma~8.19, l.~17 & \textit{for} Lemma~2752 of Volume~X \textit{read} Lemma~2756 of Volume~X \\
\dag\,1726 & proof of Lemma~8.19, (8.33) & \textit{for} $\mathrm{MG}(z,w_2) \iff \forall u\,(uRw_2 \to E(z,u) \wedge \mathrm{ME}(z,u))$ \textit{read} $\mathrm{MG}(z,w_2) \iff \forall u\,(w_2Ru \to E(z,u) \wedge \mathrm{ME}(z,u))$ \\
1727 & l.~8 & \textit{after} great-making properties \textit{insert} jointly \\
\end{errlist}

\subsection*{Chapter~9. The coherence of maximal excellence: summary \hfill \normalfont\small pp.~1729--1757}
\begin{errlist}
1730 & l.~10 & \textit{for} ovbious \textit{read} obvious \\
1730 & l.~17 & \textit{for} argeud \textit{read} argued \\
1730 & l.~23 & \textit{for} omnsicience \textit{read} omniscience \\
1731 & l.~27 & \textit{for} Volume~VI, Chapter~7 \textit{read} Volume~VI, Chapter~5 (the noseeum inference) \\
1732 & l.~7 & \textit{for} elsewherre \textit{read} elsewhere \\
1732 & l.~33 & \textit{for} Volume~VI, Chapter~7 \textit{read} Volume~VI, Chapter~6 (divine hiddenness) \\
1733 & l.~6 & \textit{for} Volume~VII, Chapter~6 \textit{read} Volume~VII, Chapter~5 (evolutionary debunking) \\
1733 & l.~16 & \textit{for} inedpendence \textit{read} independence \\
1734 & l.~1 & \textit{for} Lemma~2906 of Volume~X \textit{read} Lemma~2903 of Volume~X; the result cited is on p.~3859 \\
1734 & l.~11 & \textit{for} mentoin \textit{read} mention \\
1734 & l.~26 & \textit{for} anallogous \textit{read} analogous \\
1736 & l.~1 & \textit{for} especialy \textit{read} especially \\
1737 & l.~23 & \textit{for} relevnat \textit{read} relevant \\
1737 & l.~29 & \textit{for} regrded \textit{read} regarded \\
1739 & l.~5 & \textit{for} empircial \textit{read} empirical \\
1739 & l.~22 & \textit{for} mteaphysical \textit{read} metaphysical \\
1740 & l.~2 from foot & \textit{for} Lemma~2968 of Volume~X \textit{read} Lemma~2966 of Volume~X \\
1741 & l.~17 & \textit{for} premisse \textit{read} premise \\
1741 & l.~1 from foot & \textit{for} Volume~VI, Chapter~3 \textit{read} Volume~VI, Chapter~4 (the evidential problem of evil) \\
1742 & l.~13 & \textit{for} strnegth \textit{read} strength \\
1742 & l.~20 & \textit{for} proced \textit{read} proceed \\
1742 & l.~25 & \textit{for} Lemma~2988 of Volume~X \textit{read} Lemma~2994 of Volume~X \\
1743 & l.~2 from foot & \textit{for} Lemma~2655 of Volume~X \textit{read} Lemma~2657 of Volume~X \\
\dag\,1743 & proof of Lemma~9.3, l.~40 & \textit{for} Lemma~2769 of Volume~X \textit{read} Lemma~2773 of Volume~X; the result cited is on p.~3846 \\
1744 & l.~16 & \textit{for} p.~1474 \textit{read} p.~1484 (Lemma~3.12) \\
1745 & l.~6 & \textit{for} p.~1724 \textit{read} p.~1725 (Definition~8.18) \\
1745 & l.~36 & \textit{for} controevrsial \textit{read} controversial \\
1746 & l.~15 & \textit{for} fundamantal \textit{read} fundamental \\
1746 & l.~19 & \textit{for} Lemma~2760 of Volume~X \textit{read} Lemma~2755 of Volume~X; the result cited is on p.~3844 \\
1746 & l.~32 & \textit{for} Lemma~2778 of Volume~X \textit{read} Lemma~2774 of Volume~X \\
1747 & l.~6 & \textit{for} recongize \textit{read} recognize \\
1747 & l.~11 & \textit{for} persence \textit{read} presence \\
1747 & l.~17 & \textit{for} prniciple \textit{read} principle \\
1747 & l.~19 & \textit{for} Lemma~2792 of Volume~X \textit{read} Lemma~2796 of Volume~X; the result cited is on p.~3848 \\
1748 & l.~14 & \textit{for} clrify \textit{read} clarify \\
1748 & l.~15 & \textit{for} disgareement \textit{read} disagreement \\
1748 & l.~22 & \textit{for} cateogry \textit{read} category \\
1748 & l.~36 & \textit{for} deficinecy \textit{read} deficiency \\
1749 & l.~32 & \textit{for} p.~1417 \textit{read} p.~1418 (Proposition~2.2) \\
1749 & l.~1 from foot & \textit{for} Volume~IX, Chapter~6 \textit{read} Volume~IX, Chapter~4 (creation and conservation) \\
1751 & l.~26 & \textit{for} crterion \textit{read} criterion \\
1752 & l.~3 & \textit{for} p.~1697 \textit{read} p.~1679 (Lemma~8.2) \\
1752 & l.~5 & \textit{for} cicrumstance \textit{read} circumstance \\
1753 & l.~14 & \textit{for} Volume~III, Chapter~8 \textit{read} Volume~III, Chapter~6 (the converse premise) \\
1753 & l.~20 & \textit{for} naemly \textit{read} namely \\
1754 & l.~11 & \textit{for} matrial \textit{read} material \\
1754 & l.~6 from foot & \textit{for} systmeatic \textit{read} systematic \\
1756 & l.~20 & \textit{for} Lemma~2917 of Volume~X \textit{read} Lemma~2911 of Volume~X \\
1756 & l.~2 from foot & \textit{for} p.~1436 \textit{read} p.~1435 (Lemma~2.10) \\
1757 & l.~8 & \textit{for} wehreas \textit{read} whereas \\
1757 & l.~26 & \textit{for} moreovr \textit{read} moreover \\
\end{errlist}


\section{Volume V: Objections from Logic and Semantics}\label{s:vol5}
% generated by generate.py -- do not edit
\subsection*{Front matter \hfill \normalfont\small pp.~1758--1760}
\begin{errlist}
1759 & Contents, Chapter~5 & \textit{for} 1912 \textit{read} 1913 \\
\end{errlist}

\subsection*{Chapter~1. Introduction \hfill \normalfont\small pp.~1761--1769}
\begin{errlist}
1762 & l.~27 & \textit{for} Volume~VI, Chapter~4 \textit{read} Volume~VI, Chapter~2 (the logical problem of evil) \\
1762 & l.~30 & \textit{for} subtsantive \textit{read} substantive \\
1763 & l.~35 & \textit{for} accommoadte \textit{read} accommodate \\
1764 & l.~14 & \textit{for} delliberate \textit{read} deliberate \\
1765 & l.~25 & \textit{for} hypohteses \textit{read} hypotheses \\
1765 & l.~28 & \textit{for} Lemma~2687 of Volume~X \textit{read} Lemma~2684 of Volume~X \\
1766 & l.~13 & \textit{for} puropse \textit{read} purpose \\
1767 & l.~14 & \textit{for} exteranl \textit{read} external \\
1767 & l.~36 & \textit{for} Definition~2.1 of Volume~IV \textit{read} Definition~2.3 of Volume~IV \\
1768 & l.~13 & \textit{for} accommodatte \textit{read} accommodate \\
1768 & l.~26 & \textit{delete} the second occurrence of \emph{a} \\
1768 & l.~30 & \textit{for} spcific \textit{read} specific \\
1768 & l.~32 & \textit{for} copmlete \textit{read} complete \\
1768 & l.~4 from foot & \textit{for} rigoorus \textit{read} rigorous \\
1769 & l.~1 & \textit{for} ocasion \textit{read} occasion \\
1769 & l.~7 from foot & \textit{for} suggession \textit{read} suggestion \\
\end{errlist}

\subsection*{Chapter~2. Existence as a predicate \hfill \normalfont\small pp.~1770--1811}
\begin{errlist}
1771 & l.~27 & \textit{for} fundmaental \textit{read} fundamental \\
1772 & l.~2 & \textit{for} excluison \textit{read} exclusion \\
1772 & l.~36 & \textit{for} imediate \textit{read} immediate \\
1773 & l.~8 & \textit{for} presnece \textit{read} presence \\
1773 & l.~15 & \textit{after} exists \textit{insert} at $u$ \\
1774 & l.~32 & \textit{for} Lemma~3041 of Volume~X \textit{read} Lemma~3037 of Volume~X; the result cited is on p.~3873 \\
1776 & l.~11 & \textit{for} eexrcise \textit{read} exercise \\
1776 & l.~15 & \textit{for} Lemma~3011 of Volume~X \textit{read} Lemma~3013 of Volume~X; the result cited is on p.~3870 \\
1776 & l.~24 & \textit{for} indicatte \textit{read} indicate \\
1777 & display (2.5) & \textit{for} $\nec(\exists y\,\mathrm{MG}(y) \leftrightarrow E(g))$ \textit{read} $\nec(\exists y\,\mathrm{MG}(y) \to E(g))$; the formula is used in its correct form on p.~1779 \\
1777 & l.~37 & \textit{for} cnostant \textit{read} constant \\
1778 & l.~14 & \textit{for} cnonection \textit{read} connection \\
1779 & l.~19 & \textit{for} $\pos\exists x\,{\psi}x \to \exists x\,\nec {\psi}x$ \textit{read} $\pos\exists x\,{\psi}x \to \exists x\,\pos {\psi}x$ \\
1779 & l.~21 & \textit{for} motviation \textit{read} motivation \\
1779 & l.~31 & \textit{for} judegment \textit{read} judgement \\
1780 & l.~31 & \textit{for} notoin \textit{read} notion \\
1781 & l.~18 & \textit{for} wihtdrawn \textit{read} withdrawn \\
1781 & l.~25 & \textit{for} transitoin \textit{read} transition \\
1781 & l.~31 & \textit{for} conisderable \textit{read} considerable \\
\dag\,1782 & proof of Lemma~2.3, l.~13 & \textit{for} Definition~2.6 \textit{read} Definition~2.1; the step uses Definition~2.1 (p.~1771) \\
\dag\,1782 & proof of Lemma~2.3, l.~19 & \textit{for} Proposition~2.5 \textit{read} Proposition~2.2; the step uses Proposition~2.2 (p.~1780) \\
1783 & l.~20 & \textit{see} Corrigendum~S9 \\
1784 & l.~3 & \textit{delete} the second occurrence of \emph{that} \\
1784 & l.~18 & \textit{for} actualism \textit{read} necessitism \\
1785 & l.~5 & \textit{for} soltion \textit{read} solution \\
1786 & l.~11 & \textit{for} reaminder \textit{read} remainder \\
1787 & l.~5 & \textit{for} sysetmatic \textit{read} systematic \\
1787 & l.~29 & \textit{for} crtiical \textit{read} critical \\
1787 & l.~40 & \textit{for} benneficial \textit{read} beneficial \\
1788 & l.~6 & \textit{for} Lemma~2979 of Volume~X \textit{read} Lemma~2974 of Volume~X \\
1788 & l.~13 & \textit{for} p.~1782 \textit{read} p.~1781 (Lemma~2.3) \\
1788 & l.~26 & \textit{delete} the second occurrence of \emph{that} \\
1788 & l.~27 & \textit{for} criitcal \textit{read} critical \\
1788 & l.~31 & \textit{for} prospcet \textit{read} prospect \\
1789 & l.~17 & \textit{for} $\exists y\,\pos {\phi}y \leftrightarrow \pos\exists y\,{\phi}y$ \textit{read} $\exists y\,\pos {\phi}y \to \pos\exists y\,{\phi}y$ \\
1789 & l.~3 from foot & \textit{delete} the second occurrence of \emph{a} \\
1790 & display (2.11) & \textit{for} $\nec\forall x\,\pos\exists y\,(y = x)$ \textit{read} $\nec\forall x\,\nec\exists y\,(y = x)$ \\
1790 & display (2.12) & \textit{for} $\nec\forall x\,{\chi}x \leftrightarrow \forall x\,\nec {\chi}x$ \textit{read} $\nec\forall x\,{\chi}x \to \forall x\,\nec {\chi}x$ \\
1790 & l.~21 & \textit{for} cnosensus \textit{read} consensus \\
1790 & l.~31 & \textit{for} Lemma~2710 of Volume~X \textit{read} Lemma~2704 of Volume~X \\
\dag\,1791 & proof of Theorem~2.11, l.~7 & \textit{for} Lemma~2.3 \textit{read} Lemma~2.4 \\
1791 & l.~14 & \textit{delete} the second occurrence of \emph{the} \\
1792 & l.~9 & \textit{delete} the second occurrence of \emph{that} \\
1792 & l.~16 & \textit{for} Kant, A592/B602 \textit{read} Kant, A592/B620 \\
1792 & l.~4 from foot & \textit{for} undemrine \textit{read} undermine \\
1794 & l.~19 & \textit{for} $wRu \to D_{u} \subseteq D_{w}$ \textit{read} $uRw \to D_{u} \subseteq D_{w}$ \\
1794 & l.~20 & \textit{for} evalutae \textit{read} evaluate \\
1795 & display (2.15) & \textit{for} $w_2Ru \to D_{w_1} \subseteq D_{u}$ \textit{read} $w_2Ru \to D_{w_2} \subseteq D_{u}$; the formula is used in its correct form on p.~1798 \\
1795 & l.~16 & \textit{for} Malcolm (1959) \textit{read} Malcolm (1960) \\
1795 & l.~26 & \textit{delete} the second occurrence of \emph{the} \\
1795 & l.~27 & \textit{delete} the second occurrence of \emph{the} \\
1795 & l.~7 from foot & \textit{for} propnoent \textit{read} proponent \\
1797 & display (2.17) & \textit{for} $\nec(\forall z\,\mathrm{MG}(z) \to E(g))$ \textit{read} $\nec(\exists z\,\mathrm{MG}(z) \to E(g))$; correctly printed in (2.20), p.~1802 \\
1797 & l.~3 from foot & \textit{for} omtted \textit{read} omitted \\
1797 & l.~40 & \textit{for} actaulism \textit{read} actualism \\
1798 & l.~20 & \textit{for} refeernce \textit{read} reference \\
1798 & l.~39 & \textit{for} sentnece \textit{read} sentence \\
\dag\,1799 & proof of Proposition~2.12, l.~24 & \textit{for} Definition~2.6 \textit{read} Definition~2.9; the step uses Definition~2.9 (p.~1789) \\
\dag\,1799 & proof of Proposition~2.12, l.~40 & \textit{for} $\forall z\,\nec {\phi}z \to \nec\exists z\,{\phi}z$ \textit{read} $\forall z\,\nec {\phi}z \to \nec\forall z\,{\phi}z$ \\
1799 & n.~5 & \textit{for} Maclolm (1960) \textit{read} Malcolm (1960) \\
1800 & l.~12 & \textit{for} theroem \textit{read} theorem \\
1802 & l.~4 & \textit{for} the converse Barcan formula \textit{read} the Barcan formula \\
1804 & l.~2 from foot & \textit{for} tcehnical \textit{read} technical \\
1805 & display (2.22) & \textit{for} $\nec\forall x\,{\phi}x \leftrightarrow \forall x\,\nec {\phi}x$ \textit{read} $\nec\forall x\,{\phi}x \to \forall x\,\nec {\phi}x$ \\
1807 & l.~7 from foot & \textit{for} conesnsus \textit{read} consensus \\
1808 & l.~1 & \textit{for} p.~1867 \textit{read} p.~1767 (Definition~1.1) \\
\dag\,1808 & proof of Lemma~2.14, l.~7 & \textit{for} Theorem~2.11 \textit{read} Theorem~2.13 \\
1808 & l.~11 & \textit{after} domain \textit{insert} of the world $v'$ \\
1808 & l.~16 & \textit{delete} the second occurrence of \emph{a} \\
1808 & l.~29 & \textit{for} p.~1810 \textit{read} p.~1800 (Theorem~2.13) \\
\dag\,1809 & proof of Lemma~2.14, (2.24) & \textit{for} $uRv \to D_{u} \subseteq D_{u}$ \textit{read} $uRv \to D_{v} \subseteq D_{u}$; the formula is used in its correct form on p.~1811 \\
\dag\,1809 & proof of Lemma~2.14, l.~8 from foot & \textit{for} $E(z,w_0) \leftrightarrow z \notin D_{w_0}$ \textit{read} $E(z,w_0) \leftrightarrow z \in D_{w_0}$ \\
1810 & l.~6 & \textit{delete} the second occurrence of \emph{of} \\
1810 & display (2.25) & \textit{for} $\exists z\,\nec {\chi}z \to \nec\forall z\,{\chi}z$ \textit{read} $\forall z\,\nec {\chi}z \to \nec\forall z\,{\chi}z$ \\
1811 & l.~1 & \textit{for} obsevration \textit{read} observation \\
\end{errlist}

\subsection*{Chapter~3. The question-begging objection \hfill \normalfont\small pp.~1812--1874}
\begin{errlist}
1813 & l.~4 from foot & \textit{for} justifictaion \textit{read} justification \\
1814 & running heads, pp.~1814--1874 (even) & \textit{for} Definitions and existential import \textit{read} The question-begging objection \\
1814 & display (3.2) & \textit{for} $@ \notin \nec\,\exists z\, \mathrm{MG}(z,\cdot)$ \textit{read} $@ \in \nec\,\exists z\, \mathrm{MG}(z,\cdot)$; the formula is used in its correct form on p.~1818 \\
1816 & l.~8 from foot & \textit{for} complte \textit{read} complete \\
1817 & l.~4 & \textit{for} procses \textit{read} process \\
1817 & l.~13 & \textit{for} prmarily \textit{read} primarily \\
1817 & display (3.4) & \textit{for} $\mathrm{NME}(y,w_2) \iff \exists w\,(w_2Rw \to \mathrm{ME}(y,w))$ \textit{read} $\mathrm{NME}(y,w_2) \iff \forall w\,(w_2Rw \to \mathrm{ME}(y,w))$; correctly printed in (3.5), p.~1818 \\
1817 & l.~19 & \textit{for} notihng \textit{read} nothing \\
1818 & l.~11 & \textit{for} irreleavnt \textit{read} irrelevant \\
1818 & l.~18 & \textit{for} occasoinally \textit{read} occasionally \\
1819 & l.~33 & \textit{delete} the second occurrence of \emph{of} \\
\dag\,1820 & proof of Lemma~3.2, (3.7) & \textit{for} $\pos\exists x\,\mathrm{MG}(x) \to \nec\forall x\,\mathrm{MG}(x)$ \textit{read} $\pos\exists x\,\mathrm{MG}(x) \to \nec\exists x\,\mathrm{MG}(x)$ \\
1821 & l.~4 & \textit{after} exists \textit{insert} at every world accessible from $w_1$ \\
1821 & l.~7 & \textit{for} establsih \textit{read} establish \\
1822 & display (3.9) & \textit{for} $(\pos G \wedge \pos\neg G)$ \textit{read} $\neg(\pos G \wedge \pos\neg G)$; correctly printed in (3.10), p.~1825 \\
1823 & l.~7 from foot & \textit{for} decsive \textit{read} decisive \\
1825 & l.~26 & \textit{for} Rowe (1974) \textit{read} Rowe (1975) \\
1825 & l.~3 from foot & \textit{for} rleevance \textit{read} relevance \\
1826 & display (3.12) & \textit{for} $\pos\exists z\,\mathrm{MG}(z) \to \pos\exists z\,\mathrm{MG}(z)$ \textit{read} $\pos\exists z\,\mathrm{MG}(z) \to \nec\exists z\,\mathrm{MG}(z)$; correctly printed in (3.14), p.~1830 \\
1826 & l.~17 & \textit{for} elsewehre \textit{read} elsewhere \\
1827 & l.~8 & \textit{for} concpetion \textit{read} conception \\
1828 & l.~7 & \textit{for} Lemma~2977 of Volume~X \textit{read} Lemma~2983 of Volume~X \\
1828 & l.~22 & \textit{for} rceeive \textit{read} receive \\
1829 & l.~24 & \textit{for} evidant \textit{read} evident \\
1831 & l.~8 & \textit{for} Lemma~2995 of Volume~X \textit{read} Lemma~2997 of Volume~X; the result cited is on p.~3869 \\
1831 & l.~14 & \textit{for} p.~1809 \textit{read} p.~1808 (Lemma~2.14) \\
1832 & l.~23 & \textit{for} Lemma~2751 of Volume~X \textit{read} Lemma~2756 of Volume~X; the result cited is on p.~3844 \\
1832 & l.~31 & \textit{for} exists and is maximally excellent at some world accessible from $w_2$ \textit{read} exists and is maximally excellent at every world accessible from $w_2$ \\
1832 & l.~6 from foot & \textit{delete} the second occurrence of \emph{to} \\
1833 & l.~5 & \textit{for} Lemma~2662 of Volume~X \textit{read} Lemma~2658 of Volume~X \\
1833 & l.~8 from foot & \textit{for} proabbility \textit{read} probability \\
1834 & l.~21 & \textit{delete} the second occurrence of \emph{of} \\
1834 & l.~22 & \textit{after} conceivable \textit{insert} on ideal reflection \\
1834 & display (3.16) & \textit{for} $P(q \mid p) = P(p \mid q)\,P(q)/P(q)$ \textit{read} $P(q \mid p) = P(p \mid q)\,P(q)/P(p)$ \\
1835 & display (3.17) & \textit{for} $\mathrm{NE}(z,w_0) \iff \forall w_2\,(w_0Rw_2 \leftrightarrow E(z,w_2))$ \textit{read} $\mathrm{NE}(z,w_0) \iff \forall w_2\,(w_0Rw_2 \to E(z,w_2))$ \\
1836 & l.~2 & \textit{for} Lemma~2742 of Volume~X \textit{read} Lemma~2737 of Volume~X; the result cited is on p.~3842 \\
1836 & l.~3 from foot & \textit{for} Lemma~2681 of Volume~X \textit{read} Lemma~2682 of Volume~X; the result cited is on p.~3837 \\
1837 & l.~17 & \textit{for} embarrasmsent \textit{read} embarrassment \\
1837 & l.~29 & \textit{after} supported \textit{insert} independently of the conclusion \\
1837 & l.~35 & \textit{after} maximally excellent \textit{insert} at $v'$ \\
1837 & l.~38 & \textit{for} noiton \textit{read} notion \\
1838 & Proposition~3.4 & \textit{see} Corrigendum~S8 \\
\dag\,1839 & proof of Proposition~3.5, l.~30 & \textit{for} Proposition~3.6 \textit{read} Proposition~3.4 \\
1840 & l.~4 & \textit{for} conenction \textit{read} connection \\
\dag\,1842 & proof of Proposition~3.6, l.~19 & \textit{for} Definition~3.7 \textit{read} Definition~3.3 \\
1843 & l.~17 & \textit{for} guaratnee \textit{read} guarantee \\
\dag\,1843 & proof of Proposition~3.6, l.~23 & \textit{for} Proposition~3.4 \textit{read} Proposition~3.5; the step uses Proposition~3.5 (p.~1838) \\
1844 & display (3.22) & \textit{for} $\mathrm{MG}(y,w) \iff \forall w\,(uRw \to E(y,w) \wedge \mathrm{ME}(y,w))$ \textit{read} $\mathrm{MG}(y,u) \iff \forall w\,(uRw \to E(y,w) \wedge \mathrm{ME}(y,w))$; correctly printed in (3.24), p.~1849 \\
1845 & l.~25 & \textit{for} dpeendent \textit{read} dependent \\
1845 & l.~1 from foot & \textit{for} sligthly \textit{read} slightly \\
1846 & l.~2 & \textit{for} comaprable \textit{read} comparable \\
1846 & l.~17 & \textit{for} vraiation \textit{read} variation \\
1847 & l.~14 & \textit{for} smymetric \textit{read} symmetric \\
1848 & l.~4 & \textit{for} idetnify \textit{read} identify \\
\dag\,1848 & proof of Theorem~3.9, l.~30 & \textit{for} Proposition~3.5 \textit{read} Proposition~3.6; the step uses Proposition~3.6 (p.~1842) \\
1850 & l.~33 & \textit{for} Lemma~2953 of Volume~X \textit{read} Lemma~2949 of Volume~X \\
1850 & l.~2 from foot & \textit{for} $\pos\neg\exists z\,\mathrm{MG}(z) \to \neg\forall z\,\mathrm{MG}(z)$ \textit{read} $\pos\neg\exists z\,\mathrm{MG}(z) \to \neg\exists z\,\mathrm{MG}(z)$ \\
1851 & l.~6 & \textit{for} insatnce \textit{read} instance \\
1851 & l.~14 & \textit{for} p.~1798 \textit{read} p.~1789 (Definition~2.9) \\
1851 & l.~16 & \textit{for} defedned \textit{read} defended \\
1851 & l.~22 & \textit{for} Lemma~2981 of Volume~X \textit{read} Lemma~2976 of Volume~X \\
1851 & l.~3 from foot & \textit{for} contrsat \textit{read} contrast \\
1852 & l.~12 & \textit{for} suggesiton \textit{read} suggestion \\
1852 & l.~22 & \textit{for} controevrsy \textit{read} controversy \\
1853 & l.~2 & \textit{for} accrding \textit{read} according \\
1853 & l.~1 from foot & \textit{for} tohrough \textit{read} thorough \\
1854 & l.~17 & \textit{for} conidtional \textit{read} conditional \\
1854 & l.~19 & \textit{for} literautre \textit{read} literature \\
1855 & l.~19 & \textit{for} ctaegory \textit{read} category \\
1855 & l.~29 & \textit{for} p.~1768 \textit{read} p.~1767 (Corollary~1.3) \\
1855 & l.~8 from foot & \textit{for} Volume~I, Chapter~7 \textit{read} Volume~I, Chapter~8 (the comparison with Gödel's argument) \\
1855 & display (3.28) & \textit{for} $P(\pos G) < P(\pos\neg G)$ \textit{read} $P(\pos G) > P(\pos\neg G)$ \\
1856 & l.~14 & \textit{for} contrvoersial \textit{read} controversial \\
1856 & l.~25 & \textit{for} Lemma~2807 of Volume~X \textit{read} Lemma~2806 of Volume~X; the result cited is on p.~3849 \\
1858 & l.~1 & \textit{delete} the second occurrence of \emph{is} \\
1858 & l.~26 & \textit{for} p.~1855 \textit{read} p.~1845 (Lemma~3.8) \\
1858 & l.~28 & \textit{for} occasionnally \textit{read} occasionally \\
1858 & l.~8 from foot & \textit{for} evluate \textit{read} evaluate \\
1859 & l.~19 & \textit{for} techincal \textit{read} technical \\
1859 & l.~20 & \textit{for} exrecise \textit{read} exercise \\
1860 & l.~1 & \textit{for} irrleevant \textit{read} irrelevant \\
1860 & l.~17 & \textit{for} intneded \textit{read} intended \\
\dag\,1860 & proof of Lemma~3.13, (3.30) & \textit{for} $\mathrm{MG}(x,w') \leftrightarrow \mathrm{NE}(x,w') \wedge \mathrm{NME}(x,w')$ \textit{read} $\mathrm{MG}(x,w') \to \mathrm{NE}(x,w') \wedge \mathrm{NME}(x,w')$ \\
1860 & l.~36 & \textit{for} mnaifest \textit{read} manifest \\
\dag\,1861 & proof of Lemma~3.13, l.~5 & \textit{for} Lemma~3.8 \textit{read} Lemma~3.11; the step uses Lemma~3.11 (p.~1853) \\
1861 & l.~11 & \textit{for} symetry \textit{read} symmetry \\
1862 & l.~7 & \textit{for} lietrature \textit{read} literature \\
1862 & l.~15 & \textit{for} $\mathrm{MG}(z,u) \wedge vRu \to \mathrm{MG}(z,v)$ \textit{read} $\mathrm{MG}(z,u) \wedge uRv \to \mathrm{MG}(z,v)$ \\
1863 & l.~3 & \textit{for} $\mathrm{NME}(y,w') \iff \forall w'\,(w'Ru \to \mathrm{ME}(y,u))$ \textit{read} $\mathrm{NME}(y,w') \iff \forall u\,(w'Ru \to \mathrm{ME}(y,u))$ \\
1863 & l.~16 & \textit{for} dilema \textit{read} dilemma \\
1863 & l.~7 from foot & \textit{for} Lemma~2710 of Volume~X \textit{read} Lemma~2708 of Volume~X \\
1866 & l.~13 & \textit{delete} the second occurrence of \emph{to} \\
\dag\,1866 & proof of Lemma~3.14, l.~15 & \textit{for} Lemma~2731 of Volume~X \textit{read} Lemma~2733 of Volume~X \\
1867 & l.~5 & \textit{for} motvation \textit{read} motivation \\
1867 & l.~7 & \textit{for} Lemma~2708 of Volume~X \textit{read} Lemma~2711 of Volume~X; the result cited is on p.~3840 \\
\dag\,1867 & proof of Lemma~3.14, (3.33) & \textit{for} $@ \in \nec\,\forall z\, \mathrm{MG}(z,\cdot)$ \textit{read} $@ \in \nec\,\exists z\, \mathrm{MG}(z,\cdot)$ \\
1868 & l.~5 & \textit{for} siginficant \textit{read} significant \\
\dag\,1870 & proof of Lemma~3.17, l.~24 & \textit{for} Definition~3.3 \textit{read} Definition~3.1 \\
1870 & l.~7 from foot & \textit{for} intiution \textit{read} intuition \\
\dag\,1871 & proof of Lemma~3.17, l.~13 & \textit{for} Lemma~3.13 \textit{read} Lemma~3.8 \\
1871 & l.~18 & \textit{for} hypthesis \textit{read} hypothesis \\
1871 & l.~19 & \textit{for} spceific \textit{read} specific \\
\dag\,1871 & proof of Corollary~3.19, l.~20 & \textit{for} Lemma~3.14 \textit{read} Lemma~3.11 \\
1872 & l.~6 & \textit{for} p.~1760 \textit{read} p.~1860 (Lemma~3.13) \\
1872 & l.~10 & \textit{for} $\mathrm{MG}(z,w_2) \iff \forall w_0\,(w_2Rw_0 \leftrightarrow E(z,w_0) \wedge \mathrm{ME}(z,w_0))$ \textit{read} $\mathrm{MG}(z,w_2) \iff \forall w_0\,(w_2Rw_0 \to E(z,w_0) \wedge \mathrm{ME}(z,w_0))$ \\
1872 & l.~27 & \textit{for} p.~1860 \textit{read} p.~1870 (Lemma~3.16) \\
1873 & l.~14 & \textit{for} nothnig \textit{read} nothing \\
1873 & l.~24 & \textit{for} p.~1780 \textit{read} p.~1870 (Definition~3.18) \\
\dag\,1873 & proof of Proposition~3.20, (3.38) & \textit{for} $\exists x\,\mathrm{MG}(x,w') \leftrightarrow \exists x\,(E(x,w') \wedge \mathrm{ME}(x,w'))$ \textit{read} $\exists x\,\mathrm{MG}(x,w') \to \exists x\,(E(x,w') \wedge \mathrm{ME}(x,w'))$ \\
1874 & l.~9 & \textit{for} conseuqence \textit{read} consequence \\
1874 & l.~22 & \textit{for} p.~1782 \textit{read} p.~1872 (Proposition~3.20) \\
1874 & l.~39 & \textit{for} abstrct \textit{read} abstract \\
\end{errlist}

\subsection*{Chapter~4. Definitions and existential import \hfill \normalfont\small pp.~1875--1912}
\begin{errlist}
1876 & running heads, pp.~1876--1912 (even) & \textit{for} The Barcan formula \textit{read} Definitions and existential import \\
1877 & l.~11 & \textit{for} $\mathrm{MG}(y,v) \iff \exists v'\,(vRv' \to E(y,v') \wedge \mathrm{ME}(y,v'))$ \textit{read} $\mathrm{MG}(y,v) \iff \forall v'\,(vRv' \to E(y,v') \wedge \mathrm{ME}(y,v'))$ \\
1877 & l.~13 & \textit{for} faimliar \textit{read} familiar \\
1877 & l.~7 from foot & \textit{for} $\exists x\,\mathrm{MG}(x,w_1) \to \exists x\,(E(x,w_1) \vee \mathrm{ME}(x,w_1))$ \textit{read} $\exists x\,\mathrm{MG}(x,w_1) \to \exists x\,(E(x,w_1) \wedge \mathrm{ME}(x,w_1))$ \\
1879 & l.~24 & \textit{for} surly \textit{read} surely \\
\dag\,1880 & proof of Lemma~4.3, (4.2) & \textit{for} $\nec\forall y\,{\chi}y \to \exists y\,\nec {\chi}y$ \textit{read} $\nec\forall y\,{\chi}y \to \forall y\,\nec {\chi}y$ \\
1880 & l.~30 & \textit{for} Lemma~2756 of Volume~X \textit{read} Lemma~2750 of Volume~X; the result cited is on p.~3844 \\
\dag\,1881 & proof of Lemma~4.3, (4.4) & \textit{for} $\mathrm{ME}(x,w) \iff \mathrm{Op}(x,w) \wedge \mathrm{Os}(x,w) \vee \mathrm{MP}(x,w)$ \textit{read} $\mathrm{ME}(x,w) \iff \mathrm{Op}(x,w) \wedge \mathrm{Os}(x,w) \wedge \mathrm{MP}(x,w)$ \\
\dag\,1881 & proof of Lemma~4.3, (4.3) & \textit{for} $\mathrm{NME}(y,w') \iff \forall w_2\,(w_2Rw' \to \mathrm{ME}(y,w_2))$ \textit{read} $\mathrm{NME}(y,w') \iff \forall w_2\,(w'Rw_2 \to \mathrm{ME}(y,w_2))$ \\
1882 & display (4.5) & \textit{for} $\mathrm{NME}(z,w_1) \iff \forall u\,(uRw_1 \to \mathrm{ME}(z,u))$ \textit{read} $\mathrm{NME}(z,w_1) \iff \forall u\,(w_1Ru \to \mathrm{ME}(z,u))$; the formula is used in its correct form on p.~1883 \\
1884 & display (4.6) & \textit{for} $@ \in \pos\,\forall z\, \mathrm{MG}(z,\cdot)$ \textit{read} $@ \in \pos\,\exists z\, \mathrm{MG}(z,\cdot)$ \\
1885 & l.~26 & \textit{for} codnition \textit{read} condition \\
1885 & l.~5 from foot & \textit{for} strcuture \textit{read} structure \\
1887 & l.~16 & \textit{for} compelte \textit{read} complete \\
1887 & l.~23 & \textit{for} Lemma~2857 of Volume~X \textit{read} Lemma~2853 of Volume~X; the result cited is on p.~3854 \\
1887 & l.~34 & \textit{for} riigdity \textit{read} rigidity \\
1888 & l.~6 & \textit{for} $\nec\forall y\,\nec\forall z\,(z = y)$ \textit{read} $\nec\forall y\,\nec\exists z\,(z = y)$ \\
1888 & display (4.7) & \textit{for} $\exists y\,\pos {\psi}y \leftrightarrow \pos\exists y\,{\psi}y$ \textit{read} $\exists y\,\pos {\psi}y \to \pos\exists y\,{\psi}y$ \\
1889 & l.~1 & \textit{delete} the second occurrence of \emph{that} \\
\dag\,1889 & proof of Corollary~4.4, l.~13 & \textit{for} Lemma~4.7 \textit{read} Lemma~4.2 \\
\dag\,1889 & proof of Corollary~4.4, (4.8) & \textit{for} $\exists y\,\mathrm{MG}(y,w') \to \exists y\,(E(y,w') \vee \mathrm{ME}(y,w'))$ \textit{read} $\exists y\,\mathrm{MG}(y,w') \to \exists y\,(E(y,w') \wedge \mathrm{ME}(y,w'))$ \\
1892 & l.~9 & \textit{for} $\pos\exists y\,\mathrm{MG}(y) \to \nec\forall y\,\mathrm{MG}(y)$ \textit{read} $\pos\exists y\,\mathrm{MG}(y) \to \nec\exists y\,\mathrm{MG}(y)$ \\
1892 & display (4.10) & \textit{for} $\exists y\,\pos {\phi}y \to \nec\exists y\,{\phi}y$ \textit{read} $\exists y\,\pos {\phi}y \to \pos\exists y\,{\phi}y$ \\
1892 & l.~32 & \textit{for} greatnes \textit{read} greatness \\
1893 & display (4.11) & \textit{for} $uRw \to D_{w} \subseteq D_{u}$ \textit{read} $wRu \to D_{w} \subseteq D_{u}$ \\
1895 & l.~1 & \textit{for} Lemma~2940 of Volume~X \textit{read} Lemma~2942 of Volume~X; the result cited is on p.~3863 \\
1895 & l.~9 & \textit{for} $w_0Rw' \to D_{w_0} \subseteq D_{w'}$ \textit{read} $w'Rw_0 \to D_{w_0} \subseteq D_{w'}$ \\
\dag\,1896 & proof of Lemma~4.7, l.~10 & \textit{for} Lemma~3010 of Volume~X \textit{read} Lemma~3016 of Volume~X \\
\dag\,1897 & proof of Lemma~4.7, l.~29 & \textit{for} Lemma~2819 of Volume~X \textit{read} Lemma~2815 of Volume~X \\
1897 & l.~32 & \textit{for} $\mathrm{MG}(y,v') \iff \forall w'\,(w'Rv' \to E(y,w') \wedge \mathrm{ME}(y,w'))$ \textit{read} $\mathrm{MG}(y,v') \iff \forall w'\,(v'Rw' \to E(y,w') \wedge \mathrm{ME}(y,w'))$ \\
1897 & l.~34 & \textit{for} ditsinction \textit{read} distinction \\
\dag\,1897 & proof of Lemma~4.7, (4.12) & \textit{for} $\mathrm{NE}(y,w_0) \iff \forall v'\,(w_0Rv' \leftrightarrow E(y,v'))$ \textit{read} $\mathrm{NE}(y,w_0) \iff \forall v'\,(w_0Rv' \to E(y,v'))$; the formula is used in its correct form on p.~1899 \\
1898 & display (4.13) & \textit{for} $wRw' \to D_{w'} \subseteq D_{w'}$ \textit{read} $wRw' \to D_{w} \subseteq D_{w'}$; correctly printed in (4.14), p.~1899 \\
1898 & l.~36 & \textit{for} Lemma~2854 of Volume~X \textit{read} Lemma~2850 of Volume~X \\
1899 & l.~31 & \textit{for} qeustion \textit{read} question \\
1899 & l.~3 from foot & \textit{for} howevr \textit{read} however \\
1900 & l.~13 & \textit{for} p.~1796 \textit{read} p.~1896 (Lemma~4.7) \\
1900 & l.~14 & \textit{for} inefrence \textit{read} inference \\
1900 & l.~19 & \textit{for} scnario \textit{read} scenario \\
1901 & l.~29 & \textit{for} Lewis (1971) \textit{read} Lewis (1970) \\
1902 & display (4.16) & \textit{for} $\pos\neg\exists y\,\mathrm{MG}(y) \to \exists y\,\mathrm{MG}(y)$ \textit{read} $\pos\neg\exists y\,\mathrm{MG}(y) \to \neg\exists y\,\mathrm{MG}(y)$; correctly printed in (4.17), p.~1907 \\
1902 & l.~27 & \textit{for} alterantive \textit{read} alternative \\
1903 & l.~5 & \textit{for} accessbile \textit{read} accessible \\
\dag\,1904 & proof of Proposition~4.8, l.~6 & \textit{for} Lemma~4.2 \textit{read} Lemma~4.3 \\
\dag\,1904 & proof of Proposition~4.9, l.~22 & \textit{for} Definition~4.10 \textit{read} Definition~4.5; the step uses Definition~4.5 (p.~1894) \\
1905 & l.~2 & \textit{for} empircal \textit{read} empirical \\
\dag\,1905 & proof of Proposition~4.9, l.~4 & \textit{for} Definition~4.5 \textit{read} Definition~4.1; the step uses Definition~4.1 (p.~1876) \\
1906 & l.~25 & \textit{for} ntohing \textit{read} nothing \\
1908 & l.~6 & \textit{for} $\mathrm{MG}(x,v) \wedge uRv \to \mathrm{MG}(x,u)$ \textit{read} $\mathrm{MG}(x,v) \wedge vRu \to \mathrm{MG}(x,u)$ \\
1908 & l.~7 & \textit{for} maxiamlly \textit{read} maximally \\
\dag\,1908 & proof of Theorem~4.11, l.~12 & \textit{for} Proposition~4.12 \textit{read} Proposition~4.8 \\
\dag\,1908 & proof of Proposition~4.12, (4.18) & \textit{for} $\mathrm{NE}(x,w_2) \iff \forall v\,(w_2Rv \leftrightarrow E(x,v))$ \textit{read} $\mathrm{NE}(x,w_2) \iff \forall v\,(w_2Rv \to E(x,v))$ \\
1908 & l.~27 & \textit{for} etxension \textit{read} extension \\
\dag\,1909 & proof of Proposition~4.12, l.~3 & \textit{for} Theorem~4.11 \textit{read} Theorem~4.6; the step uses Theorem~4.6 (p.~1894) \\
1909 & l.~4 & \textit{for} p.~1847 \textit{read} p.~1848 (Theorem~3.9) \\
\dag\,1909 & proof of Corollary~4.13, l.~14 & \textit{for} Proposition~4.9 \textit{read} Proposition~4.8 \\
\dag\,1909 & proof of Corollary~4.13, l.~17 & \textit{for} $\pos\exists z\,{\phi}z \to \exists z\,\nec {\phi}z$ \textit{read} $\pos\exists z\,{\phi}z \to \exists z\,\pos {\phi}z$ \\
1909 & l.~21 & \textit{delete} the second occurrence of \emph{is} \\
1909 & l.~27 & \textit{for} presuppsoition \textit{read} presupposition \\
\dag\,1909 & proof of Corollary~4.13, (4.19) & \textit{for} $E(y,v') \to y \in D_{v'}$ \textit{read} $E(y,v') \leftrightarrow y \in D_{v'}$; correctly printed in (4.20), p.~1910 \\
1910 & l.~5 & \textit{for} motivtion \textit{read} motivation \\
1911 & display (4.22) & \textit{for} $\pos\exists y\,\mathrm{MG}(y) \to \pos\exists y\,\mathrm{MG}(y)$ \textit{read} $\pos\exists y\,\mathrm{MG}(y) \to \nec\exists y\,\mathrm{MG}(y)$ \\
1911 & display (4.23) & \textit{for} $\nec\forall z\,{\phi}z \leftrightarrow \forall z\,\nec {\phi}z$ \textit{read} $\nec\forall z\,{\phi}z \to \forall z\,\nec {\phi}z$; the formula is used in its correct form on p.~1912 \\
1911 & l.~26 & \textit{for} assertoin \textit{read} assertion \\
1911 & display (4.21) & \textit{for} $\pos(\exists x\,\mathrm{MG}(x) \to E(g))$ \textit{read} $\nec(\exists x\,\mathrm{MG}(x) \to E(g))$ \\
1912 & l.~5 & \textit{for} Proposition~5.3 of Volume~III \textit{read} Proposition~5.2 of Volume~III; the result cited is on p.~956 \\
1912 & l.~6 & \textit{for} Lemma~2802 of Volume~X \textit{read} Lemma~2799 of Volume~X \\
\end{errlist}

\subsection*{Chapter~5. Actualism, necessitism and the Barcan formula \hfill \normalfont\small pp.~1913--1967}
\begin{errlist}
1915 & l.~11 & \textit{for} belivee \textit{read} believe \\
1915 & l.~2 from foot & \textit{for} $v \models \pos q \iff \exists w_2\,(vRw_2 \vee w_2 \models q)$ \textit{read} $v \models \pos q \iff \exists w_2\,(vRw_2 \wedge w_2 \models q)$ \\
1916 & l.~23 & \textit{for} formulasion \textit{read} formulation \\
1918 & display (5.1) & \textit{for} $\nec q \wedge \pos r \leftrightarrow \pos(q \wedge r)$ \textit{read} $\nec q \wedge \pos r \to \pos(q \wedge r)$; the formula is used in its correct form on p.~1922 \\
1918 & l.~26 & \textit{for} p.~1780 \textit{read} p.~1870 (Lemma~3.17) \\
1918 & l.~27 & \textit{for} $v'Rw_1$ by transitivity \textit{read} $v'Rw_1$ by symmetry \\
1919 & l.~29 & \textit{for} Lemma~2653 of Volume~X \textit{read} Lemma~2651 of Volume~X; the result cited is on p.~3834 \\
1919 & l.~36 & \textit{for} p.~1930 \textit{read} p.~1903 (Proposition~4.8) \\
1921 & l.~6 & \textit{for} Lemma~2808 of Volume~X \textit{read} Lemma~2802 of Volume~X \\
1921 & l.~13 & \textit{for} p.~1895 \textit{read} p.~1894 (Definition~4.5) \\
1921 & l.~28 & \textit{for} Barcan (1945) \textit{read} Barcan (1946) \\
\dag\,1922 & proof of Theorem~5.2, (5.2) & \textit{for} $\nec(q \wedge p) \leftrightarrow (\pos q \wedge \nec p)$ \textit{read} $\nec(q \wedge p) \leftrightarrow (\nec q \wedge \nec p)$ \\
1923 & l.~12 & \textit{for} apaprent \textit{read} apparent \\
1924 & display (5.3) & \textit{for} $\nec\forall z\,\nec\exists z\,(x = z)$ \textit{read} $\nec\forall z\,\nec\exists x\,(x = z)$ \\
1924 & l.~21 & \textit{for} $vRw' \to D_{v} \subseteq D_{w'}$ \textit{read} $w'Rv \to D_{v} \subseteq D_{w'}$ \\
1924 & l.~32 & \textit{for} straightforwrd \textit{read} straightforward \\
1924 & l.~7 from foot & \textit{for} sonudness \textit{read} soundness \\
1924 & l.~3 from foot & \textit{for} p.~1792 \textit{read} p.~1791 (Theorem~2.11) \\
1926 & Lemma~5.3 & \textit{see} Corrigendum~S10 \\
1927 & l.~3 & \textit{for} Baracn (1946) \textit{read} Barcan (1946) \\
\dag\,1927 & proof of Lemma~5.3, l.~7 & \textit{for} $\nec(p \vee r) \leftrightarrow (\pos p \vee \pos r)$ \textit{read} $\pos(p \vee r) \leftrightarrow (\pos p \vee \pos r)$ \\
1927 & l.~11 & \textit{for} delibearte \textit{read} deliberate \\
1928 & display (5.4) & \textit{for} $\pos r \leftrightarrow \nec\pos r$ \textit{read} $\pos r \to \nec\pos r$ \\
1928 & l.~10 & \textit{delete} the second occurrence of \emph{is} \\
1928 & l.~27 & \textit{for} strutcure \textit{read} structure \\
1929 & display (5.5) & \textit{for} $(vRv' \wedge v'Rw_1) \to vRw_1$ \textit{read} $(v'Rv \wedge v'Rw_1) \to vRw_1$ \\
1929 & display (5.6) & \textit{for} $\pos\exists z\,{\phi}z \leftrightarrow \exists z\,\pos {\phi}z$ \textit{read} $\pos\exists z\,{\phi}z \to \exists z\,\pos {\phi}z$; correctly printed in (5.14), p.~1935 \\
1929 & display (5.7) & \textit{for} $w_1Rw_2 \wedge w'Rw_2 \to w_1Rw'$ \textit{read} $w_1Rw_2 \wedge w_2Rw' \to w_1Rw'$; the formula is used in its correct form on p.~1932 \\
1929 & n.~9 & \textit{for} Hughs and Cresswell (1996) \textit{read} Hughes and Cresswell (1996) \\
\dag\,1930 & proof of Lemma~5.4, (5.8) & \textit{for} $\exists x\,\pos {\chi}x \leftrightarrow \pos\exists x\,{\chi}x$ \textit{read} $\exists x\,\pos {\chi}x \to \pos\exists x\,{\chi}x$; correctly printed in (5.12), p.~1932 \\
\dag\,1930 & proof of Lemma~5.4, (5.9) & \textit{for} $w_1 \models \nec r \iff \forall w_0\,(w_0Rw_1 \to w_0 \models r)$ \textit{read} $w_1 \models \nec r \iff \forall w_0\,(w_1Rw_0 \to w_0 \models r)$ \\
\dag\,1931 & proof of Lemma~5.4, (5.11) & \textit{for} $(uRv \vee uRw_2) \to vRw_2$ \textit{read} $(uRv \wedge uRw_2) \to vRw_2$; correctly printed in (5.15), p.~1935 \\
1931 & l.~27 & \textit{for} p.~1994 \textit{read} p.~1894 (Theorem~4.6) \\
\dag\,1931 & proof of Lemma~5.4, (5.10) & \textit{for} $p \to \nec\nec p$ \textit{read} $p \to \nec\pos p$ \\
1932 & l.~9 & \textit{for} p.~1808 \textit{read} p.~1908 (Proposition~4.12) \\
1932 & l.~22 & \textit{for} cosnistent \textit{read} consistent \\
1933 & display (5.13) & \textit{for} $\pos\pos p \to p$ \textit{read} $\pos\nec p \to p$; the formula is used in its correct form on p.~1936 \\
1933 & l.~3 from foot & \textit{for} tednency \textit{read} tendency \\
\dag\,1935 & proof of Theorem~5.5, l.~19 & \textit{for} $w_2Rw' \leftrightarrow D_{w_2} \subseteq D_{w'}$ \textit{read} $w_2Rw' \to D_{w_2} \subseteq D_{w'}$ \\
\dag\,1935 & proof of Theorem~5.5, l.~21 & \textit{for} Lemma~5.4 \textit{read} Lemma~5.3; the step uses Lemma~5.3 (p.~1926) \\
\dag\,1936 & proof of Corollary~5.6, (5.17) & \textit{for} $\pos\nec p \leftrightarrow p$ \textit{read} $\pos\nec p \to p$ \\
\dag\,1936 & proof of Corollary~5.6, (5.16) & \textit{for} $\nec r \to \nec\pos r$ \textit{read} $\nec r \to \nec\nec r$ \\
\dag\,1937 & proof of Corollary~5.6, l.~5 & \textit{for} Theorem~5.2 \textit{read} Theorem~5.5 \\
1937 & l.~3 from foot & \textit{for} formualtion \textit{read} formulation \\
1938 & l.~19 & \textit{for} referrence \textit{read} reference \\
1939 & l.~7 & \textit{for} catgeory \textit{read} category \\
\dag\,1941 & proof of Corollary~5.9, l.~9 & \textit{for} Theorem~5.10 \textit{read} Theorem~5.5; the step uses Theorem~5.5 (p.~1934) \\
\dag\,1941 & proof of Corollary~5.9, (5.18) & \textit{for} $\pos q \leftrightarrow \nec\neg q$ \textit{read} $\neg\pos q \leftrightarrow \nec\neg q$ \\
1942 & l.~4 & \textit{for} pricnipal \textit{read} principal \\
\dag\,1942 & proof of Corollary~5.9, l.~23 & \textit{for} $\exists x\,\pos {\phi}x \to \pos\forall x\,{\phi}x$ \textit{read} $\exists x\,\pos {\phi}x \to \pos\exists x\,{\phi}x$ \\
1942 & l.~35 & \textit{after} true at every world \textit{insert} accessible from $u$ \\
1943 & l.~20 & \textit{for} oridnary \textit{read} ordinary \\
1944 & l.~31 & \textit{for} gneuine \textit{read} genuine \\
1945 & display (5.19) & \textit{for} $\pos\nec p \to \pos p$ \textit{read} $\pos\pos p \to \pos p$ \\
\dag\,1946 & proof of Theorem~5.10, (5.20) & \textit{for} $\nec(r \vee p) \leftrightarrow (\pos r \vee \pos p)$ \textit{read} $\pos(r \vee p) \leftrightarrow (\pos r \vee \pos p)$; correctly printed in (5.23), p.~1948 \\
\dag\,1947 & proof of Theorem~5.10, (5.21) & \textit{for} $vRw_1 \to D_{w_1} \supseteq D_{v}$ \textit{read} $vRw_1 \to D_{w_1} \subseteq D_{v}$; correctly printed in (5.25), p.~1952 \\
\dag\,1947 & proof of Theorem~5.10, (5.22) & \textit{for} $\pos\forall z\,{\chi}z \to \exists z\,\pos {\chi}z$ \textit{read} $\pos\exists z\,{\chi}z \to \exists z\,\pos {\chi}z$ \\
1947 & l.~16 & \textit{for} p.~1871 \textit{read} p.~1817 (Definition~3.1) \\
\dag\,1947 & proof of Theorem~5.10, l.~21 & \textit{for} $u \models \nec q \iff \forall w_2\,(w_2Ru \to w_2 \models q)$ \textit{read} $u \models \nec q \iff \forall w_2\,(uRw_2 \to w_2 \models q)$ \\
1948 & l.~3 & \textit{for} Lemma~2850 of Volume~X \textit{read} Lemma~2849 of Volume~X \\
1948 & l.~7 & \textit{for} respectievly \textit{read} respectively \\
1948 & l.~23 & \textit{for} $w_2 \models \pos r \iff \exists w\,(w_1Rw \wedge w \models r)$ \textit{read} $w_1 \models \pos r \iff \exists w\,(w_1Rw \wedge w \models r)$ \\
1948 & l.~5 from foot & \textit{for} Lemma~3038 of Volume~X \textit{read} Lemma~3033 of Volume~X \\
1949 & l.~8 & \textit{for} transtiivity \textit{read} transitivity \\
1949 & display (5.24) & \textit{for} $w_0Rw' \leftrightarrow w'Rw_0$ \textit{read} $w_0Rw' \to w'Rw_0$ \\
1950 & l.~15 & \textit{for} thoelogical \textit{read} theological \\
1950 & n.~18 & \textit{for} Hugehs and Cresswell (1996) \textit{read} Hughes and Cresswell (1996) \\
1951 & l.~4 from foot & \textit{for} consequnetly \textit{read} consequently \\
\dag\,1952 & proof of Lemma~5.12, l.~22 & \textit{for} Lemma~2682 of Volume~X \textit{read} Lemma~2686 of Volume~X \\
\dag\,1952 & proof of Lemma~5.12, l.~26 & \textit{for} Corollary~5.6 \textit{read} Corollary~5.9 \\
1954 & l.~3 & \textit{for} devellop \textit{read} develop \\
\dag\,1954 & proof of Proposition~5.13, l.~18 & \textit{for} Definition~5.14 of Volume~III \textit{read} Definition~5.10 of Volume~III \\
\dag\,1954 & proof of Proposition~5.13, (5.26) & \textit{for} $\pos p \to \nec\nec p$ \textit{read} $\pos p \to \nec\pos p$; the formula is used in its correct form on p.~1956 \\
\dag\,1954 & proof of Proposition~5.13, l.~6 from foot & \textit{for} $E(x,w_0) \leftrightarrow x \notin D_{w_0}$ \textit{read} $E(x,w_0) \leftrightarrow x \in D_{w_0}$ \\
\dag\,1955 & proof of Proposition~5.13, l.~12 & \textit{for} $E(z,w_2) \leftrightarrow z \notin D_{w_2}$ \textit{read} $E(z,w_2) \leftrightarrow z \in D_{w_2}$ \\
1955 & l.~19 & \textit{for} clasiscal \textit{read} classical \\
1955 & l.~31 & \textit{for} geunine \textit{read} genuine \\
1955 & l.~40 & \textit{for} exapmle \textit{read} example \\
\dag\,1956 & proof of Corollary~5.15, l.~4 & \textit{for} Proposition~5.11 \textit{read} Proposition~5.13 \\
1956 & l.~6 from foot & \textit{for} ultmately \textit{read} ultimately \\
\dag\,1957 & proof of Proposition~5.16, (5.27) & \textit{for} $\nec\nec p \to \nec p$ \textit{read} $\pos\nec p \to \nec p$; correctly printed in (5.31), p.~1961 \\
1957 & l.~14 & \textit{for} $\nec(p \to p) \to (\nec r \to \nec p)$ \textit{read} $\nec(r \to p) \to (\nec r \to \nec p)$ \\
\dag\,1957 & proof of Proposition~5.16, (5.28) & \textit{for} $\nec q \wedge \pos q \to \pos(p \wedge q)$ \textit{read} $\nec p \wedge \pos q \to \pos(p \wedge q)$; the formula is used in its correct form on p.~1960 \\
1958 & l.~1 & \textit{delete} the second occurrence of \emph{is} \\
1960 & l.~3 & \textit{for} the axiom B, $\nec p \to \nec\nec p$ \textit{read} the axiom 4, $\nec p \to \nec\nec p$ \\
1960 & display (5.30) & \textit{for} $\pos(r \wedge p) \leftrightarrow (\nec r \wedge \nec p)$ \textit{read} $\nec(r \wedge p) \leftrightarrow (\nec r \wedge \nec p)$; the formula is used in its correct form on p.~1963 \\
1964 & l.~3 & \textit{for} a reflexive and transitive frame validates B \textit{read} a reflexive and symmetric frame validates B \\
1964 & l.~20 & \textit{for} Lemma~2699 of Volume~X \textit{read} Lemma~2701 of Volume~X; the result cited is on p.~3839 \\
1965 & display (5.32) & \textit{for} $\nec\forall y\,\nec\forall x\,(x = y)$ \textit{read} $\nec\forall y\,\nec\exists x\,(x = y)$ \\
1965 & l.~7 from foot & \textit{for} indepenedntly \textit{read} independently \\
1966 & display (5.33) & \textit{for} $uRw_1 \to w_2Ru$ \textit{read} $uRw_1 \to w_1Ru$; the formula is used in its correct form on p.~1967 \\
1966 & l.~25 & \textit{for} Fitting and Mendelsohn (1997) \textit{read} Fitting and Mendelsohn (1998) \\
1966 & l.~39 & \textit{for} p.~1995 \textit{read} p.~1959 (Definition~5.17) \\
1967 & l.~2 & \textit{for} symmertic \textit{read} symmetric \\
1967 & l.~17 & \textit{for} impossble \textit{read} impossible \\
1967 & l.~19 & \textit{for} Lemma~3000 of Volume~X \textit{read} Lemma~2998 of Volume~X \\
1967 & l.~40 & \textit{for} permissilbe \textit{read} permissible \\
\end{errlist}

\subsection*{Chapter~6. The S5 objection revisited \hfill \normalfont\small pp.~1968--2028}
\begin{errlist}
1969 & display (6.1) & \textit{for} $\nec r \wedge \pos p \to \pos(p \wedge p)$ \textit{read} $\nec r \wedge \pos p \to \pos(r \wedge p)$; correctly printed in (6.2), p.~1971 \\
1969 & l.~34 & \textit{for} p.~1799 \textit{read} p.~1789 (Definition~2.8) \\
1970 & running heads, pp.~1970--2028 (even) & \textit{for} The Barcan formula \textit{read} The S5 objection revisited \\
1972 & l.~28 & \textit{for} Lemma~2835 of Volume~X \textit{read} Lemma~2837 of Volume~X \\
1973 & l.~30 & \textit{for} p.~1945 \textit{read} p.~1954 (Proposition~5.13) \\
1975 & l.~5 & \textit{for} $\nec(p \to r) \leftrightarrow (\nec p \to \nec r)$ \textit{read} $\nec(p \to r) \to (\nec p \to \nec r)$ \\
1975 & l.~25 & \textit{for} $\pos r \to \nec\nec r$ \textit{read} $\pos r \to \nec\pos r$ \\
1975 & l.~27 & \textit{for} Proposition~5.11 of Volume~III \textit{read} Proposition~5.16 of Volume~III \\
1975 & l.~39 & \textit{for} p.~1798 \textit{read} p.~1799 (Proposition~2.12) \\
1976 & l.~7 & \textit{for} Lemma~2893 of Volume~X \textit{read} Lemma~2899 of Volume~X \\
1977 & l.~3 & \textit{for} Lemma~2973 of Volume~X \textit{read} Lemma~2967 of Volume~X \\
1977 & l.~13 & \textit{for} sceanrio \textit{read} scenario \\
1977 & l.~40 & \textit{for} Lemma~2705 of Volume~X \textit{read} Lemma~2704 of Volume~X; the result cited is on p.~3839 \\
1978 & l.~20 & \textit{for} unifrom \textit{read} uniform \\
1978 & l.~32 & \textit{for} difefrent \textit{read} different \\
1978 & l.~38 & \textit{for} Chandler (1975) \textit{read} Chandler (1976) \\
1979 & l.~3 & \textit{for} neihter \textit{read} neither \\
1979 & l.~12 & \textit{for} accoridng \textit{read} according \\
1980 & l.~20 & \textit{for} $\neg\nec r \leftrightarrow \nec\neg r$ \textit{read} $\neg\pos r \leftrightarrow \nec\neg r$ \\
1980 & l.~21 & \textit{for} concption \textit{read} conception \\
1981 & l.~16 & \textit{for} emlpoyed \textit{read} employed \\
\dag\,1982 & proof of Lemma~6.6, l.~12 & \textit{for} Definition~6.4 \textit{read} Definition~6.1; the step uses Definition~6.1 (p.~1970) \\
1982 & l.~15 & \textit{for} tretament \textit{read} treatment \\
1982 & l.~40 & \textit{for} concievable \textit{read} conceivable \\
1983 & l.~2 & \textit{for} comparalbe \textit{read} comparable \\
1983 & l.~3 & \textit{for} Lemma~2992 of Volume~X \textit{read} Lemma~2987 of Volume~X \\
1983 & l.~10 & \textit{for} Volume~IV, Chapter~4 \textit{read} Volume~IV, Chapter~3 (omnipotence) \\
1983 & l.~27 & \textit{after} S5 conditions \textit{insert} on $R$ \\
\dag\,1983 & proof of Proposition~6.7, (6.6) & \textit{for} $\pos\pos p \leftrightarrow \pos p$ \textit{read} $\pos\pos p \to \pos p$ \\
1985 & l.~2 & \textit{for} clsasical \textit{read} classical \\
1985 & l.~27 & \textit{for} guarante \textit{read} guarantee \\
\dag\,1986 & proof of Proposition~6.8, l.~3 & \textit{for} Definition~6.1 \textit{read} Definition~6.3 \\
\dag\,1986 & proof of Proposition~6.8, (6.9) & \textit{for} $\pos(q \vee p) \leftrightarrow (\pos q \wedge \pos p)$ \textit{read} $\pos(q \vee p) \leftrightarrow (\pos q \vee \pos p)$; correctly printed in (6.13), p.~1991 \\
\dag\,1986 & proof of Proposition~6.8, (6.7) & \textit{for} $vRv' \to vRv'$ \textit{read} $v'Rv \to vRv'$ \\
1986 & l.~2 from foot & \textit{for} ineviatble \textit{read} inevitable \\
1987 & l.~5 & \textit{for} strentgh \textit{read} strength \\
1987 & l.~20 & \textit{for} otherwisse \textit{read} otherwise \\
1988 & l.~19 & \textit{for} Lemma~2695 of Volume~X \textit{read} Lemma~2689 of Volume~X \\
1988 & display (6.10) & \textit{for} $w_0Rw_1 \to w_0Rw_1$ \textit{read} $w_0Rw_1 \to w_1Rw_0$ \\
\dag\,1989 & proof of Corollary~6.10, (6.11) & \textit{for} $w_1 \models \pos r \iff \exists w_0\,(w_2Rw_0 \wedge w_0 \models r)$ \textit{read} $w_2 \models \pos r \iff \exists w_0\,(w_2Rw_0 \wedge w_0 \models r)$ \\
1989 & l.~23 & \textit{for} consistant \textit{read} consistent \\
1990 & l.~27 & \textit{after} accessible \textit{insert} from $v'$ \\
\dag\,1990 & proof of Lemma~6.9, l.~29 & \textit{for} Proposition~6.12 \textit{read} Proposition~6.7 \\
\dag\,1992 & proof of Proposition~6.11, l.~11 & \textit{for} Definition~6.3 \textit{read} Definition~6.4 \\
\dag\,1992 & proof of Proposition~6.12, l.~25 & \textit{for} Lemma~6.2 \textit{read} Lemma~6.6 \\
1993 & l.~6 & \textit{for} Chandler (1977) \textit{read} Chandler (1976) \\
1993 & l.~19 & \textit{for} independantly \textit{read} independently \\
\dag\,1993 & proof of Proposition~6.12, l.~30 & \textit{for} Proposition~6.11 \textit{read} Proposition~6.8 \\
1993 & l.~36 & \textit{for} indiacte \textit{read} indicate \\
1994 & l.~9 & \textit{for} $uRv$ by reflexivity \textit{read} $uRu$ by reflexivity \\
1994 & l.~29 & \textit{for} developmnet \textit{read} development \\
1995 & l.~1 from foot & \textit{delete} the second occurrence of \emph{of} \\
\dag\,1996 & proof of Theorem~6.13, l.~7 & \textit{for} Corollary~6.10 \textit{read} Corollary~6.5 \\
\dag\,1996 & proof of Theorem~6.13, (6.17) & \textit{for} $w_0Rv' \wedge v'Rw' \leftrightarrow w_0Rw'$ \textit{read} $w_0Rv' \wedge v'Rw' \to w_0Rw'$; correctly printed in (6.18), p.~1996 \\
1996 & l.~10 & \textit{for} differant \textit{read} different \\
\dag\,1996 & proof of Theorem~6.13, l.~11 & \textit{for} Lemma~6.9 \textit{read} Lemma~6.6 \\
1996 & l.~23 & \textit{for} prcisely \textit{read} precisely \\
1996 & l.~7 from foot & \textit{for} diffciulty \textit{read} difficulty \\
\dag\,1997 & proof of Theorem~6.13, l.~10 & \textit{for} Proposition~6.8 \textit{read} Proposition~6.12; the step uses Proposition~6.12 (p.~1992) \\
1997 & l.~13 & \textit{for} true at some world accessible from $u$ (the clause for $\nec$) \textit{read} true at every world accessible from $u$ (the clause for $\nec$) \\
1997 & l.~25 & \textit{for} p.~1943 \textit{read} p.~1843 (Definition~3.7) \\
\dag\,2000 & proof of Theorem~6.15, l.~20 & \textit{for} Lemma~3009 of Volume~X \textit{read} Lemma~3011 of Volume~X; the result cited is on p.~3870 \\
2001 & l.~4 & \textit{for} ambgiuous \textit{read} ambiguous \\
2003 & l.~7 & \textit{for} Kripke (1963a) \textit{read} Salmon (1989) (the rejection of S4) \\
2003 & l.~8 & \textit{for} euqivalent \textit{read} equivalent \\
2003 & l.~31 & \textit{for} Kripke (1962) \textit{read} Kripke (1963b) \\
2004 & n.~14 & \textit{for} Krpike (1963b) \textit{read} Kripke (1963b) \\
2005 & l.~17 & \textit{for} Krpke (1963b) \textit{read} Kripke (1963b) \\
2006 & l.~6 & \textit{for} aparent \textit{read} apparent \\
2007 & l.~2 & \textit{for} $vRw_2 \leftrightarrow w_2Rv$ \textit{read} $vRw_2 \to w_2Rv$ \\
2007 & l.~16 & \textit{for} Lemma~3007 of Volume~X \textit{read} Lemma~3012 of Volume~X; the result cited is on p.~3870 \\
2008 & l.~1 from foot & \textit{for} p.~1776 \textit{read} p.~1767 (Lemma~1.2) \\
2009 & l.~8 from foot & \textit{for} Lemma~2772 of Volume~X \textit{read} Lemma~2773 of Volume~X \\
2011 & l.~7 & \textit{for} eventaully \textit{read} eventually \\
2011 & l.~10 & \textit{for} Lemma~2879 of Volume~X \textit{read} Lemma~2875 of Volume~X; the result cited is on p.~3856 \\
2011 & l.~17 & \textit{for} fortiorri \textit{read} fortiori \\
2011 & l.~34 & \textit{for} ordniary \textit{read} ordinary \\
2012 & l.~21 & \textit{for} Salmon (1989) \textit{read} Chandler (1976) (the failure of S4 for metaphysical modality) \\
2012 & l.~32 & \textit{for} $\pos(q \vee r) \to (\pos q \vee \pos r)$ \textit{read} $\pos(q \vee r) \leftrightarrow (\pos q \vee \pos r)$ \\
2012 & l.~7 from foot & \textit{for} Proposition~6.9 of Volume~IV \textit{read} Proposition~6.10 of Volume~IV \\
2012 & l.~1 from foot & \textit{for} $\nec(q \vee p) \leftrightarrow (\nec q \wedge \nec p)$ \textit{read} $\nec(q \wedge p) \leftrightarrow (\nec q \wedge \nec p)$ \\
2013 & l.~27 & \textit{for} p.~2002 \textit{read} p.~2012 (Definition~6.18) \\
2014 & l.~12 & \textit{for} $\nec(r \wedge q) \leftrightarrow (\nec r \vee \nec q)$ \textit{read} $\nec(r \wedge q) \leftrightarrow (\nec r \wedge \nec q)$ \\
2014 & l.~15 & \textit{for} proecss \textit{read} process \\
2014 & l.~19 & \textit{for} dreivation \textit{read} derivation \\
2015 & l.~12 & \textit{for} refererd \textit{read} referred \\
2015 & display (6.28) & \textit{for} $w \models \nec r \iff \forall w'\,(wRw' \leftrightarrow w' \models r)$ \textit{read} $w \models \nec r \iff \forall w'\,(wRw' \to w' \models r)$; the formula is used in its correct form on p.~2019 \\
2015 & l.~37 & \textit{for} aproach \textit{read} approach \\
2016 & l.~25 & \textit{for} $\pos(p \vee q) \leftrightarrow (\pos p \vee \nec q)$ \textit{read} $\pos(p \vee q) \leftrightarrow (\pos p \vee \pos q)$ \\
2017 & l.~12 & \textit{for} coherance \textit{read} coherence \\
2018 & l.~6 & \textit{for} Kripke (1962) \textit{read} Kripke (1963a) \\
\dag\,2019 & proof of Proposition~6.19, l.~10 & \textit{for} Lemma~6.9 \textit{read} Lemma~6.14; the step uses Lemma~6.14 (p.~1999) \\
2019 & l.~18 & \textit{for} $\nec q \leftrightarrow \nec\nec q$ \textit{read} $\nec q \to \nec\nec q$ \\
2019 & l.~7 from foot & \textit{for} persusaive \textit{read} persuasive \\
\dag\,2020 & proof of Proposition~6.19, l.~13 & \textit{for} Proposition~6.11 \textit{read} Proposition~6.7 \\
\dag\,2020 & proof of Proposition~6.19, l.~16 & \textit{for} Theorem~6.15 \textit{read} Theorem~6.13 \\
\dag\,2020 & proof of Proposition~6.19, l.~18 & \textit{for} Lemma~2812 of Volume~X \textit{read} Lemma~2809 of Volume~X \\
2020 & l.~5 from foot & \textit{for} Chndler (1976) \textit{read} Chandler (1976) \\
2021 & l.~7 & \textit{for} dirctly \textit{read} directly \\
2021 & l.~26 & \textit{for} Salmon (1990) \textit{read} Salmon (1989) \\
2022 & display (6.33) & \textit{for} $w_0 \models \pos p \iff \exists u\,(uRw_0 \wedge u \models p)$ \textit{read} $w_0 \models \pos p \iff \exists u\,(w_0Ru \wedge u \models p)$ \\
2022 & l.~16 & \textit{for} Slamon (1989) \textit{read} Salmon (1989) \\
2022 & l.~22 & \textit{for} p.~1771 \textit{read} p.~1871 (Corollary~3.19) \\
2024 & l.~16 & \textit{for} frme \textit{read} frame \\
2026 & l.~11 & \textit{for} Krikpe (1963a) \textit{read} Kripke (1963a) \\
2026 & display (6.36) & \textit{for} $\pos\nec q \leftrightarrow q$ \textit{read} $\pos\nec q \to q$ \\
2027 & display (6.37) & \textit{for} $w' \models \nec q \iff \exists w_1\,(w'Rw_1 \to w_1 \models q)$ \textit{read} $w' \models \nec q \iff \forall w_1\,(w'Rw_1 \to w_1 \models q)$ \\
2028 & l.~15 & \textit{for} essetnial \textit{read} essential \\
\end{errlist}

\subsection*{Chapter~7. Modal collapse \hfill \normalfont\small pp.~2029--2080}
\begin{errlist}
2030 & l.~11 & \textit{for} p.~1779 \textit{read} p.~1789 (Definition~2.10) \\
2032 & l.~10 & \textit{for} Hájek (2003) \textit{read} Hájek (2002) \\
\dag\,2032 & proof of Proposition~7.2, (7.1) & \textit{for} $\mathcal{P}_0 \subseteq \mathcal{P} \Rightarrow \forall x\,\exists w_0\,\forall {\chi} \in \mathcal{P}_0\; {\chi}(x,w_0)$ \textit{read} $\mathcal{P}_0 \subseteq \mathcal{P} \Rightarrow \exists x\,\exists w_0\,\forall {\chi} \in \mathcal{P}_0\; {\chi}(x,w_0)$; the formula is used in its correct form on p.~2036 \\
2032 & l.~28 & \textit{for} technicl \textit{read} technical \\
2033 & display (7.3) & \textit{for} $v'Ru \leftrightarrow uRv'$ \textit{read} $v'Ru \to uRv'$; correctly printed in (7.4), p.~2038 \\
2033 & l.~13 & \textit{for} propsoal \textit{read} proposal \\
2033 & display (7.2) & \textit{for} $\mathrm{G}z \leftrightarrow \forall {\chi}\,(\mathrm{Pos}({\chi}) \leftrightarrow {\chi}z)$ \textit{read} $\mathrm{G}z \leftrightarrow \forall {\chi}\,(\mathrm{Pos}({\chi}) \to {\chi}z)$ \\
\dag\,2034 & proof of Corollary~7.3, l.~4 & \textit{for} Proposition~7.5 \textit{read} Proposition~7.2 \\
2034 & l.~33 & \textit{for} resmeblance \textit{read} resemblance \\
2035 & l.~10 & \textit{for} Hájek (2002) \textit{read} Anderson (1990) (the emended positivity axioms) \\
2035 & l.~27 & \textit{for} Lemma~3037 of Volume~X \textit{read} Lemma~3040 of Volume~X \\
2035 & l.~32 & \textit{for} plausable \textit{read} plausible \\
2037 & l.~9 & \textit{for} ntoation \textit{read} notation \\
2037 & l.~21 & \textit{for} metapyhsical \textit{read} metaphysical \\
2037 & l.~26 & \textit{for} consrtuction \textit{read} construction \\
2037 & l.~1 from foot & \textit{for} scehma \textit{read} schema \\
2039 & l.~11 & \textit{for} acocrding \textit{read} according \\
2039 & l.~3 from foot & \textit{for} $vRu$ by symmetry \textit{read} $vRu$ by transitivity \\
2040 & l.~3 & \textit{for} epxlicit \textit{read} explicit \\
2040 & display (7.5) & \textit{for} $\pos\pos r \to \nec r$ \textit{read} $\pos\nec r \to \nec r$ \\
2040 & l.~26 & \textit{for} particlarly \textit{read} particularly \\
2041 & l.~14 & \textit{for} succsesful \textit{read} successful \\
2041 & l.~7 from foot & \textit{for} Lemma~2824 of Volume~X \textit{read} Lemma~2821 of Volume~X \\
\dag\,2042 & proof of Corollary~7.8, l.~5 & \textit{for} Corollary~7.6 \textit{read} Corollary~7.3; the step uses Corollary~7.3 (p.~2034) \\
\dag\,2042 & proof of Corollary~7.8, (7.7) & \textit{for} $(v'Rw_0 \vee v'Rw_1) \to w_0Rw_1$ \textit{read} $(v'Rw_0 \wedge v'Rw_1) \to w_0Rw_1$ \\
\dag\,2043 & proof of Corollary~7.8, (7.9) & \textit{for} $\nec(q \to q) \to (\nec q \to \nec r)$ \textit{read} $\nec(q \to r) \to (\nec q \to \nec r)$ \\
2043 & l.~25 & \textit{for} p.~1880 \textit{read} p.~1780 (Proposition~2.2) \\
\dag\,2043 & proof of Corollary~7.8, (7.8) & \textit{for} $w' \models \pos r \iff \exists w'\,(w'Rv \wedge v \models r)$ \textit{read} $w' \models \pos r \iff \exists v\,(w'Rv \wedge v \models r)$ \\
\dag\,2044 & proof of Proposition~7.9, (7.10) & \textit{for} $\mathrm{Pos}({\phi}) \vee \mathrm{Pos}({\chi}) \to \mathrm{Pos}({\phi} \sqcap {\chi})$ \textit{read} $\mathrm{Pos}({\phi}) \wedge \mathrm{Pos}({\chi}) \to \mathrm{Pos}({\phi} \sqcap {\chi})$; correctly printed in (7.11), p.~2048 \\
2044 & l.~29 & \textit{for} apparntly \textit{read} apparently \\
2046 & l.~9 & \textit{for} $u \models \nec q \iff \forall v\,(uRv \leftrightarrow v \models q)$ \textit{read} $u \models \nec q \iff \forall v\,(uRv \to v \models q)$ \\
2048 & l.~11 & \textit{for} Sobel (1986) \textit{read} Sobel (1987) \\
2049 & display (7.13) & \textit{for} $w_1 \models \nec q \iff \forall v'\,(v'Rw_1 \to v' \models q)$ \textit{read} $w_1 \models \nec q \iff \forall v'\,(w_1Rv' \to v' \models q)$ \\
2050 & l.~10 & \textit{for} commitetd \textit{read} committed \\
2051 & l.~5 & \textit{for} $\nec\nec r \to \nec r$ \textit{read} $\pos\nec r \to \nec r$ \\
2051 & l.~32 & \textit{for} direclty \textit{read} directly \\
\dag\,2053 & proof of Proposition~7.10, l.~18 & \textit{for} Proposition~7.5 \textit{read} Proposition~7.9; the step uses Proposition~7.9 (p.~2044) \\
\dag\,2053 & proof of Proposition~7.10, l.~27 & \textit{for} $\nec\forall x\,\mathrm{G}x$ \textit{read} $\nec\exists x\,\mathrm{G}x$ \\
\dag\,2054 & proof of Proposition~7.10, l.~11 & \textit{for} Corollary~7.6 \textit{read} Corollary~7.8 \\
\dag\,2055 & proof of Lemma~7.11, l.~17 & \textit{for} Proposition~7.10 \textit{read} Proposition~7.9 \\
2055 & l.~21 & \textit{for} Lemma~3039 of Volume~X \textit{read} Lemma~3038 of Volume~X; the result cited is on p.~3873 \\
2056 & l.~15 & \textit{for} ${\chi} \in \mathcal{P} \wedge \nec\forall y\,({\chi}y \leftrightarrow {\psi}y) \Rightarrow {\psi} \in \mathcal{P}$ \textit{read} ${\chi} \in \mathcal{P} \wedge \nec\forall y\,({\chi}y \to {\psi}y) \Rightarrow {\psi} \in \mathcal{P}$ \\
2056 & l.~23 & \textit{for} acconut \textit{read} account \\
\dag\,2056 & proof of Proposition~7.12, (7.14) & \textit{for} $p \to \pos\pos p$ \textit{read} $p \to \nec\pos p$ \\
2056 & l.~3 from foot & \textit{for} Lemma~2929 of Volume~X \textit{read} Lemma~2923 of Volume~X \\
2056 & l.~40 & \textit{for} p.~1837 \textit{read} p.~1836 (Definition~3.3) \\
2057 & l.~5 from foot & \textit{for} substantve \textit{read} substantive \\
2057 & l.~39 & \textit{for} prceisely \textit{read} precisely \\
2059 & l.~40 & \textit{for} $uRw_0 \to uRw_0$ \textit{read} $uRw_0 \to w_0Ru$ \\
2061 & display (7.16) & \textit{for} $\pos(q \vee q) \leftrightarrow (\pos r \vee \pos q)$ \textit{read} $\pos(r \vee q) \leftrightarrow (\pos r \vee \pos q)$ \\
2061 & l.~16 & \textit{for} Lemma~2879 of Volume~X \textit{read} Lemma~2882 of Volume~X \\
2061 & l.~32 & \textit{for} procdure \textit{read} procedure \\
2061 & l.~40 & \textit{for} prvious \textit{read} previous \\
2062 & l.~4 & \textit{for} $\pos\nec r \to \pos r$ \textit{read} $\pos\pos r \to \pos r$ \\
2062 & l.~11 & \textit{for} $(w_2Rw' \wedge w_1Rv') \to w'Rv'$ \textit{read} $(w_2Rw' \wedge w_2Rv') \to w'Rv'$ \\
2062 & display (7.17) & \textit{for} $u \models \pos p \iff \forall w_2\,(uRw_2 \wedge w_2 \models p)$ \textit{read} $u \models \pos p \iff \exists w_2\,(uRw_2 \wedge w_2 \models p)$ \\
2062 & l.~17 & \textit{for} treatemnt \textit{read} treatment \\
2062 & l.~23 & \textit{for} fruther \textit{read} further \\
2062 & l.~26 & \textit{for} inavriant \textit{read} invariant \\
2062 & l.~34 & \textit{for} asessment \textit{read} assessment \\
2065 & l.~14 & \textit{for} Lemma~3006 of Volume~X \textit{read} Lemma~3011 of Volume~X; the result cited is on p.~3870 \\
2065 & l.~19 & \textit{delete} the second occurrence of \emph{to} \\
2065 & display (7.18) & \textit{for} $\nec q \to \pos\nec q$ \textit{read} $\nec q \to \nec\nec q$ \\
2066 & l.~8 from foot & \textit{for} Fitting (2002) \textit{read} Anderson (1990) (the emended positivity axioms) \\
2066 & l.~35 & \textit{for} Lemma~2671 of Volume~X \textit{read} Lemma~2675 of Volume~X \\
2067 & l.~37 & \textit{for} $wRw_1 \wedge w_1Rv \leftrightarrow wRv$ \textit{read} $wRw_1 \wedge w_1Rv \to wRv$ \\
2067 & display (7.20) & \textit{for} $\mathrm{G}y \leftrightarrow \forall {\chi}\,(\mathrm{Pos}({\chi}) \leftrightarrow \pos {\chi}y)$ \textit{read} $\mathrm{G}y \leftrightarrow \forall {\chi}\,(\mathrm{Pos}({\chi}) \leftrightarrow \nec {\chi}y)$ \\
2068 & l.~29 & \textit{for} aruged \textit{read} argued \\
2069 & l.~8 & \textit{for} ${\psi}, {\chi} \notin \mathcal{P} \Rightarrow {\psi} \sqcap {\chi} \in \mathcal{P}$ \textit{read} ${\psi}, {\chi} \in \mathcal{P} \Rightarrow {\psi} \sqcap {\chi} \in \mathcal{P}$ \\
2069 & l.~13 & \textit{for} p.~1975 \textit{read} p.~1957 (Proposition~5.16) \\
2069 & l.~30 & \textit{for} Definition~6.8 of Volume~IV \textit{read} Definition~6.2 of Volume~IV; the result cited is on p.~1589 \\
2070 & l.~11 & \textit{for} conjcture \textit{read} conjecture \\
2070 & display (7.23) & \textit{for} $\pos\pos r \leftrightarrow \pos r$ \textit{read} $\pos\pos r \to \pos r$; correctly printed in (7.26), p.~2073 \\
2071 & l.~15 & \textit{after} great-making properties \textit{insert} jointly \\
2072 & l.~4 & \textit{for} implict \textit{read} implicit \\
2072 & l.~21 & \textit{for} Lemma~2724 of Volume~X \textit{read} Lemma~2729 of Volume~X \\
2072 & display (7.24) & \textit{for} $\nec p \wedge \pos r \to \pos(p \vee r)$ \textit{read} $\nec p \wedge \pos r \to \pos(p \wedge r)$; correctly printed in (7.28), p.~2077 \\
2075 & l.~13 & \textit{for} essecne \textit{read} essence \\
2075 & l.~36 & \textit{for} orginal \textit{read} original \\
\dag\,2076 & proof of Proposition~7.15, l.~11 & \textit{for} Corollary~7.3 \textit{read} Corollary~7.8; the step uses Corollary~7.8 (p.~2042) \\
\dag\,2076 & proof of Proposition~7.15, (7.27) & \textit{for} $\mathrm{Pos}({\psi}) \leftrightarrow \pos\exists y\,{\psi}y$ \textit{read} $\mathrm{Pos}({\psi}) \to \pos\exists y\,{\psi}y$ \\
\dag\,2076 & proof of Proposition~7.15, l.~22 & \textit{for} $\mathrm{Pos}({\phi}) \vee \mathrm{Pos}({\psi}) \to \mathrm{Pos}({\phi} \sqcap {\psi})$ \textit{read} $\mathrm{Pos}({\phi}) \wedge \mathrm{Pos}({\psi}) \to \mathrm{Pos}({\phi} \sqcap {\psi})$ \\
\dag\,2076 & proof of Proposition~7.15, l.~23 & \textit{for} Proposition~7.9 \textit{read} Proposition~7.12 \\
2077 & n.~11 & \textit{for} Sobel (1988) \textit{read} Sobel (1987) \\
2078 & display (7.29) & \textit{for} $vRw \wedge wRu \leftrightarrow vRu$ \textit{read} $vRw \wedge wRu \to vRu$ \\
2079 & l.~5 & \textit{for} acheive \textit{read} achieve \\
2079 & display (7.30) & \textit{for} $\nec(p \wedge q) \to (\nec p \wedge \nec q)$ \textit{read} $\nec(p \wedge q) \leftrightarrow (\nec p \wedge \nec q)$; the formula is used in its correct form on p.~2080 \\
2079 & l.~6 from foot & \textit{for} p.~1908 \textit{read} p.~1909 (Corollary~4.13) \\
2079 & display (7.31) & \textit{for} $\nec(r \wedge q) \to (\nec r \wedge \nec q)$ \textit{read} $\nec(r \wedge q) \leftrightarrow (\nec r \wedge \nec q)$ \\
2079 & l.~1 from foot & \textit{for} completenses \textit{read} completeness \\
2079 & n.~1 & \textit{for} Fitting (2003) \textit{read} Fitting (2002) \\
\end{errlist}

\subsection*{Chapter~8. Semantic objections: rigidity and reference to God \hfill \normalfont\small pp.~2081--2126}
\begin{errlist}
2082 & running heads, pp.~2082--2126 (even) & \textit{for} Modal collapse \textit{read} Semantic objections \\
2082 & l.~2 & \textit{for} disgreement \textit{read} disagreement \\
2082 & l.~5 & \textit{for} Lemma~2893 of Volume~X \textit{read} Lemma~2888 of Volume~X \\
2082 & l.~9 & \textit{for} prblem \textit{read} problem \\
2083 & l.~12 & \textit{for} resistence \textit{read} resistance \\
2083 & l.~31 & \textit{for} Lemma~2916 of Volume~X \textit{read} Lemma~2912 of Volume~X \\
2085 & l.~2 & \textit{for} Lemma~2782 of Volume~X \textit{read} Lemma~2776 of Volume~X \\
2085 & l.~11 & \textit{for} Lemma~2825 of Volume~X \textit{read} Lemma~2823 of Volume~X \\
2085 & l.~12 & \textit{for} riogrous \textit{read} rigorous \\
2086 & l.~15 & \textit{for} p.~2000 \textit{read} p.~2010 (Lemma~6.17) \\
2086 & l.~1 from foot & \textit{for} Lweis (1970) \textit{read} Lewis (1970) \\
\dag\,2087 & proof of Lemma~8.2, l.~14 & \textit{for} Lemma~2816 of Volume~X \textit{read} Lemma~2811 of Volume~X; the result cited is on p.~3850 \\
2087 & l.~24 & \textit{for} valdiity \textit{read} validity \\
2089 & l.~2 & \textit{for} reocmmend \textit{read} recommend \\
2089 & l.~13 & \textit{for} asesssment \textit{read} assessment \\
2090 & display (8.5) & \textit{for} $E(x,w') \leftrightarrow x \notin D_{w'}$ \textit{read} $E(x,w') \leftrightarrow x \in D_{w'}$; the formula is used in its correct form on p.~2091 \\
2090 & display (8.6) & \textit{for} $wRw_1 \to D_{w_2} \subseteq D_{w}$ \textit{read} $wRw_2 \to D_{w_2} \subseteq D_{w}$ \\
2091 & l.~24 & \textit{for} paritcular \textit{read} particular \\
2091 & l.~37 & \textit{for} Lemma~2723 of Volume~X \textit{read} Lemma~2721 of Volume~X \\
2092 & l.~20 & \textit{for} p.~2034 \textit{read} p.~2035 (Corollary~7.6) \\
2092 & display (8.7) & \textit{for} $vRv' \to D_{v} \supseteq D_{v'}$ \textit{read} $vRv' \to D_{v} \subseteq D_{v'}$; correctly printed in (8.9), p.~2095 \\
\dag\,2096 & proof of Lemma~8.5, (8.10) & \textit{for} $\nec\forall x\,{\psi}x \to \exists x\,\nec {\psi}x$ \textit{read} $\nec\forall x\,{\psi}x \to \forall x\,\nec {\psi}x$ \\
\dag\,2096 & proof of Proposition~8.4, l.~26 & \textit{for} Lemma~8.6 \textit{read} Lemma~8.2 \\
2096 & l.~27 & \textit{for} prleiminary \textit{read} preliminary \\
\dag\,2097 & proof of Lemma~8.6, l.~23 & \textit{for} $vRw \to D_{v} \subseteq D_{w}$ \textit{read} $wRv \to D_{v} \subseteq D_{w}$ \\
2097 & l.~25 & \textit{for} Alston (1988) \textit{read} Alston (1989) \\
2097 & l.~26 & \textit{delete} the second occurrence of \emph{the} \\
2097 & l.~39 & \textit{for} distinugish \textit{read} distinguish \\
2098 & l.~6 from foot & \textit{for} Volume~IX, Chapter~5 \textit{read} Volume~VIII, Chapter~5 (divine simplicity) \\
\dag\,2100 & proof of Proposition~8.7, l.~23 & \textit{for} Lemma~8.9 \textit{read} Lemma~8.5 \\
\dag\,2101 & proof of Proposition~8.7, l.~20 & \textit{for} $\nec\exists y\,{\psi}y \to \forall y\,\nec {\psi}y$ \textit{read} $\nec\forall y\,{\psi}y \to \forall y\,\nec {\psi}y$ \\
2101 & l.~22 & \textit{for} Lemma~2870 of Volume~X \textit{read} Lemma~2865 of Volume~X; the result cited is on p.~3855 \\
2102 & l.~11 & \textit{for} remanider \textit{read} remainder \\
2102 & l.~27 & \textit{for} particlar \textit{read} particular \\
2103 & l.~17 & \textit{for} p.~2033 \textit{read} p.~2034 (Lemma~7.4) \\
2103 & display (8.15) & \textit{for} $\nec(\forall y\,\mathrm{MG}(y) \to E(g))$ \textit{read} $\nec(\exists y\,\mathrm{MG}(y) \to E(g))$ \\
2104 & l.~27 & \textit{for} Altson (1989) \textit{read} Alston (1989) \\
2105 & l.~2 & \textit{for} governemnt \textit{read} government \\
2105 & l.~10 & \textit{for} $vRw \to D_{v} \supseteq D_{w}$ \textit{read} $vRw \to D_{v} \subseteq D_{w}$ \\
2105 & l.~20 & \textit{for} exprssion \textit{read} expression \\
\dag\,2105 & proof of Proposition~8.8, l.~29 & \textit{for} Lemma~2890 of Volume~X \textit{read} Lemma~2885 of Volume~X; the result cited is on p.~3857 \\
\dag\,2106 & proof of Lemma~8.9, l.~3 & \textit{for} Proposition~8.7 \textit{read} Proposition~8.4 \\
2106 & l.~34 & \textit{for} p.~1828 \textit{read} p.~1838 (Proposition~3.5) \\
2107 & l.~18 & \textit{for} quessionable \textit{read} questionable \\
2108 & l.~1 & \textit{for} rigidty \textit{read} rigidity \\
2109 & l.~13 & \textit{for} Alston (1990) \textit{read} Alston (1989) \\
2109 & l.~39 & \textit{for} instence \textit{read} instance \\
2110 & l.~6 & \textit{for} farmework \textit{read} framework \\
\dag\,2110 & proof of Proposition~8.11, l.~26 & \textit{for} Lemma~8.9 \textit{read} Lemma~8.6 \\
\dag\,2110 & proof of Proposition~8.11, (8.20) & \textit{for} $\nec\forall y\,\nec\exists x\,(x = x)$ \textit{read} $\nec\forall y\,\nec\exists x\,(x = y)$; the formula is used in its correct form on p.~2113 \\
\dag\,2111 & proof of Proposition~8.11, l.~9 & \textit{for} Definition~8.15 \textit{read} Definition~8.10 \\
2111 & l.~10 & \textit{for} inferrence \textit{read} inference \\
\dag\,2111 & proof of Proposition~8.11, l.~19 & \textit{for} Lemma~8.5 \textit{read} Lemma~8.2 \\
2111 & l.~25 & \textit{for} subsatntive \textit{read} substantive \\
2112 & l.~5 & \textit{for} ohterwise \textit{read} otherwise \\
2113 & l.~8 & \textit{for} Lemma~2747 of Volume~X \textit{read} Lemma~2751 of Volume~X \\
2113 & display (8.21) & \textit{for} $\nec\forall z\,\nec\exists x\,(z = z)$ \textit{read} $\nec\forall z\,\nec\exists x\,(x = z)$ \\
2113 & l.~26 & \textit{for} Lemma~2799 of Volume~X \textit{read} Lemma~2805 of Volume~X; the result cited is on p.~3849 \\
2114 & l.~31 & \textit{for} $\forall y\,\nec {\chi}y \to \pos\forall y\,{\chi}y$ \textit{read} $\forall y\,\nec {\chi}y \to \nec\forall y\,{\chi}y$ \\
2115 & display (8.22) & \textit{for} $\pos\exists z\,{\phi}z \to \forall z\,\pos {\phi}z$ \textit{read} $\pos\exists z\,{\phi}z \to \exists z\,\pos {\phi}z$; the formula is used in its correct form on p.~2118 \\
2116 & l.~6 from foot & \textit{for} caapcity \textit{read} capacity \\
2116 & l.~4 from foot & \textit{for} agruably \textit{read} arguably \\
2117 & l.~2 & \textit{for} previuos \textit{read} previous \\
2117 & l.~11 & \textit{for} intrisnic \textit{read} intrinsic \\
2117 & l.~17 & \textit{for} Lemma~2987 of Volume~X \textit{read} Lemma~2990 of Volume~X; the result cited is on p.~3868 \\
2118 & l.~11 & \textit{for} stateemnt \textit{read} statement \\
2118 & l.~18 & \textit{for} Lemma~2798 of Volume~X \textit{read} Lemma~2794 of Volume~X \\
2120 & l.~14 & \textit{for} considerbale \textit{read} considerable \\
\dag\,2121 & proof of Corollary~8.12, l.~3 & \textit{for} Lemma~8.5 \textit{read} Lemma~8.6; the step uses Lemma~8.6 (p.~2097) \\
2121 & l.~8 & \textit{for} evdent \textit{read} evident \\
\dag\,2121 & proof of Lemma~8.13, (8.25) & \textit{for} $\pos\exists x\,{\phi}x \leftrightarrow \exists x\,\pos {\phi}x$ \textit{read} $\pos\exists x\,{\phi}x \to \exists x\,\pos {\phi}x$ \\
2122 & l.~22 & \textit{for} p.~2016 \textit{read} p.~2106 (Lemma~8.9) \\
\dag\,2122 & proof of Proposition~8.14, l.~27 & \textit{for} Proposition~8.8 \textit{read} Proposition~8.7; the step uses Proposition~8.7 (p.~2100) \\
2122 & l.~8 from foot & \textit{for} $E(x,w') \to x \in D_{w'}$ \textit{read} $E(x,w') \leftrightarrow x \in D_{w'}$ \\
2122 & n.~3 & \textit{for} Lews (1970) \textit{read} Lewis (1970) \\
2123 & l.~1 & \textit{for} p.~2196 \textit{read} p.~2096 (Proposition~8.4) \\
\dag\,2123 & proof of Proposition~8.14, (8.27) & \textit{for} $\forall y\,\pos {\psi}y \to \nec\forall y\,{\psi}y$ \textit{read} $\forall y\,\nec {\psi}y \to \nec\forall y\,{\psi}y$ \\
\dag\,2123 & proof of Proposition~8.14, (8.28) & \textit{for} $\exists z\,\nec {\chi}z \to \pos\exists z\,{\chi}z$ \textit{read} $\exists z\,\pos {\chi}z \to \pos\exists z\,{\chi}z$; the formula is used in its correct form on p.~2125 \\
2124 & l.~25 & \textit{for} contributtion \textit{read} contribution \\
2125 & l.~5 & \textit{for} embarrasment \textit{read} embarrassment \\
2125 & l.~7 from foot & \textit{for} acknolwedge \textit{read} acknowledge \\
\end{errlist}

\subsection*{Chapter~9. Paraconsistent and intuitionistic variants \hfill \normalfont\small pp.~2127--2188}
\begin{errlist}
2128 & l.~29 & \textit{for} quetsion \textit{read} question \\
2129 & l.~14 & \textit{for} delibrate \textit{read} deliberate \\
2129 & l.~29 & \textit{for} exrpession \textit{read} expression \\
2130 & l.~28 & \textit{for} Lemma~2665 of Volume~X \textit{read} Lemma~2664 of Volume~X \\
2131 & l.~21 & \textit{for} Lemma~2909 of Volume~X \textit{read} Lemma~2910 of Volume~X; the result cited is on p.~3860 \\
2132 & l.~13 & \textit{for} $\nec p \wedge \pos q \leftrightarrow \pos(p \wedge q)$ \textit{read} $\nec p \wedge \pos q \to \pos(p \wedge q)$ \\
2132 & l.~39 & \textit{for} coheernce \textit{read} coherence \\
2132 & display (9.3) & \textit{for} $\pos(q \vee r) \leftrightarrow (\pos r \vee \pos r)$ \textit{read} $\pos(q \vee r) \leftrightarrow (\pos q \vee \pos r)$ \\
2133 & l.~21 & \textit{for} Definition~5.10 of Volume~III \textit{read} Definition~5.13 of Volume~III \\
2133 & l.~5 from foot & \textit{for} Blckburn et~al. (2001) \textit{read} Blackburn et~al. (2001) \\
2134 & l.~13 & \textit{for} indepnedently \textit{read} independently \\
2135 & l.~30 & \textit{for} desirbale \textit{read} desirable \\
2135 & l.~3 from foot & \textit{for} probelm \textit{read} problem \\
2135 & l.~39 & \textit{for} instnace \textit{read} instance \\
2136 & l.~5 & \textit{for} possbiility \textit{read} possibility \\
2136 & l.~16 & \textit{for} straightfowrard \textit{read} straightforward \\
2136 & l.~27 & \textit{for} permissbile \textit{read} permissible \\
2137 & l.~15 & \textit{for} Lemma~2750 of Volume~X \textit{read} Lemma~2749 of Volume~X \\
\dag\,2138 & proof of Proposition~9.2, l.~2 & \textit{for} Lemma~2920 of Volume~X \textit{read} Lemma~2923 of Volume~X \\
2138 & l.~32 & \textit{for} p.~2022 \textit{read} p.~2122 (Definition~8.15) \\
2140 & l.~7 & \textit{for} criteroin \textit{read} criterion \\
2140 & l.~17 & \textit{for} absolutly \textit{read} absolutely \\
2140 & l.~26 & \textit{for} Lemma~6.12 of Volume~IV \textit{read} Lemma~6.16 of Volume~IV \\
2140 & l.~29 & \textit{for} p.~1835 \textit{read} p.~1853 (Lemma~3.11) \\
2141 & l.~8 & \textit{for} differnet \textit{read} different \\
2142 & l.~22 & \textit{for} p.~2095 \textit{read} p.~2096 (Lemma~8.5) \\
2143 & display (9.7) & \textit{for} $(wRw' \wedge v'Rw) \to w'Rv'$ \textit{read} $(wRw' \wedge wRv') \to w'Rv'$; the formula is used in its correct form on p.~2146 \\
2143 & display (9.6) & \textit{for} $vRw_1 \wedge w_1Rw_0 \leftrightarrow vRw_0$ \textit{read} $vRw_1 \wedge w_1Rw_0 \to vRw_0$ \\
2144 & l.~7 & \textit{for} $wRv' \leftrightarrow v'Rw$ \textit{read} $wRv' \to v'Rw$ \\
2145 & l.~12 & \textit{for} p.~1979 \textit{read} p.~1989 (Corollary~6.10) \\
2147 & l.~6 & \textit{for} Lemma~2783 of Volume~X \textit{read} Lemma~2782 of Volume~X \\
2147 & l.~33 & \textit{for} contxt \textit{read} context \\
2147 & l.~5 from foot & \textit{for} schemma \textit{read} schema \\
2147 & n.~4 & \textit{for} Chelals (1980) \textit{read} Chellas (1980) \\
\dag\,2149 & proof of Proposition~9.8, l.~3 & \textit{for} Definition~9.4 \textit{read} Definition~9.5 \\
2149 & l.~5 & \textit{for} suondness \textit{read} soundness \\
\dag\,2149 & proof of Lemma~9.7, l.~18 & \textit{for} Definition~9.6 \textit{read} Definition~9.1 \\
2149 & l.~19 & \textit{for} semantcs \textit{read} semantics \\
2149 & l.~21 & \textit{for} diivne \textit{read} divine \\
\dag\,2150 & proof of Proposition~9.8, l.~6 & \textit{for} Definition~9.10 \textit{read} Definition~9.6; the step uses Definition~9.6 (p.~2149) \\
2150 & l.~20 & \textit{for} embarrassmnet \textit{read} embarrassment \\
2150 & l.~22 & \textit{for} neecssity \textit{read} necessity \\
\dag\,2150 & proof of Lemma~9.9, l.~29 & \textit{for} Definition~9.10 \textit{read} Definition~9.5 \\
2150 & l.~7 from foot & \textit{for} p.~2103 \textit{read} p.~2130 (Definition~9.1) \\
\dag\,2151 & proof of Lemma~9.9, l.~1 & \textit{for} $\neg\nec p \leftrightarrow \nec\neg p$ \textit{read} $\neg\pos p \leftrightarrow \nec\neg p$ \\
2153 & l.~2 & \textit{for} $\nec\pos p \to \pos p$ \textit{read} $\pos\pos p \to \pos p$ \\
2153 & l.~26 & \textit{delete} the second occurrence of \emph{to} \\
2153 & l.~6 from foot & \textit{for} $\neg\nec q \leftrightarrow \nec\neg q$ \textit{read} $\neg\pos q \leftrightarrow \nec\neg q$ \\
2154 & l.~9 & \textit{for} Lemma~2777 of Volume~X \textit{read} Lemma~2774 of Volume~X \\
2154 & l.~29 & \textit{for} Lemma~2734 of Volume~X \textit{read} Lemma~2735 of Volume~X; the result cited is on p.~3842 \\
2156 & l.~11 & \textit{for} anitcipate \textit{read} anticipate \\
2156 & l.~16 & \textit{for} $\pos\pos r \to r$ \textit{read} $\pos\nec r \to r$ \\
2157 & l.~12 & \textit{for} idetnical \textit{read} identical \\
2157 & display (9.16) & \textit{for} $v'Rw_2 \leftrightarrow w_2Rv'$ \textit{read} $v'Rw_2 \to w_2Rv'$; correctly printed in (9.20), p.~2161 \\
2158 & l.~1 & \textit{for} disagreemment \textit{read} disagreement \\
2160 & l.~7 & \textit{for} p.~1971 \textit{read} p.~1970 (Definition~6.1) \\
2161 & l.~7 & \textit{for} Proposition~2.5 of Volume~IV \textit{read} Proposition~2.6 of Volume~IV; the result cited is on p.~1425 \\
\dag\,2162 & proof of Corollary~9.12, l.~2 & \textit{for} Definition~9.4 \textit{read} Definition~9.1 \\
2162 & l.~7 & \textit{for} whreas \textit{read} whereas \\
2162 & l.~16 & \textit{for} necesity \textit{read} necessity \\
\dag\,2163 & proof of Corollary~9.12, (9.22) & \textit{for} $(v'Rw \vee v'Rw') \to wRw'$ \textit{read} $(v'Rw \wedge v'Rw') \to wRw'$ \\
2163 & l.~29 & \textit{for} accuont \textit{read} account \\
2163 & l.~5 from foot & \textit{for} eqiuvalent \textit{read} equivalent \\
2164 & l.~10 & \textit{for} Blackburn et~al. (2000) \textit{read} Blackburn et~al. (2001) \\
2164 & l.~12 & \textit{for} Lemma~2864 of Volume~X \textit{read} Lemma~2863 of Volume~X \\
2164 & l.~16 & \textit{for} Lemma~2741 of Volume~X \textit{read} Lemma~2746 of Volume~X \\
2164 & display (9.23) & \textit{for} $w_1Rv \leftrightarrow vRw_1$ \textit{read} $w_1Rv \to vRw_1$ \\
2165 & l.~14 & \textit{for} $w_1Rw_2 \wedge w_2Rw' \to w_2Rw'$ \textit{read} $w_1Rw_2 \wedge w_2Rw' \to w_1Rw'$ \\
2165 & l.~28 & \textit{for} impliction \textit{read} implication \\
2167 & l.~5 & \textit{for} underminne \textit{read} undermine \\
2167 & l.~14 & \textit{for} constriant \textit{read} constraint \\
2167 & l.~15 & \textit{for} $\pos q \to \nec\nec q$ \textit{read} $\nec q \to \nec\nec q$ \\
2167 & l.~18 & \textit{for} Lemma~2822 of Volume~X \textit{read} Lemma~2817 of Volume~X; the result cited is on p.~3850 \\
2167 & n.~3 & \textit{for} Chells (1980) \textit{read} Chellas (1980) \\
2168 & l.~35 & \textit{for} dennial \textit{read} denial \\
2170 & l.~8 & \textit{for} emprical \textit{read} empirical \\
2172 & display (9.26) & \textit{for} $\nec r \wedge \pos p \leftrightarrow \pos(r \wedge p)$ \textit{read} $\nec r \wedge \pos p \to \pos(r \wedge p)$; the formula is used in its correct form on p.~2176 \\
2172 & l.~38 & \textit{for} discusison \textit{read} discussion \\
2173 & l.~7 & \textit{for} deifciency \textit{read} deficiency \\
2173 & display (9.27) & \textit{for} $w_1Rw' \vee w'Ru \to w_1Ru$ \textit{read} $w_1Rw' \wedge w'Ru \to w_1Ru$ \\
\dag\,2174 & proof of Lemma~9.15, l.~6 & \textit{for} Lemma~9.7 \textit{read} Lemma~9.9 \\
\dag\,2174 & proof of Lemma~9.15, l.~26 & \textit{for} Definition~9.11 \textit{read} Definition~9.6 \\
\dag\,2175 & proof of Lemma~9.16, l.~1 & \textit{for} Lemma~2851 of Volume~X \textit{read} Lemma~2855 of Volume~X \\
2175 & l.~3 & \textit{for} govermnent \textit{read} government \\
\dag\,2175 & proof of Lemma~9.15, l.~5 & \textit{for} Definition~9.10 \textit{read} Definition~9.14 \\
\dag\,2175 & proof of Lemma~9.16, l.~15 & \textit{for} Definition~9.5 \textit{read} Definition~9.6 \\
2175 & l.~24 & \textit{for} p.~2001 \textit{read} p.~2000 (Theorem~6.15) \\
2175 & l.~8 from foot & \textit{for} evidnet \textit{read} evident \\
2176 & l.~13 & \textit{for} ocasionally \textit{read} occasionally \\
2176 & display (9.29) & \textit{for} $\nec\nec q \to q$ \textit{read} $\pos\nec q \to q$; the formula is used in its correct form on p.~2180 \\
2176 & l.~4 from foot & \textit{for} Lemma~2896 of Volume~X \textit{read} Lemma~2902 of Volume~X \\
\dag\,2177 & proof of Lemma~9.17, l.~6 & \textit{for} Lemma~9.20 \textit{read} Lemma~9.16 \\
\dag\,2177 & proof of Lemma~9.17, (9.30) & \textit{for} $\pos(p \to q) \to (\nec p \to \nec q)$ \textit{read} $\nec(p \to q) \to (\nec p \to \nec q)$ \\
\dag\,2178 & proof of Theorem~9.18, l.~6 & \textit{for} Definition~9.4 \textit{read} Definition~9.6 \\
2178 & l.~31 & \textit{for} Volume~I, Chapter~3 \textit{read} Volume~I, Chapter~4 (what the formalization establishes) \\
2179 & l.~21 & \textit{for} supposse \textit{read} suppose \\
\dag\,2179 & proof of Theorem~9.18, (9.32) & \textit{for} $\nec p \to \nec\pos p$ \textit{read} $\pos p \to \nec\pos p$ \\
2180 & display (9.33) & \textit{for} $\neg\pos p \leftrightarrow \pos\neg p$ \textit{read} $\neg\pos p \leftrightarrow \nec\neg p$ \\
2181 & l.~12 & \textit{for} conseqeuntly \textit{read} consequently \\
2182 & l.~1 & \textit{for} porcess \textit{read} process \\
2182 & n.~15 & \textit{for} Chellas (1979) \textit{read} Chellas (1980) \\
2183 & l.~6 & \textit{for} Lemma~2746 of Volume~X \textit{read} Lemma~2747 of Volume~X \\
2183 & l.~16 & \textit{for} Definition~5.15 of Volume~III \textit{read} Definition~5.18 of Volume~III \\
2184 & l.~2 & \textit{for} paardigm \textit{read} paradigm \\
2186 & l.~4 & \textit{for} apparenlty \textit{read} apparently \\
\dag\,2186 & proof of Lemma~9.20, l.~18 & \textit{for} $\nec p \wedge \pos r \leftrightarrow \pos(p \wedge r)$ \textit{read} $\nec p \wedge \pos r \to \pos(p \wedge r)$ \\
2188 & l.~23 & \textit{for} Blackbrun et~al. (2001) \textit{read} Blackburn et~al. (2001) \\
2188 & l.~40 & \textit{for} $r \leftrightarrow \nec\pos r$ \textit{read} $r \to \nec\pos r$ \\
\end{errlist}

\subsection*{Chapter~10. Summary \hfill \normalfont\small pp.~2189--2204}
\begin{errlist}
2191 & l.~12 & \textit{for} p.~2105 \textit{read} p.~2150 (Lemma~9.9) \\
2191 & l.~28 & \textit{for} Lemma~2772 of Volume~X \textit{read} Lemma~2777 of Volume~X; the result cited is on p.~3846 \\
2191 & l.~39 & \textit{for} invarriant \textit{read} invariant \\
2193 & l.~10 & \textit{for} execrise \textit{read} exercise \\
2193 & l.~11 & \textit{for} mcehanism \textit{read} mechanism \\
2193 & l.~15 & \textit{for} assmption \textit{read} assumption \\
2193 & l.~27 & \textit{for} otehrwise \textit{read} otherwise \\
2193 & l.~31 & \textit{for} procedrue \textit{read} procedure \\
2194 & l.~12 & \textit{for} opponant \textit{read} opponent \\
2195 & l.~9 & \textit{for} initialy \textit{read} initially \\
2195 & l.~4 from foot & \textit{for} irrelevnat \textit{read} irrelevant \\
2195 & l.~39 & \textit{for} follownig \textit{read} following \\
2196 & l.~3 & \textit{for} furhter \textit{read} further \\
2197 & l.~9 & \textit{for} Volume~V, Chapter~4 \textit{read} Volume~V, Chapter~3 (the question-begging objection) \\
2197 & l.~6 from foot & \textit{for} lattr \textit{read} latter \\
2198 & l.~23 & \textit{for} manitain \textit{read} maintain \\
2198 & l.~24 & \textit{for} Volume~VI, Chapter~6 \textit{read} Volume~VI, Chapter~5 (the noseeum inference) \\
2199 & l.~32 & \textit{for} inconsistnet \textit{read} inconsistent \\
2199 & l.~7 from foot & \textit{for} varrious \textit{read} various \\
2200 & l.~18 & \textit{for} comparsion \textit{read} comparison \\
2200 & l.~37 & \textit{for} Lemma~2801 of Volume~X \textit{read} Lemma~2796 of Volume~X; the result cited is on p.~3848 \\
2201 & l.~11 & \textit{for} descirbe \textit{read} describe \\
2203 & l.~2 & \textit{for} acknowlegde \textit{read} acknowledge \\
2203 & l.~7 & \textit{for} elsewhre \textit{read} elsewhere \\
2203 & l.~8 & \textit{for} Lemma~2806 of Volume~X \textit{read} Lemma~2800 of Volume~X \\
2204 & l.~27 & \textit{for} systemtaic \textit{read} systematic \\
2204 & l.~28 & \textit{for} othrwise \textit{read} otherwise \\
2204 & l.~1 from foot & \textit{for} Lemma~2893 of Volume~X \textit{read} Lemma~2897 of Volume~X \\
\end{errlist}


\section{Volume VI: Objections from Evil and Hiddenness}\label{s:vol6}
% generated by generate.py -- do not edit
\subsection*{Front matter \hfill \normalfont\small pp.~2205--2207}
\begin{errlist}
2206 & Contents, Chapter~6 & \textit{for} 2404 \textit{read} 2414 \\
\end{errlist}

\subsection*{Chapter~1. Introduction \hfill \normalfont\small pp.~2208--2214}
\begin{errlist}
2209 & l.~30 & \textit{for} propponent \textit{read} proponent \\
2209 & l.~1 from foot & \textit{for} orgnization \textit{read} organization \\
2210 & l.~14 & \textit{for} estbalish \textit{read} establish \\
2210 & n.~10 & \textit{for} Mackie (1956) \textit{read} Mackie (1955) \\
2212 & l.~24 & \textit{for} cnoscious \textit{read} conscious \\
2212 & l.~25 & \textit{for} objectoin \textit{read} objection \\
2212 & l.~8 from foot & \textit{delete} the second occurrence of \emph{is} \\
2213 & l.~2 & \textit{for} intiially \textit{read} initially \\
2213 & l.~4 & \textit{for} assessmant \textit{read} assessment \\
2213 & l.~7 from foot & \textit{for} reagrdless \textit{read} regardless \\
2213 & l.~4 from foot & \textit{delete} the second occurrence of \emph{to} \\
2214 & l.~5 from foot & \textit{for} p.~2211 \textit{read} p.~2210 (Definition~1.1) \\
\end{errlist}

\subsection*{Chapter~2. The logical problem of evil \hfill \normalfont\small pp.~2215--2258}
\begin{errlist}
2216 & l.~5 & \textit{for} Mackie (1955) \textit{read} Pike (1969) (omnipotence and the ability to sin) \\
2216 & l.~11 & \textit{for} dfiference \textit{read} difference \\
2217 & l.~4 & \textit{for} propposal \textit{read} proposal \\
2217 & l.~31 & \textit{for} Lemma~3118 of Volume~X \textit{read} Lemma~3117 of Volume~X \\
2217 & l.~32 & \textit{for} preilminary \textit{read} preliminary \\
2217 & l.~36 & \textit{for} $v'Rw_2 \wedge uRw_2 \to v'Ru$ \textit{read} $v'Rw_2 \wedge w_2Ru \to v'Ru$ \\
2218 & l.~9 & \textit{delete} the second occurrence of \emph{the} \\
2218 & l.~12 & \textit{delete} the second occurrence of \emph{that} \\
2218 & l.~2 from foot & \textit{for} $\nec(q \to q) \to (\nec p \to \nec q)$ \textit{read} $\nec(p \to q) \to (\nec p \to \nec q)$ \\
2219 & l.~31 & \textit{for} $w_1Rw'$ by symmetry \textit{read} $w_1Rw'$ by transitivity \\
2220 & l.~5 from foot & \textit{for} Lemma~3200 of Volume~X \textit{read} Lemma~3196 of Volume~X \\
2221 & l.~31 & \textit{for} Lemma~3257 of Volume~X \textit{read} Lemma~3260 of Volume~X; the result cited is on p.~3906 \\
2223 & l.~4 from foot & \textit{for} contant \textit{read} content \\
2223 & l.~38 & \textit{for} vlaidity \textit{read} validity \\
2223 & l.~40 & \textit{after} true at every world \textit{insert} accessible from $w'$ \\
2224 & l.~24 & \textit{for} atlhough \textit{read} although \\
\dag\,2225 & proof of Proposition~2.3, l.~8 & \textit{for} Lemma~3246 of Volume~X \textit{read} Lemma~3248 of Volume~X \\
\dag\,2225 & proof of Proposition~2.3, l.~24 & \textit{for} Lemma~2.8 of Volume~IV \textit{read} Lemma~2.4 of Volume~IV; the result cited is on p.~1419 \\
2226 & l.~33 & \textit{delete} the second occurrence of \emph{is} \\
2227 & l.~25 & \textit{for} $\nec(q \wedge p) \leftrightarrow (\nec q \wedge \pos p)$ \textit{read} $\nec(q \wedge p) \leftrightarrow (\nec q \wedge \nec p)$ \\
2229 & l.~2 & \textit{for} docrtine \textit{read} doctrine \\
2229 & l.~22 & \textit{for} true at some world accessible from $w_2$ (the clause for $\nec$) \textit{read} true at every world accessible from $w_2$ (the clause for $\nec$) \\
2229 & l.~31 & \textit{delete} the second occurrence of \emph{that} \\
2229 & l.~4 from foot & \textit{for} Lemma~3249 of Volume~X \textit{read} Lemma~3246 of Volume~X \\
2229 & l.~39 & \textit{for} lanugage \textit{read} language \\
\dag\,2230 & proof of Proposition~2.4, (2.8) & \textit{for} $\nec p \wedge \pos r \to \pos(r \wedge r)$ \textit{read} $\nec p \wedge \pos r \to \pos(p \wedge r)$ \\
\dag\,2230 & proof of Proposition~2.4, l.~19 & \textit{for} Definition~2.5 \textit{read} Definition~2.1; the step uses Definition~2.1 (p.~2223) \\
2230 & l.~22 & \textit{for} startegy \textit{read} strategy \\
2231 & l.~7 & \textit{delete} the second occurrence of \emph{the} \\
2232 & l.~25 & \textit{for} the noseeum inference is warranted \textit{read} the noseeum inference is unwarranted \\
2232 & display (2.10) & \textit{for} $wRu \vee uRv' \to wRv'$ \textit{read} $wRu \wedge uRv' \to wRv'$ \\
2232 & n.~18 & \textit{for} Mcakie (1955) \textit{read} Mackie (1955) \\
2233 & l.~23 & \textit{for} copmarison \textit{read} comparison \\
2233 & l.~29 & \textit{for} $\nec(T \wedge E)$ \textit{read} $\pos(T \wedge E)$ \\
\dag\,2234 & proof of Lemma~2.7, l.~5 & \textit{for} Definition~2.5 \textit{read} Definition~2.6; the step uses Definition~2.6 (p.~2230) \\
2234 & l.~9 & \textit{for} capcity \textit{read} capacity \\
2235 & l.~18 & \textit{for} Hume, \emph{Dialogues}, Part~XI (the a priori argument) \textit{read} Hume, \emph{Dialogues}, Part~IX (the a priori argument) \\
2236 & l.~5 from foot & \textit{for} Hume, \emph{Dialogues}, Part~IX (the argument from evil) \textit{read} Hume, \emph{Dialogues}, Parts~X--XI (the argument from evil) \\
2237 & l.~25 & \textit{for} $\nec(p \to r) \to (\nec p \leftrightarrow \nec r)$ \textit{read} $\nec(p \to r) \to (\nec p \to \nec r)$ \\
2238 & l.~18 & \textit{for} p.~2226 \textit{read} p.~2225 (Proposition~2.3) \\
2238 & l.~32 & \textit{for} Lemma~3076 of Volume~X \textit{read} Lemma~3070 of Volume~X \\
2239 & l.~1 & \textit{for} Lemma~3182 of Volume~X \textit{read} Lemma~3184 of Volume~X \\
2239 & l.~12 & \textit{for} experssion \textit{read} expression \\
2240 & l.~19 & \textit{for} the logical problem of evil succeeds \textit{read} the logical problem of evil fails \\
2241 & l.~12 & \textit{for} postiion \textit{read} position \\
\dag\,2241 & proof of Proposition~2.8, (2.16) & \textit{for} $(w_0Rv \wedge w_0Rw) \leftrightarrow vRw$ \textit{read} $(w_0Rv \wedge w_0Rw) \to vRw$; the formula is used in its correct form on p.~2245 \\
2242 & l.~15 & \textit{for} $\pos\nec r \leftrightarrow r$ \textit{read} $\pos\nec r \to r$ \\
2242 & l.~20 & \textit{for} suffreing \textit{read} suffering \\
\dag\,2244 & proof of Proposition~2.9, l.~8 & \textit{for} Proposition~2.4 \textit{read} Proposition~2.3 \\
2244 & l.~25 & \textit{delete} the second occurrence of \emph{to} \\
\dag\,2244 & proof of Proposition~2.9, l.~26 & \textit{for} Definition~2.6 \textit{read} Definition~2.1 \\
\dag\,2245 & proof of Proposition~2.9, (2.20) & \textit{for} $\pos\nec p \leftrightarrow \nec p$ \textit{read} $\pos\nec p \to \nec p$ \\
\dag\,2247 & proof of Lemma~2.10, l.~9 & \textit{for} $\pos r \leftrightarrow \nec\neg r$ \textit{read} $\neg\pos r \leftrightarrow \nec\neg r$ \\
2247 & l.~22 & \textit{for} meerly \textit{read} merely \\
\dag\,2248 & proof of Lemma~2.11, l.~24 & \textit{for} Lemma~2.10 \textit{read} Lemma~2.7; the step uses Lemma~2.7 (p.~2234) \\
\dag\,2248 & proof of Lemma~2.11, l.~25 & \textit{for} Proposition~2.8 \textit{read} Proposition~2.3 \\
\dag\,2248 & proof of Lemma~2.11, (2.22) & \textit{for} $u_C(\mathrm{free}) < u_C(\mathrm{assured})$ \textit{read} $u_C(\mathrm{free}) > u_C(\mathrm{assured})$ \\
2249 & l.~39 & \textit{for} p.~2258 \textit{read} p.~2248 (Lemma~2.11) \\
2251 & l.~25 & \textit{for} $\pos q \leftrightarrow \nec\pos q$ \textit{read} $\pos q \to \nec\pos q$ \\
2251 & l.~34 & \textit{for} Plantinga (1974a) \textit{read} Mackie (1955) (the logical problem of evil) \\
2253 & l.~14 & \textit{for} evidence of a gratuitous evil \textit{read} evidence of the absence of a justifying good \\
2253 & l.~29 & \textit{for} p.~2250 \textit{read} p.~2240 (Proposition~2.8) \\
2254 & l.~2 & \textit{for} theodciy \textit{read} theodicy \\
2255 & l.~2 & \textit{for} unambiguous revelation would enlarge the freedom of response \textit{read} unambiguous revelation would constrain the freedom of response \\
\dag\,2256 & proof of Proposition~2.15, l.~14 & \textit{for} Proposition~2.8 \textit{read} Proposition~2.9 \\
2257 & l.~10 & \textit{for} p.~2266 \textit{read} p.~2256 (Proposition~2.15) \\
2257 & l.~19 & \textit{for} p.~2252 \textit{read} p.~2253 (Definition~2.13) \\
2258 & l.~3 & \textit{after} accessible \textit{insert} from $w$ \\
2258 & l.~20 & \textit{for} intedned \textit{read} intended \\
2258 & l.~25 & \textit{for} ambigiuty \textit{read} ambiguity \\
2258 & l.~26 & \textit{for} sufferng \textit{read} suffering \\
2258 & l.~34 & \textit{after} intense suffering \textit{insert} that it could prevent \\
2258 & l.~5 from foot & \textit{for} frmae \textit{read} frame \\
2258 & l.~3 from foot & \textit{for} $uRw_2$ by reflexivity \textit{read} $uRu$ by reflexivity \\
\end{errlist}

\subsection*{Chapter~3. The structure of logical space and the free-will defence \hfill \normalfont\small pp.~2259--2311}
\begin{errlist}
2260 & running heads, pp.~2260--2310 (even) & \textit{for} The logical problem of evil \textit{read} Logical space and the free-will defence \\
2260 & l.~7 & \textit{for} irreelvant \textit{read} irrelevant \\
2262 & l.~29 & \textit{for} Lemma~3219 of Volume~X \textit{read} Lemma~3214 of Volume~X; the result cited is on p.~3900 \\
2264 & l.~9 & \textit{for} convicsion \textit{read} conviction \\
2265 & l.~9 & \textit{for} requireemnt \textit{read} requirement \\
2265 & l.~18 & \textit{delete} the second occurrence of \emph{that} \\
2265 & l.~25 & \textit{delete} the second occurrence of \emph{is} \\
2265 & display (3.5) & \textit{for} $\pos p \leftrightarrow \nec\pos p$ \textit{read} $\pos p \to \nec\pos p$ \\
2267 & l.~25 & \textit{for} p.~2213 \textit{read} p.~2212 (Theorem~1.2) \\
2267 & l.~30 & \textit{for} etsablish \textit{read} establish \\
2267 & l.~1 from foot & \textit{for} familair \textit{read} familiar \\
2268 & l.~15 & \textit{for} Lemma~3190 of Volume~X \textit{read} Lemma~3184 of Volume~X \\
2268 & l.~2 from foot & \textit{for} $T \leftrightarrow \neg\exists e\,\mathrm{Grat}(e)$ \textit{read} $T \to \neg\exists e\,\mathrm{Grat}(e)$ \\
2269 & l.~24 & \textit{for} p.~2320 \textit{read} p.~2230 (Definition~2.5) \\
\dag\,2270 & proof of Lemma~3.4, l.~10 & \textit{for} Definition~3.1 \textit{read} Definition~3.3 \\
\dag\,2271 & proof of Corollary~3.5, l.~3 & \textit{for} Lemma~3.6 \textit{read} Lemma~3.4; the step uses Lemma~3.4 (p.~2270) \\
2271 & l.~22 & \textit{for} Hick (1965) \textit{read} Hick (1966) \\
2271 & l.~23 & \textit{for} defedner \textit{read} defender \\
\dag\,2271 & proof of Lemma~3.6, l.~27 & \textit{for} Definition~3.3 \textit{read} Definition~3.2 \\
2271 & l.~36 & \textit{for} acocrdingly \textit{read} accordingly \\
\dag\,2272 & proof of Lemma~3.6, l.~1 & \textit{for} Definition~3.8 \textit{read} Definition~3.3 \\
\dag\,2272 & proof of Corollary~3.5, l.~11 & \textit{for} Definition~3.1 \textit{read} Definition~3.2 \\
2272 & l.~19 & \textit{for} cnoception \textit{read} conception \\
\dag\,2272 & proof of Lemma~3.6, l.~27 & \textit{for} Lemma~3.7 \textit{read} Lemma~3.4; the step uses Lemma~3.4 (p.~2270) \\
2273 & l.~19 & \textit{for} Lemma~3179 of Volume~X \textit{read} Lemma~3182 of Volume~X \\
2274 & l.~5 & \textit{for} regardd \textit{read} regarded \\
2275 & l.~1 & \textit{for} theolgoical \textit{read} theological \\
2275 & l.~32 & \textit{for} defenedd \textit{read} defended \\
2275 & l.~35 & \textit{for} proeced \textit{read} proceed \\
2276 & l.~5 & \textit{delete} the second occurrence of \emph{the} \\
2276 & l.~10 & \textit{for} comparbale \textit{read} comparable \\
2276 & display (3.10) & \textit{for} $T \to \exists e\,\mathrm{Grat}(e)$ \textit{read} $T \to \neg\exists e\,\mathrm{Grat}(e)$; the formula is used in its correct form on p.~2279 \\
2276 & l.~18 & \textit{for} p.~2253 \textit{read} p.~2254 (Proposition~2.14) \\
2277 & l.~21 & \textit{for} requirrement \textit{read} requirement \\
2277 & l.~39 & \textit{for} obviuos \textit{read} obvious \\
2279 & l.~29 & \textit{for} Lemma~3053 of Volume~X \textit{read} Lemma~3052 of Volume~X \\
2280 & l.~4 & \textit{for} equivallence \textit{read} equivalence \\
2280 & display (3.11) & \textit{for} $uRw' \to uRw'$ \textit{read} $uRw' \to w'Ru$ \\
2280 & l.~21 & \textit{for} sucecssful \textit{read} successful \\
2280 & l.~26 & \textit{for} p.~2213 \textit{read} p.~2223 (Definition~2.1) \\
2280 & l.~4 from foot & \textit{for} the axiom 4, $\pos p \to \nec\pos p$ \textit{read} the axiom 5, $\pos p \to \nec\pos p$ \\
2282 & l.~1 & \textit{for} exericse \textit{read} exercise \\
\dag\,2282 & proof of Lemma~3.7, l.~5 & \textit{for} $q \to \nec\nec q$ \textit{read} $q \to \nec\pos q$ \\
2282 & l.~8 & \textit{for} convetnion \textit{read} convention \\
\dag\,2282 & proof of Lemma~3.7, (3.14) & \textit{for} $u_G(h, r) \leq u_G(v, r)$ \textit{read} $u_G(h, r) \geq u_G(v, r)$ \\
2282 & l.~27 & \textit{for} contradcition \textit{read} contradiction \\
2283 & l.~11 & \textit{for} Lemma~3090 of Volume~X \textit{read} Lemma~3089 of Volume~X; the result cited is on p.~3882 \\
2284 & l.~6 from foot & \textit{delete} the second occurrence of \emph{the} \\
2285 & l.~4 & \textit{for} Lemma~3186 of Volume~X \textit{read} Lemma~3189 of Volume~X \\
2285 & l.~9 & \textit{after} S5 conditions \textit{insert} on $R$ \\
2286 & l.~2 & \textit{for} Volume~III, Chapter~4 \textit{read} Volume~III, Chapter~2 (conceivability and possibility) \\
2286 & l.~25 & \textit{for} Lemma~3124 of Volume~X \textit{read} Lemma~3127 of Volume~X \\
2287 & l.~12 & \textit{delete} the second occurrence of \emph{of} \\
\dag\,2287 & proof of Theorem~3.9, (3.18) & \textit{for} $(wRv \vee wRw_1) \to vRw_1$ \textit{read} $(wRv \wedge wRw_1) \to vRw_1$; correctly printed in (3.20), p.~2292 \\
2287 & l.~22 & \textit{for} significnat \textit{read} significant \\
\dag\,2287 & proof of Theorem~3.9, l.~33 & \textit{for} $\exists w\,(T(w) \vee E(w))$ \textit{read} $\exists w\,(T(w) \wedge E(w))$ \\
2288 & l.~19 & \textit{for} $(w_2Rw' \wedge w'Rw) \to w_2Rw$ \textit{read} $(w'Rw_2 \wedge w'Rw) \to w_2Rw$ \\
2289 & l.~6 & \textit{for} Definition~6.1 of Volume~IV \textit{read} Definition~6.2 of Volume~IV \\
\dag\,2291 & proof of Proposition~3.11, l.~11 & \textit{for} Lemma~3160 of Volume~X \textit{read} Lemma~3161 of Volume~X; the result cited is on p.~3892 \\
2291 & l.~32 & \textit{for} particulrly \textit{read} particularly \\
2293 & display (3.21) & \textit{for} $P(N \mid  S) \approx P(N)$ \textit{read} $P(N \mid \neg S) \approx P(N)$ \\
2293 & l.~16 & \textit{for} conrtadiction \textit{read} contradiction \\
2293 & l.~8 from foot & \textit{for} p.~2276 \textit{read} p.~2286 (Theorem~3.9) \\
2295 & l.~14 & \textit{for} posiiton \textit{read} position \\
2296 & l.~29 & \textit{for} absoluetly \textit{read} absolutely \\
2296 & l.~5 from foot & \textit{for} dveelop \textit{read} develop \\
2297 & l.~9 & \textit{for} releavnce \textit{read} relevance \\
2297 & n.~1 & \textit{for} van Inwagen (2005) \textit{read} van Inwagen (2006) \\
2298 & display (3.25) & \textit{for} $\nec q \wedge \pos p \to \pos(q \vee p)$ \textit{read} $\nec q \wedge \pos p \to \pos(q \wedge p)$; correctly printed in (3.27), p.~2304 \\
2300 & l.~5 & \textit{for} commtment \textit{read} commitment \\
2302 & display (3.26) & \textit{for} $\pos\nec r \leftrightarrow \nec r$ \textit{read} $\pos\nec r \to \nec r$ \\
2302 & l.~31 & \textit{for} trnsferred \textit{read} transferred \\
\dag\,2303 & proof of Proposition~3.13, l.~7 & \textit{for} Proposition~3.14 \textit{read} Proposition~3.11 \\
\dag\,2304 & proof of Proposition~3.15, l.~6 & \textit{for} Proposition~3.14 \textit{read} Proposition~3.13; the step uses Proposition~3.13 (p.~2303) \\
\dag\,2304 & proof of Proposition~3.14, l.~12 & \textit{for} Lemma~3.7 \textit{read} Lemma~3.6; the step uses Lemma~3.6 (p.~2271) \\
2304 & l.~22 & \textit{for} Hick (1967) \textit{read} Hick (1966) \\
\dag\,2304 & proof of Proposition~3.15, l.~24 & \textit{for} Definition~3.10 \textit{read} Definition~3.8; the step uses Definition~3.8 (p.~2285) \\
\dag\,2304 & proof of Proposition~3.13, l.~30 & \textit{for} Lemma~3.7 \textit{read} Lemma~3.12 \\
2305 & l.~9 & \textit{for} soundnes \textit{read} soundness \\
\dag\,2305 & proof of Proposition~3.15, l.~15 & \textit{for} Lemma~3260 of Volume~X \textit{read} Lemma~3254 of Volume~X \\
2306 & l.~7 & \textit{for} eaxctly \textit{read} exactly \\
2306 & l.~24 & \textit{for} Lemma~3202 of Volume~X \textit{read} Lemma~3203 of Volume~X; the result cited is on p.~3898 \\
2307 & display (3.29) & \textit{for} $\pos(T \vee E)$ \textit{read} $\pos(T \wedge E)$ \\
2308 & l.~29 & \textit{for} $P(E \mid T) > P(E \mid \neg T)$ \textit{read} $P(E \mid T) < P(E \mid \neg T)$ \\
2308 & display (3.30) & \textit{for} $\nec(p \to q) \to (\nec p \leftrightarrow \nec q)$ \textit{read} $\nec(p \to q) \to (\nec p \to \nec q)$; correctly printed in (3.32), p.~2310 \\
2308 & l.~39 & \textit{for} the axiom 4, $\nec p \to p$ \textit{read} the axiom T, $\nec p \to p$ \\
2309 & l.~21 & \textit{for} evetnually \textit{read} eventually \\
2311 & l.~18 & \textit{for} somethnig \textit{read} something \\
2311 & l.~27 & \textit{for} $\nec(p \wedge r) \leftrightarrow (\nec r \wedge \nec r)$ \textit{read} $\nec(p \wedge r) \leftrightarrow (\nec p \wedge \nec r)$ \\
2311 & l.~5 from foot & \textit{for} sometihng \textit{read} something \\
\end{errlist}

\subsection*{Chapter~4. The evidential problem of evil \hfill \normalfont\small pp.~2312--2360}
\begin{errlist}
2313 & display (4.1) & \textit{for} $\forall w_0\,(T(w_0) \wedge E(w_0))$ \textit{read} $\exists w_0\,(T(w_0) \wedge E(w_0))$ \\
2313 & l.~8 & \textit{for} nsoeeum \textit{read} noseeum \\
2313 & l.~9 & \textit{for} explcit \textit{read} explicit \\
2315 & l.~17 & \textit{for} p.~2217 \textit{read} p.~2271 (Lemma~3.6) \\
2316 & l.~6 & \textit{for} Rowe (1978) \textit{read} Rowe (1979) \\
2316 & l.~11 & \textit{for} deficiecny \textit{read} deficiency \\
2316 & l.~30 & \textit{for} ultimmately \textit{read} ultimately \\
2318 & display (4.2) & \textit{see} Corrigendum~S11 \\
2319 & l.~4 & \textit{for} qusetion \textit{read} question \\
2319 & l.~5 & \textit{for} Lemma~3057 of Volume~X \textit{read} Lemma~3063 of Volume~X \\
2320 & l.~22 & \textit{for} Lemma~2.8 of Volume~IV \textit{read} Lemma~2.12 of Volume~IV \\
2320 & l.~27 & \textit{after} equilibrium \textit{insert} of the signalling game \\
2322 & l.~13 & \textit{for} tehodicy \textit{read} theodicy \\
2322 & l.~26 & \textit{for} Volume~V, Chapter~3 \textit{read} Volume~V, Chapter~2 (existence as a predicate) \\
2323 & l.~28 & \textit{for} p.~2300 \textit{read} p.~2310 (Definition~3.17) \\
2323 & l.~38 & \textit{for} probalbe \textit{read} probable \\
2323 & l.~1 from foot & \textit{for} referrd \textit{read} referred \\
2324 & l.~18 & \textit{for} convction \textit{read} conviction \\
2324 & l.~19 & \textit{for} Lemma~3198 of Volume~X \textit{read} Lemma~3194 of Volume~X \\
2324 & display (4.4) & \textit{for} $P(E \mid T) < P(E \mid  T)$ \textit{read} $P(E \mid T) < P(E \mid \neg T)$ \\
2325 & l.~2 & \textit{for} Lemma~3196 of Volume~X \textit{read} Lemma~3195 of Volume~X \\
2325 & l.~15 & \textit{for} Lemma~3168 of Volume~X \textit{read} Lemma~3172 of Volume~X; the result cited is on p.~3894 \\
2326 & l.~16 & \textit{for} ecnounter \textit{read} encounter \\
2326 & l.~37 & \textit{delete} the second occurrence of \emph{to} \\
2327 & l.~3 & \textit{for} indcate \textit{read} indicate \\
2327 & l.~19 & \textit{for} Lemma~3101 of Volume~X \textit{read} Lemma~3106 of Volume~X \\
2327 & l.~20 & \textit{for} $P(\neg S \mid  N) \approx 1$ \textit{read} $P(\neg S \mid \neg N) \approx 1$ \\
2327 & l.~5 from foot & \textit{delete} the second occurrence of \emph{a} \\
2328 & l.~7 & \textit{for} explict \textit{read} explicit \\
2328 & l.~13 & \textit{for} prvided \textit{read} provided \\
2328 & l.~26 & \textit{for} scetpical \textit{read} sceptical \\
2328 & l.~1 from foot & \textit{for} calimed \textit{read} claimed \\
2330 & l.~15 & \textit{for} frmework \textit{read} framework \\
\dag\,2330 & proof of Corollary~4.7, l.~21 & \textit{for} Proposition~4.8 \textit{read} Proposition~4.4 \\
2333 & l.~15 & \textit{delete} the second occurrence of \emph{of} \\
2333 & l.~18 & \textit{for} notasion \textit{read} notation \\
2333 & l.~1 from foot & \textit{for} straighftorward \textit{read} straightforward \\
2334 & l.~2 & \textit{for} infernece \textit{read} inference \\
2334 & l.~15 & \textit{for} genreally \textit{read} generally \\
2334 & l.~24 & \textit{delete} the second occurrence of \emph{of} \\
2335 & l.~11 & \textit{for} neecssary \textit{read} necessary \\
2335 & l.~23 & \textit{for} therbey \textit{read} thereby \\
2335 & l.~39 & \textit{for} tehological \textit{read} theological \\
2340 & l.~14 & \textit{for} p.~2403 \textit{read} p.~2303 (Proposition~3.13) \\
2340 & l.~23 & \textit{for} emploeyd \textit{read} employed \\
2340 & l.~27 & \textit{for} particualrly \textit{read} particularly \\
2340 & l.~8 from foot & \textit{for} positoin \textit{read} position \\
2341 & l.~7 from foot & \textit{for} dceisive \textit{read} decisive \\
2342 & l.~24 & \textit{for} p.~2263 \textit{read} p.~2264 (Definition~3.1) \\
\dag\,2343 & proof of Proposition~4.10, l.~3 & \textit{for} Corollary~4.7 \textit{read} Corollary~4.9; the step uses Corollary~4.9 (p.~2341) \\
\dag\,2343 & proof of Corollary~4.11, l.~8 & \textit{for} Proposition~4.8 \textit{read} Proposition~4.10 \\
2343 & l.~27 & \textit{for} p.~2268 \textit{read} p.~2286 (Definition~3.10) \\
2343 & l.~3 from foot & \textit{for} van Inwagen (2006) \textit{read} Draper (1989) (the hypothesis of indifference) \\
2344 & l.~8 & \textit{for} Lemma~3168 of Volume~X \textit{read} Lemma~3173 of Volume~X \\
2345 & n.~19 & \textit{for} Draper (1990) \textit{read} Draper (1989) \\
2346 & l.~14 & \textit{for} defneder \textit{read} defender \\
\dag\,2346 & proof of Lemma~4.12, l.~15 & \textit{for} $\exists w'\,(T(w') \vee E(w'))$ \textit{read} $\exists w'\,(T(w') \wedge E(w'))$ \\
\dag\,2346 & proof of Lemma~4.12, l.~20 & \textit{for} Corollary~4.11 \textit{read} Corollary~4.9; the step uses Corollary~4.9 (p.~2341) \\
2347 & l.~8 & \textit{for} quession \textit{read} question \\
2347 & l.~12 & \textit{for} p.~2313 \textit{read} p.~2213 (Proposition~1.3) \\
2347 & l.~22 & \textit{for} p.~2381 \textit{read} p.~2281 (Lemma~3.7) \\
2347 & l.~27 & \textit{for} evidance \textit{read} evidence \\
\dag\,2347 & proof of Lemma~4.12, l.~35 & \textit{for} Lemma~3136 of Volume~X \textit{read} Lemma~3142 of Volume~X; the result cited is on p.~3889 \\
2348 & l.~14 & \textit{for} contemporray \textit{read} contemporary \\
2348 & l.~16 & \textit{for} p.~2334 \textit{read} p.~2343 (Proposition~4.10) \\
2348 & l.~19 & \textit{for} paritcularly \textit{read} particularly \\
2348 & l.~25 & \textit{for} Lemma~3201 of Volume~X \textit{read} Lemma~3203 of Volume~X; the result cited is on p.~3898 \\
2349 & l.~16 & \textit{for} Rowe (1980) \textit{read} Rowe (1979) \\
2349 & l.~25 & \textit{for} Rowe (1979) \textit{read} Draper (1989) (the hypothesis of indifference) \\
2352 & l.~6 & \textit{for} Lemma~5.19 of Volume~III \textit{read} Lemma~5.21 of Volume~III \\
\dag\,2354 & proof of Corollary~4.16, l.~18 & \textit{for} Definition~4.15 \textit{read} Definition~4.13 \\
2354 & l.~21 & \textit{for} Proposition~6.13 of Volume~IV \textit{read} Proposition~6.10 of Volume~IV; the result cited is on p.~1617 \\
2355 & l.~29 & \textit{for} p.~2440 \textit{read} p.~2340 (Proposition~4.8) \\
2356 & l.~8 & \textit{for} p.~2249 \textit{read} p.~2349 (Lemma~4.14) \\
2356 & l.~32 & \textit{for} Lemma~3245 of Volume~X \textit{read} Lemma~3246 of Volume~X \\
2357 & l.~8 & \textit{for} p.~2353 \textit{read} p.~2354 (Corollary~4.16) \\
2357 & l.~27 & \textit{for} hwoever \textit{read} however \\
2357 & l.~36 & \textit{for} Draepr (1989) \textit{read} Draper (1989) \\
2358 & l.~9 & \textit{for} a separating equilibrium \textit{read} a pooling equilibrium \\
2358 & l.~14 & \textit{for} p.~2314 \textit{read} p.~2341 (Corollary~4.9) \\
2358 & l.~20 & \textit{for} ackonwledge \textit{read} acknowledge \\
2359 & l.~26 & \textit{for} Draper (1988) \textit{read} Draper (1989) \\
2360 & l.~3 from foot & \textit{for} Lemma~3234 of Volume~X \textit{read} Lemma~3232 of Volume~X \\
\end{errlist}

\subsection*{Chapter~5. The noseeum inference \hfill \normalfont\small pp.~2361--2413}
\begin{errlist}
2362 & running heads, pp.~2362--2412 (even) & \textit{for} Divine hiddenness \textit{read} The noseeum inference \\
2362 & l.~19 & \textit{for} Lemma~3204 of Volume~X \textit{read} Lemma~3205 of Volume~X; the result cited is on p.~3898 \\
2362 & l.~23 & \textit{for} fmailiar \textit{read} familiar \\
2362 & l.~39 & \textit{for} Lemma~3150 of Volume~X \textit{read} Lemma~3146 of Volume~X \\
2363 & l.~3 & \textit{for} tarnsferred \textit{read} transferred \\
2363 & l.~9 & \textit{for} Lemma~3228 of Volume~X \textit{read} Lemma~3224 of Volume~X; the result cited is on p.~3901 \\
2363 & l.~15 & \textit{for} Begrmann (2001) \textit{read} Bergmann (2001) \\
2364 & l.~25 & \textit{delete} the second occurrence of \emph{is} \\
2364 & l.~26 & \textit{for} Lemma~3128 of Volume~X \textit{read} Lemma~3129 of Volume~X \\
2364 & l.~30 & \textit{for} Lemma~3226 of Volume~X \textit{read} Lemma~3227 of Volume~X; the result cited is on p.~3901 \\
\dag\,2365 & proof of Proposition~5.2, l.~22 & \textit{for} Lemma~3236 of Volume~X \textit{read} Lemma~3233 of Volume~X; the result cited is on p.~3902 \\
2366 & l.~29 & \textit{for} Bergmann (2002) \textit{read} Bergmann (2001) \\
2367 & l.~7 & \textit{for} ambgiuity \textit{read} ambiguity \\
2367 & l.~22 & \textit{for} Definition~2.3 of Volume~IV \textit{read} Definition~2.1 of Volume~IV; the result cited is on p.~1417 \\
2367 & l.~28 & \textit{for} p.~2276 \textit{read} p.~2267 (Definition~3.2) \\
2367 & l.~32 & \textit{delete} the second occurrence of \emph{to} \\
2368 & l.~24 & \textit{for} Wykstra (1984) \textit{read} Bergmann (2001) (sceptical theism) \\
2369 & l.~23 & \textit{for} Lemma~3282 of Volume~X \textit{read} Lemma~3286 of Volume~X; the result cited is on p.~3910 \\
2370 & l.~5 from foot & \textit{for} persuasvie \textit{read} persuasive \\
\dag\,2371 & proof of Corollary~5.4, l.~9 & \textit{for} Definition~5.3 \textit{read} Definition~5.1 \\
2372 & l.~38 & \textit{delete} the second occurrence of \emph{a} \\
2373 & l.~16 & \textit{for} evdient \textit{read} evident \\
\dag\,2373 & proof of Lemma~5.6, l.~24 & \textit{for} Lemma~3169 of Volume~X \textit{read} Lemma~3165 of Volume~X; the result cited is on p.~3893 \\
2373 & l.~30 & \textit{for} p.~2295 \textit{read} p.~2296 (Lemma~3.12) \\
2374 & l.~20 & \textit{for} Lemma~3.3 of Volume~I \textit{read} Lemma~3.2 of Volume~I \\
2376 & l.~5 from foot & \textit{for} evidetnial \textit{read} evidential \\
2377 & l.~15 & \textit{for} abmiguity \textit{read} ambiguity \\
2379 & l.~16 & \textit{for} succssful \textit{read} successful \\
2379 & l.~18 & \textit{for} Lemma~3045 of Volume~X \textit{read} Lemma~3051 of Volume~X \\
2379 & l.~32 & \textit{for} Wykstra (1983) \textit{read} Wykstra (1984) \\
2380 & l.~15 & \textit{for} Rowe (1979) \textit{read} Wykstra (1984) (the condition of reasonable epistemic access) \\
2380 & l.~6 from foot & \textit{for} grtuitous \textit{read} gratuitous \\
2381 & l.~12 & \textit{for} Lemma~3110 of Volume~X \textit{read} Lemma~3113 of Volume~X \\
2381 & l.~3 from foot & \textit{for} neitehr \textit{read} neither \\
2382 & l.~19 & \textit{for} resebmlance \textit{read} resemblance \\
2383 & l.~18 & \textit{for} detterminate \textit{read} determinate \\
2383 & l.~25 & \textit{for} impossable \textit{read} impossible \\
\dag\,2383 & proof of Theorem~5.8, (5.7) & \textit{for} $\forall w'\,(T(w') \wedge E(w'))$ \textit{read} $\exists w'\,(T(w') \wedge E(w'))$ \\
2384 & l.~9 & \textit{for} Lemma~3187 of Volume~X \textit{read} Lemma~3184 of Volume~X \\
2385 & l.~5 & \textit{for} solutoin \textit{read} solution \\
2385 & l.~15 & \textit{for} recomemnd \textit{read} recommend \\
2385 & l.~6 from foot & \textit{for} intially \textit{read} initially \\
2385 & l.~2 from foot & \textit{for} employd \textit{read} employed \\
2385 & l.~1 from foot & \textit{for} opponnent \textit{read} opponent \\
2386 & l.~33 & \textit{for} p.~2286 \textit{read} p.~2268 (Definition~3.3) \\
2387 & l.~13 & \textit{for} decisve \textit{read} decisive \\
2387 & l.~19 & \textit{for} requiremnet \textit{read} requirement \\
2387 & display (5.8) & \textit{for} $\exists w_0\,(T(w_0) \vee E(w_0))$ \textit{read} $\exists w_0\,(T(w_0) \wedge E(w_0))$ \\
2387 & l.~3 from foot & \textit{delete} the second occurrence of \emph{that} \\
2388 & l.~24 & \textit{for} despte \textit{read} despite \\
2389 & l.~21 & \textit{for} Lemma~3127 of Volume~X \textit{read} Lemma~3131 of Volume~X; the result cited is on p.~3888 \\
2390 & l.~6 & \textit{for} concsious \textit{read} conscious \\
2390 & l.~7 from foot & \textit{for} assesment \textit{read} assessment \\
2391 & l.~1 & \textit{for} distnguish \textit{read} distinguish \\
2392 & l.~7 & \textit{for} Lemma~3274 of Volume~X \textit{read} Lemma~3271 of Volume~X; the result cited is on p.~3908 \\
2393 & l.~3 & \textit{for} execption \textit{read} exception \\
2393 & l.~21 & \textit{for} Rowe (1979) \textit{read} Bergmann (2001) (sceptical theism) \\
\dag\,2393 & proof of Proposition~5.10, l.~26 & \textit{for} Proposition~5.7 \textit{read} Proposition~5.2 \\
2394 & l.~32 & \textit{for} cnovention \textit{read} convention \\
2396 & l.~17 & \textit{for} epxression \textit{read} expression \\
2398 & l.~22 & \textit{for} slgihtly \textit{read} slightly \\
2400 & l.~31 & \textit{for} Wyktsra (1984) \textit{read} Wykstra (1984) \\
2401 & l.~10 & \textit{for} epsecially \textit{read} especially \\
\dag\,2402 & proof of Proposition~5.13, l.~18 & \textit{for} Proposition~5.14 \textit{read} Proposition~5.10; the step uses Proposition~5.10 (p.~2393) \\
2402 & l.~21 & \textit{for} Wykstra (1985) \textit{read} Wykstra (1984) \\
2402 & l.~23 & \textit{for} p.~2243 \textit{read} p.~2234 (Lemma~2.7) \\
2402 & l.~3 from foot & \textit{for} initilaly \textit{read} initially \\
2403 & l.~23 & \textit{for} releavnt \textit{read} relevant \\
\dag\,2403 & proof of Proposition~5.13, l.~29 & \textit{for} Theorem~5.8 \textit{read} Theorem~5.11 \\
\dag\,2404 & proof of Proposition~5.14, l.~5 & \textit{for} Lemma~3201 of Volume~X \textit{read} Lemma~3206 of Volume~X; the result cited is on p.~3898 \\
2404 & l.~12 & \textit{for} sreious \textit{read} serious \\
2404 & l.~15 & \textit{for} copmarable \textit{read} comparable \\
\dag\,2404 & proof of Proposition~5.14, l.~17 & \textit{for} Lemma~3241 of Volume~X \textit{read} Lemma~3236 of Volume~X \\
2405 & l.~11 & \textit{for} p.~2275 \textit{read} p.~2285 (Definition~3.8) \\
2405 & l.~37 & \textit{for} p.~2413 \textit{read} p.~2403 (Proposition~5.14) \\
2406 & l.~22 & \textit{for} Lemma~3053 of Volume~X \textit{read} Lemma~3059 of Volume~X \\
2406 & l.~29 & \textit{for} Lemma~3149 of Volume~X \textit{read} Lemma~3145 of Volume~X; the result cited is on p.~3890 \\
2409 & display (5.13) & \textit{for} $P( S \mid \neg N) \approx 1$ \textit{read} $P(\neg S \mid \neg N) \approx 1$ \\
2409 & l.~19 & \textit{for} Lemma~3142 of Volume~X \textit{read} Lemma~3140 of Volume~X \\
2409 & l.~25 & \textit{for} Lemma~3082 of Volume~X \textit{read} Lemma~3080 of Volume~X \\
2410 & l.~5 & \textit{for} straightforawrd \textit{read} straightforward \\
2410 & l.~8 & \textit{for} p.~2473 \textit{read} p.~2373 (Lemma~5.6) \\
2411 & l.~30 & \textit{for} p.~2396 \textit{read} p.~2395 (Corollary~5.12) \\
2412 & l.~18 & \textit{for} p.~2381 \textit{read} p.~2382 (Theorem~5.8) \\
2412 & l.~21 & \textit{for} Lemma~3174 of Volume~X \textit{read} Lemma~3172 of Volume~X; the result cited is on p.~3894 \\
2412 & l.~23 & \textit{for} Bergmann (2001) \textit{read} Rowe (1979) (the evidential argument from evil) \\
\dag\,2412 & proof of Theorem~5.17, l.~26 & \textit{for} Definition~5.5 \textit{read} Definition~5.1 \\
2412 & l.~29 & \textit{for} p.~2293 \textit{read} p.~2393 (Proposition~5.10) \\
2412 & n.~1 & \textit{for} Bergmann (2000) \textit{read} Bergmann (2001) \\
2413 & l.~13 & \textit{for} conceptoin \textit{read} conception \\
\end{errlist}

\subsection*{Chapter~6. Divine hiddenness \hfill \normalfont\small pp.~2414--2458}
\begin{errlist}
2416 & l.~29 & \textit{for} Lemma~3180 of Volume~X \textit{read} Lemma~3185 of Volume~X; the result cited is on p.~3895 \\
2417 & l.~6 & \textit{for} revelatoin \textit{read} revelation \\
2417 & l.~33 & \textit{for} $T \to \neg\forall e\,\mathrm{Grat}(e)$ \textit{read} $T \to \neg\exists e\,\mathrm{Grat}(e)$ \\
2419 & l.~1 from foot & \textit{for} rgiorous \textit{read} rigorous \\
2420 & l.~12 & \textit{for} corresponidng \textit{read} corresponding \\
2421 & l.~30 & \textit{for} Lemma~3071 of Volume~X \textit{read} Lemma~3077 of Volume~X; the result cited is on p.~3880 \\
2422 & l.~23 & \textit{for} treatmnet \textit{read} treatment \\
\dag\,2422 & proof of Lemma~6.2, l.~32 & \textit{for} Lemma~3055 of Volume~X \textit{read} Lemma~3053 of Volume~X \\
2422 & l.~3 from foot & \textit{for} evidantial \textit{read} evidential \\
2423 & l.~2 & \textit{for} impossbile \textit{read} impossible \\
2423 & l.~22 & \textit{for} epistmeic \textit{read} epistemic \\
2423 & l.~24 & \textit{for} revleation \textit{read} revelation \\
2423 & l.~30 & \textit{for} equivaelnt \textit{read} equivalent \\
2425 & l.~18 & \textit{for} Lemma~3080 of Volume~X \textit{read} Lemma~3086 of Volume~X \\
2425 & l.~26 & \textit{for} p.~2250 \textit{read} p.~2350 (Definition~4.15) \\
2426 & l.~27 & \textit{for} recommned \textit{read} recommend \\
2427 & l.~7 & \textit{for} palusible \textit{read} plausible \\
2427 & l.~8 from foot & \textit{for} excetpion \textit{read} exception \\
2428 & l.~15 & \textit{for} Lemma~3216 of Volume~X \textit{read} Lemma~3211 of Volume~X \\
2428 & l.~31 & \textit{for} Definition~6.2 of Volume~IV \textit{read} Definition~6.1 of Volume~IV; the result cited is on p.~1587 \\
2429 & l.~8 & \textit{for} p.~2325 \textit{read} p.~2225 (Theorem~2.2) \\
\dag\,2429 & proof of Lemma~6.4, l.~27 & \textit{for} Definition~6.3 \textit{read} Definition~6.1; the step uses Definition~6.1 (p.~2419) \\
\dag\,2430 & proof of Lemma~6.4, l.~4 & \textit{for} Lemma~3137 of Volume~X \textit{read} Lemma~3143 of Volume~X \\
2430 & l.~20 & \textit{for} Lemma~3050 of Volume~X \textit{read} Lemma~3053 of Volume~X; the result cited is on p.~3877 \\
2430 & l.~28 & \textit{for} Lemma~3159 of Volume~X \textit{read} Lemma~3156 of Volume~X; the result cited is on p.~3891 \\
2431 & l.~6 & \textit{delete} the second occurrence of \emph{a} \\
2431 & l.~33 & \textit{for} Lemma~5.17 of Volume~III \textit{read} Lemma~5.19 of Volume~III \\
2431 & l.~37 & \textit{for} Lemma~3079 of Volume~X \textit{read} Lemma~3075 of Volume~X; the result cited is on p.~3880 \\
\dag\,2433 & proof of Proposition~6.5, l.~7 & \textit{for} Lemma~6.7 \textit{read} Lemma~6.2 \\
\dag\,2433 & proof of Proposition~6.8, l.~17 & \textit{for} Definition~6.3 \textit{read} Definition~6.6 \\
2433 & l.~22 & \textit{for} phenomeonn \textit{read} phenomenon \\
\dag\,2433 & proof of Proposition~6.5, l.~23 & \textit{for} Lemma~6.2 \textit{read} Lemma~6.4; the step uses Lemma~6.4 (p.~2429) \\
2433 & l.~26 & \textit{for} additioanl \textit{read} additional \\
2433 & n.~3 & \textit{for} Schellenberg (1994) \textit{read} Schellenberg (1993) \\
2434 & l.~17 & \textit{for} nnoresistant \textit{read} nonresistant \\
2434 & l.~7 from foot & \textit{for} Proposition~2.5 of Volume~IV \textit{read} Proposition~2.6 of Volume~IV; the result cited is on p.~1425 \\
2437 & l.~27 & \textit{for} Lemma~3252 of Volume~X \textit{read} Lemma~3255 of Volume~X \\
2438 & l.~1 & \textit{for} straetgy \textit{read} strategy \\
2438 & l.~9 & \textit{for} presuppossition \textit{read} presupposition \\
\dag\,2438 & proof of Lemma~6.10, l.~17 & \textit{for} Definition~6.6 \textit{read} Definition~6.9 \\
\dag\,2438 & proof of Lemma~6.10, l.~3 from foot & \textit{for} Lemma~3183 of Volume~X \textit{read} Lemma~3188 of Volume~X \\
2439 & l.~23 & \textit{for} notaiton \textit{read} notation \\
\dag\,2439 & proof of Lemma~6.10, l.~27 & \textit{for} Definition~6.6 \textit{read} Definition~6.1; the step uses Definition~6.1 (p.~2419) \\
2440 & l.~24 & \textit{for} Schellenbrg (1993) \textit{read} Schellenberg (1993) \\
2440 & l.~4 from foot & \textit{for} p.~2420 \textit{read} p.~2402 (Proposition~5.13) \\
\dag\,2442 & proof of Corollary~6.11, l.~7 & \textit{for} Lemma~6.7 \textit{read} Lemma~6.4 \\
2442 & l.~10 & \textit{for} oppnent \textit{read} opponent \\
2442 & l.~17 & \textit{for} p.~2347 \textit{read} p.~2437 (Definition~6.9) \\
\dag\,2442 & proof of Theorem~6.12, l.~20 & \textit{for} Lemma~3217 of Volume~X \textit{read} Lemma~3216 of Volume~X \\
\dag\,2442 & proof of Theorem~6.12, l.~22 & \textit{for} Proposition~6.5 \textit{read} Proposition~6.8 \\
\dag\,2443 & proof of Theorem~6.12, l.~10 & \textit{for} Lemma~6.7 \textit{read} Lemma~6.10; the step uses Lemma~6.10 (p.~2438) \\
2443 & l.~23 & \textit{for} Lemma~3198 of Volume~X \textit{read} Lemma~3195 of Volume~X \\
2444 & l.~25 & \textit{for} aprpopriate \textit{read} appropriate \\
2445 & l.~21 & \textit{for} Definition~6.7 of Volume~IV \textit{read} Definition~6.1 of Volume~IV; the result cited is on p.~1587 \\
2445 & l.~7 from foot & \textit{for} threeby \textit{read} thereby \\
2446 & l.~10 & \textit{for} Lemma~3250 of Volume~X \textit{read} Lemma~3244 of Volume~X; the result cited is on p.~3904 \\
2446 & l.~37 & \textit{for} Schellenberg (1992) \textit{read} Schellenberg (1993) \\
2447 & l.~16 & \textit{for} Lemma~3142 of Volume~X \textit{read} Lemma~3137 of Volume~X \\
2447 & l.~21 & \textit{for} p.~2519 \textit{read} p.~2419 (Definition~6.1) \\
2448 & l.~1 & \textit{for} Lemma~3191 of Volume~X \textit{read} Lemma~3187 of Volume~X; the result cited is on p.~3896 \\
2448 & l.~30 & \textit{for} contardiction \textit{read} contradiction \\
2449 & l.~32 & \textit{for} anticpiate \textit{read} anticipate \\
2449 & l.~6 from foot & \textit{for} p.~2434 \textit{read} p.~2433 (Proposition~6.8) \\
2450 & l.~2 & \textit{for} whenevr \textit{read} whenever \\
2450 & l.~4 from foot & \textit{for} cosnensus \textit{read} consensus \\
2451 & l.~19 & \textit{for} p.~2208 \textit{read} p.~2308 (Definition~3.16) \\
2451 & l.~20 & \textit{for} absolutley \textit{read} absolutely \\
2451 & l.~30 & \textit{for} thoedicy \textit{read} theodicy \\
2451 & l.~1 from foot & \textit{for} discussoin \textit{read} discussion \\
2452 & l.~14 & \textit{for} contemporry \textit{read} contemporary \\
2452 & l.~23 & \textit{for} p.~2386 \textit{read} p.~2368 (Definition~5.3) \\
2453 & l.~1 & \textit{for} disticntion \textit{read} distinction \\
2454 & l.~8 & \textit{for} sentance \textit{read} sentence \\
2454 & l.~10 & \textit{for} cnotrary \textit{read} contrary \\
2457 & l.~26 & \textit{for} p.~2396 \textit{read} p.~2395 (Theorem~5.11) \\
2458 & l.~20 & \textit{for} symemtric \textit{read} symmetric \\
\end{errlist}

\subsection*{Chapter~7. Hiddenness as a signalling equilibrium \hfill \normalfont\small pp.~2459--2510}
\begin{errlist}
2460 & l.~10 & \textit{for} Lemma~5.21 of Volume~III \textit{read} Lemma~5.17 of Volume~III \\
2462 & l.~1 & \textit{for} resitsance \textit{read} resistance \\
2462 & l.~12 & \textit{for} signaling \textit{read} signalling \\
2464 & l.~19 & \textit{for} Lemma~3049 of Volume~X \textit{read} Lemma~3053 of Volume~X \\
2466 & l.~1 & \textit{for} regadred \textit{read} regarded \\
2467 & l.~19 & \textit{for} Lemma~3121 of Volume~X \textit{read} Lemma~3119 of Volume~X \\
2467 & l.~26 & \textit{for} obivous \textit{read} obvious \\
2468 & l.~22 & \textit{for} Lemma~3213 of Volume~X \textit{read} Lemma~3216 of Volume~X; the result cited is on p.~3900 \\
2468 & l.~28 & \textit{for} Lemma~3270 of Volume~X \textit{read} Lemma~3266 of Volume~X \\
2468 & l.~6 from foot & \textit{for} elsehwere \textit{read} elsewhere \\
2469 & l.~13 & \textit{for} Lemma~3176 of Volume~X \textit{read} Lemma~3172 of Volume~X \\
2469 & l.~17 & \textit{for} intenedd \textit{read} intended \\
2471 & l.~14 & \textit{for} inidvidual \textit{read} individual \\
2471 & l.~17 & \textit{delete} the second occurrence of \emph{of} \\
2471 & l.~20 & \textit{for} abtsract \textit{read} abstract \\
2471 & l.~34 & \textit{for} Lemma~3156 of Volume~X \textit{read} Lemma~3150 of Volume~X \\
2472 & l.~20 & \textit{for} inscrutalbe \textit{read} inscrutable \\
2474 & l.~12 & \textit{for} p.~2303 \textit{read} p.~2304 (Proposition~3.14) \\
2474 & l.~22 & \textit{for} ahcieve \textit{read} achieve \\
2474 & l.~24 & \textit{for} bewteen \textit{read} between \\
2474 & l.~6 from foot & \textit{for} suffernig \textit{read} suffering \\
\dag\,2474 & proof of Corollary~7.2, l.~40 & \textit{for} Lemma~3066 of Volume~X \textit{read} Lemma~3064 of Volume~X; the result cited is on p.~3878 \\
2475 & l.~8 & \textit{for} p.~2228 \textit{read} p.~2328 (Proposition~4.4) \\
2475 & l.~2 from foot & \textit{for} Lemma~3182 of Volume~X \textit{read} Lemma~3185 of Volume~X \\
2477 & l.~9 & \textit{for} obejctive \textit{read} objective \\
2477 & l.~11 & \textit{for} p.~2433 \textit{read} p.~2343 (Corollary~4.11) \\
2477 & l.~22 & \textit{for} symmtric \textit{read} symmetric \\
2477 & l.~27 & \textit{for} emplyoed \textit{read} employed \\
2479 & l.~14 & \textit{for} contmeporary \textit{read} contemporary \\
2479 & l.~20 & \textit{for} cohernece \textit{read} coherence \\
2479 & l.~21 & \textit{for} primarrily \textit{read} primarily \\
2479 & l.~28 & \textit{for} Definition~6.8 of Volume~IV \textit{read} Definition~6.7 of Volume~IV; the result cited is on p.~1610 \\
2480 & l.~3 & \textit{for} p.~2434 \textit{read} p.~2433 (Proposition~6.5) \\
2480 & l.~14 & \textit{for} subsequant \textit{read} subsequent \\
2481 & l.~32 & \textit{for} succesfsul \textit{read} successful \\
2483 & Lemma~7.5 & \textit{see} Corrigendum~S12 \\
2484 & l.~4 & \textit{for} consenuss \textit{read} consensus \\
\dag\,2484 & proof of Lemma~7.5, l.~8 & \textit{for} Corollary~5.12 of Volume~III \textit{read} Corollary~5.9 of Volume~III \\
2488 & l.~35 & \textit{for} hidenness \textit{read} hiddenness \\
2489 & l.~8 & \textit{for} objcetion \textit{read} objection \\
2489 & l.~20 & \textit{for} Lemma~3128 of Volume~X \textit{read} Lemma~3126 of Volume~X; the result cited is on p.~3887 \\
\dag\,2491 & proof of Theorem~7.7, l.~13 & \textit{for} Lemma~7.9 \textit{read} Lemma~7.5 \\
\dag\,2492 & proof of Lemma~7.8, l.~8 & \textit{for} Definition~7.1 \textit{read} Definition~7.3; the step uses Definition~7.3 (p.~2474) \\
2492 & l.~17 & \textit{for} objctive \textit{read} objective \\
2493 & l.~12 & \textit{for} p.~2413 \textit{read} p.~2412 (Definition~5.16) \\
2493 & l.~16 & \textit{for} langauge \textit{read} language \\
2493 & l.~26 & \textit{for} Proposition~5.7 of Volume~III \textit{read} Proposition~5.8 of Volume~III; the result cited is on p.~991 \\
2493 & l.~38 & \textit{for} Lemma~3107 of Volume~X \textit{read} Lemma~3102 of Volume~X \\
2494 & l.~12 & \textit{for} p.~2383 \textit{read} p.~2483 (Lemma~7.5) \\
2494 & l.~25 & \textit{for} proceudre \textit{read} procedure \\
2495 & l.~19 & \textit{for} foritori \textit{read} fortiori \\
2495 & l.~28 & \textit{for} correspodning \textit{read} corresponding \\
\dag\,2496 & proof of Lemma~7.9, l.~28 & \textit{for} Lemma~7.8 \textit{read} Lemma~7.5 \\
2496 & l.~31 & \textit{for} perhpas \textit{read} perhaps \\
\dag\,2497 & proof of Corollary~7.10, l.~17 & \textit{for} Lemma~3267 of Volume~X \textit{read} Lemma~3263 of Volume~X \\
2497 & l.~22 & \textit{for} unievrsal \textit{read} universal \\
2498 & l.~4 & \textit{delete} the second occurrence of \emph{a} \\
\dag\,2498 & proof of Lemma~7.12, l.~8 & \textit{for} Lemma~7.14 \textit{read} Lemma~7.9 \\
\dag\,2498 & proof of Lemma~7.12, l.~20 & \textit{for} Lemma~3287 of Volume~X \textit{read} Lemma~3284 of Volume~X; the result cited is on p.~3910 \\
2498 & l.~23 & \textit{for} despitte \textit{read} despite \\
2498 & l.~38 & \textit{for} equillibrium \textit{read} equilibrium \\
2499 & l.~27 & \textit{for} Lemma~3090 of Volume~X \textit{read} Lemma~3096 of Volume~X \\
\dag\,2500 & proof of Lemma~7.15, l.~6 & \textit{for} Lemma~7.8 \textit{read} Lemma~7.9; the step uses Lemma~7.9 (p.~2496) \\
\dag\,2500 & proof of Lemma~7.14, l.~21 & \textit{for} Corollary~7.10 \textit{read} Corollary~7.13 \\
2500 & l.~34 & \textit{for} Lemma~3288 of Volume~X \textit{read} Lemma~3292 of Volume~X \\
2501 & l.~3 & \textit{for} nonresitsant \textit{read} nonresistant \\
\dag\,2501 & proof of Corollary~7.13, l.~12 & \textit{for} Definition~7.16 \textit{read} Definition~7.11 \\
\dag\,2501 & proof of Lemma~7.15, l.~26 & \textit{for} Lemma~7.12 \textit{read} Lemma~7.8 \\
\dag\,2501 & proof of Lemma~7.14, l.~28 & \textit{for} Lemma~7.9 \textit{read} Lemma~7.12 \\
2501 & l.~29 & \textit{for} deliberte \textit{read} deliberate \\
2502 & l.~7 & \textit{for} signifciant \textit{read} significant \\
2502 & l.~12 & \textit{for} p.~2498 \textit{read} p.~2497 (Definition~7.11) \\
2502 & l.~38 & \textit{for} Lemma~3203 of Volume~X \textit{read} Lemma~3209 of Volume~X \\
2503 & l.~1 from foot & \textit{for} Lemma~3245 of Volume~X \textit{read} Lemma~3249 of Volume~X; the result cited is on p.~3905 \\
2504 & l.~11 & \textit{for} p.~2400 \textit{read} p.~2500 (Corollary~7.13) \\
2505 & l.~2 & \textit{for} validty \textit{read} validity \\
2505 & l.~5 & \textit{for} Lemma~3195 of Volume~X \textit{read} Lemma~3197 of Volume~X \\
2505 & l.~23 & \textit{for} Lemma~3197 of Volume~X \textit{read} Lemma~3198 of Volume~X \\
2506 & l.~7 from foot & \textit{for} somtehing \textit{read} something \\
2507 & l.~17 & \textit{for} indvidual \textit{read} individual \\
2508 & l.~13 & \textit{for} p.~2434 \textit{read} p.~2433 (Definition~6.6) \\
2508 & l.~29 & \textit{for} p.~2516 \textit{read} p.~2506 (Definition~7.16) \\
2508 & l.~32 & \textit{for} conceviable \textit{read} conceivable \\
2509 & l.~2 & \textit{for} p.~2597 \textit{read} p.~2497 (Corollary~7.10) \\
2509 & l.~24 & \textit{for} p.~2508 \textit{read} p.~2507 (Definition~7.17) \\
2509 & display (7.12) & \textit{for} $\exists u\,(T(u) \vee E(u))$ \textit{read} $\exists u\,(T(u) \wedge E(u))$; the formula is used in its correct form on p.~2510 \\
2510 & l.~27 & \textit{for} $\exists w_2\,(T(w_1) \wedge E(w_1))$ \textit{read} $\exists w_1\,(T(w_1) \wedge E(w_1))$ \\
2510 & l.~35 & \textit{for} ojbective \textit{read} objective \\
\end{errlist}

\subsection*{Chapter~8. Evil, the possibility premise and the converse \hfill \normalfont\small pp.~2511--2580}
\begin{errlist}
2512 & l.~20 & \textit{for} p.~2217 \textit{read} p.~2317 (Theorem~4.2) \\
2513 & l.~11 & \textit{for} Lemma~3193 of Volume~X \textit{read} Lemma~3195 of Volume~X; the result cited is on p.~3897 \\
2515 & l.~10 & \textit{for} p.~2592 \textit{read} p.~2492 (Lemma~7.8) \\
\dag\,2515 & proof of Lemma~8.2, (8.2) & \textit{for} $\exists z\,\mathrm{MG}(z,v) \to \exists z\,(E(z,v) \vee \mathrm{ME}(z,v))$ \textit{read} $\exists z\,\mathrm{MG}(z,v) \to \exists z\,(E(z,v) \wedge \mathrm{ME}(z,v))$; correctly printed in (8.3), p.~2517 \\
2515 & l.~23 & \textit{for} objectvie \textit{read} objective \\
2515 & l.~27 & \textit{for} Plantinga (1974b) \textit{read} Rowe (1979) (the evidential argument from evil) \\
2516 & l.~18 & \textit{for} Lemma~3175 of Volume~X \textit{read} Lemma~3181 of Volume~X \\
\dag\,2516 & proof of Lemma~8.2, l.~25 & \textit{for} Lemma~3193 of Volume~X \textit{read} Lemma~3189 of Volume~X; the result cited is on p.~3896 \\
2518 & display (8.4) & \textit{for} $\mathrm{NE}(z,u) \iff \forall v'\,(uRv' \to E(z,u))$ \textit{read} $\mathrm{NE}(z,u) \iff \forall v'\,(uRv' \to E(z,v'))$; the formula is used in its correct form on p.~2522 \\
2518 & l.~2 from foot & \textit{for} Lemma~3177 of Volume~X \textit{read} Lemma~3180 of Volume~X \\
2519 & l.~21 & \textit{for} Lemma~3170 of Volume~X \textit{read} Lemma~3168 of Volume~X; the result cited is on p.~3893 \\
2519 & l.~32 & \textit{for} commimtent \textit{read} commitment \\
2519 & l.~3 from foot & \textit{for} Lemma~3055 of Volume~X \textit{read} Lemma~3059 of Volume~X \\
2520 & l.~4 & \textit{for} maximmally \textit{read} maximally \\
2521 & l.~3 from foot & \textit{for} Lemma~3146 of Volume~X \textit{read} Lemma~3142 of Volume~X \\
2522 & l.~8 & \textit{for} Lemma~3060 of Volume~X \textit{read} Lemma~3066 of Volume~X \\
\dag\,2522 & proof of Corollary~8.4, (8.10) & \textit{for} $\forall v\,(T(v) \wedge E(v))$ \textit{read} $\exists v\,(T(v) \wedge E(v))$ \\
\dag\,2522 & proof of Corollary~8.4, l.~21 & \textit{for} Lemma~3129 of Volume~X \textit{read} Lemma~3127 of Volume~X \\
\dag\,2522 & proof of Corollary~8.4, (8.8) & \textit{for} $\pos\neg\forall y\,\mathrm{MG}(y) \to \neg\exists y\,\mathrm{MG}(y)$ \textit{read} $\pos\neg\exists y\,\mathrm{MG}(y) \to \neg\exists y\,\mathrm{MG}(y)$; correctly printed in (8.11), p.~2527 \\
2525 & l.~17 & \textit{for} univesral \textit{read} universal \\
\dag\,2525 & proof of Lemma~8.5, l.~24 & \textit{for} Definition~8.1 \textit{read} Definition~8.3; the step uses Definition~8.3 (p.~2519) \\
2525 & l.~35 & \textit{for} Volume~IV, Chapter~6 \textit{read} Volume~IV, Chapter~5 (moral perfection) \\
2525 & l.~37 & \textit{for} distingusih \textit{read} distinguish \\
2526 & l.~3 & \textit{for} p.~2499 \textit{read} p.~2498 (Lemma~7.12) \\
2526 & l.~7 & \textit{for} deficeincy \textit{read} deficiency \\
\dag\,2526 & proof of Proposition~8.6, l.~9 & \textit{for} Lemma~8.8 \textit{read} Lemma~8.5 \\
\dag\,2526 & proof of Proposition~8.6, l.~7 from foot & \textit{for} Lemma~3203 of Volume~X \textit{read} Lemma~3201 of Volume~X \\
\dag\,2526 & proof of Proposition~8.6, l.~1 from foot & \textit{for} Lemma~3194 of Volume~X \textit{read} Lemma~3191 of Volume~X \\
\dag\,2527 & proof of Proposition~8.6, l.~18 & \textit{for} Lemma~6.6 of Volume~IV \textit{read} Lemma~6.12 of Volume~IV; the result cited is on p.~1622 \\
\dag\,2527 & proof of Proposition~8.6, l.~28 & \textit{for} $\mathrm{NE}(y,w_2) \iff \forall u\,(uRw_2 \to E(y,u))$ \textit{read} $\mathrm{NE}(y,w_2) \iff \forall u\,(w_2Ru \to E(y,u))$ \\
2527 & l.~35 & \textit{for} Oppy (2006) \textit{read} Rowe (1979) (the evidential argument from evil) \\
2527 & l.~5 from foot & \textit{for} stritcly \textit{read} strictly \\
2527 & l.~3 from foot & \textit{for} Volume~V, Chapter~4 \textit{read} Volume~V, Chapter~3 (the question-begging objection) \\
2528 & l.~4 & \textit{for} greatenss \textit{read} greatness \\
2528 & l.~25 & \textit{for} distinctoin \textit{read} distinction \\
2528 & l.~8 from foot & \textit{for} Lemma~3172 of Volume~X \textit{read} Lemma~3167 of Volume~X \\
2529 & l.~5 from foot & \textit{for} p.~2254 \textit{read} p.~2524 (Lemma~8.5) \\
2529 & l.~3 from foot & \textit{for} determintae \textit{read} determinate \\
2529 & l.~2 from foot & \textit{for} consenss \textit{read} consensus \\
\dag\,2530 & proof of Theorem~8.7, l.~15 & \textit{for} Proposition~8.11 \textit{read} Proposition~8.6; the step uses Proposition~8.6 (p.~2526) \\
2530 & l.~23 & \textit{for} futrher \textit{read} further \\
2530 & l.~29 & \textit{for} aanlogous \textit{read} analogous \\
\dag\,2530 & proof of Theorem~8.7, l.~8 from foot & \textit{for} Lemma~3229 of Volume~X \textit{read} Lemma~3231 of Volume~X; the result cited is on p.~3902 \\
\dag\,2531 & proof of Theorem~8.7, l.~8 from foot & \textit{for} Lemma~3190 of Volume~X \textit{read} Lemma~3196 of Volume~X \\
2533 & l.~13 & \textit{for} Lemma~5.19 of Volume~III \textit{read} Lemma~5.17 of Volume~III; the result cited is on p.~1003 \\
\dag\,2534 & proof of Lemma~8.8, (8.13) & \textit{for} $\mathrm{MG}(z,w_2) \wedge w_1Ru \to \mathrm{MG}(z,u)$ \textit{read} $\mathrm{MG}(z,w_2) \wedge w_2Ru \to \mathrm{MG}(z,u)$; correctly printed in (8.16), p.~2539 \\
\dag\,2535 & proof of Lemma~8.8, l.~7 & \textit{for} Lemma~3274 of Volume~X \textit{read} Lemma~3279 of Volume~X \\
2535 & l.~8 & \textit{for} excption \textit{read} exception \\
\dag\,2536 & proof of Theorem~8.10, l.~12 & \textit{for} Lemma~8.5 \textit{read} Lemma~8.2; the step uses Lemma~8.2 (p.~2515) \\
\dag\,2536 & proof of Theorem~8.10, l.~17 & \textit{for} $\pos\forall y\,\mathrm{MG}(y) \to \nec\exists y\,\mathrm{MG}(y)$ \textit{read} $\pos\exists y\,\mathrm{MG}(y) \to \nec\exists y\,\mathrm{MG}(y)$ \\
\dag\,2536 & proof of Theorem~8.10, (8.14) & \textit{for} $\mathrm{MG}(z,u) \to \mathrm{NE}(z,u) \vee \mathrm{NME}(z,u)$ \textit{read} $\mathrm{MG}(z,u) \to \mathrm{NE}(z,u) \wedge \mathrm{NME}(z,u)$; correctly printed in (8.17), p.~2542 \\
2536 & l.~28 & \textit{for} nevrtheless \textit{read} nevertheless \\
\dag\,2536 & proof of Theorem~8.10, l.~30 & \textit{for} Theorem~8.12 \textit{read} Theorem~8.7 \\
\dag\,2537 & proof of Theorem~8.10, (8.15) & \textit{for} $\mathrm{NME}(x,w_0) \iff \exists v\,(w_0Rv \to \mathrm{ME}(x,v))$ \textit{read} $\mathrm{NME}(x,w_0) \iff \forall v\,(w_0Rv \to \mathrm{ME}(x,v))$ \\
2538 & l.~6 from foot & \textit{for} omnipotant \textit{read} omnipotent \\
2542 & display (8.18) & \textit{for} $\mathrm{MG}(z,u) \wedge uRw_1 \to \mathrm{MG}(z,w_2)$ \textit{read} $\mathrm{MG}(z,u) \wedge uRw_1 \to \mathrm{MG}(z,w_1)$ \\
2542 & l.~7 & \textit{for} Lemma~3059 of Volume~X \textit{read} Lemma~3060 of Volume~X; the result cited is on p.~3878 \\
2543 & l.~29 & \textit{for} Lemma~3171 of Volume~X \textit{read} Lemma~3168 of Volume~X \\
2544 & l.~17 & \textit{for} $\mathrm{NE}(y,w') \iff \forall w'\,(wRw' \to E(y,w'))$ \textit{read} $\mathrm{NE}(y,w) \iff \forall w'\,(wRw' \to E(y,w'))$ \\
\dag\,2545 & proof of Lemma~8.13, l.~5 & \textit{for} Theorem~8.10 \textit{read} Theorem~8.12 \\
2545 & l.~26 & \textit{for} p.~2423 \textit{read} p.~2433 (Lemma~6.7) \\
\dag\,2545 & proof of Theorem~8.12, l.~30 & \textit{for} Theorem~8.7 \textit{read} Theorem~8.10; the step uses Theorem~8.10 (p.~2536) \\
2545 & l.~7 from foot & \textit{for} obvoius \textit{read} obvious \\
2545 & l.~2 from foot & \textit{for} soultion \textit{read} solution \\
2546 & l.~3 & \textit{after} maximally excellent \textit{insert} at $w$ \\
\dag\,2546 & proof of Proposition~8.14, l.~19 & \textit{for} Lemma~8.8 \textit{read} Lemma~8.13 \\
\dag\,2546 & proof of Proposition~8.14, (8.20) & \textit{for} $\pos\exists z\,\mathrm{MG}(z) \to \neg\exists z\,\mathrm{MG}(z)$ \textit{read} $\pos\neg\exists z\,\mathrm{MG}(z) \to \neg\exists z\,\mathrm{MG}(z)$; the formula is used in its correct form on p.~2548 \\
2547 & l.~4 from foot & \textit{for} p.~2340 \textit{read} p.~2330 (Corollary~4.7) \\
2548 & display (8.23) & \textit{for} $\mathrm{MG}(y,w_0) \to \mathrm{NE}(y,w_0) \vee \mathrm{NME}(y,w_0)$ \textit{read} $\mathrm{MG}(y,w_0) \to \mathrm{NE}(y,w_0) \wedge \mathrm{NME}(y,w_0)$ \\
2550 & display (8.26) & \textit{for} $\mathrm{MG}(x,v') \wedge v'Ru \leftrightarrow \mathrm{MG}(x,u)$ \textit{read} $\mathrm{MG}(x,v') \wedge v'Ru \to \mathrm{MG}(x,u)$ \\
\dag\,2551 & proof of Theorem~8.16, l.~12 & \textit{for} Theorem~8.10 \textit{read} Theorem~8.15 \\
2551 & l.~19 & \textit{for} consqeuence \textit{read} consequence \\
2551 & l.~22 & \textit{for} neverhteless \textit{read} nevertheless \\
2552 & l.~6 & \textit{for} Lemma~3131 of Volume~X \textit{read} Lemma~3136 of Volume~X; the result cited is on p.~3888 \\
\dag\,2552 & proof of Theorem~8.16, l.~15 & \textit{for} Definition~5.5 of Volume~III \textit{read} Definition~5.10 of Volume~III; the result cited is on p.~991 \\
2552 & l.~24 & \textit{for} nonressitant \textit{read} nonresistant \\
2553 & l.~23 & \textit{for} Lemma~3292 of Volume~X \textit{read} Lemma~3289 of Volume~X; the result cited is on p.~3910 \\
2553 & l.~30 & \textit{for} Lemma~3158 of Volume~X \textit{read} Lemma~3157 of Volume~X \\
2554 & l.~4 from foot & \textit{for} Lemma~3056 of Volume~X \textit{read} Lemma~3054 of Volume~X; the result cited is on p.~3877 \\
\dag\,2556 & proof of Corollary~8.18, l.~19 & \textit{for} Definition~5.10 of Volume~III \textit{read} Definition~5.13 of Volume~III; the result cited is on p.~995 \\
2556 & l.~24 & \textit{for} Lemma~3099 of Volume~X \textit{read} Lemma~3103 of Volume~X \\
\dag\,2556 & proof of Corollary~8.18, l.~29 & \textit{for} Definition~8.19 \textit{read} Definition~8.17; the step uses Definition~8.17 (p.~2556) \\
2557 & l.~16 & \textit{for} unviersal \textit{read} universal \\
\dag\,2557 & proof of Corollary~8.18, (8.29) & \textit{for} $@ \in \pos\,\exists z\, \mathrm{MG}(z,\cdot)$ \textit{read} $@ \in \nec\,\exists z\, \mathrm{MG}(z,\cdot)$; the formula is used in its correct form on p.~2561 \\
2559 & l.~23 & \textit{for} Lemma~3157 of Volume~X \textit{read} Lemma~3152 of Volume~X \\
2560 & l.~33 & \textit{for} $@ \notin \pos\,\exists y\, \mathrm{MG}(y,\cdot)$ \textit{read} $@ \in \pos\,\exists y\, \mathrm{MG}(y,\cdot)$ \\
2560 & l.~35 & \textit{for} Lemma~3280 of Volume~X \textit{read} Lemma~3275 of Volume~X \\
2561 & l.~21 & \textit{for} Lemma~3242 of Volume~X \textit{read} Lemma~3243 of Volume~X \\
2561 & l.~6 from foot & \textit{for} Lemma~3238 of Volume~X \textit{read} Lemma~3232 of Volume~X; the result cited is on p.~3902 \\
2562 & display (8.34) & \textit{for} $\mathrm{MG}(y,w_2) \leftrightarrow \mathrm{NE}(y,w_2) \wedge \mathrm{NME}(y,w_2)$ \textit{read} $\mathrm{MG}(y,w_2) \to \mathrm{NE}(y,w_2) \wedge \mathrm{NME}(y,w_2)$; correctly printed in (8.35), p.~2565 \\
2562 & display (8.33) & \textit{for} $\mathrm{NME}(y,v) \iff \exists w'\,(vRw' \to \mathrm{ME}(y,w'))$ \textit{read} $\mathrm{NME}(y,v) \iff \forall w'\,(vRw' \to \mathrm{ME}(y,w'))$ \\
2562 & l.~2 from foot & \textit{for} instanitated \textit{read} instantiated \\
2564 & l.~25 & \textit{for} mechnaism \textit{read} mechanism \\
2568 & l.~27 & \textit{after} exists \textit{insert} at every world accessible from $w_1$ \\
2568 & l.~31 & \textit{for} threby \textit{read} thereby \\
\dag\,2570 & proof of Proposition~8.21, l.~27 & \textit{for} Proposition~8.11 \textit{read} Proposition~8.14 \\
\dag\,2570 & proof of Proposition~8.21, l.~29 & \textit{for} Lemma~3283 of Volume~X \textit{read} Lemma~3287 of Volume~X; the result cited is on p.~3910 \\
\dag\,2570 & proof of Proposition~8.21, l.~34 & \textit{for} Lemma~3100 of Volume~X \textit{read} Lemma~3094 of Volume~X \\
2571 & l.~2 & \textit{for} p.~2344 \textit{read} p.~2244 (Proposition~2.9) \\
\dag\,2571 & proof of Proposition~8.21, (8.38) & \textit{for} $\pos\neg\exists x\,\mathrm{MG}(x) \leftrightarrow \neg\exists x\,\mathrm{MG}(x)$ \textit{read} $\pos\neg\exists x\,\mathrm{MG}(x) \to \neg\exists x\,\mathrm{MG}(x)$; the formula is used in its correct form on p.~2573 \\
2571 & l.~24 & \textit{for} Lemma~3289 of Volume~X \textit{read} Lemma~3294 of Volume~X; the result cited is on p.~3911 \\
2571 & l.~27 & \textit{for} inscrtable \textit{read} inscrutable \\
2572 & display (8.39) & \textit{for} $\mathrm{MG}(z,w') \iff \forall w'\,(uRw' \to E(z,w') \wedge \mathrm{ME}(z,w'))$ \textit{read} $\mathrm{MG}(z,u) \iff \forall w'\,(uRw' \to E(z,w') \wedge \mathrm{ME}(z,w'))$ \\
2574 & display (8.40) & \textit{for} $\mathrm{MG}(y,w') \iff \exists u\,(w'Ru \to E(y,u) \wedge \mathrm{ME}(y,u))$ \textit{read} $\mathrm{MG}(y,w') \iff \forall u\,(w'Ru \to E(y,u) \wedge \mathrm{ME}(y,u))$; the formula is used in its correct form on p.~2575 \\
2574 & l.~7 from foot & \textit{for} Lemma~3224 of Volume~X \textit{read} Lemma~3230 of Volume~X \\
2575 & l.~4 & \textit{for} Lemma~3289 of Volume~X \textit{read} Lemma~3291 of Volume~X \\
2575 & l.~20 & \textit{for} opopnent \textit{read} opponent \\
2575 & l.~28 & \textit{for} Lemma~3283 of Volume~X \textit{read} Lemma~3282 of Volume~X; the result cited is on p.~3909 \\
2576 & l.~36 & \textit{for} Lemma~3084 of Volume~X \textit{read} Lemma~3081 of Volume~X \\
2577 & l.~5 & \textit{for} $\exists w_2\,(T(w_2) \vee E(w_2))$ \textit{read} $\exists w_2\,(T(w_2) \wedge E(w_2))$ \\
2578 & display (8.41) & \textit{for} $\pos\exists x\,\mathrm{MG}(x) \leftrightarrow \nec\exists x\,\mathrm{MG}(x)$ \textit{read} $\pos\exists x\,\mathrm{MG}(x) \to \nec\exists x\,\mathrm{MG}(x)$; the formula is used in its correct form on p.~2579 \\
2578 & l.~6 & \textit{for} phenomnon \textit{read} phenomenon \\
2578 & l.~38 & \textit{for} p.~2560 \textit{read} p.~2570 (Proposition~8.21) \\
2579 & l.~15 & \textit{for} coresponding \textit{read} corresponding \\
2579 & l.~24 & \textit{for} $\forall w_1\,(T(w_1) \wedge E(w_1))$ \textit{read} $\exists w_1\,(T(w_1) \wedge E(w_1))$ \\
2579 & display (8.42) & \textit{for} $\mathrm{MG}(z,w') \iff \forall v\,(w'Rv \leftrightarrow E(z,v) \wedge \mathrm{ME}(z,v))$ \textit{read} $\mathrm{MG}(z,w') \iff \forall v\,(w'Rv \to E(z,v) \wedge \mathrm{ME}(z,v))$ \\
2580 & l.~3 & \textit{for} p.~2576 \textit{read} p.~2567 (Definition~8.19) \\
2580 & l.~12 & \textit{for} rseemblance \textit{read} resemblance \\
\end{errlist}

\subsection*{Chapter~9. Summary \hfill \normalfont\small pp.~2581--2602}
\begin{errlist}
2583 & l.~11 & \textit{for} Lemma~3225 of Volume~X \textit{read} Lemma~3231 of Volume~X; the result cited is on p.~3902 \\
2583 & l.~25 & \textit{for} Mackie (1981) \textit{read} Mackie (1982) \\
2584 & l.~1 & \textit{for} Definition~5.1 of Volume~III \textit{read} Definition~5.5 of Volume~III; the result cited is on p.~981 \\
2584 & l.~2 & \textit{for} quetsionable \textit{read} questionable \\
2584 & l.~34 & \textit{for} Lemma~3142 of Volume~X \textit{read} Lemma~3147 of Volume~X; the result cited is on p.~3890 \\
2584 & l.~5 from foot & \textit{for} Mackei (1982) \textit{read} Mackie (1982) \\
2587 & l.~3 & \textit{for} Lemma~3208 of Volume~X \textit{read} Lemma~3204 of Volume~X; the result cited is on p.~3898 \\
2587 & l.~2 from foot & \textit{for} subsequnet \textit{read} subsequent \\
2589 & l.~8 & \textit{for} Lemma~3063 of Volume~X \textit{read} Lemma~3067 of Volume~X \\
2590 & l.~32 & \textit{for} preesntation \textit{read} presentation \\
2591 & l.~32 & \textit{for} Lemma~3073 of Volume~X \textit{read} Lemma~3071 of Volume~X \\
2591 & l.~35 & \textit{for} p.~2419 \textit{read} p.~2519 (Definition~8.3) \\
2591 & l.~37 & \textit{for} resemblnace \textit{read} resemblance \\
2593 & l.~3 & \textit{for} Definition~5.10 of Volume~III \textit{read} Definition~5.5 of Volume~III; the result cited is on p.~981 \\
2593 & l.~8 from foot & \textit{for} absloutely \textit{read} absolutely \\
2594 & l.~25 & \textit{for} preceidng \textit{read} preceding \\
2594 & l.~6 from foot & \textit{for} questionnable \textit{read} questionable \\
2594 & l.~5 from foot & \textit{for} equivalnet \textit{read} equivalent \\
2595 & l.~29 & \textit{for} aditional \textit{read} additional \\
2595 & l.~2 from foot & \textit{for} dififculty \textit{read} difficulty \\
2596 & l.~11 & \textit{for} Lemma~3232 of Volume~X \textit{read} Lemma~3229 of Volume~X; the result cited is on p.~3902 \\
2596 & l.~12 & \textit{for} p.~2444 \textit{read} p.~2544 (Theorem~8.12) \\
2596 & l.~20 & \textit{for} distincsion \textit{read} distinction \\
2596 & l.~8 from foot & \textit{for} presentatoin \textit{read} presentation \\
2597 & l.~11 & \textit{for} Lemma~3192 of Volume~X \textit{read} Lemma~3193 of Volume~X; the result cited is on p.~3897 \\
2598 & l.~5 & \textit{for} akcnowledge \textit{read} acknowledge \\
2598 & l.~13 & \textit{for} additionl \textit{read} additional \\
2598 & l.~38 & \textit{for} anlaogous \textit{read} analogous \\
2598 & l.~39 & \textit{for} conjetcure \textit{read} conjecture \\
2599 & l.~2 & \textit{for} Lemma~3286 of Volume~X \textit{read} Lemma~3290 of Volume~X \\
2599 & l.~6 from foot & \textit{for} Lemma~3092 of Volume~X \textit{read} Lemma~3093 of Volume~X \\
2601 & l.~4 & \textit{for} p.~2512 \textit{read} p.~2513 (Definition~8.1) \\
2601 & l.~14 & \textit{for} Lemma~3229 of Volume~X \textit{read} Lemma~3235 of Volume~X \\
2601 & l.~17 & \textit{for} Macike (1982) \textit{read} Mackie (1982) \\
2601 & l.~8 from foot & \textit{for} Lemma~3129 of Volume~X \textit{read} Lemma~3131 of Volume~X; the result cited is on p.~3888 \\
2602 & l.~16 & \textit{for} p.~2428 \textit{read} p.~2429 (Lemma~6.4) \\
2602 & l.~39 & \textit{for} chapetr \textit{read} chapter \\
\end{errlist}


\section{Volume VII: Objections from the Sciences}\label{s:vol7}
% generated by generate.py -- do not edit
\subsection*{Front matter \hfill \normalfont\small pp.~2603--2605}
\begin{errlist}
2604 & Contents, Chapter~3 & \textit{for} 2751 \textit{read} 2651 \\
\end{errlist}

\subsection*{Chapter~1. Introduction \hfill \normalfont\small pp.~2606--2612}
\begin{errlist}
2608 & l.~38 & \textit{for} Lemma~3092 of Volume~X \textit{read} Lemma~3098 of Volume~X; the result cited is on p.~3883 \\
2609 & l.~3 & \textit{delete} the second occurrence of \emph{to} \\
2610 & l.~26 & \textit{for} Lemma~3266 of Volume~X \textit{read} Lemma~3264 of Volume~X \\
2611 & l.~18 & \textit{delete} the second occurrence of \emph{to} \\
2611 & l.~21 & \textit{for} summray \textit{read} summary \\
2611 & l.~29 & \textit{for} vollume \textit{read} volume \\
2611 & l.~31 & \textit{delete} the second occurrence of \emph{of} \\
2611 & n.~10 & \textit{for} Platninga (2011) \textit{read} Plantinga (2011) \\
2612 & l.~8 & \textit{for} p.~2670 \textit{read} p.~2607 (Definition~1.1) \\
2612 & l.~5 from foot & \textit{for} Plantinga (2010) \textit{read} Plantinga (2011) \\
2612 & l.~3 from foot & \textit{for} p.~2612 \textit{read} p.~2611 (Lemma~1.3) \\
\end{errlist}

\subsection*{Chapter~2. Naturalism and the modal argument \hfill \normalfont\small pp.~2613--2650}
\begin{errlist}
2614 & l.~13 & \textit{delete} the second occurrence of \emph{that} \\
2615 & display (2.1) & \textit{for} $\mathrm{MG}(z,w') \iff \forall w_0\,(w_0Rw' \to E(z,w_0) \wedge \mathrm{ME}(z,w_0))$ \textit{read} $\mathrm{MG}(z,w') \iff \forall w_0\,(w'Rw_0 \to E(z,w_0) \wedge \mathrm{ME}(z,w_0))$; correctly printed in (2.3), p.~2620 \\
2616 & l.~23 & \textit{for} $\mathrm{MG}(y,v') \vee v'Ru \to \mathrm{MG}(y,u)$ \textit{read} $\mathrm{MG}(y,v') \wedge v'Ru \to \mathrm{MG}(y,u)$ \\
2618 & l.~4 & \textit{after} measure \textit{insert} over the space of physical constants \\
2619 & l.~29 & \textit{delete} the second occurrence of \emph{a} \\
2620 & l.~4 & \textit{delete} the second occurrence of \emph{a} \\
2620 & l.~7 & \textit{for} terminolgoy \textit{read} terminology \\
2620 & l.~16 & \textit{for} sloution \textit{read} solution \\
2620 & l.~18 & \textit{for} corrseponding \textit{read} corresponding \\
2620 & l.~29 & \textit{for} exists and is maximally excellent at some world accessible from $w'$ \textit{read} exists and is maximally excellent at every world accessible from $w'$ \\
2620 & l.~5 from foot & \textit{delete} the second occurrence of \emph{of} \\
2621 & display (2.5) & \textit{for} $P(B \mid T) = P(B \mid  T)$ \textit{read} $P(B \mid T) = P(B \mid \neg T)$; the formula is used in its correct form on p.~2622 \\
2621 & l.~5 from foot & \textit{delete} the second occurrence of \emph{a} \\
2624 & l.~5 & \textit{for} $\mathrm{MG}(x,v') \vee v'Ru \to \mathrm{MG}(x,u)$ \textit{read} $\mathrm{MG}(x,v') \wedge v'Ru \to \mathrm{MG}(x,u)$ \\
2624 & display (2.7) & \textit{for} $\mathrm{NME}(x,w) \iff \exists u\,(wRu \to \mathrm{ME}(x,u))$ \textit{read} $\mathrm{NME}(x,w) \iff \forall u\,(wRu \to \mathrm{ME}(x,u))$; correctly printed in (2.8), p.~2625 \\
2624 & l.~15 & \textit{for} multvierse \textit{read} multiverse \\
\dag\,2625 & proof of Lemma~2.3, l.~8 & \textit{for} Definition~2.2 \textit{read} Definition~2.1 \\
2626 & l.~2 & \textit{for} distignuish \textit{read} distinguish \\
2626 & l.~27 & \textit{delete} the second occurrence of \emph{that} \\
\dag\,2627 & proof of Theorem~2.4, (2.9) & \textit{for} $@ \in \pos\,\forall x\, \mathrm{MG}(x,\cdot)$ \textit{read} $@ \in \pos\,\exists x\, \mathrm{MG}(x,\cdot)$ \\
2627 & l.~7 from foot & \textit{for} nautralism \textit{read} naturalism \\
2628 & l.~23 & \textit{for} Lemma~3194 of Volume~X \textit{read} Lemma~3188 of Volume~X; the result cited is on p.~3896 \\
2629 & l.~2 & \textit{delete} the second occurrence of \emph{of} \\
2630 & display (2.10) & \textit{for} $P(T \mid E) > P(T)$ \textit{read} $P(T \mid E) < P(T)$; correctly printed in (2.13), p.~2635 \\
2630 & l.~12 & \textit{delete} the second occurrence of \emph{a} \\
2630 & l.~4 from foot & \textit{for} Lemma~3289 of Volume~X \textit{read} Lemma~3286 of Volume~X; the result cited is on p.~3910 \\
\dag\,2632 & proof of Theorem~2.5, l.~2 & \textit{for} $\nec\neg\exists z\,\mathrm{MG}(z) \to \neg\exists z\,\mathrm{MG}(z)$ \textit{read} $\pos\neg\exists z\,\mathrm{MG}(z) \to \neg\exists z\,\mathrm{MG}(z)$ \\
2633 & l.~5 & \textit{for} p.~2635 \textit{read} p.~2625 (Lemma~2.3) \\
2633 & l.~6 & \textit{for} Lemma~3140 of Volume~X \textit{read} Lemma~3138 of Volume~X \\
\dag\,2636 & proof of Lemma~2.9, l.~4 & \textit{for} Lemma~3237 of Volume~X \textit{read} Lemma~3242 of Volume~X; the result cited is on p.~3904 \\
\dag\,2636 & proof of Lemma~2.9, l.~7 & \textit{for} Theorem~2.5 \textit{read} Theorem~2.4; the step uses Theorem~2.4 (p.~2627) \\
2636 & l.~8 & \textit{delete} the second occurrence of \emph{is} \\
\dag\,2637 & proof of Lemma~2.9, l.~14 & \textit{for} Lemma~3218 of Volume~X \textit{read} Lemma~3222 of Volume~X; the result cited is on p.~3901 \\
2638 & l.~7 & \textit{for} debunks the possibility premise \textit{read} does not bear on the possibility premise \\
2638 & l.~5 from foot & \textit{for} p.~2618 \textit{read} p.~2617 (Definition~2.1) \\
2639 & l.~4 from foot & \textit{for} cosmolgy \textit{read} cosmology \\
2641 & l.~9 & \textit{for} p.~2625 \textit{read} p.~2635 (Definition~2.8) \\
2641 & l.~13 & \textit{delete} the second occurrence of \emph{the} \\
2641 & l.~1 from foot & \textit{for} Lemma~3224 of Volume~X \textit{read} Lemma~3226 of Volume~X \\
2642 & l.~4 & \textit{for} physcs \textit{read} physics \\
2642 & l.~7 & \textit{for} a normalizable measure over the constants \textit{read} no normalizable measure over the constants \\
2642 & l.~24 & \textit{for} p.~2626 \textit{read} p.~2636 (Lemma~2.9) \\
\dag\,2643 & proof of Proposition~2.10, l.~5 & \textit{for} Definition~2.6 \textit{read} Definition~2.8; the step uses Definition~2.8 (p.~2635) \\
\dag\,2643 & proof of Proposition~2.10, l.~26 & \textit{for} Lemma~2.11 \textit{read} Lemma~2.9 \\
\dag\,2643 & proof of Proposition~2.10, l.~29 & \textit{for} Definition~2.1 \textit{read} Definition~2.6 \\
2644 & l.~7 & \textit{for} p.~2731 \textit{read} p.~2631 (Theorem~2.5) \\
2644 & l.~12 & \textit{for} p.~2743 \textit{read} p.~2643 (Proposition~2.10) \\
2644 & l.~16 & \textit{for} a posteriori premise \textit{read} a priori premise \\
\dag\,2644 & proof of Proposition~2.10, (2.18) & \textit{for} $\pos\neg\exists x\,\mathrm{MG}(x) \to \neg\forall x\,\mathrm{MG}(x)$ \textit{read} $\pos\neg\exists x\,\mathrm{MG}(x) \to \neg\exists x\,\mathrm{MG}(x)$; the formula is used in its correct form on p.~2647 \\
\dag\,2644 & proof of Proposition~2.10, l.~29 & \textit{for} Definition~2.7 \textit{read} Definition~2.2; the step uses Definition~2.2 (p.~2623) \\
\dag\,2644 & proof of Proposition~2.10, l.~35 & \textit{for} $\exists y\,\mathrm{MG}(y,w_2) \to \forall y\,(E(y,w_2) \wedge \mathrm{ME}(y,w_2))$ \textit{read} $\exists y\,\mathrm{MG}(y,w_2) \to \exists y\,(E(y,w_2) \wedge \mathrm{ME}(y,w_2))$ \\
2645 & l.~12 & \textit{for} divne \textit{read} divine \\
2646 & display (2.19) & \textit{for} $\mathrm{NME}(x,w_1) \iff \forall u\,(w_2Ru \to \mathrm{ME}(x,u))$ \textit{read} $\mathrm{NME}(x,w_2) \iff \forall u\,(w_2Ru \to \mathrm{ME}(x,u))$ \\
\dag\,2647 & proof of Lemma~2.11, l.~21 & \textit{for} Definition~2.7 \textit{read} Definition~2.8; the step uses Definition~2.8 (p.~2635) \\
2647 & l.~24 & \textit{after} maximally excellent \textit{insert} at $u$ \\
2647 & l.~27 & \textit{for} p.~2633 \textit{read} p.~2632 (Definition~2.6) \\
2647 & l.~2 from foot & \textit{after} exists \textit{insert} at every world accessible from $w_1$ \\
2648 & l.~31 & \textit{for} particluar \textit{read} particular \\
\dag\,2648 & proof of Proposition~2.12, (2.22) & \textit{for} $\mathrm{ME}(y,u) \iff \mathrm{Op}(y,u) \wedge \mathrm{Os}(y,u) \vee \mathrm{MP}(y,u)$ \textit{read} $\mathrm{ME}(y,u) \iff \mathrm{Op}(y,u) \wedge \mathrm{Os}(y,u) \wedge \mathrm{MP}(y,u)$ \\
2649 & l.~2 & \textit{for} p.~2638 \textit{read} p.~2648 (Proposition~2.12) \\
2650 & l.~11 & \textit{for} p.~2749 \textit{read} p.~2649 (Definition~2.13) \\
2650 & l.~13 & \textit{for} divnie \textit{read} divine \\
\end{errlist}

\subsection*{Chapter~3. Cosmology and the beginning of the universe \hfill \normalfont\small pp.~2651--2698}
\begin{errlist}
2652 & running heads, pp.~2652--2698 (even) & \textit{for} Naturalism and the modal argument \textit{read} Cosmology \\
2653 & l.~14 & \textit{for} contmporary \textit{read} contemporary \\
2653 & l.~19 & \textit{delete} the second occurrence of \emph{to} \\
2653 & l.~29 & \textit{for} Lemma~3208 of Volume~X \textit{read} Lemma~3209 of Volume~X \\
2655 & l.~9 & \textit{for} mutliverse \textit{read} multiverse \\
2655 & l.~30 & \textit{for} p.~2625 \textit{read} p.~2635 (Definition~2.7) \\
2656 & l.~7 & \textit{for} Craig (1980) \textit{read} Craig (1979) \\
2656 & l.~19 & \textit{for} detreminate \textit{read} determinate \\
2657 & display (3.2) & \textit{for} $ D(y,y)$ \textit{read} $\neg D(y,y)$ \\
2657 & l.~21 & \textit{for} Lemma~3282 of Volume~X \textit{read} Lemma~3284 of Volume~X \\
2658 & l.~8 & \textit{for} achivee \textit{read} achieve \\
2658 & l.~18 & \textit{for} dbeunking \textit{read} debunking \\
2658 & l.~26 & \textit{for} Lemma~3114 of Volume~X \textit{read} Lemma~3112 of Volume~X \\
2659 & display (3.3) & \textit{for} $\forall y\,(\neg\mathrm{NE}(x) \to \exists x\,D(y,x))$ \textit{read} $\forall y\,(\neg\mathrm{NE}(y) \to \exists x\,D(y,x))$; correctly printed in (3.5), p.~2661 \\
2660 & l.~16 & \textit{for} eventulaly \textit{read} eventually \\
2660 & l.~29 & \textit{for} anlaysis \textit{read} analysis \\
2660 & l.~1 from foot & \textit{for} p.~2657 \textit{read} p.~2647 (Lemma~2.11) \\
2661 & l.~4 & \textit{for} Craig (1979) \textit{read} Koons (1997) (the cosmological argument from contingency) \\
2661 & l.~5 & \textit{after} dependee \textit{insert} in every accessible world \\
2662 & l.~15 & \textit{for} Kons (1997) \textit{read} Koons (1997) \\
2662 & l.~17 & \textit{for} distniguish \textit{read} distinguish \\
2663 & l.~9 & \textit{for} gennerally \textit{read} generally \\
2665 & l.~34 & \textit{for} Craig (1979) \textit{read} Gale and Pruss (1999) (a weak principle of sufficient reason) \\
\dag\,2666 & proof of Corollary~3.3, l.~20 & \textit{for} Corollary~3.7 \textit{read} Corollary~3.2 \\
2667 & l.~4 from foot & \textit{for} p.~2620 \textit{read} p.~2610 (Definition~1.2) \\
2668 & l.~21 & \textit{for} Craig (1978) \textit{read} Craig (1979) \\
\dag\,2669 & proof of Lemma~3.4, l.~8 & \textit{for} Lemma~3162 of Volume~X \textit{read} Lemma~3156 of Volume~X \\
\dag\,2670 & proof of Lemma~3.4, l.~8 & \textit{for} Corollary~3.2 \textit{read} Corollary~3.3 \\
2671 & l.~13 & \textit{for} cognitvie \textit{read} cognitive \\
2671 & display (3.7) & \textit{for} $P(E \mid T)/P(E \mid  T) < 1$ \textit{read} $P(E \mid T)/P(E \mid \neg T) < 1$ \\
2671 & n.~11 & \textit{for} Gale and Pruss (1998) \textit{read} Gale and Pruss (1999) \\
2673 & l.~14 & \textit{for} Lemma~3134 of Volume~X \textit{read} Lemma~3139 of Volume~X; the result cited is on p.~3889 \\
2673 & l.~32 & \textit{for} Lemma~3155 of Volume~X \textit{read} Lemma~3149 of Volume~X \\
\dag\,2674 & proof of Corollary~3.5, l.~11 & \textit{for} Corollary~3.7 \textit{read} Corollary~3.3 \\
2674 & l.~28 & \textit{for} Lemma~3070 of Volume~X \textit{read} Lemma~3067 of Volume~X \\
\dag\,2674 & proof of Proposition~3.6, (3.10) & \textit{for} $\mathrm{NE}(y) \wedge \exists x\,D(y,x) \to \mathrm{ME}(y)$ \textit{read} $\mathrm{NE}(y) \wedge \neg\exists x\,D(y,x) \to \mathrm{ME}(y)$; the formula is used in its correct form on p.~2678 \\
2675 & l.~19 & \textit{for} primmarily \textit{read} primarily \\
\dag\,2675 & proof of Corollary~3.7, (3.11) & \textit{for} $\forall z\,(\neg\mathrm{NE}(z) \leftrightarrow \exists x\,D(z,x))$ \textit{read} $\forall z\,(\neg\mathrm{NE}(z) \to \exists x\,D(z,x))$ \\
2676 & l.~1 & \textit{for} ordinnary \textit{read} ordinary \\
\dag\,2676 & proof of Corollary~3.7, l.~5 & \textit{for} Proposition~3.8 \textit{read} Proposition~3.6; the step uses Proposition~3.6 (p.~2674) \\
2677 & l.~20 & \textit{for} Konos (1997) \textit{read} Koons (1997) \\
2679 & l.~29 & \textit{for} Lemma~3171 of Volume~X \textit{read} Lemma~3177 of Volume~X \\
2680 & l.~6 & \textit{for} Proposition~2.5 of Volume~IV \textit{read} Proposition~2.2 of Volume~IV \\
\dag\,2680 & proof of Lemma~3.9, l.~12 & \textit{for} Lemma~3095 of Volume~X \textit{read} Lemma~3100 of Volume~X \\
\dag\,2681 & proof of Lemma~3.9, l.~25 & \textit{for} Corollary~3.5 \textit{read} Corollary~3.3; the step uses Corollary~3.3 (p.~2666) \\
\dag\,2682 & proof of Lemma~3.10, l.~9 & \textit{for} Corollary~3.2 \textit{read} Corollary~3.5 \\
\dag\,2683 & proof of Corollary~3.12, l.~3 & \textit{for} Lemma~3.10 \textit{read} Lemma~3.9 \\
2683 & l.~5 & \textit{for} Lemma~3049 of Volume~X \textit{read} Lemma~3052 of Volume~X; the result cited is on p.~3876 \\
\dag\,2684 & proof of Corollary~3.12, l.~17 & \textit{for} Lemma~3.9 \textit{read} Lemma~3.4; the step uses Lemma~3.4 (p.~2669) \\
2684 & l.~18 & \textit{for} p.~2682 \textit{read} p.~2681 (Lemma~3.10) \\
2684 & l.~3 from foot & \textit{for} Lemma~3239 of Volume~X \textit{read} Lemma~3242 of Volume~X \\
2684 & l.~2 from foot & \textit{for} phnomenon \textit{read} phenomenon \\
2685 & l.~7 & \textit{for} a contingent dependee \textit{read} a necessary dependee \\
2686 & l.~11 & \textit{for} Definition~6.7 of Volume~IV \textit{read} Definition~6.1 of Volume~IV; the result cited is on p.~1587 \\
2687 & l.~3 & \textit{for} $ D(x,x)$ \textit{read} $\neg D(x,x)$ \\
2688 & l.~3 & \textit{for} fine-tuinng \textit{read} fine-tuning \\
2688 & l.~20 & \textit{for} p.~2783 \textit{read} p.~2683 (Corollary~3.12) \\
\dag\,2688 & proof of Lemma~3.14, l.~5 from foot & \textit{for} Lemma~3262 of Volume~X \textit{read} Lemma~3259 of Volume~X \\
2689 & l.~26 & \textit{for} anecstral \textit{read} ancestral \\
2690 & l.~7 & \textit{for} Lemma~3181 of Volume~X \textit{read} Lemma~3184 of Volume~X; the result cited is on p.~3895 \\
2691 & l.~3 & \textit{for} $\mathrm{NE}(x) \vee \neg\exists z\,D(x,z) \to \mathrm{ME}(x)$ \textit{read} $\mathrm{NE}(x) \wedge \neg\exists z\,D(x,z) \to \mathrm{ME}(x)$ \\
2691 & l.~10 & \textit{for} Gale and Pruss (1999) \textit{read} Koons (1997) (the cosmological argument from contingency) \\
2691 & l.~17 & \textit{for} Gale and Pruss (2000) \textit{read} Gale and Pruss (1999) \\
2691 & l.~18 & \textit{delete} the second occurrence of \emph{the} \\
\dag\,2692 & proof of Corollary~3.15, l.~5 & \textit{for} Lemma~3296 of Volume~X \textit{read} Lemma~3297 of Volume~X \\
\dag\,2692 & proof of Corollary~3.15, l.~14 & \textit{for} Lemma~3160 of Volume~X \textit{read} Lemma~3166 of Volume~X \\
2692 & l.~19 & \textit{for} p.~2766 \textit{read} p.~2666 (Corollary~3.3) \\
2692 & l.~7 from foot & \textit{for} Lemma~3.2 of Volume~I \textit{read} Lemma~3.3 of Volume~I \\
2692 & l.~40 & \textit{for} commitemnt \textit{read} commitment \\
2693 & l.~8 & \textit{for} $D(y,x) \wedge D(x,z) \leftrightarrow D(y,z)$ \textit{read} $D(y,x) \wedge D(x,z) \to D(y,z)$ \\
\dag\,2693 & proof of Corollary~3.15, l.~15 & \textit{for} Lemma~3.14 \textit{read} Lemma~3.9 \\
2693 & l.~23 & \textit{for} immedaitely \textit{read} immediately \\
2693 & l.~4 from foot & \textit{for} p.~2693 \textit{read} p.~2692 (Corollary~3.15) \\
2694 & l.~6 & \textit{for} additinoal \textit{read} additional \\
2694 & display (3.25) & \textit{for} $D(x,x) \wedge D(x,z) \to D(y,z)$ \textit{read} $D(y,x) \wedge D(x,z) \to D(y,z)$; the formula is used in its correct form on p.~2696 \\
2697 & l.~19 & \textit{for} p.~2686 \textit{read} p.~2687 (Definition~3.13) \\
\dag\,2697 & proof of Lemma~3.16, l.~29 & \textit{for} Proposition~3.6 \textit{read} Proposition~3.11; the step uses Proposition~3.11 (p.~2682) \\
2698 & l.~9 & \textit{delete} the second occurrence of \emph{the} \\
2698 & l.~6 from foot & \textit{for} p.~2696 \textit{read} p.~2697 (Lemma~3.16) \\
\end{errlist}

\subsection*{Chapter~4. Fine-tuning and the multiverse \hfill \normalfont\small pp.~2699--2747}
\begin{errlist}
2700 & l.~19 & \textit{for} Lemma~2.4 of Volume~IV \textit{read} Lemma~2.8 of Volume~IV; the result cited is on p.~1427 \\
2702 & l.~7 & \textit{for} Lemma~3261 of Volume~X \textit{read} Lemma~3267 of Volume~X \\
2702 & l.~12 & \textit{for} Lemma~3072 of Volume~X \textit{read} Lemma~3078 of Volume~X \\
2702 & l.~13 & \textit{for} Lemma~3252 of Volume~X \textit{read} Lemma~3251 of Volume~X; the result cited is on p.~3905 \\
\dag\,2704 & proof of Corollary~4.3, l.~2 & \textit{for} Lemma~3136 of Volume~X \textit{read} Lemma~3131 of Volume~X; the result cited is on p.~3888 \\
2705 & l.~5 & \textit{for} Lemma~3147 of Volume~X \textit{read} Lemma~3146 of Volume~X \\
2706 & l.~5 & \textit{for} cgonitive \textit{read} cognitive \\
2706 & l.~23 & \textit{for} illustrte \textit{read} illustrate \\
2706 & l.~39 & \textit{for} p.~2700 \textit{read} p.~2701 (Definition~4.2) \\
2707 & l.~15 & \textit{for} genrally \textit{read} generally \\
2708 & l.~13 & \textit{for} Leslie (1989) \textit{read} Kraay (2010) (the theistic multiverse) \\
2708 & l.~3 from foot & \textit{for} p.~2777 \textit{read} p.~2677 (Proposition~3.8) \\
\dag\,2709 & proof of Theorem~4.4, l.~7 from foot & \textit{for} Lemma~3208 of Volume~X \textit{read} Lemma~3206 of Volume~X; the result cited is on p.~3898 \\
\dag\,2710 & proof of Theorem~4.4, l.~17 & \textit{for} Definition~4.2 \textit{read} Definition~4.1; the step uses Definition~4.1 (p.~2700) \\
\dag\,2710 & proof of Theorem~4.4, l.~23 & \textit{for} Definition~4.7 \textit{read} Definition~4.2 \\
\dag\,2710 & proof of Theorem~4.4, l.~4 from foot & \textit{for} Lemma~3162 of Volume~X \textit{read} Lemma~3160 of Volume~X \\
2711 & l.~5 & \textit{for} cosncious \textit{read} conscious \\
2712 & l.~14 & \textit{for} p.~2860 \textit{read} p.~2680 (Lemma~3.9) \\
2712 & l.~28 & \textit{for} nomralizable \textit{read} normalizable \\
2713 & l.~12 & \textit{for} Proposition~5.7 of Volume~III \textit{read} Proposition~5.8 of Volume~III \\
2714 & l.~8 from foot & \textit{for} p.~2685 \textit{read} p.~2675 (Corollary~3.7) \\
\dag\,2716 & proof of Proposition~4.6, l.~2 & \textit{for} Definition~4.2 \textit{read} Definition~4.5 \\
2717 & l.~18 & \textit{for} Karay (2010) \textit{read} Kraay (2010) \\
2718 & l.~18 & \textit{for} Swinburne (2004) \textit{read} Kraay (2010) (the theistic multiverse) \\
2719 & l.~6 & \textit{for} p.~2678 \textit{read} p.~2688 (Lemma~3.14) \\
2722 & l.~4 & \textit{for} Lemma~3077 of Volume~X \textit{read} Lemma~3080 of Volume~X \\
2722 & display (4.6) & \textit{for} $P(E \mid T)/P(E \mid \neg T) > 1$ \textit{read} $P(E \mid T)/P(E \mid \neg T) < 1$; correctly printed in (4.7), p.~2728 \\
2722 & l.~25 & \textit{for} Proposition~2.2 of Volume~IV \textit{read} Proposition~2.5 of Volume~IV; the result cited is on p.~1425 \\
2723 & l.~19 & \textit{for} cojnecture \textit{read} conjecture \\
2723 & l.~23 & \textit{for} p.~2675 \textit{read} p.~2674 (Proposition~3.6) \\
2723 & l.~2 from foot & \textit{for} p.~2684 \textit{read} p.~2674 (Corollary~3.5) \\
\dag\,2724 & proof of Proposition~4.8, l.~16 & \textit{for} Proposition~4.11 \textit{read} Proposition~4.6; the step uses Proposition~4.6 (p.~2716) \\
\dag\,2724 & proof of Proposition~4.8, l.~22 & \textit{for} Lemma~3265 of Volume~X \textit{read} Lemma~3263 of Volume~X \\
2725 & l.~18 & \textit{delete} the second occurrence of \emph{is} \\
\dag\,2725 & proof of Proposition~4.8, l.~26 & \textit{for} Lemma~3238 of Volume~X \textit{read} Lemma~3244 of Volume~X \\
2725 & l.~30 & \textit{delete} the second occurrence of \emph{is} \\
2726 & l.~16 & \textit{delete} the second occurrence of \emph{is} \\
2726 & l.~8 from foot & \textit{for} p.~2701 \textit{read} p.~2710 (Definition~4.5) \\
2727 & l.~31 & \textit{for} descrbe \textit{read} describe \\
2728 & l.~5 from foot & \textit{for} Lemma~3110 of Volume~X \textit{read} Lemma~3105 of Volume~X; the result cited is on p.~3884 \\
2728 & l.~37 & \textit{for} p.~2809 \textit{read} p.~2709 (Theorem~4.4) \\
2729 & l.~18 & \textit{for} Lemma~3200 of Volume~X \textit{read} Lemma~3196 of Volume~X \\
2729 & l.~23 & \textit{for} Lemma~3165 of Volume~X \textit{read} Lemma~3160 of Volume~X \\
2730 & l.~6 from foot & \textit{delete} the second occurrence of \emph{the} \\
2730 & l.~1 from foot & \textit{for} Lemma~3128 of Volume~X \textit{read} Lemma~3125 of Volume~X; the result cited is on p.~3887 \\
2731 & l.~21 & \textit{for} cognitve \textit{read} cognitive \\
2732 & l.~23 & \textit{delete} the second occurrence of \emph{that} \\
2732 & l.~25 & \textit{for} Lemma~3161 of Volume~X \textit{read} Lemma~3160 of Volume~X \\
2733 & l.~17 & \textit{for} Lselie (1989) \textit{read} Leslie (1989) \\
2733 & l.~21 & \textit{for} netiher \textit{read} neither \\
2734 & l.~21 & \textit{for} Lemma~3068 of Volume~X \textit{read} Lemma~3071 of Volume~X \\
2734 & l.~38 & \textit{for} naturlaism \textit{read} naturalism \\
\dag\,2736 & proof of Corollary~4.10, l.~28 & \textit{for} Proposition~4.8 \textit{read} Proposition~4.6; the step uses Proposition~4.6 (p.~2716) \\
2736 & l.~1 from foot & \textit{for} Lemma~6.6 of Volume~IV \textit{read} Lemma~6.4 of Volume~IV \\
2737 & l.~6 & \textit{for} Lemma~3060 of Volume~X \textit{read} Lemma~3055 of Volume~X \\
2737 & l.~10 & \textit{for} p.~2727 \textit{read} p.~2627 (Theorem~2.4) \\
2737 & l.~21 & \textit{for} Lemma~3117 of Volume~X \textit{read} Lemma~3113 of Volume~X \\
2738 & l.~20 & \textit{for} Leslei (1989) \textit{read} Leslie (1989) \\
2739 & l.~18 & \textit{for} Lemma~3182 of Volume~X \textit{read} Lemma~3188 of Volume~X \\
2739 & l.~30 & \textit{for} Lemma~3217 of Volume~X \textit{read} Lemma~3216 of Volume~X \\
2739 & l.~35 & \textit{delete} the second occurrence of \emph{to} \\
\dag\,2740 & proof of Proposition~4.11, l.~6 & \textit{for} Definition~4.1 \textit{read} Definition~4.5; the step uses Definition~4.5 (p.~2710) \\
\dag\,2740 & proof of Proposition~4.11, l.~13 & \textit{for} Definition~4.5 \textit{read} Definition~4.7; the step uses Definition~4.7 (p.~2717) \\
\dag\,2741 & proof of Proposition~4.12, l.~30 & \textit{for} Lemma~3075 of Volume~X \textit{read} Lemma~3076 of Volume~X; the result cited is on p.~3880 \\
2741 & l.~5 from foot & \textit{for} p.~2741 \textit{read} p.~2740 (Proposition~4.11) \\
2741 & l.~39 & \textit{for} fine-utning \textit{read} fine-tuning \\
\dag\,2742 & proof of Proposition~4.12, l.~6 & \textit{for} Proposition~4.11 \textit{read} Proposition~4.8 \\
\dag\,2742 & proof of Proposition~4.12, l.~22 & \textit{for} Lemma~3068 of Volume~X \textit{read} Lemma~3067 of Volume~X; the result cited is on p.~3879 \\
2742 & l.~27 & \textit{delete} the second occurrence of \emph{of} \\
2743 & l.~7 from foot & \textit{for} Lemma~3080 of Volume~X \textit{read} Lemma~3084 of Volume~X \\
2745 & l.~31 & \textit{for} Lemma~2.8 of Volume~IV \textit{read} Lemma~2.12 of Volume~IV \\
2746 & l.~4 & \textit{for} p.~2845 \textit{read} p.~2745 (Theorem~4.15) \\
\dag\,2746 & proof of Proposition~4.16, l.~15 & \textit{for} Proposition~4.12 \textit{read} Proposition~4.11 \\
2746 & l.~16 & \textit{for} pyhsics \textit{read} physics \\
\dag\,2746 & proof of Proposition~4.16, l.~20 & \textit{for} Definition~4.14 \textit{read} Definition~4.13; the step uses Definition~4.13 (p.~2742) \\
\dag\,2746 & proof of Proposition~4.16, l.~28 & \textit{for} Proposition~4.8 \textit{read} Proposition~4.12 \\
2747 & l.~27 & \textit{for} derivtion \textit{read} derivation \\
2747 & l.~38 & \textit{for} p.~2736 \textit{read} p.~2746 (Proposition~4.16) \\
\dag\,2747 & proof of Proposition~4.16, l.~2 from foot & \textit{for} Lemma~3174 of Volume~X \textit{read} Lemma~3175 of Volume~X \\
\end{errlist}

\subsection*{Chapter~5. Evolutionary debunking of theistic belief \hfill \normalfont\small pp.~2748--2796}
\begin{errlist}
2749 & l.~4 & \textit{for} Lemma~3212 of Volume~X \textit{read} Lemma~3215 of Volume~X \\
2749 & l.~30 & \textit{for} Lemma~3263 of Volume~X \textit{read} Lemma~3258 of Volume~X \\
2750 & n.~6 & \textit{for} Street (2005) \textit{read} Street (2006) \\
2751 & l.~30 & \textit{for} Lemma~3085 of Volume~X \textit{read} Lemma~3080 of Volume~X; the result cited is on p.~3880 \\
2752 & l.~29 & \textit{for} p.~2726 \textit{read} p.~2716 (Proposition~4.6) \\
2755 & l.~33 & \textit{for} p.~2681 \textit{read} p.~2682 (Proposition~3.11) \\
2755 & l.~2 from foot & \textit{for} Barrett (2005) \textit{read} Barrett (2004) \\
2756 & l.~3 & \textit{for} Plantinga (2011) \textit{read} Street (2006) (the Darwinian dilemma) \\
2756 & l.~5 from foot & \textit{for} Barret (2004) \textit{read} Barrett (2004) \\
2757 & l.~6 from foot & \textit{after} supported \textit{insert} independently of the conclusion \\
2758 & l.~10 & \textit{for} Lemma~3216 of Volume~X \textit{read} Lemma~3212 of Volume~X; the result cited is on p.~3899 \\
2758 & l.~16 & \textit{for} conceivaiblity \textit{read} conceivability \\
2759 & l.~5 from foot & \textit{for} Lemma~2.10 of Volume~IV \textit{read} Lemma~2.12 of Volume~IV \\
2760 & l.~25 & \textit{for} cerdence \textit{read} credence \\
2761 & l.~8 & \textit{for} Lemma~3275 of Volume~X \textit{read} Lemma~3269 of Volume~X \\
2762 & l.~5 & \textit{for} Lemma~3166 of Volume~X \textit{read} Lemma~3172 of Volume~X \\
2762 & l.~19 & \textit{for} Definition~2.1 of Volume~IV \textit{read} Definition~2.3 of Volume~IV \\
2762 & l.~27 & \textit{for} consequecne \textit{read} consequence \\
2763 & l.~33 & \textit{for} Lemma~3155 of Volume~X \textit{read} Lemma~3158 of Volume~X \\
2763 & l.~6 from foot & \textit{for} likeilhood \textit{read} likelihood \\
\dag\,2764 & proof of Corollary~5.5, l.~27 & \textit{for} Lemma~5.6 \textit{read} Lemma~5.4 \\
\dag\,2764 & proof of Corollary~5.5, l.~32 & \textit{for} Lemma~3110 of Volume~X \textit{read} Lemma~3108 of Volume~X \\
2765 & l.~15 & \textit{for} Lemma~3183 of Volume~X \textit{read} Lemma~3184 of Volume~X \\
\dag\,2765 & proof of Corollary~5.5, l.~25 & \textit{for} Proposition~5.7 \textit{read} Proposition~5.3 \\
2765 & l.~32 & \textit{for} Lemma~3075 of Volume~X \textit{read} Lemma~3078 of Volume~X \\
2765 & l.~7 from foot & \textit{for} conceivabliity \textit{read} conceivability \\
2766 & l.~1 & \textit{after} conceivable \textit{insert} on ideal reflection \\
2766 & l.~14 & \textit{for} derivattion \textit{read} derivation \\
\dag\,2766 & proof of Lemma~5.6, l.~16 & \textit{for} Definition~5.1 \textit{read} Definition~5.2 \\
2767 & l.~19 & \textit{for} plausble \textit{read} plausible \\
2768 & l.~11 & \textit{for} p.~2864 \textit{read} p.~2764 (Corollary~5.5) \\
2768 & l.~12 & \textit{for} p.~2731 \textit{read} p.~2741 (Proposition~4.12) \\
2769 & l.~10 & \textit{for} symemtry \textit{read} symmetry \\
2769 & l.~40 & \textit{for} Lemma~3212 of Volume~X \textit{read} Lemma~3210 of Volume~X \\
\dag\,2770 & proof of Proposition~5.7, l.~28 & \textit{for} Lemma~3278 of Volume~X \textit{read} Lemma~3273 of Volume~X; the result cited is on p.~3908 \\
2771 & display (5.6) & \textit{see} Corrigendum~S13 \\
2772 & l.~7 & \textit{for} Lemma~3225 of Volume~X \textit{read} Lemma~3222 of Volume~X \\
2772 & l.~25 & \textit{for} Lemma~3223 of Volume~X \textit{read} Lemma~3229 of Volume~X \\
2774 & l.~29 & \textit{for} p.~2707 \textit{read} p.~2770 (Proposition~5.7) \\
\dag\,2775 & proof of Corollary~5.8, l.~8 & \textit{for} Lemma~5.11 \textit{read} Lemma~5.6 \\
\dag\,2775 & proof of Corollary~5.8, l.~11 & \textit{for} Lemma~3168 of Volume~X \textit{read} Lemma~3169 of Volume~X; the result cited is on p.~3893 \\
2775 & l.~17 & \textit{for} equvialent \textit{read} equivalent \\
2779 & l.~27 & \textit{for} Lemma~3085 of Volume~X \textit{read} Lemma~3084 of Volume~X \\
\dag\,2780 & proof of Theorem~5.9, l.~5 & \textit{for} Corollary~5.8 \textit{read} Corollary~5.5 \\
2780 & l.~27 & \textit{for} p.~2656 \textit{read} p.~2756 (Definition~5.2) \\
2781 & l.~13 & \textit{for} fine-tnuing \textit{read} fine-tuning \\
2782 & l.~14 & \textit{for} Lemma~3094 of Volume~X \textit{read} Lemma~3090 of Volume~X \\
2783 & l.~4 & \textit{for} Lemma~3122 of Volume~X \textit{read} Lemma~3128 of Volume~X \\
2783 & l.~21 & \textit{for} Lemma~3074 of Volume~X \textit{read} Lemma~3075 of Volume~X \\
2784 & l.~2 & \textit{for} Lemma~3267 of Volume~X \textit{read} Lemma~3269 of Volume~X \\
2784 & l.~13 & \textit{for} p.~2865 \textit{read} p.~2765 (Lemma~5.6) \\
2785 & l.~9 & \textit{for} deubnking \textit{read} debunking \\
2786 & l.~10 & \textit{for} Lemma~3220 of Volume~X \textit{read} Lemma~3217 of Volume~X \\
2786 & l.~33 & \textit{for} fien-tuning \textit{read} fine-tuning \\
\dag\,2787 & proof of Lemma~5.11, l.~4 & \textit{for} Lemma~3189 of Volume~X \textit{read} Lemma~3195 of Volume~X; the result cited is on p.~3897 \\
2787 & l.~9 & \textit{for} Lemma~3197 of Volume~X \textit{read} Lemma~3194 of Volume~X; the result cited is on p.~3897 \\
\dag\,2787 & proof of Lemma~5.11, l.~11 & \textit{for} Lemma~3127 of Volume~X \textit{read} Lemma~3125 of Volume~X \\
2787 & l.~29 & \textit{delete} the second occurrence of \emph{that} \\
2788 & n.~12 & \textit{for} Barrett (2003) \textit{read} Barrett (2004) \\
\dag\,2789 & proof of Corollary~5.12, l.~15 & \textit{for} Lemma~5.14 \textit{read} Lemma~5.11; the step uses Lemma~5.11 (p.~2787) \\
2789 & l.~20 & \textit{for} naturalsim \textit{read} naturalism \\
\dag\,2789 & proof of Corollary~5.12, l.~24 & \textit{for} Proposition~5.10 \textit{read} Proposition~5.7; the step uses Proposition~5.7 (p.~2770) \\
2789 & l.~28 & \textit{for} divinne \textit{read} divine \\
\dag\,2789 & proof of Corollary~5.12, l.~40 & \textit{for} Lemma~3283 of Volume~X \textit{read} Lemma~3278 of Volume~X; the result cited is on p.~3909 \\
2790 & l.~28 & \textit{for} defaeter \textit{read} defeater \\
\dag\,2791 & proof of Lemma~5.15, l.~12 & \textit{for} Corollary~5.8 \textit{read} Corollary~5.12; the step uses Corollary~5.12 (p.~2789) \\
\dag\,2791 & proof of Lemma~5.15, l.~15 & \textit{for} Lemma~3137 of Volume~X \textit{read} Lemma~3142 of Volume~X \\
\dag\,2791 & proof of Lemma~5.15, l.~25 & \textit{for} Lemma~3076 of Volume~X \textit{read} Lemma~3075 of Volume~X; the result cited is on p.~3880 \\
2791 & l.~1 from foot & \textit{for} p.~2889 \textit{read} p.~2789 (Corollary~5.12) \\
2793 & l.~4 & \textit{for} Lemma~2.12 of Volume~IV \textit{read} Lemma~2.11 of Volume~IV \\
2793 & l.~6 & \textit{for} Barrett (2004) \textit{read} Street (2006) (the Darwinian dilemma) \\
2793 & l.~11 & \textit{for} p.~2790 \textit{read} p.~2791 (Lemma~5.15) \\
2793 & l.~14 & \textit{for} Lemma~3094 of Volume~X \textit{read} Lemma~3099 of Volume~X \\
2793 & l.~19 & \textit{for} Street (2007) \textit{read} Street (2006) \\
2794 & l.~2 & \textit{for} Lemma~3179 of Volume~X \textit{read} Lemma~3175 of Volume~X \\
2794 & l.~20 & \textit{for} ratoinal \textit{read} rational \\
\dag\,2795 & proof of Proposition~5.16, l.~15 & \textit{for} Lemma~3144 of Volume~X \textit{read} Lemma~3149 of Volume~X \\
2796 & l.~21 & \textit{for} p.~2794 \textit{read} p.~2795 (Proposition~5.16) \\
\dag\,2796 & proof of Proposition~5.16, l.~27 & \textit{for} Lemma~3062 of Volume~X \textit{read} Lemma~3056 of Volume~X \\
\dag\,2796 & proof of Proposition~5.16, l.~8 from foot & \textit{for} Lemma~3107 of Volume~X \textit{read} Lemma~3112 of Volume~X \\
\end{errlist}

\subsection*{Chapter~6. The cognitive science of religion \hfill \normalfont\small pp.~2797--2839}
\begin{errlist}
2798 & l.~17 & \textit{for} Lemma~3169 of Volume~X \textit{read} Lemma~3174 of Volume~X \\
2798 & l.~28 & \textit{for} Lemma~3210 of Volume~X \textit{read} Lemma~3215 of Volume~X \\
2798 & l.~29 & \textit{for} Lemma~3087 of Volume~X \textit{read} Lemma~3084 of Volume~X \\
2799 & l.~20 & \textit{for} Lemma~3145 of Volume~X \textit{read} Lemma~3149 of Volume~X; the result cited is on p.~3890 \\
2801 & l.~26 & \textit{for} Lemma~3216 of Volume~X \textit{read} Lemma~3213 of Volume~X \\
2801 & l.~29 & \textit{for} Lemma~3149 of Volume~X \textit{read} Lemma~3148 of Volume~X \\
2802 & l.~4 & \textit{for} Lemma~3120 of Volume~X \textit{read} Lemma~3123 of Volume~X \\
2803 & l.~3 & \textit{for} Lemma~3221 of Volume~X \textit{read} Lemma~3222 of Volume~X \\
2805 & l.~10 & \textit{for} Lemma~3292 of Volume~X \textit{read} Lemma~3295 of Volume~X; the result cited is on p.~3911 \\
2805 & l.~26 & \textit{for} anticipte \textit{read} anticipate \\
2806 & l.~33 & \textit{for} Proposition~5.8 of Volume~III \textit{read} Proposition~5.2 of Volume~III \\
2807 & l.~7 & \textit{for} p.~2774 \textit{read} p.~2764 (Lemma~5.4) \\
2807 & l.~8 & \textit{for} Lemma~3139 of Volume~X \textit{read} Lemma~3133 of Volume~X; the result cited is on p.~3888 \\
2808 & l.~20 & \textit{for} Lemma~3139 of Volume~X \textit{read} Lemma~3142 of Volume~X \\
2808 & l.~27 & \textit{for} Lemma~3180 of Volume~X \textit{read} Lemma~3182 of Volume~X \\
\dag\,2809 & proof of Theorem~6.3, l.~9 & \textit{for} Lemma~3052 of Volume~X \textit{read} Lemma~3058 of Volume~X \\
\dag\,2809 & proof of Theorem~6.3, l.~14 & \textit{for} Definition~6.6 \textit{read} Definition~6.1 \\
\dag\,2810 & proof of Proposition~6.4, l.~20 & \textit{for} Lemma~3278 of Volume~X \textit{read} Lemma~3277 of Volume~X \\
2813 & l.~6 from foot & \textit{for} p.~2801 \textit{read} p.~2810 (Proposition~6.4) \\
2814 & l.~15 & \textit{for} Definition~6.2 of Volume~IV \textit{read} Definition~6.8 of Volume~IV \\
2815 & l.~2 & \textit{for} Lemma~3126 of Volume~X \textit{read} Lemma~3122 of Volume~X; the result cited is on p.~3886 \\
2815 & l.~23 & \textit{for} Lemma~3261 of Volume~X \textit{read} Lemma~3258 of Volume~X \\
\dag\,2816 & proof of Lemma~6.7, l.~10 & \textit{for} Definition~6.5 \textit{read} Definition~6.1; the step uses Definition~6.1 (p.~2804) \\
\dag\,2816 & proof of Lemma~6.7, l.~15 & \textit{for} Lemma~3101 of Volume~X \textit{read} Lemma~3106 of Volume~X; the result cited is on p.~3884 \\
\dag\,2816 & proof of Lemma~6.7, l.~27 & \textit{for} Definition~6.1 \textit{read} Definition~6.2 \\
\dag\,2817 & proof of Lemma~6.8, l.~24 & \textit{for} Lemma~6.9 \textit{read} Lemma~6.7 \\
\dag\,2818 & proof of Lemma~6.8, l.~7 & \textit{for} Definition~6.5 \textit{read} Definition~6.2; the step uses Definition~6.2 (p.~2804) \\
\dag\,2818 & proof of Lemma~6.8, l.~21 & \textit{for} Lemma~3107 of Volume~X \textit{read} Lemma~3113 of Volume~X \\
2820 & l.~20 & \textit{for} Lemma~3102 of Volume~X \textit{read} Lemma~3107 of Volume~X; the result cited is on p.~3884 \\
2822 & l.~22 & \textit{for} Lemma~3224 of Volume~X \textit{read} Lemma~3220 of Volume~X \\
2822 & l.~3 from foot & \textit{for} Lemma~3080 of Volume~X \textit{read} Lemma~3083 of Volume~X; the result cited is on p.~3881 \\
2823 & l.~9 & \textit{for} consquence \textit{read} consequence \\
2823 & l.~28 & \textit{for} Lemma~3157 of Volume~X \textit{read} Lemma~3163 of Volume~X; the result cited is on p.~3892 \\
\dag\,2824 & proof of Lemma~6.9, l.~7 & \textit{for} Definition~6.6 \textit{read} Definition~6.5 \\
\dag\,2824 & proof of Lemma~6.9, l.~10 & \textit{for} Lemma~6.8 \textit{read} Lemma~6.7; the step uses Lemma~6.7 (p.~2816) \\
2824 & l.~22 & \textit{for} p.~2687 \textit{read} p.~2787 (Lemma~5.11) \\
2824 & l.~26 & \textit{for} Definition~5.18 of Volume~III \textit{read} Definition~5.20 of Volume~III; the result cited is on p.~1010 \\
2826 & l.~21 & \textit{for} p.~2842 \textit{read} p.~2824 (Lemma~6.9) \\
\dag\,2827 & proof of Proposition~6.11, l.~4 & \textit{for} Lemma~6.13 \textit{read} Lemma~6.8 \\
\dag\,2828 & proof of Proposition~6.11, l.~1 & \textit{for} Lemma~6.8 \textit{read} Lemma~6.9 \\
2828 & l.~16 & \textit{for} p.~2817 \textit{read} p.~2827 (Definition~6.10) \\
\dag\,2828 & proof of Proposition~6.11, l.~26 & \textit{for} Definition~6.2 \textit{read} Definition~6.6 \\
\dag\,2828 & proof of Proposition~6.11, l.~37 & \textit{for} Definition~5.15 of Volume~III \textit{read} Definition~5.18 of Volume~III; the result cited is on p.~1004 \\
2828 & l.~3 from foot & \textit{for} p.~2819 \textit{read} p.~2809 (Theorem~6.3) \\
2829 & l.~6 & \textit{for} p.~2732 \textit{read} p.~2742 (Definition~4.13) \\
2829 & l.~30 & \textit{for} Lemma~3134 of Volume~X \textit{read} Lemma~3130 of Volume~X \\
\dag\,2830 & proof of Lemma~6.13, l.~11 & \textit{for} Lemma~3064 of Volume~X \textit{read} Lemma~3069 of Volume~X \\
\dag\,2830 & proof of Lemma~6.13, l.~23 & \textit{for} Lemma~6.6 of Volume~IV \textit{read} Lemma~6.12 of Volume~IV \\
\dag\,2830 & proof of Lemma~6.13, l.~24 & \textit{for} Definition~6.6 \textit{read} Definition~6.10; the step uses Definition~6.10 (p.~2827) \\
2830 & l.~25 & \textit{for} Lemma~3099 of Volume~X \textit{read} Lemma~3100 of Volume~X \\
2831 & l.~6 & \textit{for} Lemma~3046 of Volume~X \textit{read} Lemma~3049 of Volume~X; the result cited is on p.~3876 \\
2831 & l.~22 & \textit{for} p.~2817 \textit{read} p.~2816 (Lemma~6.7) \\
2831 & l.~24 & \textit{for} p.~2804 \textit{read} p.~2814 (Definition~6.6) \\
2831 & l.~2 from foot & \textit{for} conveniton \textit{read} convention \\
2832 & l.~16 & \textit{for} p.~2803 \textit{read} p.~2830 (Lemma~6.13) \\
2833 & l.~5 & \textit{for} Lemma~3049 of Volume~X \textit{read} Lemma~3046 of Volume~X \\
2833 & l.~23 & \textit{for} p.~2665 \textit{read} p.~2656 (Definition~3.1) \\
2833 & l.~27 & \textit{for} phsyics \textit{read} physics \\
2834 & l.~3 & \textit{for} p.~2624 \textit{read} p.~2724 (Proposition~4.8) \\
\dag\,2834 & proof of Lemma~6.14, l.~11 & \textit{for} Definition~6.5 \textit{read} Definition~6.10 \\
\dag\,2834 & proof of Lemma~6.14, l.~32 & \textit{for} Lemma~3102 of Volume~X \textit{read} Lemma~3104 of Volume~X \\
\dag\,2835 & proof of Lemma~6.14, l.~5 & \textit{for} Lemma~3201 of Volume~X \textit{read} Lemma~3203 of Volume~X \\
2835 & l.~36 & \textit{for} cnosequence \textit{read} consequence \\
2835 & l.~40 & \textit{for} p.~2843 \textit{read} p.~2834 (Lemma~6.14) \\
2837 & l.~2 & \textit{for} p.~2811 \textit{read} p.~2810 (Definition~6.5) \\
2837 & l.~8 from foot & \textit{for} Lemma~3061 of Volume~X \textit{read} Lemma~3064 of Volume~X \\
2838 & l.~7 & \textit{for} Lemma~5.21 of Volume~III \textit{read} Lemma~5.17 of Volume~III \\
2838 & l.~22 & \textit{for} p.~2763 \textit{read} p.~2663 (Corollary~3.2) \\
2839 & l.~15 & \textit{for} likelhiood \textit{read} likelihood \\
2839 & l.~32 & \textit{for} Lemma~3269 of Volume~X \textit{read} Lemma~3275 of Volume~X \\
\end{errlist}

\subsection*{Chapter~7. Probability, evidence and a priori premises \hfill \normalfont\small pp.~2840--2896}
\begin{errlist}
2841 & l.~21 & \textit{for} Lemma~3069 of Volume~X \textit{read} Lemma~3063 of Volume~X; the result cited is on p.~3878 \\
2841 & l.~30 & \textit{for} Proposition~5.8 of Volume~III \textit{read} Proposition~5.11 of Volume~III \\
2842 & running heads, pp.~2842--2896 (even) & \textit{for} The cognitive science of religion \textit{read} Probability and a priori premises \\
2842 & l.~10 & \textit{for} Lemma~3096 of Volume~X \textit{read} Lemma~3098 of Volume~X \\
2843 & l.~26 & \textit{for} Definition~5.13 of Volume~III \textit{read} Definition~5.14 of Volume~III \\
2843 & l.~4 from foot & \textit{for} Lemma~3061 of Volume~X \textit{read} Lemma~3057 of Volume~X; the result cited is on p.~3877 \\
2844 & l.~7 & \textit{for} Lemma~3051 of Volume~X \textit{read} Lemma~3045 of Volume~X; the result cited is on p.~3876 \\
2844 & l.~25 & \textit{for} $P(\pos G) > P(\nec\neg G)$ \textit{read} $P(\pos G) > P(\pos\neg G)$ \\
2845 & l.~11 & \textit{for} $P(q \mid r) = P(q \mid q)\,P(q)/P(r)$ \textit{read} $P(q \mid r) = P(r \mid q)\,P(q)/P(r)$ \\
2846 & display (7.2) & \textit{for} $\neg\mathrm{Conc}(\neg r) \vee \mathrm{Conc}(r)$ \textit{read} $\neg\mathrm{Conc}(\neg r) \wedge \mathrm{Conc}(r)$; correctly printed in (7.3), p.~2850 \\
2847 & l.~31 & \textit{for} Lemma~3103 of Volume~X \textit{read} Lemma~3107 of Volume~X \\
\dag\,2848 & proof of Corollary~7.4, l.~14 & \textit{for} Definition~7.5 \textit{read} Definition~7.1 \\
2848 & l.~38 & \textit{for} p.~2856 \textit{read} p.~2756 (Definition~5.1) \\
2849 & l.~36 & \textit{for} Lemma~3252 of Volume~X \textit{read} Lemma~3247 of Volume~X \\
2850 & l.~15 & \textit{for} Lemma~3219 of Volume~X \textit{read} Lemma~3215 of Volume~X \\
2850 & display (7.4) & \textit{for} $P(\pos G \mid A) = P(\pos\neg G)$ \textit{read} $P(\pos\neg G \mid A) = P(\pos\neg G)$; correctly printed in (7.5), p.~2850 \\
2851 & l.~32 & \textit{for} Lemma~3292 of Volume~X \textit{read} Lemma~3289 of Volume~X \\
2854 & l.~23 & \textit{for} Lemma~3247 of Volume~X \textit{read} Lemma~3241 of Volume~X; the result cited is on p.~3903 \\
2855 & l.~28 & \textit{for} likelihod \textit{read} likelihood \\
2856 & l.~29 & \textit{for} p.~2791 \textit{read} p.~2790 (Lemma~5.14) \\
2857 & l.~10 & \textit{for} Lemma~3053 of Volume~X \textit{read} Lemma~3052 of Volume~X \\
2857 & l.~12 & \textit{for} Lemma~3261 of Volume~X \textit{read} Lemma~3262 of Volume~X \\
2858 & l.~1 & \textit{for} p.~2622 \textit{read} p.~2623 (Definition~2.2) \\
2859 & l.~31 & \textit{for} Lemma~3187 of Volume~X \textit{read} Lemma~3185 of Volume~X \\
2860 & l.~1 & \textit{for} p.~2837 \textit{read} p.~2847 (Definition~7.3) \\
2860 & l.~4 from foot & \textit{for} Lemma~3173 of Volume~X \textit{read} Lemma~3168 of Volume~X \\
2861 & l.~20 & \textit{for} Lemma~3182 of Volume~X \textit{read} Lemma~3176 of Volume~X; the result cited is on p.~3894 \\
\dag\,2862 & proof of Lemma~7.6, l.~12 & \textit{for} Definition~7.3 \textit{read} Definition~7.5 \\
\dag\,2863 & proof of Corollary~7.7, l.~7 & \textit{for} Lemma~3235 of Volume~X \textit{read} Lemma~3238 of Volume~X \\
2863 & l.~15 & \textit{for} normalziable \textit{read} normalizable \\
2863 & l.~25 & \textit{for} Lemma~3227 of Volume~X \textit{read} Lemma~3223 of Volume~X \\
2864 & l.~3 & \textit{for} Lemma~3147 of Volume~X \textit{read} Lemma~3153 of Volume~X \\
\dag\,2864 & proof of Proposition~7.8, l.~22 & \textit{for} Lemma~3144 of Volume~X \textit{read} Lemma~3148 of Volume~X \\
\dag\,2864 & proof of Proposition~7.8, l.~30 & \textit{for} Definition~7.1 \textit{read} Definition~7.3 \\
\dag\,2865 & proof of Proposition~7.8, l.~12 & \textit{for} Lemma~3074 of Volume~X \textit{read} Lemma~3077 of Volume~X \\
2865 & l.~4 from foot & \textit{for} Definition~2.3 of Volume~I \textit{read} Definition~2.2 of Volume~I \\
\dag\,2866 & proof of Lemma~7.9, l.~3 & \textit{for} Lemma~3077 of Volume~X \textit{read} Lemma~3071 of Volume~X \\
2866 & l.~11 & \textit{for} Lemma~3162 of Volume~X \textit{read} Lemma~3164 of Volume~X; the result cited is on p.~3892 \\
\dag\,2866 & proof of Lemma~7.9, l.~27 & \textit{for} Proposition~7.11 \textit{read} Proposition~7.8; the step uses Proposition~7.8 (p.~2864) \\
2866 & l.~37 & \textit{for} naturalsm \textit{read} naturalism \\
2869 & l.~34 & \textit{for} Lemma~3202 of Volume~X \textit{read} Lemma~3200 of Volume~X \\
2870 & l.~2 from foot & \textit{for} multievrse \textit{read} multiverse \\
\dag\,2871 & proof of Proposition~7.12, l.~10 & \textit{for} Lemma~7.13 \textit{read} Lemma~7.9; the step uses Lemma~7.9 (p.~2866) \\
\dag\,2872 & proof of Lemma~7.13, l.~5 & \textit{for} Lemma~7.14 \textit{read} Lemma~7.9; the step uses Lemma~7.9 (p.~2866) \\
2872 & l.~19 & \textit{for} Lemma~3049 of Volume~X \textit{read} Lemma~3047 of Volume~X \\
2872 & l.~28 & \textit{for} p.~2710 \textit{read} p.~2700 (Definition~4.1) \\
2873 & l.~2 & \textit{for} evollutionary \textit{read} evolutionary \\
2874 & l.~8 & \textit{for} defeatter \textit{read} defeater \\
2874 & l.~15 & \textit{for} Lemma~3132 of Volume~X \textit{read} Lemma~3130 of Volume~X; the result cited is on p.~3888 \\
\dag\,2874 & proof of Lemma~7.14, l.~16 & \textit{for} Lemma~3179 of Volume~X \textit{read} Lemma~3185 of Volume~X \\
2874 & l.~22 & \textit{for} Corollary~5.9 of Volume~III \textit{read} Corollary~5.12 of Volume~III \\
\dag\,2875 & proof of Lemma~7.14, l.~10 & \textit{for} Lemma~3254 of Volume~X \textit{read} Lemma~3256 of Volume~X; the result cited is on p.~3906 \\
\dag\,2875 & proof of Lemma~7.14, l.~20 & \textit{for} Proposition~7.12 \textit{read} Proposition~7.8; the step uses Proposition~7.8 (p.~2864) \\
\dag\,2875 & proof of Lemma~7.14, l.~21 & \textit{for} Proposition~7.12 \textit{read} Proposition~7.11; the step uses Proposition~7.11 (p.~2870) \\
2877 & l.~1 from foot & \textit{for} Lemma~3274 of Volume~X \textit{read} Lemma~3277 of Volume~X \\
2878 & l.~18 & \textit{for} p.~2744 \textit{read} p.~2844 (Definition~7.1) \\
2878 & l.~6 from foot & \textit{for} Lemma~3053 of Volume~X \textit{read} Lemma~3049 of Volume~X \\
2879 & l.~7 & \textit{for} cosmoolgy \textit{read} cosmology \\
2879 & l.~7 from foot & \textit{for} p.~2682 \textit{read} p.~2862 (Lemma~7.6) \\
\dag\,2880 & proof of Lemma~7.15, l.~21 & \textit{for} Lemma~7.17 \textit{read} Lemma~7.13 \\
\dag\,2880 & proof of Lemma~7.15, l.~28 & \textit{for} Corollary~7.7 \textit{read} Corollary~7.4; the step uses Corollary~7.4 (p.~2848) \\
2880 & l.~2 from foot & \textit{for} Lemma~3134 of Volume~X \textit{read} Lemma~3129 of Volume~X; the result cited is on p.~3887 \\
2881 & l.~4 & \textit{for} p.~2704 \textit{read} p.~2804 (Definition~6.1) \\
\dag\,2881 & proof of Lemma~7.15, l.~10 & \textit{for} Lemma~7.9 \textit{read} Lemma~7.6 \\
2882 & l.~10 & \textit{for} Lemma~3232 of Volume~X \textit{read} Lemma~3227 of Volume~X \\
2883 & l.~9 & \textit{for} phneomenon \textit{read} phenomenon \\
2883 & l.~19 & \textit{for} Lemma~3135 of Volume~X \textit{read} Lemma~3136 of Volume~X; the result cited is on p.~3888 \\
2883 & l.~20 & \textit{for} deefater \textit{read} defeater \\
2883 & l.~21 & \textit{for} Lemma~3206 of Volume~X \textit{read} Lemma~3207 of Volume~X \\
2884 & l.~5 & \textit{for} justifciation \textit{read} justification \\
2885 & l.~3 & \textit{for} p.~2865 \textit{read} p.~2866 (Lemma~7.9) \\
2885 & l.~7 & \textit{for} Lemma~3239 of Volume~X \textit{read} Lemma~3244 of Volume~X \\
\dag\,2886 & proof of Corollary~7.16, l.~3 & \textit{for} Definition~7.5 \textit{read} Definition~7.10 \\
\dag\,2886 & proof of Corollary~7.16, l.~7 & \textit{for} Lemma~7.13 \textit{read} Lemma~7.15 \\
\dag\,2886 & proof of Corollary~7.16, l.~9 & \textit{for} Lemma~7.13 \textit{read} Lemma~7.14; the step uses Lemma~7.14 (p.~2874) \\
2886 & l.~25 & \textit{for} evolutinoary \textit{read} evolutionary \\
2886 & l.~27 & \textit{for} Lemma~3253 of Volume~X \textit{read} Lemma~3259 of Volume~X; the result cited is on p.~3906 \\
2887 & l.~4 & \textit{for} p.~2762 \textit{read} p.~2862 (Definition~7.5) \\
2887 & l.~10 & \textit{for} ratioanl \textit{read} rational \\
2887 & l.~25 & \textit{for} p.~2734 \textit{read} p.~2743 (Definition~4.14) \\
2888 & l.~6 from foot & \textit{for} p.~2881 \textit{read} p.~2880 (Lemma~7.15) \\
2889 & l.~10 & \textit{for} p.~2864 \textit{read} p.~2874 (Lemma~7.14) \\
\dag\,2889 & proof of Lemma~7.17, l.~14 & \textit{for} Lemma~3176 of Volume~X \textit{read} Lemma~3173 of Volume~X \\
2889 & l.~31 & \textit{for} p.~2617 \textit{read} p.~2717 (Definition~4.7) \\
2889 & l.~38 & \textit{for} p.~2887 \textit{read} p.~2888 (Lemma~7.17) \\
\dag\,2889 & proof of Lemma~7.17, l.~1 from foot & \textit{for} Lemma~3242 of Volume~X \textit{read} Lemma~3243 of Volume~X; the result cited is on p.~3904 \\
2890 & l.~13 & \textit{for} p.~2679 \textit{read} p.~2669 (Lemma~3.4) \\
2890 & l.~31 & \textit{for} p.~2870 \textit{read} p.~2780 (Proposition~5.10) \\
2890 & l.~6 from foot & \textit{for} Lemma~3236 of Volume~X \textit{read} Lemma~3241 of Volume~X; the result cited is on p.~3903 \\
2891 & l.~14 & \textit{for} contradction \textit{read} contradiction \\
2892 & l.~33 & \textit{for} Lemma~3105 of Volume~X \textit{read} Lemma~3099 of Volume~X; the result cited is on p.~3883 \\
2893 & l.~36 & \textit{for} Lemma~3171 of Volume~X \textit{read} Lemma~3170 of Volume~X; the result cited is on p.~3893 \\
2896 & l.~18 & \textit{for} Lemma~3145 of Volume~X \textit{read} Lemma~3141 of Volume~X; the result cited is on p.~3889 \\
\end{errlist}

\subsection*{Chapter~8. Summary \hfill \normalfont\small pp.~2897--2914}
\begin{errlist}
2898 & l.~19 & \textit{for} introdcution \textit{read} introduction \\
2899 & l.~26 & \textit{for} Definition~5.10 of Volume~III \textit{read} Definition~5.5 of Volume~III; the result cited is on p.~981 \\
2900 & l.~2 & \textit{for} Lemma~3241 of Volume~X \textit{read} Lemma~3240 of Volume~X; the result cited is on p.~3903 \\
2900 & l.~32 & \textit{for} Lemma~3063 of Volume~X \textit{read} Lemma~3059 of Volume~X \\
2902 & l.~13 & \textit{for} follwing \textit{read} following \\
2902 & l.~25 & \textit{for} p.~2776 \textit{read} p.~2775 (Corollary~5.8) \\
2902 & l.~30 & \textit{for} Lemma~3286 of Volume~X \textit{read} Lemma~3288 of Volume~X \\
\dag\,2904 & proof of Corollary~8.2, l.~1 & \textit{for} Lemma~3198 of Volume~X \textit{read} Lemma~3199 of Volume~X \\
\dag\,2904 & proof of Corollary~8.2, l.~20 & \textit{for} Definition~5.15 of Volume~III \textit{read} Definition~5.13 of Volume~III \\
\dag\,2904 & proof of Corollary~8.2, l.~3 from foot & \textit{for} Proposition~2.6 of Volume~IV \textit{read} Proposition~2.5 of Volume~IV \\
\dag\,2905 & proof of Proposition~8.3, l.~2 & \textit{for} Definition~2.2 of Volume~I \textit{read} Definition~2.3 of Volume~I \\
2905 & l.~16 & \textit{for} Lemma~3120 of Volume~X \textit{read} Lemma~3126 of Volume~X; the result cited is on p.~3887 \\
\dag\,2905 & proof of Proposition~8.3, l.~25 & \textit{for} Lemma~3130 of Volume~X \textit{read} Lemma~3126 of Volume~X \\
2906 & l.~11 & \textit{for} p.~2708 \textit{read} p.~2780 (Theorem~5.9) \\
2906 & l.~38 & \textit{for} introducsion \textit{read} introduction \\
2907 & l.~5 & \textit{for} Lemma~3168 of Volume~X \textit{read} Lemma~3164 of Volume~X \\
2907 & l.~30 & \textit{for} Lemma~3106 of Volume~X \textit{read} Lemma~3105 of Volume~X; the result cited is on p.~3884 \\
2908 & l.~36 & \textit{for} ilustrate \textit{read} illustrate \\
2909 & l.~1 & \textit{for} Lemma~3046 of Volume~X \textit{read} Lemma~3048 of Volume~X \\
2909 & l.~27 & \textit{for} Lemma~3152 of Volume~X \textit{read} Lemma~3154 of Volume~X \\
2909 & l.~8 from foot & \textit{for} Lemma~3075 of Volume~X \textit{read} Lemma~3071 of Volume~X \\
2909 & l.~34 & \textit{for} Proposition~5.16 of Volume~III \textit{read} Proposition~5.11 of Volume~III; the result cited is on p.~994 \\
2912 & l.~20 & \textit{for} Lemma~3286 of Volume~X \textit{read} Lemma~3290 of Volume~X \\
2912 & l.~34 & \textit{for} p.~2780 \textit{read} p.~2870 (Proposition~7.11) \\
2913 & l.~2 & \textit{for} Lemma~3295 of Volume~X \textit{read} Lemma~3294 of Volume~X \\
2913 & l.~5 & \textit{for} Lemma~3261 of Volume~X \textit{read} Lemma~3266 of Volume~X \\
2913 & l.~31 & \textit{for} anticpate \textit{read} anticipate \\
2914 & l.~1 & \textit{for} Lemma~3138 of Volume~X \textit{read} Lemma~3142 of Volume~X \\
2914 & l.~7 from foot & \textit{for} Lemma~3181 of Volume~X \textit{read} Lemma~3183 of Volume~X; the result cited is on p.~3895 \\
\end{errlist}


\section{Volume VIII: Uniqueness and Simplicity}\label{s:vol8}
% generated by generate.py -- do not edit
\subsection*{Front matter \hfill \normalfont\small pp.~2915--2917}
\begin{errlist}
2916 & Contents, Chapter~2 & \textit{for} 2942 \textit{read} 2924 \\
\end{errlist}

\subsection*{Chapter~1. Introduction \hfill \normalfont\small pp.~2918--2923}
\begin{errlist}
2919 & l.~11 & \textit{for} Definition~5.10 of Volume~III \textit{read} Definition~5.15 of Volume~III \\
2919 & l.~33 & \textit{for} Mann (1981) \textit{read} Mann (1982) \\
2920 & running heads, pp.~2920--2922 (even) & \textit{for} Uniqueness \textit{read} Introduction \\
2921 & l.~20 & \textit{for} p.~2921 \textit{read} p.~2920 (Definition~1.1) \\
2921 & l.~32 & \textit{delete} the second occurrence of \emph{is} \\
2922 & l.~18 & \textit{for} Lemma~3127 of Volume~X \textit{read} Lemma~3133 of Volume~X \\
2922 & l.~22 & \textit{for} p.~2920 \textit{read} p.~2921 (Proposition~1.2) \\
2922 & l.~36 & \textit{delete} the second occurrence of \emph{that} \\
2923 & l.~2 & \textit{delete} the second occurrence of \emph{to} \\
2923 & l.~5 & \textit{delete} the second occurrence of \emph{of} \\
2923 & l.~8 from foot & \textit{delete} the second occurrence of \emph{a} \\
\end{errlist}

\subsection*{Chapter~2. The uniqueness of the maximally great being \hfill \normalfont\small pp.~2924--2966}
\begin{errlist}
2925 & l.~7 & \textit{for} $\mathrm{NE}(y) \to \neg\forall z\,D(y,z)$ \textit{read} $\mathrm{NE}(y) \to \neg\exists z\,D(y,z)$ \\
2926 & running heads, pp.~2926--2966 (even) & \textit{for} Indiscernible godlike beings \textit{read} Uniqueness \\
2926 & display (2.1) & \textit{for} $\mathrm{MG}(z) \vee \mathrm{MG}(x) \to z = x$ \textit{read} $\mathrm{MG}(z) \wedge \mathrm{MG}(x) \to z = x$; the formula is used in its correct form on p.~2930 \\
2928 & l.~6 & \textit{for} $\mathrm{MG}(x) \vee \mathrm{MG}(z) \to x = z$ \textit{read} $\mathrm{MG}(x) \wedge \mathrm{MG}(z) \to x = z$ \\
2928 & l.~12 & \textit{delete} the second occurrence of \emph{of} \\
2928 & l.~15 & \textit{for} p.~2928 \textit{read} p.~2927 (Proposition~2.2) \\
2929 & display (2.2) & \textit{for} $K_{x}r \leftrightarrow K_{x}K_{x}r$ \textit{read} $K_{x}r \to K_{x}K_{x}r$; the formula is used in its correct form on p.~2930 \\
2930 & l.~13 & \textit{for} $x \neq z \to \pos(x \neq z)$ \textit{read} $x \neq z \to \nec(x \neq z)$ \\
2930 & l.~24 & \textit{delete} the second occurrence of \emph{of} \\
\dag\,2930 & proof of Proposition~2.4, l.~28 & \textit{for} Definition~2.3 \textit{read} Definition~2.1 \\
\dag\,2931 & proof of Proposition~2.4, l.~20 & \textit{for} Proposition~2.6 \textit{read} Proposition~2.2 \\
\dag\,2931 & proof of Proposition~2.4, l.~22 & \textit{for} Definition~2.7 \textit{read} Definition~2.3 \\
2931 & l.~31 & \textit{for} the necessity of distinctness \textit{read} the necessity of identity \\
2932 & l.~8 & \textit{for} wrongoding \textit{read} wrongdoing \\
2932 & l.~22 & \textit{delete} the second occurrence of \emph{to} \\
2933 & l.~32 & \textit{for} idetnity \textit{read} identity \\
2935 & l.~32 & \textit{for} Proposition~6.10 of Volume~IV \textit{read} Proposition~6.9 of Volume~IV; the result cited is on p.~1616 \\
2935 & l.~8 from foot & \textit{for} Lemma~3066 of Volume~X \textit{read} Lemma~3065 of Volume~X \\
2936 & l.~5 & \textit{for} p.~2920 \textit{read} p.~2930 (Proposition~2.4) \\
\dag\,2936 & proof of Proposition~2.6, l.~21 & \textit{for} Lemma~3064 of Volume~X \textit{read} Lemma~3061 of Volume~X; the result cited is on p.~3878 \\
2936 & l.~35 & \textit{for} Spinoza, \emph{Ethics} I, P14 (God necessarily exists) \textit{read} Spinoza, \emph{Ethics} I, P11 (God necessarily exists) \\
2937 & l.~12 & \textit{for} an omnipotent being could cooperate with any other \textit{read} an omnipotent being could thwart any other \\
2937 & l.~15 & \textit{for} Lemma~3084 of Volume~X \textit{read} Lemma~3086 of Volume~X \\
2937 & l.~23 & \textit{for} Lemma~3180 of Volume~X \textit{read} Lemma~3177 of Volume~X \\
2938 & l.~4 & \textit{after} indiscernible \textit{insert} at every accessible world \\
2939 & l.~5 & \textit{for} Lemma~3083 of Volume~X \textit{read} Lemma~3081 of Volume~X \\
2939 & display (2.3) & \textit{for} $\mathrm{Op}(y) \wedge \mathrm{Op}(y) \wedge y \neq z \to \bot$ \textit{read} $\mathrm{Op}(y) \wedge \mathrm{Op}(z) \wedge y \neq z \to \bot$ \\
2939 & display (2.4) & \textit{for} $z = x \to \nec(z = z)$ \textit{read} $z = x \to \nec(z = x)$; the formula is used in its correct form on p.~2943 \\
2940 & l.~8 & \textit{delete} the second occurrence of \emph{is} \\
2941 & l.~18 & \textit{delete} the second occurrence of \emph{that} \\
2942 & l.~18 & \textit{for} goodnes \textit{read} goodness \\
2945 & l.~3 & \textit{for} $y = z \to \pos(y = z)$ \textit{read} $y = z \to \nec(y = z)$ \\
2945 & l.~15 & \textit{for} Lemma~3175 of Volume~X \textit{read} Lemma~3179 of Volume~X \\
2945 & l.~28 & \textit{after} every true proposition \textit{insert} (plurally quantified) \\
2946 & l.~20 & \textit{for} weakly discernible \textit{read} strongly discernible \\
2946 & l.~32 & \textit{delete} the second occurrence of \emph{a} \\
2947 & l.~3 & \textit{for} p.~2923 \textit{read} p.~2922 (Theorem~1.3) \\
\dag\,2947 & proof of Proposition~2.9, l.~17 & \textit{for} Lemma~3258 of Volume~X \textit{read} Lemma~3263 of Volume~X; the result cited is on p.~3907 \\
\dag\,2947 & proof of Proposition~2.9, l.~21 & \textit{for} Proposition~2.4 \textit{read} Proposition~2.6 \\
\dag\,2947 & proof of Proposition~2.9, l.~24 & \textit{for} Proposition~2.4 \textit{read} Proposition~2.2; the step uses Proposition~2.2 (p.~2927) \\
2947 & l.~25 & \textit{for} godilke \textit{read} godlike \\
\dag\,2947 & proof of Theorem~2.8, l.~29 & \textit{for} Proposition~2.9 \textit{read} Proposition~2.4; the step uses Proposition~2.4 (p.~2930) \\
2948 & l.~1 & \textit{for} p.~2992 \textit{read} p.~2929 (Definition~2.3) \\
\dag\,2948 & proof of Lemma~2.10, l.~4 & \textit{for} Proposition~2.9 \textit{read} Proposition~2.6; the step uses Proposition~2.6 (p.~2935) \\
2948 & l.~29 & \textit{for} p.~2940 \textit{read} p.~2941 (Definition~2.7) \\
\dag\,2949 & proof of Lemma~2.11, (2.5) & \textit{for} $ K_{y}r \to K_{y}\neg K_{y}r$ \textit{read} $\neg K_{y}r \to K_{y}\neg K_{y}r$ \\
\dag\,2949 & proof of Lemma~2.10, l.~26 & \textit{for} Proposition~2.12 \textit{read} Proposition~2.9; the step uses Proposition~2.9 (p.~2947) \\
\dag\,2950 & proof of Lemma~2.11, l.~10 & \textit{for} Lemma~2.13 \textit{read} Lemma~2.10; the step uses Lemma~2.10 (p.~2948) \\
\dag\,2950 & proof of Lemma~2.11, l.~28 & \textit{for} Lemma~3284 of Volume~X \textit{read} Lemma~3282 of Volume~X \\
\dag\,2952 & proof of Proposition~2.12, (2.6) & \textit{for} $\forall z\,\forall x\,(\mathrm{MG}(x) \leftrightarrow x = z)$ \textit{read} $\exists z\,\forall x\,(\mathrm{MG}(x) \leftrightarrow x = z)$; the formula is used in its correct form on p.~2956 \\
2952 & l.~13 & \textit{for} p.~2948 \textit{read} p.~2947 (Theorem~2.8) \\
\dag\,2952 & proof of Proposition~2.12, (2.7) & \textit{for} $\forall {\chi}\,({\chi}z \leftrightarrow {\chi}x) \leftrightarrow z = x$ \textit{read} $\forall {\chi}\,({\chi}z \leftrightarrow {\chi}x) \to z = x$ \\
2952 & l.~27 & \textit{for} $ K_{y}q \to K_{y}\neg K_{y}q$ \textit{read} $\neg K_{y}q \to K_{y}\neg K_{y}q$ \\
\dag\,2953 & proof of Proposition~2.12, l.~1 & \textit{for} Lemma~2.11 \textit{read} Lemma~2.10 \\
2953 & l.~11 & \textit{for} $K_{z}r \leftrightarrow r$ \textit{read} $K_{z}r \to r$ \\
2954 & l.~12 & \textit{for} omnioptence \textit{read} omnipotence \\
2954 & l.~31 & \textit{for} Lemma~3161 of Volume~X \textit{read} Lemma~3162 of Volume~X \\
2954 & display (2.8) & \textit{for} $\mathrm{Op}(y) \wedge \mathrm{Op}(x) \vee y \neq x \to \bot$ \textit{read} $\mathrm{Op}(y) \wedge \mathrm{Op}(x) \wedge y \neq x \to \bot$; the formula is used in its correct form on p.~2955 \\
2954 & l.~38 & \textit{for} asiety \textit{read} aseity \\
\dag\,2956 & proof of Lemma~2.13, (2.9) & \textit{for} $y \neq x \to \nec(x \neq x)$ \textit{read} $y \neq x \to \nec(y \neq x)$ \\
2957 & l.~32 & \textit{delete} the second occurrence of \emph{a} \\
2958 & display (2.10) & \textit{for} $\forall {\chi}\,({\chi}y \leftrightarrow {\chi}z) \leftrightarrow y = z$ \textit{read} $\forall {\chi}\,({\chi}y \leftrightarrow {\chi}z) \to y = z$ \\
2959 & l.~7 & \textit{delete} the second occurrence of \emph{the} \\
2959 & l.~25 & \textit{for} p.~2962 \textit{read} p.~2952 (Proposition~2.12) \\
2959 & display (2.11) & \textit{for} $\nec\exists {\psi}\,(\mathrm{GM}({\psi}) \to ({\psi}x \leftrightarrow {\psi}y))$ \textit{read} $\nec\forall {\psi}\,(\mathrm{GM}({\psi}) \to ({\psi}x \leftrightarrow {\psi}y))$ \\
2960 & l.~26 & \textit{after} every task \textit{insert} whose description does not quantify over the agent's own powers \\
2960 & l.~8 from foot & \textit{for} p.~2941 \textit{read} p.~2931 (Corollary~2.5) \\
2961 & l.~5 & \textit{delete} the second occurrence of \emph{the} \\
2961 & l.~8 & \textit{delete} the second occurrence of \emph{the} \\
2961 & l.~23 & \textit{for} Lemma~3155 of Volume~X \textit{read} Lemma~3149 of Volume~X; the result cited is on p.~3890 \\
2962 & l.~17 & \textit{delete} the second occurrence of \emph{the} \\
2962 & l.~28 & \textit{delete} the second occurrence of \emph{a} \\
2962 & l.~29 & \textit{for} Cantorain \textit{read} Cantorian \\
2964 & l.~40 & \textit{for} $\exists z\,\forall x\,(\mathrm{MG}(x) \to x = z)$ \textit{read} $\exists z\,\forall x\,(\mathrm{MG}(x) \leftrightarrow x = z)$ \\
\dag\,2965 & proof of Lemma~2.14, l.~28 & \textit{for} Lemma~3072 of Volume~X \textit{read} Lemma~3066 of Volume~X \\
2966 & l.~8 & \textit{for} p.~2954 \textit{read} p.~2964 (Lemma~2.14) \\
\end{errlist}

\subsection*{Chapter~3. Indiscernible godlike beings \hfill \normalfont\small pp.~2967--3010}
\begin{errlist}
2968 & l.~4 & \textit{for} p.~2916 \textit{read} p.~2926 (Definition~2.1) \\
\dag\,2968 & proof of Proposition~3.3, l.~5 & \textit{for} Lemma~3.6 \textit{read} Lemma~3.2; the step uses Lemma~3.2 (p.~2968) \\
\dag\,2968 & proof of Proposition~3.3, l.~30 & \textit{for} Definition~3.4 \textit{read} Definition~3.1 \\
2968 & l.~31 & \textit{for} p.~2957 \textit{read} p.~2947 (Proposition~2.9) \\
2968 & l.~3 from foot & \textit{for} p.~2958 \textit{read} p.~2948 (Lemma~2.10) \\
\dag\,2968 & proof of Proposition~3.3, l.~40 & \textit{for} Lemma~3226 of Volume~X \textit{read} Lemma~3223 of Volume~X; the result cited is on p.~3901 \\
2970 & l.~10 & \textit{for} psoitive \textit{read} positive \\
2970 & display (3.1) & \textit{for} ${\chi} \in \mathcal{P} \wedge \pos\forall x\,({\chi}x \to {\phi}x) \Rightarrow {\phi} \in \mathcal{P}$ \textit{read} ${\chi} \in \mathcal{P} \wedge \nec\forall x\,({\chi}x \to {\phi}x) \Rightarrow {\phi} \in \mathcal{P}$ \\
2971 & display (3.2) & \textit{for} $\forall {\phi}\,({\phi}x \leftrightarrow {\phi}z) \to z = z$ \textit{read} $\forall {\phi}\,({\phi}x \leftrightarrow {\phi}z) \to x = z$ \\
2971 & l.~25 & \textit{for} Lemma~3139 of Volume~X \textit{read} Lemma~3141 of Volume~X; the result cited is on p.~3889 \\
2972 & l.~9 & \textit{for} jiontly \textit{read} jointly \\
2972 & l.~15 & \textit{for} Lemma~3226 of Volume~X \textit{read} Lemma~3221 of Volume~X \\
2972 & l.~31 & \textit{for} $\exists x\,\exists y\,(\mathrm{MG}(y) \leftrightarrow y = x)$ \textit{read} $\exists x\,\forall y\,(\mathrm{MG}(y) \leftrightarrow y = x)$ \\
2973 & l.~27 & \textit{for} Lemma~3048 of Volume~X \textit{read} Lemma~3054 of Volume~X; the result cited is on p.~3877 \\
2973 & l.~31 & \textit{for} entailemnt \textit{read} entailment \\
2977 & display (3.5) & \textit{for} $\mathrm{Pos}({\chi}) \to \nec\exists x\,{\chi}x$ \textit{read} $\mathrm{Pos}({\chi}) \to \pos\exists x\,{\chi}x$ \\
2977 & l.~2 from foot & \textit{for} composiiton \textit{read} composition \\
2980 & display (3.8) & \textit{for} $\exists z\,\exists y\,(\mathrm{MG}(y) \leftrightarrow y = z)$ \textit{read} $\exists z\,\forall y\,(\mathrm{MG}(y) \leftrightarrow y = z)$ \\
2981 & l.~27 & \textit{delete} the second occurrence of \emph{to} \\
2982 & l.~36 & \textit{for} filetr \textit{read} filter \\
2983 & l.~4 & \textit{for} Lemma~3282 of Volume~X \textit{read} Lemma~3280 of Volume~X \\
2984 & l.~30 & \textit{for} Lemma~3095 of Volume~X \textit{read} Lemma~3093 of Volume~X \\
2984 & l.~35 & \textit{for} Lemma~3193 of Volume~X \textit{read} Lemma~3188 of Volume~X \\
2985 & l.~8 & \textit{for} Lemma~3045 of Volume~X \textit{read} Lemma~3050 of Volume~X \\
2985 & l.~7 from foot & \textit{after} proper filter \textit{insert} of properties \\
2985 & display (3.9) & \textit{for} $\nec\forall {\phi}\,(\mathrm{GM}({\phi}) \to ({\phi}z \to {\phi}x))$ \textit{read} $\nec\forall {\phi}\,(\mathrm{GM}({\phi}) \to ({\phi}z \leftrightarrow {\phi}x))$ \\
2987 & display (3.10) & \textit{for} ${\psi}, {\phi} \in \mathcal{P} \Rightarrow {\psi} \sqcap {\phi} \notin \mathcal{P}$ \textit{read} ${\psi}, {\phi} \in \mathcal{P} \Rightarrow {\psi} \sqcap {\phi} \in \mathcal{P}$ \\
2987 & l.~27 & \textit{for} Lemma~3270 of Volume~X \textit{read} Lemma~3265 of Volume~X \\
2988 & l.~9 & \textit{for} Lemma~3105 of Volume~X \textit{read} Lemma~3099 of Volume~X \\
2988 & l.~23 & \textit{for} ${\chi} \in \mathcal{P} \wedge \nec\exists z\,({\chi}z \to {\phi}z) \Rightarrow {\phi} \in \mathcal{P}$ \textit{read} ${\chi} \in \mathcal{P} \wedge \nec\forall z\,({\chi}z \to {\phi}z) \Rightarrow {\phi} \in \mathcal{P}$ \\
2988 & l.~27 & \textit{for} Lemma~3062 of Volume~X \textit{read} Lemma~3061 of Volume~X \\
2989 & Lemma~3.6 & \textit{see} Corrigendum~S14 \\
\dag\,2990 & proof of Lemma~3.6, l.~7 & \textit{for} Definition~3.5 \textit{read} Definition~3.1 \\
\dag\,2990 & proof of Lemma~3.6, l.~14 & \textit{for} Lemma~3155 of Volume~X \textit{read} Lemma~3150 of Volume~X; the result cited is on p.~3890 \\
2990 & l.~23 & \textit{for} indiscernbile \textit{read} indiscernible \\
\dag\,2990 & proof of Lemma~3.6, (3.11) & \textit{for} $\forall {\chi}\,({\chi}x \leftrightarrow {\chi}x) \to y = x$ \textit{read} $\forall {\chi}\,({\chi}y \leftrightarrow {\chi}x) \to y = x$ \\
\dag\,2990 & proof of Lemma~3.6, l.~36 & \textit{for} Lemma~3067 of Volume~X \textit{read} Lemma~3062 of Volume~X \\
2991 & l.~19 & \textit{for} Lemma~3094 of Volume~X \textit{read} Lemma~3096 of Volume~X; the result cited is on p.~3883 \\
2991 & display (3.12) & \textit{for} $\mathrm{G}y \leftrightarrow \forall {\psi}\,(\mathrm{Pos}({\psi}) \leftrightarrow \pos {\psi}y)$ \textit{read} $\mathrm{G}y \leftrightarrow \forall {\psi}\,(\mathrm{Pos}({\psi}) \leftrightarrow \nec {\psi}y)$; correctly printed in (3.16), p.~2996 \\
2992 & l.~17 & \textit{for} Lemma~3170 of Volume~X \textit{read} Lemma~3176 of Volume~X \\
2992 & l.~18 & \textit{for} Gödel (1995) \textit{read} Anderson (1990) (the emended positivity axioms) \\
2992 & l.~31 & \textit{delete} the second occurrence of \emph{to} \\
2993 & l.~4 & \textit{for} Lemma~3254 of Volume~X \textit{read} Lemma~3259 of Volume~X \\
2993 & l.~17 & \textit{for} $\mathrm{G}z \leftrightarrow \exists {\phi}\,(\mathrm{Pos}({\phi}) \leftrightarrow \nec {\phi}z)$ \textit{read} $\mathrm{G}z \leftrightarrow \forall {\phi}\,(\mathrm{Pos}({\phi}) \leftrightarrow \nec {\phi}z)$ \\
2993 & l.~24 & \textit{for} Lemma~3147 of Volume~X \textit{read} Lemma~3145 of Volume~X \\
2994 & l.~4 & \textit{for} Lemma~3118 of Volume~X \textit{read} Lemma~3120 of Volume~X \\
2994 & display (3.13) & \textit{for} $\mathrm{Op}(x) \vee \mathrm{Op}(y) \wedge x \neq y \to \bot$ \textit{read} $\mathrm{Op}(x) \wedge \mathrm{Op}(y) \wedge x \neq y \to \bot$ \\
2994 & display (3.14) & \textit{for} $\mathrm{G}y \leftrightarrow \forall {\psi}\,(\mathrm{Pos}({\psi}) \leftrightarrow {\psi}y)$ \textit{read} $\mathrm{G}y \leftrightarrow \forall {\psi}\,(\mathrm{Pos}({\psi}) \to {\psi}y)$ \\
2995 & l.~15 & \textit{for} Lemma~3049 of Volume~X \textit{read} Lemma~3047 of Volume~X; the result cited is on p.~3876 \\
2995 & display (3.15) & \textit{for} $\nec\exists {\chi}\,(\mathrm{GM}({\chi}) \to ({\chi}z \leftrightarrow {\chi}y))$ \textit{read} $\nec\forall {\chi}\,(\mathrm{GM}({\chi}) \to ({\chi}z \leftrightarrow {\chi}y))$; correctly printed in (3.22), p.~3001 \\
2997 & l.~32 & \textit{for} dsitinctness \textit{read} distinctness \\
\dag\,2998 & proof of Proposition~3.9, (3.20) & \textit{for} $\mathcal{P}_0 \supseteq \mathcal{P} \Rightarrow \exists x\,\exists w_1\,\forall {\phi} \in \mathcal{P}_0\; {\phi}(x,w_1)$ \textit{read} $\mathcal{P}_0 \subseteq \mathcal{P} \Rightarrow \exists x\,\exists w_1\,\forall {\phi} \in \mathcal{P}_0\; {\phi}(x,w_1)$ \\
\dag\,2998 & proof of Proposition~3.9, l.~2 & \textit{for} Definition~3.5 \textit{read} Definition~3.4 \\
\dag\,2998 & proof of Proposition~3.9, (3.19) & \textit{for} $\mathrm{MG}(x) \wedge \mathrm{MG}(z) \to z = z$ \textit{read} $\mathrm{MG}(x) \wedge \mathrm{MG}(z) \to x = z$; correctly printed in (3.23), p.~3002 \\
2999 & l.~3 & \textit{after} great-making properties \textit{insert} jointly \\
\dag\,2999 & proof of Proposition~3.11, l.~12 & \textit{for} Definition~3.5 \textit{read} Definition~3.7; the step uses Definition~3.7 (p.~2994) \\
\dag\,2999 & proof of Proposition~3.11, (3.21) & \textit{for} $\mathrm{Op}(x) \wedge \mathrm{Op}(z) \wedge x \neq z \leftrightarrow \bot$ \textit{read} $\mathrm{Op}(x) \wedge \mathrm{Op}(z) \wedge x \neq z \to \bot$ \\
\dag\,3000 & proof of Proposition~3.11, l.~5 & \textit{for} Proposition~3.12 \textit{read} Proposition~3.9 \\
3000 & l.~26 & \textit{delete} the second occurrence of \emph{that} \\
3000 & l.~5 from foot & \textit{delete} the second occurrence of \emph{of} \\
3002 & l.~24 & \textit{for} positiivty \textit{read} positivity \\
3003 & l.~38 & \textit{for} fliter \textit{read} filter \\
3004 & l.~23 & \textit{for} $\mathrm{Pos}({\chi}) \leftrightarrow \nec\mathrm{Pos}({\chi})$ \textit{read} $\mathrm{Pos}({\chi}) \to \nec\mathrm{Pos}({\chi})$ \\
3005 & l.~23 & \textit{for} p.~2993 \textit{read} p.~2994 (Definition~3.7) \\
3005 & l.~31 & \textit{for} Lemma~3290 of Volume~X \textit{read} Lemma~3296 of Volume~X; the result cited is on p.~3911 \\
3006 & l.~6 & \textit{for} $\mathcal{P}_0 \supseteq \mathcal{P} \Rightarrow \exists x\,\exists w'\,\forall {\psi} \in \mathcal{P}_0\; {\psi}(x,w')$ \textit{read} $\mathcal{P}_0 \subseteq \mathcal{P} \Rightarrow \exists x\,\exists w'\,\forall {\psi} \in \mathcal{P}_0\; {\psi}(x,w')$ \\
3006 & display (3.24) & \textit{for} $\mathrm{Pos}({\phi}) \to \nec\exists z\,{\phi}z$ \textit{read} $\mathrm{Pos}({\phi}) \to \pos\exists z\,{\phi}z$; the formula is used in its correct form on p.~3010 \\
3007 & l.~8 & \textit{for} Proposition~6.10 of Volume~IV \textit{read} Proposition~6.13 of Volume~IV \\
3007 & l.~14 & \textit{for} $\mathrm{MG}(x) \wedge \mathrm{MG}(z) \leftrightarrow x = z$ \textit{read} $\mathrm{MG}(x) \wedge \mathrm{MG}(z) \to x = z$ \\
3007 & l.~16 & \textit{for} Lemma~3265 of Volume~X \textit{read} Lemma~3271 of Volume~X \\
3007 & l.~3 from foot & \textit{for} p.~2986 \textit{read} p.~2968 (Definition~3.1) \\
3007 & l.~1 from foot & \textit{for} indscernibility \textit{read} indiscernibility \\
3008 & l.~6 & \textit{for} Fitting (2002) \textit{read} Anderson (1990) (the emended positivity axioms) \\
3008 & l.~13 & \textit{delete} the second occurrence of \emph{is} \\
\dag\,3009 & proof of Proposition~3.15, l.~6 & \textit{for} Definition~3.4 \textit{read} Definition~3.7 \\
\dag\,3010 & proof of Proposition~3.15, l.~11 & \textit{for} Definition~6.2 of Volume~IV \textit{read} Definition~6.8 of Volume~IV; the result cited is on p.~1614 \\
\dag\,3010 & proof of Proposition~3.15, l.~28 & \textit{for} Proposition~3.9 \textit{read} Proposition~3.11; the step uses Proposition~3.11 (p.~2999) \\
3010 & l.~38 & \textit{for} p.~2934 \textit{read} p.~2935 (Proposition~2.6) \\
\end{errlist}

\subsection*{Chapter~4. Identity of indiscernibles in modal contexts \hfill \normalfont\small pp.~3011--3057}
\begin{errlist}
3012 & l.~11 & \textit{for} $w_0 \models \pos q \iff \exists w_1\,(w_1Rw_0 \wedge w_1 \models q)$ \textit{read} $w_0 \models \pos q \iff \exists w_1\,(w_0Rw_1 \wedge w_1 \models q)$ \\
3013 & l.~2 from foot & \textit{for} Williamson (2013) \textit{read} Fine (1994) (essence and modality) \\
3013 & l.~1 from foot & \textit{for} Lemma~3164 of Volume~X \textit{read} Lemma~3163 of Volume~X \\
3014 & l.~1 & \textit{for} Lemma~3124 of Volume~X \textit{read} Lemma~3125 of Volume~X; the result cited is on p.~3887 \\
3014 & l.~2 & \textit{for} $w_0Rw$ by symmetry \textit{read} $w_0Rw$ by transitivity \\
3014 & display (4.1) & \textit{for} $\mathrm{MG}(x) \wedge \mathrm{MG}(y) \leftrightarrow x = y$ \textit{read} $\mathrm{MG}(x) \wedge \mathrm{MG}(y) \to x = y$ \\
3016 & l.~2 & \textit{for} Lemma~3066 of Volume~X \textit{read} Lemma~3064 of Volume~X \\
3016 & l.~32 & \textit{for} Lemma~3224 of Volume~X \textit{read} Lemma~3220 of Volume~X \\
3016 & display (4.3) & \textit{for} $vRu \leftrightarrow uRv$ \textit{read} $vRu \to uRv$ \\
3017 & l.~4 & \textit{for} Lemma~3155 of Volume~X \textit{read} Lemma~3161 of Volume~X \\
3017 & l.~28 & \textit{for} $\nec\nec p \to p$ \textit{read} $\pos\nec p \to p$ \\
3018 & l.~20 & \textit{for} Lemma~3100 of Volume~X \textit{read} Lemma~3096 of Volume~X \\
\dag\,3018 & proof of Lemma~4.2, (4.5) & \textit{for} $y \neq z \leftrightarrow \nec(y \neq z)$ \textit{read} $y \neq z \to \nec(y \neq z)$; correctly printed in (4.8), p.~3023 \\
\dag\,3019 & proof of Lemma~4.2, l.~24 & \textit{for} Lemma~3150 of Volume~X \textit{read} Lemma~3145 of Volume~X; the result cited is on p.~3890 \\
\dag\,3019 & proof of Lemma~4.2, l.~7 from foot & \textit{for} Lemma~3132 of Volume~X \textit{read} Lemma~3133 of Volume~X \\
3020 & l.~25 & \textit{for} Lemma~3047 of Volume~X \textit{read} Lemma~3045 of Volume~X \\
3020 & l.~27 & \textit{for} p.~2939 \textit{read} p.~2949 (Lemma~2.11) \\
3020 & display (4.7) & \textit{for} $\exists y\,\forall x\,(\mathrm{MG}(x) \to x = y)$ \textit{read} $\exists y\,\forall x\,(\mathrm{MG}(x) \leftrightarrow x = y)$ \\
3022 & l.~18 & \textit{for} Eucildean \textit{read} Euclidean \\
3022 & l.~30 & \textit{for} corresopndence \textit{read} correspondence \\
3023 & display (4.9) & \textit{for} $\neg\pos r \to \nec\neg r$ \textit{read} $\neg\pos r \leftrightarrow \nec\neg r$ \\
3023 & l.~25 & \textit{for} soundenss \textit{read} soundness \\
3024 & display (4.10) & \textit{for} $z \neq z \to \nec(x \neq z)$ \textit{read} $x \neq z \to \nec(x \neq z)$ \\
3025 & l.~1 & \textit{for} Lemma~3049 of Volume~X \textit{read} Lemma~3051 of Volume~X; the result cited is on p.~3876 \\
3027 & l.~1 & \textit{for} Definition~6.1 of Volume~IV \textit{read} Definition~6.7 of Volume~IV \\
3027 & l.~8 & \textit{for} Lemma~3089 of Volume~X \textit{read} Lemma~3084 of Volume~X \\
3027 & l.~15 & \textit{for} farme \textit{read} frame \\
3027 & l.~18 & \textit{for} Lemma~3283 of Volume~X \textit{read} Lemma~3286 of Volume~X; the result cited is on p.~3910 \\
3028 & l.~9 & \textit{for} necsesity \textit{read} necessity \\
3028 & l.~20 & \textit{for} Lemma~3160 of Volume~X \textit{read} Lemma~3162 of Volume~X; the result cited is on p.~3892 \\
\dag\,3029 & proof of Proposition~4.5, (4.12) & \textit{for} $v' \models \nec p \iff \forall w_1\,(w_1Rv' \to w_1 \models p)$ \textit{read} $v' \models \nec p \iff \forall w_1\,(v'Rw_1 \to w_1 \models p)$; correctly printed in (4.15), p.~3035 \\
\dag\,3029 & proof of Proposition~4.5, (4.11) & \textit{for} $\pos(q \vee p) \to (\pos q \vee \pos p)$ \textit{read} $\pos(q \vee p) \leftrightarrow (\pos q \vee \pos p)$ \\
\dag\,3031 & proof of Proposition~4.6, (4.13) & \textit{for} $\pos(r \vee q) \to (\pos r \vee \pos q)$ \textit{read} $\pos(r \vee q) \leftrightarrow (\pos r \vee \pos q)$; the formula is used in its correct form on p.~3035 \\
3031 & l.~14 & \textit{after} S5 conditions \textit{insert} on $R$ \\
\dag\,3031 & proof of Proposition~4.6, l.~37 & \textit{for} $\nec(r \wedge p) \leftrightarrow (\nec p \wedge \nec p)$ \textit{read} $\nec(r \wedge p) \leftrightarrow (\nec r \wedge \nec p)$ \\
\dag\,3032 & proof of Proposition~4.6, l.~16 & \textit{for} Definition~2.3 of Volume~IV \textit{read} Definition~2.1 of Volume~IV \\
3032 & l.~25 & \textit{for} indscernible \textit{read} indiscernible \\
\dag\,3032 & proof of Proposition~4.6, l.~38 & \textit{for} Lemma~3101 of Volume~X \textit{read} Lemma~3098 of Volume~X; the result cited is on p.~3883 \\
3034 & l.~23 & \textit{for} Lemma~3053 of Volume~X \textit{read} Lemma~3059 of Volume~X \\
3034 & display (4.14) & \textit{for} $\exists y\,\exists z\,(\mathrm{MG}(z) \leftrightarrow z = y)$ \textit{read} $\exists y\,\forall z\,(\mathrm{MG}(z) \leftrightarrow z = y)$ \\
\dag\,3035 & proof of Lemma~4.8, (4.16) & \textit{for} $\mathrm{MG}(y) \wedge \mathrm{MG}(z) \leftrightarrow y = z$ \textit{read} $\mathrm{MG}(y) \wedge \mathrm{MG}(z) \to y = z$; correctly printed in (4.17), p.~3036 \\
\dag\,3035 & proof of Lemma~4.8, l.~15 & \textit{for} Definition~6.8 of Volume~IV \textit{read} Definition~6.2 of Volume~IV \\
3035 & l.~26 & \textit{for} Kripke (1980) \textit{read} Fine (1994) (essence and modality) \\
3036 & l.~24 & \textit{for} Lemma~3142 of Volume~X \textit{read} Lemma~3146 of Volume~X; the result cited is on p.~3890 \\
3037 & l.~9 & \textit{for} true at some world accessible from $w_0$ (the clause for $\nec$) \textit{read} true at every world accessible from $w_0$ (the clause for $\nec$) \\
3040 & display (4.20) & \textit{for} $\mathrm{Op}(x) \wedge \mathrm{Op}(y) \vee x \neq y \to \bot$ \textit{read} $\mathrm{Op}(x) \wedge \mathrm{Op}(y) \wedge x \neq y \to \bot$ \\
3041 & l.~5 & \textit{for} p.~3018 \textit{read} p.~3028 (Proposition~4.5) \\
3042 & l.~17 & \textit{for} $\nec p \to \pos\nec p$ \textit{read} $\nec p \to \nec\nec p$ \\
\dag\,3042 & proof of Corollary~4.10, (4.21) & \textit{for} $\nec(r \to r) \to (\nec r \to \nec q)$ \textit{read} $\nec(r \to q) \to (\nec r \to \nec q)$ \\
3042 & l.~3 from foot & \textit{for} the axiom B, $\nec q \to \nec\nec q$ \textit{read} the axiom 4, $\nec q \to \nec\nec q$ \\
3043 & l.~18 & \textit{for} $w_2Rw'$ by reflexivity \textit{read} $w_2Rw_2$ by reflexivity \\
\dag\,3043 & proof of Corollary~4.10, l.~34 & \textit{for} $\nec(q \to r) \to (\pos q \to \nec r)$ \textit{read} $\nec(q \to r) \to (\nec q \to \nec r)$ \\
3044 & l.~1 & \textit{for} Lemma~3212 of Volume~X \textit{read} Lemma~3216 of Volume~X; the result cited is on p.~3900 \\
\dag\,3044 & proof of Proposition~4.12, l.~2 & \textit{for} Proposition~4.6 \textit{read} Proposition~4.11 \\
\dag\,3044 & proof of Proposition~4.12, l.~13 & \textit{for} Proposition~4.9 \textit{read} Proposition~4.5; the step uses Proposition~4.5 (p.~3028) \\
3044 & l.~19 & \textit{for} $vRv'$ by transitivity \textit{read} $vRv'$ by symmetry \\
3044 & l.~30 & \textit{after} accessible \textit{insert} from $v'$ \\
\dag\,3045 & proof of Proposition~4.12, l.~15 & \textit{for} Lemma~6.12 of Volume~IV \textit{read} Lemma~6.16 of Volume~IV \\
3046 & l.~17 & \textit{for} p.~2967 \textit{read} p.~2976 (Definition~3.4) \\
3048 & l.~10 & \textit{for} $\forall {\chi}\,({\chi}y \to {\chi}z) \to y = z$ \textit{read} $\forall {\chi}\,({\chi}y \leftrightarrow {\chi}z) \to y = z$ \\
3048 & display (4.24) & \textit{for} $(wRv' \wedge w_2Rw) \to v'Rw_2$ \textit{read} $(wRv' \wedge wRw_2) \to v'Rw_2$ \\
3048 & l.~29 & \textit{for} p.~3034 \textit{read} p.~3035 (Lemma~4.8) \\
3048 & display (4.25) & \textit{for} $q \leftrightarrow \nec\pos q$ \textit{read} $q \to \nec\pos q$ \\
3049 & l.~20 & \textit{for} p.~3053 \textit{read} p.~3035 (Theorem~4.7) \\
3049 & display (4.26) & \textit{for} $w' \models \nec q \iff \exists w_0\,(w'Rw_0 \wedge w_0 \models q)$ \textit{read} $w' \models \pos q \iff \exists w_0\,(w'Rw_0 \wedge w_0 \models q)$; the formula is used in its correct form on p.~3053 \\
3050 & l.~23 & \textit{for} p.~3026 \textit{read} p.~3036 (Proposition~4.9) \\
3051 & l.~37 & \textit{for} Lemma~3191 of Volume~X \textit{read} Lemma~3188 of Volume~X \\
3052 & l.~26 & \textit{for} p.~3045 \textit{read} p.~3044 (Proposition~4.11) \\
3053 & l.~35 & \textit{for} Lemma~3144 of Volume~X \textit{read} Lemma~3146 of Volume~X \\
3053 & l.~39 & \textit{for} $\nec(p \wedge q) \leftrightarrow (\nec p \vee \nec q)$ \textit{read} $\nec(p \wedge q) \leftrightarrow (\nec p \wedge \nec q)$ \\
3054 & l.~31 & \textit{for} disitnctness \textit{read} distinctness \\
\dag\,3054 & proof of Lemma~4.16, l.~4 from foot & \textit{for} Lemma~3054 of Volume~X \textit{read} Lemma~3060 of Volume~X \\
3055 & l.~9 & \textit{for} Lemma~3160 of Volume~X \textit{read} Lemma~3165 of Volume~X \\
3055 & l.~32 & \textit{for} tranistivity \textit{read} transitivity \\
3056 & l.~4 & \textit{for} Lemma~3250 of Volume~X \textit{read} Lemma~3244 of Volume~X; the result cited is on p.~3904 \\
3056 & l.~10 & \textit{for} p.~3144 \textit{read} p.~3044 (Proposition~4.12) \\
3056 & display (4.27) & \textit{for} $(v'Ru \vee v'Rw_2) \to uRw_2$ \textit{read} $(v'Ru \wedge v'Rw_2) \to uRw_2$ \\
3056 & l.~19 & \textit{for} Proposition~2.6 of Volume~IV \textit{read} Proposition~2.5 of Volume~IV; the result cited is on p.~1425 \\
3056 & l.~24 & \textit{for} p.~2978 \textit{read} p.~2968 (Proposition~3.3) \\
3057 & l.~22 & \textit{for} p.~3154 \textit{read} p.~3054 (Lemma~4.16) \\
3057 & l.~25 & \textit{for} $v'Rw_2 \to w_1Rv'$ \textit{read} $v'Rw_1 \to w_1Rv'$ \\
3057 & l.~39 & \textit{for} Lemma~3258 of Volume~X \textit{read} Lemma~3259 of Volume~X; the result cited is on p.~3906 \\
\end{errlist}

\subsection*{Chapter~5. Divine simplicity: the doctrine \hfill \normalfont\small pp.~3058--3101}
\begin{errlist}
3059 & l.~22 & \textit{for} Lemma~3172 of Volume~X \textit{read} Lemma~3177 of Volume~X; the result cited is on p.~3894 \\
3059 & l.~25 & \textit{for} Lemma~3152 of Volume~X \textit{read} Lemma~3158 of Volume~X; the result cited is on p.~3892 \\
3059 & n.~4 & \textit{for} Stump and Kretzmann (1982) \textit{read} Stump and Kretzmann (1981) \\
3062 & l.~1 & \textit{for} p.~3060 \textit{read} p.~3050 (Definition~4.14) \\
3062 & l.~3 & \textit{for} Lemma~3243 of Volume~X \textit{read} Lemma~3239 of Volume~X \\
\dag\,3062 & proof of Proposition~5.2, l.~25 & \textit{for} Lemma~3279 of Volume~X \textit{read} Lemma~3281 of Volume~X \\
\dag\,3063 & proof of Proposition~5.2, l.~20 & \textit{for} Lemma~3179 of Volume~X \textit{read} Lemma~3184 of Volume~X \\
\dag\,3063 & proof of Proposition~5.2, l.~3 from foot & \textit{for} Lemma~3199 of Volume~X \textit{read} Lemma~3193 of Volume~X \\
3064 & l.~20 & \textit{delete} the second occurrence of \emph{is} \\
3066 & l.~25 & \textit{see} Corrigendum~S15 \\
3068 & l.~2 from foot & \textit{for} Lemma~3059 of Volume~X \textit{read} Lemma~3064 of Volume~X \\
3069 & l.~3 & \textit{for} Lemma~3086 of Volume~X \textit{read} Lemma~3091 of Volume~X; the result cited is on p.~3882 \\
3069 & l.~11 & \textit{for} $y = x \to \pos(y = x)$ \textit{read} $y = x \to \nec(y = x)$ \\
3069 & l.~16 & \textit{for} Lemma~3275 of Volume~X \textit{read} Lemma~3281 of Volume~X \\
3070 & l.~5 & \textit{for} Mann (1983) \textit{read} Mann (1982) \\
3070 & l.~29 & \textit{for} p.~3145 \textit{read} p.~3045 (Definition~4.13) \\
3071 & l.~3 & \textit{for} p.~3026 \textit{read} p.~3062 (Proposition~5.2) \\
3072 & l.~16 & \textit{for} p.~3008 \textit{read} p.~3018 (Lemma~4.2) \\
3072 & l.~38 & \textit{for} Lemma~3192 of Volume~X \textit{read} Lemma~3187 of Volume~X; the result cited is on p.~3896 \\
\dag\,3073 & proof of Corollary~5.5, (5.1) & \textit{for} $\forall {\chi}\,({\chi}x \leftrightarrow {\chi}y) \leftrightarrow x = y$ \textit{read} $\forall {\chi}\,({\chi}x \leftrightarrow {\chi}y) \to x = y$; the formula is used in its correct form on p.~3075 \\
\dag\,3073 & proof of Corollary~5.5, l.~4 from foot & \textit{for} Lemma~3238 of Volume~X \textit{read} Lemma~3237 of Volume~X; the result cited is on p.~3903 \\
3075 & l.~14 & \textit{for} indiscenrible \textit{read} indiscernible \\
3076 & l.~5 & \textit{for} Lemma~3091 of Volume~X \textit{read} Lemma~3085 of Volume~X \\
3076 & l.~3 from foot & \textit{for} Lemma~3166 of Volume~X \textit{read} Lemma~3170 of Volume~X \\
3077 & display (5.2) & \textit{for} $\mathrm{Op}(x) \wedge \mathrm{Op}(z) \wedge z \neq z \to \bot$ \textit{read} $\mathrm{Op}(x) \wedge \mathrm{Op}(z) \wedge x \neq z \to \bot$; correctly printed in (5.4), p.~3081 \\
3077 & display (5.3) & \textit{for} $\forall y\,\forall z\,(\mathrm{MG}(z) \leftrightarrow z = y)$ \textit{read} $\exists y\,\forall z\,(\mathrm{MG}(z) \leftrightarrow z = y)$; the formula is used in its correct form on p.~3081 \\
3079 & l.~16 & \textit{for} indiscernbiility \textit{read} indiscernibility \\
3080 & l.~6 & \textit{for} disticntness \textit{read} distinctness \\
3080 & l.~27 & \textit{for} Lemma~3131 of Volume~X \textit{read} Lemma~3128 of Volume~X; the result cited is on p.~3887 \\
3080 & l.~28 & \textit{for} Lemma~6.6 of Volume~IV \textit{read} Lemma~6.4 of Volume~IV; the result cited is on p.~1601 \\
3081 & display (5.5) & \textit{for} $\nec\exists {\chi}\,(\mathrm{GM}({\chi}) \to ({\chi}z \leftrightarrow {\chi}x))$ \textit{read} $\nec\forall {\chi}\,(\mathrm{GM}({\chi}) \to ({\chi}z \leftrightarrow {\chi}x))$ \\
3084 & l.~3 from foot & \textit{for} indiscerniiblity \textit{read} indiscernibility \\
3084 & l.~2 from foot & \textit{for} Lemma~3132 of Volume~X \textit{read} Lemma~3137 of Volume~X \\
\dag\,3085 & proof of Proposition~5.11, l.~16 & \textit{for} Lemma~3198 of Volume~X \textit{read} Lemma~3196 of Volume~X \\
\dag\,3085 & proof of Proposition~5.11, l.~19 & \textit{for} Definition~5.15 \textit{read} Definition~5.10 \\
\dag\,3085 & proof of Proposition~5.11, l.~28 & \textit{for} Proposition~5.6 \textit{read} Proposition~5.2; the step uses Proposition~5.2 (p.~3062) \\
\dag\,3085 & proof of Proposition~5.11, l.~29 & \textit{for} Definition~5.12 \textit{read} Definition~5.8 \\
3085 & l.~34 & \textit{for} Lemma~2.10 of Volume~IV \textit{read} Lemma~2.8 of Volume~IV \\
\dag\,3085 & proof of Proposition~5.11, l.~2 from foot & \textit{for} $\mathrm{Op}(x) \wedge \mathrm{Op}(x) \wedge z \neq x \to \bot$ \textit{read} $\mathrm{Op}(z) \wedge \mathrm{Op}(x) \wedge z \neq x \to \bot$ \\
\dag\,3086 & proof of Proposition~5.11, l.~12 & \textit{for} Definition~5.12 \textit{read} Definition~5.9; the step uses Definition~5.9 (p.~3082) \\
\dag\,3086 & proof of Proposition~5.11, l.~24 & \textit{for} $\mathrm{MG}(z) \wedge \mathrm{MG}(x) \leftrightarrow z = x$ \textit{read} $\mathrm{MG}(z) \wedge \mathrm{MG}(x) \to z = x$ \\
\dag\,3086 & proof of Proposition~5.11, l.~35 & \textit{for} $\mathrm{MG}(z) \wedge \mathrm{MG}(y) \leftrightarrow z = y$ \textit{read} $\mathrm{MG}(z) \wedge \mathrm{MG}(y) \to z = y$ \\
3086 & l.~38 & \textit{for} Lemma~3247 of Volume~X \textit{read} Lemma~3252 of Volume~X \\
3087 & l.~2 & \textit{for} Lemma~3179 of Volume~X \textit{read} Lemma~3173 of Volume~X \\
3087 & l.~3 & \textit{for} Lemma~3210 of Volume~X \textit{read} Lemma~3215 of Volume~X \\
3087 & l.~19 & \textit{delete} the second occurrence of \emph{that} \\
3088 & l.~23 & \textit{for} p.~3052 \textit{read} p.~3042 (Corollary~4.10) \\
3088 & l.~29 & \textit{for} p.~3063 \textit{read} p.~3053 (Definition~4.15) \\
3089 & l.~4 & \textit{for} Lemma~3142 of Volume~X \textit{read} Lemma~3145 of Volume~X; the result cited is on p.~3890 \\
3089 & l.~31 & \textit{for} Lemma~3293 of Volume~X \textit{read} Lemma~3292 of Volume~X; the result cited is on p.~3911 \\
3090 & l.~25 & \textit{for} Proposition~5.8 of Volume~III \textit{read} Proposition~5.2 of Volume~III; the result cited is on p.~956 \\
3091 & display (5.6) & \textit{for} $y \neq x \to \pos(y \neq x)$ \textit{read} $y \neq x \to \nec(y \neq x)$ \\
3093 & l.~6 & \textit{for} Lemma~3288 of Volume~X \textit{read} Lemma~3293 of Volume~X \\
3094 & l.~7 & \textit{for} Plantinga (1980) \textit{read} Mann (1982) (the property account of simplicity) \\
3094 & n.~1 & \textit{for} Stump and Kretzmann (1980) \textit{read} Stump and Kretzmann (1981) \\
\dag\,3096 & proof of Proposition~5.13, l.~1 & \textit{for} Definition~5.9 \textit{read} Definition~5.8 \\
\dag\,3096 & proof of Proposition~5.13, l.~6 & \textit{for} Lemma~5.4 \textit{read} Lemma~5.3 \\
\dag\,3096 & proof of Proposition~5.13, l.~10 & \textit{for} Lemma~3110 of Volume~X \textit{read} Lemma~3104 of Volume~X \\
\dag\,3096 & proof of Proposition~5.13, l.~17 & \textit{for} Definition~5.15 of Volume~III \textit{read} Definition~5.10 of Volume~III \\
3096 & l.~23 & \textit{for} $y = y \to \nec(x = y)$ \textit{read} $x = y \to \nec(x = y)$ \\
\dag\,3096 & proof of Proposition~5.13, (5.8) & \textit{for} $\exists {\psi}\,({\psi}y \leftrightarrow {\psi}z) \to y = z$ \textit{read} $\forall {\psi}\,({\psi}y \leftrightarrow {\psi}z) \to y = z$ \\
3096 & l.~27 & \textit{for} Stump and Kretzmann (1981) \textit{read} Mann (1982) (the property account of simplicity) \\
3096 & l.~30 & \textit{for} Lemma~3111 of Volume~X \textit{read} Lemma~3110 of Volume~X; the result cited is on p.~3885 \\
\dag\,3097 & proof of Lemma~5.14, l.~3 & \textit{for} Proposition~5.11 \textit{read} Proposition~5.13 \\
\dag\,3097 & proof of Proposition~5.13, l.~13 & \textit{for} Definition~5.8 \textit{read} Definition~5.10 \\
\dag\,3097 & proof of Lemma~5.14, l.~14 & \textit{for} Proposition~5.11 \textit{read} Proposition~5.6; the step uses Proposition~5.6 (p.~3074) \\
\dag\,3097 & proof of Lemma~5.14, (5.9) & \textit{for} $\exists x\,\forall y\,(\mathrm{MG}(y) \to y = x)$ \textit{read} $\exists x\,\forall y\,(\mathrm{MG}(y) \leftrightarrow y = x)$; correctly printed in (5.11), p.~3100 \\
\dag\,3098 & proof of Lemma~5.14, l.~13 & \textit{for} Definition~5.9 \textit{read} Definition~5.10 \\
3098 & l.~19 & \textit{for} p.~3090 \textit{read} p.~3009 (Proposition~3.15) \\
3099 & l.~7 & \textit{for} p.~2989 \textit{read} p.~2999 (Proposition~3.11) \\
3099 & display (5.10) & \textit{for} $\nec\forall {\phi}\,(\mathrm{GM}({\phi}) \leftrightarrow ({\phi}y \leftrightarrow {\phi}z))$ \textit{read} $\nec\forall {\phi}\,(\mathrm{GM}({\phi}) \to ({\phi}y \leftrightarrow {\phi}z))$; the formula is used in its correct form on p.~3101 \\
3099 & l.~25 & \textit{for} p.~3083 \textit{read} p.~3082 (Definition~5.9) \\
3099 & l.~31 & \textit{for} Lemma~3102 of Volume~X \textit{read} Lemma~3104 of Volume~X; the result cited is on p.~3884 \\
3100 & l.~6 & \textit{for} p.~3170 \textit{read} p.~3070 (Lemma~5.4) \\
3100 & l.~18 & \textit{for} p.~3095 \textit{read} p.~3096 (Definition~5.12) \\
3101 & l.~4 & \textit{for} Lemma~3295 of Volume~X \textit{read} Lemma~3296 of Volume~X \\
3101 & l.~10 & \textit{for} Mann (1982) \textit{read} Stump and Kretzmann (1981) (atemporal eternity) \\
3101 & l.~40 & \textit{for} godlkie \textit{read} godlike \\
\end{errlist}

\subsection*{Chapter~6. Simplicity and the attributes \hfill \normalfont\small pp.~3102--3140}
\begin{errlist}
3103 & l.~22 & \textit{for} $z = y \leftrightarrow \nec(z = y)$ \textit{read} $z = y \to \nec(z = y)$ \\
3103 & l.~32 & \textit{for} Lemma~3291 of Volume~X \textit{read} Lemma~3290 of Volume~X; the result cited is on p.~3910 \\
\dag\,3104 & proof of Proposition~6.2, l.~24 & \textit{for} $\mathrm{Os}(y) \wedge \mathrm{Os}(z) \leftrightarrow \forall q\,(K_{y}q \leftrightarrow K_{z}q)$ \textit{read} $\mathrm{Os}(y) \wedge \mathrm{Os}(z) \to \forall q\,(K_{y}q \leftrightarrow K_{z}q)$ \\
\dag\,3105 & proof of Proposition~6.2, l.~7 & \textit{for} Lemma~3119 of Volume~X \textit{read} Lemma~3121 of Volume~X \\
3105 & l.~10 & \textit{for} haeceity \textit{read} haecceity \\
3107 & l.~6 & \textit{for} $x = y \to \nec(x = x)$ \textit{read} $x = y \to \nec(x = y)$ \\
3108 & l.~2 & \textit{for} $\mathrm{Can}(y,t) \to \nec\,\mathrm{Does}(y,t)$ \textit{read} $\mathrm{Can}(y,t) \to \pos\,\mathrm{Does}(y,t)$ \\
3109 & l.~11 & \textit{for} omnisicence \textit{read} omniscience \\
3110 & l.~16 & \textit{for} Kretzmann (1966) \textit{read} Mann (1982) (the property account of simplicity) \\
3110 & l.~23 & \textit{for} Lemma~3054 of Volume~X \textit{read} Lemma~3050 of Volume~X; the result cited is on p.~3876 \\
3110 & l.~31 & \textit{for} p.~3071 \textit{read} p.~3070 (Lemma~5.3) \\
3111 & l.~5 & \textit{for} Definition~5.5 of Volume~III \textit{read} Definition~5.1 of Volume~III \\
3112 & l.~24 & \textit{for} Lemma~3160 of Volume~X \textit{read} Lemma~3161 of Volume~X; the result cited is on p.~3892 \\
3112 & l.~29 & \textit{for} $\mathrm{Op}(x) \leftrightarrow \forall t\,(\mathrm{Task}(t) \wedge \neg Q(t,x) \leftrightarrow \mathrm{Can}(x,t))$ \textit{read} $\mathrm{Op}(x) \leftrightarrow \forall t\,(\mathrm{Task}(t) \wedge \neg Q(t,x) \to \mathrm{Can}(x,t))$ \\
3114 & l.~14 & \textit{for} Lemma~3178 of Volume~X \textit{read} Lemma~3172 of Volume~X; the result cited is on p.~3894 \\
3114 & l.~32 & \textit{for} p.~2965 \textit{read} p.~2956 (Lemma~2.13) \\
3114 & l.~8 from foot & \textit{for} p.~3174 \textit{read} p.~3074 (Proposition~5.6) \\
3115 & display (6.1) & \textit{for} $\forall {\chi}\,({\chi}z \leftrightarrow {\chi}y) \leftrightarrow z = y$ \textit{read} $\forall {\chi}\,({\chi}z \leftrightarrow {\chi}y) \to z = y$; the formula is used in its correct form on p.~3118 \\
3115 & display (6.2) & \textit{for} $K_{y}q \leftrightarrow K_{y}K_{y}q$ \textit{read} $K_{y}q \to K_{y}K_{y}q$ \\
3116 & l.~8 & \textit{for} indiscerinble \textit{read} indiscernible \\
3116 & l.~16 & \textit{for} $z \neq x \to \nec(x \neq x)$ \textit{read} $z \neq x \to \nec(z \neq x)$ \\
3117 & l.~7 & \textit{for} Lemma~3192 of Volume~X \textit{read} Lemma~3196 of Volume~X \\
3117 & l.~15 & \textit{for} p.~3079 \textit{read} p.~3097 (Lemma~5.14) \\
3117 & l.~26 & \textit{for} $\mathrm{MG}(x) \wedge \mathrm{MG}(x) \to z = x$ \textit{read} $\mathrm{MG}(z) \wedge \mathrm{MG}(x) \to z = x$ \\
\dag\,3119 & proof of Proposition~6.6, l.~5 & \textit{for} Proposition~6.7 \textit{read} Proposition~6.2 \\
\dag\,3119 & proof of Proposition~6.8, l.~9 & \textit{for} Proposition~6.9 \textit{read} Proposition~6.7 \\
3119 & l.~20 & \textit{for} Lemma~3146 of Volume~X \textit{read} Lemma~3150 of Volume~X \\
\dag\,3119 & proof of Proposition~6.8, l.~26 & \textit{for} Definition~6.4 \textit{read} Definition~6.1 \\
3119 & l.~29 & \textit{for} p.~3130 \textit{read} p.~3103 (Definition~6.1) \\
3119 & l.~6 from foot & \textit{for} compossibliity \textit{read} compossibility \\
\dag\,3119 & proof of Proposition~6.8, l.~1 from foot & \textit{for} Lemma~3293 of Volume~X \textit{read} Lemma~3295 of Volume~X; the result cited is on p.~3911 \\
\dag\,3120 & proof of Proposition~6.8, l.~4 & \textit{for} Proposition~6.6 \textit{read} Proposition~6.2; the step uses Proposition~6.2 (p.~3104) \\
\dag\,3120 & proof of Proposition~6.7, l.~19 & \textit{for} Proposition~6.10 \textit{read} Proposition~6.6 \\
3121 & l.~15 & \textit{for} Definition~6.7 of Volume~IV \textit{read} Definition~6.2 of Volume~IV; the result cited is on p.~1589 \\
3124 & l.~24 & \textit{for} Mann (1982) \textit{read} Bergmann (2001) (sceptical theism) \\
3124 & l.~28 & \textit{for} Lemma~3147 of Volume~X \textit{read} Lemma~3152 of Volume~X \\
3124 & l.~29 & \textit{for} Krtezmann (1966) \textit{read} Kretzmann (1966) \\
3126 & display (6.3) & \textit{for} $K_{x}q \leftrightarrow K_{x}K_{x}q$ \textit{read} $K_{x}q \to K_{x}K_{x}q$ \\
\dag\,3127 & proof of Proposition~6.9, l.~30 & \textit{for} Proposition~6.8 \textit{read} Proposition~6.6 \\
3129 & l.~19 & \textit{for} p.~3128 \textit{read} p.~3127 (Proposition~6.9) \\
3130 & l.~27 & \textit{for} Lemma~3223 of Volume~X \textit{read} Lemma~3222 of Volume~X \\
3131 & l.~21 & \textit{for} $\mathrm{Os}(x) \to \forall r\,(r \to K_{x}r)$ \textit{read} $\mathrm{Os}(x) \leftrightarrow \forall r\,(r \to K_{x}r)$ \\
3131 & display (6.4) & \textit{for} $\mathrm{Can}(x,t) \leftrightarrow \pos\,\mathrm{Does}(x,t)$ \textit{read} $\mathrm{Can}(x,t) \to \pos\,\mathrm{Does}(x,t)$; the formula is used in its correct form on p.~3133 \\
\dag\,3132 & proof of Proposition~6.10, l.~8 & \textit{for} Definition~6.3 \textit{read} Definition~6.4 \\
\dag\,3132 & proof of Proposition~6.10, l.~13 & \textit{for} Proposition~6.6 \textit{read} Proposition~6.7; the step uses Proposition~6.7 (p.~3119) \\
\dag\,3132 & proof of Proposition~6.10, (6.5) & \textit{for} $\mathrm{MP}(x) \to \neg\nec\,\mathrm{Wrong}(x,a)$ \textit{read} $\mathrm{MP}(x) \to \neg\pos\,\mathrm{Wrong}(x,a)$ \\
\dag\,3133 & proof of Proposition~6.10, (6.6) & \textit{for} $\nec\forall {\chi}\,(\mathrm{GM}({\chi}) \leftrightarrow ({\chi}z \leftrightarrow {\chi}y))$ \textit{read} $\nec\forall {\chi}\,(\mathrm{GM}({\chi}) \to ({\chi}z \leftrightarrow {\chi}y))$ \\
\dag\,3133 & proof of Proposition~6.10, l.~10 & \textit{for} Definition~6.3 \textit{read} Definition~6.1 \\
3133 & l.~20 & \textit{for} $\mathrm{Op}(z) \to \forall t\,(\mathrm{Task}(t) \wedge \neg Q(t,z) \to \mathrm{Can}(z,t))$ \textit{read} $\mathrm{Op}(z) \leftrightarrow \forall t\,(\mathrm{Task}(t) \wedge \neg Q(t,z) \to \mathrm{Can}(z,t))$ \\
\dag\,3133 & proof of Proposition~6.10, (6.7) & \textit{for} $\neg K_{y}p \to K_{y} K_{y}p$ \textit{read} $\neg K_{y}p \to K_{y}\neg K_{y}p$ \\
3133 & l.~35 & \textit{for} $ Q(t,x) \wedge \mathrm{Task}(t) \to \mathrm{Can}(x,t)$ \textit{read} $\neg Q(t,x) \wedge \mathrm{Task}(t) \to \mathrm{Can}(x,t)$ \\
3133 & l.~3 from foot & \textit{for} $\mathrm{NE}(z) \leftrightarrow \neg\exists x\,D(z,x)$ \textit{read} $\mathrm{NE}(z) \to \neg\exists x\,D(z,x)$ \\
3134 & display (6.8) & \textit{for} $\mathrm{Os}(z) \wedge \mathrm{Os}(x) \to \forall q\,(K_{z}q \leftrightarrow K_{z}q)$ \textit{read} $\mathrm{Os}(z) \wedge \mathrm{Os}(x) \to \forall q\,(K_{z}q \leftrightarrow K_{x}q)$ \\
3135 & display (6.9) & \textit{for} $\mathrm{Os}(z) \leftrightarrow \exists p\,(p \to K_{z}p)$ \textit{read} $\mathrm{Os}(z) \leftrightarrow \forall p\,(p \to K_{z}p)$; the formula is used in its correct form on p.~3138 \\
3135 & l.~26 & \textit{for} Lemma~3217 of Volume~X \textit{read} Lemma~3211 of Volume~X; the result cited is on p.~3899 \\
3137 & l.~8 & \textit{for} Bergmann (2001) \textit{read} Mann (1982) (the property account of simplicity) \\
3137 & l.~21 & \textit{for} $\pos\forall {\phi}\,(\mathrm{GM}({\phi}) \to ({\phi}x \leftrightarrow {\phi}z))$ \textit{read} $\nec\forall {\phi}\,(\mathrm{GM}({\phi}) \to ({\phi}x \leftrightarrow {\phi}z))$ \\
\dag\,3138 & proof of Lemma~6.11, (6.10) & \textit{for} $\forall {\phi}\,({\phi}x \leftrightarrow {\phi}x) \to x = z$ \textit{read} $\forall {\phi}\,({\phi}x \leftrightarrow {\phi}z) \to x = z$; the formula is used in its correct form on p.~3140 \\
3138 & l.~1 from foot & \textit{for} Lemma~3267 of Volume~X \textit{read} Lemma~3270 of Volume~X \\
\dag\,3139 & proof of Proposition~6.13, l.~11 & \textit{for} Proposition~6.6 \textit{read} Proposition~6.9; the step uses Proposition~6.9 (p.~3127) \\
\dag\,3139 & proof of Theorem~6.12, l.~13 & \textit{for} Proposition~6.13 \textit{read} Proposition~6.8 \\
\dag\,3139 & proof of Proposition~6.13, l.~19 & \textit{for} Proposition~6.8 \textit{read} Proposition~6.10; the step uses Proposition~6.10 (p.~3132) \\
\dag\,3139 & proof of Proposition~6.13, l.~22 & \textit{for} $\neg Q(t,y) \wedge \mathrm{Task}(t) \leftrightarrow \mathrm{Can}(y,t)$ \textit{read} $\neg Q(t,y) \wedge \mathrm{Task}(t) \to \mathrm{Can}(y,t)$ \\
3139 & l.~23 & \textit{for} Plantinga (1980) \textit{read} Bergmann (2001) (sceptical theism) \\
\dag\,3139 & proof of Proposition~6.13, l.~5 from foot & \textit{for} Lemma~3161 of Volume~X \textit{read} Lemma~3156 of Volume~X; the result cited is on p.~3891 \\
3140 & l.~6 & \textit{for} indiscrenibility \textit{read} indiscernibility \\
\dag\,3140 & proof of Theorem~6.12, l.~19 & \textit{for} Proposition~6.7 \textit{read} Proposition~6.10 \\
3140 & l.~26 & \textit{for} Plantinga (1980) \textit{read} Kretzmann (1966) (omniscience and change) \\
\end{errlist}

\subsection*{Chapter~7. The objection from modal collapse to simplicity \hfill \normalfont\small pp.~3141--3188}
\begin{errlist}
3143 & display (7.2) & \textit{for} $z \neq y \to \nec(z \neq z)$ \textit{read} $z \neq y \to \nec(z \neq y)$ \\
3143 & l.~2 from foot & \textit{for} Lemma~3144 of Volume~X \textit{read} Lemma~3145 of Volume~X \\
3144 & l.~8 & \textit{for} Lemma~2.10 of Volume~IV \textit{read} Lemma~2.4 of Volume~IV \\
3145 & l.~16 & \textit{for} Lemma~3153 of Volume~X \textit{read} Lemma~3149 of Volume~X \\
3146 & l.~10 & \textit{for} Lemma~3221 of Volume~X \textit{read} Lemma~3219 of Volume~X; the result cited is on p.~3900 \\
3146 & display (7.3) & \textit{for} $\forall {\psi}\,({\psi}y \leftrightarrow {\psi}y) \to y = x$ \textit{read} $\forall {\psi}\,({\psi}y \leftrightarrow {\psi}x) \to y = x$; correctly printed in (7.4), p.~3148 \\
3147 & l.~2 from foot & \textit{for} Lemma~3077 of Volume~X \textit{read} Lemma~3083 of Volume~X; the result cited is on p.~3881 \\
3148 & display (7.5) & \textit{for} $w'Rv' \wedge v'Rw_2 \to w_2Rw'$ \textit{read} $w'Rv' \wedge v'Rw_2 \to w'Rw_2$ \\
3148 & l.~29 & \textit{for} Lemma~3096 of Volume~X \textit{read} Lemma~3102 of Volume~X \\
3149 & l.~9 & \textit{for} Lemma~3147 of Volume~X \textit{read} Lemma~3141 of Volume~X \\
3149 & display (7.7) & \textit{for} $w_1Rv \wedge uRv \to w_1Ru$ \textit{read} $w_1Rv \wedge vRu \to w_1Ru$ \\
3149 & l.~14 & \textit{for} Lemma~6.16 of Volume~IV \textit{read} Lemma~6.12 of Volume~IV \\
3150 & l.~2 & \textit{for} characteirstic \textit{read} characteristic \\
3150 & l.~3 & \textit{for} Lemma~3158 of Volume~X \textit{read} Lemma~3164 of Volume~X \\
3150 & l.~7 & \textit{for} $x = z \leftrightarrow \nec(x = z)$ \textit{read} $x = z \to \nec(x = z)$ \\
3150 & l.~14 & \textit{for} Lemma~3137 of Volume~X \textit{read} Lemma~3141 of Volume~X \\
3150 & l.~15 & \textit{for} Proposition~5.8 of Volume~III \textit{read} Proposition~5.3 of Volume~III \\
3151 & l.~1 & \textit{for} reflxivity \textit{read} reflexivity \\
3152 & l.~32 & \textit{for} Lemma~3209 of Volume~X \textit{read} Lemma~3211 of Volume~X \\
3153 & l.~11 & \textit{for} p.~3123 \textit{read} p.~3132 (Proposition~6.10) \\
3153 & l.~12 & \textit{for} Lemma~3269 of Volume~X \textit{read} Lemma~3267 of Volume~X \\
3153 & l.~33 & \textit{for} $\nec(q \wedge r) \leftrightarrow (\nec r \wedge \nec r)$ \textit{read} $\nec(q \wedge r) \leftrightarrow (\nec q \wedge \nec r)$ \\
3154 & l.~34 & \textit{for} p.~3019 \textit{read} p.~3119 (Proposition~6.8) \\
3155 & l.~3 & \textit{for} Lemma~3200 of Volume~X \textit{read} Lemma~3194 of Volume~X; the result cited is on p.~3897 \\
3155 & display (7.9) & \textit{for} $\pos\forall {\psi}\,(\mathrm{GM}({\psi}) \to ({\psi}y \leftrightarrow {\psi}x))$ \textit{read} $\nec\forall {\psi}\,(\mathrm{GM}({\psi}) \to ({\psi}y \leftrightarrow {\psi}x))$; correctly printed in (7.10), p.~3156 \\
3156 & l.~32 & \textit{for} compsoition \textit{read} composition \\
3157 & l.~17 & \textit{for} $v' \models \nec p \iff \exists v\,(v'Rv \to v \models p)$ \textit{read} $v' \models \nec p \iff \forall v\,(v'Rv \to v \models p)$ \\
\dag\,3157 & proof of Theorem~7.2, l.~31 & \textit{for} Lemma~3203 of Volume~X \textit{read} Lemma~3200 of Volume~X \\
3157 & l.~3 from foot & \textit{for} tarnsitivity \textit{read} transitivity \\
3158 & l.~28 & \textit{for} heacceity \textit{read} haecceity \\
\dag\,3159 & proof of Lemma~7.6, l.~13 & \textit{for} Definition~7.3 \textit{read} Definition~7.1 \\
\dag\,3160 & proof of Lemma~7.6, l.~16 & \textit{for} $\pos(r \vee p) \leftrightarrow (\pos r \wedge \pos p)$ \textit{read} $\pos(r \vee p) \leftrightarrow (\pos r \vee \pos p)$ \\
3161 & l.~5 & \textit{for} Lemma~3121 of Volume~X \textit{read} Lemma~3125 of Volume~X \\
3161 & l.~25 & \textit{for} Lemma~3184 of Volume~X \textit{read} Lemma~3186 of Volume~X; the result cited is on p.~3896 \\
3161 & display (7.11) & \textit{for} $z \neq x \leftrightarrow \nec(z \neq x)$ \textit{read} $z \neq x \to \nec(z \neq x)$; correctly printed in (7.13), p.~3163 \\
3162 & l.~21 & \textit{after} true at every world \textit{insert} accessible from $w_0$ \\
3164 & l.~24 & \textit{for} $(w'Rv \vee w'Rw_2) \to vRw_2$ \textit{read} $(w'Rv \wedge w'Rw_2) \to vRw_2$ \\
3164 & l.~35 & \textit{for} Lemma~3279 of Volume~X \textit{read} Lemma~3282 of Volume~X; the result cited is on p.~3909 \\
3166 & l.~4 & \textit{for} compositoin \textit{read} composition \\
3166 & l.~19 & \textit{for} Theorem~5.6 of Volume~III \textit{read} Theorem~5.4 of Volume~III \\
3167 & l.~6 from foot & \textit{for} p.~3158 \textit{read} p.~3159 (Corollary~7.5) \\
3168 & l.~9 & \textit{for} Proposition~2.2 of Volume~IV \textit{read} Proposition~2.5 of Volume~IV; the result cited is on p.~1425 \\
3168 & l.~16 & \textit{for} $\mathrm{Op}(z) \vee \mathrm{Op}(y) \wedge z \neq y \to \bot$ \textit{read} $\mathrm{Op}(z) \wedge \mathrm{Op}(y) \wedge z \neq y \to \bot$ \\
3168 & display (7.14) & \textit{for} $\neg\pos r \leftrightarrow \pos\neg r$ \textit{read} $\neg\pos r \leftrightarrow \nec\neg r$ \\
3168 & l.~7 from foot & \textit{for} p.~3108 \textit{read} p.~3118 (Lemma~6.5) \\
3168 & l.~6 from foot & \textit{for} Lemma~3246 of Volume~X \textit{read} Lemma~3244 of Volume~X; the result cited is on p.~3904 \\
3171 & display (7.15) & \textit{for} $y = z \leftrightarrow \nec(y = z)$ \textit{read} $y = z \to \nec(y = z)$ \\
\dag\,3172 & proof of Theorem~7.13, l.~5 & \textit{for} Corollary~7.5 \textit{read} Corollary~7.4 \\
\dag\,3172 & proof of Theorem~7.13, l.~23 & \textit{for} Lemma~3275 of Volume~X \textit{read} Lemma~3279 of Volume~X \\
\dag\,3172 & proof of Theorem~7.13, (7.18) & \textit{for} $(w_0Rw_1 \wedge w_0Rw') \to w_2Rw'$ \textit{read} $(w_0Rw_2 \wedge w_0Rw') \to w_2Rw'$ \\
3173 & l.~32 & \textit{for} p.~3193 \textit{read} p.~3139 (Theorem~6.12) \\
3174 & l.~21 & \textit{for} Proposition~5.8 of Volume~III \textit{read} Proposition~5.11 of Volume~III \\
3174 & l.~36 & \textit{for} Lemma~3171 of Volume~X \textit{read} Lemma~3173 of Volume~X; the result cited is on p.~3894 \\
3175 & display (7.21) & \textit{for} $\pos\nec q \to \pos q$ \textit{read} $\pos\nec q \to \nec q$; correctly printed in (7.24), p.~3180 \\
3175 & display (7.20) & \textit{for} $\forall {\phi}\,({\phi}z \leftrightarrow {\phi}y) \leftrightarrow z = y$ \textit{read} $\forall {\phi}\,({\phi}z \leftrightarrow {\phi}y) \to z = y$; the formula is used in its correct form on p.~3177 \\
3175 & l.~2 from foot & \textit{for} Lemma~3231 of Volume~X \textit{read} Lemma~3234 of Volume~X \\
3176 & display (7.22) & \textit{for} $\nec\exists {\psi}\,(\mathrm{GM}({\psi}) \to ({\psi}x \leftrightarrow {\psi}z))$ \textit{read} $\nec\forall {\psi}\,(\mathrm{GM}({\psi}) \to ({\psi}x \leftrightarrow {\psi}z))$; correctly printed in (7.25), p.~3180 \\
3176 & l.~18 & \textit{for} $\pos(q \to p) \to (\nec q \to \nec p)$ \textit{read} $\nec(q \to p) \to (\nec q \to \nec p)$ \\
\dag\,3177 & proof of Lemma~7.14, l.~27 & \textit{for} Definition~7.9 \textit{read} Definition~7.10 \\
3178 & l.~8 & \textit{for} p.~3165 \textit{read} p.~3156 (Theorem~7.2) \\
3178 & l.~32 & \textit{for} p.~3138 \textit{read} p.~3137 (Lemma~6.11) \\
3179 & l.~8 & \textit{for} $\exists y\,\forall z\,(\mathrm{MG}(z) \to z = y)$ \textit{read} $\exists y\,\forall z\,(\mathrm{MG}(z) \leftrightarrow z = y)$ \\
3179 & l.~13 & \textit{for} Lemma~3086 of Volume~X \textit{read} Lemma~3083 of Volume~X \\
3179 & l.~29 & \textit{for} Definition~5.14 of Volume~III \textit{read} Definition~5.20 of Volume~III \\
3179 & l.~30 & \textit{for} Lemma~3063 of Volume~X \textit{read} Lemma~3066 of Volume~X; the result cited is on p.~3878 \\
3180 & l.~16 & \textit{for} $\nec\pos q \to \pos q$ \textit{read} $\pos\pos q \to \pos q$ \\
3180 & l.~30 & \textit{for} $\nec p \wedge \pos p \to \pos(r \wedge p)$ \textit{read} $\nec r \wedge \pos p \to \pos(r \wedge p)$ \\
3181 & l.~5 & \textit{for} p.~3172 \textit{read} p.~3072 (Corollary~5.5) \\
\dag\,3181 & proof of Theorem~7.15, (7.26) & \textit{for} $\nec r \to \pos\nec r$ \textit{read} $\nec r \to \nec\nec r$ \\
3181 & l.~31 & \textit{for} p.~3001 \textit{read} p.~3000 (Proposition~3.12) \\
\dag\,3182 & proof of Theorem~7.15, l.~4 & \textit{for} Lemma~5.19 of Volume~III \textit{read} Lemma~5.21 of Volume~III; the result cited is on p.~1010 \\
3182 & l.~22 & \textit{for} p.~3109 \textit{read} p.~3108 (Definition~6.3) \\
\dag\,3182 & proof of Theorem~7.15, l.~23 & \textit{for} Lemma~3234 of Volume~X \textit{read} Lemma~3237 of Volume~X \\
3182 & l.~29 & \textit{for} Lemma~2.8 of Volume~IV \textit{read} Lemma~2.12 of Volume~IV; the result cited is on p.~1442 \\
3183 & l.~2 & \textit{for} Lemma~3289 of Volume~X \textit{read} Lemma~3294 of Volume~X \\
3183 & display (7.27) & \textit{for} $u \models \nec q \iff \forall u\,(w'Ru \to u \models q)$ \textit{read} $w' \models \nec q \iff \forall u\,(w'Ru \to u \models q)$ \\
3184 & display (7.29) & \textit{for} $\pos r \to \pos\pos r$ \textit{read} $\pos r \to \nec\pos r$ \\
3184 & l.~30 & \textit{for} p.~3087 \textit{read} p.~3097 (Definition~5.15) \\
3184 & l.~4 from foot & \textit{for} semmantics \textit{read} semantics \\
3184 & display (7.28) & \textit{for} $w_1 \models \pos r \iff \exists w_0\,(w_0Rw_1 \wedge w_0 \models r)$ \textit{read} $w_1 \models \pos r \iff \exists w_0\,(w_1Rw_0 \wedge w_0 \models r)$ \\
3185 & l.~37 & \textit{for} compelteness \textit{read} completeness \\
3186 & l.~8 & \textit{for} p.~3171 \textit{read} p.~3172 (Theorem~7.13) \\
3187 & l.~4 & \textit{for} Lemma~3.3 of Volume~I \textit{read} Lemma~3.2 of Volume~I \\
3187 & l.~13 & \textit{for} Lemma~3176 of Volume~X \textit{read} Lemma~3181 of Volume~X; the result cited is on p.~3895 \\
\end{errlist}

\subsection*{Chapter~8. Summary \hfill \normalfont\small pp.~3189--3203}
\begin{errlist}
3190 & l.~37 & \textit{for} p.~3061 \textit{read} p.~3062 (Definition~5.1) \\
3190 & l.~3 from foot & \textit{for} p.~3083 \textit{read} p.~3084 (Definition~5.10) \\
3191 & l.~15 & \textit{for} Lemma~3204 of Volume~X \textit{read} Lemma~3198 of Volume~X \\
3191 & l.~28 & \textit{for} Lemma~3254 of Volume~X \textit{read} Lemma~3248 of Volume~X \\
3191 & l.~34 & \textit{for} Lemma~3179 of Volume~X \textit{read} Lemma~3174 of Volume~X \\
3191 & l.~39 & \textit{for} Lemma~3094 of Volume~X \textit{read} Lemma~3091 of Volume~X \\
3193 & l.~19 & \textit{for} p.~3118 \textit{read} p.~3119 (Proposition~6.7) \\
3194 & l.~18 & \textit{for} fololwing \textit{read} following \\
3194 & l.~28 & \textit{for} Lemma~3193 of Volume~X \textit{read} Lemma~3195 of Volume~X; the result cited is on p.~3897 \\
3194 & l.~1 from foot & \textit{for} Lemma~3109 of Volume~X \textit{read} Lemma~3112 of Volume~X \\
3195 & l.~12 & \textit{for} Lemma~3116 of Volume~X \textit{read} Lemma~3119 of Volume~X \\
3195 & l.~14 & \textit{for} Definition~5.18 of Volume~III \textit{read} Definition~5.14 of Volume~III \\
3195 & l.~31 & \textit{for} p.~3187 \textit{read} p.~3186 (Definition~7.16) \\
3196 & l.~29 & \textit{for} Lemma~3218 of Volume~X \textit{read} Lemma~3223 of Volume~X \\
3196 & l.~40 & \textit{for} Lemma~3176 of Volume~X \textit{read} Lemma~3180 of Volume~X \\
3197 & l.~17 & \textit{for} Lemma~3143 of Volume~X \textit{read} Lemma~3141 of Volume~X \\
3197 & l.~28 & \textit{for} Lemma~3268 of Volume~X \textit{read} Lemma~3272 of Volume~X; the result cited is on p.~3908 \\
3197 & l.~5 from foot & \textit{for} Lemma~5.21 of Volume~III \textit{read} Lemma~5.19 of Volume~III \\
3198 & l.~10 & \textit{for} Lemma~3174 of Volume~X \textit{read} Lemma~3178 of Volume~X \\
3198 & l.~35 & \textit{for} p.~3057 \textit{read} p.~3157 (Definition~7.3) \\
3199 & l.~15 & \textit{for} Lemma~3227 of Volume~X \textit{read} Lemma~3222 of Volume~X \\
\dag\,3200 & proof of Proposition~8.2, l.~18 & \textit{for} Lemma~3292 of Volume~X \textit{read} Lemma~3293 of Volume~X \\
\dag\,3200 & proof of Proposition~8.2, l.~26 & \textit{for} Lemma~3197 of Volume~X \textit{read} Lemma~3196 of Volume~X; the result cited is on p.~3897 \\
3202 & l.~5 & \textit{for} Lemma~3151 of Volume~X \textit{read} Lemma~3149 of Volume~X; the result cited is on p.~3890 \\
3202 & l.~6 & \textit{for} p.~2986 \textit{read} p.~2996 (Theorem~3.8) \\
3202 & l.~13 & \textit{for} Lemma~3092 of Volume~X \textit{read} Lemma~3098 of Volume~X; the result cited is on p.~3883 \\
\end{errlist}


\section{Volume IX: Divine Agency}\label{s:vol9}
% generated by generate.py -- do not edit
\subsection*{Front matter \hfill \normalfont\small pp.~3204--3206}
\begin{errlist}
3205 & Contents, Chapter~3 & \textit{for} 3261 \textit{read} 3251 \\
\end{errlist}

\subsection*{Chapter~1. Introduction \hfill \normalfont\small pp.~3207--3213}
\begin{errlist}
3208 & running heads, pp.~3208--3212 (even) & \textit{for} The main result and action \textit{read} Introduction \\
3208 & l.~15 & \textit{for} Lemma~3067 of Volume~X \textit{read} Lemma~3066 of Volume~X \\
3208 & l.~17 & \textit{for} Lemma~3273 of Volume~X \textit{read} Lemma~3276 of Volume~X; the result cited is on p.~3908 \\
3208 & l.~20 & \textit{delete} the second occurrence of \emph{that} \\
3208 & l.~4 from foot & \textit{delete} the second occurrence of \emph{is} \\
3209 & l.~11 & \textit{delete} the second occurrence of \emph{of} \\
3209 & n.~14 & \textit{for} Kenny (1980) \textit{read} Kenny (1979) \\
3210 & n.~7 & \textit{for} Keny (1979) \textit{read} Kenny (1979) \\
\dag\,3211 & proof of Proposition~1.2, l.~13 & \textit{for} Lemma~3082 of Volume~X \textit{read} Lemma~3080 of Volume~X; the result cited is on p.~3880 \\
3213 & l.~29 & \textit{delete} the second occurrence of \emph{that} \\
3213 & l.~36 & \textit{delete} the second occurrence of \emph{a} \\
3213 & l.~39 & \textit{delete} the second occurrence of \emph{is} \\
3213 & n.~15 & \textit{for} Kenny (1978) \textit{read} Kenny (1979) \\
\end{errlist}

\subsection*{Chapter~2. What the main result says about action \hfill \normalfont\small pp.~3214--3250}
\begin{errlist}
3215 & l.~11 & \textit{for} cuasal \textit{read} causal \\
3216 & display (2.2) & \textit{for} $\mathrm{MG}(x,w_0) \wedge w_0Rw \leftrightarrow \mathrm{MG}(x,w)$ \textit{read} $\mathrm{MG}(x,w_0) \wedge w_0Rw \to \mathrm{MG}(x,w)$; correctly printed in (2.4), p.~3221 \\
3217 & l.~14 & \textit{delete} the second occurrence of \emph{is} \\
3219 & l.~20 & \textit{delete} the second occurrence of \emph{of} \\
3219 & l.~36 & \textit{for} Lemma~3158 of Volume~X \textit{read} Lemma~3156 of Volume~X \\
3220 & l.~8 & \textit{for} omnisciant \textit{read} omniscient \\
3221 & l.~8 & \textit{for} Lemma~3065 of Volume~X \textit{read} Lemma~3060 of Volume~X \\
3222 & l.~6 & \textit{for} $\exists v'\,\exists w'\,V(w') > V(v')$ \textit{read} $\forall v'\,\exists w'\,V(w') > V(v')$ \\
\dag\,3222 & proof of Lemma~2.3, l.~26 & \textit{for} $\mathrm{MG}(x,v) \wedge vRu \to \mathrm{MG}(x,v)$ \textit{read} $\mathrm{MG}(x,v) \wedge vRu \to \mathrm{MG}(x,u)$ \\
\dag\,3223 & proof of Lemma~2.3, l.~3 & \textit{for} Definition~2.4 \textit{read} Definition~2.1 \\
3223 & l.~28 & \textit{after} exists \textit{insert} at every world accessible from $w_0$ \\
3224 & l.~17 & \textit{for} $\mathrm{ME}(x,w_0) \iff \mathrm{Op}(x,w_0) \wedge \mathrm{Os}(x,w_0) \vee \mathrm{MP}(x,w_0)$ \textit{read} $\mathrm{ME}(x,w_0) \iff \mathrm{Op}(x,w_0) \wedge \mathrm{Os}(x,w_0) \wedge \mathrm{MP}(x,w_0)$ \\
3226 & l.~1 & \textit{for} instantiatted \textit{read} instantiated \\
3226 & l.~17 & \textit{after} maximally excellent \textit{insert} at $w_0$ \\
3226 & l.~20 & \textit{for} Lemma~3119 of Volume~X \textit{read} Lemma~3117 of Volume~X \\
3226 & display (2.6) & \textit{for} $\mathrm{NME}(x,w) \iff \exists w_2\,(wRw_2 \to \mathrm{ME}(x,w_2))$ \textit{read} $\mathrm{NME}(x,w) \iff \forall w_2\,(wRw_2 \to \mathrm{ME}(x,w_2))$ \\
3229 & l.~5 & \textit{for} the main result fixes the causal role \textit{read} the main result leaves the causal role undetermined \\
3229 & l.~7 & \textit{delete} the second occurrence of \emph{the} \\
3229 & l.~12 & \textit{for} $\mathrm{Causes}(g,e) \leftrightarrow \mathrm{Wills}(g,e)$ \textit{read} $\mathrm{Causes}(g,e) \to \mathrm{Wills}(g,e)$ \\
3230 & display (2.9) & \textit{for} $\exists x\,\mathrm{MG}(x,w_0) \to \forall x\,(E(x,w_0) \wedge \mathrm{ME}(x,w_0))$ \textit{read} $\exists x\,\mathrm{MG}(x,w_0) \to \exists x\,(E(x,w_0) \wedge \mathrm{ME}(x,w_0))$; correctly printed in (2.10), p.~3231 \\
3230 & l.~28 & \textit{for} moraly \textit{read} morally \\
3231 & l.~2 & \textit{for} p.~3211 \textit{read} p.~3212 (Proposition~1.3) \\
3231 & l.~20 & \textit{for} $\pos\neg\exists y\,\mathrm{MG}(y) \to \neg\forall y\,\mathrm{MG}(y)$ \textit{read} $\pos\neg\exists y\,\mathrm{MG}(y) \to \neg\exists y\,\mathrm{MG}(y)$ \\
3232 & l.~1 & \textit{for} $\mathrm{ME}(y,v') \iff \mathrm{Op}(y,v') \wedge \mathrm{Os}(y,v') \vee \mathrm{MP}(y,v')$ \textit{read} $\mathrm{ME}(y,v') \iff \mathrm{Op}(y,v') \wedge \mathrm{Os}(y,v') \wedge \mathrm{MP}(y,v')$ \\
3232 & l.~10 & \textit{for} choices depend counterfactually on foreknowledge \textit{read} foreknowledge depends counterfactually on free choices \\
\dag\,3232 & proof of Proposition~2.6, l.~12 & \textit{for} Definition~2.7 \textit{read} Definition~2.2 \\
\dag\,3233 & proof of Proposition~2.6, l.~10 & \textit{for} Lemma~3074 of Volume~X \textit{read} Lemma~3078 of Volume~X \\
\dag\,3233 & proof of Proposition~2.6, l.~13 & \textit{for} Lemma~3176 of Volume~X \textit{read} Lemma~3182 of Volume~X \\
\dag\,3233 & proof of Proposition~2.6, (2.11) & \textit{for} $\pos\,\mathrm{Creates}(g,w) \wedge \pos\,\mathrm{Creates}(g,w)$ \textit{read} $\pos\,\mathrm{Creates}(g,w) \wedge \pos\,\neg\mathrm{Creates}(g,w)$; correctly printed in (2.13), p.~3237 \\
\dag\,3233 & proof of Proposition~2.6, l.~23 & \textit{for} Definition~2.8 \textit{read} Definition~2.4; the step uses Definition~2.4 (p.~3224) \\
3234 & l.~7 & \textit{for} Lemma~3191 of Volume~X \textit{read} Lemma~3197 of Volume~X \\
3234 & l.~31 & \textit{after} causal role \textit{insert} of the maximally great being \\
3234 & l.~8 from foot & \textit{for} p.~3212 \textit{read} p.~3222 (Lemma~2.3) \\
3234 & l.~4 from foot & \textit{for} p.~3212 \textit{read} p.~3221 (Definition~2.2) \\
3236 & l.~16 & \textit{for} Lemma~3286 of Volume~X \textit{read} Lemma~3282 of Volume~X \\
3237 & l.~7 & \textit{for} Lemma~3279 of Volume~X \textit{read} Lemma~3274 of Volume~X; the result cited is on p.~3908 \\
3238 & l.~11 & \textit{delete} the second occurrence of \emph{a} \\
3239 & l.~1 & \textit{for} p.~3225 \textit{read} p.~3224 (Definition~2.4) \\
3239 & display (2.14) & \textit{for} $\forall u\,\exists w_0\,V(w_0) < V(u)$ \textit{read} $\forall u\,\exists w_0\,V(w_0) > V(u)$ \\
3239 & l.~26 & \textit{for} Lemma~6.4 of Volume~IV \textit{read} Lemma~6.6 of Volume~IV; the result cited is on p.~1609 \\
3240 & l.~39 & \textit{delete} the second occurrence of \emph{of} \\
3241 & display (2.15) & \textit{for} $(C \mathrel{\Box\!\!\to} K_g p) \wedge (K_g p \mathrel{\Box\!\!\to} C)$ \textit{read} $(C \mathrel{\Box\!\!\to} K_g p) \wedge \neg(K_g p \mathrel{\Box\!\!\to} C)$ \\
3242 & l.~26 & \textit{for} in time \textit{read} timeless \\
3242 & l.~30 & \textit{for} p.~3211 \textit{read} p.~3210 (Proposition~1.2) \\
3242 & l.~32 & \textit{for} Lemma~3070 of Volume~X \textit{read} Lemma~3064 of Volume~X \\
3243 & display (2.19) & \textit{for} $\mathcal{M}_1 \models \mathrm{Causes}(g,e) \wedge \mathcal{M}_1 \models \neg\mathrm{Causes}(g,e)$ \textit{read} $\mathcal{M}_1 \models \mathrm{Causes}(g,e) \wedge \mathcal{M}_2 \models \neg\mathrm{Causes}(g,e)$ \\
3243 & l.~16 & \textit{for} Lemma~3086 of Volume~X \textit{read} Lemma~3081 of Volume~X \\
3243 & display (2.18) & \textit{for} $\mathcal{M}_2 \models \mathrm{Causes}(g,e) \wedge \mathcal{M}_2 \models \neg\mathrm{Causes}(g,e)$ \textit{read} $\mathcal{M}_1 \models \mathrm{Causes}(g,e) \wedge \mathcal{M}_2 \models \neg\mathrm{Causes}(g,e)$; correctly printed in (2.21), p.~3246 \\
\dag\,3244 & proof of Corollary~2.10, l.~1 & \textit{for} Definition~2.7 \textit{read} Definition~2.9 \\
\dag\,3244 & proof of Corollary~2.10, l.~21 & \textit{for} Definition~2.7 \textit{read} Definition~2.8 \\
3244 & l.~28 & \textit{delete} the second occurrence of \emph{the} \\
\dag\,3245 & proof of Corollary~2.10, l.~6 & \textit{for} Definition~2.2 \textit{read} Definition~2.1 \\
\dag\,3245 & proof of Corollary~2.10, (2.20) & \textit{for} $\pos\,\mathrm{Creates}(g,w) \wedge \nec\,\neg\mathrm{Creates}(g,w)$ \textit{read} $\pos\,\mathrm{Creates}(g,w) \wedge \pos\,\neg\mathrm{Creates}(g,w)$ \\
3245 & l.~2 from foot & \textit{for} p.~3250 \textit{read} p.~3240 (Definition~2.7) \\
3246 & l.~16 & \textit{delete} the second occurrence of \emph{the} \\
\dag\,3249 & proof of Lemma~2.12, l.~2 & \textit{for} Definition~2.9 \textit{read} Definition~2.11; the step uses Definition~2.11 (p.~3248) \\
\dag\,3249 & proof of Lemma~2.12, l.~18 & \textit{for} Definition~2.4 \textit{read} Definition~2.2 \\
\dag\,3249 & proof of Lemma~2.12, l.~25 & \textit{for} Definition~2.7 \textit{read} Definition~2.4 \\
3249 & l.~32 & \textit{delete} the second occurrence of \emph{a} \\
3250 & l.~4 & \textit{for} p.~3344 \textit{read} p.~3244 (Definition~2.9) \\
3250 & l.~9 & \textit{for} p.~3349 \textit{read} p.~3249 (Lemma~2.12) \\
3250 & l.~13 & \textit{for} p.~3243 \textit{read} p.~3244 (Corollary~2.10) \\
3250 & l.~29 & \textit{delete} the second occurrence of \emph{the} \\
\end{errlist}

\subsection*{Chapter~3. Models agreeing on maximal greatness \hfill \normalfont\small pp.~3251--3297}
\begin{errlist}
3252 & l.~5 & \textit{for} $\nec(q \vee p) \leftrightarrow (\pos q \vee \pos p)$ \textit{read} $\pos(q \vee p) \leftrightarrow (\pos q \vee \pos p)$ \\
3252 & l.~13 & \textit{for} $w' \models \nec p \iff \exists u\,(w'Ru \wedge u \models p)$ \textit{read} $w' \models \pos p \iff \exists u\,(w'Ru \wedge u \models p)$ \\
3252 & display (3.1) & \textit{for} $v' \models \nec p \iff \forall w'\,(w'Rv' \to w' \models p)$ \textit{read} $v' \models \nec p \iff \forall w'\,(v'Rw' \to w' \models p)$ \\
3252 & l.~31 & \textit{for} correspondecne \textit{read} correspondence \\
3253 & display (3.2) & \textit{for} $w_0 \models \nec p \iff \exists v'\,(w_0Rv' \wedge v' \models p)$ \textit{read} $w_0 \models \pos p \iff \exists v'\,(w_0Rv' \wedge v' \models p)$ \\
3253 & l.~8 & \textit{for} the axiom 4, $\nec r \to r$ \textit{read} the axiom T, $\nec r \to r$ \\
3254 & l.~6 & \textit{for} Lemma~6.6 of Volume~IV \textit{read} Lemma~6.12 of Volume~IV; the result cited is on p.~1622 \\
3257 & l.~12 & \textit{after} accessible \textit{insert} from $w_1$ \\
3258 & l.~18 & \textit{for} accesibility \textit{read} accessibility \\
3259 & l.~7 & \textit{for} forenkowledge \textit{read} foreknowledge \\
3259 & l.~25 & \textit{for} creattion \textit{read} creation \\
3260 & l.~12 & \textit{for} feredom \textit{read} freedom \\
3260 & l.~17 & \textit{delete} the second occurrence of \emph{is} \\
3260 & l.~36 & \textit{for} transtivity \textit{read} transitivity \\
3260 & l.~38 & \textit{for} Lemma~3046 of Volume~X \textit{read} Lemma~3048 of Volume~X \\
3261 & l.~7 & \textit{for} Proposition~5.11 of Volume~III \textit{read} Proposition~5.16 of Volume~III \\
3261 & l.~20 & \textit{for} $w'Rw_1 \leftrightarrow w_1Rw'$ \textit{read} $w'Rw_1 \to w_1Rw'$ \\
3261 & l.~25 & \textit{for} Lemma~3143 of Volume~X \textit{read} Lemma~3137 of Volume~X \\
\dag\,3264 & proof of Corollary~3.3, l.~16 & \textit{for} Lemma~2.10 of Volume~IV \textit{read} Lemma~2.11 of Volume~IV \\
\dag\,3264 & proof of Corollary~3.3, l.~17 & \textit{for} Definition~3.4 \textit{read} Definition~3.2; the step uses Definition~3.2 (p.~3264) \\
3264 & l.~26 & \textit{for} Lemma~3099 of Volume~X \textit{read} Lemma~3093 of Volume~X \\
\dag\,3265 & proof of Corollary~3.3, l.~7 & \textit{for} Definition~3.2 \textit{read} Definition~3.1 \\
3265 & l.~16 & \textit{for} Lemma~3165 of Volume~X \textit{read} Lemma~3169 of Volume~X \\
3265 & l.~21 & \textit{for} Theorem~5.6 of Volume~III \textit{read} Theorem~5.4 of Volume~III; the result cited is on p.~965 \\
3266 & display (3.8) & \textit{for} $V(u) < V(w_0) \to \mathrm{Chooses}(g,u)$ \textit{read} $V(u) < V(w_0) \to \neg\mathrm{Chooses}(g,u)$; the formula is used in its correct form on p.~3267 \\
3266 & l.~33 & \textit{for} Lemma~3147 of Volume~X \textit{read} Lemma~3152 of Volume~X \\
3267 & l.~20 & \textit{after} S5 conditions \textit{insert} on $R$ \\
3267 & l.~6 from foot & \textit{for} p.~3209 \textit{read} p.~3208 (Definition~1.1) \\
3270 & display (3.11) & \textit{for} $\nec(q \leftrightarrow r) \to (\nec q \to \nec r)$ \textit{read} $\nec(q \to r) \to (\nec q \to \nec r)$; correctly printed in (3.13), p.~3272 \\
3271 & display (3.12) & \textit{for} $K_g r \leftrightarrow r$ \textit{read} $K_g r \to r$ \\
3271 & l.~25 & \textit{for} Lemma~3.3 of Volume~I \textit{read} Lemma~3.2 of Volume~I \\
3272 & l.~29 & \textit{for} $w_0Ru$ by reflexivity \textit{read} $w_0Rw_0$ by reflexivity \\
3273 & l.~12 & \textit{for} $v'Rw_2$ by transitivity \textit{read} $v'Rw_2$ by symmetry \\
3273 & display (3.17) & \textit{for} $(C \mathrel{\Box\!\!\to} K_g p) \vee \neg(K_g p \mathrel{\Box\!\!\to} C)$ \textit{read} $(C \mathrel{\Box\!\!\to} K_g p) \wedge \neg(K_g p \mathrel{\Box\!\!\to} C)$; correctly printed in (3.19), p.~3275 \\
3273 & display (3.16) & \textit{for} $\mathcal{M}_1 \models \mathrm{Causes}(g,e) \vee \mathcal{M}_2 \models \neg\mathrm{Causes}(g,e)$ \textit{read} $\mathcal{M}_1 \models \mathrm{Causes}(g,e) \wedge \mathcal{M}_2 \models \neg\mathrm{Causes}(g,e)$ \\
3273 & l.~24 & \textit{for} accssibility \textit{read} accessibility \\
3273 & display (3.15) & \textit{for} $\nec\,\mathrm{Creates}(g,w_0) \wedge \pos\,\neg\mathrm{Creates}(g,w_0)$ \textit{read} $\pos\,\mathrm{Creates}(g,w_0) \wedge \pos\,\neg\mathrm{Creates}(g,w_0)$ \\
3275 & display (3.20) & \textit{for} $\nec(p \wedge p) \leftrightarrow (\nec q \wedge \nec p)$ \textit{read} $\nec(q \wedge p) \leftrightarrow (\nec q \wedge \nec p)$; the formula is used in its correct form on p.~3276 \\
3275 & l.~23 & \textit{delete} the second occurrence of \emph{that} \\
3276 & l.~1 & \textit{for} Lemma~3258 of Volume~X \textit{read} Lemma~3253 of Volume~X \\
3277 & l.~19 & \textit{for} p.~3238 \textit{read} p.~3248 (Definition~2.11) \\
3277 & l.~23 & \textit{for} p.~3364 \textit{read} p.~3264 (Corollary~3.3) \\
3278 & l.~1 & \textit{for} timelless \textit{read} timeless \\
3278 & l.~7 & \textit{for} p.~3254 \textit{read} p.~3264 (Definition~3.2) \\
3278 & display (3.22) & \textit{for} $\nec q \wedge \pos p \leftrightarrow \pos(q \wedge p)$ \textit{read} $\nec q \wedge \pos p \to \pos(q \wedge p)$; correctly printed in (3.23), p.~3280 \\
3280 & proof of Theorem~3.5 & \textit{see} Corrigendum~S16 \\
3281 & l.~19 & \textit{for} Proposition~6.10 of Volume~IV \textit{read} Proposition~6.13 of Volume~IV \\
3281 & l.~32 & \textit{delete} the second occurrence of \emph{to} \\
\dag\,3282 & proof of Corollary~3.6, l.~17 & \textit{for} Lemma~2.10 of Volume~IV \textit{read} Lemma~2.4 of Volume~IV; the result cited is on p.~1419 \\
3283 & l.~15 & \textit{for} modaltiy \textit{read} modality \\
\dag\,3283 & proof of Corollary~3.6, (3.25) & \textit{for} $(w'Rw_1 \wedge w_1Rv) \to w'Rv$ \textit{read} $(w_1Rw' \wedge w_1Rv) \to w'Rv$ \\
\dag\,3283 & proof of Corollary~3.6, l.~18 & \textit{for} Lemma~3218 of Volume~X \textit{read} Lemma~3216 of Volume~X \\
3284 & l.~14 & \textit{for} Lemma~3162 of Volume~X \textit{read} Lemma~3160 of Volume~X \\
3284 & l.~29 & \textit{for} p.~3379 \textit{read} p.~3279 (Theorem~3.5) \\
3285 & l.~21 & \textit{for} accsesibility \textit{read} accessibility \\
3285 & l.~27 & \textit{for} Lemma~3159 of Volume~X \textit{read} Lemma~3158 of Volume~X \\
3286 & l.~3 & \textit{for} $\forall u\,\exists v\,V(u) > V(u)$ \textit{read} $\forall u\,\exists v\,V(v) > V(u)$ \\
3286 & l.~9 & \textit{for} $w_2Rv$ by symmetry \textit{read} $w_2Rv$ by transitivity \\
\dag\,3287 & proof of Proposition~3.8, l.~26 & \textit{for} Lemma~3089 of Volume~X \textit{read} Lemma~3094 of Volume~X; the result cited is on p.~3882 \\
\dag\,3287 & proof of Proposition~3.8, (3.26) & \textit{for} $K_g p \leftrightarrow p$ \textit{read} $K_g p \to p$ \\
3289 & l.~19 & \textit{for} $(uRw_2 \wedge uRv') \to w_1Rv'$ \textit{read} $(uRw_2 \wedge uRv') \to w_2Rv'$ \\
3289 & l.~7 from foot & \textit{for} Lemma~3193 of Volume~X \textit{read} Lemma~3189 of Volume~X \\
\dag\,3290 & proof of Lemma~3.10, l.~6 & \textit{for} Corollary~3.3 \textit{read} Corollary~3.6; the step uses Corollary~3.6 (p.~3282) \\
\dag\,3290 & proof of Lemma~3.10, l.~9 & \textit{for} Definition~2.3 of Volume~IV \textit{read} Definition~2.1 of Volume~IV \\
3290 & l.~25 & \textit{for} p.~3231 \textit{read} p.~3232 (Proposition~2.6) \\
3291 & l.~15 & \textit{for} p.~3288 \textit{read} p.~3287 (Proposition~3.8) \\
3291 & l.~16 & \textit{for} Lemma~3292 of Volume~X \textit{read} Lemma~3295 of Volume~X \\
3291 & l.~28 & \textit{for} p.~3242 \textit{read} p.~3241 (Definition~2.8) \\
3292 & l.~6 & \textit{for} p.~3291 \textit{read} p.~3290 (Lemma~3.10) \\
\dag\,3292 & proof of Theorem~3.11, l.~8 & \textit{for} Lemma~3.13 \textit{read} Lemma~3.10 \\
\dag\,3293 & proof of Theorem~3.11, l.~31 & \textit{for} Assumption~3.1 of Volume~I \textit{read} Assumption~2.1 of Volume~I \\
3293 & l.~5 from foot & \textit{for} casual \textit{read} causal \\
\dag\,3294 & proof of Lemma~3.13, l.~16 & \textit{for} Definition~3.12 \textit{read} Definition~3.9; the step uses Definition~3.9 (p.~3287) \\
3295 & l.~12 & \textit{for} p.~3293 \textit{read} p.~3294 (Lemma~3.13) \\
\dag\,3295 & proof of Lemma~3.13, l.~16 & \textit{for} Definition~3.4 \textit{read} Definition~3.1 \\
\dag\,3295 & proof of Corollary~3.14, (3.27) & \textit{for} $\pos\,\mathrm{Creates}(g,w_0) \vee \pos\,\neg\mathrm{Creates}(g,w_0)$ \textit{read} $\pos\,\mathrm{Creates}(g,w_0) \wedge \pos\,\neg\mathrm{Creates}(g,w_0)$ \\
3295 & l.~33 & \textit{for} modlaity \textit{read} modality \\
3296 & l.~5 & \textit{for} p.~3229 \textit{read} p.~3292 (Theorem~3.11) \\
3296 & l.~6 & \textit{for} intantion \textit{read} intention \\
3296 & l.~33 & \textit{for} p.~3249 \textit{read} p.~3294 (Corollary~3.14) \\
3297 & l.~9 & \textit{for} Lemma~3143 of Volume~X \textit{read} Lemma~3149 of Volume~X \\
\end{errlist}

\subsection*{Chapter~4. Creation and conservation \hfill \normalfont\small pp.~3298--3345}
\begin{errlist}
3299 & l.~10 & \textit{for} Lemma~3210 of Volume~X \textit{read} Lemma~3206 of Volume~X; the result cited is on p.~3898 \\
3299 & l.~24 & \textit{for} Lemma~3289 of Volume~X \textit{read} Lemma~3286 of Volume~X \\
3299 & l.~5 from foot & \textit{for} $V(v') < V(w_0) \to \mathrm{Chooses}(g,v')$ \textit{read} $V(v') < V(w_0) \to \neg\mathrm{Chooses}(g,v')$ \\
3301 & l.~37 & \textit{for} p.~3283 \textit{read} p.~3293 (Definition~3.12) \\
3303 & l.~11 & \textit{for} Lemma~3098 of Volume~X \textit{read} Lemma~3103 of Volume~X \\
3303 & l.~2 from foot & \textit{for} Lemma~3117 of Volume~X \textit{read} Lemma~3119 of Volume~X; the result cited is on p.~3886 \\
3303 & l.~40 & \textit{for} conesrvation \textit{read} conservation \\
3304 & display (4.1) & \textit{for} $\mathcal{M}_1 \models \mathrm{Causes}(g,e) \wedge \mathcal{M}_2 \models \mathrm{Causes}(g,e)$ \textit{read} $\mathcal{M}_1 \models \mathrm{Causes}(g,e) \wedge \mathcal{M}_2 \models \neg\mathrm{Causes}(g,e)$; the formula is used in its correct form on p.~3307 \\
3305 & l.~10 & \textit{for} intension \textit{read} intention \\
3305 & l.~22 & \textit{for} Lemma~3185 of Volume~X \textit{read} Lemma~3181 of Volume~X \\
3305 & display (4.4) & \textit{for} $K_g q \leftrightarrow q$ \textit{read} $K_g q \to q$ \\
3306 & l.~31 & \textit{delete} the second occurrence of \emph{of} \\
3309 & l.~11 & \textit{for} p.~3380 \textit{read} p.~3308 (Proposition~4.3) \\
\dag\,3309 & proof of Proposition~4.4, l.~15 & \textit{for} Proposition~4.6 \textit{read} Proposition~4.3 \\
\dag\,3309 & proof of Proposition~4.4, l.~21 & \textit{for} Lemma~3075 of Volume~X \textit{read} Lemma~3076 of Volume~X \\
\dag\,3310 & proof of Theorem~4.5, l.~5 & \textit{for} Proposition~4.9 \textit{read} Proposition~4.4 \\
\dag\,3310 & proof of Theorem~4.5, l.~6 & \textit{for} Proposition~4.4 \textit{read} Proposition~4.3; the step uses Proposition~4.3 (p.~3308) \\
3311 & l.~14 & \textit{for} p.~3278 \textit{read} p.~3287 (Definition~3.9) \\
\dag\,3311 & proof of Theorem~4.5, (4.6) & \textit{for} $\pos\,\mathrm{Creates}(g,v') \vee \pos\,\neg\mathrm{Creates}(g,v')$ \textit{read} $\pos\,\mathrm{Creates}(g,v') \wedge \pos\,\neg\mathrm{Creates}(g,v')$ \\
\dag\,3311 & proof of Theorem~4.5, l.~4 from foot & \textit{for} Lemma~3118 of Volume~X \textit{read} Lemma~3116 of Volume~X \\
3312 & l.~24 & \textit{for} p.~3281 \textit{read} p.~3282 (Corollary~3.6) \\
3313 & l.~18 & \textit{for} ancetsral \textit{read} ancestral \\
3313 & l.~27 & \textit{after} dependee \textit{insert} in every accessible world \\
\dag\,3315 & proof of Proposition~4.6, l.~1 & \textit{for} Lemma~3205 of Volume~X \textit{read} Lemma~3208 of Volume~X; the result cited is on p.~3899 \\
\dag\,3315 & proof of Proposition~4.6, l.~10 & \textit{for} Definition~6.7 of Volume~IV \textit{read} Definition~6.8 of Volume~IV \\
\dag\,3315 & proof of Proposition~4.6, (4.11) & \textit{for} $\exists w'\,\exists w_0\,V(w_0) > V(w')$ \textit{read} $\forall w'\,\exists w_0\,V(w_0) > V(w')$ \\
\dag\,3315 & proof of Proposition~4.6, l.~30 & \textit{for} Lemma~3260 of Volume~X \textit{read} Lemma~3263 of Volume~X; the result cited is on p.~3907 \\
3315 & l.~32 & \textit{for} Proposition~2.2 of Volume~IV \textit{read} Proposition~2.6 of Volume~IV; the result cited is on p.~1425 \\
3317 & l.~12 & \textit{for} Lemma~3185 of Volume~X \textit{read} Lemma~3182 of Volume~X \\
3318 & l.~10 & \textit{for} p.~3380 \textit{read} p.~3308 (Proposition~4.4) \\
3318 & display (4.14) & \textit{for} $\mathrm{NE}(y) \wedge \neg\exists x\,D(y,x) \leftrightarrow \mathrm{ME}(y)$ \textit{read} $\mathrm{NE}(y) \wedge \neg\exists x\,D(y,x) \to \mathrm{ME}(y)$ \\
3319 & l.~17 & \textit{for} Lemma~3179 of Volume~X \textit{read} Lemma~3177 of Volume~X \\
3319 & l.~18 & \textit{for} $\mathrm{NE}(z) \wedge \neg\forall x\,D(z,x) \to \mathrm{ME}(z)$ \textit{read} $\mathrm{NE}(z) \wedge \neg\exists x\,D(z,x) \to \mathrm{ME}(z)$ \\
3320 & l.~11 & \textit{for} fredom \textit{read} freedom \\
\dag\,3320 & proof of Proposition~4.9, l.~15 & \textit{for} Lemma~3121 of Volume~X \textit{read} Lemma~3115 of Volume~X \\
3320 & l.~3 from foot & \textit{for} Lemma~3219 of Volume~X \textit{read} Lemma~3216 of Volume~X \\
3321 & l.~7 from foot & \textit{for} $ D(z,z)$ \textit{read} $\neg D(z,z)$ \\
\dag\,3322 & proof of Proposition~4.10, l.~28 & \textit{for} Definition~4.11 \textit{read} Definition~4.7; the step uses Definition~4.7 (p.~3317) \\
3322 & l.~30 & \textit{for} Lemma~3263 of Volume~X \textit{read} Lemma~3267 of Volume~X \\
\dag\,3323 & proof of Proposition~4.10, l.~3 & \textit{for} $\forall x\,(\neg\mathrm{NE}(z) \to \exists z\,D(x,z))$ \textit{read} $\forall x\,(\neg\mathrm{NE}(x) \to \exists z\,D(x,z))$ \\
\dag\,3324 & proof of Theorem~4.12, l.~10 & \textit{for} Lemma~3187 of Volume~X \textit{read} Lemma~3189 of Volume~X \\
\dag\,3324 & proof of Theorem~4.12, (4.17) & \textit{for} $\forall w_1\,\exists w\,V(w) > V(w_2)$ \textit{read} $\forall w_1\,\exists w\,V(w) > V(w_1)$ \\
3324 & l.~25 & \textit{for} agnecy \textit{read} agency \\
3325 & l.~9 & \textit{for} Lemma~3096 of Volume~X \textit{read} Lemma~3098 of Volume~X; the result cited is on p.~3883 \\
3326 & l.~4 & \textit{for} p.~3234 \textit{read} p.~3324 (Theorem~4.12) \\
3326 & l.~5 & \textit{delete} the second occurrence of \emph{to} \\
3326 & l.~11 & \textit{for} p.~3422 \textit{read} p.~3322 (Proposition~4.10) \\
3328 & display (4.19) & \textit{for} $(C \mathrel{\Box\!\!\to} K_g r) \vee \neg(K_g r \mathrel{\Box\!\!\to} C)$ \textit{read} $(C \mathrel{\Box\!\!\to} K_g r) \wedge \neg(K_g r \mathrel{\Box\!\!\to} C)$ \\
3328 & l.~14 & \textit{for} Lemma~3220 of Volume~X \textit{read} Lemma~3218 of Volume~X \\
3328 & display (4.18) & \textit{for} $\forall z\,(\neg\mathrm{NE}(z) \to \forall y\,D(z,y))$ \textit{read} $\forall z\,(\neg\mathrm{NE}(z) \to \exists y\,D(z,y))$ \\
3329 & l.~14 & \textit{delete} the second occurrence of \emph{to} \\
\dag\,3330 & proof of Lemma~4.13, l.~7 & \textit{for} Proposition~4.10 \textit{read} Proposition~4.9 \\
3330 & l.~14 & \textit{for} foreknwoledge \textit{read} foreknowledge \\
3330 & l.~15 & \textit{for} well-fuonded \textit{read} well-founded \\
3330 & l.~23 & \textit{for} grounidng \textit{read} grounding \\
\dag\,3331 & proof of Lemma~4.13, l.~20 & \textit{for} Proposition~4.10 \textit{read} Proposition~4.6; the step uses Proposition~4.6 (p.~3315) \\
\dag\,3331 & proof of Lemma~4.13, l.~25 & \textit{for} Definition~4.11 \textit{read} Definition~4.8; the step uses Definition~4.8 (p.~3318) \\
3332 & display (4.22) & \textit{for} $D(y,z) \wedge D(z,x) \leftrightarrow D(y,x)$ \textit{read} $D(y,z) \wedge D(z,x) \to D(y,x)$; correctly printed in (4.24), p.~3333 \\
3332 & display (4.21) & \textit{for} $(C \mathrel{\Box\!\!\to} K_g r) \wedge (K_g r \mathrel{\Box\!\!\to} C)$ \textit{read} $(C \mathrel{\Box\!\!\to} K_g r) \wedge \neg(K_g r \mathrel{\Box\!\!\to} C)$; the formula is used in its correct form on p.~3336 \\
3332 & display (4.23) & \textit{for} $V(w_1) < V(w_0) \to \neg\mathrm{Chooses}(g,w_2)$ \textit{read} $V(w_2) < V(w_0) \to \neg\mathrm{Chooses}(g,w_2)$ \\
3333 & l.~17 & \textit{for} Lemma~3062 of Volume~X \textit{read} Lemma~3067 of Volume~X \\
3333 & l.~30 & \textit{delete} the second occurrence of \emph{that} \\
\dag\,3334 & proof of Lemma~4.15, l.~17 & \textit{for} Definition~4.11 \textit{read} Definition~4.14 \\
3334 & l.~19 & \textit{for} Definition~5.5 of Volume~III \textit{read} Definition~5.1 of Volume~III; the result cited is on p.~952 \\
\dag\,3335 & proof of Lemma~4.15, l.~1 & \textit{for} Definition~4.8 \textit{read} Definition~4.7 \\
3335 & l.~9 & \textit{for} intentoin \textit{read} intention \\
3335 & l.~3 from foot & \textit{for} p.~3277 \textit{read} p.~3267 (Definition~3.4) \\
3336 & l.~29 & \textit{for} Lemma~3113 of Volume~X \textit{read} Lemma~3115 of Volume~X \\
3338 & l.~4 & \textit{for} Lemma~3224 of Volume~X \textit{read} Lemma~3220 of Volume~X; the result cited is on p.~3900 \\
\dag\,3339 & proof of Proposition~4.16, l.~1 & \textit{for} Proposition~4.6 \textit{read} Proposition~4.4; the step uses Proposition~4.4 (p.~3308) \\
\dag\,3339 & proof of Proposition~4.16, l.~26 & \textit{for} Lemma~4.13 \textit{read} Lemma~4.15 \\
\dag\,3339 & proof of Proposition~4.16, l.~29 & \textit{for} Proposition~4.6 \textit{read} Proposition~4.9; the step uses Proposition~4.9 (p.~3319) \\
3339 & l.~2 from foot & \textit{for} Lemma~3295 of Volume~X \textit{read} Lemma~3296 of Volume~X \\
\dag\,3340 & proof of Proposition~4.16, l.~14 & \textit{for} Lemma~3062 of Volume~X \textit{read} Lemma~3058 of Volume~X; the result cited is on p.~3877 \\
3341 & l.~4 & \textit{for} Lemma~3211 of Volume~X \textit{read} Lemma~3216 of Volume~X; the result cited is on p.~3900 \\
3341 & l.~8 from foot & \textit{for} p.~3333 \textit{read} p.~3334 (Definition~4.14) \\
3343 & l.~27 & \textit{for} Definition~5.15 of Volume~III \textit{read} Definition~5.14 of Volume~III \\
3343 & l.~39 & \textit{for} Definition~5.13 of Volume~III \textit{read} Definition~5.18 of Volume~III \\
3344 & l.~17 & \textit{for} p.~3335 \textit{read} p.~3334 (Lemma~4.15) \\
3344 & l.~5 from foot & \textit{for} Lemma~3215 of Volume~X \textit{read} Lemma~3212 of Volume~X; the result cited is on p.~3899 \\
3345 & l.~2 & \textit{for} p.~3333 \textit{read} p.~3323 (Definition~4.11) \\
3345 & l.~10 & \textit{for} p.~3418 \textit{read} p.~3318 (Definition~4.8) \\
\end{errlist}

\subsection*{Chapter~5. The freedom of the maximally great being \hfill \normalfont\small pp.~3346--3392}
\begin{errlist}
3347 & l.~4 & \textit{for} Lemma~3196 of Volume~X \textit{read} Lemma~3194 of Volume~X \\
3347 & l.~39 & \textit{for} Lemma~3110 of Volume~X \textit{read} Lemma~3105 of Volume~X \\
3348 & display (5.1) & \textit{for} $(C \mathrel{\Box\!\!\leftrightarrow} K_g p) \wedge \neg(K_g p \mathrel{\Box\!\!\to} C)$ \textit{read} $(C \mathrel{\Box\!\!\to} K_g p) \wedge \neg(K_g p \mathrel{\Box\!\!\to} C)$ \\
3349 & l.~17 & \textit{delete} the second occurrence of \emph{a} \\
3352 & l.~4 & \textit{for} Lemma~3.2 of Volume~I \textit{read} Lemma~3.3 of Volume~I \\
3352 & l.~17 & \textit{for} Lemma~3228 of Volume~X \textit{read} Lemma~3230 of Volume~X; the result cited is on p.~3902 \\
3352 & display (5.3) & \textit{for} $\pos\,\mathrm{Creates}(g,w_0) \wedge \nec\,\neg\mathrm{Creates}(g,w_0)$ \textit{read} $\pos\,\mathrm{Creates}(g,w_0) \wedge \pos\,\neg\mathrm{Creates}(g,w_0)$; the formula is used in its correct form on p.~3356 \\
\dag\,3354 & proof of Corollary~5.3, l.~18 & \textit{for} $\forall w'\,\exists w'\,V(w') > V(u)$ \textit{read} $\forall u\,\exists w'\,V(w') > V(u)$ \\
3354 & l.~27 & \textit{for} Lemma~3193 of Volume~X \textit{read} Lemma~3190 of Volume~X \\
\dag\,3354 & proof of Corollary~5.3, l.~36 & \textit{for} Lemma~3277 of Volume~X \textit{read} Lemma~3271 of Volume~X \\
3355 & l.~2 & \textit{for} Lemma~3107 of Volume~X \textit{read} Lemma~3105 of Volume~X \\
3355 & l.~3 & \textit{for} Proposition~5.3 of Volume~III \textit{read} Proposition~5.2 of Volume~III \\
3355 & l.~39 & \textit{for} Lemma~3189 of Volume~X \textit{read} Lemma~3188 of Volume~X \\
3356 & l.~4 & \textit{for} Rowe (2004) \textit{read} Adams (1972) (whether God must create the best) \\
3357 & l.~29 & \textit{for} prvidence \textit{read} providence \\
3360 & l.~22 & \textit{for} $V(w_2) < V(u) \to \mathrm{Chooses}(g,w_2)$ \textit{read} $V(w_2) < V(u) \to \neg\mathrm{Chooses}(g,w_2)$ \\
3360 & l.~30 & \textit{for} Kenny (1979) \textit{read} Adams (1972) (whether God must create the best) \\
3361 & l.~25 & \textit{for} Lemma~3208 of Volume~X \textit{read} Lemma~3211 of Volume~X \\
3362 & display (5.6) & \textit{for} $\forall w_2\,\exists v'\,V(v') < V(w_2)$ \textit{read} $\forall w_2\,\exists v'\,V(v') > V(w_2)$; correctly printed in (5.7), p.~3366 \\
3363 & l.~15 & \textit{for} Lemma~3048 of Volume~X \textit{read} Lemma~3046 of Volume~X \\
3363 & l.~22 & \textit{for} Lemma~3160 of Volume~X \textit{read} Lemma~3155 of Volume~X \\
3363 & l.~5 from foot & \textit{for} Lemma~3188 of Volume~X \textit{read} Lemma~3182 of Volume~X \\
3363 & l.~37 & \textit{for} Theorem~5.4 of Volume~III \textit{read} Theorem~5.6 of Volume~III; the result cited is on p.~983 \\
3364 & l.~7 & \textit{for} Lemma~3216 of Volume~X \textit{read} Lemma~3217 of Volume~X \\
3364 & l.~23 & \textit{for} foreknowldge \textit{read} foreknowledge \\
3364 & l.~6 from foot & \textit{for} conservattion \textit{read} conservation \\
3365 & l.~15 & \textit{for} Lemma~3079 of Volume~X \textit{read} Lemma~3074 of Volume~X; the result cited is on p.~3880 \\
3366 & l.~3 & \textit{for} Lemma~3225 of Volume~X \textit{read} Lemma~3226 of Volume~X \\
3367 & l.~15 & \textit{for} p.~3230 \textit{read} p.~3220 (Definition~2.1) \\
3368 & l.~1 & \textit{for} Lemma~3212 of Volume~X \textit{read} Lemma~3209 of Volume~X; the result cited is on p.~3899 \\
3369 & l.~14 & \textit{for} craetion \textit{read} creation \\
3369 & l.~15 & \textit{for} Lemma~3240 of Volume~X \textit{read} Lemma~3237 of Volume~X; the result cited is on p.~3903 \\
\dag\,3369 & proof of Proposition~5.6, l.~25 & \textit{for} Lemma~3197 of Volume~X \textit{read} Lemma~3200 of Volume~X; the result cited is on p.~3898 \\
3371 & l.~21 & \textit{for} Lemma~3075 of Volume~X \textit{read} Lemma~3074 of Volume~X \\
3371 & l.~29 & \textit{for} p.~3318 \textit{read} p.~3319 (Proposition~4.9) \\
3372 & l.~10 & \textit{for} Lemma~3068 of Volume~X \textit{read} Lemma~3069 of Volume~X; the result cited is on p.~3879 \\
3372 & l.~14 & \textit{for} Lemma~3134 of Volume~X \textit{read} Lemma~3139 of Volume~X \\
3373 & l.~24 & \textit{for} Lemma~3110 of Volume~X \textit{read} Lemma~3107 of Volume~X \\
\dag\,3375 & proof of Theorem~5.7, l.~12 & \textit{for} Proposition~5.10 \textit{read} Proposition~5.6; the step uses Proposition~5.6 (p.~3368) \\
3375 & l.~18 & \textit{for} p.~3364 \textit{read} p.~3374 (Theorem~5.7) \\
3376 & l.~14 & \textit{for} p.~3357 \textit{read} p.~3356 (Lemma~5.4) \\
3376 & l.~29 & \textit{for} Lemma~3293 of Volume~X \textit{read} Lemma~3296 of Volume~X \\
3377 & l.~28 & \textit{for} forknowledge \textit{read} foreknowledge \\
3379 & l.~32 & \textit{for} consevration \textit{read} conservation \\
3380 & l.~3 & \textit{for} Lemma~3247 of Volume~X \textit{read} Lemma~3243 of Volume~X \\
3380 & l.~21 & \textit{for} Lemma~3054 of Volume~X \textit{read} Lemma~3058 of Volume~X \\
\dag\,3381 & proof of Corollary~5.11, l.~2 & \textit{for} Lemma~3071 of Volume~X \textit{read} Lemma~3068 of Volume~X \\
3381 & l.~17 & \textit{for} Lemma~3255 of Volume~X \textit{read} Lemma~3251 of Volume~X; the result cited is on p.~3905 \\
3382 & l.~22 & \textit{for} Lemma~3296 of Volume~X \textit{read} Lemma~3290 of Volume~X \\
3382 & l.~25 & \textit{for} ceration \textit{read} creation \\
3383 & l.~31 & \textit{for} Lemma~3242 of Volume~X \textit{read} Lemma~3243 of Volume~X \\
3383 & l.~7 from foot & \textit{for} Lemma~3153 of Volume~X \textit{read} Lemma~3152 of Volume~X \\
3383 & l.~38 & \textit{for} Definition~5.15 of Volume~III \textit{read} Definition~5.18 of Volume~III \\
3384 & l.~14 & \textit{for} Rowe (2005) \textit{read} Rowe (2004) \\
3384 & l.~16 & \textit{for} p.~3210 \textit{read} p.~3310 (Theorem~4.5) \\
3385 & l.~30 & \textit{for} Lemma~3074 of Volume~X \textit{read} Lemma~3071 of Volume~X \\
3385 & l.~35 & \textit{for} p.~3281 \textit{read} p.~3381 (Corollary~5.11) \\
3386 & l.~19 & \textit{for} Lemma~3227 of Volume~X \textit{read} Lemma~3230 of Volume~X; the result cited is on p.~3902 \\
3386 & l.~21 & \textit{for} p.~3367 \textit{read} p.~3368 (Proposition~5.6) \\
3386 & l.~27 & \textit{for} p.~3278 \textit{read} p.~3378 (Proposition~5.10) \\
3386 & l.~3 from foot & \textit{for} Proposition~6.13 of Volume~IV \textit{read} Proposition~6.9 of Volume~IV; the result cited is on p.~1616 \\
\dag\,3387 & proof of Lemma~5.12, l.~29 & \textit{for} Proposition~5.6 \textit{read} Proposition~5.2 \\
3387 & l.~6 from foot & \textit{for} itnention \textit{read} intention \\
\dag\,3388 & proof of Lemma~5.14, l.~10 & \textit{for} Corollary~5.15 \textit{read} Corollary~5.11 \\
3388 & l.~30 & \textit{for} caussal \textit{read} causal \\
3388 & l.~3 from foot & \textit{for} coutnerfactual \textit{read} counterfactual \\
3389 & l.~30 & \textit{for} Lemma~3081 of Volume~X \textit{read} Lemma~3077 of Volume~X \\
3389 & l.~32 & \textit{for} Lemma~3258 of Volume~X \textit{read} Lemma~3254 of Volume~X \\
3389 & l.~6 from foot & \textit{for} p.~3437 \textit{read} p.~3347 (Definition~5.1) \\
3389 & l.~37 & \textit{for} Lemma~3177 of Volume~X \textit{read} Lemma~3175 of Volume~X; the result cited is on p.~3894 \\
\dag\,3390 & proof of Corollary~5.15, l.~7 & \textit{for} Lemma~5.14 \textit{read} Lemma~5.12 \\
\dag\,3390 & proof of Corollary~5.15, l.~27 & \textit{for} Corollary~5.12 of Volume~III \textit{read} Corollary~5.9 of Volume~III; the result cited is on p.~991 \\
3391 & l.~2 & \textit{for} Lemma~3146 of Volume~X \textit{read} Lemma~3152 of Volume~X \\
3391 & l.~9 & \textit{for} p.~3287 \textit{read} p.~3387 (Lemma~5.14) \\
3391 & l.~17 & \textit{for} Lemma~3096 of Volume~X \textit{read} Lemma~3094 of Volume~X; the result cited is on p.~3882 \\
3391 & l.~19 & \textit{for} p.~3363 \textit{read} p.~3353 (Corollary~5.3) \\
3391 & l.~39 & \textit{for} counterfctual \textit{read} counterfactual \\
3392 & l.~13 & \textit{for} Lemma~3130 of Volume~X \textit{read} Lemma~3134 of Volume~X \\
3392 & l.~24 & \textit{for} Lemma~3105 of Volume~X \textit{read} Lemma~3107 of Volume~X; the result cited is on p.~3884 \\
3392 & l.~32 & \textit{for} Lemma~3279 of Volume~X \textit{read} Lemma~3281 of Volume~X \\
\end{errlist}

\subsection*{Chapter~6. Must a maximally great being create the best? \hfill \normalfont\small pp.~3393--3439}
\begin{errlist}
3394 & l.~3 & \textit{for} p.~3307 \textit{read} p.~3317 (Definition~4.7) \\
3395 & l.~30 & \textit{for} providnce \textit{read} providence \\
3397 & l.~17 & \textit{for} Lemma~3116 of Volume~X \textit{read} Lemma~3112 of Volume~X \\
3397 & l.~8 from foot & \textit{for} Lemma~3052 of Volume~X \textit{read} Lemma~3048 of Volume~X; the result cited is on p.~3876 \\
3398 & l.~3 & \textit{for} undetermiend \textit{read} undetermined \\
3398 & l.~8 & \textit{for} Lemma~5.17 of Volume~III \textit{read} Lemma~5.19 of Volume~III \\
\dag\,3398 & proof of Proposition~6.4, l.~23 & \textit{for} Lemma~3185 of Volume~X \textit{read} Lemma~3188 of Volume~X \\
\dag\,3399 & proof of Proposition~6.4, l.~15 & \textit{for} Lemma~3097 of Volume~X \textit{read} Lemma~3096 of Volume~X; the result cited is on p.~3883 \\
\dag\,3400 & proof of Proposition~6.6, l.~15 & \textit{for} Lemma~3175 of Volume~X \textit{read} Lemma~3172 of Volume~X \\
3400 & l.~29 & \textit{for} consrvation \textit{read} conservation \\
3402 & l.~4 & \textit{for} Adams (1972) \textit{read} Kraay (2010) (the theistic multiverse) \\
3402 & l.~12 & \textit{for} Lemma~3160 of Volume~X \textit{read} Lemma~3154 of Volume~X \\
3403 & l.~31 & \textit{for} Kraay (2010) \textit{read} Adams (1972) (whether God must create the best) \\
\dag\,3405 & proof of Corollary~6.7, l.~13 & \textit{for} Lemma~3162 of Volume~X \textit{read} Lemma~3161 of Volume~X \\
3406 & l.~28 & \textit{for} foerknowledge \textit{read} foreknowledge \\
\dag\,3407 & proof of Corollary~6.8, l.~1 & \textit{for} Corollary~6.9 \textit{read} Corollary~6.7; the step uses Corollary~6.7 (p.~3405) \\
\dag\,3407 & proof of Corollary~6.8, l.~11 & \textit{for} Lemma~3098 of Volume~X \textit{read} Lemma~3092 of Volume~X \\
\dag\,3407 & proof of Corollary~6.8, l.~12 & \textit{for} Lemma~3243 of Volume~X \textit{read} Lemma~3242 of Volume~X; the result cited is on p.~3904 \\
\dag\,3408 & proof of Corollary~6.8, l.~16 & \textit{for} Proposition~6.4 \textit{read} Proposition~6.6 \\
\dag\,3409 & proof of Corollary~6.9, l.~26 & \textit{for} Lemma~3065 of Volume~X \textit{read} Lemma~3061 of Volume~X \\
\dag\,3410 & proof of Corollary~6.9, l.~14 & \textit{for} Lemma~3234 of Volume~X \textit{read} Lemma~3239 of Volume~X \\
\dag\,3410 & proof of Corollary~6.9, l.~21 & \textit{for} Proposition~6.2 \textit{read} Proposition~6.6 \\
3410 & l.~35 & \textit{for} Lemma~3149 of Volume~X \textit{read} Lemma~3150 of Volume~X \\
3411 & display (6.2) & \textit{for} $(C \mathrel{\Box\!\!\to} K_g q) \vee \neg(K_g q \mathrel{\Box\!\!\to} C)$ \textit{read} $(C \mathrel{\Box\!\!\to} K_g q) \wedge \neg(K_g q \mathrel{\Box\!\!\to} C)$ \\
3412 & l.~3 & \textit{for} Lemma~3197 of Volume~X \textit{read} Lemma~3201 of Volume~X \\
3412 & l.~15 & \textit{for} Lemma~3185 of Volume~X \textit{read} Lemma~3180 of Volume~X \\
3412 & l.~23 & \textit{delete} the second occurrence of \emph{to} \\
3413 & l.~1 & \textit{for} Lemma~3209 of Volume~X \textit{read} Lemma~3211 of Volume~X; the result cited is on p.~3899 \\
3413 & l.~12 & \textit{for} Lemma~3083 of Volume~X \textit{read} Lemma~3077 of Volume~X; the result cited is on p.~3880 \\
3415 & l.~4 from foot & \textit{for} Lemma~3278 of Volume~X \textit{read} Lemma~3273 of Volume~X; the result cited is on p.~3908 \\
3416 & l.~24 & \textit{for} Lemma~3122 of Volume~X \textit{read} Lemma~3128 of Volume~X; the result cited is on p.~3887 \\
3416 & l.~30 & \textit{for} Lemma~3098 of Volume~X \textit{read} Lemma~3104 of Volume~X; the result cited is on p.~3884 \\
3417 & l.~9 & \textit{for} Lemma~3140 of Volume~X \textit{read} Lemma~3135 of Volume~X; the result cited is on p.~3888 \\
3417 & l.~27 & \textit{for} Lemma~3076 of Volume~X \textit{read} Lemma~3071 of Volume~X; the result cited is on p.~3879 \\
3418 & l.~17 & \textit{for} $\pos\,\mathrm{Creates}(g,w_2) \wedge \nec\,\neg\mathrm{Creates}(g,w_2)$ \textit{read} $\pos\,\mathrm{Creates}(g,w_2) \wedge \pos\,\neg\mathrm{Creates}(g,w_2)$ \\
3419 & l.~4 from foot & \textit{for} Lemma~3118 of Volume~X \textit{read} Lemma~3113 of Volume~X; the result cited is on p.~3885 \\
3420 & display (6.5) & \textit{for} $\forall w_1\,\exists w_2\,V(w_2) < V(w_1)$ \textit{read} $\forall w_1\,\exists w_2\,V(w_2) > V(w_1)$ \\
3420 & display (6.4) & \textit{for} $\exists w'\,\exists u\,V(u) > V(w')$ \textit{read} $\forall w'\,\exists u\,V(u) > V(w')$ \\
3421 & l.~16 & \textit{for} Lemma~3164 of Volume~X \textit{read} Lemma~3161 of Volume~X; the result cited is on p.~3892 \\
3421 & l.~28 & \textit{for} Lemma~3218 of Volume~X \textit{read} Lemma~3220 of Volume~X \\
3423 & l.~6 & \textit{for} Lemma~3247 of Volume~X \textit{read} Lemma~3241 of Volume~X \\
3423 & l.~23 & \textit{for} Lemma~5.21 of Volume~III \textit{read} Lemma~5.17 of Volume~III; the result cited is on p.~1003 \\
3425 & l.~16 & \textit{for} Lemma~3082 of Volume~X \textit{read} Lemma~3086 of Volume~X \\
3425 & l.~27 & \textit{for} p.~3350 \textit{read} p.~3305 (Definition~4.1) \\
3426 & l.~15 & \textit{for} Lemma~3193 of Volume~X \textit{read} Lemma~3198 of Volume~X; the result cited is on p.~3897 \\
3426 & l.~3 from foot & \textit{for} Lemma~3287 of Volume~X \textit{read} Lemma~3293 of Volume~X; the result cited is on p.~3911 \\
3427 & l.~30 & \textit{for} Lemma~6.16 of Volume~IV \textit{read} Lemma~6.12 of Volume~IV \\
3428 & l.~26 & \textit{for} Lemma~3076 of Volume~X \textit{read} Lemma~3072 of Volume~X \\
3428 & display (6.7) & \textit{for} $V(w) > V(w') \to \neg\mathrm{Chooses}(g,w)$ \textit{read} $V(w) < V(w') \to \neg\mathrm{Chooses}(g,w)$; the formula is used in its correct form on p.~3432 \\
3429 & l.~25 & \textit{for} Lemma~3157 of Volume~X \textit{read} Lemma~3161 of Volume~X \\
3429 & l.~1 from foot & \textit{for} Lemma~3077 of Volume~X \textit{read} Lemma~3079 of Volume~X; the result cited is on p.~3880 \\
3430 & l.~22 & \textit{for} Lemma~3135 of Volume~X \textit{read} Lemma~3134 of Volume~X \\
3430 & l.~27 & \textit{for} Lemma~3211 of Volume~X \textit{read} Lemma~3207 of Volume~X \\
3431 & l.~21 & \textit{for} Lemma~3230 of Volume~X \textit{read} Lemma~3229 of Volume~X \\
3431 & l.~2 from foot & \textit{for} timleess \textit{read} timeless \\
3432 & l.~3 & \textit{for} Rowe (2004) \textit{read} Kraay (2010) (the theistic multiverse) \\
\dag\,3432 & proof of Corollary~6.12, l.~4 & \textit{for} Definition~6.14 \textit{read} Definition~6.11 \\
\dag\,3432 & proof of Corollary~6.12, l.~23 & \textit{for} Definition~6.1 \textit{read} Definition~6.5; the step uses Definition~6.5 (p.~3398) \\
\dag\,3432 & proof of Corollary~6.12, l.~27 & \textit{for} Proposition~6.4 \textit{read} Proposition~6.2; the step uses Proposition~6.2 (p.~3395) \\
3432 & l.~28 & \textit{for} p.~3306 \textit{read} p.~3307 (Lemma~4.2) \\
\dag\,3433 & proof of Corollary~6.12, l.~8 & \textit{for} Corollary~6.7 \textit{read} Corollary~6.8 \\
3434 & l.~20 & \textit{for} $V(v') < V(v) \leftrightarrow \neg\mathrm{Chooses}(g,v')$ \textit{read} $V(v') < V(v) \to \neg\mathrm{Chooses}(g,v')$ \\
3434 & display (6.9) & \textit{for} $\pos\,\mathrm{Creates}(g,w_1) \wedge \pos\,\mathrm{Creates}(g,w_1)$ \textit{read} $\pos\,\mathrm{Creates}(g,w_1) \wedge \pos\,\neg\mathrm{Creates}(g,w_1)$ \\
3435 & l.~7 & \textit{for} p.~3344 \textit{read} p.~3434 (Definition~6.13) \\
3435 & l.~9 & \textit{for} undettermined \textit{read} undetermined \\
3435 & l.~10 & \textit{for} Lemma~2.12 of Volume~IV \textit{read} Lemma~2.11 of Volume~IV \\
3435 & l.~23 & \textit{for} timmeless \textit{read} timeless \\
3435 & l.~28 & \textit{for} consrevation \textit{read} conservation \\
3436 & l.~19 & \textit{for} Lemma~3150 of Volume~X \textit{read} Lemma~3148 of Volume~X \\
3436 & l.~33 & \textit{for} Lemma~3060 of Volume~X \textit{read} Lemma~3062 of Volume~X; the result cited is on p.~3878 \\
3437 & l.~21 & \textit{for} p.~3397 \textit{read} p.~3398 (Definition~6.5) \\
3437 & l.~24 & \textit{for} p.~3298 \textit{read} p.~3398 (Lemma~6.3) \\
\dag\,3437 & proof of Lemma~6.15, l.~28 & \textit{for} Definition~6.13 \textit{read} Definition~6.14 \\
\dag\,3438 & proof of Lemma~6.15, l.~21 & \textit{for} Corollary~6.7 \textit{read} Corollary~6.12 \\
3438 & l.~29 & \textit{for} Lemma~3131 of Volume~X \textit{read} Lemma~3137 of Volume~X; the result cited is on p.~3889 \\
\dag\,3438 & proof of Corollary~6.16, l.~30 & \textit{for} Definition~6.11 \textit{read} Definition~6.13 \\
3438 & l.~32 & \textit{for} Lemma~3219 of Volume~X \textit{read} Lemma~3225 of Volume~X; the result cited is on p.~3901 \\
\dag\,3439 & proof of Corollary~6.16, l.~3 & \textit{for} Corollary~6.9 \textit{read} Corollary~6.12 \\
\dag\,3439 & proof of Corollary~6.16, l.~25 & \textit{for} Lemma~3195 of Volume~X \textit{read} Lemma~3190 of Volume~X; the result cited is on p.~3896 \\
3439 & l.~30 & \textit{for} p.~3437 \textit{read} p.~3438 (Corollary~6.16) \\
\end{errlist}

\subsection*{Chapter~7. Omniscience and foreknowledge \hfill \normalfont\small pp.~3440--3487}
\begin{errlist}
3441 & l.~15 & \textit{for} Lemma~3137 of Volume~X \textit{read} Lemma~3139 of Volume~X \\
\dag\,3441 & proof of Proposition~7.2, (7.1) & \textit{for} $\neg K_{z}p \leftrightarrow K_{z}\neg K_{z}p$ \textit{read} $\neg K_{z}p \to K_{z}\neg K_{z}p$; correctly printed in (7.2), p.~3443 \\
3441 & l.~27 & \textit{for} Zagzbski (1991) \textit{read} Zagzebski (1991) \\
\dag\,3442 & proof of Proposition~7.2, l.~2 & \textit{for} $K_{x}p \leftrightarrow K_{x}K_{x}p$ \textit{read} $K_{x}p \to K_{x}K_{x}p$ \\
3442 & l.~5 & \textit{for} Hasker (1989) \textit{read} Pike (1965) (foreknowledge and voluntary action) \\
3442 & l.~26 & \textit{for} Lemma~3270 of Volume~X \textit{read} Lemma~3269 of Volume~X \\
\dag\,3442 & proof of Proposition~7.2, l.~39 & \textit{for} Lemma~3283 of Volume~X \textit{read} Lemma~3282 of Volume~X; the result cited is on p.~3909 \\
3443 & l.~1 & \textit{for} p.~3309 \textit{read} p.~3409 (Corollary~6.9) \\
3444 & l.~29 & \textit{for} Lemma~3259 of Volume~X \textit{read} Lemma~3254 of Volume~X; the result cited is on p.~3905 \\
3444 & l.~30 & \textit{for} Pike (1966) \textit{read} Pike (1965) \\
\dag\,3445 & proof of Corollary~7.3, l.~10 & \textit{for} Lemma~3263 of Volume~X \textit{read} Lemma~3269 of Volume~X \\
\dag\,3445 & proof of Corollary~7.3, (7.3) & \textit{for} $\mathrm{Os}(z) \wedge \mathrm{Os}(z) \to \forall p\,(K_{z}p \leftrightarrow K_{x}p)$ \textit{read} $\mathrm{Os}(z) \wedge \mathrm{Os}(x) \to \forall p\,(K_{z}p \leftrightarrow K_{x}p)$ \\
3447 & l.~6 & \textit{for} Haskr (1989) \textit{read} Hasker (1989) \\
\dag\,3447 & proof of Theorem~7.6, l.~12 & \textit{for} Definition~7.10 \textit{read} Definition~7.5 \\
\dag\,3447 & proof of Theorem~7.6, l.~35 & \textit{for} Lemma~3188 of Volume~X \textit{read} Lemma~3191 of Volume~X \\
\dag\,3448 & proof of Theorem~7.6, (7.4) & \textit{for} $V(w') < V(w') \to \neg\mathrm{Chooses}(g,v')$ \textit{read} $V(v') < V(w') \to \neg\mathrm{Chooses}(g,v')$ \\
3448 & l.~30 & \textit{for} Canotrian \textit{read} Cantorian \\
3452 & display (7.5) & \textit{for} $\mathrm{Os}(z) \to \forall q\,(q \to K_{z}q)$ \textit{read} $\mathrm{Os}(z) \leftrightarrow \forall q\,(q \to K_{z}q)$ \\
3452 & l.~14 & \textit{for} Lemma~3175 of Volume~X \textit{read} Lemma~3179 of Volume~X; the result cited is on p.~3895 \\
3453 & display (7.6) & \textit{for} $\nec\,\mathrm{Creates}(g,v) \wedge \pos\,\neg\mathrm{Creates}(g,v)$ \textit{read} $\pos\,\mathrm{Creates}(g,v) \wedge \pos\,\neg\mathrm{Creates}(g,v)$ \\
3454 & l.~7 & \textit{for} omnisceince \textit{read} omniscience \\
3456 & l.~23 & \textit{for} Stump and Kretzmann (1981) \textit{read} Pike (1965) (foreknowledge and voluntary action) \\
3456 & l.~30 & \textit{for} Lemma~3082 of Volume~X \textit{read} Lemma~3076 of Volume~X \\
3457 & l.~21 & \textit{for} Lemma~3110 of Volume~X \textit{read} Lemma~3104 of Volume~X \\
3458 & l.~9 & \textit{for} p.~3347 \textit{read} p.~3437 (Definition~6.14) \\
3458 & l.~25 & \textit{for} Lemma~3184 of Volume~X \textit{read} Lemma~3187 of Volume~X \\
3458 & l.~26 & \textit{for} Lemma~3280 of Volume~X \textit{read} Lemma~3275 of Volume~X; the result cited is on p.~3908 \\
3458 & l.~32 & \textit{for} $ K_{z}q \to K_{z}\neg K_{z}q$ \textit{read} $\neg K_{z}q \to K_{z}\neg K_{z}q$ \\
3459 & l.~3 & \textit{for} fredeom \textit{read} freedom \\
3460 & display (7.7) & \textit{for} $\mathrm{Os}(z) \wedge \mathrm{Os}(x) \to \exists p\,(K_{z}p \leftrightarrow K_{x}p)$ \textit{read} $\mathrm{Os}(z) \wedge \mathrm{Os}(x) \to \forall p\,(K_{z}p \leftrightarrow K_{x}p)$; the formula is used in its correct form on p.~3462 \\
3460 & l.~18 & \textit{for} Lemma~3254 of Volume~X \textit{read} Lemma~3260 of Volume~X; the result cited is on p.~3906 \\
3461 & l.~3 & \textit{for} Lemma~3141 of Volume~X \textit{read} Lemma~3140 of Volume~X; the result cited is on p.~3889 \\
\dag\,3462 & proof of Proposition~7.8, l.~15 & \textit{for} Definition~7.5 \textit{read} Definition~7.7 \\
\dag\,3463 & proof of Proposition~7.9, (7.8) & \textit{for} $\mathrm{NE}(z) \leftrightarrow \neg\exists y\,D(z,y)$ \textit{read} $\mathrm{NE}(z) \to \neg\exists y\,D(z,y)$ \\
3463 & n.~23 & \textit{for} Pike (1964) \textit{read} Pike (1965) \\
3464 & l.~33 & \textit{for} Lemma~3051 of Volume~X \textit{read} Lemma~3057 of Volume~X \\
3465 & l.~18 & \textit{after} every true proposition \textit{insert} (plurally quantified) \\
3465 & l.~22 & \textit{for} p.~3463 \textit{read} p.~3464 (Definition~7.10) \\
\dag\,3465 & proof of Lemma~7.12, l.~23 & \textit{for} Proposition~7.13 \textit{read} Proposition~7.9; the step uses Proposition~7.9 (p.~3463) \\
\dag\,3465 & proof of Lemma~7.12, l.~30 & \textit{for} Proposition~7.13 \textit{read} Proposition~7.8; the step uses Proposition~7.8 (p.~3461) \\
\dag\,3465 & proof of Lemma~7.12, l.~2 from foot & \textit{for} Lemma~3120 of Volume~X \textit{read} Lemma~3121 of Volume~X \\
3466 & l.~37 & \textit{for} p.~3353 \textit{read} p.~3253 (Definition~3.1) \\
3467 & display (7.9) & \textit{for} $\mathrm{Op}(x) \leftrightarrow \forall t\,(\mathrm{Task}(t) \vee \neg Q(t,x) \to \mathrm{Can}(x,t))$ \textit{read} $\mathrm{Op}(x) \leftrightarrow \forall t\,(\mathrm{Task}(t) \wedge \neg Q(t,x) \to \mathrm{Can}(x,t))$; the formula is used in its correct form on p.~3471 \\
3467 & l.~34 & \textit{for} p.~3462 \textit{read} p.~3463 (Proposition~7.9) \\
3469 & l.~3 & \textit{for} providnece \textit{read} providence \\
\dag\,3469 & proof of Proposition~7.13, l.~5 & \textit{for} Lemma~7.16 \textit{read} Lemma~7.12 \\
\dag\,3469 & proof of Proposition~7.13, l.~10 & \textit{for} Definition~7.5 \textit{read} Definition~7.1 \\
\dag\,3469 & proof of Proposition~7.13, l.~16 & \textit{for} Definition~7.11 \textit{read} Definition~7.10 \\
\dag\,3469 & proof of Proposition~7.13, l.~19 & \textit{for} Proposition~7.9 \textit{read} Proposition~7.8 \\
3470 & l.~26 & \textit{for} Lemma~3203 of Volume~X \textit{read} Lemma~3208 of Volume~X; the result cited is on p.~3899 \\
3470 & l.~2 from foot & \textit{for} Lemma~3225 of Volume~X \textit{read} Lemma~3221 of Volume~X \\
3471 & l.~20 & \textit{for} Lemma~3267 of Volume~X \textit{read} Lemma~3269 of Volume~X; the result cited is on p.~3907 \\
3472 & l.~14 & \textit{for} cretaion \textit{read} creation \\
3472 & l.~27 & \textit{for} Definition~5.13 of Volume~III \textit{read} Definition~5.15 of Volume~III \\
3472 & l.~1 from foot & \textit{for} p.~3410 \textit{read} p.~3400 (Proposition~6.6) \\
\dag\,3473 & proof of Proposition~7.14, l.~29 & \textit{for} Lemma~3098 of Volume~X \textit{read} Lemma~3101 of Volume~X; the result cited is on p.~3883 \\
3475 & display (7.10) & \textit{for} $\neg Q(t,z) \wedge \mathrm{Task}(t) \leftrightarrow \mathrm{Can}(z,t)$ \textit{read} $\neg Q(t,z) \wedge \mathrm{Task}(t) \to \mathrm{Can}(z,t)$; the formula is used in its correct form on p.~3478 \\
3475 & l.~9 & \textit{for} Hasker (1988) \textit{read} Hasker (1989) \\
3475 & l.~12 & \textit{for} $\pos\,\mathrm{Creates}(g,u) \wedge \pos\,\mathrm{Creates}(g,u)$ \textit{read} $\pos\,\mathrm{Creates}(g,u) \wedge \pos\,\neg\mathrm{Creates}(g,u)$ \\
3475 & l.~15 & \textit{for} Lemma~3186 of Volume~X \textit{read} Lemma~3191 of Volume~X; the result cited is on p.~3896 \\
3475 & l.~34 & \textit{for} Hasker (1989) \textit{read} Stump and Kretzmann (1981) (atemporal eternity) \\
3476 & l.~7 from foot & \textit{for} Lemma~3058 of Volume~X \textit{read} Lemma~3053 of Volume~X; the result cited is on p.~3877 \\
3477 & l.~2 from foot & \textit{for} Lemma~3206 of Volume~X \textit{read} Lemma~3201 of Volume~X; the result cited is on p.~3898 \\
3478 & l.~29 & \textit{for} Lemma~3198 of Volume~X \textit{read} Lemma~3194 of Volume~X; the result cited is on p.~3897 \\
3479 & l.~5 from foot & \textit{for} Lemma~3046 of Volume~X \textit{read} Lemma~3050 of Volume~X \\
3480 & display (7.11) & \textit{for} $\mathrm{MP}(y) \to \neg\nec\,\mathrm{Wrong}(y,a)$ \textit{read} $\mathrm{MP}(y) \to \neg\pos\,\mathrm{Wrong}(y,a)$ \\
3480 & l.~40 & \textit{for} Zagzebksi (1991) \textit{read} Zagzebski (1991) \\
3481 & l.~38 & \textit{for} Proposition~2.6 of Volume~IV \textit{read} Proposition~2.2 of Volume~IV; the result cited is on p.~1418 \\
3482 & l.~10 & \textit{after} every task \textit{insert} whose description does not quantify over the agent's own powers \\
3482 & l.~18 & \textit{for} $V(v') < V(w_1) \leftrightarrow \neg\mathrm{Chooses}(g,v')$ \textit{read} $V(v') < V(w_1) \to \neg\mathrm{Chooses}(g,v')$ \\
3483 & l.~14 & \textit{for} Lemma~3290 of Volume~X \textit{read} Lemma~3286 of Volume~X; the result cited is on p.~3910 \\
3483 & l.~18 & \textit{for} p.~3331 \textit{read} p.~3330 (Lemma~4.13) \\
3483 & display (7.12) & \textit{for} $(C \mathrel{\Box\!\!\to} K_g q) \wedge \neg(K_g q \mathrel{\Box\!\!\leftrightarrow} C)$ \textit{read} $(C \mathrel{\Box\!\!\to} K_g q) \wedge \neg(K_g q \mathrel{\Box\!\!\to} C)$; the formula is used in its correct form on p.~3486 \\
\dag\,3484 & proof of Lemma~7.16, l.~17 & \textit{for} Definition~7.10 \textit{read} Definition~7.15 \\
\dag\,3484 & proof of Lemma~7.16, l.~19 & \textit{for} Proposition~7.9 \textit{read} Proposition~7.14 \\
\dag\,3484 & proof of Lemma~7.16, l.~23 & \textit{for} Definition~7.7 \textit{read} Definition~7.11 \\
\dag\,3485 & proof of Lemma~7.16, l.~1 & \textit{for} $\mathrm{Os}(y) \leftrightarrow \exists q\,(q \to K_{y}q)$ \textit{read} $\mathrm{Os}(y) \leftrightarrow \forall q\,(q \to K_{y}q)$ \\
3485 & l.~4 & \textit{for} perfecsion \textit{read} perfection \\
3486 & l.~19 & \textit{for} providecne \textit{read} providence \\
3487 & l.~2 & \textit{for} Lemma~3155 of Volume~X \textit{read} Lemma~3156 of Volume~X \\
3487 & l.~13 & \textit{for} Zagzbeski (1991) \textit{read} Zagzebski (1991) \\
3487 & l.~30 & \textit{for} Lemma~3070 of Volume~X \textit{read} Lemma~3071 of Volume~X; the result cited is on p.~3879 \\
\end{errlist}

\subsection*{Chapter~8. Particular providence and the limits of the result \hfill \normalfont\small pp.~3488--3520}
\begin{errlist}
3489 & l.~16 & \textit{for} Lemma~3220 of Volume~X \textit{read} Lemma~3223 of Volume~X \\
3489 & l.~26 & \textit{for} p.~3450 \textit{read} p.~3405 (Corollary~6.7) \\
3489 & l.~3 from foot & \textit{for} Lemma~3240 of Volume~X \textit{read} Lemma~3238 of Volume~X \\
3489 & display (8.1) & \textit{for} $\forall w'\,\exists v\,V(v) < V(w')$ \textit{read} $\forall w'\,\exists v\,V(v) > V(w')$; the formula is used in its correct form on p.~3491 \\
3490 & l.~28 & \textit{for} p.~3397 \textit{read} p.~3398 (Proposition~6.4) \\
3490 & l.~5 from foot & \textit{for} tmieless \textit{read} timeless \\
3491 & l.~26 & \textit{for} $\exists v\,\exists u\,V(u) > V(v)$ \textit{read} $\forall v\,\exists u\,V(u) > V(v)$ \\
3491 & l.~31 & \textit{for} Lemma~3206 of Volume~X \textit{read} Lemma~3202 of Volume~X \\
3492 & l.~4 from foot & \textit{for} Lemma~3170 of Volume~X \textit{read} Lemma~3174 of Volume~X \\
\dag\,3494 & proof of Proposition~8.5, l.~14 & \textit{for} Lemma~8.3 \textit{read} Lemma~8.4 \\
\dag\,3495 & proof of Lemma~8.4, l.~18 & \textit{for} Definition~8.1 \textit{read} Definition~8.2; the step uses Definition~8.2 (p.~3490) \\
3495 & l.~22 & \textit{for} Lemma~3153 of Volume~X \textit{read} Lemma~3149 of Volume~X \\
3495 & l.~3 from foot & \textit{for} p.~3274 \textit{read} p.~3374 (Definition~5.8) \\
3496 & l.~4 & \textit{for} Lemma~3059 of Volume~X \textit{read} Lemma~3057 of Volume~X; the result cited is on p.~3877 \\
3496 & l.~17 & \textit{for} Lemma~3110 of Volume~X \textit{read} Lemma~3111 of Volume~X; the result cited is on p.~3885 \\
3496 & l.~20 & \textit{for} Lemma~3244 of Volume~X \textit{read} Lemma~3241 of Volume~X \\
3496 & l.~34 & \textit{for} Lemma~3148 of Volume~X \textit{read} Lemma~3143 of Volume~X; the result cited is on p.~3889 \\
3499 & l.~1 & \textit{for} Lemma~3079 of Volume~X \textit{read} Lemma~3076 of Volume~X \\
3500 & l.~22 & \textit{for} Lemma~3231 of Volume~X \textit{read} Lemma~3228 of Volume~X \\
3501 & l.~11 & \textit{for} Lemma~3085 of Volume~X \textit{read} Lemma~3081 of Volume~X; the result cited is on p.~3881 \\
3501 & l.~13 & \textit{for} counetrfactual \textit{read} counterfactual \\
\dag\,3501 & proof of Corollary~8.6, l.~21 & \textit{for} Proposition~8.9 \textit{read} Proposition~8.5; the step uses Proposition~8.5 (p.~3494) \\
3501 & l.~27 & \textit{for} p.~3414 \textit{read} p.~3441 (Definition~7.1) \\
3501 & l.~28 & \textit{for} p.~3329 \textit{read} p.~3339 (Proposition~4.16) \\
\dag\,3502 & proof of Corollary~8.6, l.~2 & \textit{for} $V(v') < V(w_2) \to \neg\mathrm{Chooses}(g,v')$ \textit{read} $V(v') < V(w_1) \to \neg\mathrm{Chooses}(g,v')$ \\
3503 & l.~6 & \textit{for} p.~3359 \textit{read} p.~3395 (Proposition~6.2) \\
3503 & l.~11 & \textit{for} Lemma~3123 of Volume~X \textit{read} Lemma~3121 of Volume~X; the result cited is on p.~3886 \\
3504 & l.~4 from foot & \textit{for} cauasl \textit{read} causal \\
3506 & l.~3 & \textit{for} $V(w_1) > V(v) \to \neg\mathrm{Chooses}(g,w_1)$ \textit{read} $V(w_1) < V(v) \to \neg\mathrm{Chooses}(g,w_1)$ \\
3506 & l.~23 & \textit{for} Lemma~3059 of Volume~X \textit{read} Lemma~3064 of Volume~X \\
3507 & l.~3 & \textit{for} Lemma~3087 of Volume~X \textit{read} Lemma~3089 of Volume~X; the result cited is on p.~3882 \\
\dag\,3507 & proof of Theorem~8.8, l.~21 & \textit{for} Definition~6.2 of Volume~IV \textit{read} Definition~6.8 of Volume~IV \\
3507 & l.~22 & \textit{for} agecny \textit{read} agency \\
3507 & l.~27 & \textit{for} p.~3394 \textit{read} p.~3494 (Proposition~8.5) \\
\dag\,3508 & proof of Theorem~8.8, l.~8 & \textit{for} Lemma~8.7 \textit{read} Lemma~8.3 \\
\dag\,3508 & proof of Lemma~8.10, l.~26 & \textit{for} Lemma~8.4 \textit{read} Lemma~8.7 \\
3509 & l.~6 from foot & \textit{for} Lemma~3207 of Volume~X \textit{read} Lemma~3210 of Volume~X \\
3511 & l.~4 & \textit{for} Lemma~3170 of Volume~X \textit{read} Lemma~3175 of Volume~X \\
3511 & l.~12 & \textit{for} Lemma~3292 of Volume~X \textit{read} Lemma~3293 of Volume~X; the result cited is on p.~3911 \\
3511 & l.~39 & \textit{for} p.~3449 \textit{read} p.~3494 (Lemma~8.4) \\
3512 & l.~37 & \textit{for} Lemma~3207 of Volume~X \textit{read} Lemma~3211 of Volume~X \\
3514 & l.~32 & \textit{for} Lemma~3275 of Volume~X \textit{read} Lemma~3273 of Volume~X \\
3515 & l.~4 & \textit{for} Definition~5.5 of Volume~III \textit{read} Definition~5.10 of Volume~III \\
3515 & l.~5 & \textit{for} Lemma~3107 of Volume~X \textit{read} Lemma~3104 of Volume~X \\
3515 & l.~8 & \textit{for} Lemma~3154 of Volume~X \textit{read} Lemma~3155 of Volume~X; the result cited is on p.~3891 \\
3516 & l.~5 & \textit{for} p.~3500 \textit{read} p.~3501 (Corollary~8.6) \\
3517 & l.~3 & \textit{for} Lemma~3207 of Volume~X \textit{read} Lemma~3206 of Volume~X; the result cited is on p.~3898 \\
3517 & l.~9 & \textit{for} p.~3454 \textit{read} p.~3445 (Corollary~7.3) \\
3517 & l.~22 & \textit{for} Lemma~3171 of Volume~X \textit{read} Lemma~3169 of Volume~X; the result cited is on p.~3893 \\
3517 & l.~28 & \textit{for} p.~3364 \textit{read} p.~3464 (Definition~7.11) \\
3518 & l.~14 & \textit{for} Lemma~3279 of Volume~X \textit{read} Lemma~3283 of Volume~X; the result cited is on p.~3909 \\
3518 & display (8.8) & \textit{for} $\pos\,\mathrm{Creates}(g,u) \wedge \nec\,\neg\mathrm{Creates}(g,u)$ \textit{read} $\pos\,\mathrm{Creates}(g,u) \wedge \pos\,\neg\mathrm{Creates}(g,u)$ \\
3518 & l.~40 & \textit{for} Lemma~3155 of Volume~X \textit{read} Lemma~3152 of Volume~X; the result cited is on p.~3891 \\
3520 & l.~6 & \textit{for} Proposition~5.2 of Volume~III \textit{read} Proposition~5.3 of Volume~III; the result cited is on p.~965 \\
3520 & l.~18 & \textit{for} Lemma~5.21 of Volume~III \textit{read} Lemma~5.19 of Volume~III; the result cited is on p.~1005 \\
3520 & l.~8 from foot & \textit{for} p.~3475 \textit{read} p.~3465 (Lemma~7.12) \\
\end{errlist}

\subsection*{Chapter~9. Summary \hfill \normalfont\small pp.~3521--3537}
\begin{errlist}
3522 & l.~30 & \textit{for} Lemma~3055 of Volume~X \textit{read} Lemma~3060 of Volume~X; the result cited is on p.~3878 \\
3523 & l.~17 & \textit{for} Lemma~3249 of Volume~X \textit{read} Lemma~3245 of Volume~X \\
3524 & l.~20 & \textit{for} follwoing \textit{read} following \\
3524 & l.~29 & \textit{for} Lemma~3119 of Volume~X \textit{read} Lemma~3114 of Volume~X \\
3524 & l.~37 & \textit{for} ovrview \textit{read} overview \\
3525 & l.~10 & \textit{for} ovreview \textit{read} overview \\
3525 & l.~37 & \textit{for} inrtoduction \textit{read} introduction \\
3527 & l.~26 & \textit{for} Lemma~3061 of Volume~X \textit{read} Lemma~3057 of Volume~X \\
3527 & l.~4 from foot & \textit{for} Lemma~3120 of Volume~X \textit{read} Lemma~3126 of Volume~X; the result cited is on p.~3887 \\
3528 & l.~3 & \textit{for} oevrview \textit{read} overview \\
3528 & l.~16 & \textit{for} Lemma~3202 of Volume~X \textit{read} Lemma~3206 of Volume~X; the result cited is on p.~3898 \\
3529 & l.~7 & \textit{for} Lemma~3169 of Volume~X \textit{read} Lemma~3173 of Volume~X \\
3529 & l.~32 & \textit{for} Lemma~3058 of Volume~X \textit{read} Lemma~3064 of Volume~X \\
3530 & l.~3 & \textit{for} presenttaion \textit{read} presentation \\
\dag\,3530 & proof of Corollary~9.2, l.~30 & \textit{for} Lemma~3081 of Volume~X \textit{read} Lemma~3082 of Volume~X; the result cited is on p.~3881 \\
3530 & l.~5 from foot & \textit{for} p.~3427 \textit{read} p.~3472 (Proposition~7.14) \\
3531 & l.~6 & \textit{for} p.~3531 \textit{read} p.~3530 (Corollary~9.2) \\
3531 & l.~6 from foot & \textit{for} Lemma~3063 of Volume~X \textit{read} Lemma~3062 of Volume~X \\
3532 & l.~14 & \textit{for} Lemma~3190 of Volume~X \textit{read} Lemma~3186 of Volume~X; the result cited is on p.~3896 \\
3532 & l.~5 from foot & \textit{for} Lemma~3218 of Volume~X \textit{read} Lemma~3214 of Volume~X \\
3533 & l.~7 & \textit{for} Lemma~3156 of Volume~X \textit{read} Lemma~3152 of Volume~X \\
3533 & l.~11 & \textit{for} Lemma~3244 of Volume~X \textit{read} Lemma~3249 of Volume~X \\
3533 & l.~15 & \textit{for} Lemma~2.10 of Volume~IV \textit{read} Lemma~2.12 of Volume~IV \\
3534 & l.~7 & \textit{for} introdutcion \textit{read} introduction \\
3534 & l.~17 & \textit{for} Lemma~3207 of Volume~X \textit{read} Lemma~3208 of Volume~X \\
3534 & l.~1 from foot & \textit{for} Definition~2.4 of Volume~I \textit{read} Definition~2.2 of Volume~I \\
3535 & l.~16 & \textit{for} Lemma~3133 of Volume~X \textit{read} Lemma~3137 of Volume~X \\
3536 & l.~17 & \textit{for} presentaiton \textit{read} presentation \\
\dag\,3537 & proof of Lemma~9.3, l.~32 & \textit{for} Lemma~3183 of Volume~X \textit{read} Lemma~3182 of Volume~X; the result cited is on p.~3895 \\
\end{errlist}


\section{Volume X: Formal Appendices}\label{s:vol10}
% generated by generate.py -- do not edit
\subsection*{Front matter \hfill \normalfont\small pp.~3538--3540}
\begin{errlist}
3539 & Contents, Chapter~6 & \textit{for} 3843 \textit{read} 3833 \\
\end{errlist}

\subsection*{Chapter~1. The modal language and its semantics \hfill \normalfont\small pp.~3541--3581}
\begin{errlist}
\dag\,3542 & proof of Lemma~3, l.~21 & \textit{for} Lemma~2 \textit{read} Lemma~1; the step uses Lemma~1 (p.~3542) \\
\dag\,3543 & proof of Lemma~15, l.~19 & \textit{for} Lemma~2 \textit{read} Lemma~4 \\
\dag\,3544 & proof of Lemma~30, (1.3) & \textit{for} $u \models \nec q \iff \forall w_1\,(uRw_1 \leftrightarrow w_1 \models q)$ \textit{read} $u \models \nec q \iff \forall w_1\,(uRw_1 \to w_1 \models q)$; correctly printed in (1.5), p.~3547 \\
\dag\,3544 & proof of Lemma~30, (1.2) & \textit{for} $w \models \nec p \iff \forall w_1\,(wRw_1 \leftrightarrow w_1 \models p)$ \textit{read} $w \models \nec p \iff \forall w_1\,(wRw_1 \to w_1 \models p)$ \\
\dag\,3545 & proof of Lemma~29, l.~6 & \textit{for} Lemma~20 \textit{read} Lemma~25 \\
3545 & l.~25 & \textit{delete} the second occurrence of \emph{to} \\
\dag\,3545 & proof of Lemma~40, (1.4) & \textit{for} $(w'Rw_0 \vee w'Rw_2) \to w_0Rw_2$ \textit{read} $(w'Rw_0 \wedge w'Rw_2) \to w_0Rw_2$; the formula is used in its correct form on p.~3549 \\
\dag\,3545 & proof of Lemma~40, l.~35 & \textit{for} Definition~2.4 of Volume~I \textit{read} Definition~2.2 of Volume~I \\
3546 & l.~13 & \textit{delete} the second occurrence of \emph{is} \\
\dag\,3546 & proof of Lemma~47, l.~16 & \textit{for} Lemma~42 \textit{read} Lemma~46 \\
\dag\,3547 & proof of Lemma~42, l.~5 & \textit{for} Lemma~22 \textit{read} Lemma~24 \\
\dag\,3547 & proof of Lemma~58, l.~6 & \textit{for} Lemma~22 \textit{read} Lemma~18 \\
\dag\,3547 & proof of Lemma~60, (1.6) & \textit{for} $\nec(p \wedge r) \to (\nec p \wedge \nec r)$ \textit{read} $\nec(p \wedge r) \leftrightarrow (\nec p \wedge \nec r)$ \\
\dag\,3548 & proof of Lemma~55, l.~11 & \textit{for} Lemma~21 \textit{read} Lemma~23 \\
\dag\,3548 & proof of Lemma~62, l.~25 & \textit{for} Lemma~22 \textit{read} Lemma~26 \\
3548 & l.~27 & \textit{delete} the second occurrence of \emph{of} \\
3549 & l.~5 & \textit{for} refleixvity \textit{read} reflexivity \\
\dag\,3549 & proof of Lemma~79, l.~13 & \textit{for} Lemma~5 \textit{read} Lemma~2 \\
3549 & l.~24 & \textit{for} semnatics \textit{read} semantics \\
3549 & l.~28 & \textit{for} p.~3549 \textit{read} p.~3548 (Lemma~69) \\
\dag\,3549 & proof of Lemma~80, l.~29 & \textit{for} Assumption~2.1 of Volume~I \textit{read} Assumption~3.1 of Volume~I \\
\dag\,3550 & proof of Lemma~85, l.~2 & \textit{for} Lemma~84 \textit{read} Lemma~82 \\
\dag\,3550 & proof of Lemma~83, l.~6 & \textit{for} Lemma~28 \textit{read} Lemma~25 \\
\dag\,3550 & proof of Lemma~83, l.~8 & \textit{for} Lemma~33 \textit{read} Lemma~36 \\
\dag\,3550 & proof of Lemma~72, l.~9 & \textit{for} Lemma~57 \textit{read} Lemma~56; the step uses Lemma~56 (p.~3547) \\
\dag\,3550 & proof of Lemma~86, l.~12 & \textit{for} Lemma~41 \textit{read} Lemma~44; the step uses Lemma~44 (p.~3546) \\
\dag\,3550 & proof of Lemma~90, (1.7) & \textit{for} $\nec q \wedge \pos p \to \pos(q \wedge q)$ \textit{read} $\nec q \wedge \pos p \to \pos(q \wedge p)$ \\
\dag\,3551 & proof of Lemma~99, l.~16 & \textit{for} Definition~2.3 of Volume~I \textit{read} Definition~2.2 of Volume~I \\
\dag\,3551 & proof of Lemma~95, l.~18 & \textit{for} Lemma~61 \textit{read} Lemma~66; the step uses Lemma~66 (p.~3548) \\
\dag\,3551 & proof of Lemma~87, l.~21 & \textit{for} Lemma~11 \textit{read} Lemma~10 \\
\dag\,3551 & proof of Lemma~85, l.~29 & \textit{for} Lemma~26 \textit{read} Lemma~21 \\
3551 & l.~7 from foot & \textit{for} p.~3594 \textit{read} p.~3549 (Lemma~80) \\
3552 & l.~21 & \textit{for} $vRv'$ by symmetry \textit{read} $vRv'$ by transitivity \\
\dag\,3552 & proof of Lemma~107, l.~28 & \textit{for} Lemma~41 \textit{read} Lemma~37; the step uses Lemma~37 (p.~3545) \\
\dag\,3553 & proof of Lemma~119, l.~20 & \textit{for} Lemma~63 \textit{read} Lemma~59 \\
\dag\,3553 & proof of Lemma~116, l.~26 & \textit{for} Lemma~98 \textit{read} Lemma~100; the step uses Lemma~100 (p.~3552) \\
3553 & l.~4 from foot & \textit{after} true at every world \textit{insert} accessible from $v'$ \\
\dag\,3554 & proof of Lemma~129, l.~1 & \textit{for} Definition~2.3 of Volume~I \textit{read} Definition~2.4 of Volume~I \\
3554 & l.~8 & \textit{delete} the second occurrence of \emph{is} \\
\dag\,3555 & proof of Lemma~130, l.~11 & \textit{for} Lemma~29 \textit{read} Lemma~26; the step uses Lemma~26 (p.~3544) \\
\dag\,3556 & proof of Lemma~139, l.~8 & \textit{for} Lemma~127 \textit{read} Lemma~123 \\
\dag\,3557 & proof of Lemma~159, l.~3 & \textit{for} Definition~2.2 of Volume~I \textit{read} Definition~2.4 of Volume~I \\
3557 & l.~24 & \textit{for} $\nec(r \to q) \to (\nec r \to \nec r)$ \textit{read} $\nec(r \to q) \to (\nec r \to \nec q)$ \\
\dag\,3557 & proof of Lemma~145, l.~26 & \textit{for} Lemma~89 \textit{read} Lemma~85; the step uses Lemma~85 (p.~3550) \\
\dag\,3558 & proof of Lemma~154, l.~4 & \textit{for} Lemma~88 \textit{read} Lemma~92; the step uses Lemma~92 (p.~3551) \\
\dag\,3558 & proof of Lemma~162, l.~18 & \textit{for} Lemma~17 \textit{read} Lemma~16 \\
3559 & l.~29 & \textit{for} $v'Rw'$ by reflexivity \textit{read} $v'Rv'$ by reflexivity \\
\dag\,3560 & proof of Lemma~189, (1.11) & \textit{for} $w_1Rw \to wRw_2$ \textit{read} $w_2Rw \to wRw_2$ \\
\dag\,3560 & proof of Lemma~183, l.~13 & \textit{for} Lemma~67 \textit{read} Lemma~69 \\
\dag\,3560 & proof of Lemma~173, l.~18 & \textit{for} Lemma~145 \textit{read} Lemma~141; the step uses Lemma~141 (p.~3556) \\
\dag\,3560 & proof of Lemma~189, (1.12) & \textit{for} $w_0Rw' \vee w'Rw_1 \to w_0Rw_1$ \textit{read} $w_0Rw' \wedge w'Rw_1 \to w_0Rw_1$ \\
\dag\,3560 & proof of Lemma~186, l.~24 & \textit{for} Lemma~157 \textit{read} Lemma~155 \\
\dag\,3560 & proof of Lemma~189, (1.10) & \textit{for} $\pos p \leftrightarrow \nec\neg p$ \textit{read} $\neg\pos p \leftrightarrow \nec\neg p$ \\
3560 & l.~5 from foot & \textit{delete} the second occurrence of \emph{that} \\
\dag\,3561 & proof of Lemma~198, (1.13) & \textit{for} $(vRv' \wedge vRw) \to wRv'$ \textit{read} $(vRv' \wedge vRw) \to v'Rw$; correctly printed in (1.15), p.~3564 \\
\dag\,3561 & proof of Lemma~193, l.~25 & \textit{for} Lemma~29 \textit{read} Lemma~33 \\
3561 & l.~4 from foot & \textit{after} S5 conditions \textit{insert} on $R$ \\
\dag\,3562 & proof of Lemma~202, l.~21 & \textit{for} Lemma~150 \textit{read} Lemma~146 \\
3563 & Lemma~214 & \textit{see} Corrigendum~S17 \\
3564 & l.~3 & \textit{for} p.~3595 \textit{read} p.~3559 (Lemma~170) \\
3564 & l.~15 & \textit{delete} the second occurrence of \emph{a} \\
3564 & l.~31 & \textit{for} p.~3557 \textit{read} p.~3556 (Lemma~144) \\
\dag\,3565 & proof of Lemma~230, l.~4 & \textit{for} Lemma~125 \textit{read} Lemma~126; the step uses Lemma~126 (p.~3554) \\
\dag\,3566 & proof of Lemma~240, l.~6 & \textit{for} Lemma~90 \textit{read} Lemma~86; the step uses Lemma~86 (p.~3550) \\
3566 & l.~11 & \textit{for} $uRv \to uRv$ \textit{read} $vRu \to uRv$ \\
\dag\,3566 & proof of Lemma~248, l.~17 & \textit{for} Definition~2.2 of Volume~I \textit{read} Definition~2.3 of Volume~I \\
\dag\,3566 & proof of Lemma~242, l.~19 & \textit{for} Lemma~80 \textit{read} Lemma~79 \\
3567 & l.~1 & \textit{after} accessible \textit{insert} from $w'$ \\
\dag\,3567 & proof of Lemma~258, l.~2 & \textit{for} Lemma~3.2 of Volume~I \textit{read} Lemma~3.3 of Volume~I \\
3567 & l.~23 & \textit{delete} the second occurrence of \emph{is} \\
\dag\,3567 & proof of Lemma~258, l.~25 & \textit{for} $w'Rv' \vee v'Rw_1 \to w'Rw_1$ \textit{read} $w'Rv' \wedge v'Rw_1 \to w'Rw_1$ \\
\dag\,3568 & proof of Lemma~267, l.~6 & \textit{for} Lemma~256 \textit{read} Lemma~261; the step uses Lemma~261 (p.~3568) \\
\dag\,3568 & proof of Lemma~261, l.~18 & \textit{for} Lemma~262 \textit{read} Lemma~259; the step uses Lemma~259 (p.~3568) \\
3568 & l.~19 & \textit{for} p.~3556 \textit{read} p.~3566 (Lemma~239) \\
\dag\,3568 & proof of Lemma~268, l.~28 & \textit{for} $\pos(q \wedge r) \leftrightarrow (\pos q \vee \pos r)$ \textit{read} $\pos(q \vee r) \leftrightarrow (\pos q \vee \pos r)$ \\
3569 & l.~35 & \textit{for} p.~3557 \textit{read} p.~3567 (Lemma~252) \\
\dag\,3570 & proof of Lemma~285, l.~25 & \textit{for} Lemma~274 \textit{read} Lemma~277 \\
\dag\,3570 & proof of Lemma~288, l.~8 from foot & \textit{for} $\nec(p \vee q) \leftrightarrow (\nec p \wedge \nec q)$ \textit{read} $\nec(p \wedge q) \leftrightarrow (\nec p \wedge \nec q)$ \\
\dag\,3571 & proof of Lemma~294, l.~5 & \textit{for} Lemma~191 \textit{read} Lemma~194 \\
\dag\,3571 & proof of Lemma~286, l.~8 & \textit{for} Lemma~2 \textit{read} Lemma~3; the step uses Lemma~3 (p.~3542) \\
\dag\,3571 & proof of Lemma~285, l.~9 & \textit{for} Lemma~17 \textit{read} Lemma~12 \\
3571 & l.~11 & \textit{delete} the second occurrence of \emph{is} \\
\dag\,3571 & proof of Lemma~297, l.~26 & \textit{for} Lemma~152 \textit{read} Lemma~154 \\
\dag\,3572 & proof of Lemma~306, l.~18 & \textit{for} Lemma~178 \textit{read} Lemma~181 \\
3572 & l.~26 & \textit{for} p.~3562 \textit{read} p.~3563 (Lemma~214) \\
\dag\,3572 & proof of Lemma~307, (1.18) & \textit{for} $\pos q \wedge \pos r \to \pos(q \wedge r)$ \textit{read} $\nec q \wedge \pos r \to \pos(q \wedge r)$; correctly printed in (1.20), p.~3578 \\
\dag\,3574 & proof of Lemma~327, l.~12 & \textit{for} Lemma~248 \textit{read} Lemma~253; the step uses Lemma~253 (p.~3567) \\
3574 & l.~35 & \textit{for} p.~3560 \textit{read} p.~3570 (Lemma~286) \\
\dag\,3576 & proof of Lemma~342, l.~7 & \textit{for} Lemma~216 \textit{read} Lemma~218; the step uses Lemma~218 (p.~3563) \\
\dag\,3576 & proof of Lemma~344, l.~11 & \textit{for} Lemma~167 \textit{read} Lemma~170 \\
\dag\,3576 & proof of Lemma~334, l.~23 & \textit{for} Lemma~319 \textit{read} Lemma~316 \\
\dag\,3577 & proof of Lemma~347, l.~17 & \textit{for} Lemma~281 \textit{read} Lemma~278; the step uses Lemma~278 (p.~3569) \\
3577 & l.~23 & \textit{for} p.~3562 \textit{read} p.~3552 (Lemma~102) \\
3578 & l.~2 & \textit{for} p.~3543 \textit{read} p.~3553 (Lemma~114) \\
\dag\,3578 & proof of Lemma~367, (1.19) & \textit{for} $v \models \pos q \iff \exists w_1\,(vRw_2 \wedge w_2 \models q)$ \textit{read} $v \models \pos q \iff \exists w_2\,(vRw_2 \wedge w_2 \models q)$ \\
3579 & l.~8 & \textit{for} p.~3578 \textit{read} p.~3577 (Lemma~354) \\
\dag\,3580 & proof of Lemma~386, l.~3 from foot & \textit{for} $w \models \pos q \iff \exists w_2\,(w_2Rw \wedge w_2 \models q)$ \textit{read} $w \models \pos q \iff \exists w_2\,(wRw_2 \wedge w_2 \models q)$ \\
3581 & l.~4 & \textit{for} $uRv'$ by transitivity \textit{read} $uRv'$ by symmetry \\
\dag\,3581 & proof of Lemma~385, l.~22 & \textit{for} Lemma~59 \textit{read} Lemma~60; the step uses Lemma~60 (p.~3547) \\
3581 & l.~8 from foot & \textit{for} relfexivity \textit{read} reflexivity \\
3581 & l.~4 from foot & \textit{for} p.~3564 \textit{read} p.~3563 (Lemma~211) \\
\end{errlist}

\subsection*{Chapter~2. Formal material for Volume I \hfill \normalfont\small pp.~3582--3627}
\begin{errlist}
\dag\,3583 & proof of Lemma~392, l.~5 & \textit{for} Lemma~394 \textit{read} Lemma~391; the step uses Lemma~391 (p.~3583) \\
\dag\,3583 & proof of Lemma~390, l.~16 & \textit{for} Lemma~388 \textit{read} Lemma~389; the step uses Lemma~389 (p.~3583) \\
3583 & l.~24 & \textit{for} p.~3560 \textit{read} p.~3550 (Lemma~86) \\
\dag\,3585 & proof of Lemma~414, l.~3 & \textit{for} Lemma~411 \textit{read} Lemma~413 \\
\dag\,3585 & proof of Lemma~406, l.~11 & \textit{for} Lemma~388 \textit{read} Lemma~390 \\
\dag\,3585 & proof of Lemma~405, l.~16 & \textit{for} Lemma~392 \textit{read} Lemma~389; the step uses Lemma~389 (p.~3583) \\
\dag\,3585 & proof of Lemma~414, (2.1) & \textit{for} $\mathrm{NME}(x,w_0) \iff \forall w'\,(w_0Rw' \leftrightarrow \mathrm{ME}(x,w'))$ \textit{read} $\mathrm{NME}(x,w_0) \iff \forall w'\,(w_0Rw' \to \mathrm{ME}(x,w'))$ \\
\dag\,3586 & proof of Lemma~416, l.~8 & \textit{for} Lemma~397 \textit{read} Lemma~400; the step uses Lemma~400 (p.~3584) \\
\dag\,3586 & proof of Lemma~422, l.~14 & \textit{for} Lemma~409 \textit{read} Lemma~410 \\
\dag\,3586 & proof of Lemma~423, (2.3) & \textit{for} $\exists y\,\mathrm{MG}(y,v') \to \exists y\,(E(y,v') \vee \mathrm{ME}(y,v'))$ \textit{read} $\exists y\,\mathrm{MG}(y,v') \to \exists y\,(E(y,v') \wedge \mathrm{ME}(y,v'))$ \\
\dag\,3586 & proof of Lemma~408, l.~27 & \textit{for} Lemma~406 \textit{read} Lemma~407 \\
\dag\,3586 & proof of Lemma~423, (2.2) & \textit{for} $w' \models \nec r \iff \forall w\,(w'Rw \leftrightarrow w \models r)$ \textit{read} $w' \models \nec r \iff \forall w\,(w'Rw \to w \models r)$; the formula is used in its correct form on p.~3588 \\
\dag\,3587 & proof of Lemma~425, l.~5 & \textit{for} Lemma~401 \textit{read} Lemma~404 \\
\dag\,3587 & proof of Lemma~432, (2.4) & \textit{for} $v' \models \nec r \iff \forall v\,(vRv' \to v \models r)$ \textit{read} $v' \models \nec r \iff \forall v\,(v'Rv \to v \models r)$ \\
\dag\,3587 & proof of Lemma~417, l.~26 & \textit{for} Lemma~410 \textit{read} Lemma~405; the step uses Lemma~405 (p.~3584) \\
\dag\,3588 & proof of Lemma~442, (2.5) & \textit{for} $v'Rw \leftrightarrow wRv'$ \textit{read} $v'Rw \to wRv'$ \\
\dag\,3588 & proof of Lemma~434, l.~29 & \textit{for} Lemma~397 \textit{read} Lemma~392; the step uses Lemma~392 (p.~3583) \\
\dag\,3589 & proof of Lemma~434, l.~17 & \textit{for} Lemma~406 \textit{read} Lemma~409 \\
3589 & l.~22 & \textit{delete} the second occurrence of \emph{that} \\
\dag\,3589 & proof of Lemma~451, l.~1 from foot & \textit{for} $w' \models \pos q \iff \exists w\,(w'Rw \vee w \models q)$ \textit{read} $w' \models \pos q \iff \exists w\,(w'Rw \wedge w \models q)$ \\
\dag\,3589 & proof of Lemma~451, (2.6) & \textit{for} $\mathrm{NE}(z,u) \iff \exists w_1\,(uRw_1 \to E(z,w_1))$ \textit{read} $\mathrm{NE}(z,u) \iff \forall w_1\,(uRw_1 \to E(z,w_1))$; the formula is used in its correct form on p.~3593 \\
\dag\,3590 & proof of Lemma~454, l.~2 & \textit{for} Lemma~440 \textit{read} Lemma~445 \\
\dag\,3590 & proof of Lemma~451, l.~23 & \textit{for} Lemma~445 \textit{read} Lemma~446 \\
\dag\,3590 & proof of Lemma~460, (2.7) & \textit{for} $\mathrm{MG}(x,v) \iff \forall w\,(vRw \to E(x,w) \vee \mathrm{ME}(x,w))$ \textit{read} $\mathrm{MG}(x,v) \iff \forall w\,(vRw \to E(x,w) \wedge \mathrm{ME}(x,w))$; the formula is used in its correct form on p.~3593 \\
\dag\,3590 & proof of Lemma~454, l.~24 & \textit{for} Lemma~444 \textit{read} Lemma~441; the step uses Lemma~441 (p.~3588) \\
\dag\,3591 & proof of Lemma~469, (2.8) & \textit{for} $\mathrm{NME}(y,w_0) \iff \exists w'\,(w_0Rw' \to \mathrm{ME}(y,w'))$ \textit{read} $\mathrm{NME}(y,w_0) \iff \forall w'\,(w_0Rw' \to \mathrm{ME}(y,w'))$ \\
\dag\,3592 & proof of Lemma~472, l.~27 & \textit{for} Lemma~451 \textit{read} Lemma~454; the step uses Lemma~454 (p.~3590) \\
\dag\,3593 & proof of Lemma~480, l.~9 & \textit{for} Lemma~410 \textit{read} Lemma~406 \\
3593 & l.~38 & \textit{delete} the second occurrence of \emph{that} \\
\dag\,3593 & proof of Lemma~487, (2.9) & \textit{for} $\pos(q \vee p) \leftrightarrow (\nec q \vee \pos p)$ \textit{read} $\pos(q \vee p) \leftrightarrow (\pos q \vee \pos p)$ \\
\dag\,3595 & proof of Lemma~499, l.~2 & \textit{for} Lemma~487 \textit{read} Lemma~491 \\
\dag\,3595 & proof of Lemma~506, l.~5 & \textit{for} $\mathrm{NE}(y,w_2) \iff \exists u\,(w_2Ru \to E(y,u))$ \textit{read} $\mathrm{NE}(y,w_2) \iff \forall u\,(w_2Ru \to E(y,u))$ \\
\dag\,3595 & proof of Lemma~506, l.~16 & \textit{for} $\neg(\pos G \wedge \nec\neg G)$ \textit{read} $\neg(\pos G \wedge \pos\neg G)$ \\
\dag\,3595 & proof of Lemma~496, l.~27 & \textit{for} Lemma~471 \textit{read} Lemma~472; the step uses Lemma~472 (p.~3592) \\
\dag\,3596 & proof of Lemma~512, l.~9 & \textit{for} Lemma~407 \textit{read} Lemma~402 \\
\dag\,3596 & proof of Lemma~515, (2.10) & \textit{for} $\pos\neg\exists y\,\mathrm{MG}(y) \leftrightarrow \neg\exists y\,\mathrm{MG}(y)$ \textit{read} $\pos\neg\exists y\,\mathrm{MG}(y) \to \neg\exists y\,\mathrm{MG}(y)$; the formula is used in its correct form on p.~3600 \\
\dag\,3596 & proof of Lemma~511, l.~23 & \textit{for} Lemma~465 \textit{read} Lemma~470 \\
\dag\,3596 & proof of Lemma~502, l.~25 & \textit{for} Lemma~427 \textit{read} Lemma~428 \\
\dag\,3596 & proof of Lemma~498, l.~30 & \textit{for} Lemma~460 \textit{read} Lemma~465; the step uses Lemma~465 (p.~3591) \\
3597 & l.~13 & \textit{delete} the second occurrence of \emph{the} \\
\dag\,3597 & proof of Lemma~511, l.~14 & \textit{for} Lemma~509 \textit{read} Lemma~510; the step uses Lemma~510 (p.~3596) \\
\dag\,3597 & proof of Lemma~524, l.~32 & \textit{for} $\nec(p \to p) \to (\nec p \to \nec q)$ \textit{read} $\nec(p \to q) \to (\nec p \to \nec q)$ \\
3598 & l.~27 & \textit{for} omnsicient \textit{read} omniscient \\
\dag\,3599 & proof of Lemma~543, (2.11) & \textit{for} $\neg\pos q \leftrightarrow \nec q$ \textit{read} $\neg\pos q \leftrightarrow \nec\neg q$ \\
\dag\,3600 & proof of Lemma~552, l.~2 & \textit{for} Assumption~3.1 of Volume~I \textit{read} Assumption~2.1 of Volume~I \\
\dag\,3600 & proof of Lemma~547, l.~3 & \textit{for} Lemma~444 \textit{read} Lemma~440 \\
3600 & l.~13 & \textit{for} p.~3548 \textit{read} p.~3549 (Lemma~72) \\
\dag\,3600 & proof of Lemma~549, l.~21 & \textit{for} Lemma~483 \textit{read} Lemma~481 \\
\dag\,3601 & proof of Lemma~554, l.~1 & \textit{for} Lemma~399 \textit{read} Lemma~401; the step uses Lemma~401 (p.~3584) \\
\dag\,3601 & proof of Lemma~553, l.~10 & \textit{for} Lemma~529 \textit{read} Lemma~526; the step uses Lemma~526 (p.~3598) \\
\dag\,3601 & proof of Lemma~561, (2.13) & \textit{for} $v \models \pos q \iff \exists w_0\,(vRw_0 \vee w_0 \models q)$ \textit{read} $v \models \pos q \iff \exists w_0\,(vRw_0 \wedge w_0 \models q)$; the formula is used in its correct form on p.~3605 \\
\dag\,3601 & proof of Lemma~561, (2.12) & \textit{for} $\nec p \wedge \nec r \to \pos(p \wedge r)$ \textit{read} $\nec p \wedge \pos r \to \pos(p \wedge r)$ \\
3602 & l.~18 & \textit{for} moarlly \textit{read} morally \\
\dag\,3603 & proof of Lemma~579, (2.14) & \textit{for} $\pos\exists y\,\mathrm{MG}(y) \leftrightarrow \nec\exists y\,\mathrm{MG}(y)$ \textit{read} $\pos\exists y\,\mathrm{MG}(y) \to \nec\exists y\,\mathrm{MG}(y)$ \\
3603 & l.~8 & \textit{for} p.~3699 \textit{read} p.~3599 (Lemma~535) \\
\dag\,3603 & proof of Lemma~579, (2.15) & \textit{for} $\nec r \leftrightarrow \nec\nec r$ \textit{read} $\nec r \to \nec\nec r$; the formula is used in its correct form on p.~3605 \\
\dag\,3604 & proof of Lemma~586, l.~17 & \textit{for} Lemma~481 \textit{read} Lemma~479 \\
\dag\,3604 & proof of Lemma~588, (2.17) & \textit{for} $vRw \vee wRv' \to vRv'$ \textit{read} $vRw \wedge wRv' \to vRv'$ \\
\dag\,3604 & proof of Lemma~588, (2.18) & \textit{for} $\mathrm{MG}(y,u) \to \mathrm{NE}(y,u) \vee \mathrm{NME}(y,u)$ \textit{read} $\mathrm{MG}(y,u) \to \mathrm{NE}(y,u) \wedge \mathrm{NME}(y,u)$ \\
\dag\,3604 & proof of Lemma~588, (2.16) & \textit{for} $@ \in \nec\,\exists y\, \mathrm{MG}(y,\cdot)$ \textit{read} $@ \in \pos\,\exists y\, \mathrm{MG}(y,\cdot)$ \\
\dag\,3605 & proof of Lemma~581, l.~18 & \textit{for} Lemma~461 \textit{read} Lemma~464 \\
\dag\,3606 & proof of Lemma~599, l.~5 & \textit{for} Lemma~590 \textit{read} Lemma~593; the step uses Lemma~593 (p.~3605) \\
3606 & l.~20 & \textit{for} ominscient \textit{read} omniscient \\
\dag\,3606 & proof of Lemma~607, (2.19) & \textit{for} $\mathrm{MG}(y,w') \wedge vRw' \to \mathrm{MG}(y,v)$ \textit{read} $\mathrm{MG}(y,w') \wedge w'Rv \to \mathrm{MG}(y,v)$ \\
3607 & l.~10 & \textit{for} p.~3556 \textit{read} p.~3546 (Lemma~47) \\
3607 & l.~34 & \textit{for} p.~3566 \textit{read} p.~3576 (Lemma~344) \\
3608 & l.~12 & \textit{for} Definition~2.4 of Volume~I \textit{read} Definition~2.3 of Volume~I \\
\dag\,3608 & proof of Lemma~617, l.~28 & \textit{for} Lemma~559 \textit{read} Lemma~560; the step uses Lemma~560 (p.~3601) \\
\dag\,3609 & proof of Lemma~634, l.~9 & \textit{for} Lemma~497 \textit{read} Lemma~500; the step uses Lemma~500 (p.~3595) \\
\dag\,3609 & proof of Lemma~631, l.~16 & \textit{for} Lemma~422 \textit{read} Lemma~417 \\
3610 & l.~2 & \textit{for} exists and is maximally excellent at some world accessible from $w_1$ \textit{read} exists and is maximally excellent at every world accessible from $w_1$ \\
\dag\,3610 & proof of Lemma~644, l.~26 & \textit{for} $\nec q \wedge \pos r \to \pos(q \vee r)$ \textit{read} $\nec q \wedge \pos r \to \pos(q \wedge r)$ \\
3610 & l.~4 from foot & \textit{for} $\pos p \to \pos\pos p$ \textit{read} $\pos p \to \nec\pos p$ \\
\dag\,3611 & proof of Lemma~653, l.~1 & \textit{for} Lemma~576 \textit{read} Lemma~577 \\
\dag\,3611 & proof of Lemma~649, l.~3 & \textit{for} Lemma~446 \textit{read} Lemma~442 \\
\dag\,3611 & proof of Lemma~650, l.~19 & \textit{for} Lemma~424 \textit{read} Lemma~419 \\
3611 & l.~28 & \textit{for} p.~3617 \textit{read} p.~3607 (Lemma~609) \\
3612 & l.~7 & \textit{for} omnisicent \textit{read} omniscient \\
\dag\,3612 & proof of Lemma~656, l.~12 & \textit{for} Lemma~541 \textit{read} Lemma~544 \\
\dag\,3612 & proof of Lemma~661, l.~29 & \textit{for} Lemma~636 \textit{read} Lemma~640 \\
\dag\,3612 & proof of Lemma~662, (2.20) & \textit{for} $\pos\pos p \to \nec p$ \textit{read} $\pos\pos p \to \pos p$ \\
\dag\,3613 & proof of Lemma~659, l.~21 & \textit{for} Lemma~484 \textit{read} Lemma~489 \\
\dag\,3614 & proof of Lemma~680, (2.21) & \textit{for} $\mathrm{MG}(z,w_0) \to \mathrm{NE}(z,w_0) \vee \mathrm{NME}(z,w_0)$ \textit{read} $\mathrm{MG}(z,w_0) \to \mathrm{NE}(z,w_0) \wedge \mathrm{NME}(z,w_0)$ \\
\dag\,3614 & proof of Lemma~667, l.~27 & \textit{for} Lemma~575 \textit{read} Lemma~579; the step uses Lemma~579 (p.~3603) \\
\dag\,3614 & proof of Lemma~668, l.~28 & \textit{for} Lemma~646 \textit{read} Lemma~649 \\
3615 & l.~6 & \textit{for} p.~3711 \textit{read} p.~3611 (Lemma~646) \\
3616 & l.~8 from foot & \textit{for} neecssarily \textit{read} necessarily \\
\dag\,3617 & proof of Lemma~708, l.~8 & \textit{for} Lemma~564 \textit{read} Lemma~565; the step uses Lemma~565 (p.~3602) \\
\dag\,3618 & proof of Lemma~717, (2.22) & \textit{for} $\nec(q \wedge q) \leftrightarrow (\nec r \wedge \nec q)$ \textit{read} $\nec(r \wedge q) \leftrightarrow (\nec r \wedge \nec q)$; the formula is used in its correct form on p.~3619 \\
\dag\,3618 & proof of Lemma~710, l.~24 & \textit{for} Lemma~505 \textit{read} Lemma~502 \\
\dag\,3619 & proof of Lemma~713, l.~7 & \textit{for} Lemma~700 \textit{read} Lemma~701; the step uses Lemma~701 (p.~3617) \\
3620 & l.~1 & \textit{for} p.~3617 \textit{read} p.~3618 (Lemma~709) \\
\dag\,3620 & proof of Lemma~729, l.~4 & \textit{for} Lemma~578 \textit{read} Lemma~581; the step uses Lemma~581 (p.~3604) \\
\dag\,3620 & proof of Lemma~735, (2.23) & \textit{for} $w_0Rw_2 \wedge w_2Ru \leftrightarrow w_0Ru$ \textit{read} $w_0Rw_2 \wedge w_2Ru \to w_0Ru$; the formula is used in its correct form on p.~3624 \\
\dag\,3620 & proof of Lemma~734, l.~13 & \textit{for} Lemma~679 \textit{read} Lemma~680 \\
\dag\,3621 & proof of Lemma~732, l.~4 & \textit{for} Lemma~691 \textit{read} Lemma~693 \\
\dag\,3621 & proof of Lemma~731, l.~5 & \textit{for} Lemma~519 \textit{read} Lemma~523 \\
3621 & l.~14 & \textit{for} p.~3573 \textit{read} p.~3572 (Lemma~302) \\
\dag\,3623 & proof of Lemma~760, l.~27 & \textit{for} Lemma~612 \textit{read} Lemma~609; the step uses Lemma~609 (p.~3607) \\
3623 & l.~28 & \textit{delete} the second occurrence of \emph{that} \\
3623 & l.~2 from foot & \textit{for} p.~3581 \textit{read} p.~3580 (Lemma~385) \\
\dag\,3624 & proof of Lemma~766, l.~11 & \textit{for} Lemma~688 \textit{read} Lemma~690 \\
\dag\,3624 & proof of Lemma~765, l.~12 & \textit{for} Lemma~679 \textit{read} Lemma~676; the step uses Lemma~676 (p.~3614) \\
\dag\,3624 & proof of Lemma~772, l.~20 & \textit{for} $\nec\exists y\,\mathrm{MG}(y) \to \nec\exists y\,\mathrm{MG}(y)$ \textit{read} $\pos\exists y\,\mathrm{MG}(y) \to \nec\exists y\,\mathrm{MG}(y)$ \\
\dag\,3624 & proof of Lemma~757, l.~27 & \textit{for} Lemma~469 \textit{read} Lemma~474 \\
\dag\,3624 & proof of Lemma~770, l.~28 & \textit{for} Lemma~450 \textit{read} Lemma~451; the step uses Lemma~451 (p.~3589) \\
3624 & l.~3 from foot & \textit{for} $\mathrm{MG}(x,v) \wedge w_0Rv \to \mathrm{MG}(x,w_0)$ \textit{read} $\mathrm{MG}(x,v) \wedge vRw_0 \to \mathrm{MG}(x,w_0)$ \\
3625 & l.~31 & \textit{delete} the second occurrence of \emph{to} \\
3626 & l.~8 & \textit{after} maximally excellent \textit{insert} at $v'$ \\
\dag\,3626 & proof of Lemma~788, l.~10 & \textit{for} Lemma~641 \textit{read} Lemma~640 \\
3626 & l.~30 & \textit{for} p.~3724 \textit{read} p.~3624 (Lemma~766) \\
3627 & l.~9 & \textit{for} $\mathrm{MG}(x,w) \iff \forall u\,(wRu \leftrightarrow E(x,u) \wedge \mathrm{ME}(x,u))$ \textit{read} $\mathrm{MG}(x,w) \iff \forall u\,(wRu \to E(x,u) \wedge \mathrm{ME}(x,u))$ \\
3627 & l.~19 & \textit{delete} the second occurrence of \emph{a} \\
\end{errlist}

\subsection*{Chapter~3. Higher-order modal logic and property abstraction \hfill \normalfont\small pp.~3628--3682}
\begin{errlist}
3629 & l.~16 & \textit{for} p.~3720 \textit{read} p.~3620 (Lemma~728) \\
3629 & l.~30 & \textit{for} gdolike \textit{read} godlike \\
\dag\,3630 & proof of Lemma~794, l.~21 & \textit{for} Lemma~791 \textit{read} Lemma~792 \\
\dag\,3631 & proof of Lemma~814, l.~4 & \textit{for} Lemma~802 \textit{read} Lemma~803; the step uses Lemma~803 (p.~3630) \\
\dag\,3631 & proof of Lemma~819, (3.1) & \textit{for} $\mathcal{P}_0 \supseteq \mathcal{P} \Rightarrow \exists x\,\exists w'\,\forall {\chi} \in \mathcal{P}_0\; {\chi}(x,w')$ \textit{read} $\mathcal{P}_0 \subseteq \mathcal{P} \Rightarrow \exists x\,\exists w'\,\forall {\chi} \in \mathcal{P}_0\; {\chi}(x,w')$ \\
\dag\,3631 & proof of Lemma~819, (3.2) & \textit{for} $\nec\forall z\,({\chi}z \to {\phi}z) \leftrightarrow ({\chi} \sqcap {\phi} = {\chi})$ \textit{read} $\nec\forall z\,({\chi}z \to {\phi}z) \to ({\chi} \sqcap {\phi} = {\chi})$ \\
\dag\,3632 & proof of Lemma~818, l.~7 & \textit{for} Lemma~807 \textit{read} Lemma~810; the step uses Lemma~810 (p.~3630) \\
\dag\,3632 & proof of Lemma~814, l.~9 & \textit{for} Lemma~796 \textit{read} Lemma~793 \\
\dag\,3632 & proof of Lemma~824, l.~13 & \textit{for} Lemma~805 \textit{read} Lemma~807; the step uses Lemma~807 (p.~3630) \\
3633 & l.~4 & \textit{after} great-making properties \textit{insert} jointly \\
\dag\,3633 & proof of Lemma~826, l.~25 & \textit{for} Lemma~807 \textit{read} Lemma~809; the step uses Lemma~809 (p.~3630) \\
3634 & l.~15 & \textit{for} generral model \textit{read} general model \\
\dag\,3634 & proof of Lemma~844, l.~27 & \textit{for} Lemma~804 \textit{read} Lemma~806 \\
\dag\,3635 & proof of Lemma~855, l.~10 & \textit{for} Lemma~828 \textit{read} Lemma~830 \\
\dag\,3635 & proof of Lemma~842, l.~14 & \textit{for} Lemma~829 \textit{read} Lemma~825 \\
3635 & l.~19 & \textit{for} tpye \textit{read} type \\
\dag\,3635 & proof of Lemma~850, l.~30 & \textit{for} Lemma~794 \textit{read} Lemma~797 \\
\dag\,3636 & proof of Lemma~862, l.~17 & \textit{for} Lemma~830 \textit{read} Lemma~827; the step uses Lemma~827 (p.~3632) \\
3637 & l.~4 & \textit{for} Fitting (2002) \textit{read} Henkin (1950) (general models of type theory) \\
\dag\,3637 & proof of Lemma~868, l.~8 & \textit{for} Lemma~863 \textit{read} Lemma~867 \\
\dag\,3637 & proof of Lemma~860, l.~18 & \textit{for} Lemma~839 \textit{read} Lemma~840 \\
\dag\,3638 & proof of Lemma~883, l.~9 & \textit{for} Lemma~880 \textit{read} Lemma~877; the step uses Lemma~877 (p.~3637) \\
\dag\,3638 & proof of Lemma~873, l.~19 & \textit{for} Lemma~875 \textit{read} Lemma~872 \\
\dag\,3638 & proof of Lemma~884, l.~23 & \textit{for} Lemma~864 \textit{read} Lemma~860 \\
\dag\,3639 & proof of Lemma~891, l.~7 & \textit{for} Lemma~885 \textit{read} Lemma~890 \\
\dag\,3640 & proof of Lemma~906, (3.6) & \textit{for} ${\phi} \in \mathcal{P} \wedge \nec\forall z\,({\phi}z \to {\psi}z) \Rightarrow {\psi} \notin \mathcal{P}$ \textit{read} ${\phi} \in \mathcal{P} \wedge \nec\forall z\,({\phi}z \to {\psi}z) \Rightarrow {\psi} \in \mathcal{P}$; correctly printed in (3.7), p.~3644 \\
3641 & l.~24 & \textit{for} p.~3585 \textit{read} p.~3558 (Lemma~161) \\
\dag\,3642 & proof of Lemma~926, l.~13 & \textit{for} ${\phi} \sqcap {\chi} = \lambda z.\,{\phi}z \vee {\chi}z$ \textit{read} ${\phi} \sqcap {\chi} = \lambda z.\,{\phi}z \wedge {\chi}z$ \\
3643 & l.~2 & \textit{for} p.~3631 \textit{read} p.~3632 (Lemma~827) \\
\dag\,3643 & proof of Lemma~926, l.~22 & \textit{for} Lemma~906 \textit{read} Lemma~911; the step uses Lemma~911 (p.~3641) \\
3643 & l.~27 & \textit{for} p.~3644 \textit{read} p.~3544 (Lemma~22) \\
\dag\,3643 & proof of Lemma~923, l.~30 & \textit{for} Lemma~855 \textit{read} Lemma~856; the step uses Lemma~856 (p.~3635) \\
3644 & l.~9 & \textit{for} $(\lambda z.\,{\chi}z \wedge {\phi}z)(a) \to {\chi}a \wedge {\phi}a$ \textit{read} $(\lambda z.\,{\chi}z \wedge {\phi}z)(a) \leftrightarrow {\chi}a \wedge {\phi}a$ \\
\dag\,3644 & proof of Lemma~944, l.~18 & \textit{for} Lemma~916 \textit{read} Lemma~920 \\
3644 & l.~23 & \textit{for} p.~3548 \textit{read} p.~3584 (Lemma~400) \\
\dag\,3644 & proof of Lemma~938, l.~25 & \textit{for} Lemma~895 \textit{read} Lemma~891 \\
\dag\,3645 & proof of Lemma~955, (3.8) & \textit{for} $\nec\exists x\,({\psi}x \to {\chi}x) \to ({\psi} \sqcap {\chi} = {\psi})$ \textit{read} $\nec\forall x\,({\psi}x \to {\chi}x) \to ({\psi} \sqcap {\chi} = {\psi})$; the formula is used in its correct form on p.~3647 \\
3646 & l.~8 & \textit{delete} the second occurrence of \emph{to} \\
3646 & l.~9 & \textit{for} extensional \textit{read} intensional \\
\dag\,3646 & proof of Lemma~951, l.~11 & \textit{for} Lemma~831 \textit{read} Lemma~828 \\
3646 & l.~18 & \textit{for} p.~3722 \textit{read} p.~3622 (Lemma~749) \\
\dag\,3648 & proof of Lemma~984, (3.9) & \textit{for} $\mathrm{Pos}({\chi}) \to \pos\forall x\,{\chi}x$ \textit{read} $\mathrm{Pos}({\chi}) \to \pos\exists x\,{\chi}x$; the formula is used in its correct form on p.~3649 \\
\dag\,3648 & proof of Lemma~979, l.~13 & \textit{for} Lemma~969 \textit{read} Lemma~974; the step uses Lemma~974 (p.~3647) \\
3648 & l.~21 & \textit{for} p.~3567 \textit{read} p.~3566 (Lemma~242) \\
\dag\,3648 & proof of Lemma~980, l.~28 & \textit{for} Lemma~853 \textit{read} Lemma~856; the step uses Lemma~856 (p.~3635) \\
3648 & l.~39 & \textit{for} psoitivity \textit{read} positivity \\
\dag\,3652 & proof of Lemma~1011, l.~20 & \textit{for} Lemma~892 \textit{read} Lemma~891; the step uses Lemma~891 (p.~3639) \\
\dag\,3652 & proof of Lemma~1008, l.~25 & \textit{for} Lemma~897 \textit{read} Lemma~901 \\
3652 & l.~36 & \textit{for} p.~3543 \textit{read} p.~3553 (Lemma~116) \\
\dag\,3653 & proof of Lemma~1029, l.~1 & \textit{for} Lemma~801 \textit{read} Lemma~799; the step uses Lemma~799 (p.~3629) \\
\dag\,3653 & proof of Lemma~1029, l.~14 & \textit{for} Lemma~819 \textit{read} Lemma~815 \\
\dag\,3653 & proof of Lemma~1032, (3.13) & \textit{for} $\forall {\phi}\,\nec\forall y\,({\phi}y \leftrightarrow {\chi}y)$ \textit{read} $\exists {\phi}\,\nec\forall y\,({\phi}y \leftrightarrow {\chi}y)$ \\
\dag\,3654 & proof of Lemma~1033, l.~8 & \textit{for} Lemma~828 \textit{read} Lemma~825 \\
3654 & l.~10 & \textit{for} p.~3747 \textit{read} p.~3647 (Lemma~967) \\
\dag\,3654 & proof of Lemma~1035, l.~27 & \textit{for} Lemma~877 \textit{read} Lemma~876; the step uses Lemma~876 (p.~3637) \\
\dag\,3656 & proof of Lemma~1044, l.~23 & \textit{for} Lemma~832 \textit{read} Lemma~833; the step uses Lemma~833 (p.~3633) \\
3657 & l.~32 & \textit{for} Hnekin (1950) \textit{read} Henkin (1950) \\
\dag\,3657 & proof of Lemma~1071, (3.16) & \textit{for} $\mathrm{G}z \leftrightarrow \exists {\psi}\,(\mathrm{Pos}({\psi}) \to {\psi}z)$ \textit{read} $\mathrm{G}z \leftrightarrow \forall {\psi}\,(\mathrm{Pos}({\psi}) \to {\psi}z)$; correctly printed in (3.17), p.~3659 \\
\dag\,3658 & proof of Lemma~1073, l.~2 & \textit{for} Lemma~1006 \textit{read} Lemma~1011 \\
\dag\,3658 & proof of Lemma~1065, l.~4 & \textit{for} Lemma~992 \textit{read} Lemma~991 \\
\dag\,3658 & proof of Lemma~1063, l.~10 & \textit{for} Lemma~818 \textit{read} Lemma~816 \\
\dag\,3658 & proof of Lemma~1077, l.~13 & \textit{for} Lemma~967 \textit{read} Lemma~966 \\
3658 & l.~15 & \textit{for} posittivity \textit{read} positivity \\
\dag\,3658 & proof of Lemma~1064, l.~22 & \textit{for} Lemma~898 \textit{read} Lemma~903; the step uses Lemma~903 (p.~3640) \\
\dag\,3658 & proof of Lemma~1080, l.~26 & \textit{for} $\mathrm{G}z \to \forall {\chi}\,(\mathrm{Pos}({\chi}) \to {\chi}z)$ \textit{read} $\mathrm{G}z \leftrightarrow \forall {\chi}\,(\mathrm{Pos}({\chi}) \to {\chi}z)$ \\
\dag\,3658 & proof of Lemma~1074, l.~28 & \textit{for} Lemma~841 \textit{read} Lemma~837 \\
\dag\,3658 & proof of Lemma~1062, l.~29 & \textit{for} Lemma~806 \textit{read} Lemma~802; the step uses Lemma~802 (p.~3630) \\
\dag\,3659 & proof of Lemma~1087, l.~25 & \textit{for} Lemma~1089 \textit{read} Lemma~1085; the step uses Lemma~1085 (p.~3659) \\
3659 & l.~26 & \textit{delete} the second occurrence of \emph{of} \\
\dag\,3660 & proof of Lemma~1092, l.~9 & \textit{for} Lemma~1037 \textit{read} Lemma~1036 \\
\dag\,3660 & proof of Lemma~1095, l.~14 & \textit{for} Lemma~1048 \textit{read} Lemma~1044 \\
\dag\,3660 & proof of Lemma~1096, l.~23 & \textit{for} Lemma~909 \textit{read} Lemma~906; the step uses Lemma~906 (p.~3640) \\
3660 & l.~40 & \textit{for} joinlty \textit{read} jointly \\
\dag\,3661 & proof of Lemma~1109, (3.18) & \textit{for} $\mathrm{Pos}({\chi}) \vee \mathrm{Pos}({\phi}) \to \mathrm{Pos}({\chi} \sqcap {\phi})$ \textit{read} $\mathrm{Pos}({\chi}) \wedge \mathrm{Pos}({\phi}) \to \mathrm{Pos}({\chi} \sqcap {\phi})$; the formula is used in its correct form on p.~3664 \\
\dag\,3661 & proof of Lemma~1101, l.~10 & \textit{for} Lemma~1040 \textit{read} Lemma~1041; the step uses Lemma~1041 (p.~3654) \\
\dag\,3661 & proof of Lemma~1094, l.~26 & \textit{for} Lemma~1089 \textit{read} Lemma~1091; the step uses Lemma~1091 (p.~3660) \\
3662 & l.~14 & \textit{for} Henikn \textit{read} Henkin \\
\dag\,3663 & proof of Lemma~1129, l.~15 & \textit{for} $\mathrm{Pos}({\chi}) \leftrightarrow \pos\exists x\,{\chi}x$ \textit{read} $\mathrm{Pos}({\chi}) \to \pos\exists x\,{\chi}x$ \\
\dag\,3663 & proof of Lemma~1123, l.~28 & \textit{for} Lemma~854 \textit{read} Lemma~855 \\
3664 & l.~5 & \textit{for} standard models \textit{read} general (Henkin) models \\
3664 & l.~8 & \textit{for} p.~3553 \textit{read} p.~3554 (Lemma~129) \\
\dag\,3664 & proof of Lemma~1125, l.~18 & \textit{for} Lemma~848 \textit{read} Lemma~851; the step uses Lemma~851 (p.~3635) \\
\dag\,3665 & proof of Lemma~1145, l.~22 & \textit{for} Lemma~994 \textit{read} Lemma~989 \\
3665 & l.~8 from foot & \textit{for} comprehesnion \textit{read} comprehension \\
3666 & l.~11 & \textit{for} great-mking \textit{read} great-making \\
\dag\,3666 & proof of Lemma~1142, l.~30 & \textit{for} Lemma~1054 \textit{read} Lemma~1056; the step uses Lemma~1056 (p.~3656) \\
3667 & l.~7 & \textit{for} p.~3631 \textit{read} p.~3613 (Lemma~667) \\
\dag\,3668 & proof of Lemma~1162, l.~7 & \textit{for} Lemma~1022 \textit{read} Lemma~1020; the step uses Lemma~1020 (p.~3652) \\
\dag\,3668 & proof of Lemma~1177, (3.20) & \textit{for} $\mathcal{P}_0 \subseteq \mathcal{P} \Rightarrow \exists y\,\exists u\,\forall {\phi} \notin \mathcal{P}_0\; {\phi}(y,u)$ \textit{read} $\mathcal{P}_0 \subseteq \mathcal{P} \Rightarrow \exists y\,\exists u\,\forall {\phi} \in \mathcal{P}_0\; {\phi}(y,u)$ \\
3668 & l.~23 & \textit{for} prperty \textit{read} property \\
\dag\,3669 & proof of Lemma~1181, l.~10 & \textit{for} Lemma~1044 \textit{read} Lemma~1047 \\
\dag\,3669 & proof of Lemma~1173, l.~27 & \textit{for} Lemma~1141 \textit{read} Lemma~1144 \\
\dag\,3669 & proof of Lemma~1185, l.~29 & \textit{for} Lemma~1114 \textit{read} Lemma~1118 \\
3669 & l.~32 & \textit{for} p.~3540 \textit{read} p.~3640 (Lemma~899) \\
3669 & l.~5 from foot & \textit{delete} the second occurrence of \emph{of} \\
\dag\,3670 & proof of Lemma~1180, l.~6 & \textit{for} Lemma~831 \textit{read} Lemma~834 \\
\dag\,3670 & proof of Lemma~1183, l.~10 & \textit{for} Lemma~836 \textit{read} Lemma~834; the step uses Lemma~834 (p.~3633) \\
\dag\,3670 & proof of Lemma~1196, l.~20 & \textit{for} Lemma~934 \textit{read} Lemma~930; the step uses Lemma~930 (p.~3643) \\
3671 & l.~30 & \textit{delete} the second occurrence of \emph{a} \\
\dag\,3671 & proof of Lemma~1206, (3.21) & \textit{for} $\mathrm{G}z \leftrightarrow \forall {\psi}\,(\mathrm{Pos}({\psi}) \leftrightarrow \pos {\psi}z)$ \textit{read} $\mathrm{G}z \leftrightarrow \forall {\psi}\,(\mathrm{Pos}({\psi}) \leftrightarrow \nec {\psi}z)$; the formula is used in its correct form on p.~3675 \\
\dag\,3672 & proof of Lemma~1211, l.~8 & \textit{for} Lemma~815 \textit{read} Lemma~814 \\
\dag\,3672 & proof of Lemma~1212, l.~11 & \textit{for} Lemma~933 \textit{read} Lemma~934 \\
\dag\,3672 & proof of Lemma~1216, (3.22) & \textit{for} ${\phi}, {\chi} \notin \mathcal{P} \Rightarrow {\phi} \sqcap {\chi} \in \mathcal{P}$ \textit{read} ${\phi}, {\chi} \in \mathcal{P} \Rightarrow {\phi} \sqcap {\chi} \in \mathcal{P}$ \\
3672 & l.~22 & \textit{for} p.~3665 \textit{read} p.~3666 (Lemma~1149) \\
\dag\,3673 & proof of Lemma~1222, l.~21 & \textit{for} Lemma~1205 \textit{read} Lemma~1206 \\
3673 & l.~24 & \textit{for} p.~3651 \textit{read} p.~3652 (Lemma~1018) \\
\dag\,3673 & proof of Lemma~1225, l.~25 & \textit{for} ${\phi} \sqcap {\psi} = \lambda z.\,{\phi}z \vee {\psi}z$ \textit{read} ${\phi} \sqcap {\psi} = \lambda z.\,{\phi}z \wedge {\psi}z$ \\
\dag\,3673 & proof of Lemma~1212, l.~30 & \textit{for} Lemma~831 \textit{read} Lemma~836; the step uses Lemma~836 (p.~3633) \\
3675 & l.~7 & \textit{for} p.~3606 \textit{read} p.~3660 (Lemma~1094) \\
\dag\,3675 & proof of Lemma~1227, l.~8 & \textit{for} Lemma~1029 \textit{read} Lemma~1024 \\
3675 & l.~10 & \textit{for} comprehensoin \textit{read} comprehension \\
3675 & l.~24 & \textit{after} comprehension \textit{insert} for intensional properties \\
\dag\,3675 & proof of Lemma~1243, l.~25 & \textit{for} Lemma~1041 \textit{read} Lemma~1045; the step uses Lemma~1045 (p.~3655) \\
3675 & l.~29 & \textit{for} Gallin (1974) \textit{read} Gallin (1975) \\
\dag\,3676 & proof of Lemma~1249, l.~1 & \textit{for} Lemma~1088 \textit{read} Lemma~1085 \\
3676 & l.~22 & \textit{for} hgiher-order \textit{read} higher-order \\
3676 & l.~27 & \textit{for} Henkn \textit{read} Henkin \\
\dag\,3676 & proof of Lemma~1254, (3.24) & \textit{for} $\forall {\psi}\,\nec\forall y\,({\psi}y \leftrightarrow {\chi}y)$ \textit{read} $\exists {\psi}\,\nec\forall y\,({\psi}y \leftrightarrow {\chi}y)$ \\
\dag\,3677 & proof of Lemma~1248, l.~7 & \textit{for} Lemma~982 \textit{read} Lemma~987; the step uses Lemma~987 (p.~3649) \\
3677 & l.~15 & \textit{for} Henkin (1949) \textit{read} Henkin (1950) \\
\dag\,3677 & proof of Lemma~1249, l.~20 & \textit{for} Lemma~1136 \textit{read} Lemma~1137; the step uses Lemma~1137 (p.~3664) \\
3677 & l.~29 & \textit{delete} the second occurrence of \emph{of} \\
\dag\,3678 & proof of Lemma~1259, l.~7 & \textit{for} Lemma~1087 \textit{read} Lemma~1090; the step uses Lemma~1090 (p.~3659) \\
\dag\,3678 & proof of Lemma~1259, l.~9 & \textit{for} Lemma~827 \textit{read} Lemma~828; the step uses Lemma~828 (p.~3632) \\
3678 & l.~39 & \textit{delete} the second occurrence of \emph{the} \\
\dag\,3679 & proof of Lemma~1274, l.~14 & \textit{for} Lemma~1002 \textit{read} Lemma~1004 \\
\dag\,3679 & proof of Lemma~1269, l.~16 & \textit{for} Lemma~972 \textit{read} Lemma~976 \\
3679 & l.~1 from foot & \textit{for} p.~3673 \textit{read} p.~3674 (Lemma~1228) \\
3680 & l.~23 & \textit{for} p.~3689 \textit{read} p.~3679 (Lemma~1280) \\
\dag\,3681 & proof of Lemma~1301, l.~1 & \textit{for} Lemma~892 \textit{read} Lemma~893 \\
\dag\,3681 & proof of Lemma~1302, (3.28) & \textit{for} $(\lambda x.\,{\phi}x \vee {\psi}x)(a) \leftrightarrow {\phi}a \wedge {\psi}a$ \textit{read} $(\lambda x.\,{\phi}x \wedge {\psi}x)(a) \leftrightarrow {\phi}a \wedge {\psi}a$; the formula is used in its correct form on p.~3682 \\
\dag\,3681 & proof of Lemma~1302, (3.27) & \textit{for} $\mathrm{Pos}( {\chi}) \leftrightarrow \neg\mathrm{Pos}({\chi})$ \textit{read} $\mathrm{Pos}(\neg {\chi}) \leftrightarrow \neg\mathrm{Pos}({\chi})$ \\
3682 & l.~34 & \textit{for} p.~3603 \textit{read} p.~3630 (Lemma~806) \\
3682 & l.~4 from foot & \textit{for} p.~3546 \textit{read} p.~3556 (Lemma~140) \\
3682 & n.~13 & \textit{for} Gallin (1976) \textit{read} Gallin (1975) \\
\end{errlist}

\subsection*{Chapter~4. Formal material for Volume II \hfill \normalfont\small pp.~3683--3728}
\begin{errlist}
\dag\,3684 & proof of Lemma~1311, l.~8 & \textit{for} Lemma~1306 \textit{read} Lemma~1310 \\
\dag\,3685 & proof of Lemma~1314, l.~7 & \textit{for} Lemma~1307 \textit{read} Lemma~1304; the step uses Lemma~1304 (p.~3684) \\
\dag\,3685 & proof of Lemma~1321, l.~8 & \textit{for} Lemma~1305 \textit{read} Lemma~1304 \\
3685 & l.~13 & \textit{for} p.~3598 \textit{read} p.~3588 (Lemma~441) \\
\dag\,3685 & proof of Lemma~1311, l.~19 & \textit{for} Lemma~1313 \textit{read} Lemma~1308; the step uses Lemma~1308 (p.~3684) \\
\dag\,3685 & proof of Lemma~1313, l.~23 & \textit{for} Lemma~1308 \textit{read} Lemma~1307 \\
\dag\,3686 & proof of Lemma~1323, l.~12 & \textit{for} Lemma~1305 \textit{read} Lemma~1303 \\
\dag\,3686 & proof of Lemma~1323, l.~21 & \textit{for} Lemma~1310 \textit{read} Lemma~1305 \\
\dag\,3686 & proof of Lemma~1327, l.~27 & \textit{for} Lemma~1330 \textit{read} Lemma~1325 \\
3687 & l.~1 & \textit{for} p.~3767 \textit{read} p.~3677 (Lemma~1258) \\
\dag\,3687 & proof of Lemma~1347, l.~9 & \textit{for} Lemma~1329 \textit{read} Lemma~1327 \\
\dag\,3687 & proof of Lemma~1343, l.~12 & \textit{for} Lemma~1331 \textit{read} Lemma~1328; the step uses Lemma~1328 (p.~3686) \\
\dag\,3687 & proof of Lemma~1346, l.~14 & \textit{for} Lemma~1334 \textit{read} Lemma~1337 \\
\dag\,3687 & proof of Lemma~1344, l.~27 & \textit{for} Lemma~1330 \textit{read} Lemma~1333; the step uses Lemma~1333 (p.~3686) \\
\dag\,3687 & proof of Lemma~1347, l.~30 & \textit{for} Lemma~1328 \textit{read} Lemma~1324 \\
\dag\,3688 & proof of Lemma~1338, l.~11 & \textit{for} Lemma~1339 \textit{read} Lemma~1335 \\
\dag\,3689 & proof of Lemma~1354, l.~16 & \textit{for} Lemma~1320 \textit{read} Lemma~1318; the step uses Lemma~1318 (p.~3685) \\
3690 & l.~7 & \textit{for} p.~3611 \textit{read} p.~3601 (Lemma~560) \\
\dag\,3691 & proof of Lemma~1386, l.~13 & \textit{for} Lemma~1370 \textit{read} Lemma~1373; the step uses Lemma~1373 (p.~3690) \\
\dag\,3691 & proof of Lemma~1379, l.~19 & \textit{for} Lemma~1319 \textit{read} Lemma~1317; the step uses Lemma~1317 (p.~3685) \\
3692 & l.~29 & \textit{for} p.~3656 \textit{read} p.~3566 (Lemma~243) \\
3692 & l.~32 & \textit{for} p.~3595 \textit{read} p.~3594 (Lemma~493) \\
3693 & l.~11 & \textit{after} dependee \textit{insert} in every accessible world \\
\dag\,3693 & proof of Lemma~1409, l.~23 & \textit{for} Lemma~1332 \textit{read} Lemma~1330; the step uses Lemma~1330 (p.~3686) \\
3693 & l.~32 & \textit{for} $\mathrm{NE}(y) \vee \neg\exists z\,D(y,z) \to \mathrm{ME}(y)$ \textit{read} $\mathrm{NE}(y) \wedge \neg\exists z\,D(y,z) \to \mathrm{ME}(y)$ \\
\dag\,3694 & proof of Lemma~1423, l.~19 & \textit{for} Lemma~1316 \textit{read} Lemma~1319; the step uses Lemma~1319 (p.~3685) \\
\dag\,3694 & proof of Lemma~1418, l.~30 & \textit{for} Lemma~1318 \textit{read} Lemma~1317; the step uses Lemma~1317 (p.~3685) \\
\dag\,3695 & proof of Lemma~1437, l.~15 & \textit{for} Lemma~1362 \textit{read} Lemma~1357; the step uses Lemma~1357 (p.~3688) \\
\dag\,3696 & proof of Lemma~1442, l.~22 & \textit{for} Lemma~1435 \textit{read} Lemma~1430 \\
\dag\,3696 & proof of Lemma~1436, l.~23 & \textit{for} Lemma~1322 \textit{read} Lemma~1323 \\
\dag\,3696 & proof of Lemma~1439, l.~27 & \textit{for} Lemma~1420 \textit{read} Lemma~1424 \\
3696 & l.~32 & \textit{for} p.~3675 \textit{read} p.~3575 (Lemma~332) \\
\dag\,3697 & proof of Lemma~1461, (4.4) & \textit{for} $D(y,x) \wedge D(x,z) \to D(y,y)$ \textit{read} $D(y,x) \wedge D(x,z) \to D(y,z)$ \\
\dag\,3698 & proof of Lemma~1470, l.~1 & \textit{for} Lemma~1350 \textit{read} Lemma~1349 \\
\dag\,3698 & proof of Lemma~1461, l.~15 & \textit{for} Lemma~1383 \textit{read} Lemma~1384; the step uses Lemma~1384 (p.~3691) \\
\dag\,3699 & proof of Lemma~1484, (4.5) & \textit{for} $\forall y\,(\neg\mathrm{NE}(y) \leftrightarrow \exists x\,D(y,x))$ \textit{read} $\forall y\,(\neg\mathrm{NE}(y) \to \exists x\,D(y,x))$ \\
3699 & l.~4 & \textit{for} p.~3712 \textit{read} p.~3612 (Lemma~661) \\
3699 & l.~27 & \textit{for} p.~3573 \textit{read} p.~3572 (Lemma~298) \\
\dag\,3700 & proof of Lemma~1479, l.~22 & \textit{for} Lemma~1328 \textit{read} Lemma~1330; the step uses Lemma~1330 (p.~3686) \\
\dag\,3700 & proof of Lemma~1489, l.~30 & \textit{for} Lemma~1306 \textit{read} Lemma~1307 \\
\dag\,3701 & proof of Lemma~1496, l.~3 & \textit{for} Lemma~1365 \textit{read} Lemma~1362 \\
3701 & l.~27 & \textit{for} p.~3629 \textit{read} p.~3619 (Lemma~726) \\
3701 & l.~29 & \textit{for} p.~3667 \textit{read} p.~3677 (Lemma~1259) \\
\dag\,3702 & proof of Lemma~1509, l.~7 & \textit{for} Lemma~1392 \textit{read} Lemma~1393 \\
\dag\,3702 & proof of Lemma~1499, l.~21 & \textit{for} Lemma~1361 \textit{read} Lemma~1356 \\
\dag\,3702 & proof of Lemma~1515, l.~26 & \textit{for} Lemma~1401 \textit{read} Lemma~1399 \\
\dag\,3704 & proof of Lemma~1540, l.~4 & \textit{for} Lemma~1379 \textit{read} Lemma~1374 \\
\dag\,3704 & proof of Lemma~1532, l.~9 & \textit{for} Lemma~1348 \textit{read} Lemma~1346; the step uses Lemma~1346 (p.~3687) \\
\dag\,3704 & proof of Lemma~1523, l.~15 & \textit{for} Lemma~1304 \textit{read} Lemma~1303 \\
\dag\,3704 & proof of Lemma~1523, l.~20 & \textit{for} Lemma~1469 \textit{read} Lemma~1471; the step uses Lemma~1471 (p.~3698) \\
3704 & l.~8 from foot & \textit{for} Rowe (1975) \textit{read} Koons (1997) (the cosmological argument from contingency) \\
\dag\,3705 & proof of Lemma~1535, l.~3 & \textit{for} Lemma~1314 \textit{read} Lemma~1316 \\
\dag\,3705 & proof of Lemma~1552, (4.8) & \textit{for} $D(z,x) \wedge D(x,y) \leftrightarrow D(z,y)$ \textit{read} $D(z,x) \wedge D(x,y) \to D(z,y)$ \\
\dag\,3706 & proof of Lemma~1556, l.~3 & \textit{for} Lemma~1327 \textit{read} Lemma~1328 \\
\dag\,3706 & proof of Lemma~1563, (4.9) & \textit{for} $\forall y\,(\mathrm{NE}(y) \to \exists x\,D(y,x))$ \textit{read} $\forall y\,(\neg\mathrm{NE}(y) \to \exists x\,D(y,x))$ \\
\dag\,3706 & proof of Lemma~1556, l.~17 & \textit{for} Lemma~1495 \textit{read} Lemma~1491 \\
\dag\,3706 & proof of Lemma~1558, l.~19 & \textit{for} Lemma~1358 \textit{read} Lemma~1356; the step uses Lemma~1356 (p.~3688) \\
\dag\,3706 & proof of Lemma~1556, l.~25 & \textit{for} Lemma~1340 \textit{read} Lemma~1337 \\
\dag\,3706 & proof of Lemma~1554, l.~30 & \textit{for} Lemma~1362 \textit{read} Lemma~1367 \\
\dag\,3707 & proof of Lemma~1567, l.~7 & \textit{for} Lemma~1551 \textit{read} Lemma~1554; the step uses Lemma~1554 (p.~3706) \\
\dag\,3707 & proof of Lemma~1569, l.~9 & \textit{for} Lemma~1379 \textit{read} Lemma~1383; the step uses Lemma~1383 (p.~3691) \\
\dag\,3707 & proof of Lemma~1568, l.~27 & \textit{for} Lemma~1367 \textit{read} Lemma~1371 \\
\dag\,3708 & proof of Lemma~1581, l.~9 & \textit{for} Lemma~1460 \textit{read} Lemma~1457 \\
\dag\,3708 & proof of Lemma~1586, l.~31 & \textit{for} $\mathrm{NE}(x) \wedge \neg\forall z\,D(x,z) \to \mathrm{ME}(x)$ \textit{read} $\mathrm{NE}(x) \wedge \neg\exists z\,D(x,z) \to \mathrm{ME}(x)$ \\
\dag\,3709 & proof of Lemma~1589, l.~11 & \textit{for} Lemma~1420 \textit{read} Lemma~1418; the step uses Lemma~1418 (p.~3694) \\
\dag\,3709 & proof of Lemma~1585, l.~16 & \textit{for} Lemma~1406 \textit{read} Lemma~1404 \\
\dag\,3709 & proof of Lemma~1584, l.~19 & \textit{for} Lemma~1373 \textit{read} Lemma~1374; the step uses Lemma~1374 (p.~3690) \\
3709 & l.~23 & \textit{for} p.~3561 \textit{read} p.~3551 (Lemma~93) \\
3710 & l.~39 & \textit{for} Pruss (2006) \textit{read} Koons (1997) (the cosmological argument from contingency) \\
3713 & l.~2 & \textit{for} p.~3552 \textit{read} p.~3553 (Lemma~117) \\
\dag\,3713 & proof of Lemma~1625, l.~6 & \textit{for} Lemma~1415 \textit{read} Lemma~1414 \\
\dag\,3713 & proof of Lemma~1629, l.~23 & \textit{for} Lemma~1621 \textit{read} Lemma~1626 \\
\dag\,3713 & proof of Lemma~1638, l.~28 & \textit{for} Lemma~1440 \textit{read} Lemma~1444; the step uses Lemma~1444 (p.~3696) \\
3715 & l.~15 & \textit{for} p.~3662 \textit{read} p.~3663 (Lemma~1125) \\
3715 & l.~18 & \textit{for} p.~3701 \textit{read} p.~3702 (Lemma~1512) \\
\dag\,3715 & proof of Lemma~1655, l.~24 & \textit{for} Lemma~1636 \textit{read} Lemma~1640 \\
\dag\,3716 & proof of Lemma~1664, l.~8 & \textit{for} Lemma~1553 \textit{read} Lemma~1549 \\
\dag\,3716 & proof of Lemma~1668, l.~18 & \textit{for} Lemma~1558 \textit{read} Lemma~1553; the step uses Lemma~1553 (p.~3706) \\
\dag\,3716 & proof of Lemma~1676, l.~23 & \textit{for} Lemma~1619 \textit{read} Lemma~1614; the step uses Lemma~1614 (p.~3711) \\
3716 & l.~28 & \textit{for} p.~3567 \textit{read} p.~3577 (Lemma~353) \\
\dag\,3717 & proof of Lemma~1667, l.~14 & \textit{for} Lemma~1355 \textit{read} Lemma~1354 \\
3717 & l.~28 & \textit{for} p.~3661 \textit{read} p.~3662 (Lemma~1119) \\
\dag\,3718 & proof of Lemma~1683, l.~10 & \textit{for} Lemma~1331 \textit{read} Lemma~1329 \\
\dag\,3719 & proof of Lemma~1704, l.~19 & \textit{for} Lemma~1652 \textit{read} Lemma~1651 \\
\dag\,3720 & proof of Lemma~1716, l.~20 & \textit{for} Lemma~1701 \textit{read} Lemma~1705; the step uses Lemma~1705 (p.~3719) \\
3720 & l.~5 from foot & \textit{for} p.~3580 \textit{read} p.~3570 (Lemma~282) \\
3720 & l.~1 from foot & \textit{for} p.~3649 \textit{read} p.~3694 (Lemma~1426) \\
\dag\,3721 & proof of Lemma~1713, l.~3 & \textit{for} Lemma~1605 \textit{read} Lemma~1601 \\
3721 & l.~18 & \textit{for} p.~3550 \textit{read} p.~3551 (Lemma~91) \\
\dag\,3721 & proof of Lemma~1727, l.~27 & \textit{for} Lemma~1676 \textit{read} Lemma~1679; the step uses Lemma~1679 (p.~3717) \\
\dag\,3721 & proof of Lemma~1713, l.~30 & \textit{for} Lemma~1466 \textit{read} Lemma~1462; the step uses Lemma~1462 (p.~3698) \\
3721 & l.~2 from foot & \textit{for} p.~3574 \textit{read} p.~3584 (Lemma~404) \\
\dag\,3722 & proof of Lemma~1736, l.~5 & \textit{for} Lemma~1376 \textit{read} Lemma~1377 \\
\dag\,3723 & proof of Lemma~1741, l.~6 & \textit{for} Lemma~1680 \textit{read} Lemma~1681 \\
\dag\,3723 & proof of Lemma~1756, l.~13 & \textit{for} Lemma~1722 \textit{read} Lemma~1719; the step uses Lemma~1719 (p.~3720) \\
\dag\,3723 & proof of Lemma~1755, l.~23 & \textit{for} Lemma~1714 \textit{read} Lemma~1712 \\
3723 & l.~24 & \textit{for} p.~3715 \textit{read} p.~3716 (Lemma~1677) \\
\dag\,3723 & proof of Lemma~1751, l.~26 & \textit{for} Lemma~1533 \textit{read} Lemma~1531; the step uses Lemma~1531 (p.~3704) \\
3723 & l.~37 & \textit{for} p.~3548 \textit{read} p.~3549 (Lemma~73) \\
\dag\,3724 & proof of Lemma~1758, l.~1 & \textit{for} Lemma~1543 \textit{read} Lemma~1539 \\
\dag\,3724 & proof of Lemma~1764, l.~19 & \textit{for} Lemma~1539 \textit{read} Lemma~1540; the step uses Lemma~1540 (p.~3704) \\
\dag\,3724 & proof of Lemma~1757, l.~25 & \textit{for} Lemma~1729 \textit{read} Lemma~1731; the step uses Lemma~1731 (p.~3721) \\
\dag\,3725 & proof of Lemma~1772, l.~10 & \textit{for} Lemma~1609 \textit{read} Lemma~1612 \\
\dag\,3726 & proof of Lemma~1783, l.~6 & \textit{for} Lemma~1744 \textit{read} Lemma~1742 \\
\dag\,3726 & proof of Lemma~1786, l.~24 & \textit{for} Lemma~1325 \textit{read} Lemma~1323 \\
\dag\,3726 & proof of Lemma~1787, l.~29 & \textit{for} Lemma~1614 \textit{read} Lemma~1610 \\
3727 & l.~13 & \textit{delete} the second occurrence of \emph{a} \\
\dag\,3727 & proof of Lemma~1801, l.~18 & \textit{for} Lemma~1702 \textit{read} Lemma~1706 \\
\dag\,3727 & proof of Lemma~1797, l.~20 & \textit{for} Lemma~1573 \textit{read} Lemma~1576 \\
3727 & l.~29 & \textit{for} p.~3727 \textit{read} p.~3726 (Lemma~1785) \\
\dag\,3727 & proof of Lemma~1782, l.~30 & \textit{for} Lemma~1673 \textit{read} Lemma~1675 \\
3727 & l.~8 from foot & \textit{delete} the second occurrence of \emph{to} \\
3727 & l.~4 from foot & \textit{for} p.~3625 \textit{read} p.~3615 (Lemma~685) \\
\end{errlist}

\subsection*{Chapter~5. Formal material for Volume III \hfill \normalfont\small pp.~3729--3832}
\begin{errlist}
\dag\,3730 & proof of Lemma~1810, l.~2 & \textit{for} $\mathrm{MG}(y,u) \vee uRv' \to \mathrm{MG}(y,v')$ \textit{read} $\mathrm{MG}(y,u) \wedge uRv' \to \mathrm{MG}(y,v')$ \\
3730 & l.~24 & \textit{for} p.~3709 \textit{read} p.~3609 (Lemma~633) \\
\dag\,3730 & proof of Lemma~1810, l.~4 from foot & \textit{for} Corollary~7.15 of Volume~III \textit{read} Corollary~7.13 of Volume~III \\
3730 & l.~1 from foot & \textit{for} p.~3647 \textit{read} p.~3637 (Lemma~869) \\
\dag\,3731 & proof of Lemma~1814, l.~14 & \textit{for} Lemma~1816 \textit{read} Lemma~1813; the step uses Lemma~1813 (p.~3731) \\
3731 & l.~18 & \textit{for} p.~3622 \textit{read} p.~3623 (Lemma~758) \\
3732 & l.~15 & \textit{for} p.~3659 \textit{read} p.~3658 (Lemma~1076) \\
\dag\,3732 & proof of Lemma~1825, l.~18 & \textit{for} Lemma~1807 \textit{read} Lemma~1811; the step uses Lemma~1811 (p.~3731) \\
\dag\,3733 & proof of Lemma~1819, l.~10 & \textit{for} Lemma~1809 \textit{read} Lemma~1813 \\
\dag\,3733 & proof of Lemma~1828, l.~22 & \textit{for} Lemma~1810 \textit{read} Lemma~1809 \\
\dag\,3733 & proof of Lemma~1835, l.~24 & \textit{for} Proposition~7.20 of Volume~III \textit{read} Proposition~7.17 of Volume~III \\
\dag\,3733 & proof of Lemma~1827, l.~29 & \textit{for} Lemma~1802 \textit{read} Lemma~1803; the step uses Lemma~1803 (p.~3730) \\
3733 & l.~37 & \textit{for} $\mathrm{Pos}({\psi}) \wedge \mathrm{Pos}({\phi}) \leftrightarrow \mathrm{Pos}({\psi} \sqcap {\phi})$ \textit{read} $\mathrm{Pos}({\psi}) \wedge \mathrm{Pos}({\phi}) \to \mathrm{Pos}({\psi} \sqcap {\phi})$ \\
3734 & l.~2 & \textit{for} p.~3596 \textit{read} p.~3586 (Lemma~421) \\
\dag\,3734 & proof of Lemma~1843, (5.1) & \textit{for} $\mathrm{G}y \leftrightarrow \exists {\psi}\,(\mathrm{Pos}({\psi}) \to {\psi}y)$ \textit{read} $\mathrm{G}y \leftrightarrow \forall {\psi}\,(\mathrm{Pos}({\psi}) \to {\psi}y)$ \\
\dag\,3735 & proof of Lemma~1845, l.~1 & \textit{for} Lemma~1814 \textit{read} Lemma~1812 \\
\dag\,3735 & proof of Lemma~1846, l.~21 & \textit{for} Lemma~1803 \textit{read} Lemma~1806; the step uses Lemma~1806 (p.~3730) \\
3736 & l.~2 & \textit{for} p.~3691 \textit{read} p.~3692 (Lemma~1399) \\
3736 & l.~13 & \textit{for} p.~3661 \textit{read} p.~3662 (Lemma~1113) \\
\dag\,3736 & proof of Lemma~1855, l.~18 & \textit{for} Lemma~1847 \textit{read} Lemma~1846 \\
\dag\,3736 & proof of Lemma~1860, l.~27 & \textit{for} Lemma~1810 \textit{read} Lemma~1806; the step uses Lemma~1806 (p.~3730) \\
3737 & l.~1 & \textit{for} Sobel (1987) \textit{read} Anderson (1990) (the emended positivity axioms) \\
\dag\,3738 & proof of Lemma~1868, l.~13 & \textit{for} Lemma~1816 \textit{read} Lemma~1821; the step uses Lemma~1821 (p.~3732) \\
\dag\,3738 & proof of Lemma~1876, l.~15 & \textit{for} Lemma~7.7 of Volume~III \textit{read} Lemma~7.9 of Volume~III; the result cited is on p.~1105 \\
\dag\,3738 & proof of Lemma~1876, (5.4) & \textit{for} $\mathrm{NME}(y,w) \iff \forall v'\,(v'Rw \to \mathrm{ME}(y,v'))$ \textit{read} $\mathrm{NME}(y,w) \iff \forall v'\,(wRv' \to \mathrm{ME}(y,v'))$ \\
\dag\,3739 & proof of Lemma~1881, l.~2 & \textit{for} Lemma~1874 \textit{read} Lemma~1870 \\
\dag\,3739 & proof of Lemma~1882, l.~14 & \textit{for} Lemma~1815 \textit{read} Lemma~1813 \\
\dag\,3739 & proof of Lemma~1885, l.~39 & \textit{for} Lemma~7.4 of Volume~III \textit{read} Lemma~7.10 of Volume~III \\
\dag\,3740 & proof of Lemma~1886, l.~1 & \textit{for} Lemma~1869 \textit{read} Lemma~1871 \\
3740 & l.~2 & \textit{for} Lemma~5.19 of Volume~III \textit{read} Lemma~5.21 of Volume~III; the result cited is on p.~1010 \\
\dag\,3740 & proof of Lemma~1885, l.~20 & \textit{for} Lemma~1865 \textit{read} Lemma~1862 \\
\dag\,3741 & proof of Lemma~1889, l.~23 & \textit{for} Lemma~1869 \textit{read} Lemma~1865 \\
3741 & l.~40 & \textit{for} $\mathrm{MG}(z,w_0) \vee w_0Rv' \to \mathrm{MG}(z,v')$ \textit{read} $\mathrm{MG}(z,w_0) \wedge w_0Rv' \to \mathrm{MG}(z,v')$ \\
\dag\,3742 & proof of Lemma~1907, l.~20 & \textit{for} Lemma~1852 \textit{read} Lemma~1849; the step uses Lemma~1849 (p.~3735) \\
\dag\,3742 & proof of Lemma~1900, l.~25 & \textit{for} Lemma~1873 \textit{read} Lemma~1872; the step uses Lemma~1872 (p.~3738) \\
\dag\,3742 & proof of Lemma~1895, l.~27 & \textit{for} Lemma~1827 \textit{read} Lemma~1831 \\
\dag\,3742 & proof of Lemma~1910, (5.5) & \textit{for} $\mathrm{G}x \leftrightarrow \exists {\phi}\,(\mathrm{Pos}({\phi}) \leftrightarrow \nec {\phi}x)$ \textit{read} $\mathrm{G}x \leftrightarrow \forall {\phi}\,(\mathrm{Pos}({\phi}) \leftrightarrow \nec {\phi}x)$ \\
\dag\,3743 & proof of Lemma~1916, l.~1 & \textit{for} Lemma~1802 \textit{read} Lemma~1805 \\
3743 & l.~31 & \textit{for} p.~3711 \textit{read} p.~3701 (Lemma~1505) \\
\dag\,3744 & proof of Lemma~1918, l.~6 & \textit{for} Lemma~1892 \textit{read} Lemma~1893; the step uses Lemma~1893 (p.~3740) \\
3744 & l.~17 & \textit{for} positivtiy \textit{read} positivity \\
\dag\,3744 & proof of Lemma~1916, l.~25 & \textit{for} Lemma~1890 \textit{read} Lemma~1894 \\
\dag\,3745 & proof of Lemma~1928, l.~12 & \textit{for} Lemma~1884 \textit{read} Lemma~1880 \\
\dag\,3745 & proof of Lemma~1935, l.~5 from foot & \textit{for} Definition~7.11 of Volume~III \textit{read} Definition~7.14 of Volume~III; the result cited is on p.~1114 \\
\dag\,3746 & proof of Lemma~1936, l.~14 & \textit{for} Lemma~1804 \textit{read} Lemma~1803 \\
3746 & l.~20 & \textit{delete} the second occurrence of \emph{the} \\
\dag\,3747 & proof of Lemma~1938, l.~7 & \textit{for} Lemma~1819 \textit{read} Lemma~1821; the step uses Lemma~1821 (p.~3732) \\
\dag\,3747 & proof of Lemma~1951, l.~15 & \textit{for} Proposition~7.17 of Volume~III \textit{read} Proposition~7.20 of Volume~III \\
3747 & l.~30 & \textit{for} p.~3642 \textit{read} p.~3632 (Lemma~825) \\
\dag\,3748 & proof of Lemma~1960, l.~6 & \textit{for} Definition~5.18 of Volume~III \textit{read} Definition~5.20 of Volume~III; the result cited is on p.~1010 \\
\dag\,3748 & proof of Lemma~1959, l.~13 & \textit{for} Lemma~1835 \textit{read} Lemma~1834 \\
\dag\,3748 & proof of Lemma~1949, l.~17 & \textit{for} Lemma~1894 \textit{read} Lemma~1897; the step uses Lemma~1897 (p.~3741) \\
\dag\,3749 & proof of Lemma~1968, (5.9) & \textit{for} $\mathrm{G}z \to \forall {\psi}\,(\mathrm{Pos}({\psi}) \to {\psi}z)$ \textit{read} $\mathrm{G}z \leftrightarrow \forall {\psi}\,(\mathrm{Pos}({\psi}) \to {\psi}z)$ \\
\dag\,3749 & proof of Lemma~1968, (5.6) & \textit{for} $\mathrm{NE}(y,v') \iff \forall w_0\,(w_0Rv' \to E(y,w_0))$ \textit{read} $\mathrm{NE}(y,v') \iff \forall w_0\,(v'Rw_0 \to E(y,w_0))$; correctly printed in (5.7), p.~3749 \\
\dag\,3750 & proof of Lemma~1976, (5.10) & \textit{for} $\pos G \to \neg\pos G$ \textit{read} $\pos G \to \neg\pos\neg G$; correctly printed in (5.14), p.~3755 \\
\dag\,3750 & proof of Lemma~1976, l.~19 & \textit{for} Corollary~7.13 of Volume~III \textit{read} Corollary~7.15 of Volume~III \\
\dag\,3750 & proof of Lemma~1974, l.~23 & \textit{for} Lemma~1824 \textit{read} Lemma~1825 \\
3750 & l.~39 & \textit{for} p.~3606 \textit{read} p.~3616 (Lemma~691) \\
3751 & l.~5 & \textit{for} p.~3720 \textit{read} p.~3721 (Lemma~1729) \\
\dag\,3751 & proof of Lemma~1982, l.~8 & \textit{for} Lemma~1977 \textit{read} Lemma~1979; the step uses Lemma~1979 (p.~3751) \\
\dag\,3751 & proof of Lemma~1985, l.~12 & \textit{for} Theorem~5.4 of Volume~III \textit{read} Theorem~5.6 of Volume~III; the result cited is on p.~983 \\
\dag\,3751 & proof of Lemma~1981, l.~29 & \textit{for} Lemma~1892 \textit{read} Lemma~1891 \\
\dag\,3752 & proof of Lemma~1990, l.~10 & \textit{for} Lemma~1937 \textit{read} Lemma~1935 \\
\dag\,3752 & proof of Lemma~1984, l.~16 & \textit{for} Lemma~1878 \textit{read} Lemma~1879 \\
3752 & l.~17 & \textit{for} Definition~5.14 of Volume~III \textit{read} Definition~5.20 of Volume~III; the result cited is on p.~1010 \\
3752 & l.~20 & \textit{for} p.~3542 \textit{read} p.~3552 (Lemma~106) \\
\dag\,3752 & proof of Lemma~1990, l.~21 & \textit{for} Lemma~1804 \textit{read} Lemma~1802; the step uses Lemma~1802 (p.~3730) \\
3752 & l.~22 & \textit{for} p.~3732 \textit{read} p.~3722 (Lemma~1734) \\
\dag\,3753 & proof of Lemma~2001, (5.12) & \textit{for} $\mathrm{MG}(y,v') \iff \forall u\,(v'Ru \to E(y,v') \wedge \mathrm{ME}(y,u))$ \textit{read} $\mathrm{MG}(y,v') \iff \forall u\,(v'Ru \to E(y,u) \wedge \mathrm{ME}(y,u))$ \\
\dag\,3753 & proof of Lemma~2001, l.~24 & \textit{for} Lemma~7.10 of Volume~III \textit{read} Lemma~7.7 of Volume~III \\
3754 & l.~4 & \textit{for} p.~3565 \textit{read} p.~3665 (Lemma~1140) \\
\dag\,3754 & proof of Lemma~2010, l.~22 & \textit{for} Corollary~7.5 of Volume~III \textit{read} Corollary~7.3 of Volume~III \\
\dag\,3755 & proof of Lemma~2018, (5.13) & \textit{for} $\mathrm{NE}(x,v') \iff \exists w\,(v'Rw \to E(x,w))$ \textit{read} $\mathrm{NE}(x,v') \iff \forall w\,(v'Rw \to E(x,w))$ \\
\dag\,3756 & proof of Lemma~2012, l.~23 & \textit{for} Lemma~1807 \textit{read} Lemma~1803 \\
\dag\,3757 & proof of Lemma~2035, l.~28 & \textit{for} Lemma~1902 \textit{read} Lemma~1899; the step uses Lemma~1899 (p.~3741) \\
3758 & l.~33 & \textit{for} p.~3613 \textit{read} p.~3713 (Lemma~1642) \\
\dag\,3759 & proof of Lemma~2043, l.~22 & \textit{for} Lemma~1889 \textit{read} Lemma~1890 \\
\dag\,3760 & proof of Lemma~2060, (5.17) & \textit{for} ${\phi} \in \mathcal{P} \wedge \pos\forall y\,({\phi}y \to {\chi}y) \Rightarrow {\chi} \in \mathcal{P}$ \textit{read} ${\phi} \in \mathcal{P} \wedge \nec\forall y\,({\phi}y \to {\chi}y) \Rightarrow {\chi} \in \mathcal{P}$; correctly printed in (5.18), p.~3763 \\
\dag\,3760 & proof of Lemma~2058, l.~24 & \textit{for} Lemma~1847 \textit{read} Lemma~1844 \\
\dag\,3761 & proof of Lemma~2064, l.~13 & \textit{for} Lemma~2042 \textit{read} Lemma~2039 \\
3761 & l.~19 & \textit{for} p.~3579 \textit{read} p.~3578 (Lemma~364) \\
3761 & l.~29 & \textit{for} execllence \textit{read} excellence \\
\dag\,3761 & proof of Lemma~2068, l.~34 & \textit{for} Proposition~5.7 of Volume~III \textit{read} Proposition~5.8 of Volume~III; the result cited is on p.~991 \\
\dag\,3762 & proof of Lemma~2067, l.~11 & \textit{for} Lemma~2011 \textit{read} Lemma~2008 \\
\dag\,3762 & proof of Lemma~2076, l.~12 & \textit{for} $\mathrm{MG}(y,u) \leftrightarrow \mathrm{NE}(y,u) \wedge \mathrm{NME}(y,u)$ \textit{read} $\mathrm{MG}(y,u) \to \mathrm{NE}(y,u) \wedge \mathrm{NME}(y,u)$ \\
\dag\,3762 & proof of Lemma~2074, l.~15 & \textit{for} Lemma~1943 \textit{read} Lemma~1946 \\
3762 & l.~23 & \textit{for} p.~3606 \textit{read} p.~3607 (Lemma~616) \\
\dag\,3762 & proof of Lemma~2076, l.~28 & \textit{for} $\mathrm{Pos}({\phi}) \to \pos\forall z\,{\phi}z$ \textit{read} $\mathrm{Pos}({\phi}) \to \pos\exists z\,{\phi}z$ \\
3763 & l.~8 & \textit{for} Definition~5.10 of Volume~III \textit{read} Definition~5.14 of Volume~III \\
\dag\,3763 & proof of Lemma~2082, l.~18 & \textit{for} Lemma~1891 \textit{read} Lemma~1889 \\
3763 & l.~21 & \textit{for} p.~3622 \textit{read} p.~3722 (Lemma~1738) \\
\dag\,3763 & proof of Lemma~2082, l.~23 & \textit{for} Lemma~1897 \textit{read} Lemma~1893 \\
3763 & l.~27 & \textit{for} p.~3633 \textit{read} p.~3643 (Lemma~931) \\
\dag\,3764 & proof of Lemma~2087, l.~9 & \textit{for} Lemma~1906 \textit{read} Lemma~1911 \\
\dag\,3764 & proof of Lemma~2093, l.~18 & \textit{for} $\mathrm{G}x \to \forall {\chi}\,(\mathrm{Pos}({\chi}) \to {\chi}x)$ \textit{read} $\mathrm{G}x \leftrightarrow \forall {\chi}\,(\mathrm{Pos}({\chi}) \to {\chi}x)$ \\
\dag\,3765 & proof of Lemma~2098, l.~4 & \textit{for} Lemma~1811 \textit{read} Lemma~1814; the step uses Lemma~1814 (p.~3731) \\
\dag\,3766 & proof of Lemma~2102, l.~2 & \textit{for} Lemma~1973 \textit{read} Lemma~1969; the step uses Lemma~1969 (p.~3750) \\
\dag\,3766 & proof of Lemma~2104, l.~15 & \textit{for} Lemma~1858 \textit{read} Lemma~1854 \\
\dag\,3766 & proof of Lemma~2102, l.~17 & \textit{for} Lemma~2037 \textit{read} Lemma~2042 \\
3766 & l.~32 & \textit{for} p.~3721 \textit{read} p.~3722 (Lemma~1739) \\
\dag\,3767 & proof of Lemma~2105, l.~3 & \textit{for} Lemma~2021 \textit{read} Lemma~2023 \\
3768 & l.~23 & \textit{for} p.~3621 \textit{read} p.~3721 (Lemma~1730) \\
3768 & l.~31 & \textit{for} p.~3852 \textit{read} p.~3752 (Lemma~1986) \\
\dag\,3768 & proof of Lemma~2126, l.~35 & \textit{for} Lemma~5.19 of Volume~III \textit{read} Lemma~5.17 of Volume~III \\
\dag\,3769 & proof of Lemma~2134, l.~7 & \textit{for} Definition~5.13 of Volume~III \textit{read} Definition~5.10 of Volume~III \\
\dag\,3770 & proof of Lemma~2140, l.~3 & \textit{for} Lemma~1817 \textit{read} Lemma~1813 \\
\dag\,3770 & proof of Lemma~2132, l.~18 & \textit{for} Lemma~1951 \textit{read} Lemma~1946 \\
3770 & l.~27 & \textit{for} p.~3738 \textit{read} p.~3748 (Lemma~1960) \\
\dag\,3770 & proof of Lemma~2138, l.~29 & \textit{for} Lemma~1981 \textit{read} Lemma~1977 \\
\dag\,3771 & proof of Lemma~2147, l.~24 & \textit{for} Lemma~1833 \textit{read} Lemma~1832; the step uses Lemma~1832 (p.~3733) \\
3771 & l.~30 & \textit{for} Definition~5.14 of Volume~III \textit{read} Definition~5.10 of Volume~III; the result cited is on p.~991 \\
3772 & l.~4 & \textit{for} conjnuction \textit{read} conjunction \\
3773 & l.~7 & \textit{for} p.~3714 \textit{read} p.~3614 (Lemma~677) \\
\dag\,3774 & proof of Lemma~2176, l.~3 from foot & \textit{for} Proposition~5.7 of Volume~III \textit{read} Proposition~5.3 of Volume~III; the result cited is on p.~965 \\
3775 & l.~1 & \textit{for} p.~3716 \textit{read} p.~3616 (Lemma~694) \\
\dag\,3775 & proof of Lemma~2169, l.~18 & \textit{for} Lemma~2012 \textit{read} Lemma~2017 \\
\dag\,3775 & proof of Lemma~2181, l.~22 & \textit{for} Lemma~1802 \textit{read} Lemma~1806; the step uses Lemma~1806 (p.~3730) \\
\dag\,3775 & proof of Lemma~2172, l.~23 & \textit{for} Lemma~1944 \textit{read} Lemma~1946 \\
\dag\,3776 & proof of Lemma~2183, l.~6 & \textit{for} Lemma~1864 \textit{read} Lemma~1869 \\
\dag\,3776 & proof of Lemma~2186, l.~7 & \textit{for} Lemma~2107 \textit{read} Lemma~2105 \\
\dag\,3777 & proof of Lemma~2191, l.~21 & \textit{for} Lemma~2123 \textit{read} Lemma~2127 \\
\dag\,3779 & proof of Lemma~2205, l.~2 & \textit{for} Lemma~1881 \textit{read} Lemma~1878 \\
3779 & l.~5 & \textit{for} p.~3650 \textit{read} p.~3660 (Lemma~1095) \\
\dag\,3780 & proof of Lemma~2211, l.~24 & \textit{for} Lemma~1896 \textit{read} Lemma~1891 \\
\dag\,3780 & proof of Lemma~2220, l.~28 & \textit{for} Lemma~2030 \textit{read} Lemma~2035; the step uses Lemma~2035 (p.~3757) \\
\dag\,3781 & proof of Lemma~2234, (5.28) & \textit{for} $\mathrm{MG}(y,w_0) \wedge w_0Rv' \leftrightarrow \mathrm{MG}(y,v')$ \textit{read} $\mathrm{MG}(y,w_0) \wedge w_0Rv' \to \mathrm{MG}(y,v')$; the formula is used in its correct form on p.~3782 \\
\dag\,3781 & proof of Lemma~2225, l.~30 & \textit{for} Lemma~2134 \textit{read} Lemma~2130 \\
\dag\,3782 & proof of Lemma~2238, l.~16 & \textit{for} Lemma~2000 \textit{read} Lemma~1995; the step uses Lemma~1995 (p.~3753) \\
\dag\,3782 & proof of Lemma~2240, l.~29 & \textit{for} Lemma~2186 \textit{read} Lemma~2183; the step uses Lemma~2183 (p.~3775) \\
\dag\,3784 & proof of Lemma~2257, l.~30 & \textit{for} Lemma~2175 \textit{read} Lemma~2178 \\
\dag\,3784 & proof of Lemma~2259, (5.29) & \textit{for} $\mathrm{NME}(z,w') \iff \forall w_2\,(w'Rw_1 \to \mathrm{ME}(z,w_1))$ \textit{read} $\mathrm{NME}(z,w') \iff \forall w_1\,(w'Rw_1 \to \mathrm{ME}(z,w_1))$ \\
3784 & l.~1 from foot & \textit{for} compssible \textit{read} compossible \\
3785 & l.~15 & \textit{for} p.~3769 \textit{read} p.~3679 (Lemma~1282) \\
\dag\,3785 & proof of Lemma~2261, l.~19 & \textit{for} Lemma~2112 \textit{read} Lemma~2114 \\
\dag\,3786 & proof of Lemma~2273, l.~3 & \textit{for} Lemma~1896 \textit{read} Lemma~1892 \\
\dag\,3786 & proof of Lemma~2276, l.~7 & \textit{for} Lemma~7.7 of Volume~III \textit{read} Lemma~7.2 of Volume~III \\
\dag\,3786 & proof of Lemma~2268, l.~22 & \textit{for} Lemma~2023 \textit{read} Lemma~2028; the step uses Lemma~2028 (p.~3757) \\
3786 & l.~6 from foot & \textit{for} Definition~7.8 of Volume~III \textit{read} Definition~7.11 of Volume~III \\
\dag\,3787 & proof of Lemma~2284, (5.30) & \textit{for} $\mathcal{P}_0 \supseteq \mathcal{P} \Rightarrow \exists x\,\exists w_0\,\forall {\psi} \in \mathcal{P}_0\; {\psi}(x,w_0)$ \textit{read} $\mathcal{P}_0 \subseteq \mathcal{P} \Rightarrow \exists x\,\exists w_0\,\forall {\psi} \in \mathcal{P}_0\; {\psi}(x,w_0)$ \\
3787 & l.~12 & \textit{for} p.~3742 \textit{read} p.~3741 (Lemma~1894) \\
\dag\,3787 & proof of Lemma~2278, l.~16 & \textit{for} Lemma~1859 \textit{read} Lemma~1863 \\
\dag\,3788 & proof of Lemma~2277, l.~6 & \textit{for} Lemma~2102 \textit{read} Lemma~2103 \\
\dag\,3788 & proof of Lemma~2286, l.~18 & \textit{for} Lemma~1979 \textit{read} Lemma~1984; the step uses Lemma~1984 (p.~3751) \\
\dag\,3788 & proof of Lemma~2283, l.~27 & \textit{for} Lemma~1945 \textit{read} Lemma~1944 \\
\dag\,3789 & proof of Lemma~2301, l.~2 & \textit{for} $\mathrm{Pos}({\psi}) \to \pos\mathrm{Pos}({\psi})$ \textit{read} $\mathrm{Pos}({\psi}) \to \nec\mathrm{Pos}({\psi})$ \\
\dag\,3789 & proof of Lemma~2301, (5.32) & \textit{for} $\mathrm{MG}(z,v) \iff \forall w'\,(vRw' \to E(z,w') \vee \mathrm{ME}(z,w'))$ \textit{read} $\mathrm{MG}(z,v) \iff \forall w'\,(vRw' \to E(z,w') \wedge \mathrm{ME}(z,w'))$ \\
3789 & l.~10 & \textit{for} Gödel (1995) \textit{read} Anderson (1990) (the emended positivity axioms) \\
\dag\,3789 & proof of Lemma~2289, l.~15 & \textit{for} Lemma~1942 \textit{read} Lemma~1941 \\
\dag\,3790 & proof of Lemma~2309, l.~17 & \textit{for} Proposition~5.11 of Volume~III \textit{read} Proposition~5.7 of Volume~III \\
\dag\,3791 & proof of Lemma~2318, l.~17 & \textit{for} Proposition~5.8 of Volume~III \textit{read} Proposition~5.2 of Volume~III \\
\dag\,3792 & proof of Lemma~2326, l.~3 & \textit{for} Lemma~2041 \textit{read} Lemma~2038; the step uses Lemma~2038 (p.~3758) \\
\dag\,3792 & proof of Lemma~2320, l.~8 & \textit{for} Lemma~1893 \textit{read} Lemma~1888 \\
3792 & l.~11 & \textit{for} p.~3740 \textit{read} p.~3730 (Lemma~1805) \\
\dag\,3792 & proof of Lemma~2310, l.~13 & \textit{for} Lemma~2108 \textit{read} Lemma~2106 \\
\dag\,3792 & proof of Lemma~2324, l.~15 & \textit{for} Lemma~2306 \textit{read} Lemma~2309 \\
\dag\,3793 & proof of Lemma~2319, l.~4 & \textit{for} Lemma~2216 \textit{read} Lemma~2220; the step uses Lemma~2220 (p.~3780) \\
\dag\,3793 & proof of Lemma~2325, l.~10 & \textit{for} Lemma~1968 \textit{read} Lemma~1969 \\
\dag\,3794 & proof of Lemma~2343, l.~13 & \textit{for} Lemma~5.21 of Volume~III \textit{read} Lemma~5.17 of Volume~III; the result cited is on p.~1003 \\
\dag\,3794 & proof of Lemma~2328, l.~27 & \textit{for} Lemma~1850 \textit{read} Lemma~1845; the step uses Lemma~1845 (p.~3735) \\
\dag\,3794 & proof of Lemma~2343, (5.35) & \textit{for} $\exists z\,\mathrm{MG}(z,w_1) \leftrightarrow \exists z\,(E(z,w_1) \wedge \mathrm{ME}(z,w_1))$ \textit{read} $\exists z\,\mathrm{MG}(z,w_1) \to \exists z\,(E(z,w_1) \wedge \mathrm{ME}(z,w_1))$; correctly printed in (5.36), p.~3797 \\
3795 & l.~1 & \textit{for} Definition~7.14 of Volume~III \textit{read} Definition~7.8 of Volume~III; the result cited is on p.~1104 \\
3795 & l.~4 & \textit{after} proper filter \textit{insert} of properties \\
3795 & l.~20 & \textit{for} p.~3739 \textit{read} p.~3639 (Lemma~895) \\
\dag\,3795 & proof of Lemma~2335, l.~24 & \textit{for} Lemma~2309 \textit{read} Lemma~2304 \\
\dag\,3796 & proof of Lemma~2353, l.~2 & \textit{for} Lemma~2031 \textit{read} Lemma~2035 \\
\dag\,3796 & proof of Lemma~2358, l.~5 & \textit{for} Lemma~2356 \textit{read} Lemma~2354 \\
\dag\,3796 & proof of Lemma~2350, l.~7 & \textit{for} Lemma~1990 \textit{read} Lemma~1991; the step uses Lemma~1991 (p.~3752) \\
\dag\,3796 & proof of Lemma~2351, l.~25 & \textit{for} Lemma~1961 \textit{read} Lemma~1958 \\
\dag\,3797 & proof of Lemma~2360, l.~1 & \textit{for} Lemma~1904 \textit{read} Lemma~1902 \\
\dag\,3797 & proof of Lemma~2365, l.~3 & \textit{for} Lemma~2059 \textit{read} Lemma~2060; the step uses Lemma~2060 (p.~3760) \\
3797 & l.~12 & \textit{for} p.~3615 \textit{read} p.~3715 (Lemma~1663) \\
3797 & l.~15 & \textit{for} p.~3753 \textit{read} p.~3653 (Lemma~1028) \\
\dag\,3798 & proof of Lemma~2376, l.~3 & \textit{for} Lemma~1974 \textit{read} Lemma~1971; the step uses Lemma~1971 (p.~3750) \\
\dag\,3799 & proof of Lemma~2368, l.~12 & \textit{for} Lemma~2167 \textit{read} Lemma~2165 \\
3799 & l.~15 & \textit{for} Definition~5.10 of Volume~III \textit{read} Definition~5.5 of Volume~III; the result cited is on p.~981 \\
\dag\,3799 & proof of Lemma~2382, l.~29 & \textit{for} Lemma~1807 \textit{read} Lemma~1806 \\
\dag\,3801 & proof of Lemma~2400, l.~12 & \textit{for} Lemma~2335 \textit{read} Lemma~2340 \\
\dag\,3801 & proof of Lemma~2401, l.~15 & \textit{for} Definition~5.15 of Volume~III \textit{read} Definition~5.18 of Volume~III \\
3801 & l.~16 & \textit{for} interesction \textit{read} intersection \\
3802 & l.~5 & \textit{for} p.~3716 \textit{read} p.~3715 (Lemma~1661) \\
\dag\,3802 & proof of Lemma~2399, l.~14 & \textit{for} Lemma~2074 \textit{read} Lemma~2070 \\
\dag\,3802 & proof of Lemma~2402, l.~19 & \textit{for} Lemma~1859 \textit{read} Lemma~1861 \\
\dag\,3802 & proof of Lemma~2402, l.~29 & \textit{for} Lemma~2155 \textit{read} Lemma~2154 \\
3803 & Lemma~2417 and following paragraph & \textit{see} Corrigendum~S2 \\
3805 & l.~5 & \textit{for} Definition~5.20 of Volume~III \textit{read} Definition~5.15 of Volume~III; the result cited is on p.~996 \\
\dag\,3805 & proof of Lemma~2434, l.~6 & \textit{for} Lemma~2049 \textit{read} Lemma~2048 \\
\dag\,3805 & proof of Lemma~2433, l.~13 & \textit{for} Lemma~2112 \textit{read} Lemma~2116; the step uses Lemma~2116 (p.~3767) \\
\dag\,3805 & proof of Lemma~2426, l.~29 & \textit{for} Lemma~1962 \textit{read} Lemma~1958; the step uses Lemma~1958 (p.~3748) \\
3806 & l.~32 & \textit{for} Theorem~5.6 of Volume~III \textit{read} Theorem~5.4 of Volume~III; the result cited is on p.~965 \\
3806 & l.~1 from foot & \textit{for} p.~3736 \textit{read} p.~3737 (Lemma~1866) \\
\dag\,3807 & proof of Lemma~2448, l.~21 & \textit{for} Lemma~1966 \textit{read} Lemma~1964; the step uses Lemma~1964 (p.~3749) \\
\dag\,3807 & proof of Lemma~2438, l.~24 & \textit{for} Lemma~2213 \textit{read} Lemma~2218; the step uses Lemma~2218 (p.~3779) \\
3808 & l.~5 & \textit{for} p.~3878 \textit{read} p.~3778 (Lemma~2206) \\
3808 & l.~11 & \textit{for} p.~3757 \textit{read} p.~3767 (Lemma~2118) \\
3808 & l.~17 & \textit{for} p.~3638 \textit{read} p.~3637 (Lemma~873) \\
3808 & l.~31 & \textit{for} p.~3697 \textit{read} p.~3687 (Lemma~1347) \\
\dag\,3809 & proof of Lemma~2467, l.~8 & \textit{for} Lemma~7.4 of Volume~III \textit{read} Lemma~7.9 of Volume~III; the result cited is on p.~1105 \\
\dag\,3809 & proof of Lemma~2467, (5.39) & \textit{for} $\pos\forall z\,\mathrm{MG}(z) \to \nec\exists z\,\mathrm{MG}(z)$ \textit{read} $\pos\exists z\,\mathrm{MG}(z) \to \nec\exists z\,\mathrm{MG}(z)$ \\
\dag\,3810 & proof of Lemma~2469, l.~17 & \textit{for} Lemma~2053 \textit{read} Lemma~2051 \\
3810 & l.~7 from foot & \textit{for} p.~3785 \textit{read} p.~3795 (Lemma~2350) \\
\dag\,3811 & proof of Lemma~2484, (5.40) & \textit{for} $\mathrm{ME}(z,w_2) \iff \mathrm{Op}(z,w_2) \wedge \mathrm{Os}(z,w_2) \vee \mathrm{MP}(z,w_2)$ \textit{read} $\mathrm{ME}(z,w_2) \iff \mathrm{Op}(z,w_2) \wedge \mathrm{Os}(z,w_2) \wedge \mathrm{MP}(z,w_2)$ \\
\dag\,3811 & proof of Lemma~2470, l.~19 & \textit{for} Lemma~1901 \textit{read} Lemma~1896 \\
3812 & l.~5 & \textit{for} p.~3750 \textit{read} p.~3570 (Lemma~287) \\
\dag\,3813 & proof of Lemma~2494, l.~9 & \textit{for} Lemma~2118 \textit{read} Lemma~2122; the step uses Lemma~2122 (p.~3768) \\
\dag\,3813 & proof of Lemma~2499, l.~23 & \textit{for} Lemma~2282 \textit{read} Lemma~2283; the step uses Lemma~2283 (p.~3787) \\
\dag\,3813 & proof of Lemma~2494, l.~29 & \textit{for} Lemma~1886 \textit{read} Lemma~1883; the step uses Lemma~1883 (p.~3739) \\
3815 & l.~5 & \textit{for} p.~3559 \textit{read} p.~3595 (Lemma~506) \\
\dag\,3815 & proof of Lemma~2510, l.~9 & \textit{for} Lemma~2252 \textit{read} Lemma~2256 \\
\dag\,3815 & proof of Lemma~2517, (5.43) & \textit{for} $\mathrm{NE}(y,v) \iff \forall v\,(wRv \to E(y,v))$ \textit{read} $\mathrm{NE}(y,w) \iff \forall v\,(wRv \to E(y,v))$; correctly printed in (5.44), p.~3819 \\
\dag\,3815 & proof of Lemma~2513, l.~27 & \textit{for} Lemma~2379 \textit{read} Lemma~2384 \\
3815 & l.~1 from foot & \textit{for} p.~3622 \textit{read} p.~3722 (Lemma~1735) \\
\dag\,3816 & proof of Lemma~2519, l.~8 & \textit{for} Lemma~2359 \textit{read} Lemma~2361; the step uses Lemma~2361 (p.~3797) \\
3816 & l.~10 & \textit{for} p.~3647 \textit{read} p.~3747 (Lemma~1944) \\
\dag\,3816 & proof of Lemma~2511, l.~15 & \textit{for} Lemma~1858 \textit{read} Lemma~1859 \\
3816 & l.~25 & \textit{for} p.~3805 \textit{read} p.~3806 (Lemma~2437) \\
\dag\,3817 & proof of Lemma~2522, l.~28 & \textit{for} Lemma~2254 \textit{read} Lemma~2249 \\
\dag\,3818 & proof of Lemma~2532, l.~2 & \textit{for} Lemma~2049 \textit{read} Lemma~2053; the step uses Lemma~2053 (p.~3760) \\
\dag\,3818 & proof of Lemma~2542, l.~31 & \textit{for} Definition~5.5 of Volume~III \textit{read} Definition~5.1 of Volume~III; the result cited is on p.~952 \\
3819 & l.~4 & \textit{for} Corollary~5.9 of Volume~III \textit{read} Corollary~5.12 of Volume~III \\
\dag\,3819 & proof of Lemma~2548, l.~24 & \textit{for} Lemma~2010 \textit{read} Lemma~2015 \\
3819 & l.~4 from foot & \textit{for} p.~3661 \textit{read} p.~3561 (Lemma~192) \\
3820 & l.~4 & \textit{for} p.~3823 \textit{read} p.~3813 (Lemma~2497) \\
3820 & l.~6 & \textit{for} essance \textit{read} essence \\
3820 & l.~26 & \textit{for} p.~3830 \textit{read} p.~3803 (Lemma~2414) \\
\dag\,3820 & proof of Lemma~2559, l.~34 & \textit{for} Corollary~5.12 of Volume~III \textit{read} Corollary~5.9 of Volume~III \\
3821 & l.~19 & \textit{for} p.~3586 \textit{read} p.~3585 (Lemma~409) \\
\dag\,3821 & proof of Lemma~2555, l.~29 & \textit{for} Lemma~2006 \textit{read} Lemma~2001 \\
\dag\,3821 & proof of Lemma~2567, (5.45) & \textit{for} ${\phi} \in \mathcal{P} \wedge \pos\forall x\,({\phi}x \to {\chi}x) \Rightarrow {\chi} \in \mathcal{P}$ \textit{read} ${\phi} \in \mathcal{P} \wedge \nec\forall x\,({\phi}x \to {\chi}x) \Rightarrow {\chi} \in \mathcal{P}$ \\
3823 & l.~4 & \textit{for} p.~3704 \textit{read} p.~3705 (Lemma~1551) \\
\dag\,3823 & proof of Lemma~2581, l.~10 & \textit{for} Lemma~2447 \textit{read} Lemma~2449 \\
\dag\,3824 & proof of Lemma~2592, l.~7 & \textit{for} Lemma~7.7 of Volume~III \textit{read} Lemma~7.4 of Volume~III; the result cited is on p.~1088 \\
\dag\,3824 & proof of Lemma~2588, l.~11 & \textit{for} Lemma~2268 \textit{read} Lemma~2267 \\
\dag\,3824 & proof of Lemma~2580, l.~19 & \textit{for} Lemma~2387 \textit{read} Lemma~2388 \\
\dag\,3824 & proof of Lemma~2592, l.~22 & \textit{for} Lemma~2059 \textit{read} Lemma~2055 \\
\dag\,3825 & proof of Lemma~2592, l.~4 & \textit{for} Lemma~2352 \textit{read} Lemma~2347 \\
\dag\,3825 & proof of Lemma~2593, l.~14 & \textit{for} Lemma~2591 \textit{read} Lemma~2588 \\
3825 & l.~27 & \textit{for} Proposition~5.11 of Volume~III \textit{read} Proposition~5.16 of Volume~III; the result cited is on p.~998 \\
\dag\,3825 & proof of Lemma~2601, l.~4 from foot & \textit{for} $\mathcal{P}_0 \supseteq \mathcal{P} \Rightarrow \exists z\,\exists v'\,\forall {\phi} \in \mathcal{P}_0\; {\phi}(z,v')$ \textit{read} $\mathcal{P}_0 \subseteq \mathcal{P} \Rightarrow \exists z\,\exists v'\,\forall {\phi} \in \mathcal{P}_0\; {\phi}(z,v')$ \\
\dag\,3826 & proof of Lemma~2598, l.~9 & \textit{for} Lemma~2122 \textit{read} Lemma~2119 \\
\dag\,3826 & proof of Lemma~2603, l.~12 & \textit{for} Lemma~2581 \textit{read} Lemma~2583 \\
3826 & l.~20 & \textit{for} Corollary~7.3 of Volume~III \textit{read} Corollary~7.5 of Volume~III; the result cited is on p.~1090 \\
3826 & l.~24 & \textit{for} p.~3713 \textit{read} p.~3813 (Lemma~2498) \\
\dag\,3827 & proof of Lemma~2617, (5.47) & \textit{for} $\mathrm{G}z \to \forall {\psi}\,(\mathrm{Pos}({\psi}) \leftrightarrow \nec {\psi}z)$ \textit{read} $\mathrm{G}z \leftrightarrow \forall {\psi}\,(\mathrm{Pos}({\psi}) \leftrightarrow \nec {\psi}z)$ \\
\dag\,3827 & proof of Lemma~2603, l.~19 & \textit{for} Lemma~1828 \textit{read} Lemma~1825 \\
3827 & l.~26 & \textit{for} p.~3564 \textit{read} p.~3565 (Lemma~235) \\
\dag\,3828 & proof of Lemma~2626, (5.49) & \textit{for} $\mathrm{NME}(y,w_0) \iff \exists v\,(w_0Rv \to \mathrm{ME}(y,v))$ \textit{read} $\mathrm{NME}(y,w_0) \iff \forall v\,(w_0Rv \to \mathrm{ME}(y,v))$ \\
\dag\,3828 & proof of Lemma~2626, (5.50) & \textit{for} $\mathrm{MG}(z,u) \iff \exists w_2\,(uRw_2 \to E(z,w_2) \wedge \mathrm{ME}(z,w_2))$ \textit{read} $\mathrm{MG}(z,u) \iff \forall w_2\,(uRw_2 \to E(z,w_2) \wedge \mathrm{ME}(z,w_2))$ \\
\dag\,3829 & proof of Lemma~2626, l.~17 & \textit{for} Lemma~2167 \textit{read} Lemma~2162; the step uses Lemma~2162 (p.~3773) \\
3829 & l.~28 & \textit{for} p.~3735 \textit{read} p.~3753 (Lemma~1999) \\
\dag\,3830 & proof of Lemma~2642, (5.52) & \textit{for} ${\phi} \in \mathcal{P} \vee \nec\forall z\,({\phi}z \to {\chi}z) \Rightarrow {\chi} \in \mathcal{P}$ \textit{read} ${\phi} \in \mathcal{P} \wedge \nec\forall z\,({\phi}z \to {\chi}z) \Rightarrow {\chi} \in \mathcal{P}$ \\
\dag\,3830 & proof of Lemma~2636, l.~17 & \textit{for} Lemma~2493 \textit{read} Lemma~2488; the step uses Lemma~2488 (p.~3812) \\
3830 & l.~2 from foot & \textit{for} p.~3766 \textit{read} p.~3676 (Lemma~1246) \\
\dag\,3831 & proof of Lemma~2646, l.~23 & \textit{for} Lemma~2088 \textit{read} Lemma~2093 \\
\dag\,3831 & proof of Lemma~2650, l.~28 & \textit{for} Definition~5.13 of Volume~III \textit{read} Definition~5.18 of Volume~III \\
\dag\,3832 & proof of Lemma~2649, l.~19 & \textit{for} Lemma~2126 \textit{read} Lemma~2131 \\
3832 & l.~38 & \textit{for} p.~3830 \textit{read} p.~3831 (Lemma~2645) \\
\end{errlist}

\subsection*{Chapter~6. Formal material for Volumes IV and V \hfill \normalfont\small pp.~3833--3874}
\begin{errlist}
3834 & running heads, pp.~3834--3874 (even) & \textit{for} Formal material for Volumes VI to IX \textit{read} Formal material for Volumes IV and V \\
\dag\,3835 & proof of Lemma~2665, l.~10 & \textit{for} Lemma~2655 \textit{read} Lemma~2654 \\
3835 & l.~35 & \textit{for} $\neg K_{z}q \leftrightarrow K_{z}\neg K_{z}q$ \textit{read} $\neg K_{z}q \to K_{z}\neg K_{z}q$ \\
\dag\,3836 & proof of Lemma~2680, l.~6 from foot & \textit{for} Definition~6.2 of Volume~IV \textit{read} Definition~6.8 of Volume~IV \\
3837 & l.~14 & \textit{for} p.~3743 \textit{read} p.~3753 (Lemma~1998) \\
\dag\,3837 & proof of Lemma~2690, l.~24 & \textit{for} Lemma~2672 \textit{read} Lemma~2667 \\
3838 & l.~24 & \textit{for} p.~3564 \textit{read} p.~3664 (Lemma~1133) \\
3838 & l.~32 & \textit{for} $\mathrm{NE}(z) \to \exists x\,D(z,x)$ \textit{read} $\mathrm{NE}(z) \to \neg\exists x\,D(z,x)$ \\
3838 & l.~8 from foot & \textit{after} domain \textit{insert} of the world $v$ \\
\dag\,3839 & proof of Lemma~2705, l.~9 & \textit{for} Lemma~2680 \textit{read} Lemma~2685 \\
\dag\,3839 & proof of Lemma~2710, (6.1) & \textit{for} $w_1Ru \leftrightarrow D_{u} \subseteq D_{w_1}$ \textit{read} $w_1Ru \to D_{u} \subseteq D_{w_1}$; correctly printed in (6.2), p.~3840 \\
\dag\,3839 & proof of Lemma~2700, l.~26 & \textit{for} Lemma~2658 \textit{read} Lemma~2655 \\
\dag\,3840 & proof of Lemma~2701, l.~6 & \textit{for} Lemma~2657 \textit{read} Lemma~2661 \\
\dag\,3840 & proof of Lemma~2717, l.~8 & \textit{for} Lemma~2680 \textit{read} Lemma~2678 \\
\dag\,3840 & proof of Lemma~2719, (6.3) & \textit{for} $\nec\forall y\,{\psi}y \leftrightarrow \forall y\,\nec {\psi}y$ \textit{read} $\nec\forall y\,{\psi}y \to \forall y\,\nec {\psi}y$ \\
\dag\,3841 & proof of Lemma~2725, l.~2 & \textit{for} Lemma~2661 \textit{read} Lemma~2658 \\
\dag\,3841 & proof of Lemma~2725, l.~20 & \textit{for} Lemma~2714 \textit{read} Lemma~2716; the step uses Lemma~2716 (p.~3840) \\
3843 & l.~2 & \textit{for} p.~3724 \textit{read} p.~3725 (Lemma~1776) \\
\dag\,3843 & proof of Lemma~2749, l.~4 & \textit{for} Lemma~2738 \textit{read} Lemma~2733 \\
3844 & l.~39 & \textit{for} $uRv' \to D_{u} \subseteq D_{u}$ \textit{read} $uRv' \to D_{u} \subseteq D_{v'}$ \\
\dag\,3845 & proof of Lemma~2767, l.~13 & \textit{for} Lemma~2722 \textit{read} Lemma~2720 \\
\dag\,3845 & proof of Lemma~2765, l.~15 & \textit{for} Lemma~2731 \textit{read} Lemma~2734; the step uses Lemma~2734 (p.~3842) \\
\dag\,3845 & proof of Lemma~2758, l.~22 & \textit{for} Lemma~2741 \textit{read} Lemma~2744; the step uses Lemma~2744 (p.~3843) \\
\dag\,3845 & proof of Lemma~2759, l.~25 & \textit{for} Lemma~2685 \textit{read} Lemma~2687 \\
\dag\,3846 & proof of Lemma~2761, l.~15 & \textit{for} Lemma~2705 \textit{read} Lemma~2700; the step uses Lemma~2700 (p.~3838) \\
\dag\,3847 & proof of Lemma~2776, l.~6 & \textit{for} Lemma~2718 \textit{read} Lemma~2716 \\
\dag\,3847 & proof of Lemma~2776, l.~28 & \textit{for} Lemma~2703 \textit{read} Lemma~2698 \\
3848 & l.~7 & \textit{for} $\mathrm{Os}(z) \to \forall p\,(p \to K_{z}p)$ \textit{read} $\mathrm{Os}(z) \leftrightarrow \forall p\,(p \to K_{z}p)$ \\
\dag\,3848 & proof of Lemma~2798, l.~22 & \textit{for} $\mathrm{Os}(y) \wedge \mathrm{Os}(z) \to \forall p\,(K_{y}p \to K_{z}p)$ \textit{read} $\mathrm{Os}(y) \wedge \mathrm{Os}(z) \to \forall p\,(K_{y}p \leftrightarrow K_{z}p)$ \\
\dag\,3848 & proof of Lemma~2798, (6.5) & \textit{for} $K_{z}p \leftrightarrow K_{z}K_{z}p$ \textit{read} $K_{z}p \to K_{z}K_{z}p$; correctly printed in (6.9), p.~3853 \\
\dag\,3849 & proof of Lemma~2799, l.~24 & \textit{for} Lemma~2672 \textit{read} Lemma~2675 \\
\dag\,3849 & proof of Lemma~2796, l.~30 & \textit{for} Lemma~2752 \textit{read} Lemma~2753 \\
3849 & l.~33 & \textit{for} p.~3707 \textit{read} p.~3807 (Lemma~2444) \\
3849 & l.~39 & \textit{for} $\mathrm{Os}(z) \vee \mathrm{Os}(x) \to \forall q\,(K_{z}q \leftrightarrow K_{x}q)$ \textit{read} $\mathrm{Os}(z) \wedge \mathrm{Os}(x) \to \forall q\,(K_{z}q \leftrightarrow K_{x}q)$ \\
\dag\,3850 & proof of Lemma~2818, (6.6) & \textit{for} $\nec\exists x\,{\chi}x \to \exists x\,\pos {\chi}x$ \textit{read} $\pos\exists x\,{\chi}x \to \exists x\,\pos {\chi}x$; correctly printed in (6.10), p.~3853 \\
\dag\,3850 & proof of Lemma~2814, l.~9 & \textit{for} Lemma~2711 \textit{read} Lemma~2710 \\
\dag\,3850 & proof of Lemma~2811, l.~10 & \textit{for} Lemma~2803 \textit{read} Lemma~2798 \\
3850 & l.~13 & \textit{for} $\forall z\,\nec {\chi}z \to \pos\forall z\,{\chi}z$ \textit{read} $\forall z\,\nec {\chi}z \to \nec\forall z\,{\chi}z$ \\
\dag\,3850 & proof of Lemma~2818, l.~39 & \textit{for} Definition~6.1 of Volume~IV \textit{read} Definition~6.2 of Volume~IV; the result cited is on p.~1589 \\
3850 & n.~22 & \textit{for} Grim (1990) \textit{read} Grim (1991) \\
\dag\,3851 & proof of Lemma~2828, (6.8) & \textit{for} $\exists x\,\pos {\phi}x \leftrightarrow \pos\exists x\,{\phi}x$ \textit{read} $\exists x\,\pos {\phi}x \to \pos\exists x\,{\phi}x$; the formula is used in its correct form on p.~3853 \\
3851 & l.~32 & \textit{for} p.~3559 \textit{read} p.~3659 (Lemma~1083) \\
\dag\,3851 & proof of Lemma~2828, (6.7) & \textit{for} $K_{y}q \leftrightarrow q$ \textit{read} $K_{y}q \to q$; correctly printed in (6.13), p.~3856 \\
\dag\,3852 & proof of Lemma~2832, l.~2 & \textit{for} Lemma~2700 \textit{read} Lemma~2696; the step uses Lemma~2696 (p.~3838) \\
\dag\,3852 & proof of Lemma~2837, l.~8 & \textit{for} Lemma~2684 \textit{read} Lemma~2679 \\
3852 & l.~13 & \textit{delete} the second occurrence of \emph{the} \\
\dag\,3853 & proof of Lemma~2836, l.~6 & \textit{for} Lemma~2725 \textit{read} Lemma~2727 \\
\dag\,3853 & proof of Lemma~2843, l.~13 & \textit{for} Lemma~2699 \textit{read} Lemma~2698 \\
\dag\,3854 & proof of Lemma~2851, l.~10 & \textit{for} Lemma~2761 \textit{read} Lemma~2762 \\
\dag\,3854 & proof of Lemma~2857, l.~24 & \textit{for} Definition~3.7 of Volume~V \textit{read} Definition~3.3 of Volume~V \\
3855 & l.~2 & \textit{after} every task \textit{insert} whose description does not quantify over the agent's own powers \\
3855 & l.~4 & \textit{after} every true proposition \textit{insert} (plurally quantified) \\
\dag\,3855 & proof of Lemma~2859, l.~5 & \textit{for} Lemma~2831 \textit{read} Lemma~2827 \\
\dag\,3856 & proof of Lemma~2875, l.~5 & \textit{for} Lemma~2655 \textit{read} Lemma~2651; the step uses Lemma~2651 (p.~3834) \\
\dag\,3856 & proof of Lemma~2877, (6.11) & \textit{for} $K_{z}p \leftrightarrow p$ \textit{read} $K_{z}p \to p$ \\
\dag\,3856 & proof of Lemma~2877, (6.12) & \textit{for} $\mathrm{Op}(y) \leftrightarrow \forall t\,(\mathrm{Task}(t) \vee \neg Q(t,y) \to \mathrm{Can}(y,t))$ \textit{read} $\mathrm{Op}(y) \leftrightarrow \forall t\,(\mathrm{Task}(t) \wedge \neg Q(t,y) \to \mathrm{Can}(y,t))$ \\
\dag\,3856 & proof of Lemma~2875, l.~27 & \textit{for} Lemma~2832 \textit{read} Lemma~2837 \\
\dag\,3856 & proof of Lemma~2877, l.~37 & \textit{for} Lemma~3.14 of Volume~V \textit{read} Lemma~3.17 of Volume~V; the result cited is on p.~1870 \\
\dag\,3857 & proof of Lemma~2883, l.~2 & \textit{for} Lemma~2848 \textit{read} Lemma~2851; the step uses Lemma~2851 (p.~3854) \\
\dag\,3857 & proof of Lemma~2885, l.~4 & \textit{for} Lemma~2734 \textit{read} Lemma~2732 \\
\dag\,3857 & proof of Lemma~2883, l.~7 & \textit{for} Lemma~2812 \textit{read} Lemma~2809 \\
3857 & l.~8 & \textit{for} p.~3782 \textit{read} p.~3781 (Lemma~2229) \\
3857 & l.~11 & \textit{for} Theorem~5.5 of Volume~V \textit{read} Theorem~5.2 of Volume~V \\
\dag\,3857 & proof of Lemma~2887, l.~22 & \textit{for} Lemma~2654 \textit{read} Lemma~2651 \\
3857 & l.~8 from foot & \textit{for} Lemma~4.9 of Volume~IV \textit{read} Lemma~4.3 of Volume~IV; the result cited is on p.~1510 \\
3858 & l.~6 & \textit{for} Grim (1992) \textit{read} Grim (1991) \\
3859 & l.~4 & \textit{for} Definition~3.2 of Volume~IV \textit{read} Definition~3.4 of Volume~IV \\
\dag\,3859 & proof of Lemma~2900, l.~9 & \textit{for} Lemma~2706 \textit{read} Lemma~2702; the step uses Lemma~2702 (p.~3839) \\
\dag\,3859 & proof of Lemma~2907, (6.15) & \textit{for} $E(z,v') \to z \in D_{v'}$ \textit{read} $E(z,v') \leftrightarrow z \in D_{v'}$ \\
\dag\,3859 & proof of Lemma~2907, l.~8 from foot & \textit{for} $\nec(\mathrm{Op}(z) \wedge \mathrm{Os}(z) \vee \mathrm{MP}(z))$ \textit{read} $\nec(\mathrm{Op}(z) \wedge \mathrm{Os}(z) \wedge \mathrm{MP}(z))$ \\
\dag\,3860 & proof of Lemma~2900, l.~5 & \textit{for} Lemma~2880 \textit{read} Lemma~2884; the step uses Lemma~2884 (p.~3857) \\
3861 & l.~31 & \textit{for} p.~3935 \textit{read} p.~3835 (Lemma~2661) \\
\dag\,3862 & proof of Lemma~2928, l.~28 & \textit{for} Lemma~2839 \textit{read} Lemma~2844 \\
\dag\,3863 & proof of Lemma~2943, l.~4 & \textit{for} Lemma~2891 \textit{read} Lemma~2893 \\
\dag\,3863 & proof of Lemma~2939, l.~14 & \textit{for} Lemma~2690 \textit{read} Lemma~2695 \\
\dag\,3863 & proof of Lemma~2942, l.~18 & \textit{for} Lemma~2876 \textit{read} Lemma~2879; the step uses Lemma~2879 (p.~3857) \\
\dag\,3866 & proof of Lemma~2976, (6.17) & \textit{for} $uRw_0 \leftrightarrow D_{u} \subseteq D_{w_0}$ \textit{read} $uRw_0 \to D_{u} \subseteq D_{w_0}$; correctly printed in (6.20), p.~3872 \\
\dag\,3867 & proof of Lemma~2971, l.~2 & \textit{for} Lemma~2758 \textit{read} Lemma~2757 \\
\dag\,3867 & proof of Lemma~2973, l.~4 & \textit{for} Lemma~2849 \textit{read} Lemma~2854 \\
3867 & l.~12 & \textit{for} p.~3738 \textit{read} p.~3739 (Lemma~1879) \\
3867 & l.~14 & \textit{for} $\pos\exists y\,{\psi}y \to \exists y\,\nec {\psi}y$ \textit{read} $\pos\exists y\,{\psi}y \to \exists y\,\pos {\psi}y$ \\
\dag\,3867 & proof of Lemma~2985, l.~17 & \textit{for} Theorem~3.7 of Volume~IV \textit{read} Theorem~3.13 of Volume~IV \\
\dag\,3867 & proof of Lemma~2984, l.~24 & \textit{for} Lemma~2697 \textit{read} Lemma~2699 \\
\dag\,3867 & proof of Lemma~2985, (6.18) & \textit{for} $\mathrm{NE}(x) \to \neg\exists x\,D(x,z)$ \textit{read} $\mathrm{NE}(x) \to \neg\exists z\,D(x,z)$ \\
\dag\,3868 & proof of Lemma~2985, l.~14 & \textit{for} Lemma~2912 \textit{read} Lemma~2907 \\
3868 & l.~32 & \textit{for} p.~3762 \textit{read} p.~3726 (Lemma~1787) \\
3869 & l.~5 & \textit{for} Lemma~6.6 of Volume~IV \textit{read} Lemma~6.4 of Volume~IV \\
\dag\,3870 & proof of Lemma~3006, l.~9 & \textit{for} Lemma~2998 \textit{read} Lemma~3003 \\
\dag\,3870 & proof of Lemma~3004, l.~11 & \textit{for} Lemma~2800 \textit{read} Lemma~2798 \\
3870 & l.~20 & \textit{for} $\pos\forall z\,\nec\exists x\,(x = z)$ \textit{read} $\nec\forall z\,\nec\exists x\,(x = z)$ \\
3870 & l.~23 & \textit{for} Barcan (1947) \textit{read} Barcan (1946) \\
3870 & l.~30 & \textit{for} Definition~3.3 of Volume~V \textit{read} Definition~3.7 of Volume~V; the result cited is on p.~1843 \\
3871 & l.~16 & \textit{for} p.~3585 \textit{read} p.~3558 (Lemma~162) \\
\dag\,3871 & proof of Lemma~3022, l.~22 & \textit{for} Lemma~2736 \textit{read} Lemma~2739 \\
\dag\,3872 & proof of Lemma~3035, l.~4 & \textit{for} Proposition~3.10 of Volume~IV \textit{read} Proposition~3.6 of Volume~IV; the result cited is on p.~1466 \\
3872 & l.~12 & \textit{for} Corollary~5.9 of Volume~V \textit{read} Corollary~5.6 of Volume~V \\
3872 & l.~20 & \textit{for} p.~3748 \textit{read} p.~3749 (Lemma~1967) \\
\dag\,3872 & proof of Lemma~3035, (6.21) & \textit{for} $w_1Rw_0 \to D_{w_1} \subseteq D_{w_0}$ \textit{read} $w_0Rw_1 \to D_{w_1} \subseteq D_{w_0}$ \\
3872 & l.~29 & \textit{for} p.~3861 \textit{read} p.~3871 (Lemma~3023) \\
3873 & l.~11 & \textit{for} p.~3774 \textit{read} p.~3773 (Lemma~2168) \\
\dag\,3873 & proof of Lemma~3044, l.~12 & \textit{for} $\mathrm{Op}(z) \leftrightarrow \exists t\,(\mathrm{Task}(t) \wedge \neg Q(t,z) \to \mathrm{Can}(z,t))$ \textit{read} $\mathrm{Op}(z) \leftrightarrow \forall t\,(\mathrm{Task}(t) \wedge \neg Q(t,z) \to \mathrm{Can}(z,t))$ \\
\dag\,3873 & proof of Lemma~3031, l.~21 & \textit{for} Lemma~2967 \textit{read} Lemma~2962 \\
3874 & l.~31 & \textit{for} powr \textit{read} power \\
\dag\,3874 & proof of Lemma~3044, l.~33 & \textit{for} $\forall x\,\nec {\chi}x \to \pos\forall x\,{\chi}x$ \textit{read} $\forall x\,\nec {\chi}x \to \nec\forall x\,{\chi}x$ \\
\end{errlist}

\subsection*{Chapter~7. Formal material for Volumes VI to IX \hfill \normalfont\small pp.~3875--3912}
\begin{errlist}
\dag\,3876 & proof of Lemma~3051, l.~7 & \textit{for} Lemma~3048 \textit{read} Lemma~3045 \\
\dag\,3876 & proof of Lemma~3052, l.~22 & \textit{for} Lemma~3050 \textit{read} Lemma~3049 \\
\dag\,3877 & proof of Lemma~3054, l.~7 & \textit{for} Lemma~3049 \textit{read} Lemma~3045 \\
\dag\,3877 & proof of Lemma~3059, l.~10 & \textit{for} $\pos\,\mathrm{Creates}(g,w_2) \vee \pos\,\neg\mathrm{Creates}(g,w_2)$ \textit{read} $\pos\,\mathrm{Creates}(g,w_2) \wedge \pos\,\neg\mathrm{Creates}(g,w_2)$ \\
3877 & l.~22 & \textit{for} p.~3659 \textit{read} p.~3658 (Lemma~1077) \\
\dag\,3878 & proof of Lemma~3058, l.~16 & \textit{for} Lemma~3055 \textit{read} Lemma~3057; the step uses Lemma~3057 (p.~3877) \\
3878 & l.~17 & \textit{for} sceptcal \textit{read} sceptical \\
\dag\,3878 & proof of Lemma~3059, l.~27 & \textit{for} Lemma~3054 \textit{read} Lemma~3056 \\
3878 & l.~8 from foot & \textit{after} intense suffering \textit{insert} that it could prevent \\
\dag\,3879 & proof of Lemma~3073, l.~22 & \textit{for} Lemma~7.9 of Volume~VI \textit{read} Lemma~7.5 of Volume~VI \\
\dag\,3879 & proof of Lemma~3065, l.~25 & \textit{for} Lemma~3063 \textit{read} Lemma~3058 \\
3879 & l.~27 & \textit{for} Kripke (1980) \textit{read} Rowe (1979) (the evidential argument from evil) \\
3879 & l.~30 & \textit{for} p.~3804 \textit{read} p.~3704 (Lemma~1533) \\
\dag\,3880 & proof of Lemma~3077, l.~30 & \textit{for} Lemma~3054 \textit{read} Lemma~3057; the step uses Lemma~3057 (p.~3877) \\
3881 & l.~5 & \textit{for} undeetrmined \textit{read} undetermined \\
3881 & l.~19 & \textit{for} p.~3690 \textit{read} p.~3590 (Lemma~460) \\
\dag\,3882 & proof of Lemma~3089, l.~4 & \textit{for} Lemma~3075 \textit{read} Lemma~3076 \\
\dag\,3882 & proof of Lemma~3091, l.~13 & \textit{for} Lemma~3049 \textit{read} Lemma~3046 \\
\dag\,3882 & proof of Lemma~3086, l.~15 & \textit{for} Lemma~3046 \textit{read} Lemma~3050; the step uses Lemma~3050 (p.~3876) \\
3882 & l.~32 & \textit{for} grautitous \textit{read} gratuitous \\
3883 & l.~4 & \textit{for} p.~3587 \textit{read} p.~3577 (Lemma~356) \\
\dag\,3883 & proof of Lemma~3098, l.~25 & \textit{for} Lemma~3069 \textit{read} Lemma~3068 \\
\dag\,3884 & proof of Lemma~3103, l.~20 & \textit{for} Lemma~3051 \textit{read} Lemma~3046 \\
\dag\,3884 & proof of Lemma~3104, l.~25 & \textit{for} Lemma~3086 \textit{read} Lemma~3091; the step uses Lemma~3091 (p.~3882) \\
\dag\,3884 & proof of Lemma~3107, l.~29 & \textit{for} Lemma~3079 \textit{read} Lemma~3076 \\
3884 & l.~2 from foot & \textit{for} Definition~3.5 of Volume~VIII \textit{read} Definition~3.7 of Volume~VIII \\
3885 & l.~24 & \textit{for} undeterimned \textit{read} undetermined \\
3885 & l.~6 from foot & \textit{for} Definition~3.7 of Volume~VIII \textit{read} Definition~3.1 of Volume~VIII; the result cited is on p.~2968 \\
3887 & l.~28 & \textit{after} causal role \textit{insert} of the maximally great being \\
\dag\,3889 & proof of Lemma~3132, l.~21 & \textit{for} Lemma~3080 \textit{read} Lemma~3085 \\
\dag\,3890 & proof of Lemma~3147, l.~8 & \textit{for} Lemma~3131 \textit{read} Lemma~3126 \\
\dag\,3890 & proof of Lemma~3145, l.~22 & \textit{for} Lemma~3148 \textit{read} Lemma~3143 \\
\dag\,3891 & proof of Lemma~3157, (7.9) & \textit{for} $\exists v'\,\exists w_0\,V(w_0) > V(v')$ \textit{read} $\forall v'\,\exists w_0\,V(w_0) > V(v')$; correctly printed in (7.13), p.~3895 \\
\dag\,3892 & proof of Lemma~3162, l.~2 & \textit{for} Lemma~3047 \textit{read} Lemma~3046 \\
\dag\,3892 & proof of Lemma~3153, l.~26 & \textit{for} Lemma~3062 \textit{read} Lemma~3063 \\
\dag\,3893 & proof of Lemma~3166, l.~4 & \textit{for} Lemma~3122 \textit{read} Lemma~3127; the step uses Lemma~3127 (p.~3887) \\
\dag\,3893 & proof of Lemma~3158, l.~7 & \textit{for} Lemma~3066 \textit{read} Lemma~3064; the step uses Lemma~3064 (p.~3878) \\
\dag\,3893 & proof of Lemma~3161, l.~14 & \textit{for} Lemma~3062 \textit{read} Lemma~3066 \\
\dag\,3894 & proof of Lemma~3176, l.~1 & \textit{for} Lemma~3077 \textit{read} Lemma~3073 \\
\dag\,3894 & proof of Lemma~3178, l.~7 & \textit{for} Definition~7.16 of Volume~VI \textit{read} Definition~7.11 of Volume~VI \\
\dag\,3894 & proof of Lemma~3169, l.~18 & \textit{for} Lemma~3160 \textit{read} Lemma~3163 \\
3894 & l.~25 & \textit{after} indiscernible \textit{insert} at every accessible world \\
3894 & l.~1 from foot & \textit{for} p.~3708 \textit{read} p.~3707 (Lemma~1573) \\
\dag\,3895 & proof of Lemma~3183, l.~15 & \textit{for} Lemma~3048 \textit{read} Lemma~3047 \\
\dag\,3895 & proof of Lemma~3185, l.~37 & \textit{for} Theorem~7.4 of Volume~VI \textit{read} Theorem~7.7 of Volume~VI \\
3896 & l.~6 & \textit{for} evidnetial \textit{read} evidential \\
3897 & l.~10 & \textit{for} Definition~7.17 of Volume~VI \textit{read} Definition~7.11 of Volume~VI; the result cited is on p.~2497 \\
\dag\,3897 & proof of Lemma~3199, l.~11 & \textit{for} $x \neq x \to \nec(y \neq x)$ \textit{read} $y \neq x \to \nec(y \neq x)$ \\
\dag\,3897 & proof of Lemma~3197, l.~28 & \textit{for} Lemma~3111 \textit{read} Lemma~3107 \\
\dag\,3898 & proof of Lemma~3206, l.~14 & \textit{for} Lemma~3054 \textit{read} Lemma~3049; the step uses Lemma~3049 (p.~3876) \\
3898 & l.~4 from foot & \textit{for} p.~3713 \textit{read} p.~3703 (Lemma~1528) \\
\dag\,3899 & proof of Lemma~3203, l.~2 & \textit{for} Lemma~3088 \textit{read} Lemma~3092 \\
3900 & l.~9 & \textit{after} equilibrium \textit{insert} of the signalling game \\
\dag\,3900 & proof of Lemma~3220, l.~10 & \textit{for} Lemma~3083 \textit{read} Lemma~3080 \\
3900 & l.~32 & \textit{for} p.~3881 \textit{read} p.~3880 (Lemma~3078) \\
3901 & l.~12 & \textit{for} p.~3717 \textit{read} p.~3716 (Lemma~1669) \\
3901 & l.~15 & \textit{for} $x \neq x \to \nec(x \neq y)$ \textit{read} $x \neq y \to \nec(x \neq y)$ \\
\dag\,3901 & proof of Lemma~3227, l.~19 & \textit{for} Definition~7.11 of Volume~VI \textit{read} Definition~7.16 of Volume~VI; the result cited is on p.~2506 \\
3901 & l.~39 & \textit{for} $V(w_0) < V(w_1) \to \mathrm{Chooses}(g,w_0)$ \textit{read} $V(w_0) < V(w_1) \to \neg\mathrm{Chooses}(g,w_0)$ \\
3902 & l.~7 & \textit{for} Definition~4.15 of Volume~VI \textit{read} Definition~4.13 of Volume~VI \\
\dag\,3903 & proof of Lemma~3240, l.~2 & \textit{for} Lemma~3102 \textit{read} Lemma~3101; the step uses Lemma~3101 (p.~3883) \\
\dag\,3903 & proof of Lemma~3241, l.~3 & \textit{for} Lemma~3211 \textit{read} Lemma~3206; the step uses Lemma~3206 (p.~3898) \\
\dag\,3903 & proof of Lemma~3230, l.~24 & \textit{for} Lemma~3127 \textit{read} Lemma~3124 \\
\dag\,3904 & proof of Lemma~3239, l.~2 & \textit{for} Lemma~3059 \textit{read} Lemma~3061 \\
3904 & l.~4 & \textit{for} intetnion \textit{read} intention \\
\dag\,3904 & proof of Lemma~3248, l.~19 & \textit{for} Lemma~3227 \textit{read} Lemma~3228; the step uses Lemma~3228 (p.~3902) \\
3904 & l.~24 & \textit{for} p.~3892 \textit{read} p.~3829 (Lemma~2631) \\
3905 & l.~6 & \textit{for} Definition~3.1 of Volume~IX \textit{read} Definition~3.2 of Volume~IX \\
\dag\,3905 & proof of Lemma~3255, l.~25 & \textit{for} Lemma~3131 \textit{read} Lemma~3128; the step uses Lemma~3128 (p.~3887) \\
\dag\,3905 & proof of Lemma~3255, l.~26 & \textit{for} Lemma~3190 \textit{read} Lemma~3194 \\
\dag\,3905 & proof of Lemma~3244, l.~28 & \textit{for} Lemma~3111 \textit{read} Lemma~3113 \\
\dag\,3906 & proof of Lemma~3259, l.~28 & \textit{for} Lemma~3069 \textit{read} Lemma~3074 \\
3906 & l.~31 & \textit{for} p.~3740 \textit{read} p.~3741 (Lemma~1898) \\
3907 & l.~2 & \textit{for} p.~3669 \textit{read} p.~3659 (Lemma~1088) \\
\dag\,3907 & proof of Lemma~3268, l.~3 & \textit{for} Lemma~3093 \textit{read} Lemma~3090 \\
3908 & l.~2 & \textit{for} Definition~3.4 of Volume~IX \textit{read} Definition~3.2 of Volume~IX \\
3908 & l.~5 & \textit{for} $z = y \to \nec(y = y)$ \textit{read} $z = y \to \nec(z = y)$ \\
3908 & l.~20 & \textit{for} p.~3542 \textit{read} p.~3543 (Lemma~18) \\
\dag\,3908 & proof of Lemma~3274, l.~22 & \textit{for} Lemma~3276 \textit{read} Lemma~3271 \\
3908 & l.~28 & \textit{for} Lemma~3.13 of Volume~IX \textit{read} Lemma~3.10 of Volume~IX; the result cited is on p.~3290 \\
\dag\,3908 & proof of Lemma~3276, (7.16) & \textit{for} $\mathrm{MG}(z) \vee \mathrm{MG}(y) \to z = y$ \textit{read} $\mathrm{MG}(z) \wedge \mathrm{MG}(y) \to z = y$; the formula is used in its correct form on p.~3909 \\
\dag\,3909 & proof of Lemma~3272, l.~1 & \textit{for} Lemma~3140 \textit{read} Lemma~3145; the step uses Lemma~3145 (p.~3890) \\
\dag\,3909 & proof of Lemma~3283, l.~8 & \textit{for} Lemma~3073 \textit{read} Lemma~3076 \\
3909 & l.~10 & \textit{for} p.~3859 \textit{read} p.~3849 (Lemma~2801) \\
3909 & l.~29 & \textit{for} indiscrnibility \textit{read} indiscernibility \\
\dag\,3910 & proof of Lemma~3287, l.~16 & \textit{for} Lemma~3167 \textit{read} Lemma~3169 \\
\dag\,3910 & proof of Lemma~3282, l.~26 & \textit{for} Lemma~3241 \textit{read} Lemma~3242; the step uses Lemma~3242 (p.~3904) \\
3910 & l.~39 & \textit{for} p.~3662 \textit{read} p.~3663 (Lemma~1129) \\
\dag\,3911 & proof of Lemma~3297, (7.17) & \textit{for} $\forall {\phi}\,({\phi}y \to {\phi}z) \to y = z$ \textit{read} $\forall {\phi}\,({\phi}y \leftrightarrow {\phi}z) \to y = z$ \\
\dag\,3911 & proof of Lemma~3297, (7.19) & \textit{for} $\exists w_1\,(T(w_1) \vee E(w_1))$ \textit{read} $\exists w_1\,(T(w_1) \wedge E(w_1))$; the formula is used in its correct form on p.~3912 \\
\dag\,3911 & proof of Lemma~3297, l.~20 & \textit{for} $V(w_0) < V(w') \to \mathrm{Chooses}(g,w_0)$ \textit{read} $V(w_0) < V(w') \to \neg\mathrm{Chooses}(g,w_0)$ \\
3911 & l.~30 & \textit{for} nonreisstant \textit{read} nonresistant \\
\dag\,3912 & proof of Lemma~3297, l.~32 & \textit{for} Lemma~7.8 of Volume~VI \textit{read} Lemma~7.12 of Volume~VI; the result cited is on p.~2498 \\
\end{errlist}

\subsection*{Chapter~8. The Lean modules \hfill \normalfont\small pp.~3913--3939}
\begin{errlist}
3914 & l.~29 & \textit{for} tacitc \textit{read} tactic \\
\dag\,3915 & proof of Lemma~3301, l.~5 & \textit{for} Lemma~3304 \textit{read} Lemma~3299 \\
3916 & l.~14 & \textit{for} p.~3809 \textit{read} p.~3890 (Lemma~3150) \\
3916 & l.~31 & \textit{for} p.~3732 \textit{read} p.~3742 (Lemma~1908) \\
\dag\,3916 & proof of Lemma~3305, (8.1) & \textit{for} $w' \models \pos r \iff \exists w_0\,(w_0Rw' \wedge w_0 \models r)$ \textit{read} $w' \models \pos r \iff \exists w_0\,(w'Rw_0 \wedge w_0 \models r)$ \\
\dag\,3917 & proof of Lemma~3308, l.~6 & \textit{for} Lemma~3303 \textit{read} Lemma~3300 \\
\dag\,3918 & proof of Lemma~3307, l.~1 & \textit{for} Lemma~3301 \textit{read} Lemma~3300 \\
\dag\,3918 & proof of Lemma~3310, (8.4) & \textit{for} $w_2Rv \to w_2Rv$ \textit{read} $vRw_2 \to w_2Rv$ \\
3918 & l.~30 & \textit{for} necessary\_exstence \textit{read} necessary\_existence \\
3919 & l.~22 & \textit{for} p.~3643 \textit{read} p.~3743 (Lemma~1916) \\
\dag\,3920 & proof of Lemma~3312, l.~2 & \textit{for} Lemma~3307 \textit{read} Lemma~3311 \\
\dag\,3920 & proof of Lemma~3315, l.~11 & \textit{for} Lemma~3304 \textit{read} Lemma~3302; the step uses Lemma~3302 (p.~3915) \\
\dag\,3920 & proof of Lemma~3311, l.~14 & \textit{for} Lemma~3307 \textit{read} Lemma~3308 \\
\dag\,3920 & proof of Lemma~3311, l.~22 & \textit{for} Lemma~3306 \textit{read} Lemma~3302 \\
\dag\,3920 & proof of Lemma~3314, l.~27 & \textit{for} Lemma~3311 \textit{read} Lemma~3312 \\
\dag\,3921 & proof of Lemma~3318, l.~4 & \textit{for} $\pos(q \vee q) \leftrightarrow (\pos p \vee \pos q)$ \textit{read} $\pos(p \vee q) \leftrightarrow (\pos p \vee \pos q)$ \\
3921 & l.~12 & \textit{for} $\nec p \wedge \pos p \to \pos(q \wedge p)$ \textit{read} $\nec q \wedge \pos p \to \pos(q \wedge p)$ \\
\dag\,3921 & proof of Lemma~3316, l.~24 & \textit{for} Lemma~3307 \textit{read} Lemma~3303; the step uses Lemma~3303 (p.~3915) \\
\dag\,3922 & proof of Lemma~3321, l.~10 & \textit{for} $wRv' \wedge v'Rw_0 \leftrightarrow wRw_0$ \textit{read} $wRv' \wedge v'Rw_0 \to wRw_0$ \\
\dag\,3922 & proof of Lemma~3316, l.~12 & \textit{for} Lemma~3318 \textit{read} Lemma~3315; the step uses Lemma~3315 (p.~3920) \\
\dag\,3922 & proof of Lemma~3319, l.~27 & \textit{for} Lemma~3313 \textit{read} Lemma~3312 \\
3923 & l.~21 & \textit{for} Oppneheimer and Zalta (2011) \textit{read} Oppenheimer and Zalta (2011) \\
\dag\,3924 & proof of Lemma~3324, l.~10 & \textit{for} Lemma~3303 \textit{read} Lemma~3302 \\
\dag\,3924 & proof of Lemma~3326, (8.6) & \textit{for} $\pos(q \wedge p) \leftrightarrow (\nec q \wedge \nec p)$ \textit{read} $\nec(q \wedge p) \leftrightarrow (\nec q \wedge \nec p)$; correctly printed in (8.7), p.~3925 \\
\dag\,3924 & proof of Lemma~3326, l.~20 & \textit{for} Lemma~3322 \textit{read} Lemma~3317; the step uses Lemma~3317 (p.~3921) \\
\dag\,3924 & proof of Lemma~3324, l.~27 & \textit{for} Lemma~3305 \textit{read} Lemma~3300; the step uses Lemma~3300 (p.~3914) \\
3925 & l.~1 & \textit{for} $\nec(q \to r) \leftrightarrow (\nec q \to \nec r)$ \textit{read} $\nec(q \to r) \to (\nec q \to \nec r)$ \\
3925 & l.~5 & \textit{for} p.~3790 \textit{read} p.~3890 (Lemma~3147) \\
\dag\,3925 & proof of Lemma~3328, l.~15 & \textit{for} Lemma~3301 \textit{read} Lemma~3302 \\
\dag\,3925 & proof of Lemma~3328, l.~26 & \textit{for} Lemma~3308 \textit{read} Lemma~3305 \\
\dag\,3925 & proof of Lemma~3328, (8.8) & \textit{for} $w_2Rw \wedge wRw_1 \leftrightarrow w_2Rw_1$ \textit{read} $w_2Rw \wedge wRw_1 \to w_2Rw_1$; the formula is used in its correct form on p.~3928 \\
\dag\,3926 & proof of Lemma~3327, l.~15 & \textit{for} Lemma~3325 \textit{read} Lemma~3324 \\
\dag\,3927 & proof of Lemma~3332, l.~6 & \textit{for} Lemma~3314 \textit{read} Lemma~3313 \\
\dag\,3927 & proof of Lemma~3329, l.~12 & \textit{for} Lemma~3308 \textit{read} Lemma~3313 \\
3928 & l.~3 & \textit{for} Benzmüller and Woltzenlogel Paleo (2013) \textit{read} Benzmüller and Woltzenlogel Paleo (2014) \\
\dag\,3928 & proof of Lemma~3336, l.~29 & \textit{for} Lemma~3312 \textit{read} Lemma~3308 \\
\dag\,3928 & proof of Lemma~3336, l.~34 & \textit{for} $\nec r \wedge \pos q \leftrightarrow \pos(r \wedge q)$ \textit{read} $\nec r \wedge \pos q \to \pos(r \wedge q)$ \\
3929 & l.~12 & \textit{for} p.~3835 \textit{read} p.~3836 (Lemma~2671) \\
\dag\,3929 & proof of Lemma~3336, l.~18 & \textit{for} Lemma~3329 \textit{read} Lemma~3328 \\
\dag\,3929 & proof of Lemma~3334, l.~19 & \textit{for} Lemma~3323 \textit{read} Lemma~3318; the step uses Lemma~3318 (p.~3921) \\
3929 & l.~29 & \textit{for} Benzmüller and Woltzenlogel Paleo (2014) \textit{read} Oppenheimer and Zalta (2011) (the computational simplification of Anselm's argument) \\
3929 & l.~31 & \textit{for} Oppenheimer and Zalta (2010) \textit{read} Oppenheimer and Zalta (2011) \\
\dag\,3930 & proof of Lemma~3341, l.~13 & \textit{for} Lemma~3317 \textit{read} Lemma~3312 \\
\dag\,3930 & proof of Lemma~3341, l.~27 & \textit{for} Lemma~3324 \textit{read} Lemma~3327 \\
\dag\,3931 & proof of Lemma~3343, l.~24 & \textit{for} Lemma~3334 \textit{read} Lemma~3339; the step uses Lemma~3339 (p.~3930) \\
\dag\,3932 & proof of Lemma~3344, l.~2 & \textit{for} Lemma~3308 \textit{read} Lemma~3303 \\
3932 & l.~11 & \textit{for} $w_1 \models \pos p \iff \forall u\,(w_1Ru \to u \models p)$ \textit{read} $w_1 \models \nec p \iff \forall u\,(w_1Ru \to u \models p)$ \\
\dag\,3932 & proof of Lemma~3345, l.~24 & \textit{for} Lemma~3314 \textit{read} Lemma~3311 \\
\dag\,3933 & proof of Lemma~3349, (8.11) & \textit{for} $v'Rv \to v'Rv$ \textit{read} $vRv' \to v'Rv$; the formula is used in its correct form on p.~3935 \\
\dag\,3933 & proof of Lemma~3349, (8.13) & \textit{for} $\nec(p \wedge q) \leftrightarrow (\pos p \wedge \nec q)$ \textit{read} $\nec(p \wedge q) \leftrightarrow (\nec p \wedge \nec q)$ \\
3933 & l.~23 & \textit{for} p.~3625 \textit{read} p.~3725 (Lemma~1773) \\
\dag\,3933 & proof of Lemma~3349, (8.10) & \textit{for} $(w_0Rw_1 \wedge w_0Rv') \leftrightarrow w_1Rv'$ \textit{read} $(w_0Rw_1 \wedge w_0Rv') \to w_1Rv'$ \\
\dag\,3933 & proof of Lemma~3349, (8.12) & \textit{for} $u \models \pos p \iff \forall w_1\,(uRw_1 \to w_1 \models p)$ \textit{read} $u \models \nec p \iff \forall w_1\,(uRw_1 \to w_1 \models p)$ \\
3934 & l.~8 & \textit{for} $(wRw_0 \wedge wRw_2) \to w_0Rw_1$ \textit{read} $(wRw_0 \wedge wRw_2) \to w_0Rw_2$ \\
\dag\,3934 & proof of Lemma~3351, l.~9 & \textit{for} Lemma~3299 \textit{read} Lemma~3302 \\
3934 & l.~22 & \textit{for} Oppenheimr and Zalta (2011) \textit{read} Oppenheimer and Zalta (2011) \\
\dag\,3934 & proof of Lemma~3351, l.~25 & \textit{for} Lemma~3324 \textit{read} Lemma~3321 \\
3935 & l.~6 & \textit{for} p.~3795 \textit{read} p.~3796 (Lemma~2354) \\
\dag\,3935 & proof of Lemma~3354, l.~14 & \textit{for} $w_1 \models \pos p \iff \exists w_0\,(w_1Rw_0 \vee w_0 \models p)$ \textit{read} $w_1 \models \pos p \iff \exists w_0\,(w_1Rw_0 \wedge w_0 \models p)$ \\
\dag\,3935 & proof of Lemma~3353, l.~16 & \textit{for} Lemma~3308 \textit{read} Lemma~3311 \\
3935 & l.~27 & \textit{for} p.~3610 \textit{read} p.~3620 (Lemma~734) \\
\dag\,3935 & proof of Lemma~3354, (8.14) & \textit{for} $\pos\nec q \leftrightarrow \nec q$ \textit{read} $\pos\nec q \to \nec q$; the formula is used in its correct form on p.~3936 \\
3936 & l.~4 & \textit{for} Benzmüller and Woltzenlogel Paleo (2015) \textit{read} Benzmüller and Woltzenlogel Paleo (2014) \\
\dag\,3936 & proof of Lemma~3356, l.~7 & \textit{for} Lemma~3321 \textit{read} Lemma~3317; the step uses Lemma~3317 (p.~3921) \\
\dag\,3936 & proof of Lemma~3354, l.~26 & \textit{for} Lemma~3320 \textit{read} Lemma~3322 \\
\dag\,3938 & proof of Lemma~3360, l.~21 & \textit{for} Lemma~3300 \textit{read} Lemma~3299 \\
\end{errlist}


\section{Volume XI: Bibliography and Indices}\label{s:vol11}
% generated by generate.py -- do not edit
\subsection*{Front matter \hfill \normalfont\small pp.~3940--3942}
\begin{errlist}
3941 & Contents, Chapter~4 & \textit{for} 4168 \textit{read} 4158 \\
\end{errlist}

\subsection*{Bibliography \hfill \normalfont\small pp.~3943--4062}
\begin{errlist}
3944 & Adams, R.~M. (1972) & \textit{for} 81(2) \textit{read} 81(3) \\
3945 & Alston, W.~P. (1989) & \textit{for} Clarendon Press \textit{read} Cornell University Press \\
3946 & Anderson, C.~A. (1990) & \textit{for} 291--330 \textit{read} 291--303 \\
3951 & Barcan, R.~C. (1946) & \textit{for} fucntional \textit{read} functional \\
3952 & Barrett, J.~L. (2004) & \textit{for} 2003 \textit{read} 2004 \\
3954 & Benzmüller, C. and Woltzenlogel Paleo, B. (2014) & \textit{for} 2015 \textit{read} 2014 \\
3955 & Bergmann, M. (2001) & \textit{for} T. \textit{read} M. \\
3957 & Blackburn, P. et~al. (2001) & \textit{for} Harvard University Press \textit{read} Cambridge University Press \\
3958 & Chandler, H.~S. (1976) & \textit{for} 35(2) \textit{read} 36(2) \\
3960 & Chellas, B.~F. (1980) & \textit{for} 1979 \textit{read} 1980 \\
3961 & Craig, W.~L. (1979) & \textit{for} Routledge \textit{read} Macmillan \\
3964 & Draper, P. (1989) & \textit{for} 32(3) \textit{read} 23(3) \\
3966 & Fine, K. (1994) & \textit{for} 9 \textit{read} 8 \\
3967 & Fitting, M. and Mendelsohn, R.~L. (1998) & \textit{for} 1997 \textit{read} 1998 \\
3968 & Fitting, M. (2002) & \textit{for} 2001 \textit{read} 2002 \\
3970 & Frankfurt, H.~G. (1964) & \textit{for} 1965 \textit{read} 1964 \\
3971 & Gale, R.~M. and Pruss, A.~R. (1999) & \textit{for} K.~M. \textit{read} R.~M. \\
3973 & Gallin, D. (1975) & \textit{for} P. \textit{read} D. \\
3976 & Gödel, K. (1995) & \textit{see} Corrigendum~S18 \\
3977 & Grim, P. (1983) & \textit{for} nelgected \textit{read} neglected \\
3979 & Grim, P. (1991) & \textit{for} 1992 \textit{read} 1991 \\
3980 & Hájek, P. (2002) & \textit{for} 2001 \textit{read} 2002 \\
3982 & Hartshorne, C. (1962) & \textit{for} S. \textit{read} C. \\
3983 & Hasker, W. (1989) & \textit{for} Knoweldge \textit{read} Knowledge \\
3985 & Henkin, L. (1950) & \textit{for} 15(3) \textit{read} 15(2) \\
3986 & Hick, J. (1966) & \textit{for} 1967 \textit{read} 1966 \\
3988 & Hoffman, J. and Rosenkrantz, G.~S. (2002) & \textit{for} Routledge \textit{read} Blackwell \\
3989 & Hughes, G.~E. and Cresswell, M.~J. (1996) & \textit{for} 1995 \textit{read} 1996 \\
3993 & Kenny, A. (1979) & \textit{for} C. \textit{read} A. \\
3995 & Koons, R.~C. (1997) & \textit{for} 34(1) \textit{read} 34(2) \\
3996 & Kraay, K.~J. (2010) & \textit{for} 2009 \textit{read} 2010 \\
3998 & Kretzmann, N. (1966) & \textit{for} W. \textit{read} N. \\
3999 & Kripke, S.~A. (1963a) & \textit{for} 1964 \textit{read} 1963 \\
4001 & Kripke, S.~A. (1963b) & \textit{for} N.~A. \textit{read} S.~A. \\
4002 & Kripke, S.~A. (1980) & \textit{for} L.~A. \textit{read} S.~A. \\
4004 & Leftow, B. (2012) & \textit{for} R. \textit{read} B. \\
4007 & Leslie, J. (1989) & \textit{for} Blackwell \textit{read} Routledge \\
4008 & Lewis, D. (1970) & \textit{for} Asnelm \textit{read} Anselm \\
4010 & Mackie, J.~L. (1955) & \textit{for} 210--212 \textit{read} 200--212 \\
4011 & Mackie, J.~L. (1982) & \textit{for} Tehism \textit{read} Theism \\
4013 & Malcolm, N. (1960) & \textit{for} B. \textit{read} N. \\
4014 & Mann, W.~E. (1982) & \textit{for} 18(3) \textit{read} 18(4) \\
4015 & Mavrodes, G.~I. (1963) & \textit{for} 1962 \textit{read} 1963 \\
4017 & Morris, T.~V. (1987) & \textit{for} Aneslmian \textit{read} Anselmian \\
4018 & Oppenheimer, P.~E. and Zalta, E.~N. (2011) & \textit{for} 2012 \textit{read} 2011 \\
4020 & Oppy, G. (1995) & \textit{for} 1996 \textit{read} 1995 \\
4021 & Oppy, G. (2006) & \textit{for} Agruing \textit{read} Arguing \\
4023 & Pike, N. (1965) & \textit{for} 1964 \textit{read} 1965 \\
4024 & Pike, N. (1969) & \textit{for} 7(3) \textit{read} 6(3) \\
4026 & Plantinga, A. (1967) & \textit{for} Oxford University Press \textit{read} Cornell University Press \\
4027 & Plantinga, A. (1974a) & \textit{for} E. \textit{read} A. \\
4029 & Plantinga, A. (1974b) & \textit{for} Necssity \textit{read} Necessity \\
4030 & Plantinga, A. (1980) & \textit{for} Natrue \textit{read} Nature \\
4032 & Plantinga, A. (2011) & \textit{for} Cambridge University Press \textit{read} Oxford University Press \\
4033 & Pruss, A.~R. (2006) & \textit{for} K.~R. \textit{read} A.~R. \\
4035 & Rowe, W.~L. (1975) & \textit{for} Clarendon Press \textit{read} Princeton University Press \\
4036 & Rowe, W.~L. (1979) & \textit{for} 16(3) \textit{read} 16(4) \\
4038 & Rowe, W.~L. (2004) & \textit{for} 2003 \textit{read} 2004 \\
4039 & Salmon, N. (1989) & \textit{for} W. \textit{read} N. \\
4040 & Savage, C.~W. (1967) & \textit{for} paardox \textit{read} paradox \\
4042 & Schellenberg, J.~L. (1993) & \textit{for} Dviine \textit{read} Divine \\
4043 & Sobel, J.~H. (1987) & \textit{for} ontloogical \textit{read} ontological \\
4045 & Sobel, J.~H. (2004) & \textit{for} Belliefs \textit{read} Beliefs \\
4048 & Street, S. (2006) & \textit{for} relaist \textit{read} realist \\
4049 & Stump, E. and Kretzmann, N. (1981) & \textit{for} 1982 \textit{read} 1981 \\
4051 & Swinburne, R. (1993) & \textit{for} M. \textit{read} R. \\
4052 & Swinburne, R. (2004) & \textit{for} Exstence \textit{read} Existence \\
4054 & van Inwagen, P. (1977) & \textit{for} T. \textit{read} P. \\
4055 & van Inwagen, P. (2006) & \textit{for} Prblem \textit{read} Problem \\
4057 & Wierenga, E.~R. (1989) & \textit{for} Harvard University Press \textit{read} Cornell University Press \\
4058 & Williamson, T. (2013) & \textit{for} L. \textit{read} T. \\
4060 & Wykstra, S.~J. (1984) & \textit{for} 16(1) \textit{read} 16(2) \\
4061 & Zagzebski, L.~T. (1991) & \textit{for} S.~T. \textit{read} L.~T. \\
\end{errlist}

\subsection*{Index of names \hfill \normalfont\small pp.~4063--4101}
\begin{errlist}
4064 & s.v.\ Adams, R.~M. & \textit{for} 3396 \textit{read} 3386 \\
4064 & s.v.\ Alston, W.~P. & \textit{add} 3492 \\
4065 & s.v.\ Anderson, C.~A. & \textit{for} 344 \textit{read} 345 \\
4065 & s.v.\ Anselm & \textit{for} 1145 \textit{read} 1154 \\
4066 & s.v.\ Aquinas, T. & \textit{for} 1368 \textit{read} 1468 \\
4066 & s.v.\ Barcan, R.~C. & \textit{for} 1956 \textit{read} 1946 \\
4067 & s.v.\ Barrett, J.~L. & \textit{add} 2821 \\
4067 & s.v.\ Benzmüller, C. & \textit{add} 3934 \\
4068 & s.v.\ Bergmann, M. & \textit{for} 1315 \textit{read} 3115 \\
4068 & s.v.\ Blackburn, P. & \textit{for} 3459 \textit{read} 3549 \\
4069 & s.v.\ Chandler, H.~S. & \textit{for} 206 \textit{read} 207 \\
4069 & s.v.\ Chellas, B.~F. & \textit{for} 3460 \textit{read} 3560 \\
4070 & s.v.\ Craig, W.~L. & \textit{for} 2569 \textit{read} 2669 \\
4070 & s.v.\ Cresswell, M.~J. & \textit{for} 3534 \textit{read} 3543 \\
4071 & s.v.\ de Rijke, M. & \textit{add} 3294 \\
4071 & s.v.\ Descartes, R. & \textit{for} 3 \textit{read} 4 \\
4072 & s.v.\ Draper, P. & \textit{for} 2313 \textit{read} 2312 \\
4072 & s.v.\ Fine, K. & \textit{for} 1462 \textit{read} 1642 \\
4073 & s.v.\ Fitting, M. & \textit{for} 326 \textit{read} 327 \\
4074 & s.v.\ Frankfurt, H.~G. & \textit{for} 1451 \textit{read} 1452 \\
4074 & s.v.\ Gale, R.~M. & \textit{for} 2680 \textit{read} 2690 \\
4075 & s.v.\ Gallin, D. & \textit{for} 3663 \textit{read} 3653 \\
4075 & s.v.\ Gaunilo & \textit{add} 1885 \\
4076 & s.v.\ Gödel, K. & \textit{for} 1733 \textit{read} 1723 \\
4076 & s.v.\ Grim, P. & \textit{for} 3838 \textit{read} 3839 \\
4077 & s.v.\ Hájek, P. & \textit{add} 258 \\
4077 & s.v.\ Hartshorne, C. & \textit{add} 1283 \\
4078 & s.v.\ Hasker, W. & \textit{add} 3467 \\
4078 & s.v.\ Henkin, L. & \textit{add} 3640 \\
4079 & s.v.\ Hick, J. & \textit{for} 2339 \textit{read} 2439 \\
4079 & s.v.\ Hoffman, J. & \textit{for} 4160 \textit{read} 1460 \\
4080 & s.v.\ Hughes, G.~E. & \textit{for} 1228 \textit{read} 1227 \\
4080 & s.v.\ Hume, D. & \textit{add} 2228 \\
4081 & s.v.\ Kant, I. & \textit{for} 15 \textit{read} 14 \\
4081 & s.v.\ Kenny, A. & \textit{add} 3212 \\
4082 & s.v.\ Koons, R.~C. & \textit{for} 3813 \textit{read} 3713 \\
4083 & s.v.\ Kraay, K.~J. & \textit{for} 2969 \textit{read} 2699 \\
4083 & s.v.\ Kretzmann, N. & \textit{for} 3109 \textit{read} 3108 \\
4084 & s.v.\ Kripke, S.~A. & \textit{for} 2003 \textit{read} 2002 \\
4084 & s.v.\ Leftow, B. & \textit{for} 3083 \textit{read} 3183 \\
4085 & s.v.\ Leibniz, G.~W. & \textit{for} 1267 \textit{read} 1276 \\
4085 & s.v.\ Leslie, J. & \textit{for} 2618 \textit{read} 2718 \\
4086 & s.v.\ Lewis, D. & \textit{for} 1916 \textit{read} 1906 \\
4086 & s.v.\ Mackie, J.~L. & \textit{for} 2619 \textit{read} 2609 \\
4087 & s.v.\ Malcolm, N. & \textit{for} 1975 \textit{read} 1795 \\
4087 & s.v.\ Mann, W.~E. & \textit{for} 3087 \textit{read} 3086 \\
4088 & s.v.\ Mavrodes, G.~I. & \textit{add} 1482 \\
4088 & s.v.\ Mendelsohn, R.~L. & \textit{for} 1916 \textit{read} 1926 \\
4089 & s.v.\ Morris, T.~V. & \textit{for} 1618 \textit{read} 1628 \\
4089 & s.v.\ Oppenheimer, P.~E. & \textit{for} 3922 \textit{read} 3923 \\
4090 & s.v.\ Oppy, G. & \textit{for} 1820 \textit{read} 1802 \\
4090 & s.v.\ Pike, N. & \textit{for} 1585 \textit{read} 1595 \\
4091 & s.v.\ Plantinga, A. & \textit{for} 51 \textit{read} 50 \\
4091 & s.v.\ Pruss, A.~R. & \textit{add} 2666 \\
4092 & s.v.\ Rosenkrantz, G.~S. & \textit{for} 1436 \textit{read} 1446 \\
4093 & s.v.\ Rowe, W.~L. & \textit{for} 2648 \textit{read} 2658 \\
4093 & s.v.\ Salmon, N. & \textit{add} 233 \\
4094 & s.v.\ Savage, C.~W. & \textit{add} 1455 \\
4094 & s.v.\ Schellenberg, J.~L. & \textit{for} 2491 \textit{read} 2492 \\
4095 & s.v.\ Sobel, J.~H. & \textit{for} 2047 \textit{read} 2048 \\
4095 & s.v.\ Spinoza, B. & \textit{for} 2966 \textit{read} 2965 \\
4096 & s.v.\ Street, S. & \textit{for} 2873 \textit{read} 2773 \\
4096 & s.v.\ Stump, E. & \textit{add} 3470 \\
4097 & s.v.\ Swinburne, R. & \textit{add} 2717 \\
4097 & s.v.\ van Inwagen, P. & \textit{add} 2333 \\
4098 & s.v.\ Venema, Y. & \textit{for} 3450 \textit{read} 3550 \\
4098 & s.v.\ Wierenga, E.~R. & \textit{for} 1366 \textit{read} 1466 \\
4099 & s.v.\ Williamson, T. & \textit{for} 781 \textit{read} 871 \\
4099 & s.v.\ Woltzenlogel Paleo, B. & \textit{for} 5 \textit{read} 4 \\
4100 & s.v.\ Wykstra, S.~J. & \textit{for} 3271 \textit{read} 2371 \\
4100 & s.v.\ Zagzebski, L.~T. & \textit{for} 3448 \textit{read} 3449 \\
4101 & s.v.\ Zalta, E.~N. & \textit{for} 3915 \textit{read} 3916 \\
\end{errlist}

\subsection*{Index of subjects \hfill \normalfont\small pp.~4102--4157}
\begin{errlist}
4103 & s.v.\ accessibility relation & \textit{for} 3031 \textit{read} 3013 \\
4103 & s.v.\ actualism & \textit{for} 2069 \textit{read} 2096 \\
4104 & s.v.\ ancestral premise & \textit{for} 1141 \textit{read} 1140 \\
4105 & s.v.\ aseity & \textit{for} 3027 \textit{read} 3127 \\
4106 & s.v.\ axiom 4 & \textit{for} 1939 \textit{read} 1949 \\
4107 & s.v.\ axiom 5 & \textit{for} 2424 \textit{read} 2244 \\
4108 & s.v.\ axiom B & \textit{add} 3541 \\
4108 & s.v.\ Barcan formula & \textit{for} 2102 \textit{read} 2101 \\
4109 & s.v.\ best possible world & \textit{add} 3490 \\
4110 & s.v.\ Cantorian argument & \textit{for} 1743 \textit{read} 1744 \\
4111 & s.v.\ cognitive science of religion & \textit{for} 2701 \textit{read} 2711 \\
4112 & s.v.\ compossibility & \textit{for} 3864 \textit{read} 3846 \\
4113 & s.v.\ comprehension & \textit{for} 3622 \textit{read} 3632 \\
4114 & s.v.\ conceivability & \textit{for} 1139 \textit{read} 1129 \\
4114 & s.v.\ conservation & \textit{for} 3474 \textit{read} 3475 \\
4115 & s.v.\ converse Barcan formula & \textit{for} 1961 \textit{read} 1916 \\
4116 & s.v.\ converse premise & \textit{add} 1680 \\
4117 & s.v.\ creation & \textit{add} 3398 \\
4118 & s.v.\ dependence & \textit{add} 3336 \\
4119 & s.v.\ divine agency & \textit{for} 3511 \textit{read} 3411 \\
4119 & s.v.\ divine hiddenness & \textit{for} 2326 \textit{read} 2236 \\
4120 & s.v.\ divine simplicity & \textit{for} 3073 \textit{read} 3072 \\
4121 & s.v.\ Euclidean relation & \textit{for} 1270 \textit{read} 1207 \\
4122 & s.v.\ evidence of absence & \textit{add} 889 \\
4123 & s.v.\ evolutionary debunking & \textit{for} 2767 \textit{read} 2768 \\
4124 & s.v.\ existence predicate & \textit{add} 1875 \\
4125 & s.v.\ fine-tuning & \textit{for} 2682 \textit{read} 2692 \\
4125 & s.v.\ finite intersection property & \textit{add} 1723 \\
4126 & s.v.\ foreknowledge & \textit{for} 3239 \textit{read} 3238 \\
4127 & s.v.\ frame correspondence & \textit{for} 3162 \textit{read} 3172 \\
4128 & s.v.\ free-will defence & \textit{for} 2403 \textit{read} 2303 \\
4129 & s.v.\ general models & \textit{for} 3641 \textit{read} 3640 \\
4130 & s.v.\ godlikeness & \textit{add} 2981 \\
4130 & s.v.\ great-making property & \textit{for} 3747 \textit{read} 3757 \\
4131 & s.v.\ Identification premise & \textit{for} 3430 \textit{read} 3330 \\
4132 & s.v.\ identity of indiscernibles & \textit{for} 2925 \textit{read} 2924 \\
4133 & s.v.\ Kripke semantics & \textit{for} 1264 \textit{read} 2164 \\
4134 & s.v.\ Lean modules & \textit{for} 3915 \textit{read} 3916 \\
4135 & s.v.\ maximal excellence & \textit{for} 1824 \textit{read} 1724 \\
4136 & s.v.\ maximal greatness & \textit{for} 4146 \textit{read} 1446 \\
4136 & s.v.\ moral perfection & \textit{for} 1356 \textit{read} 1456 \\
4137 & s.v.\ multiverse & \textit{for} 2785 \textit{read} 2786 \\
4138 & s.v.\ naturalism & \textit{for} 2622 \textit{read} 2623 \\
4139 & s.v.\ necessary existence & \textit{add} 1009 \\
4140 & s.v.\ necessitism & \textit{for} 1910 \textit{read} 1900 \\
4141 & s.v.\ necessity of identity & \textit{for} 3124 \textit{read} 3024 \\
4141 & s.v.\ noseeum inference & \textit{for} 3913 \textit{read} 3903 \\
4142 & s.v.\ omnipotence & \textit{for} 1404 \textit{read} 1504 \\
4143 & s.v.\ omniscience & \textit{for} 1452 \textit{read} 1451 \\
4144 & s.v.\ paradox of the stone & \textit{add} 1517 \\
4145 & s.v.\ parody objections & \textit{for} 1073 \textit{read} 1037 \\
4146 & s.v.\ positive property & \textit{for} 1263 \textit{read} 1273 \\
4147 & s.v.\ possibility premise & \textit{for} 1038 \textit{read} 1039 \\
4147 & s.v.\ print axioms & \textit{for} 3912 \textit{read} 3922 \\
4148 & s.v.\ problem of evil & \textit{for} 2432 \textit{read} 2532 \\
4149 & s.v.\ proper filter & \textit{for} 1761 \textit{read} 1716 \\
4150 & s.v.\ property abstraction & \textit{for} 3528 \textit{read} 3628 \\
4151 & s.v.\ providence & \textit{for} 3492 \textit{read} 3482 \\
4152 & s.v.\ S5 & \textit{add} 3264 \\
4152 & s.v.\ sceptical theism & \textit{for} 2455 \textit{read} 2454 \\
4153 & s.v.\ signalling game & \textit{for} 3881 \textit{read} 3891 \\
4154 & s.v.\ symmetry objection & \textit{for} 1357 \textit{read} 1347 \\
4155 & s.v.\ uniqueness & \textit{for} 3057 \textit{read} 3047 \\
4156 & s.v.\ variable domains & \textit{add} 1665 \\
4157 & s.v.\ well-foundedness & \textit{for} 3691 \textit{read} 3692 \\
\end{errlist}

\subsection*{Index of numbered results \hfill \normalfont\small pp.~4158--4180}
\begin{errlist}
4159 & Lemma~3.3 (Vol.~III) & \textit{for} 687 \textit{read} 867 \\
4159 & Proposition~4.9 (Vol.~III) & \textit{for} 291 \textit{read} 921 \\
4159 & Proposition~4.16 (Vol.~III) & \textit{for} 942 \textit{read} 941 \\
4159 & Proposition~5.3 (Vol.~III) & \textit{for} 956 \textit{read} 965 \\
4159 & Theorem~9.7 (Vol.~III) & \textit{for} 2119 \textit{read} 1219 \\
4160 & Definition~5.10 (Vol.~IV) & \textit{for} 1658 \textit{read} 1568 \\
4160 & Proposition~3.4 (Vol.~V) & \textit{for} 1738 \textit{read} 1838 \\
4160 & Theorem~5.10 (Vol.~V) & \textit{for} 1947 \textit{read} 1946 \\
4160 & Lemma~7.4 (Vol.~V) & \textit{for} 2033 \textit{read} 2034 \\
4161 & Definition~3.8 (Vol.~VI) & \textit{for} 2258 \textit{read} 2285 \\
4161 & Proposition~5.14 (Vol.~VI) & \textit{for} 2043 \textit{read} 2403 \\
4161 & Lemma~6.7 (Vol.~VI) & \textit{for} 2432 \textit{read} 2433 \\
4161 & Lemma~3.14 (Vol.~VII) & \textit{for} 2588 \textit{read} 2688 \\
4162 & Lemma~7.15 (Vol.~VII) & \textit{for} 2808 \textit{read} 2880 \\
4162 & Corollary~7.16 (Vol.~VII) & \textit{for} 2885 \textit{read} 2886 \\
4162 & Corollary~8.2 (Vol.~VII) & \textit{for} 2803 \textit{read} 2903 \\
4162 & Definition~3.5 (Vol.~VIII) & \textit{for} 2880 \textit{read} 2980 \\
4162 & Proposition~6.6 (Vol.~VIII) & \textit{for} 3181 \textit{read} 3118 \\
4162 & Definition~2.8 (Vol.~IX) & \textit{for} 2341 \textit{read} 3241 \\
4163 & Lemma~2 (Vol.~X) & \textit{for} 3452 \textit{read} 3542 \\
4163 & Lemma~57 (Vol.~X) & \textit{for} 3457 \textit{read} 3547 \\
4163 & Lemma~76 (Vol.~X) & \textit{for} 3649 \textit{read} 3549 \\
4163 & Lemma~122 (Vol.~X) & \textit{for} 3454 \textit{read} 3554 \\
4164 & Lemma~158 (Vol.~X) & \textit{for} 3547 \textit{read} 3557 \\
4164 & Lemma~198 (Vol.~X) & \textit{for} 3560 \textit{read} 3561 \\
4164 & Lemma~214 (Vol.~X) & \textit{for} 3463 \textit{read} 3563 \\
4164 & Lemma~225 (Vol.~X) & \textit{for} 3563 \textit{read} 3564 \\
4164 & Lemma~269 (Vol.~X) & \textit{for} 3596 \textit{read} 3569 \\
4164 & Lemma~278 (Vol.~X) & \textit{for} 3579 \textit{read} 3569 \\
4164 & Lemma~294 (Vol.~X) & \textit{for} 3471 \textit{read} 3571 \\
4165 & Lemma~355 (Vol.~X) & \textit{for} 3477 \textit{read} 3577 \\
4165 & Lemma~379 (Vol.~X) & \textit{for} 3850 \textit{read} 3580 \\
4165 & Lemma~392 (Vol.~X) & \textit{for} 3483 \textit{read} 3583 \\
4165 & Lemma~424 (Vol.~X) & \textit{for} 3487 \textit{read} 3587 \\
4165 & Lemma~439 (Vol.~X) & \textit{for} 3578 \textit{read} 3588 \\
4165 & Lemma~440 (Vol.~X) & \textit{for} 3688 \textit{read} 3588 \\
4165 & Lemma~458 (Vol.~X) & \textit{for} 3490 \textit{read} 3590 \\
4165 & Lemma~476 (Vol.~X) & \textit{for} 3591 \textit{read} 3592 \\
4165 & Lemma~481 (Vol.~X) & \textit{for} 3583 \textit{read} 3593 \\
4166 & Lemma~512 (Vol.~X) & \textit{for} 3696 \textit{read} 3596 \\
4166 & Lemma~552 (Vol.~X) & \textit{for} 3601 \textit{read} 3600 \\
4166 & Lemma~576 (Vol.~X) & \textit{for} 3503 \textit{read} 3603 \\
4166 & Lemma~604 (Vol.~X) & \textit{for} 3660 \textit{read} 3606 \\
4166 & Lemma~606 (Vol.~X) & \textit{for} 3506 \textit{read} 3606 \\
4166 & Lemma~621 (Vol.~X) & \textit{for} 3068 \textit{read} 3608 \\
4167 & Lemma~709 (Vol.~X) & \textit{for} 3168 \textit{read} 3618 \\
4167 & Lemma~710 (Vol.~X) & \textit{for} 3718 \textit{read} 3618 \\
4167 & Lemma~746 (Vol.~X) & \textit{for} 3612 \textit{read} 3622 \\
4167 & Lemma~781 (Vol.~X) & \textit{for} 3725 \textit{read} 3625 \\
4167 & Lemma~790 (Vol.~X) & \textit{for} 3616 \textit{read} 3626 \\
4167 & Lemma~812 (Vol.~X) & \textit{for} 3630 \textit{read} 3631 \\
4167 & Lemma~823 (Vol.~X) & \textit{for} 3622 \textit{read} 3632 \\
4167 & Lemma~867 (Vol.~X) & \textit{for} 3366 \textit{read} 3636 \\
4168 & Lemma~902 (Vol.~X) & \textit{for} 3740 \textit{read} 3640 \\
4168 & Lemma~940 (Vol.~X) & \textit{for} 3654 \textit{read} 3644 \\
4168 & Lemma~984 (Vol.~X) & \textit{for} 3468 \textit{read} 3648 \\
4168 & Lemma~991 (Vol.~X) & \textit{for} 3469 \textit{read} 3649 \\
4168 & Lemma~993 (Vol.~X) & \textit{for} 3648 \textit{read} 3649 \\
4168 & Lemma~995 (Vol.~X) & \textit{for} 3605 \textit{read} 3650 \\
4169 & Lemma~1086 (Vol.~X) & \textit{for} 3569 \textit{read} 3659 \\
4169 & Lemma~1097 (Vol.~X) & \textit{for} 3661 \textit{read} 3660 \\
4169 & Lemma~1106 (Vol.~X) & \textit{for} 3561 \textit{read} 3661 \\
4169 & Lemma~1190 (Vol.~X) & \textit{for} 3760 \textit{read} 3670 \\
4169 & Lemma~1238 (Vol.~X) & \textit{for} 3575 \textit{read} 3675 \\
4169 & Lemma~1246 (Vol.~X) & \textit{for} 3666 \textit{read} 3676 \\
4170 & Lemma~1322 (Vol.~X) & \textit{for} 3686 \textit{read} 3685 \\
4170 & Lemma~1352 (Vol.~X) & \textit{for} 3788 \textit{read} 3688 \\
4170 & Lemma~1401 (Vol.~X) & \textit{for} 3629 \textit{read} 3692 \\
4170 & Lemma~1432 (Vol.~X) & \textit{for} 3795 \textit{read} 3695 \\
4170 & Lemma~1460 (Vol.~X) & \textit{for} 3967 \textit{read} 3697 \\
4171 & Lemma~1461 (Vol.~X) & \textit{for} 3679 \textit{read} 3697 \\
4171 & Lemma~1463 (Vol.~X) & \textit{for} 3798 \textit{read} 3698 \\
4171 & Lemma~1481 (Vol.~X) & \textit{for} 3799 \textit{read} 3699 \\
4171 & Lemma~1492 (Vol.~X) & \textit{for} 3701 \textit{read} 3700 \\
4171 & Lemma~1564 (Vol.~X) & \textit{for} 3607 \textit{read} 3707 \\
4171 & Lemma~1626 (Vol.~X) & \textit{for} 3702 \textit{read} 3712 \\
4172 & Lemma~1654 (Vol.~X) & \textit{for} 3704 \textit{read} 3714 \\
4172 & Lemma~1660 (Vol.~X) & \textit{for} 3175 \textit{read} 3715 \\
4172 & Lemma~1762 (Vol.~X) & \textit{for} 3624 \textit{read} 3724 \\
4172 & Lemma~1779 (Vol.~X) & \textit{for} 3275 \textit{read} 3725 \\
4172 & Lemma~1784 (Vol.~X) & \textit{for} 3736 \textit{read} 3726 \\
4173 & Lemma~1868 (Vol.~X) & \textit{for} 3747 \textit{read} 3737 \\
4173 & Lemma~1892 (Vol.~X) & \textit{for} 3750 \textit{read} 3740 \\
4173 & Lemma~1893 (Vol.~X) & \textit{for} 3640 \textit{read} 3740 \\
4173 & Lemma~1899 (Vol.~X) & \textit{for} 3742 \textit{read} 3741 \\
4173 & Lemma~1915 (Vol.~X) & \textit{for} 3643 \textit{read} 3743 \\
4173 & Lemma~1998 (Vol.~X) & \textit{for} 3743 \textit{read} 3753 \\
4174 & Lemma~2041 (Vol.~X) & \textit{for} 3759 \textit{read} 3758 \\
4174 & Lemma~2113 (Vol.~X) & \textit{for} 3776 \textit{read} 3767 \\
4174 & Lemma~2152 (Vol.~X) & \textit{for} 3771 \textit{read} 3772 \\
4174 & Lemma~2214 (Vol.~X) & \textit{for} 3778 \textit{read} 3779 \\
4175 & Lemma~2231 (Vol.~X) & \textit{for} 3718 \textit{read} 3781 \\
4175 & Lemma~2237 (Vol.~X) & \textit{for} 3882 \textit{read} 3782 \\
4175 & Lemma~2252 (Vol.~X) & \textit{for} 3774 \textit{read} 3784 \\
4175 & Lemma~2275 (Vol.~X) & \textit{for} 3886 \textit{read} 3786 \\
4175 & Lemma~2286 (Vol.~X) & \textit{for} 3798 \textit{read} 3788 \\
4175 & Lemma~2301 (Vol.~X) & \textit{for} 3879 \textit{read} 3789 \\
4175 & Lemma~2317 (Vol.~X) & \textit{for} 3971 \textit{read} 3791 \\
4175 & Lemma~2380 (Vol.~X) & \textit{for} 3899 \textit{read} 3799 \\
4175 & Lemma~2402 (Vol.~X) & \textit{for} 3902 \textit{read} 3802 \\
4176 & Lemma~2419 (Vol.~X) & \textit{for} 3803 \textit{read} 3804 \\
4176 & Lemma~2438 (Vol.~X) & \textit{for} 3086 \textit{read} 3806 \\
4176 & Lemma~2530 (Vol.~X) & \textit{for} 3807 \textit{read} 3817 \\
4176 & Lemma~2571 (Vol.~X) & \textit{for} 3812 \textit{read} 3822 \\
4176 & Lemma~2575 (Vol.~X) & \textit{for} 3832 \textit{read} 3822 \\
4176 & Lemma~2591 (Vol.~X) & \textit{for} 3825 \textit{read} 3824 \\
4177 & Lemma~2616 (Vol.~X) & \textit{for} 3817 \textit{read} 3827 \\
4177 & Lemma~2637 (Vol.~X) & \textit{for} 3840 \textit{read} 3830 \\
4177 & Lemma~2654 (Vol.~X) & \textit{for} 3844 \textit{read} 3834 \\
4177 & Lemma~2676 (Vol.~X) & \textit{for} 3936 \textit{read} 3836 \\
4177 & Lemma~2691 (Vol.~X) & \textit{for} 3738 \textit{read} 3838 \\
4177 & Lemma~2728 (Vol.~X) & \textit{for} 3831 \textit{read} 3841 \\
4177 & Lemma~2747 (Vol.~X) & \textit{for} 3943 \textit{read} 3843 \\
4177 & Lemma~2755 (Vol.~X) & \textit{for} 3744 \textit{read} 3844 \\
4177 & Lemma~2761 (Vol.~X) & \textit{for} 3855 \textit{read} 3845 \\
4177 & Lemma~2764 (Vol.~X) & \textit{for} 3854 \textit{read} 3845 \\
4178 & Lemma~2797 (Vol.~X) & \textit{for} 3748 \textit{read} 3848 \\
4178 & Lemma~2809 (Vol.~X) & \textit{for} 3950 \textit{read} 3850 \\
4178 & Lemma~2818 (Vol.~X) & \textit{for} 3805 \textit{read} 3850 \\
4178 & Lemma~2823 (Vol.~X) & \textit{for} 3815 \textit{read} 3851 \\
4178 & Lemma~2864 (Vol.~X) & \textit{for} 3856 \textit{read} 3855 \\
4178 & Lemma~2869 (Vol.~X) & \textit{for} 3857 \textit{read} 3856 \\
4178 & Lemma~2881 (Vol.~X) & \textit{for} 3587 \textit{read} 3857 \\
4178 & Lemma~2961 (Vol.~X) & \textit{for} 3856 \textit{read} 3865 \\
4178 & Lemma~2974 (Vol.~X) & \textit{for} 3686 \textit{read} 3866 \\
4179 & Lemma~2999 (Vol.~X) & \textit{for} 3868 \textit{read} 3869 \\
4179 & Lemma~3013 (Vol.~X) & \textit{for} 3880 \textit{read} 3870 \\
4179 & Lemma~3047 (Vol.~X) & \textit{for} 3886 \textit{read} 3876 \\
4179 & Lemma~3089 (Vol.~X) & \textit{for} 3883 \textit{read} 3882 \\
4179 & Lemma~3124 (Vol.~X) & \textit{for} 3987 \textit{read} 3887 \\
4179 & Lemma~3139 (Vol.~X) & \textit{for} 3879 \textit{read} 3889 \\
4179 & Lemma~3151 (Vol.~X) & \textit{for} 3991 \textit{read} 3891 \\
4180 & Lemma~3184 (Vol.~X) & \textit{for} 3894 \textit{read} 3895 \\
4180 & Lemma~3251 (Vol.~X) & \textit{for} 3904 \textit{read} 3905 \\
4180 & Lemma~3289 (Vol.~X) & \textit{for} 3810 \textit{read} 3910 \\
4180 & Lemma~3303 (Vol.~X) & \textit{for} 3916 \textit{read} 3915 \\
4180 & Lemma~3318 (Vol.~X) & \textit{for} 3931 \textit{read} 3921 \\
4180 & Lemma~3327 (Vol.~X) & \textit{for} 3915 \textit{read} 3925 \\
\end{errlist}


\begin{thebibliography}{99}\small
\bibitem[Anderson(1990)]{Anderson1990} Anderson, C.~A. (1990). Some emendations of Gödel's ontological proof. \emph{Faith and Philosophy}, 7(3), 291--303.
\bibitem[Aquinas(1265--1274)]{Aquinas} Thomas Aquinas (1265--1274). \emph{Summa Theologiae}, Part I, question 3.
\bibitem[Chellas(1980)]{Chellas1980} Chellas, B.~F. (1980). \emph{Modal Logic: An Introduction}. Cambridge University Press.
\bibitem[Gödel(1995)]{Godel1995} Gödel, K. (1995). Ontological proof. In S.~Feferman et al. (Eds.), \emph{Collected Works, Volume III: Unpublished Essays and Lectures} (pp.~403--404). Oxford University Press.
\bibitem[Grim(1983)]{Grim1983} Grim, P. (1983). Some neglected problems of omniscience. \emph{American Philosophical Quarterly}, 20(3), 265--276.
\bibitem[Grim(1991)]{Grim1991} Grim, P. (1991). \emph{The Incomplete Universe: Totality, Knowledge, and Truth}. MIT Press.
\bibitem[Hughes and Cresswell(1996)]{HughesCresswell1996} Hughes, G.~E. and Cresswell, M.~J. (1996). \emph{A New Introduction to Modal Logic}. Routledge.
\bibitem[Kant(1781)]{Kant} Kant, I. (1781/1787). \emph{Kritik der reinen Vernunft}, A592--602/B620--630.
\bibitem[Mavrodes(1963)]{Mavrodes1963} Mavrodes, G.~I. (1963). Some puzzles concerning omnipotence. \emph{The Philosophical Review}, 72(2), 221--223.
\bibitem[Pike(1965)]{Pike1965} Pike, N. (1965). Divine omniscience and voluntary action. \emph{The Philosophical Review}, 74(1), 27--46.
\bibitem[Pike(1969)]{Pike1969} Pike, N. (1969). Omnipotence and God's ability to sin. \emph{American Philosophical Quarterly}, 6(3), 208--216.
\bibitem[Rowe(1979)]{Rowe1979} Rowe, W.~L. (1979). The problem of evil and some varieties of atheism. \emph{American Philosophical Quarterly}, 16(4), 335--341.
\bibitem[Savage(1967)]{Savage1967} Savage, C.~W. (1967). The paradox of the stone. \emph{The Philosophical Review}, 76(1), 74--79.
\bibitem[Wykstra(1984)]{Wykstra1984} Wykstra, S.~J. (1984). The Humean obstacle to evidential arguments from suffering: On avoiding the evils of \enquote{appearance}. \emph{International Journal for Philosophy of Religion}, 16(2), 73--93.
\end{thebibliography}

\end{document}
